% =====================================================================
%  THE DAEDALUS CODE MANUAL
%  Automated MSR-JD Feynman-diagram computation for stochastic (S)PDEs
%
%  This is the MASTER file.  Each chapter lives in sections/NN-slug.tex
%  and is pulled in with \input below.  To build: compile this file with
%  pdflatex (twice, for the table of contents and cross-references).
%  Everything here is standard CTAN / Overleaf — no local packages.
% =====================================================================
\documentclass[11pt,oneside]{report}

% ---- page geometry & fonts -----------------------------------------
\usepackage[a4paper,margin=1in]{geometry}
\usepackage[T1]{fontenc}
\usepackage[utf8]{inputenc}
\usepackage{lmodern}
\usepackage{microtype}

% ---- math -----------------------------------------------------------
\usepackage{amsmath,amssymb,amsthm,mathtools}
\usepackage{bm}

% ---- graphics, colour, tables --------------------------------------
\usepackage{graphicx}
\usepackage{xcolor}
\usepackage{booktabs}
\usepackage{longtable}
\usepackage{array}
\usepackage{caption}
\usepackage{enumitem}

% ---- code listings --------------------------------------------------
\usepackage{listings}

% ---- callout boxes (tcolorbox is standard on Overleaf) -------------
\usepackage[most]{tcolorbox}

% ---- diagrams -------------------------------------------------------
\usepackage{tikz}
\usetikzlibrary{arrows.meta,positioning,calc}

% ---- hyperlinks (load late) ----------------------------------------
\usepackage[colorlinks=true,linkcolor=NavyBlue,citecolor=NavyBlue,
            urlcolor=Maroon,linktoc=all]{hyperref}

% ---------------------------------------------------------------------
%  Colour palette
% ---------------------------------------------------------------------
\definecolor{NavyBlue}{RGB}{20,60,120}
\definecolor{Maroon}{RGB}{130,30,40}
\definecolor{codebg}{RGB}{247,247,249}
\definecolor{codeframe}{RGB}{210,210,216}
\definecolor{kw}{RGB}{0,90,160}
\definecolor{str}{RGB}{160,60,40}
\definecolor{cmt}{RGB}{100,120,100}
\definecolor{notebg}{RGB}{235,243,250}
\definecolor{noteframe}{RGB}{120,160,200}
\definecolor{warnbg}{RGB}{252,243,233}
\definecolor{warnframe}{RGB}{210,150,80}
\definecolor{defbg}{RGB}{238,246,238}
\definecolor{defframe}{RGB}{140,180,140}

% ---------------------------------------------------------------------
%  Listing styles : Python  and  console
% ---------------------------------------------------------------------
\lstdefinestyle{py}{
  language=Python,
  backgroundcolor=\color{codebg},
  rulecolor=\color{codeframe},
  basicstyle=\ttfamily\footnotesize,
  keywordstyle=\color{kw}\bfseries,
  stringstyle=\color{str},
  commentstyle=\color{cmt}\itshape,
  showstringspaces=false,
  breaklines=true,
  frame=single,
  framesep=5pt,
  numbers=none,
  columns=fullflexible,
  keepspaces=true,
  upquote=true,
}
\lstdefinestyle{console}{
  backgroundcolor=\color{codebg},
  rulecolor=\color{codeframe},
  basicstyle=\ttfamily\footnotesize,
  breaklines=true,
  frame=single,
  framesep=5pt,
  columns=fullflexible,
  upquote=true,
}
\lstset{style=py}

% ---------------------------------------------------------------------
%  Inline helpers
% ---------------------------------------------------------------------
\newcommand{\code}[1]{\texttt{\small #1}}                 % inline code
\newcommand{\file}[1]{\texttt{\small\color{NavyBlue}#1}}  % a file path
\newcommand{\term}[1]{\textbf{#1}}                        % first use of a term
\newcommand{\Daedalus}{\textsc{Daedalus}}
\newcommand{\msrjd}{MSR\nobreakdash-JD}

% ---------------------------------------------------------------------
%  Callout environments
% ---------------------------------------------------------------------
\newtcolorbox{note}[1][]{colback=notebg,colframe=noteframe,
  fonttitle=\bfseries,title={Note},sharp corners,boxrule=0.6pt,#1}
\newtcolorbox[auto counter,number within=chapter]{gotcha}[1][]{colback=warnbg,colframe=warnframe,
  fonttitle=\bfseries,title={Gotcha~\thetcbcounter},sharp corners,boxrule=0.6pt,#1}
\newtcolorbox{defn}[1][]{colback=defbg,colframe=defframe,
  fonttitle=\bfseries,title={Definition},sharp corners,boxrule=0.6pt,#1}

% ---------------------------------------------------------------------
%  Math shorthands used throughout the manual
% ---------------------------------------------------------------------
\newcommand{\dd}{\mathrm{d}}
\newcommand{\ii}{\mathrm{i}}
\newcommand{\ee}{\mathrm{e}}
\newcommand{\Dt}{\partial_t}
\newcommand{\Lap}{\nabla^2}
\newcommand{\avg}[1]{\left\langle #1 \right\rangle}
\newcommand{\Gret}{G_{\mathrm R}}
\newcommand{\Sym}{\mathcal{S}}            % symmetry factor S(Gamma)
\newcommand{\Aut}{\mathrm{Aut}}
\newcommand{\Usym}{\mathcal{U}}           % first Symanzik polynomial
\newcommand{\Fsym}{\mathcal{F}}           % second Symanzik polynomial

\theoremstyle{definition}
\newtheorem{example}{Example}[chapter]

% ---------------------------------------------------------------------
%  Title
% ---------------------------------------------------------------------
\title{\Huge\bfseries The \Daedalus{} Code Manual\\[0.4em]
  \LARGE Automated \msrjd{} Feynman-Diagram Computation\\
  \large for Stochastic (S)PDEs}
\author{Daedalus framework --- Ocker Lab}
\date{Generated \today}

% =====================================================================
\begin{document}
\maketitle

\begin{abstract}
\noindent
\Daedalus{} is an automated, \emph{model-independent} framework that takes a
generic semilinear Langevin stochastic (partial) differential equation ---
equivalently a Martin--Siggia--Rose--Janssen--De~Dominicis (\msrjd) response-field
action --- and computes its connected correlation functions (cumulants) to a
chosen loop order, fully symbolically and with no per-model hand tuning.  This
manual documents \emph{every} stage of that pipeline, from the text you type to
describe a model, through the symbolic algebra, diagram enumeration, and loop
integration, to the numbers and plots that come out the other end.

\medskip\noindent
The manual is written for a reader who understands the \emph{physics}
(stochastic field theory, Feynman diagrams, the loop expansion) but who may have
\emph{little} experience with the specific software tooling the pipeline is built
on --- SageMath, the \code{nauty} graph generator, \code{sympy}, \code{numba}, and
so on.  Wherever such a tool appears we explain what it is from the ground up.
When in doubt we over-explain.
\end{abstract}

\tableofcontents

% ---- How to read this manual ---------------------------------------
% ---------------------------------------------------------------------
%  How to read this manual  (unnumbered front-matter chapter)
% ---------------------------------------------------------------------
\chapter*{How to Read This Manual}
\addcontentsline{toc}{chapter}{How to Read This Manual}
\markboth{How to Read This Manual}{How to Read This Manual}

\section*{Who this is for}

This manual documents the internals of \Daedalus{}, the automated
\msrjd{} Feynman-diagram engine.  It is written for someone who
understands the \emph{physics} --- stochastic differential equations, the
response-field (\msrjd) formalism, Feynman diagrams, and the loop expansion ---
but who has \emph{not} necessarily worked with the specific software libraries
the engine is built on.  The two audiences we keep in mind are:

\begin{itemize}[nosep]
  \item a new lab member who wants to \emph{use} the pipeline to compute
        cumulants for their own model, and
  \item a developer who needs to \emph{modify or extend} a specific stage and
        must understand exactly what the code does and why.
\end{itemize}

Wherever a piece of external tooling appears --- SageMath, the \code{nauty}
graph generator, \code{sympy}, \code{numba}, \code{scipy} --- we explain what it
is and how the engine uses it, from the ground up.  When we must choose between
saying too much and too little, we say too much.

\section*{How the manual is organised}

The book is divided into six parts that follow the data as it flows through the
pipeline, plus appendices:

\begin{description}[style=nextline,leftmargin=1.2em]
  \item[Part I --- Orientation] A bird's-eye view of the whole pipeline and a
        worked ``hello world'': from a text description of a model to a plotted
        cumulant in a dozen lines.
  \item[Part II --- Specifying a Model] How you describe a stochastic (S)PDE to
        the engine: the \code{TheoryBuilder} authoring API, the on-disk
        \file{.theory.py} file format, and the automatic rewriting of
        colored noise into an auxiliary Markovian field.
  \item[Part III --- Building the Perturbation Theory] The symbolic core:
        expanding the action into vertices, constructing the propagator,
        solving for the mean-field saddle, enumerating diagrams, assigning
        edge types, computing symmetry factors, and caching all of it.
  \item[Part IV --- Evaluating Diagrams: Temporal Models] The time-domain loop
        integrator (``Phase~J''), both the per-diagram and the grouped analytic
        mode-sum paths.
  \item[Part V --- Evaluating Diagrams: Spatial Models] The momentum-space
        integrator for stochastic PDEs: Symanzik polynomials, causal chambers,
        the heat-kernel inverse Fourier transform, and coupled multi-field
        dressing.
  \item[Part VI --- Assembly and the User Surface] How the per-diagram numbers
        are assembled into cumulants and moments, the user-facing
        \code{daedalus} API, the validation simulators, and the graphical
        theory-builder UI.
  \item[Appendices] A glossary, a consolidated primer on the external tools, a
        file-by-file index of the source tree, and a register of known issues
        and open questions.
\end{description}

You do not have to read linearly.  If you only want to \emph{use} the pipeline,
Part~I and the \code{daedalus} API chapter are enough.  If you are debugging a
specific stage, jump straight to its chapter; each is written to stand on its
own, with back-references where context is needed.

\section*{Typographic conventions}

\begin{itemize}[nosep]
  \item \code{Inline monospace} denotes code: a function, variable, keyword
        argument, or short snippet.
  \item \file{src/path/to/file.py} (navy monospace) denotes a path in the
        repository, always relative to the repository root.
  \item The \emph{first} appearance of a technical \term{term} is set in bold
        and defined nearby; every such term also appears in the glossary
        (Appendix~A).
  \item Multi-line code is shown in framed listings.  Python listings use
        syntax highlighting; console sessions and output are shown verbatim.
\end{itemize}

Three kinds of callout box punctuate the text:

\begin{note}
A clarification or a useful aside --- safe to skim, helpful to read.
\end{note}

\begin{gotcha}
A sharp edge: a non-obvious requirement, a platform-specific hazard, or a place
where the obvious thing is wrong.  Read these.
\end{gotcha}

\begin{defn}
A precise definition of a term or object the surrounding text depends on.
\end{defn}

\section*{Running the examples}

Every runnable example assumes the engine is imported through its convenience
layer:

\begin{lstlisting}
import daedalus as dd                 # notebooks/daedalus.py
model, mod = dd.load_theory('ou_quartic')
cfg = dd.Config(k=2, max_ell=2,
                external_fields=[('dx', 1), ('dx', 1)],
                parameters={'mu': 1.0, 'eps': 0.02, 'D': 1.0})
res = dd.run(model, cfg, mod)
\end{lstlisting}

\begin{gotcha}
The engine runs inside \textbf{SageMath}, not a plain CPython interpreter.  All
scripts must be launched with \code{sage -python yourscript.py}, and notebooks
must use the \emph{SageMath} Jupyter kernel.  Sage provides the exact symbolic
algebra and the multivariate polynomial rings the engine relies on; a stock
\code{python3} will not have them.  See Appendix~B for what Sage is and why it is
required.
\end{gotcha}

\section*{Getting the engine running}

Everything in this manual assumes a working \textbf{SageMath} with the project's
dependencies importable. If you are starting from nothing:
\begin{enumerate}
  \item \textbf{Install SageMath} (10.x). On macOS the binary app or a conda
        install both work; the engine only needs \code{import sage.all} to
        succeed.
  \item \textbf{Make the repository importable.} The example notebooks do this
        automatically by walking up to the directory that contains the
        \file{pipeline/} package and adding it (and \file{notebooks/}) to
        \code{sys.path}.
  \item \textbf{Use the project's environment and launcher.} The pipeline is
        developed in a dedicated conda environment (named \code{MSRJD\_diagrams}
        in this project) layered on Sage, and \emph{every} script is run with
        \code{sage -python}, never a bare \code{python3} (Appendix~B explains
        why).
  \item \textbf{Smoke-test} in under a second from the repository root:
\begin{lstlisting}[style=console]
$ sage -python -c "import sys; sys.path.insert(0,'notebooks'); import daedalus as dd; print(dd.list_theories()[:3])"
\end{lstlisting}
        If that prints a few theory names, the environment is green and you can
        go straight to the Quickstart (Chapter~\ref{ch:quickstart}).
\end{enumerate}

\begin{gotcha}
A stock miniforge/conda \emph{base} environment can ship a \textbf{broken NumPy}
(a \code{libgfortran} rpath problem) that makes \code{import numpy} fail under
Sage. If you hit an opaque NumPy import error, you are almost certainly in the
wrong environment --- activate the project's \code{MSRJD\_diagrams} env and try
again.
\end{gotcha}

\section*{A note on accuracy}

This manual was assembled by reading the source directly, and each chapter was
cross-checked against the code it describes.  Discrepancies discovered during
that process --- places where the code and its documentation or intent did not
line up --- are not silently papered over; they are recorded in
Appendix~D (Known Issues and Open Questions) for follow-up.  If you find the code
doing something this manual does not describe, Appendix~D is the first place to
look.


% =====================================================================
%  PART I  --  ORIENTATION
% =====================================================================
\part{Orientation}
% ---------------------------------------------------------------------
%  Chapter 1 : Overview and Architecture
% ---------------------------------------------------------------------
\chapter{Overview and Architecture}
\label{ch:overview}

\section{What the engine computes}

You hand \Daedalus{} a stochastic differential equation --- an ordinary one in
time, or a partial one in space and time --- of the general semilinear
(Langevin) form
\begin{equation}
  \Dt \phi_a(\mathbf x,t)
   \;=\; -\,\mu_a\,\phi_a + D_a \Lap \phi_a
          \;+\; V_a\big(\{\phi\}\big) \;+\; \xi_a(\mathbf x,t),
  \label{eq:langevin}
\end{equation}
where $\phi_a$ are the fields, $V_a$ collects the (poly)nonlinear interactions,
and $\xi_a$ is a Gaussian noise with a prescribed covariance (white, or
exponentially correlated in time).  In return the engine computes the connected
\term{cumulants} of the fields,
\begin{equation}
  \kappa_k(t_1,\dots,t_k)
    \;=\; \big\langle \phi(t_1)\cdots\phi(t_k)\big\rangle_{\mathrm c},
\end{equation}
to a chosen order in the loop expansion, as exact symbolic-then-numeric
expressions --- with \emph{no} model-specific code.  The same machinery handles
an Ornstein--Uhlenbeck particle, a multi-population nonlinear Hawkes network, a
reaction--diffusion field, and the KPZ equation, because each is just a different
instance of \eqref{eq:langevin}.

\begin{defn}
\term{Model-independence} is the central design principle: every algorithm in the
pipeline reads the structure of the action symbolically and does the right thing,
rather than special-casing a particular equation.  Adding a new model means
writing a short text description (Chapter on theory files); it never means
touching the integrator.
\end{defn}

\section{The MSR--JD formalism in one page}

The engine never works with the SDE \eqref{eq:langevin} directly.  It works with
the equivalent field theory obtained by the
Martin--Siggia--Rose--Janssen--De~Dominicis (\msrjd) construction, which trades a
stochastic equation for a deterministic path integral over \emph{two} fields per
degree of freedom.

For each physical field $\phi_a$ one introduces a \term{response field}
$\tilde\phi_a$ (the engine writes it \code{xt} for a physical field \code{x}).
Averaging the SDE's noise produces the action
\begin{equation}
  S[\phi,\tilde\phi]
   \;=\; \int \dd t \; \Big[\,
      \tilde\phi_a\big(\Dt + \mu_a - D_a\Lap\big)\phi_a
      \;-\; \tilde\phi_a\,V_a(\{\phi\})
      \;-\; \tfrac12 \tilde\phi_a\,\Sigma_{ab}\,\tilde\phi_b
   \Big],
  \label{eq:msrjd}
\end{equation}
where $\Sigma_{ab}$ is the noise covariance (for white noise of strength $D$ this
is the local $-D\,\tilde\phi^2$ term).  Correlation and response functions of the
original SDE become correlation functions in this field theory:
\begin{equation}
  \avg{\phi_a(t)\,\phi_b(t')} \quad\text{and}\quad
  \avg{\phi_a(t)\,\tilde\phi_b(t')} = \text{(response/propagator)}.
\end{equation}
From \eqref{eq:msrjd} one reads off the two ingredients every diagrammatic
calculation needs:

\begin{itemize}
  \item the \term{bare propagators}, from the quadratic part --- a
        \emph{response} (retarded) propagator $\Gret$ that is causal
        ($\propto \theta(t-t')$), and a \emph{correlation} propagator from the
        noise vertex $\Sigma$;
  \item the \term{interaction vertices}, from the nonlinear terms
        $\tilde\phi_a V_a$ --- each monomial in the fields is a vertex with a
        definite number of legs and a prefactor.
\end{itemize}

Connected cumulants are then sums of connected Feynman diagrams built from these
propagators and vertices, organised by the number of independent loops $\ell$.
Tree diagrams ($\ell=0$) give the mean-field/linear answer; each extra loop is a
higher-order fluctuation correction.  \emph{Everything} the engine does is in
service of (i) enumerating those diagrams and (ii) evaluating their integrals.

\begin{note}
Because the response propagator is causal, the loop integrals are not the
relativistic Feynman integrals of high-energy physics; they are
\emph{time-ordered}, retarded integrals.  This is why the engine has its own
integrator (Parts IV and V) rather than reusing an off-the-shelf high-energy
loop-integral library.
\end{note}

\section{The pipeline at a glance}

A single call --- \code{compute\_cumulants} in \file{pipeline/compute.py},
wrapped for users by \code{dd.run} in \file{notebooks/daedalus.py} --- drives the
whole computation.  Internally it proceeds through seven numbered phases, which
you can watch scroll past when \code{verbose=True}:

\begin{enumerate}
  \item \textbf{expand} --- Taylor/functional-expand the symbolic action about the
        mean-field saddle and extract the interaction vertices
        (Chapter on the action and vertices).
  \item \textbf{propagator} --- build the frequency-domain propagator matrix
        $G(\omega) = K(\omega)^{-1}$ over the (physical, response) field basis
        (Chapter on the propagator).
  \item \textbf{mean\_field} --- solve the deterministic saddle-point equations
        for the steady state $\phi^{*}$ about which the expansion is taken
        (Chapter on the mean-field solver).
  \item \textbf{poles} --- find the retarded poles of $G(\omega)$ and the residue
        matrices that give the time-domain propagator (same chapter as
        propagator).
  \item \textbf{diagrams} --- enumerate the distinct diagram topologies for the
        requested $(k,\ell)$ and assign causal edge types, loading from cache
        when possible (Chapters on enumeration and typing).
  \item \textbf{classify} --- for each diagram compute its combinatorial prefactor
        and symmetry factor (Chapter on typing and symmetry).
  \item \textbf{phase\_j} --- evaluate every diagram's loop integral and sum them
        into the cumulant (Part IV for temporal models, Part V for spatial).
\end{enumerate}

\begin{figure}[t]
\centering
\begin{tikzpicture}[
   node distance=6.2mm,
   box/.style={draw,rounded corners,align=center,inner sep=4pt,
               minimum width=58mm,font=\small,fill=notebg},
   io/.style ={draw,align=center,inner sep=4pt,minimum width=58mm,
               font=\small,fill=defbg},
   arr/.style={-{Latex[length=2mm]},thick}]
  \node[io]  (txt)  {Model as text \;/\; \file{.theory.py}};
  \node[box,below=of txt] (build){\textbf{build}: parse action $\to$ symbolic model};
  \node[box,below=of build](mark) {\textbf{markovianize} (if colored noise)};
  \node[box,below=of mark] (exp)  {\textbf{expand}: vertices about the saddle};
  \node[box,below=of exp]  (prop) {\textbf{propagator} \& \textbf{mean\_field} \& \textbf{poles}};
  \node[box,below=of prop] (enum) {\textbf{diagrams}: enumerate \& type \& symmetry};
  \node[box,below=of enum] (intg) {\textbf{phase\_j}: loop integration\\(temporal \emph{or} spatial)};
  \node[box,below=of intg] (asm)  {\textbf{assemble}: cumulants / moments};
  \node[io, below=of asm]  (out)  {$\kappa_k(\tau)$, plots, saved results};
  \foreach \a/\b in {txt/build,build/mark,mark/exp,exp/prop,prop/enum,enum/intg,intg/asm,asm/out}
     \draw[arr] (\a) -- (\b);
\end{tikzpicture}
\caption{The end-to-end pipeline.  Green boxes are user-facing input/output;
blue boxes are internal phases.  The branch into a temporal or spatial integrator
happens at \textbf{phase\_j}; everything above it is shared.}
\label{fig:pipeline}
\end{figure}

Figure~\ref{fig:pipeline} shows the flow.  The crucial structural fact is that
\emph{everything above the integrator is shared between temporal and spatial
models}: the same action parser, the same vertex extraction, the same diagram
enumeration.  Only the final loop-integration phase splits into a time-domain
engine (Part~IV) and a momentum-space engine (Part~V).

\section{How the source tree is organised}

The code lives in three layers, from user-facing down to mathematical core:

\begin{description}[style=nextline,leftmargin=1.2em]
  \item[\file{notebooks/daedalus.py} --- the engine API] A single
        $\sim$1100-line convenience module that users import as \code{dd}.  It
        defines the \code{Config} dataclass, \code{load\_theory}, \code{run},
        \code{describe\_model}, and the plotting helpers.  It \emph{orchestrates};
        it contains no physics of its own (Chapter on the API).
  \item[\file{pipeline/} --- orchestration] The high-level stages:
        \file{compute.py} (the seven-phase driver), the theory-authoring API
        (\file{theory.py}, \file{theory\_compiler.py}), the propagator and
        mean-field builders (\file{\_propagator.py}, \file{\_mean\_field*.py}),
        the diagram-cache front end (\file{\_diagrams.py}), the colored-noise
        preprocessor (\file{colored\_to\_markovian.py}), and the caching
        machinery (\file{\_expand\_cache.py}, \file{\_precompute.py}).
  \item[\file{msrjd/} --- the mathematical engine] The lower-level core:
        \file{core/} (the symbolic action, vertices, convolutions),
        \file{enumeration/} and \file{diagrams/} (diagram generation, typing,
        causality, symmetry factors), and \file{integration/} which splits into
        \file{time\_domain/} (the temporal Phase-J integrator, dominated by the
        5000-line \file{final\_integral.py}) and \file{spatial/} (the
        fourteen-module momentum-space integrator).
\end{description}

Two further directories complete the picture: \file{theories/} holds the on-disk
\file{*.theory.py} model descriptions, and \file{models/} holds the
\emph{independent} simulators used to validate the diagrammatic output --- direct
Euler--Maruyama and Gillespie integrators that share no code with the pipeline
(Chapter on simulators).

\section{Caching: why a second run is instant}

Several phases are expensive but deterministic functions of a small key, so the
engine caches them on disk under \file{saved\_theories/<model>/}.  The most
important is the \term{diagram-enumeration cache}: enumerating and typing all
diagrams for a given $(k,\ell)$ is combinatorially heavy at two loops, but it
depends only on the model's topology, not on parameter values, so it is computed
once and reloaded from a Sage-object (\file{.sobj}) file thereafter.  A second
cache stores the expanded vertices.  Together they are why changing
$\mu$ and re-running is fast even though the first run of a new two-loop model is
not (Chapter on caching).

\section{Validation philosophy}

Every nontrivial result is checked against an independent Monte-Carlo
simulation.  The simulators in \file{models/} integrate the \emph{original} SDE
directly --- no diagrams, no propagators --- and estimate the same cumulant the
pipeline predicts.  Agreement between the two is the engine's ground truth; the
example notebooks (\file{notebooks/examples/}) each end with such an overlay.  As
a concrete instance, the white-noise quartic oscillator of
\eqref{eq:langevin} with $\mu=1,\varepsilon=0.02,D=1$ has an equal-time variance
the linear theory puts at $D/\mu = 1$; the two-loop pipeline corrects this to
$0.950$, and a direct simulation gives $0.946\pm0.001$ --- agreement to a fraction
of a percent.

\section{Where to go next}

The next chapter walks through a complete worked example end to end.  After that,
Part~II shows how to describe your own model, and Parts~III--V open up each stage
of the machinery in full detail.

% ---------------------------------------------------------------------
%  Chapter 2 : A Worked Example, End to End
% ---------------------------------------------------------------------
\chapter{A Worked Example, End to End}\label{ch:quickstart}

This chapter is your first hands-on contact with the engine. We will take a
single, small, genuinely nonlinear stochastic model --- a quartic
Ornstein--Uhlenbeck particle driven by white noise --- and carry it the whole
way from a one-line model name to a plotted correlation function with an
independent simulator drawn on top. Nothing is hidden: every line of code is
shown, every printed line of output is quoted, and the few numbers that matter
(the equal-time variance the theory predicts, and the one the simulation
measures) are compared side by side.

The point of starting here, before any of the machinery chapters, is to give
you a mental ``spine'' for the rest of the manual. By the end of this chapter
you will have seen the four verbs that drive everything --- \code{load\_theory},
\code{Config}, \code{run}, \code{plot\_cumulant} --- and you will have watched
the engine narrate its own seven internal phases. Later chapters open up each of
those phases in turn; this one just runs them.

The model we use is the one shipped in the example notebook
\file{notebooks/examples/temporal\_ou\_quartic\_white.ipynb}, with its theory
text in \file{theories/ou\_quartic.theory.py}. You can follow along by opening
that notebook, or by pasting the cells below into a fresh one.

\section{The physics in one equation}

The model is an Ornstein--Uhlenbeck (OU) variable with a cubic restoring force,
driven by white Gaussian noise:
\begin{equation}
  \dot x \;=\; -\,\mu\,x \;-\; \varepsilon\,x^{3} \;+\; \xi,
  \qquad
  \avg{\xi(t)\,\xi(t')} \;=\; 2D\,\delta(t-t').
  \label{eq:ouquartic}
\end{equation}
This is the simplest \emph{nonlinear} process the pipeline handles. Set
$\varepsilon = 0$ and it collapses to a plain linear OU process, whose
autocorrelation is the textbook exponential $C_{xx}(\tau) = (D/\mu)\,
\ee^{-\mu|\tau|}$ with equal-time variance $C_{xx}(0) = D/\mu$. The cubic term
$-\varepsilon x^{3}$ is what makes the problem interesting: it is a genuine
interaction vertex, it cannot be solved in closed form, and it pulls the
equal-time variance \emph{below} the linear value $D/\mu$. Capturing that shift
is exactly what the loop expansion is for.

We run in the mildly perturbative single-well regime
$\mu = 1,\ \varepsilon = 0.02,\ D = 1$. With $\mu > 0$ the mean-field saddle is
$x^{*} = 0$ (a single well; for $\mu < 0$ the same theory becomes a genuine
double well, which is a different example). The natural small parameter of the
loop expansion is
\[
  g_{\mathrm{eff}} \;\approx\; \frac{\varepsilon D}{\mu^{2}} \;=\; 0.02,
\]
so the tree, one-loop, and two-loop contributions converge quickly and the whole
calculation finishes in a fraction of a second --- which is precisely why this is
the model to learn on.

\begin{note}
``White'' noise here means $\delta$-correlated in time. In the \msrjd{} action
this is the local term $-D\,\tilde x^{2}$ (a single noise vertex). There is no
finite correlation time $\tau_c$ to embed, so there is no auxiliary noise field
and the diagram count stays small. The coloured-noise sibling of this model adds
a Markovian embedding field; comparing the two isolates exactly what colour
costs. That comparison is the subject of a later chapter.
\end{note}

\section{Setup: importing the front desk}

Everything you will touch in this chapter lives in one module,
\file{notebooks/daedalus.py}, which the notebooks import as \code{dd}. It is the
\emph{front desk} of the pipeline: it loads a theory, packages your choices into
a \code{Config}, calls the heavy engine, and plots the result. It does almost no
physics itself --- all the hard work happens behind a single
\code{from pipeline import compute\_cumulants}.

The setup cell does three small jobs: find the repository root (so the import
works no matter which subdirectory the notebook lives in), put it on the import
path, and \code{import daedalus as dd}.

\begin{lstlisting}
%matplotlib inline
import os, sys, time
import numpy as np
import matplotlib.pyplot as plt
# depth-robust repo root: walk up until the 'pipeline' package is found
_root = os.path.abspath('')
while _root != os.path.dirname(_root) and not os.path.isdir(os.path.join(_root, 'pipeline')):
    _root = os.path.dirname(_root)
sys.path.insert(0, _root)
sys.path.insert(0, os.path.join(_root, 'notebooks'))
os.chdir(os.path.join(_root, 'notebooks'))  # cwd=notebooks/ for models/ data paths
import daedalus as dd
\end{lstlisting}

The \code{while} loop is worth a sentence because the same idea appears inside
\code{daedalus} itself. It starts from the current directory and walks
\emph{up} the tree (\code{\_root = os.path.dirname(\_root)}) until it finds one
that contains a \code{pipeline/} folder; that directory is the repository root.
This is why the example notebooks run unchanged whether you launch them from the
repository top, from \code{notebooks/}, or from \code{notebooks/examples/}. The
module does the very same walk in its own \code{repo\_root} function
(\file{notebooks/daedalus.py:40}), so even if you forget the setup cell, the
import still resolves \code{pipeline}.

\begin{gotcha}
The line \code{os.chdir(...)} changes the working directory to
\code{notebooks/}. Some simulators load fixture data by relative path
(\code{models/...}), and they expect to be run from there. If you skip this and
later hit a ``file not found'' from a simulator, a stale working directory is the
usual cause.
\end{gotcha}

\section{Step 1 --- Load the theory}

A \term{theory file} is a small, declarative Python module under
\file{theories/} that describes a model: its fields, its parameters, its action,
and a few run recommendations. You never edit the engine to add a model; you
write one of these. Loading one is a single call:

\begin{lstlisting}
THEORY = 'ou_quartic'
model, mod = dd.load_theory(THEORY)
\end{lstlisting}

\code{load\_theory} returns a \emph{pair}, and the distinction matters
throughout the rest of the chapter:

\begin{itemize}
  \item \code{model} is a plain Python \code{dict} --- the fully built model
        description. It is the output of the theory's own \code{build()}
        function and it is what the engine actually consumes.
  \item \code{mod} is the live theory \emph{module} itself. The model dict does
        \emph{not} carry the run recommendations; those live on the module, in
        two module-level names: \code{DEFAULT\_FUNDAMENTAL} (default numeric
        parameter values) and \code{METADATA} (suggested \code{k}, loop order,
        external legs, grid extents).
\end{itemize}

The mechanism behind \code{load\_theory} is Python's
\code{importlib.util}, which can import a module \emph{from an arbitrary file
path} rather than by a name on \code{sys.path}. The implementation
(\file{notebooks/daedalus.py:77}) is the canonical three-step idiom:

\begin{lstlisting}
spec = importlib.util.spec_from_file_location(f'theories.{name}', path)
mod = importlib.util.module_from_spec(spec)
spec.loader.exec_module(mod)
return mod.build(), mod
\end{lstlisting}

The first line describes the module to be created, the second creates an empty
module object, and the third actually \emph{runs} the file's top-level code ---
which defines \code{build}, \code{DEFAULT\_FUNDAMENTAL}, and \code{METADATA}. The
function then returns \code{(mod.build(), mod)}: the built model dict and the
live module.

\begin{note}
To see every theory you can load, call \code{dd.list\_theories()}. It lists the
\file{theories/} directory, keeps the files ending in \code{.theory.py}, and
strips that suffix --- so \file{ou\_quartic.theory.py} becomes the loadable name
\code{'ou\_quartic'}. If you mistype a name, \code{load\_theory} raises a
\code{FileNotFoundError} whose message includes that same list, so the menu of
valid names is always one error away.
\end{note}

\subsection{What the theory file actually says}

Before going further, look at the theory text we just loaded
(\file{theories/ou\_quartic.theory.py}). It is short enough to read in full, and
reading it now makes every later number traceable:

\begin{lstlisting}
from pipeline.theory import TemporalTheoryBuilder

def build():
    return (
        TemporalTheoryBuilder('OU Quartic (white noise)')
        .population('pop', size=1)
        .physical_field('x', population='pop', description='variable')
        .parameter('mu', default=1.0, domain='positive')
        .parameter('eps', default=0.02, domain='positive')
        .parameter('D', default=1.0, domain='positive')
        .set_action_text(
            'sum(xt[i]*((Dt+mu)*x[i] + eps*x[i]^3) - D*xt[i]^2 for i in pop)')
        .equation(lhs='(Dt+mu)*x[i]', rhs='-eps*x[i]^3', population='pop')
        .build()
    )

DEFAULT_FUNDAMENTAL = {'mu': 1.0, 'eps': 0.02, 'D': 1.0}

METADATA = {
    'k_default': 2,
    'ell_default': 0,
    'recommended_external_fields': [('dx', 1), ('dx', 1)],
    'tau_max': 8.0,
    'tau_step': 0.5,
    'fixed_point_index_default': 0,
}
\end{lstlisting}

Three things to notice, each of which we will see surface later:

\begin{itemize}
  \item The action text is the \msrjd{} action of \eqref{eq:ouquartic}. The
        response field is written \code{xt} (the engine's convention for the
        conjugate of \code{x}). The piece \code{xt[i]*((Dt+mu)*x[i] +
        eps*x[i]\^{}3)} is the deterministic dynamics, and \code{- D*xt[i]\^{}2}
        is the white-noise vertex. The cubic \code{eps*x[i]\^{}3} is the single
        interaction vertex that all the loop diagrams are built from.
  \item \code{METADATA['recommended\_external\_fields']} is
        \code{[('dx', 1), ('dx', 1)]} --- two copies of the physical-field leg
        \code{dx}. This is the autocorrelator $C_{xx}(\tau)$: a two-point
        correlator ($k = 2$) of the physical field with itself.
  \item \code{METADATA['ell\_default']} is \code{0} --- the theory's
        \emph{suggested} run is tree-level only. We will deliberately override
        that to two loops in a moment, which is the whole point of the example.
\end{itemize}

\section{Step 2 --- Introspect the model}

It is good practice to print the model before you run it, so you can confirm the
engine is about to compute what you think it is. \code{describe\_model} does
exactly that --- it reads the model dict (and, optionally, the module's
docstring and \code{METADATA}) and prints a tidy report:

\begin{lstlisting}
model, mod = dd.load_theory(THEORY)
dd.describe_model(model, mod)
print('\nphysical fields:', dd.field_names(model))
\end{lstlisting}

The output is:

\begin{lstlisting}[style=console]
--------------------------------------------------------------------------
  OU Quartic (white noise)
--------------------------------------------------------------------------
Domain         : temporal ODE (time-only)
Fields         : x -- variable
Response fields: xt
Parameters     :
    mu = 1.0  (positive)
    eps = 0.02  (positive)
    D = 1.0  (positive)
Mean-field saddle (solved by the pipeline): xstar
Governing eqn  : (Dt+mu)*x[i] = -eps*x[i]^3
Suggested run  : k=2, max_ell=0

The simplest nonlinear stochastic process in the framework: an
Ornstein-Uhlenbeck variable with a cubic restoring nonlinearity, driven by
*white* Gaussian noise,
    dx/dt = -mu*x - eps*x^3 + xi,    <xi(t) xi(t')> = 2 D delta(t - t').
...

physical fields: ['dx']
\end{lstlisting}

Read the report top to bottom. The \emph{Domain} line tells you this is a
temporal (time-only) model, not a spatial one --- which will route the
calculation down the temporal branch and produce a $C(\tau)$ curve rather than a
$C(x,\tau)$ surface. The \emph{Parameters} block lists the three numeric
parameters with their defaults and domains. The line \emph{Mean-field saddle
(solved by the pipeline): xstar} is important: \code{xstar} is \emph{not} a
parameter you supply --- it is the deterministic fixed point $x^{*}$ that the
engine solves for. Parameters whose names end in \code{star} are always
saddle quantities, solved, never passed in. The \emph{Governing eqn} line echoes
the dynamics in $\code{lhs} = \code{rhs}$ form, and the \emph{Suggested run} line
surfaces the theory's own \code{METADATA} recommendation ($k=2$, tree-level).

The trailing paragraph is the theory's docstring, trimmed down to its physics
prose --- \code{describe\_model} deliberately strips the changelog and
provenance boilerplate so you see the model, not its history.

\begin{gotcha}
\code{describe\_model} both \emph{prints} the report and \emph{returns} it as a
string. In a notebook cell that means you will see the text twice (once from the
print, once as the cell's echoed return value) unless you assign or suppress the
return. This is harmless, just a touch noisy.
\end{gotcha}

Two more one-line helpers are handy here. \code{dd.field\_names(model)} returns
the list of field-leg names you can use in \code{external\_fields} --- here it
prints \code{['dx']}. \code{dd.is\_spatial(model)} returns \code{False} for this
model and would return \code{True} for a PDE; it is the single predicate that
routes everything spatial versus temporal.

\section{Step 3 --- Configure the run}

Every choice a notebook can make lives in one object: a \code{Config}. It is a
\code{@dataclass}, which means Python auto-generates its constructor from the
declared fields, so you set only what you care about and inherit theory defaults
for the rest. Here is our configuration:

\begin{lstlisting}
cfg = dd.Config(
    k=2, max_ell=2,                          # C_xx(tau), tree + 1-loop + 2-loop
    external_fields=[('dx', 1), ('dx', 1)],  # physical field autocorrelator
    parameters={'mu': 1.0, 'eps': 0.02, 'D': 1.0},
    tau_max=8.0, tau_step=0.5,
    parallel=False,                          # serial (no fork in a notebook)
)
\end{lstlisting}

Field by field:

\begin{description}
  \item[\code{k=2}] The correlator order --- the number of external legs. Two
        legs means a two-point function.
  \item[\code{max\_ell=2}] The loop order $L$. \code{0} is tree-level,
        \code{1} adds one-loop, \code{2} adds two-loop. We override the theory's
        suggested \code{0} here precisely so we can watch the loop corrections.
        The engine returns each loop order separately and the running sum.
  \item[\code{external\_fields=[('dx', 1), ('dx', 1)]}] The two legs, each a
        \code{(field\_name, slot)} tuple. Both are \code{dx}, so this is the
        physical-field autocorrelator $C_{xx}$. This matches the theory's
        \code{recommended\_external\_fields}; we pass it explicitly for clarity.
  \item[\code{parameters=\{...\}}] The numeric parameter values, by canonical
        name. These override the theory's defaults. (For backward compatibility
        there is also a \code{fundamental=} alias that means the same thing; the
        \code{Config.\_\_post\_init\_\_} mirrors whichever you set onto the
        other, with \code{parameters} winning if you somehow set both.)
  \item[\code{tau\_max=8.0, tau\_step=0.5}] The temporal lag grid: the engine
        builds $\tau \in [-8, 8]$ in steps of $0.5$. These default to the
        theory's \code{METADATA} values (also $8.0$ and $0.5$ here), but we set
        them so the grid is explicit.
  \item[\code{parallel=False}] Run serially. This is deliberate and important on
        a Mac: the engine's parallel path forks, and forking a Jupyter kernel
        after the GUI/BLAS libraries have initialised can crash the kernel. Serial
        is the safe default in a notebook, and for this tiny model it is plenty
        fast.
\end{description}

The remaining \code{Config} fields all have sensible defaults and are not needed
here. (\code{output} defaults to \code{'cumulant'}; setting it to \code{'moment'}
or \code{'central\_moment'} would assemble full moments from the cumulants, at the
cost of a few extra backend runs --- the subject of a later chapter. The spatial
fields, the $k\!\ge\!3$ slicing fields, and the Dyson-dressing fields are all
inert for a temporal $k=2$ run.)

\section{Step 4 --- Run the pipeline}

Now the single call that drives the entire \msrjd{} chain --- enumerate
diagrams, build the propagator, solve the mean-field saddle, perform the loop
integrals, and assemble the cumulant:

\begin{lstlisting}
res = dd.run(model, cfg, mod)
print(dd.summary(res))
dd.plot_cumulant(res, cfg, model)   # theory only
plt.show()
\end{lstlisting}

\code{dd.run} is thin orchestration. Reading
\file{notebooks/daedalus.py:487} onward, what it does before calling the engine
is:

\begin{enumerate}
  \item \textbf{Resolve the loop order} (\file{notebooks/daedalus.py:494}):
        \code{cfg.max\_ell} if set, else \code{METADATA['ell\_default']}, else
        \code{0}. Here that is \code{2}, from our \code{Config}.
  \item \textbf{Resolve $k$} (\file{notebooks/daedalus.py:499}): an explicit
        \code{cfg.k} wins, but if you \emph{also} passed \code{external\_fields}
        and the leg count disagrees with \code{k}, \code{run} raises a
        \code{ValueError} rather than silently picking one. Here $k=2$ and we
        passed two legs, so they agree.
  \item \textbf{Resolve the external legs} (\file{notebooks/daedalus.py:517}):
        an explicit \code{external\_fields} (ours) is used verbatim. Only if you
        left it unset would \code{run} fall back to the theory's recommendation,
        and only if \emph{that} is stale or names a missing field would it
        auto-build $k$ copies of the first field's leg.
  \item \textbf{Layer the parameters} (\file{notebooks/daedalus.py:538}): start
        from the model's own parameter defaults, update with the module's
        \code{DEFAULT\_FUNDAMENTAL}, then update with \code{cfg.parameters}. Each
        layer overrides the previous one, so a freshly loaded theory runs out of
        the box and your \code{Config} can override any subset. Here all three
        layers agree on $\mu=1,\ \varepsilon=0.02,\ D=1$.
  \item \textbf{Wire the temporal grid} (\file{notebooks/daedalus.py:578}):
        \code{tau\_max} and \code{tau\_step} from the \code{Config} (or
        \code{METADATA}).
  \item \textbf{Call the engine} (\file{notebooks/daedalus.py:586}):
        \code{res = compute\_cumulants(**kw)}.
\end{enumerate}

After the engine returns, \code{run} stamps three convenience keys onto the
result so downstream code never has to thread the config around separately:
\code{res['\_cfg']} (the \code{Config}), \code{res['\_model']} (the model dict),
and \code{res['\_resolved']} (a small dict recording the resolved \code{k},
\code{max\_ell}, \code{external\_fields}, and \code{parameters}). We will reach
for \code{res['\_resolved']['parameters']} in the simulation step.

\subsection{The seven-phase trace}

We ran with \code{verbose=False}, so the run is quiet and \code{summary} prints a
two-line digest:

\begin{lstlisting}[style=console]
theory : 'OU Quartic (white noise)'
k      : 2    max_ell : 2
fields : ['dx']   spatial_dim : 0
\end{lstlisting}

That digest comes from \code{summary} (\file{notebooks/daedalus.py:1087}), which
reads the \code{\_resolved} and \code{\_model} stamps: it confirms we ran the OU
quartic theory at $k=2$, two loops, on the single field \code{dx}, with spatial
dimension $0$ (a temporal model).

If you instead pass \code{verbose=True} in the \code{Config}, the engine narrates
its work as a numbered seven-phase trace. These seven labels are the backbone of
the entire pipeline, and most of the remaining chapters of this manual are a deep
dive into one or more of them. The labels are emitted from
\file{pipeline/compute.py}; here is the sequence with what each phase does:

\begin{lstlisting}[style=console]
[1/7] FieldTheory.expand (taylor_order=...)...
[2/7] Build propagator (K_ker -> K_ft -> G_ft -> D_delta)...
[3/7] Solve MF self-consistency...
[4/7] Compute numerical poles + residue matrices...
[5/7] Enumerate prediagrams (k=2, max_ell=2)...
[6/7] Classify coefficient factors per diagram...
[7/7] Phase J: compute_correction_td per ell (0..2)...

Done.  k=2, max_ell=2, <N> unique diagrams, <M> tau points.
\end{lstlisting}

Phase by phase:

\begin{enumerate}
  \item[{[1/7]}] \textbf{FieldTheory expand}
        (\file{pipeline/compute.py:312}). Taylor-expand the symbolic action
        around the saddle and classify every term by its ``bigrade'' (how many
        response and physical field factors it carries). This is what separates
        the free quadratic part from the interaction vertices.
  \item[{[2/7]}] \textbf{Build propagator}
        (\file{pipeline/compute.py:355}). Construct the free two-point function
        $\Gret$ from the bilinear part of the action. For the OU model this is
        the familiar response propagator with its single pole.
  \item[{[3/7]}] \textbf{Solve mean-field self-consistency}
        (\file{pipeline/compute.py:362}). Find the saddle $x^{*}$. For our
        single-well $\mu > 0$ case the saddle is $x^{*} = 0$.
  \item[{[4/7]}] \textbf{Numerical poles + residues}
        (\file{pipeline/compute.py:630}). Substitute the numeric parameter
        values and find the propagator's poles and residue matrices in the
        frequency variable $\omega$.
  \item[{[5/7]}] \textbf{Enumerate prediagrams}
        (\file{pipeline/compute.py:637}). Generate every topologically distinct
        Feynman diagram with $k=2$ external legs up to two loops, deduplicating
        isomorphic graphs.
  \item[{[6/7]}] \textbf{Classify coefficient factors}
        (\file{pipeline/compute.py:668}). Attach to each diagram its scalar
        prefactor --- the symmetry factor $\Sym$ and the product of vertex
        couplings.
  \item[{[7/7]}] \textbf{Phase J time-domain integration}
        (\file{pipeline/compute.py:709}). Perform the actual loop integrals, in
        the time domain, separately for each loop order $\ell = 0, 1, 2$. This is
        the numerical heart of the temporal pipeline.
\end{enumerate}

When verbose, the engine closes with a one-line \code{Done.} summary and, if it
tracked them, a per-phase wall-time breakdown
(\file{pipeline/compute.py:947}). The phase that dominates the clock here is
\code{[7/7]}, the loop integration --- which is why so much of this manual is
devoted to it.

\begin{note}
The seven phases run in this order for the \emph{temporal} branch. A spatial
(PDE) model takes a different route after phase~3 --- it keeps the momentum
symbolic (the Laplacian cannot be folded into the $\omega$-pole finder) and
integrates the loops in momentum space instead. The spatial chapters cover that
fork; for a temporal model like this one, the seven-phase trace above is the
complete story.
\end{note}

\section{The result: $C_{xx}(\tau)$}

\code{dd.run} returns a result dictionary, and for a temporal $k=2$ run the keys
that matter are:

\begin{description}
  \item[\code{res['tau\_grid']}] the lag grid $\tau$, here $[-8, 8]$ in steps of
        $0.5$.
  \item[\code{res['C\_tau']}] the full correlator $C_{xx}(\tau)$ summed over all
        requested loop orders --- a NumPy array over \code{tau\_grid}.
  \item[\code{res['C\_tau\_by\_ell']}] a dict \code{\{ell: array\}} giving the
        per-loop-order pieces. \code{C\_tau\_by\_ell[0]} is the tree-level curve;
        the full \code{C\_tau} equals the sum over the available orders.
  \item[\code{res['total\_C']}] a callable \code{f(*tau) -> complex} that
        evaluates the correlator at an arbitrary lag, with
        \code{total\_C\_by\_ell} its per-order version.
\end{description}

\code{dd.plot\_cumulant(res, cfg, model)} dispatches on the shape of the result.
Here it sees a temporal, non-spatial, $k=2$ run with a real \code{C\_tau}
array, so it routes to the temporal line plotter, which draws
$C_{xx}(\tau)$ versus $\tau$ with one curve per cumulative loop order (tree,
tree+1loop, tree+1loop+2loop). The curve is symmetric in $\tau$, peaks at
$\tau = 0$, and decays on the relaxation timescale $1/\mu = 1$. At this point we
have the \emph{theory only}; the simulation overlay is the next section.

The single most physically meaningful number on the plot is the equal-time value
$C_{xx}(0)$ --- the variance of $x$. For the \emph{linear} OU process it would be
exactly $D/\mu = 1$. The cubic term pushes it down, and reading it off the
per-order arrays:

\begin{lstlisting}
i0 = int(np.argmin(np.abs(np.asarray(res['tau_grid']))))   # the tau = 0 index
C0_tree = np.real(res['C_tau_by_ell'][0])[i0]              # tree level
C0_loop = np.real(res['C_tau'])[i0]                        # tree + 1- + 2-loop
\end{lstlisting}

The first line finds the array index where $\tau$ is closest to $0$ (a robust way
to locate the equal-time row without assuming the grid is symmetric to the
floating-point bit). The tree value is \code{C0\_tree} $= 1.0000$ --- the bare
linear variance $D/\mu$. The loop-corrected value is \code{C0\_loop} $= 0.9496$:
the one- and two-loop self-energy corrections have shifted the variance down by
about $5\%$, exactly the suppression the cubic nonlinearity is supposed to
produce.

\section{Step 5 --- An independent simulation}

A number from the diagrammatics is only as trustworthy as the check against an
\emph{independent} computation. The example notebook integrates the SDE
\eqref{eq:ouquartic} directly with an Euler--Maruyama scheme --- code that knows
nothing about Feynman diagrams --- and estimates $C_{xx}(\tau)$ from the
time series. The crucial part is that the simulation must run the \emph{same
physics} as the theory; it pulls the resolved parameters straight off the result
dict:

\begin{lstlisting}
from models.ou_langevin_sim_numba import sim_ou_quartic_numba
from models.cumulant_estimator import estimate_kpoint_slices
fp = res['_resolved']['parameters']          # same physics as the theory run
# Cast to plain Python floats -- under the Sage kernel the resolved params are
# Sage ring elements, which numba's njit cannot type.
mu = float(fp['mu']); eps = float(fp['eps']); D = float(fp['D'])
\end{lstlisting}

\begin{gotcha}
The explicit \code{float(...)} casts are not cosmetic. When the pipeline runs
under the SageMath kernel, the resolved parameters come back as elements of
Sage's symbolic ring, not Python floats. The simulator is compiled with
\code{numba}'s \code{njit}, and numba cannot infer a machine type for a Sage ring
element --- so without the cast the JIT compilation fails. Reading the parameters
off \code{res['\_resolved']['parameters']} (rather than re-typing them) also
guarantees the simulation uses \emph{exactly} the values the theory used.
\end{gotcha}

The discretisation choices follow from the physics: the integration step
\code{dt\_sim} must be well below the relaxation time $1/\mu$, and the time
series is binned \emph{finely} (\code{dt\_bin} $\ll 1/\mu$) so that the
bin-averaged variance is effectively the instantaneous $C_{xx}(0)$ the theory
reports --- a coarse bin would smooth the fastest fluctuations and bias
$C_{xx}(0)$ low. The notebook runs three independent realisations of total
length $T = 2\times 10^{5}$ each, so it can report an error bar:

\begin{lstlisting}
dt_sim, dt_bin = float(0.01), float(0.02)
T_sim          = float(2.0e5)              # ~1-min run for clean error bars
N_RUNS         = int(3)
# ... build x_bins via sim_ou_quartic_numba, estimate slices, average over runs:
sim = {'tau': tau_sim, 'C': C_sim, 'C_err': C_err}
\end{lstlisting}

The estimator \code{estimate\_kpoint\_slices} is written for general $k$; for our
$k=2$ case it returns a single autocorrelation slice. The three runs are averaged
to give \code{C\_sim}, and their spread gives the standard error \code{C\_err}.
The result is packaged into a small \code{sim} dict with exactly the keys the
plotter expects for a temporal overlay: \code{'tau'}, \code{'C'}, and
\code{'C\_err'}.

The headline comparison the notebook prints is:

\begin{lstlisting}[style=console]
sim took 7.7s  (3 runs x T=2e+05)
C_xx(0):  tree = 1.0000   tree+loops = 0.9496   sim = 0.9464 +/- 0.0009
\end{lstlisting}

This is the payoff of the whole chapter. The bare tree-level (linear) prediction
is $1.0000$. The two-loop-corrected theory says $0.9496$. The independent
simulation measures $0.9464 \pm 0.0009$. The loop corrections move the theory
from the linear value onto the simulation: the residual gap of about $0.003$ is
roughly three standard errors and reflects the truncation at two loops plus the
finite-bin and finite-$T$ biases of the estimator. For a calculation that
required no model-specific code --- just a theory file and a \code{Config} --- the
agreement is excellent.

\section{Step 6 --- Overlay theory and simulation}

The last call draws both on one figure. \code{plot\_cumulant} takes an optional
\code{sim} argument; when present it overlays the simulator points (with error
bars) on top of the theory curves:

\begin{lstlisting}
dd.plot_cumulant(res, cfg, model, sim=sim)
plt.show()
\end{lstlisting}

The dispatcher routes exactly as before (temporal, $k=2$) to the temporal line
plotter, but now it also reads \code{sim['tau']}, \code{sim['C']}, and
\code{sim['C\_err']} and adds them as error-bar points. The result is the
theory's loop-resolved $C_{xx}(\tau)$ curves with the simulation's measured
points sitting on the fully-corrected (tree+1loop+2loop) curve --- a visual
confirmation of the $0.9496$ versus $0.9464$ agreement at $\tau = 0$, holding
across the whole lag range.

Conceptually the figure looks like this --- the tree curve sitting above the
loop-corrected curve at the origin, with the simulation points landing on the
loop-corrected curve:

\begin{center}
\begin{tikzpicture}[scale=1.0]
  % axes
  \draw[->] (-4.4,0) -- (4.4,0) node[right] {$\tau$};
  \draw[->] (0,-0.2) -- (0,3.2) node[above] {$C_{xx}(\tau)$};
  \node[below left] at (0,0) {$0$};
  % tree curve: 1.0 * exp(-|tau|), peak at 1.0 -> scaled to 3.0 box (1.0 == 3.0)
  \draw[thick, dashed]
    plot[domain=-4:4, samples=80] (\x, {3.0*exp(-abs(\x))});
  \node[right, font=\small] at (1.6,{3.0*exp(-1.6)+0.25}) {tree ($C_{xx}(0)=1.000$)};
  % loop-corrected curve: 0.9496 * exp(-|tau|)
  \draw[thick]
    plot[domain=-4:4, samples=80] (\x, {2.849*exp(-abs(\x))});
  \node[right, font=\small] at (2.2,{2.849*exp(-2.2)+0.2}) {tree+1+2 loop ($0.9496$)};
  % simulation points on the loop-corrected curve
  \foreach \x in {-3,-2,-1,0,1,2,3} {
    \fill (\x, {2.839*exp(-abs(\x))}) circle (1.6pt);
  }
  \node[font=\small] at (2.7,1.9) {sim $0.9464\pm0.0009$};
  \draw (2.45,1.75) -- (0.05,{2.839});
\end{tikzpicture}
\end{center}

(The real plot is a Matplotlib figure with proper axes and a legend; this sketch
just fixes the relationship between the three curves in your mind.)

\section{What you just did, and where to go next}

In six steps and a few dozen lines of code you:
\begin{enumerate}
  \item loaded a declarative theory file into a \code{(model, module)} pair;
  \item printed the model to confirm its fields, parameters, and saddle;
  \item packaged your choices --- $k=2$, two loops, the autocorrelator legs, the
        parameter values --- into a single \code{Config};
  \item drove the entire seven-phase \msrjd{} pipeline with one \code{run} call;
  \item read the loop-corrected equal-time variance ($0.9496$) off the result;
  \item validated it against an independent SDE simulation ($0.9464 \pm 0.0009$)
        and overlaid the two.
\end{enumerate}

Every step here was deliberately the easy case: a single temporal field, white
noise, a small loop parameter, a symmetric saddle at the origin. The rest of the
manual generalises each axis of that easy case. The chapters on the theory
specification and theory files explain how \code{build()},
\code{DEFAULT\_FUNDAMENTAL}, and \code{METADATA} are authored. The chapters on the
field-theory core, the propagator, the mean-field solve, diagram enumeration,
type assignment, and Phase~J open up the seven phases one at a time. The chapter
on cumulants and moments explains the \code{output='moment'} path we skipped. And
the spatial chapters take the same \code{load $\to$ Config $\to$ run $\to$ plot}
spine into PDEs, where the result is a $C(x,\tau)$ surface rather than a
$C(\tau)$ curve.

If you want to experiment right now, the most instructive single change is to push
$\mu$ toward $0$ (or below it, into the double well) in \code{Config.parameters}:
that grows the effective loop parameter $g_{\mathrm{eff}} = \varepsilon D/\mu^{2}$
and stresses the perturbative series, so you can watch tree, one-loop, and
two-loop pull apart on the plot --- and watch the agreement with the simulation
degrade as truncation starts to bite.


% =====================================================================
%  PART II  --  SPECIFYING A MODEL
% =====================================================================
\part{Specifying a Model}
% ---------------------------------------------------------------------
%  Chapter 3 : Specifying a Model --- the TheoryBuilder
% ---------------------------------------------------------------------
\chapter{Specifying a Model: the \texttt{TheoryBuilder}}\label{ch:theory-spec}

In the previous chapter you loaded a theory with a single call and ran it. That
theory came from a small file under \file{theories/}, and the heart of that file
was a chain of method calls ending in \code{.build()}. This chapter opens that
chain up. By the end you will know what every method does, why the model you
describe in a few readable lines becomes a large Python dictionary full of
symbolic functions, and exactly where the bodies of those functions come from.

This is the \emph{front door} of the whole pipeline. Everything downstream ---
the propagator, the mean-field solve, the diagram enumeration, the loop integrals
--- consumes one object, the \term{model dict}, and this subsystem is the only
thing that produces it. It never integrates a diagram or solves an equation; its
single job is \emph{input specification and compilation}. Understand it well and
the rest of the manual is the story of what happens to its output.

\begin{note}
The three source files this chapter documents are
\file{pipeline/theory.py} (the builder classes and \code{build()}),
\file{pipeline/theory\_compiler.py} (which turns the text you type into Sage
lambda functions), and \file{pipeline/theory\_templates.py} (pre-canned actions
and kernels). Wherever a claim here can be checked against the code, the
\code{file:line} is given so you can read the original.
\end{note}

\section{Why a builder at all?}

To compute cumulants, the rest of the pipeline needs a \term{model dict}: a large
Python \code{dict} whose values are not merely numbers and metadata but several
\emph{Sage-aware lambda functions}. The keys with names like \code{action},
\code{mf\_bg\_conditions}, \code{kernel\_ft\_image}, \code{phi\_concrete}, and
\code{mf\_equations} all map to callables. Each such lambda takes a single runtime
argument conventionally called \code{ns} --- a ``namespace'' object that carries
the symbolic field and parameter variables --- and returns a \emph{SageMath
symbolic expression}.

\begin{defn}
\textbf{SageMath} is a large open-source mathematics system built on Python. Its
relevant feature here is a \emph{computer algebra system}: it can hold and
manipulate \emph{symbolic} expressions (variables, functions, integrals) rather
than only floating-point numbers. The central object is the \term{Symbolic Ring},
abbreviated \code{SR}. \code{SR.var('mu')} creates a symbolic variable \code{mu};
arithmetic on \code{SR} objects builds expression trees, so
\code{SR.var('a')*SR.var('x')\^{}2} is the \emph{symbolic} $a\,x^{2}$, not a
number. Sage can differentiate, Taylor-expand, Fourier-transform, and substitute
on these trees.
\end{defn}

Writing one of these dicts by hand means writing those lambdas by hand, which
demands fluency in both SageMath's symbolic algebra and a thicket of private
naming conventions (which internal name is the response field, which is the
saddle, how a transfer function's derivatives are named, and so on). The module
docstring of \file{pipeline/theory.py:1} opens by saying exactly this: the
pipeline ``takes a Python dict (\code{HAWKES\_MODEL}) with several Sage-aware
lambdas \ldots Writing one of these by hand requires familiarity with Sage's
symbolic algebra and the framework's conventions.''

The builder is the human-friendly alternative. Instead of hand-writing the dict,
you \emph{describe} the theory with a fluent chain of method calls --- declare a
field, declare a parameter, type the action as a string, type the mean-field
equation as a string --- and call \code{.build()}. The builder accumulates your
declarations, compiles the text strings into the Sage lambdas the pipeline
expects, and assembles the dict for you.

\begin{gotcha}
The top docstring of \file{pipeline/theory.py} (lines~1--62) still reads like a
roadmap for a half-finished prototype: it lists an older API
(\code{transfer\_function}, \code{use\_action\_template}, \code{set\_action}) and
describes text compilation as future work. Much of that roadmap is now done.
\emph{Do not trust the top docstring over the method bodies} --- this chapter
follows the bodies.
\end{gotcha}

\section{Choosing a builder class}

There are three concrete builder classes, all defined near the end of
\file{pipeline/theory.py}:

\begin{description}
  \item[\code{TemporalTheoryBuilder}] (\file{pipeline/theory.py:2087}) ---
        time-only theories: a scalar variable $x(t)$, or several coupled
        time-only variables, with no spatial extent. It carries the
        temporal-only methods (populations, convolution kernels, coloured-noise
        embedding). Its \code{build()} asserts that \emph{every} field is
        non-spatial.
  \item[\code{SpatialTheoryBuilder}] (\file{pipeline/theory.py:2102}) --- spatial
        PDE / field theories, a field $\varphi(x,t)$. It carries the spatial-only
        methods (\code{spatial\_dim}, \code{boundary}, \code{initial},
        \code{dyson}, \code{reference\_diffusion}). Its \code{build()} asserts
        that \emph{at least one} field is spatial. Note its constructor defaults
        \code{n\_populations=0} (\file{pipeline/theory.py:2108}).
  \item[\code{TheoryBuilder}] (\file{pipeline/theory.py:2121}) --- a
        back-compatibility class that mixes in \emph{both} method sets and
        auto-detects spatial-versus-temporal at \code{build()} time. New code
        should prefer the two specific classes.
\end{description}

The split is done with Python \emph{mixins} --- small classes that contribute
methods. \code{TemporalTheoryBuilder} inherits from \code{\_TemporalMethods} and a
shared base \code{\_BaseTheoryBuilder}; \code{SpatialTheoryBuilder} from
\code{\_SpatialMethods} and the same base; \code{TheoryBuilder} from all three.
The auto-detect lives in a method each subclass overrides,
\code{\_resolve\_spatial\_dim}. The temporal subclass raises if any field declared
a non-zero spatial dimension; the spatial subclass raises if none did
(\file{pipeline/theory.py:2092}, \code{:2111}):

\begin{lstlisting}
class TemporalTheoryBuilder(_TemporalMethods, _BaseTheoryBuilder):
    def _resolve_spatial_dim(self) -> int:
        d = super()._resolve_spatial_dim()
        if d > 0:
            raise ValueError(...'Use SpatialTheoryBuilder...')
        return 0
\end{lstlisting}

Every builder method returns \code{self}, which is what makes the fluent
chain-of-dots style work: each call hands the builder back so the next dot can
land on it.

\section{Declaring fields}

\msrjd{} field theory doubles every observable field $\varphi$ with a
\term{response field} $\tilde\varphi$ (the ``tilde'' field). The physics is
recalled in the next section; for now the point is that declaring one physical
field implies a small family of related symbols, and the builder generates that
family for you.

\subsection{\texttt{physical\_field}: one call, three symbols}

\code{physical\_field} (\file{pipeline/theory.py:806}) is the single most
important field method. It supports two calling styles, distinguished by whether
you pass \code{natural\_name}:

\begin{description}
  \item[New style] (\code{natural\_name} omitted) --- the \code{name} you pass
        \emph{is} the natural letter. \code{.physical\_field('n')} means: the
        user-facing letter is \code{n}, and the internal \term{fluctuation field}
        is derived as \code{dn}. So you write \code{n} (or \code{dn}; both refer
        to the fluctuation) and the framework knows the internal name is
        \code{dn} (\file{pipeline/theory.py:858}).
  \item[Legacy style] (\code{natural\_name} given) --- \code{name} is the literal
        internal name and \code{natural\_name} is the separate user letter, as in
        \code{.physical\_field('dn', natural\_name='n')}. This is how the older
        hand-written theory files read.
\end{description}

Whichever style you use, two further symbols are generated automatically unless
you opt out:

\begin{itemize}
  \item \textbf{Auto-response} (unless \code{auto\_response=False}): a conjugate
        response field named \code{<natural>t}. So \code{.physical\_field('n')}
        also declares \code{nt} (\file{pipeline/theory.py:885}). It is skipped if
        a response field of that name already exists.
  \item \textbf{Auto-saddle} (unless \code{auto\_saddle=False}): a saddle-point
        parameter named \code{<natural>star} carrying \code{mean\_field=True}, so
        \code{.physical\_field('n')} also declares \code{nstar}
        (\file{pipeline/theory.py:899}). Its default domain is \code{'positive'}
        for the rate field \code{n} and free for everything else.
\end{itemize}

So the single line \code{.physical\_field('x')} quietly creates three things: the
fluctuation field \code{dx}, the response field \code{xt}, and the saddle
parameter \code{xstar}. The next two sections explain why those three exist.

\begin{defn}
A \textbf{saddle} (or \textbf{mean field}) is the deterministic steady state
$\bar\varphi$ around which the theory is expanded; internally it is the
\code{*star} parameter (\code{xstar}, \code{nstar}, \ldots). It is not a number
you supply --- the pipeline solves for it. The \textbf{fluctuation field} is
$\delta\varphi = \varphi - \bar\varphi$, the internal field the action is
Taylor-expanded in (\code{dx}, \code{dn}, \ldots).
\end{defn}

For a spatial field, pass \code{spatial\_dim=d}. The dimension is a
\emph{per-field} property by design (\file{pipeline/theory.py:79}); the
bulk-setter \code{.spatial\_dim(d)} (\file{pipeline/theory.py:122}) just stamps it
onto every field at once and becomes the default for fields declared afterwards.
Values $d \in \{1,2,3\}$ are validated end-to-end; $d=0$ is time-only. All spatial
fields in one theory must share a single non-zero dimension (a v1 constraint
enforced in \code{build()}).

\code{response\_field} (\file{pipeline/theory.py:796}) declares a response field
directly, which you rarely need given auto-response, but it is the escape hatch
when a response field has no physical partner.

\section{Declaring parameters}

\code{parameter} (\file{pipeline/theory.py:914}) records one model parameter as a
\code{ParameterSpec}. The signature and its record type are:

\begin{lstlisting}
@dataclass
class ParameterSpec:
    name: str
    indexed: bool = False         # True = vector
    matrix: bool = False          # True = N x N matrix
    domain: Optional[str] = None  # e.g. 'positive'
    default: Any = None
    mean_field: bool = False      # True for saddle quantities (nstar, ...)
    natural_name: Optional[str] = None
    indexed_by: Optional[list[str]] = None  # [] scalar, ['E'] vec, ['E','I'] mat
\end{lstlisting}

The fields that matter day-to-day are:

\begin{description}
  \item[\code{default}] the numeric value used when a run does not override it
        (scalar, list, or list-of-lists).
  \item[\code{domain}] an assumption Sage's simplifier and Fourier transform can
        exploit, most commonly \code{'positive'}. Passing
        \code{domain='positive'} creates the variable as
        \code{SR.var(name, domain='positive')}, which lets Sage's underlying CAS
        dispatch the right simplification.
  \item[\code{mean\_field}] \code{True} flags the parameter as a saddle quantity
        (the \code{*star} variables). Auto-saddle sets this for you.
  \item[\code{indexed\_by}] the modern heterogeneous-population annotation:
        \code{[]} is a scalar, \code{['E']} a vector, \code{['E','I']} a matrix.
        The legacy flags \code{indexed} / \code{matrix} say the same thing; both
        are normalised at declaration time into an internal
        \code{(is\_vec, is\_mat)} pair.
\end{description}

A \emph{matrix} parameter (say a connectivity \code{w}) becomes, at build time, a
grid of \code{SR} variables \code{w11, w12, \ldots}, each carrying the declared
domain so the CAS can reason about them (\file{pipeline/theory.py:1634}).

\section{Writing the action as text}

The action is the centrepiece of the specification, and it is typed as a string.
Before the syntax, a one-page recap of what the action \emph{is}.

\subsection{The \msrjd{} action in one page}

A stochastic equation of motion for a field $\varphi$ (a scalar $x(t)$ or a field
$\varphi(x,t)$) driven by noise $\eta$,
\begin{equation}
  \Dt\varphi \;=\; F[\varphi] + \eta,
  \qquad \avg{\eta(t)\,\eta(t')} = 2D\,\delta(t-t'),
  \label{eq:eom}
\end{equation}
is recast as a path integral by introducing the response field $\tilde\varphi$.
The partition function becomes $Z = \int\!\mathcal{D}\varphi\,\mathcal{D}\tilde\varphi\,
\ee^{-S[\varphi,\tilde\varphi]}$ with the \msrjd{} action
\begin{equation}
  S[\varphi,\tilde\varphi] \;=\; \int\!\dd t\;
  \Big[\, \tilde\varphi\,\big(\Dt\varphi - F[\varphi]\big)
        \;-\; D\,\tilde\varphi^{2} \,\Big].
  \label{eq:action}
\end{equation}
The term $\tilde\varphi\,\Dt\varphi$ carries the causal/kinetic structure,
$\tilde\varphi\,F[\varphi]$ the drift, and $-D\,\tilde\varphi^{2}$ encodes the
white Gaussian noise as a \emph{second cumulant} on the response field. This last
fact generalises: non-Gaussian white noise of $n$-th cumulant $\kappa^{(n)}$ adds
a response-field monomial of degree $n$, namely
$-(\kappa^{(n)}/n!)\,\tilde\varphi^{n}$. That is exactly why
\code{set\_action\_text}'s docstring states that ``non-Gaussian white noise is a
response-field monomial of degree $\ge 3$'' (\file{pipeline/theory.py:1040}).

\subsection{The saddle--fluctuation split}

The cumulants are computed perturbatively around the deterministic steady state,
the saddle $\bar\varphi$. Every physical field is split
$\varphi = \bar\varphi + \delta\varphi$, and the action is Taylor-expanded in the
fluctuation $\delta\varphi$. The expansion produces a constant term, a linear
(\term{tadpole}) term, a quadratic (propagator) term, and cubic-and-higher
(interaction-\term{vertex}) terms.

\begin{defn}
The \textbf{saddle equation} is the condition that the linear tadpole term
vanishes --- physically, that $\bar\varphi$ is the deterministic fixed point. For
the quartic OU model \eqref{eq:ouquartic-here} below it reads
$(\Dt+\mu)\bar x = -\varepsilon\bar x^{3}$, which at steady state ($\Dt\to 0$) is
$\mu\bar x = -\varepsilon\bar x^{3}$.
\end{defn}

\subsection{The text DSL}

\code{set\_action\_text} (\file{pipeline/theory.py:1008}) stores the action as a
Sage-syntax string; its long docstring (lines~1008--1049) is the canonical
reference. The rules are:

\begin{itemize}
  \item You type the \emph{per-population integrand} $S_i$. The index \code{i} is
        an implicit free index; at build time the integrand is summed over
        \code{i in range(n\_populations)}. For a scalar theory that sum has one
        term.
  \item Inner sums (e.g.\ over a connectivity) use Python comprehension syntax:
        \code{sum(w[i, j] * g * dn[j] for j in pop)}, where
        \code{pop = range(n\_populations)} is pre-bound.
  \item \code{Dt} is an inert symbol standing for $\Dt$, used multiplicatively
        (\code{(Dt+mu)*x[i]}).
  \item The response field is the \code{<letter>t} name (\code{xt}, \code{nt}).
  \item You may write the \emph{full physical field} \code{x[i]} and the framework
        expands it to \code{xstar[i] + dx[i]} for you, or write the fluctuation
        \code{dx[i]} explicitly. (The mechanism behind the convenient \code{x[i]}
        form is the \code{\_FullPhysicalField} wrapper, covered in
        \S\ref{sec:textcompile}.)
\end{itemize}

The quartic OU action of the previous chapter,
\begin{equation}
  S \;=\; \int\!\dd t\;\Big[\, \tilde x\,\big((\Dt+\mu)\,x + \varepsilon x^{3}\big)
        \;-\; D\,\tilde x^{2} \,\Big],
  \label{eq:ouquartic-here}
\end{equation}
is typed exactly as it reads (\file{theories/ou\_quartic.theory.py}):

\begin{lstlisting}
.set_action_text('sum(xt[i]*((Dt+mu)*x[i] + eps*x[i]^3) - D*xt[i]^2 for i in pop)')
\end{lstlisting}

Here \code{xt} is $\tilde x$, the \code{(Dt+mu)*x[i] + eps*x[i]\^{}3} piece is the
deterministic dynamics (with the cubic \code{eps*x[i]\^{}3} the single interaction
vertex), and \code{- D*xt[i]\^{}2} is the white-noise vertex. Coupled multi-field
theories carry one such term per field, with cross-coupling appearing as
off-diagonal terms.

\section{Functions and the transfer function}

Some drifts contain a nonlinear function whose \emph{Taylor coefficients at the
saddle} are the interaction couplings. The archetype is a neural \term{transfer
function} $\phi(v)$ (the firing-rate nonlinearity). Note the unfortunate name
clash: in this codebase the transfer function is the user-declared function
literally named \code{phi}, distinct from the field $\varphi$.

\code{define\_function} (\file{pipeline/theory.py:969}) records a function by name,
formal argument list, and a Sage-syntax body, e.g.\
\code{.define\_function('phi', ['v'], 'a*v\^{}2')}. The function body may reference
any declared parameter. \code{set\_transfer\_function}
(\file{pipeline/theory.py:1220}) marks which declared function plays the special
\code{phi\_concrete} role (it defaults to \code{'phi'}).

\begin{defn}
A SageMath \textbf{formal} (or \emph{undefined}) function is created by
\code{function('phi\_1')(v)}: an \emph{unapplied} symbolic function evaluated at
\code{v}, which Sage keeps as an opaque node that \code{taylor()} can
differentiate \emph{formally}, without ever needing a concrete body.
\end{defn}

The crucial design choice is that the action is authored with \code{phi} as a
\emph{formal} Sage function. When the field-theory core Taylor-expands the action,
Sage produces the formal-derivative symbols $\phi(\bar v),\,\phi'(\bar v),\,
\phi''(\bar v),\ldots$. The framework names these \code{phi0\_<i>},
\code{phi1\_<i>}, \code{phi2\_<i>}, \ldots (the \code{\_<i>} suffix is the 1-based
population index), and they become the \emph{vertex couplings} of the
diagrammatic expansion. The concrete numeric values of those derivatives are
substituted only later, by the \code{specializations} and
\code{mf\_bg\_conditions} lambdas.

\begin{gotcha}
If you inline the concrete body where the formal function should be, you
short-circuit the entire auto-Taylor machinery. The Hawkes template docstring is
explicit (\file{pipeline/theory\_templates.py:319}--\code{:345}): using the
concrete expression makes the auto-expander ``see \code{phi(0)=0} etc.\ as numeric
and never register symbols, leaving an uncancelled \code{(1,0)} tadpole in the
bigrade analysis.'' This is why \code{functions\_list()} returns the function as
the formal \code{lambda i, v: function(f'phi\_\{i+1\}')(v)}
(\file{pipeline/theory\_templates.py:344}).
\end{gotcha}

\section{Mean field: the two surfaces}

The saddle $\bar\varphi$ must be found. There are two authoring surfaces for the
mean-field equations: a modern one preferred for new theories, and a legacy one.

\subsection{The modern surface: \texttt{.equation()}}

\code{equation} (\file{pipeline/theory.py:1087}) declares one residual of a
\term{DAE} system: a \emph{differential-algebraic equation} system. Each call adds
one equation $\text{LHS} = \text{RHS}$, i.e.\ the residual
$\text{LHS}-\text{RHS}=0$. At the mean-field point the framework substitutes
$\Dt\to 0$ on the LHS and solves the resulting algebraic system by
\term{multi-start Newton}: seed the solver from many random starting points,
collect the converged roots, sort them, and pick the \code{fixed\_point\_index}-th.
The quartic OU saddle equation is declared as

\begin{lstlisting}
.equation(lhs='(Dt+mu)*x[i]', rhs='-eps*x[i]^3', population='pop')
\end{lstlisting}

The method does real validation (\file{pipeline/theory.py:1143}--\code{:1184}),
which is worth knowing because the error messages are precise:

\begin{enumerate}
  \item \code{lhs}/\code{rhs} must be non-empty strings.
  \item \code{Conv(...)} on \emph{either} side raises (regex
        \code{\textbackslash bConv\textbackslash s*\textbackslash(},
        case-insensitive): the DAE assumes a \emph{stationary} mean field, so any
        kernel convolution of a constant must be pre-collapsed by you (for a
        normalised kernel, that is just the constant).
  \item \code{Dt} in the \code{rhs} raises (regex
        \code{\textbackslash bDt\textbackslash b}): derivatives belong on the LHS
        only.
  \item The \code{kind} is classified \code{'differential'} if \code{Dt} appears
        as a word in the LHS, else \code{'algebraic'}.
\end{enumerate}

\code{stability\_analysis(True)} (\file{pipeline/theory.py:1187}) opts into a
linear-stability filter for bistable / differential theories: each converged root
is linearised (keeping $\Dt\to -\ii\omega$), the generalised eigenvalue problem
$(B - \ii\omega A)\,\delta x = 0$ is assembled, and \code{fixed\_point\_index} is
restricted to the linearly stable subset. It is off by default (vacuous for
all-algebraic theories).

\subsection{The legacy surface: \texttt{set\_mf\_equation}}

\code{set\_mf\_equation(saddle\_name, rhs\_text)} (\file{pipeline/theory.py:1069})
is the older per-saddle surface: you give the RHS string for each saddle directly,
as in \code{set\_mf\_equation('nstar', 'phi(vstar[i])')}. The framework
auto-emits the EOM-closure relation \code{phi0\_<i> = nstar[i]} so you need not
declare it.

The two surfaces are bridged automatically. At \code{build()} time,
\code{\_auto\_populate\_mf\_eqs\_from\_equations} (\file{pipeline/theory.py:1276})
walks each \code{.equation(...)} you declared, substitutes $\Dt\to 0$ in the LHS,
finds the single state-variable reference, infers the saddle name as
\code{<var>star}, and textually rewrites every state-variable name in the RHS to
its \code{*star} counterpart (using word-boundary regexes so \code{Em} is not hit
when rewriting \code{E}). It is a no-op if you used the legacy surface directly.
This bridge is why a theory that declares only \code{.equation(...)} still
populates the legacy machinery the field-theory sanity check depends on.

\section{Kernels}

A \term{kernel} $g(t)$ is a convolution filter (e.g.\ an exponential synapse
$g(t) = (1/\tau_g)\,\ee^{-t/\tau_g}\,\Theta(t)$). Diagram propagators live in
frequency space, so the pipeline needs the kernel's Fourier image
$\hat g(\omega)$. For the exponential synapse,
\[
  \hat g(\omega) \;=\; \frac{1}{1 + \ii\,\omega\,\tau_g}.
\]

\code{define\_kernel} (\file{pipeline/theory.py:370}) is the text surface. You must
supply at least one of two keyword strings (it raises if both are \code{None}):

\begin{description}
  \item[\code{freq\_image}] the frequency image directly, e.g.\
        \code{'1/(1 + I*omega*tau\_g)'}. This is the fast path --- no transform is
        attempted.
  \item[\code{time\_expr}] the time-domain $g(t)$, e.g.\
        \code{'(1/tau\_g)*exp(-t/tau\_g)*heaviside(t)'}. The compiler then
        \emph{symbolically} Fourier-transforms it.
\end{description}

The transform is done by \code{make\_kernel\_ft\_image\_lambda}
(\file{pipeline/theory\_compiler.py:1493}), which calls
\code{msrjd.core.field\_theory.fourier\_transform} --- a wrapper around Sage's
\code{integrate} computing $\int g(t)\,\ee^{-\ii\omega t}\,\dd t$ (the
$\ee^{-\ii\omega t}$ convention, consistent with the exponential-synapse image
$\hat g(\omega)=1/(1+\ii\omega\tau_g)$ below). The text may
reference \code{i}/\code{j} for vector/matrix kernels, e.g.\
\code{'1/(1 + I*omega*tau\_g[i, j])'}.

\begin{gotcha}
The symbolic Fourier transform is best-effort. Sage handles
\code{heaviside(t)*exp(-t/tau)} for typical neural kernels, but harder forms fail;
\code{make\_kernel\_ft\_image\_lambda} then raises a \code{ValueError} telling you
to supply \code{freq\_image} directly (\file{pipeline/theory\_compiler.py:1530}).
When in doubt, give the frequency image.
\end{gotcha}

\section{Noise beyond white Gaussian}

White Gaussian noise is the $-D\,\tilde\varphi^{2}$ term and needs no special
machinery. Two generalisations have dedicated surfaces.

\subsection{Higher cumulants: \texttt{declare\_cgf\_term}}

\code{declare\_cgf\_term} (\file{pipeline/theory.py:461}) declares one row of a
\term{CGF} --- cumulant generating functional --- term: a non-closed-form noise
cumulant. The key arguments are \code{order} (the cumulant order: 2 for
$\kappa^{(2)}$, 3 for $\kappa^{(3)}$, \ldots), a Sage-syntax \code{coefficient},
an optional time-domain \code{kernel} factor, and either \code{response\_field}
(legacy single field) or \code{response\_legs} (a per-leg list for cross-field
cumulants). Rows sharing a \code{name} and \code{order} are summed into one
cumulant.

\subsection{Coloured noise and the Markovian embedding}

\begin{defn}
\textbf{White} noise is $\delta$-correlated,
$\avg{\eta\,\eta'}\propto\delta(\tau)$; \textbf{coloured} noise has a finite
correlation time, e.g.\ $\avg{\eta\,\eta'}\propto c\,\ee^{-|\tau|/\tau_c}$ (a
single Lorentzian in frequency). Coloured noise makes the loop integrals
expensive --- they would hang in \code{scipy.nquad} at loop order $\ge 1$.
\end{defn}

The standard cure is a \term{Markovian embedding}: introduce an auxiliary
Ornstein--Uhlenbeck field driven by \emph{white} noise whose stationary
autocorrelation is exactly the desired exponential, and couple the original field
to it linearly. \code{markovianize(True)} (\file{pipeline/theory.py:429}, default
on) toggles this. When on, \code{build()} walks the declared CGF terms and
rewrites every row whose kernel matches the \code{c*exp(-abs(tau)/tauc)} template
into a white-noise auxiliary field plus a linear filter, \emph{before} the
compiler runs. The rewrite itself lives in
\file{pipeline/colored\_to\_markovian.py}; the builder only flips the switch.

\section{Templates}

For recurring model families there are pre-canned templates in
\file{pipeline/theory\_templates.py}, so you do not retype a standard action.

\begin{description}
  \item[\code{HawkesAction}] (\file{pipeline/theory\_templates.py:69}) --- the
        canonical 2-population Hawkes action generator. Its methods each return a
        lambda: \code{action()} (the full \msrjd{} action),
        \code{phi\_concrete()}, \code{specializations()},
        \code{mf\_bg\_conditions()}, \code{mf\_equations()},
        \code{mf\_substitutions()}, and
        \code{functions\_list()} (the formal $\phi_i$).
  \item[\code{ExpSynapticKernel}] (\file{pipeline/theory\_templates.py:353}) ---
        the exponential synapse with image
        $\hat g(\omega)=1/(1+\ii\omega\tau_g)$.
  \item[\code{DeltaKernel}] (\file{pipeline/theory\_templates.py:374}) --- an
        instantaneous $g(t)=\delta(t)$. Its \code{kernel\_ft\_image()} returns
        \code{None}; its \code{extra\_specializations()} returns
        \code{lambda ns: \{ns.g: ns.delta\_D\}} so the kernel symbol is replaced
        by the Dirac $\delta$ at propagator-construction time.
  \item[\code{GTaSNoise}] (\file{pipeline/theory\_templates.py:400}) --- a
        Bernoulli + Gaussian external-rate process supplying $\kappa^{(2)}$,
        $\kappa^{(3)}$, $\kappa^{(4)}$.
\end{description}

\code{use\_action\_template} (\file{pipeline/theory.py:665}) wires a template in by
pulling its seven companion lambdas off it (\code{action}, \code{phi\_concrete},
\code{specializations}, \code{mf\_bg\_conditions}, \code{mf\_equations},
\code{mf\_substitutions}, \code{functions\_list}). \code{add\_gtas\_noise}
(\file{pipeline/theory.py:591}) is a heavier convenience: it auto-declares the
GTaS fields and parameters and re-emits the action with the noise hook installed.

\begin{gotcha}
\code{add\_gtas\_noise} declares the response field \code{mt} with
\code{auto\_response=False} \emph{on purpose}. Leaving auto-response on would
declare \code{mt} twice, which produces a duplicate \code{mt1} generator in the
bigrade \code{PolynomialRing} and the hard error \code{ValueError: variable name
'mt1' appears more than once} (\file{pipeline/theory.py:606}).
\end{gotcha}

\section{Spatial theories and the operator IR}

For a spatial field $\varphi(x,t)$ the action contains spatial differential
operators. On a diagram leg of wavevector $k$ and frequency $\omega$ their Fourier
images --- the \term{form factors} the spatial integrator attaches to each leg ---
are
\[
  \Lap \to -k^{2}, \qquad \Dt \to -\ii\omega, \qquad \partial_{x_i} \to \ii k_i,
\]
and compositions multiply ($\nabla^{4}\to k^{4}$). There are two authoring styles.

\subsection{Plain (v1): the Laplacian as a bare symbol}

In a spatial theory \code{build()} registers a \code{Laplacian} operator symbol
that is used multiplicatively, exactly like \code{Dt}. A reaction--diffusion
action with a $g\,p^{2}$ vertex reads
\code{pt*(Dt(p) + mu*p - DD*Lap(p) + g*p\^{}2) - T*pt\^{}2}. This style is enough
whenever the derivative acts on a single leg.

\subsection{Operator IR (v2): \texttt{.operator\_ir()}}

\begin{defn}
The \textbf{operator IR} (intermediate representation) makes \code{Lap(phi)},
\code{Dt(phi)}, \code{Dx(phi, i)} \emph{argument-binding calls} rather than bare
multiplicative symbols. Turn it on with \code{.operator\_ir()}
(\file{pipeline/theory.py:1053}). It is required for
\emph{derivative-inside-nonlinearity} vertices, where a bare multiplicative symbol
could not say \emph{which} leg the derivative acts on.
\end{defn}

The motivating example is KPZ: $(\partial_x h)^{2}$ is \emph{not} the same vertex
as $\partial_x(h^{2})$, and only a call that binds its argument can capture the
difference. The KPZ action (\file{theories/kpz\_1d.theory.py}) is

\begin{lstlisting}
.set_action_text('ht*(Dt(h) + mu*h - D*Lap(h) - (lam/2)*Dx(h, 0)^2) - T*ht^2')
.operator_ir()
\end{lstlisting}

The lowering machinery lives in \file{pipeline/spatial\_operator\_ir.py}, driven
from the compiler's \code{\_lower\_operator\_ir\_action}
(\file{pipeline/theory\_compiler.py:720}). It applies operator linearity,
annihilates operators acting on the homogeneous saddle (\code{kill\_means}),
replaces each atomic \code{Op(fluctuation)} with a fresh ring generator (the
$u=\delta\varphi,\,v=\nabla^{2}\delta\varphi$ trick), and classifies the
generators into \term{bilinear} (folded into the propagator) versus \term{vertex}
(carrying per-leg form factors).

\begin{gotcha}
A derivative \emph{vertex} is restricted to base field-degree 1 (\emph{per-leg},
like KPZ $(\partial\varphi)^{2}$) or 2 (\emph{composite}, like Model~B
$\nabla^{2}(\varphi^{2})$ or Burgers $\partial_x(\varphi^{2})$); any other degree
raises \code{NotImplementedError} (\file{pipeline/theory\_compiler.py:784}). A
bilinear \code{Dx} on a transverse axis ($\ne 0$ for $d\ge 2$) also raises
(\code{:837}).
\end{gotcha}

\subsection{Boundary, initial conditions, and coupled diffusion}

\code{boundary(mode, ...)} (\file{pipeline/theory.py:148}) accepts
\code{'infinite'} or \code{'periodic'} (Dirichlet/Neumann/Robin are explicitly
v2); periodic requires a \code{length=}. \code{initial(mode)}
(\file{pipeline/theory.py:182}) supports only \code{'stationary'} in v1. For
coupled spatial theories with \emph{unequal} diffusion constants there is a
\term{Dyson--Duhamel} dressing of the retarded edges: \code{dyson(mode, order)}
(\file{pipeline/theory.py:206}, with sugar \code{dyson\_order}) sets the
fixed-order truncation, and \code{reference\_diffusion(D0)}
(\file{pipeline/theory.py:276}) sets the scalar reference $D_0$ in the split
$\mathcal{D} = D_0\,I + \hat{\mathcal{D}}$. Declaring any of these on a non-spatial
theory is a hard error.

\section{How text becomes lambdas}\label{sec:textcompile}

This is the machinery section-within-the-chapter: how a string like
\code{'sum(xt[i]*((Dt+mu)*x[i] + ...) ...)'} becomes a callable that returns an
\code{SR} expression. It all happens in \file{pipeline/theory\_compiler.py}, driven
from \code{\_compile\_text\_declarations} (\file{pipeline/theory.py:1367}).

\subsection{The preparse-then-eval core}

User text uses the mathematician's \code{\^{}} for powers and bare integer
literals. Sage normally rewrites such source with a \term{preparser} before
running it (most importantly \code{\^{}}~$\to$~\code{**}, and lifting integers to
Sage \code{Integer} objects). The compiler invokes that preparser by hand in
\code{\_safe\_eval} (\file{pipeline/theory\_compiler.py:682}):

\begin{lstlisting}
import sage.all as _sage_all
from sage.repl.preparse import preparse
text = _normalize_expr_text(text)
parsed = preparse(text)
eval_globals = {**_sage_all.__dict__, **locals_dict}
return eval(parsed, eval_globals)
\end{lstlisting}

Two subtleties are spelled out in the docstring (lines~682--701):

\begin{itemize}
  \item \textbf{Why preparse:} it performs \code{\^{}}~$\to$~\code{**} (so
        \code{a*v\^{}2} means $a\,v^{2}$, not the Python bitwise XOR) and the
        integer lift.
  \item \textbf{Why globals, not locals:} the namespace is passed as \code{eval}'s
        \emph{globals}, because Python's \code{eval} only feeds its \code{locals}
        argument to the \emph{outermost} scope. A generator expression inside
        \code{sum(... for j in pop)} opens an inner scope that would \emph{not}
        see \code{locals}; putting everything in globals fixes the comprehension
        scoping.
\end{itemize}

On any failure \code{\_safe\_eval} raises a rich \code{ValueError} quoting the
text, the preparsed form, the original error, and the sorted list of available
symbols --- the single most useful debugging affordance in the subsystem.

\begin{gotcha}
A consequence of globals-not-locals: user names \emph{shadow} Sage builtins. A
parameter literally named \code{I} would shadow the imaginary unit. This is
documented as intentional (\file{pipeline/theory\_compiler.py:697}) but is a
foot-gun.
\end{gotcha}

\subsection{The wrapper classes}

Before \code{\_safe\_eval} runs, the compiler builds a namespace dict mapping each
name you used to an object that makes the text behave. These wrapper objects live
only transiently inside that dict --- they never appear in the model dict --- but
they are the heart of the DSL:

\begin{description}
  \item[\code{\_FullPhysicalField}] (\file{pipeline/theory\_compiler.py:206}) ---
        exposes \code{n[i] = nstar[i] + dn[i]}, so writing \code{nt[i] * n[i]} is
        shorthand for \code{nt[i] * (nstar[i] + dn[i])}. Its
        \code{\_\_getitem\_\_} returns \code{saddle[i] + fluct[i]}.
  \item[\code{\_FieldScalar}] (\file{pipeline/theory\_compiler.py:268}) --- wraps a
        list of \code{SR} variables so a size-1 field supports \emph{both} bare
        scalar arithmetic (\code{xt * x}) and indexed access
        (\code{xt[0] * x[0]}); size $>1$ raises a clear \code{TypeError}.
  \item[\code{\_MatrixView}] (\file{pipeline/theory\_compiler.py:316}) --- lets a
        matrix parameter accept both \code{w[i, j]} and \code{w[i][j]}.
  \item[\code{\_IndexedFormalFunction}] (\file{pipeline/theory\_compiler.py:73})
        --- turns \code{phi[i](dv[i])} into \code{function('phi\_<i+1>')(dv[i])},
        the formal symbol the auto-Taylor pass needs.
\end{description}

\subsection{The unwrap dance}

There is one Sage wrinkle these wrappers force. Sage's strict argument coercion
inside \code{BuiltinFunction.\_\_call\_\_} (used by \code{exp}, \code{log}, and
formal functions) \emph{ignores} a Python object's \code{\_sage\_()} method. So
\code{exp(xt)} in scalar mode, where \code{xt} is a \code{\_FieldScalar}, raises
\code{no canonical coercion from \_FieldScalar to SR}. The fix is a pair of
helpers --- \code{\_unwrap\_field\_arg} (\file{pipeline/theory\_compiler.py:49})
and \code{\_unwrap\_builtin} (\file{pipeline/theory\_compiler.py:342}) --- that
unwrap a size-1 field wrapper to its bare \code{SR} expression \emph{before} it
reaches Sage's coercion. Every builtin exposed to user text is wrapped this way in
\code{\_builtin\_namespace} (\file{pipeline/theory\_compiler.py:359}); any new
builtin must be wrapped the same way or scalar-mode calls will break.

\subsection{The \texttt{ns} object the lambdas receive}

\begin{defn}
The lambdas all take one argument, \code{ns} --- a \code{NamespaceForFields}
object \emph{built by the consumer} (\code{msrjd.core.field\_theory}), not by this
subsystem. It carries the symbolic variables as attributes the lambdas read:
\code{ns.<fieldname>} (lists of \code{SR} vars), \code{ns.<paramname>},
\code{ns.<kernelname>}, \code{ns.pop}, \code{ns.Dt}, \code{ns.Laplacian}, and
several private bookkeeping attributes (\code{ns.\_deriv\_rename\_subs},
\code{ns.\_all\_field\_sr\_vars}, \code{ns.\_operator\_ir}, \code{ns.delta\_D}).
\end{defn}

So \code{make\_action\_lambda} (\file{pipeline/theory\_compiler.py:851}) does not
return an expression --- it returns a \emph{closure} \code{\_action(ns)} that,
when later called with a concrete namespace, builds the eval namespace from
\code{ns}, runs \code{\_safe\_eval} on your text, optionally lowers the operator
IR, kills any ``operator applied to a constant saddle'' terms (e.g.\
$\Dt\cdot\bar v = 0$, $\Lap\cdot\bar\varphi = 0$), and returns the symbolic action.

\begin{gotcha}
The ``operator-times-saddle $\to 0$'' kill rule is \emph{duplicated}: it appears
in both \code{make\_action\_lambda} and \code{make\_mf\_bg\_conditions\_lambda},
and they walk the expression tree independently. The comment at
\file{pipeline/theory\_compiler.py:1378} is explicit that the two must stay in
sync: ``a missed operator here silently leaves a non-zero saddle term and the MF
solver fails to converge.'' It is a fragile, manually-mirrored invariant.
\end{gotcha}

\section{The \texttt{build()} method and the model dict}

\code{build()} (\file{pipeline/theory.py:1688}) is the terminal call. Its sequence
is:

\begin{enumerate}
  \item \textbf{Scalar-mode autopop} --- if no populations were declared but
        physical fields exist (and \code{n\_populations <= 1}), append a single
        size-1 population \code{pop} and bind every field to it.
  \item \textbf{Markovianize} --- if there are CGF terms, rewrite the spec via
        \code{markovianize\_spec} \emph{before} compiling.
  \item \textbf{Compile text declarations} ---
        \code{\_compile\_text\_declarations()} turns every text string into a
        lambda.
  \item \textbf{Validate required hooks} --- \code{action} is always required;
        \code{phi\_concrete} is required only when a single-argument function
        exists; \code{kernel\_ft\_image} is optional. A missing hook raises a
        \code{ValueError} listing what is absent.
  \item \textbf{Reserved-name validation} (see below).
  \item \textbf{Assemble and return} the model dict.
\end{enumerate}

\begin{gotcha}
Scalar-mode autopop is gated on \code{n\_populations <= 1}
(\file{pipeline/theory.py:1697}). Firing it for a legacy 2-population Hawkes model
would clobber the population count to 1, drop the second population's response
fields, and zero every cross-population correlator. Such theories keep
\code{populations} empty and use the flat legacy index path instead.
\end{gotcha}

\subsection{Reserved names: a scoped hard error}

Some names are owned by the framework's symbol machinery; reusing one as a field,
parameter, function, or kernel name collides --- often \emph{silently} --- with a
Fourier-transform, operator, or spatial-coordinate symbol. \code{build()}
(\file{pipeline/theory.py:1850}) checks this and raises with a rename hint. The
reserved set is \emph{scoped}:

\begin{center}
\begin{tabular}{@{}lll@{}}
\toprule
\textbf{Names} & \textbf{Reserved in} & \textbf{Why} \\
\midrule
\code{t}, \code{omega} & all theories & Fourier-transform variables \\
\code{Dt}, \code{delta\_D}, \code{delta\_Dp} & all theories & operator symbols \\
\code{k} & spatial theories & wavevector (\code{SR.var('k')}) \\
\code{Laplacian} & spatial theories & diffusion operator \\
\code{x}, \code{y}, \code{z} & spatial theories & spatial coordinates / output axis \\
\bottomrule
\end{tabular}
\end{center}

\begin{gotcha}
Note the irony: \code{x} is the conventional default field name in a
\emph{temporal} theory (as in the quartic OU model), but it is reserved in a
\emph{spatial} theory, where it is the spatial coordinate. The comment at
\file{pipeline/theory.py:1842} flags that a parameter colliding with \code{k}
would be a \emph{silent wrong-mass bug} (because \file{heat\_kernel.py} uses
\code{SR.var('k')} for the wavevector), which is precisely why \code{k} is a hard
error in spatial theories. The build raises with a rename hint such as
\code{x}~$\to$~\code{phi}.
\end{gotcha}

\subsection{The keys of the model dict}

\code{build()} assembles the dict at \file{pipeline/theory.py:1943}. The
\emph{always-present} keys are:

\begin{lstlisting}[style=console]
name, populations, iteration_saddles, index_sets, response_fields,
physical_fields, parameters, kernels, equations, operators, functions,
mf_substitutions, phi_concrete, action, naming_convention,
stability_analysis, operator_ir
\end{lstlisting}

Conditionally attached lambda keys (present only when the corresponding text was
declared): \code{mf\_bg\_conditions\_action}, \code{mf\_bg\_conditions},
\code{mf\_equations}, \code{kernel\_ft\_image}, \code{kernel\_td\_image},
\code{specializations}, \code{correlated\_noises}. Spatial-only keys (present only
when a field is spatial): \code{spatial}, \code{boundary}, \code{initial}.

\begin{gotcha}
The mean-field conditions are stored \emph{twice} under two keys, and the reason
is subtle (\file{pipeline/theory.py:1974}--\code{:1986}).
\code{mf\_bg\_conditions\_action} is the closure-baked version, pre-substituted so
the \code{(1,0)} tadpole cancels symbolically --- it is what
\code{FieldTheory.expand} uses. \code{mf\_bg\_conditions} is the solver-friendly
version that keeps raw \code{nstar} so the numerical iteration can run. If only
one exists (the legacy template path), it is used for both.
\end{gotcha}

The \code{naming\_convention} key deserves a look, because it is how the rest of
the pipeline translates between the natural letters you typed and the internal
names:

\begin{lstlisting}
{
  'fluctuation_fields': {'n': 'dn', 'v': 'dv'},     # natural -> internal
  'mean_field_saddles': {'n': 'nstar', 'v': 'vstar'},  # natural -> internal
  'mf_parameters':      ['nstar', 'vstar'],          # internal names
}
\end{lstlisting}

\section{End to end: who consumes the dict}

The model dict is read by \code{compute\_cumulants}
(\file{pipeline/compute.py:108}), the function the quickstart chapter drove
through \code{dd.run}. Its docstring lists the \emph{required} hooks ---
\code{response\_fields}, \code{physical\_fields}, \code{parameters},
\code{kernels}, \code{operators}, \code{functions}, \code{mf\_substitutions},
\code{mf\_bg\_conditions}, \code{specializations}, \code{kernel\_ft\_image},
\code{phi\_concrete}, \code{mf\_equations}, \code{action} (with
\code{correlated\_noises} optional) --- and the builder produces exactly these.
The numeric parameter values come not from the dict but from the theory file's
\code{DEFAULT\_FUNDAMENTAL} constant, and the run recommendations from its
\code{METADATA}, as the previous chapter showed.

The whole flow, from the file you write to the numbers downstream, is:

\begin{center}
\begin{tikzpicture}[
    node distance=9mm and 16mm,
    box/.style={draw=noteframe, fill=notebg, rounded corners=2pt,
                inner sep=5pt, align=center, font=\small},
    op/.style={font=\footnotesize\itshape, align=center},
    >={Stealth[length=2.4mm]}]
  \node[box] (file) {\code{theories/<name>.theory.py}\\\code{build()} function};
  \node[box, right=of file] (builder) {\code{TheoryBuilder}\\(accumulates\\declarations)};
  \node[box, below=12mm of builder] (model) {model dict\\(numbers + Sage\\lambdas)};
  \node[box, left=of model] (compute) {\code{compute\_cumulants}\\(enumerate $\to$\\integrate)};
  \draw[->] (file) -- node[op, above]{text} (builder);
  \draw[->] (builder) -- node[op, right]{\code{.build()}} (model);
  \draw[->] (model) -- node[op, above]{required\\hooks} (compute);
\end{tikzpicture}
\end{center}

\section{Two complete examples}

The two theory files behind this and the previous chapter are short enough to read
whole. The temporal one (\file{theories/ou\_quartic.theory.py}) is the quartic OU
model of \eqref{eq:ouquartic-here}:

\begin{lstlisting}
from pipeline.theory import TemporalTheoryBuilder

def build():
    return (
        TemporalTheoryBuilder('OU Quartic (white noise)')
        .population('pop', size=1)
        .physical_field('x', population='pop', description='variable')
        .parameter('mu', default=1.0, domain='positive')
        .parameter('eps', default=0.02, domain='positive')
        .parameter('D', default=1.0, domain='positive')
        .set_action_text('sum(xt[i]*((Dt+mu)*x[i] + eps*x[i]^3) - D*xt[i]^2 for i in pop)')
        .equation(lhs='(Dt+mu)*x[i]', rhs='-eps*x[i]^3', population='pop')
        .build()
    )
\end{lstlisting}

Reading it as a chain: \code{.physical\_field('x')} declares the fluctuation
\code{dx}, the response \code{xt}, and the saddle \code{xstar}; the three
\code{.parameter(...)} calls add $\mu,\varepsilon,D$ (all positive); the action
text is \eqref{eq:ouquartic-here} verbatim; and \code{.equation(...)} is the
saddle condition $(\Dt+\mu)x = -\varepsilon x^{3}$. At \code{.build()},
scalar-mode autopop is skipped (a \code{pop} population was declared), the
equation is bridged into the legacy mean-field surface, and every text string is
compiled to a lambda.

The spatial one (\file{theories/kpz\_1d.theory.py}) is 1D KPZ, and shows the
operator IR in anger:

\begin{lstlisting}
from pipeline.theory import SpatialTheoryBuilder

def build():
    return (
        SpatialTheoryBuilder('1D KPZ (per-leg gradient vertex)', n_populations=0)
        .physical_field('h', spatial_dim=1, description='interface height')
        .parameter('mu',  default=1.0, domain='positive')
        .parameter('D',   default=1.0, domain='positive')
        .parameter('lam', default=0.3, domain='real')
        .parameter('T',   default=1.0, domain='positive')
        .equation(lhs='(Dt + mu - D*Laplacian)*h', rhs='0')
        .set_action_text('ht*(Dt(h) + mu*h - D*Lap(h) - (lam/2)*Dx(h, 0)^2) - T*ht^2')
        .operator_ir()
        .boundary('infinite')
        .initial('stationary')
        .build()
    )
\end{lstlisting}

Here \code{spatial\_dim=1} marks \code{h} as a 1D field (so \code{build()} will
register a \code{Laplacian} operator and emit a \code{spatial} block); the action
uses the binding-call form \code{Lap(h)}, \code{Dt(h)}, \code{Dx(h, 0)} because
the KPZ vertex $(\partial_x h)^{2}$ is a per-leg derivative vertex that the plain
multiplicative style cannot express; \code{.operator\_ir()} switches on the
lowering; and \code{.boundary('infinite')}, \code{.initial('stationary')} pin the
spatial setup. Note that the field is named \code{h}, not \code{x} --- because in
a spatial theory \code{x} is the reserved spatial coordinate.

\section{Where to go next}

You now know how the model dict is born. The next chapter, on theory files, covers
the \code{DEFAULT\_FUNDAMENTAL} and \code{METADATA} constants that sit alongside
\code{build()}, and a later chapter develops the coloured/Markovian noise embedding
sketched here. From the other side, the field-theory core chapter picks up the
model dict and starts the seven-phase pipeline you watched run in the quickstart:
it is the first consumer of the \code{action} lambda this subsystem so carefully
compiled.

% ---------------------------------------------------------------------
%  Chapter 4 : Theory Files and Serialization
% ---------------------------------------------------------------------
\chapter{Theory Files and Serialization}\label{ch:theory-files}

The previous chapter introduced the four verbs that drive the engine, and the
first of them --- \code{load\_theory} --- read a small Python file off disk to
get its model. This chapter is about that file: what a \file{.theory.py} is,
exactly what it may contain, how it becomes the in-memory \emph{model dict} the
engine consumes, and the machinery that converts it back and forth between the
graphical theory-builder form and clean, version-controllable source text.

A theory file is the \emph{single source of truth} for one physical model. You
never edit the engine to add a model; you write one of these declarative
modules, and everything downstream --- the field-theory expansion, the
propagator, the diagram enumeration, the loop integrals --- consumes the dict
its \code{build()} produces. Because the file is the authoritative description,
it pays to understand it from the ground up: every chainable builder method,
every module-level constant, and the two distinct ``serialization'' layers that
unfortunately share a name and a directory.

We will work almost entirely from the running example of
Chapter~\ref{ch:quickstart}, \file{theories/ou\_quartic.theory.py}, with a
couple of richer files (a coloured-noise sibling and a spatial KPZ model) brought
in to show the features the simple file does not exercise.

\begin{note}
Two things in this subsystem are both called ``serialization,'' and conflating
them is the single most common source of confusion. The \emph{first} is the
\file{.theory.py} \emph{text} format --- plain, human-readable Python source,
documented in this chapter and handled by \file{pipeline/theory\_serialize.py}.
The \emph{second} is the \file{saved\_theories/<slug>/} \emph{binary} cache ---
pickled SageMath objects (\code{.sobj}) plus a \code{manifest.json}, written when
a theory is actually \emph{run}. You never need Sage to read a \file{.theory.py};
you always need Sage to load a \code{.sobj}. The last section of the chapter
disentangles the two; until then, ``theory file'' always means the text format.
\end{note}

\section{Why theory files exist}

The design promise of a theory file is twofold. First, it is the \term{single
source of truth}: one file fully describes one model --- its fields, parameters,
\msrjd{} action, mean-field equations, and (for spatial models) the boundary and
initial conditions --- with no per-model code anywhere in the engine. Second, it
is the \term{authoring and persistence boundary}: the same content can be
produced either by a human typing the file directly, or by a graphical form in a
Jupyter notebook, and it can be read back into that form for editing.

That second promise is what \file{pipeline/theory\_serialize.py} delivers. The
module is a \emph{bidirectional bridge} between a flat ``spec'' dictionary --- the
internal data the UI form collects --- and \file{.theory.py} source text:

\begin{itemize}
  \item \textbf{Forward} (\code{render\_theory\_file} / \code{save\_theory\_to\_file}):
        the theory-builder UI collects your input as a flat dict and emits a
        clean, diff-friendly \file{.theory.py} file.
  \item \textbf{Reverse} (\code{load\_spec\_from\_file}): given an existing
        \file{.theory.py}, parse its source \emph{as text} --- via Python's
        \code{ast} module, \emph{not} by executing it --- and reconstruct the
        original spec dict, so the UI's ``Load existing theory'' button can
        re-populate the form.
\end{itemize}

The arrangement is worth a diagram. There are two ways a theory file is born (a
UI form, or a human), one canonical way it is loaded (\code{load\_theory}), and a
separate binary cache that is written only when the model is run:

\begin{center}
\begin{tikzpicture}[
    >=Stealth, node distance=6mm,
    box/.style={draw, rounded corners=2pt, align=center,
                font=\footnotesize, inner sep=4pt, fill=codebg},
    file/.style={draw, align=center, font=\footnotesize\ttfamily,
                 inner sep=4pt, fill=notebg},
    lbl/.style={font=\scriptsize\itshape, color=cmt}]
  \node[box] (ui)   {theory-builder UI\\(\code{ipywidgets} form)};
  \node[box, right=28mm of ui] (hand) {hand-authored\\\file{.theory.py}};
  \node[box, below=12mm of ui] (render) {\code{render\_theory\_file(spec)}\\emits source text};
  \node[file, below=24mm of $(ui)!0.5!(hand)$] (tf) {theories/<name>.theory.py};
  \node[box, below=10mm of tf] (model) {model dict\\(\code{mod.build()})};
  \node[box, below=10mm of model] (calc) {FieldTheory / compute\_cumulants\\full\_integrator};
  \node[file, right=14mm of calc] (cache) {saved\_theories/<slug>/\\*.sobj + manifest.json};

  \draw[->] (ui)   -- node[lbl,left]{\code{\_collect()}} (render);
  \draw[->] (render) -- (tf);
  \draw[->] (hand) -- (tf);
  \draw[->] (tf)   -- node[lbl,right]{\code{load\_theory}} (model);
  \draw[->] (model) -- (calc);
  \draw[->] (calc) -- node[lbl,above,align=center]{heavy\\symbolic}(cache);
  % reverse load arrow
  \draw[->, dashed] (tf.west) .. controls +(-3,0) and +(-3,0) ..
       node[lbl,left,align=center]{\code{load\_spec}\\\code{\_from\_file}} (render.west);
\end{tikzpicture}
\end{center}

The dashed arrow is the reverse path: it reads the file back into a spec so the
form can show it. The solid path down the middle is the run-time path you met in
Chapter~\ref{ch:quickstart}.

The key correctness property of the forward/reverse pair is \term{round-trip
fidelity}: a cycle \code{spec} $\to$ \code{render} $\to$ \code{file} $\to$
\code{load} $\to$ \code{spec'} reproduces the \emph{meaningful} content of the
original spec. Because the reverse path reads the abstract syntax tree rather
than running the module, string literals (action text, equation right-hand
sides, kernel expressions) round-trip \emph{verbatim}, while comments and exact
formatting are lost. We will return to the precise limits of that guarantee in
\S\ref{sec:tf-reverse}.

\section{Anatomy of a \texttt{.theory.py}}

A theory file is an ordinary, importable Python module with exactly three public
symbols:

\begin{description}
  \item[\code{build()}] a function returning the \term{model dict} --- the
        in-memory object the rest of the pipeline consumes. Internally it
        constructs the model with a fluent \emph{builder} object whose chained
        method calls declare the fields, parameters, action, and equations.
  \item[\code{DEFAULT\_FUNDAMENTAL}] a plain dict of default numeric parameter
        values --- the ``working point'' the calculation is evaluated at.
  \item[\code{METADATA}] a plain dict of \emph{run recommendations}: the default
        cumulant order $k$, the default loop order $\ell$, which external fields
        to plot, the time grid, and (for spatial models) the spatial grid.
\end{description}

Here is the entire white-noise quartic file again, this time as the object of
study rather than a passing reference (\file{theories/ou\_quartic.theory.py},
docstring elided):

\begin{lstlisting}
from pipeline.theory import TemporalTheoryBuilder


def build():
    return (
        TemporalTheoryBuilder('OU Quartic (white noise)')
        .population('pop', size=1)
        .physical_field('x', population='pop', description='variable')
        .parameter('mu', default=1.0, domain='positive')
        .parameter('eps', default=0.02, domain='positive')
        .parameter('D', default=1.0, domain='positive')
        .set_action_text(
            'sum(xt[i]*((Dt+mu)*x[i] + eps*x[i]^3) - D*xt[i]^2 for i in pop)')
        .equation(lhs='(Dt+mu)*x[i]', rhs='-eps*x[i]^3', population='pop')
        .build()
    )


DEFAULT_FUNDAMENTAL = {'mu': 1.0, 'eps': 0.02, 'D': 1.0}

METADATA = {
    'k_default': 2,
    'ell_default': 0,
    'recommended_external_fields': [('dx', 1), ('dx', 1)],
    'tau_max': 8.0,
    'tau_step': 0.5,
    'fixed_point_index_default': 0,
}
\end{lstlisting}

Read top to bottom this is almost self-documenting. The \code{import} line names
the builder class (here \code{TemporalTheoryBuilder}; \S\ref{sec:tf-split}
explains the temporal/spatial choice). \code{build()} returns a single
parenthesised expression: a constructor call followed by a chain of
\code{.method(...)} calls, terminated by \code{.build()}. Each method declares
one ingredient of the model. \code{DEFAULT\_FUNDAMENTAL} and \code{METADATA} are
module-level constants the loader hands back alongside the model dict --- they are
\emph{not} part of \code{build()}'s return value, which is why
\code{load\_theory} returns the \emph{pair} \code{(model, module)} you saw in
Chapter~\ref{ch:quickstart}: the model dict carries the physics, the module
carries the run recommendations.

\subsection{The math the file encodes}

A theory file encodes a stochastic field theory in the \msrjd{} formalism. We
assume you know that physics; here is only the minimal vocabulary needed to read
the code.

Each physical field \code{x} (or \code{phi}, \code{h}, \code{n}, \dots) gets a
conjugate \term{response field} \code{xt} (the ``tilde'' field $\tilde x$). A
Langevin process
\[
  \frac{\dd x}{\dd t} = F[x] + \xi,
  \qquad \avg{\xi(t)\,\xi(t')} = 2D\,\delta(t-t'),
\]
is rewritten as a path integral with the \msrjd{} action $S[x,\tilde x]$. For the
running example,
\[
  \frac{\dd x}{\dd t} = -\mu x - \varepsilon x^{3} + \xi,
\]
the action --- exactly as it appears in \code{set\_action\_text} --- is
\begin{equation}
  S \;=\; \int \dd t \,\Bigl[\, \tilde x\bigl((\Dt + \mu)\,x + \varepsilon x^{3}\bigr)
       \;-\; D\,\tilde x^{2} \,\Bigr].
  \label{eq:tf-action}
\end{equation}
The string passed to \code{.set\_action\_text(...)} is this expression with
\code{xt} for $\tilde x$ and \code{\^{}} for powers
(\file{theories/ou\_quartic.theory.py:27--28}):

\begin{lstlisting}
'sum(xt[i]*((Dt+mu)*x[i] + eps*x[i]^3) - D*xt[i]^2 for i in pop)'
\end{lstlisting}

The time-derivative $\partial_t$ is the operator \code{Dt}; the quadratic
response term $-D\,\tilde x^{2}$ is the white-noise vertex; the cubic
$\varepsilon x^{3}$ is the single interaction. The \code{sum(\dots\ for i in pop)}
wrapper is a \emph{population comprehension} --- see \S\ref{sec:tf-indexing}.

\paragraph{Mean field / saddle point.} The classical (noise-free) trajectory is
the \term{mean-field saddle} $x^{*}$. Theory files declare it one of two ways.
The preferred, newer form is the \term{DAE} (differential-algebraic equation)
\code{.equation(lhs=..., rhs=..., population=...)}; in the running example
(\file{theories/ou\_quartic.theory.py:29})

\begin{lstlisting}
.equation(lhs='(Dt+mu)*x[i]', rhs='-eps*x[i]^3', population='pop')
\end{lstlisting}

reads as ``$(\partial_t + \mu)\,x = -\varepsilon x^{3}$,'' a DAE the solver
integrates to a fixed point. The legacy alternative names the saddle symbol
explicitly with \code{.set\_mf\_equation('xstar', 'rhs')}; you will meet it in
heterogeneous Hawkes models. The forward serializer prefers \code{.equation(...)}
when the spec carries an \code{equations} list and falls back to
\code{.set\_mf\_equation(...)} otherwise
(\file{pipeline/theory\_serialize.py:424--433}).

\paragraph{Propagator, cumulants, loop expansion.} After choosing $x^{*}$, the
theory is expanded around it. The \term{propagator} is the inverse of the
quadratic part of $S$; \term{cumulants} (connected correlation functions) are
computed order by order in a \term{loop expansion}. Two integers parametrise the
request, and you met both as \code{Config} fields in Chapter~\ref{ch:quickstart}:
$k$, the cumulant order (number of external legs), and $\ell$ (the
\code{max\_ell} / loop order), the highest loop order included. In the theory
file these appear in \code{METADATA} as \code{k\_default} and \code{ell\_default},
and they later become file-name keys in the binary cache (e.g.
\code{...\_k2\_l1.sobj}).

\begin{note}
The small parameter of the loop expansion for this model is roughly
$g_{\mathrm{eff}} \approx \varepsilon D / \mu^{2}$. The comment at
\file{theories/ou\_quartic.theory.py:34} records $g_{\mathrm{eff}} \approx 0.02$
at the default working point, which is why the series converges fast. This is the
same number that drove the regime discussion in Chapter~\ref{ch:quickstart}; the
theory file is simply where that working point is written down.
\end{note}

\subsection{\texttt{DEFAULT\_FUNDAMENTAL} and \texttt{METADATA}}

\code{DEFAULT\_FUNDAMENTAL} is the numeric \emph{working point}. It may contain
scalars, vectors (lists), or matrices (lists of lists). Three styles appear
across the shipped theories:

\begin{itemize}
  \item \code{\{'mu': 1.0, 'eps': 0.02, 'D': 1.0\}} --- scalar defaults
        (\code{ou\_quartic}).
  \item \code{\{\}} --- \emph{empty}, when defaults are carried on each
        \code{.parameter(default=...)} instead (as in \code{ou\_quartic\_colored}
        and \code{quadratic\_hawkes\_alpha}).
  \item \code{\{'E': [0.78, 0.81], 'w': [[0.30,0.25],[0.30,0.35]], ...\}} ---
        vector and matrix entries for heterogeneous populations
        (\code{linear\_hawkes}).
\end{itemize}

\begin{gotcha}
\code{DEFAULT\_FUNDAMENTAL} \emph{can be empty}. Several theories leave it
\code{\{\}} and carry the numbers on the individual \code{.parameter(default=...)}
calls. Both styles are valid; \code{run} merges the model's own parameter
defaults, then \code{DEFAULT\_FUNDAMENTAL}, then your \code{Config.parameters}
(the layering you saw in Chapter~\ref{ch:quickstart}), so a freshly loaded theory
runs out of the box regardless of which style its author chose.
\end{gotcha}

\code{METADATA} is the bundle of \emph{run recommendations}. The keys observed
across the shipped theories are:

\begin{center}
\begin{longtable}{@{}llp{0.46\textwidth}@{}}
\toprule
\textbf{Key} & \textbf{Example} & \textbf{Meaning} \\
\midrule
\endhead
\code{k\_default} & \code{2} & Default cumulant order. \\
\code{ell\_default} & \code{0}, \code{1} & Default loop order. \\
\code{recommended\_external\_fields} & \code{[('dx',1),('dx',1)]} &
  \code{(field, population\_index)} legs to compute/plot. \\
\code{tau\_max}, \code{tau\_step} & \code{8.0}, \code{0.5} & Time grid. \\
\code{fixed\_point\_index\_default} & \code{0} & Which saddle root to use
  (temporal only). \\
\code{spatial\_grid} & \code{[-8.0, 8.0, 65]} & Spatial sampling, as an explicit
  list or a \code{(start, stop, n)} triple. \\
\code{description} & string & Optional short label. \\
\bottomrule
\end{longtable}
\end{center}

\begin{gotcha}
The entries of \code{recommended\_external\_fields} are \emph{tuples}, e.g.
\code{('dx', 1)}, and downstream code expects the shape
\code{(field\_name, population\_index)}. They render as tuples and reload as
tuples (the reverse path recovers them with \code{ast.literal\_eval}). If you
hand-author \code{METADATA} and write them as \emph{lists}, the type would change
on reload --- so keep them tuples.
\end{gotcha}

\section{The builder DSL}

The body of \code{build()} is a fluent chain on a builder object. Each method is
a small declaration; the chain reads top to bottom like a specification. This
section walks every chainable method you are likely to write, grouped by what it
declares. The methods live on \code{TemporalTheoryBuilder} /
\code{SpatialTheoryBuilder} (imported from \code{pipeline.theory}); the forward
serializer's \code{\_emit\_*} helpers, cited per method below, are the authority
on exactly which keyword arguments each call accepts.

\subsection{Fields and parameters}

\paragraph{\code{.physical\_field('x', \dots)}.} Declares one physical field. The
positional \code{'x'} is the user-facing \emph{natural name}; the builder
auto-prefixes \code{d} to form the internal fluctuation name (\code{dx}) and
auto-creates both the conjugate response field (\code{xt}) and the saddle
(\code{xstar}). This is why the model in Chapter~\ref{ch:quickstart} reported its
field as \code{dx} even though the file wrote \code{x}, and why you almost never
write a \code{.response\_field(...)} line by hand. Optional keywords, emitted by
\code{\_emit\_field} (\file{pipeline/theory\_serialize.py:112}), are
\code{population=} (heterogeneous style) \emph{or} \code{indexed=} (legacy
style), \code{natural\_name=}, \code{latex=}, \code{description=}, and a per-field
\code{spatial\_dim=}.

\paragraph{\code{.parameter('mu', default=1.0, domain='positive')}.} Declares one
numeric parameter. The keywords (emitted by \code{\_emit\_parameter},
\file{pipeline/theory\_serialize.py:151}) are \code{default=}, the indexing
declaration (\code{indexed\_by=} new-style or \code{indexed=} legacy),
\code{domain=} (e.g. \code{'positive'}, \code{'real'}), \code{natural\_name=}, and
\code{description=}.

\begin{gotcha}
\textbf{Saddle parameters are never written by hand and never emitted.} A
parameter whose name matches the auto-generated \code{<natural>star} convention
(e.g. \code{xstar}), or one flagged \code{mean\_field}, is \emph{filtered out} by
the forward serializer (\file{pipeline/theory\_serialize.py:334--339}) because the
framework re-creates it from the physical-field declaration. Emitting it would
double-declare the saddle. The filter is even applied twice --- belt and braces ---
since \code{\_emit\_parameter} guards \code{mean\_field} again at line~181.
\end{gotcha}

\paragraph{\code{.response\_field('xt', \dots)}.} Rarely needed --- only emitted
when you \emph{explicitly} declared a response field rather than relying on the
auto-creation a \code{.physical\_field(...)} performs. Most files therefore carry
no \code{.response\_field(...)} line at all.

\subsection{Functions and kernels}

\paragraph{\code{.define\_function('phi', args=[...], expression='...')}.} Declares
a named scalar function used in the action or saddle (handler
\code{\_emit\_function}, \file{pipeline/theory\_serialize.py:189}). \code{args}
and \code{expression} are mandatory; \code{population=}, \code{latex=}, and
\code{description=} are optional. For example
\code{.define\_function('phi', args=['v'], expression='a[i]*v\^{}2', population='E')}
declares a per-population quadratic transfer function. Functions exist precisely
to support the saddle indirection some models need.

\paragraph{\code{.define\_kernel(\dots)}.} Declares a convolution kernel (handler
\code{\_emit\_kernel}, \file{pipeline/theory\_serialize.py:200}). A kernel may be
given in the time domain via \code{time\_expr=} (e.g.
\code{'(t/taug[i,j]\^{}2)*exp(-t/taug[i,j])*heaviside(t)'}) \emph{or} directly as
a frequency image via \code{freq\_image=} (e.g.
\code{'1 / (1 + I*omega*tau\_g)'}). It also accepts \code{latex\_name=},
\code{sage\_name=}, \code{description=}, and the same indexing duality as
parameters.

\subsection{The action and the mean-field equations}

\paragraph{\code{.set\_action\_text('...')}.} The heart of the model: the \msrjd{}
action of \eqref{eq:tf-action} as a single string, with \code{xt} for the
response field and \code{\^{}} for powers. The handler \code{\_emit\_action}
(\file{pipeline/theory\_serialize.py:277}) emits single-line bodies as a plain
\code{repr}; multi-line bodies are first run through \code{textwrap.dedent} and
then re-indented four spaces inside a triple-quoted block. The dedent is
load-bearing for whitespace stability --- see the gotcha in \S\ref{sec:tf-forward}.

\paragraph{\code{.equation(lhs=..., rhs=..., population=...)} (preferred) and
\code{.set\_mf\_equation('xstar', 'rhs')} (legacy).} The two saddle-declaration
forms introduced above. The new DAE form is emitted by \code{\_emit\_equation}
(\file{pipeline/theory\_serialize.py:297}), which omits \code{population=} when it
is falsy; the legacy form by \code{\_emit\_mf\_equation}
(\file{pipeline/theory\_serialize.py:292}).

\subsection{Toggles}

A handful of methods set boolean toggles or scalar policies. The forward
serializer emits each \emph{only when it departs from the builder's default}, so
files stay ``textually quiet'' --- a recurring design phrase meaning a file
contains only what is non-trivial, which keeps diffs clean. These are:

\begin{itemize}
  \item \code{.stability\_analysis(True)} --- emitted only when on
        (\file{pipeline/theory\_serialize.py:439--440});
  \item \code{.markovianize(False)} --- emitted only when explicitly off
        (\file{pipeline/theory\_serialize.py:448--449}); see
        \S\ref{sec:tf-colored};
  \item \code{.operator\_ir()} --- emitted only when on, and \emph{load-bearing}
        for derivative-vertex spatial theories
        (\file{pipeline/theory\_serialize.py:452--459}); see
        \S\ref{sec:tf-spatial};
  \item \code{.boundary(mode, **params)} / \code{.initial(mode, **params)} ---
        spatial boundary and initial conditions, emitted only when present
        (\file{pipeline/theory\_serialize.py:466--477});
  \item \code{.dyson\_order(N)} / \code{.reference\_diffusion(D0)} --- the
        Dyson--Duhamel policy and reference-diffusion split, emitted only for a
        non-\code{'off'} policy or a set $D_{0}$
        (\file{pipeline/theory\_serialize.py:484--489}).
\end{itemize}

\section{Temporal versus spatial: the builder split}\label{sec:tf-split}

The \code{import} line at the top of a theory file names one of two builder
classes. \code{TemporalTheoryBuilder} is for time-only (ODE) models like the
running example; \code{SpatialTheoryBuilder} is for spatial (PDE) models where
the field lives on $x \in \mathbb{R}^{d}$ and the action carries the Laplacian.

The forward serializer does \emph{not} ask you which one you want --- it infers it.
The choice is driven purely by whether any physical field carries a non-zero
\code{spatial\_dim} (\file{pipeline/theory\_serialize.py:370--372}):

\begin{lstlisting}
_is_spatial = any(int(f.get('spatial_dim', 0) or 0) > 0
                  for f in physical_fields)
builder_cls = 'SpatialTheoryBuilder' if _is_spatial else 'TemporalTheoryBuilder'
\end{lstlisting}

If any field declares $\code{spatial\_dim} \ge 1$, the file is emitted as spatial;
otherwise it is temporal.

\begin{gotcha}
The builder split is \emph{inferred from \code{spatial\_dim}, not declared}. A
spatial theory whose fields all forgot \code{spatial\_dim} would be emitted as a
\emph{temporal} file --- and would then compute the wrong physics. The reverse
(loading) path is forgiving: it accepts \code{TheoryBuilder} (a legacy
back-compat shim), \code{SpatialTheoryBuilder}, \emph{and}
\code{TemporalTheoryBuilder} as constructor names
(\file{pipeline/theory\_serialize.py:651--652}). But the forward path can only
emit the right class if at least one field carries the dimension.
\end{gotcha}

\section{Indexing and populations}\label{sec:tf-indexing}

A theory can be \term{homogeneous} (one population, scalar fields) or
\term{heterogeneous} (multiple populations of different sizes, with vector or
matrix parameters). The action's \code{sum(... for i in pop)} comprehension you
saw above is how a field equation is written once and applied across a
population: \code{i} ranges over the population \code{pop}, and \code{x[i]},
\code{xt[i]} are the per-member fields. For the running example \code{pop} has
size~1, so the sum is a single term --- but the same syntax scales to genuine
networks.

Indexing is encoded in one of two styles, and the serializer round-trips both:

\begin{description}
  \item[New population style.] \code{.population('E', size=2)} declares a named
        population of size $N$; fields attach via \code{population='E'} and
        parameters via \code{indexed\_by=['E']} (vector) or
        \code{indexed\_by=['E','E']} (matrix). The running example uses the
        single-member form \code{.population('pop', size=1)}.
  \item[Legacy style.] \code{n\_populations=N} in the constructor, plus
        \code{indexed=True} (vector) or \code{indexed='matrix'} on parameters. The
        coloured-noise file in \S\ref{sec:tf-colored} uses
        \code{TemporalTheoryBuilder('OU Quartic Colored', n\_populations=0)} with
        \code{.physical\_field('x', indexed=True)}.
\end{description}

\begin{gotcha}
\code{indexed\_by} \emph{wins} over the legacy \code{indexed}/\code{type}. Both
\code{\_emit\_parameter} and \code{\_emit\_kernel} accept either; if a spec
somehow carries both, the new population-style \code{indexed\_by} is used and the
legacy \code{indexed=} is suppressed. Mixing the two styles in a single
hand-written file is best avoided --- pick one.
\end{gotcha}

\section{Colored noise and CGF terms}\label{sec:tf-colored}

White noise enters the action as the local term $-D\,\tilde x^{2}$ --- a single
noise vertex with no memory. \term{Colored} (finite-correlation-time) noise
instead enters as a \term{CGF term} (cumulant-generating-function term), declared
with \code{.declare\_cgf\_term(...)} and carrying an \code{order}, a
\code{coefficient}, and a temporal \code{kernel}. The coloured sibling of the
running model, \file{theories/ou\_quartic\_colored.theory.py}, is the canonical
example. Its \code{build()} reads (verbatim, \file{...:10--23}):

\begin{lstlisting}
def build():
    return (
        TemporalTheoryBuilder('OU Quartic Colored', n_populations=0)
        .physical_field('x', indexed=True, description='variable')
        .parameter('mu', default=0.1, domain='real')
        .parameter('eps', default=0.1, domain='positive')
        .parameter('D', default=1.0, domain='positive')
        .parameter('tauc', default=1.0, domain='positive')
        .declare_cgf_term('CXX', response_legs=['xt', 'xt'], order=2,
                          coefficient='2*D/tauc', kernel='exp(-abs(tau)/tauc)')
        .set_action_text('xt*((Dt+mu)*x + eps*x^3)')
        .equation(lhs='(Dt+mu)*x', rhs='-eps*x^3')
        .stability_analysis(True)
        .build()
    )
\end{lstlisting}

Two differences from the white-noise file are the whole point. First, the action
\code{'xt*((Dt+mu)*x + eps*x\^{}3)'} has \emph{no} $-D\,\tilde x^{2}$ term --- the
noise is no longer local. Second, that noise is supplied by the CGF row

\begin{lstlisting}
.declare_cgf_term('CXX', response_legs=['xt', 'xt'], order=2,
                  coefficient='2*D/tauc', kernel='exp(-abs(tau)/tauc)')
\end{lstlisting}

which encodes the Ornstein--Uhlenbeck noise spectrum
$\avg{\xi\,\xi} \propto (D/\tau_c)\,\ee^{-|\tau|/\tau_c}$. The \code{response\_legs}
list says the term couples two response fields ($\tilde x$ with $\tilde x$), the
\code{order} is the cumulant order of the term, \code{coefficient} its prefactor,
and \code{kernel} its time profile.

The handler \code{\_emit\_cgf\_term} (\file{pipeline/theory\_serialize.py:233})
renders two shapes. The multi-leg shape above is used when \code{response\_legs}
has more than one entry; a legacy single-leg shape
\code{.declare\_cgf\_term('name', 'mt', order=..., ...)} passes the single
response field as the second positional argument. \code{order} is forced to
\code{int}, and a per-row \code{markovianize=} is emitted only when it is an
explicit boolean.

\paragraph{Markovianization.} By default the pipeline \term{Markovianizes} a
coloured kernel --- it embeds the finite-correlation-time noise in an auxiliary
Markovian field so the system becomes local in time again (the embedding lives in
\file{pipeline/colored\_to\_markovian.py}). The builder-level
\code{.markovianize(False)} opts the whole theory out, and a per-term
\code{markovianize=True/False} overrides individual rows. The forward serializer
emits the builder-level \code{.markovianize(False)} only when the spec explicitly
turned it off (\file{pipeline/theory\_serialize.py:448--449}), so theories that
accept the default stay quiet.

\begin{note}
Chapter~\ref{ch:quickstart}'s note on ``what colour costs'' is exactly this
mechanism: the white-noise model has no $\tau_c$ to embed, hence no auxiliary
field and a small diagram count, while the coloured sibling adds the embedding.
The two files differ in precisely the action term and the one CGF row shown
above. The dedicated chapter on coloured and Markovian noise
(Chapter~\ref{ch:colored-markovian}) opens up the embedding itself.
\end{note}

\section{Spatial derivative vertices and the operator IR}\label{sec:tf-spatial}

For spatial theories the action carries the Laplacian $\nabla^{2}$ and, for
derivative-vertex models, gradient operators. There are two distinct authoring
styles, and the difference between them is not cosmetic --- it changes the physics
the loop integrals compute.

\begin{description}
  \item[Bare-multiplicative Laplacian (v1).] The inert symbol \code{Laplacian} is
        used multiplicatively, e.g. \code{(Dt + mu - D*Laplacian)*phi}. No
        operator IR is needed.
  \item[Operator IR (v2).] The action uses \emph{call syntax} --- \code{Dt(h)},
        \code{Lap(h)}, \code{Dx(h, 0)} --- for genuine derivative \emph{vertices}.
        A theory using call syntax \emph{must} carry \code{.operator\_ir()}.
\end{description}

\begin{defn}
The \term{operator IR} is the v2 intermediate representation in which derivative
operators are call-syntax nodes (\code{Dt(h)}, \code{Lap(h)}, \code{Dx(h, axis)})
rather than inert multiplicative symbols. Because a derivative \emph{vertex}
contributes a momentum-dependent \term{form factor} to a loop integral, the IR is
what lets the engine attach the right form factor to each vertex. The integer in
\code{Dx(h, 0)} is the \term{spatial axis index} (here $d=1$, so axis~0).
\end{defn}

Why the IR is load-bearing is best seen on KPZ. The Kardar--Parisi--Zhang
equation in \file{theories/kpz\_1d.theory.py} is
\[
  \partial_t h = -\mu h + D\nabla^{2}h + \tfrac{\lambda}{2}(\partial_x h)^{2} + \eta,
  \qquad \avg{\eta\,\eta} = 2T\,\delta(x-x')\,\delta(t-t'),
\]
written as

\begin{lstlisting}
.set_action_text(
    'ht*(Dt(h) + mu*h - D*Lap(h) - (lam/2)*Dx(h, 0)^2) - T*ht^2')
.operator_ir()
.boundary('infinite')
.initial('stationary')
\end{lstlisting}

The KPZ nonlinearity $\tfrac{\lambda}{2}(\partial_x h)^{2}$ is a \emph{per-leg}
gradient vertex, written \code{Dx(h, 0)\^{}2}. Contrast Model B's conserving flux
$g\nabla^{2}(\phi^{2})$, a \emph{composite} derivative vertex written
\code{g*Lap(phi\^{}2)}. These differ --- $(\partial_x h)^{2} \neq
\partial_x(h^{2})$ --- and they produce different momentum form factors in the
loop integral.

\begin{gotcha}
\code{.operator\_ir()} is \textbf{load-bearing and silent}. The comment at
\file{pipeline/theory\_serialize.py:452--457} is explicit: a re-saved
KPZ/Burgers/Model-B theory that loses the \code{.operator\_ir()} line will
``silently lose the \code{Dt()/Lap()/Dx()} lowering and compute the wrong
physics.'' The serializer therefore round-trips the toggle precisely, and emits it
only when on (so v1 bare-Laplacian theories stay quiet). \emph{If you hand-edit a
derivative-vertex theory, keep that line.}
\end{gotcha}

For spatial models you will also see \code{.boundary(mode, **params)} (e.g.
\code{.boundary('infinite')} or a periodic box with a \code{length=}) and
\code{.initial(mode, **params)} (e.g. \code{.initial('stationary')}), and for
coupled multi-field models the \code{.dyson\_order(N)} /
\code{.reference\_diffusion(D0)} policy pair. The spatial chapters of Part~V cover
what those declarations route to.

\section{The spec dict: the contract with the UI}\label{sec:tf-spec}

Everything above describes the \emph{file}. Internally, the UI form and the
serializer speak a flat dictionary --- the \term{spec dict}. It is produced by the
form's \code{\_collect()}, rendered to text by \code{render\_theory\_file}, and
reconstructed from text by \code{load\_spec\_from\_file}. Knowing its shape makes
both the forward and reverse paths legible, because every \code{\_emit\_*} helper
is really ``read this key, write that line,'' and every reverse-path branch is the
mirror image. The defaults below are exactly those the reverse path initialises
(\file{pipeline/theory\_serialize.py:601--637}):

\begin{center}
\begin{longtable}{@{}llp{0.50\textwidth}@{}}
\toprule
\textbf{Key} & \textbf{Type} & \textbf{Meaning} \\
\midrule
\endhead
\code{name} & str & Human theory name (also the slug source). \\
\code{description} & str & User prose (boilerplate-stripped on reload). \\
\code{populations} & list[dict] & \code{[\{'name','size','description'?\}, ...]}
  (heterogeneous style). \\
\code{n\_populations} & int & Legacy population count; synced to
  \code{len(populations)} when populations are present. \\
\code{response\_fields} & list[dict] & Explicitly-declared response fields;
  usually empty. \\
\code{physical\_fields} & list[dict] & Field records (name, indexing, natural
  name, latex, description, \code{spatial\_dim}). \\
\code{parameters} & list[dict] & Parameter records (name, default, indexing,
  domain, \dots). \\
\code{functions} & list[dict] & \code{define\_function} records (name, args,
  expression, \dots). \\
\code{kernels} & list[dict] & \code{define\_kernel} records (time/freq exprs,
  names, indexing). \\
\code{cgf\_terms} & list[dict] & Coloured-noise rows (legs, order, coefficient,
  kernel, markovianize). \\
\code{action\_text} & str & The \msrjd{} action expression (verbatim). \\
\code{mf\_equations} & list[dict] & Legacy \code{[\{'saddle','rhs'\}, ...]}. \\
\code{equations} & list[dict] & DAE form \code{[\{'lhs','rhs','population'\}, ...]}. \\
\code{stability\_analysis} & bool & Default \code{False}. \\
\code{markovianize\_default} & bool & Default \code{True}. \\
\code{operator\_ir} & bool & Default \code{False}. \textbf{Load-bearing} for
  derivative-vertex spatial theories. \\
\code{boundary} & dict & \code{\{'mode', **params\}}, e.g.
  \code{\{'mode':'infinite'\}}. \\
\code{initial} & dict & \code{\{'mode', **params\}}, e.g.
  \code{\{'mode':'stationary'\}}. \\
\code{dyson} & dict & \code{\{'mode':'fixed','order':N\}}. \\
\code{reference\_diffusion} & float & Scalar $D_{0}$ for the
  $\mathcal{D} = D_{0}\,I + \hat{\mathcal{D}}$ split. \\
\code{default\_fundamental} & dict & Emitted as \code{DEFAULT\_FUNDAMENTAL}. \\
\code{metadata} & dict & Emitted as \code{METADATA}. \\
\bottomrule
\end{longtable}
\end{center}

One transient key, \code{\_default\_spatial\_dim}, is set during reverse parsing
by a \code{.spatial\_dim(d)} call and is not normally re-emitted; it exists only
so a bulk \code{.spatial\_dim(d)} can back-fill fields that lacked a per-field
\code{spatial\_dim=}. We return to it in \S\ref{sec:tf-reverse}.

\section{Forward serialization: spec to file}\label{sec:tf-forward}

The forward heart is \code{render\_theory\_file(spec) -> str}
(\file{pipeline/theory\_serialize.py:310}), which takes a spec dict and returns
the full \file{.theory.py} source string. \code{save\_theory\_to\_file(spec, path)}
(\file{...:504}) is the thin wrapper that picks a filename and writes it. The
module itself is deliberately \emph{dependency-light}: it touches only the Python
standard library (\code{ast}, \code{re}, \code{textwrap}, \code{os}, \code{typing}).
No SageMath, no \code{ipywidgets}, no \code{numpy} --- those enter only in the
companion loader and the binary cache.

\code{render\_theory\_file} proceeds in a fixed order. Abbreviated, it:

\begin{enumerate}
  \item reads the spec fields with \code{.get(key, default)} so partial specs
        still render (\file{...:320--351});
  \item filters out saddle parameters --- those flagged \code{mean\_field} or named
        \code{<natural>star} (\file{...:334--339});
  \item picks the builder class from \code{spatial\_dim} (\S\ref{sec:tf-split},
        \file{...:370--372});
  \item emits the docstring header --- a \code{"<name> --- text-driven theory
        file."} line plus the optional description and the ``Generated by\dots\ /
        Loaded by\dots'' boilerplate the reverse path later strips
        (\file{...:354--363});
  \item emits \code{from pipeline.theory import <BuilderCls>} and
        \code{def build(): return (} (\file{...:376--380});
  \item emits the constructor: with \code{populations} present, a bare
        \code{BuilderCls('name')} followed by one \code{.population(...)} per
        population; otherwise \code{BuilderCls('name', n\_populations=N)}
        (\file{...:381--396});
  \item emits, in fixed order, response fields, physical fields, parameters,
        functions, kernels, CGF terms (\file{...:401--412});
  \item emits the action, then the equations (preferring \code{.equation(...)},
        else \code{.set\_mf\_equation(...)}) (\file{...:414--433});
  \item emits the non-default toggles --- stability, \code{markovianize(False)},
        \code{operator\_ir()}, boundary/initial, Dyson/reference-diffusion
        (\file{...:439--489});
  \item closes the chain with \code{.build()} and the \code{)}
        (\file{...:491--492}), then appends \code{DEFAULT\_FUNDAMENTAL = \{...\}}
        and \code{METADATA = \{...\}} (\file{...:495--498}).
\end{enumerate}

The low-level work is done by a handful of helpers worth knowing by name.
\code{\_py\_repr(value)} (\file{...:46}) pretty-prints a Python value as a
source-code literal: a 2D list (a matrix) becomes a multi-line block with one row
per line, a short 1D list stays on one line, a dict becomes one
\code{key: value} per line. This is what makes a matrix default print as a
readable block rather than one long line. \code{\_kw\_chain(*pairs)}
(\file{...:85}) joins \code{name=value} pairs but \emph{omits any pair whose value
is \code{None} or \code{''}} --- this is the mechanism behind the ``textually
quiet'' principle, so \code{.parameter('mu', default=1.0)} does not emit
\code{description=None}. \code{\_slugify(name)} (\file{...:96}) is the
filesystem-safe identifier
\code{re.sub(r'[\^{}A-Za-z0-9]+', '\_', name).strip('\_').lower()}.

\begin{gotcha}
In \code{\_py\_repr}, \code{bool} is checked \emph{before} \code{int}
(\file{...:54}). In Python \code{bool} is a subclass of \code{int}, so a naive
\code{isinstance(value, int)} test would catch \code{True}/\code{False} first and
print them as \code{1}/\code{0}. The order is deliberate; do not swap it.
\end{gotcha}

\begin{gotcha}
Whitespace stability depends on the dedent in \code{\_emit\_action}
(\file{...:281--289}). The action body is run through \code{textwrap.dedent}
\emph{before} being re-indented four spaces. Without that dedent, every
load$\to$save cycle would accumulate four more leading spaces on each line of a
multi-line action. Do not ``simplify'' it away.
\end{gotcha}

\code{save\_theory\_to\_file} normalises the path before writing
(\file{...:513--518}): if \code{path} is a directory, the filename is
\code{<slugify(name)>.theory.py} inside it; otherwise, if \code{path} does not
already end in \code{.theory.py}, the suffix is appended \emph{unless} it already
ends in \code{.py}.

\begin{gotcha}
The extension logic has a subtlety. A path ending in \code{.py} but \emph{not}
\code{.theory.py} is left untouched, so \code{save\_theory\_to\_file(spec, 'foo.py')}
writes \code{foo.py} --- which \code{dd.list\_theories()} would then \emph{not}
discover, because discovery requires the \code{.theory.py} suffix
(\file{notebooks/daedalus.py:63--65}). If you want a loadable theory, give it a
bare name or the full \code{.theory.py} suffix.
\end{gotcha}

\section{Reverse serialization: file to editable form}\label{sec:tf-reverse}

The reverse path is \code{load\_spec\_from\_file(path) -> dict}
(\file{pipeline/theory\_serialize.py:528}). Its job is to reconstruct the spec
dict from an existing \file{.theory.py} so the UI's ``Load'' button can
re-populate the form. The defining choice is that it reads the file's
\emph{abstract syntax tree} and never executes it.

\begin{defn}
The Python standard-library module \term{\code{ast}} (Abstract Syntax Tree)
parses source \emph{text} into a tree of node objects --- the program's grammar,
before execution. Inspecting this tree lets you learn what a program \emph{says}
without ever \emph{running} it, which is exactly what you want when reading a
configuration file you would rather not execute as a side effect of loading the
form.
\end{defn}

The whole reverse path is an \code{ast} walk. The load-bearing calls are
(\file{pipeline/theory\_serialize.py}, lines noted):

\begin{lstlisting}
import ast                                   # 545
tree = ast.parse(src)                        # 550  text -> AST
description = ast.get_docstring(tree)        # 553  module docstring
val = ast.literal_eval(node.value)           # 568  a *safe* eval of a literal
\end{lstlisting}

\code{ast.parse} turns the file into a module node (a malformed file raises
\code{SyntaxError} here). \code{ast.get\_docstring} pulls the leading
triple-quoted string. The critical one is \code{ast.literal\_eval}:

\begin{defn}
\term{\code{ast.literal\_eval}} evaluates a node \emph{only if it is a literal} ---
a number, string, list, dict, tuple, boolean, or \code{None}. It refuses function
calls and bare names, so it can \emph{never} run arbitrary code. This is why the
reverse path can safely recover a matrix default like \code{[[2,3],[1,3]]} or a
metadata tuple like \code{('dx', 1)} from the source text.
\end{defn}

The code wraps it in a local helper \code{\_lit(node, default=None)}
(\file{...:639}) that returns the fallback on
\code{ValueError}/\code{SyntaxError}/\code{TypeError}, and a sibling
\code{\_kwargs(call)} (\file{...:646}) that reads a call's keyword arguments into
a dict. With those in hand the walk proceeds:

\begin{enumerate}
  \item read the file and \code{ast.parse} it (\file{...:548--550});
  \item recover the description from the docstring, then strip the auto-header
        with \code{\_strip\_doc\_boilerplate} (\file{...:553--556});
  \item scan \code{tree.body} for the module-level \code{ast.Assign} nodes that
        target \code{DEFAULT\_FUNDAMENTAL} / \code{METADATA} and
        \code{literal\_eval} their values (\file{...:561--574});
  \item find the \code{build} function (an \code{ast.FunctionDef}), and its
        \code{ast.Return}; a missing \code{build()} or a return that is not a call
        raises \code{ValueError} (\file{...:577--585});
  \item \textbf{unwind the method chain}: starting at the return-call, step into
        \code{cur.func.value} while the node is an \code{ast.Call}, stopping at
        the constructor \code{ast.Name}; then \emph{reverse} so the constructor is
        first and methods follow in source order (\file{...:588--598});
  \item initialise the spec with all expected keys at their defaults, including
        the four toggle defaults that mirror the builder's own
        (\file{...:601--637});
  \item dispatch over the chain --- a long \code{if/elif} on the method name
        (\code{population}, \code{physical\_field}, \code{parameter},
        \code{define\_function}, \code{define\_kernel}, \code{declare\_cgf\_term},
        \code{set\_action\_text}, \code{set\_mf\_equation}, \code{equation}, the
        toggles, \code{boundary}/\code{initial}, \code{dyson}/\code{dyson\_order},
        \code{reference\_diffusion}, \code{build}) --- each branch reconstructing
        one spec sub-list (\file{...:649--865});
  \item sync \code{n\_populations} to \code{len(populations)} when the populations
        list is non-empty (\file{...:868--869}).
\end{enumerate}

A few dispatch branches deserve a sentence. The constructor branch accepts all
three class names (\file{...:651--652}). The \code{physical\_field} branch
\emph{inherits} a builder-level \code{\_default\_spatial\_dim} (set by an earlier
\code{.spatial\_dim(d)} call) when no per-field \code{spatial\_dim=} is given
(\file{...:682--683}) --- this is why the transient spec key from
\S\ref{sec:tf-spec} exists. The \code{declare\_cgf\_term} branch handles both the
legacy positional \code{response\_field} and the new \code{response\_legs} (which
it accepts as a list literal \emph{or} a comma-separated string). The toggle
branches use sensible no-argument defaults: \code{markovianize()} reads as
\code{True}, \code{operator\_ir()} reads as \code{True}.

Finally, \code{\_strip\_doc\_boilerplate(doc)} (\file{...:873}) removes the
auto-generated header so a reloaded theory shows only the user's own prose: it
drops the leading \code{"<name> --- text-driven theory file."} line (recognised by
its em-dash), stops at the ``Generated by\dots\ / Loaded by\dots'' lines, and trims
surrounding blanks. This is the same trimming \code{describe\_model} performed in
Chapter~\ref{ch:quickstart}.

\subsection{The round-trip guarantee and its limits}

The reverse path is an \emph{AST} round-trip, not a textual one, which fixes
exactly what survives:

\begin{itemize}
  \item String literals --- action text, equation right-hand sides, function and
        kernel expressions --- survive \emph{verbatim}.
  \item Python literals --- numeric defaults, \code{indexed\_by} lists, metadata
        tuples --- survive exactly via \code{literal\_eval}.
  \item Comments and exact formatting do \emph{not} survive.
\end{itemize}

\begin{gotcha}
\textbf{Non-literal arguments are lost on reload.} \code{\_lit} returns the
fallback for anything that is not \code{literal\_eval}-able. If a hand-written
theory passes a \emph{computed} expression as a builder keyword --- say
\code{default=2*0.4} instead of \code{default=0.8} --- then
\code{load\_spec\_from\_file} recovers \code{None} for it. Defaults are normally
plain literals, so this rarely bites, but a hand-authored file that puts logic in
its \code{build()} can carry meaning the AST round-trip silently flattens. The
loader (next section) \emph{executes} the file, so it sees the computed value; the
UI ``Load'' button, reading the AST, does not.
\end{gotcha}

\begin{gotcha}
\code{load\_spec\_from\_file} \emph{raises} on a malformed file: \code{ast.parse}
raises \code{SyntaxError}, and a missing \code{build()} or a \code{build()} that
does not return a call expression raises \code{ValueError}
(\file{...:581}, \file{...:585}). The UI must catch these.
\end{gotcha}

\section{Discovery and loading by name}

You already used the canonical loader in Chapter~\ref{ch:quickstart}. It lives in
\file{notebooks/daedalus.py} and is worth re-reading now that you know what it
loads. \code{daedalus.load\_theory(name)} (\file{notebooks/daedalus.py:68})
resolves \file{theories/<name>.theory.py}, imports it via \code{importlib.util}
--- \emph{executing} it --- and returns \code{(module.build(), module)}:

\begin{lstlisting}
spec = importlib.util.spec_from_file_location(f'theories.{name}', path)
mod = importlib.util.module_from_spec(spec)
spec.loader.exec_module(mod)
return mod.build(), mod
\end{lstlisting}

\code{spec\_from\_file\_location} builds an import ``spec'' for an arbitrary file
path (so the file need not be on \code{sys.path}); \code{module\_from\_spec}
creates an empty module object; \code{exec\_module} actually runs the file's
top-level code, defining \code{build}, \code{DEFAULT\_FUNDAMENTAL}, and
\code{METADATA}.

\begin{gotcha}
\textbf{The reverse path reads; the loader executes.} \code{load\_spec\_from\_file}
uses \code{ast} and never runs the file --- safe, but it cannot recover
non-literal values. \code{daedalus.load\_theory} (and the binary serializer's
\code{reload\_model}) \emph{do} execute the file. So a \file{.theory.py} runs real
Python at run time: a hand-authored theory may legitimately contain logic the
AST-only reverse path would flatten. If you put logic in \code{build()}, the
engine will honour it but the UI ``Load'' button may not faithfully reflect it.
\end{gotcha}

The companions are small. \code{list\_theories()} (\file{notebooks/daedalus.py:61})
lists every \file{theories/*.theory.py} and strips the suffix to give the loadable
name. \code{THEORIES\_DIR} (\file{...:55}) is \code{<repo\_root>/theories}, where
\code{repo\_root()} (\file{...:40}) walks up the directory tree until it finds a
\code{pipeline/} folder --- the same depth-robust walk the setup cell performed in
Chapter~\ref{ch:quickstart}.

\section{The binary cache (\texttt{saved\_theories/})}

We close where the chapter's opening note began: the \emph{other} thing called
serialization. When a theory is actually \emph{run}, the heavy symbolic objects
--- the propagator, the Taylor-expanded action, the per-diagram integrands --- are
pickled to disk under \file{saved\_theories/<slug>/} as SageMath \code{.sobj}
files plus a \code{manifest.json}. This cache is managed by
\file{pipeline/\_expand\_cache.py} and \file{msrjd/core/serialize.py} --- \emph{not}
by \code{theory\_serialize.py}.

\begin{defn}
\term{SageMath} is a large open-source computer-algebra system built on Python; it
provides exact symbolic objects (polynomials over chosen rings, symbolic matrices,
\dots). Its \code{save}/\code{load} functions pickle these objects to
\term{\code{.sobj}} files (the extension is appended automatically). You need Sage
to load a \code{.sobj}; you do \emph{not} need it to read a \file{.theory.py}.
\end{defn}

The on-disk \term{slug} is \code{model['name']} lowercased with non-alphanumerics
collapsed to underscores --- the \emph{same} pattern \code{\_slugify} uses for the
text-format filename (\file{pipeline/\_expand\_cache.py} uses an identical regex),
so the human-facing filename and the cache directory agree. Thus
\code{'OU Quartic (white noise)'} caches under
\file{saved\_theories/ou\_quartic\_white\_noise/}. The directory's contents:

\begin{center}
\begin{longtable}{@{}lll@{}}
\toprule
\textbf{File} & \textbf{Producer} & \textbf{Contents} \\
\midrule
\endhead
\code{manifest.json} & \code{\_expand\_cache.py} & index of cached artefacts
  (\code{key}, \code{stage}, \code{k}, \code{loop\_order}, \code{saved\_at}). \\
\code{propagator.sobj} & \code{\_expand\_cache.py} & Sage-pickled propagator
  (Taylor-order-independent). \\
\code{expand\_taylor<N>.sobj} & \code{\_expand\_cache.py} & Sage-pickled
  Taylor-expanded action at order $N$. \\
\code{unique\_typed\_mult\_v3\_...\_k<k>\_l<l>.sobj} &
  \code{\_expand\_cache.py} & Per-diagram integrand bundle, keyed by external legs,
  Taylor order, $k$, loop order. \\
\code{metadata.json} + \code{symbolic\_data.sobj} & \code{msrjd/core/serialize.py} &
  the older (separate) save format. \\
\bottomrule
\end{longtable}
\end{center}

A representative \code{manifest.json} entry (from
\file{saved\_theories/ou\_quartic\_white\_noise/}) shows how the cache key encodes
the request --- this is where the $k$ and $\ell$ from \code{METADATA}/\code{Config}
end up as file-name keys:

\begin{lstlisting}[style=console]
{ "key": "unique_typed_mult_v3_dx1_dx1_taylor6_k2_l2",
  "stage": "unique_typed_mult_v3_dx1_dx1_taylor6",
  "k": 2, "loop_order": 2,
  "saved_at": "2026-06-17T23:03:39.062450" }
\end{lstlisting}

\begin{gotcha}
\file{msrjd/core/serialize.py} defines its \emph{own} \code{save\_theory} /
\code{load\_theory} functions --- distinct from both \code{theory\_serialize.py}
and \code{daedalus.load\_theory}. They persist expanded symbolic objects, not
\file{.theory.py} text. Its \code{reload\_model} even uses raw \code{exec} to
re-import the model \code{.py}, so it runs arbitrary code. Three different
functions across the codebase are spelled ``\code{load\_theory}''; always check
which module you are in.
\end{gotcha}

\section{Authoring rules, in one place}

Most of the chapter's gotchas reduce to a short checklist for anyone hand-writing
or hand-editing a theory file:

\begin{enumerate}
  \item Write the action with \code{xt} for the response field and \code{\^{}} for
        powers; wrap per-population dynamics in \code{sum(... for i in pop)}.
  \item Declare the saddle with \code{.equation(lhs, rhs, population)} (preferred)
        or \code{.set\_mf\_equation('xstar', rhs)} (legacy). Never declare a
        \code{<natural>star} parameter --- the framework creates it.
  \item For a spatial model, give at least one physical field a non-zero
        \code{spatial\_dim} so the builder split picks
        \code{SpatialTheoryBuilder}; and if the action uses \code{Dt()/Lap()/Dx()}
        call syntax, \emph{keep} the \code{.operator\_ir()} line --- dropping it
        silently computes the wrong physics.
  \item Keep \code{recommended\_external\_fields} entries as \emph{tuples}.
  \item Prefer plain literals for builder keyword values; a computed expression
        works at run time but is lost when the UI reloads the file via the AST
        path.
  \item Remember the two serialization layers: a \file{.theory.py} is plain
        text you can read and diff; a \file{saved\_theories/<slug>/*.sobj} is a
        Sage pickle written when the model runs. Deleting the cache directory only
        forces a recompute; it never touches your source.
\end{enumerate}

With the model fully specified --- by hand or by form --- the next chapter opens up
the one feature this running example deliberately avoided: coloured noise and the
Markovian embedding that turns it back into a local-in-time field theory.

% ---------------------------------------------------------------------
%  Chapter 5 : Colored Noise and the Markovian Embedding
% ---------------------------------------------------------------------
\chapter{Colored Noise and the Markovian Embedding}\label{ch:colored-markovian}

Every model in the previous two chapters drove its fields with \emph{white}
noise --- a random force that is uncorrelated across any two distinct instants.
White noise is the easy case, and not by accident: the temporal integrators that
do the real work at the end of the pipeline have fast, fully \emph{analytic}
mode-sum routines that exist precisely because the noise two-point function is a
Dirac delta in time. The moment you ask for a noise with \emph{memory} --- a
finite correlation time $\tau_c$, a smooth kernel like $\ee^{-|\tau|/\tau_c}$ ---
those analytic routines no longer apply, and the pipeline is forced onto a
generic numerical quadrature that, at one loop and beyond, either crawls or
hangs.

This chapter is about the one preprocessor that makes colored noise tractable
anyway. It is a \term{builder-level rewrite}: before your model is ever compiled,
it silently replaces your colored-noise term with an equivalent \emph{white}-noise
problem by introducing one extra dynamical field --- an auxiliary
Ornstein--Uhlenbeck (OU) process --- per noisy field. The output statistics are
identical; the representation is larger but every integral in it is again the
fast white-noise kind. The whole subsystem lives in one file,
\file{pipeline/colored\_to\_markovian.py}, and this chapter walks it from the
physics down to the exact text strings it emits.

It sits in Part~II (\emph{Specifying a Model}) rather than in the
diagram-evaluation parts because it happens entirely at authoring time: it is a
transformation \emph{of the model spec}, run inside \code{TheoryBuilder.build()}
before any symbolic compilation. By the time the field-theory core (Chapter~6
onward) sees your model, the colored noise is gone and only white noise remains.

\section{Motivation: why colored noise breaks the fast integrators}

A noisy field $x(t)$ is driven by a random force $\xi(t)$. For a Gaussian noise,
that force is fully characterised by its mean (zero) and its two-time
autocovariance --- the \emph{second cumulant}:
\begin{equation}
  \avg{\xi(t)\,\xi(t')} \;=\; \kappa(t-t') \;=\; \kappa(\tau),
  \qquad \tau = t - t'.
  \label{eq:cm-secondcumulant}
\end{equation}
Two regimes matter here.

\begin{description}
  \item[White noise.] $\kappa(\tau) = 2D\,\delta(\tau)$. The force is
        uncorrelated across any two distinct instants. This is what the fast
        analytic integrators handle: a delta in time collapses one of the loop
        integrals exactly, leaving a closed-form mode sum.
  \item[Colored noise (single Lorentzian).] $\kappa(\tau) = c\,\ee^{-|\tau|/\tau_c}$.
        The force ``remembers'' itself for a correlation time $\tau_c$; as
        $\tau_c \to 0$ with $c/\tau_c$ held fixed it limits back to white. The
        name \term{Lorentzian} is because the Fourier transform of
        $\ee^{-|\tau|/\tau_c}$ is a Lorentzian, $\propto 1/(\omega^2 + 1/\tau_c^2)$.
\end{description}

The trouble is the smooth factor. \Daedalus's temporal integrator ---
\code{msrjd.integration.time\_domain.final\_integral} --- cannot absorb a smooth
$\ee^{-|\tau|/\tau_c}$ into its analytic mode-sum routines. The module docstring
states the pain exactly (\file{pipeline/colored\_to\_markovian.py:6--10}):

\begin{quote}\itshape
\code{msrjd.integration.time\_domain.final\_integral} cannot absorb a smooth
\code{exp(-|tau|/tauc)} factor coming from a CGF-tab \code{NoiseSourceType}
vertex into its analytic mode-sum integrators. At \code{max\_ell >= 1} the scipy
\code{nquad} fallback path either hangs or runs for orders of magnitude longer
than the equivalent white-noise theory.
\end{quote}

Here \code{nquad} is SciPy's nested numerical multi-dimensional quadrature --- the
slow generic fallback. And \code{max\_ell} is the loop-order cutoff $\ell$:
$\ell = 0$ is tree level, $\ell \ge 1$ adds the loop corrections. So the bottom
line is blunt: at $\ell \ge 1$ a colored theory is effectively unusable.

\begin{note}
This is purely a \emph{performance} cliff, not a correctness one. At tree level
($\ell = 0$) a colored theory will still run through the \code{nquad} fallback
and give the right answer. The embedding in this chapter exists so that the
\emph{loop} corrections become affordable --- it turns an intractable
quadrature into the same fast analytic mode sum every white-noise model uses.
\end{note}

\subsection{How a colored cumulant is declared}

In \Daedalus{} a colored second cumulant is declared as a \term{CGF row} ---
a cumulant-generating-functional term --- with \code{order=2}, a
\code{coefficient} text (the amplitude $c$), and a \code{kernel} text (the
$\tau$-shape). The canonical example, the OU process with a cubic restoring force
and colored driving, declares it like this
(\file{theories/ou\_quartic\_colored.theory.py:18}):

\begin{lstlisting}
.declare_cgf_term('CXX', response_legs=['xt', 'xt'], order=2,
                  coefficient='2*D/tauc', kernel='exp(-abs(tau)/tauc)')
\end{lstlisting}

so here $c = 2D/\tau_c$ and the kernel is $\ee^{-|\tau|/\tau_c}$ --- the full
second cumulant is $(2D/\tau_c)\,\ee^{-|\tau|/\tau_c}$. The response legs
\code{['xt', 'xt']} say the cumulant connects the response field \code{xt}
(the \msrjd{} conjugate of \code{x}) to itself: it is an \emph{auto}-cumulant
on a single field. The accompanying action text is the plain deterministic
dynamics,
\begin{lstlisting}
.set_action_text('xt*((Dt+mu)*x + eps*x^3)')
\end{lstlisting}
--- notice there is \emph{no} noise vertex in the action text. For colored noise
the noise lives entirely in the CGF row, not as a $-D\tilde x^2$ term in the
action (contrast the white-noise OU quartic of Chapter~2, whose action carried
\code{- D*xt[i]\^{}2} explicitly).

\begin{defn}
A \term{CGF row} is a declared contribution to the noise's cumulants, given as a
4-tuple \code{(coefficient, kernel, order, response\_legs)}. \code{order=2} is the
noise \emph{covariance}. The \code{coefficient} is the amplitude, the
\code{kernel} is the function of the lag $\tau$; \code{response\_legs} names which
response fields the cumulant connects. Colored noise sets a smooth kernel like
\code{exp(-abs(tau)/tauc)}; white noise sets \code{dirac\_delta(tau)}.
\end{defn}

\section{The fix in one idea: a colored noise is a filtered white noise}

There is a classic trick in stochastic dynamics. A colored noise whose
autocovariance is a single decaying exponential is \emph{exactly} the output of a
first-order linear filter --- an Ornstein--Uhlenbeck process --- driven by white
noise. So instead of feeding colored noise $\xi(t)$ directly into your equation,
you do two things:

\begin{enumerate}
  \item drive $\xi$ \emph{itself} with white noise through an OU relaxation
        equation, and
  \item feed the OU \emph{output} $\xi$ into your field as an ordinary linear
        coupling.
\end{enumerate}

Your field now sees a deterministic linear coupling to a new auxiliary field, and
the only stochastic forcing left anywhere in the system is white. The white-noise
analytic integrators apply again. The price is one extra field (and its \msrjd{}
response partner) per noisy field: a bigger theory, but a cheaper-to-integrate
one. The module docstring writes the two-line system out
(\file{pipeline/colored\_to\_markovian.py:13--23}):
\begin{align}
  \frac{\dd x}{\dd t}  \;&=\; f(x) \;+\; \xi,
  \label{eq:cm-toptline}\\[2pt]
  \frac{\dd \xi}{\dd t} \;&=\; -\frac{\xi}{\tau_c}
        \;+\; \sqrt{\tfrac{2c}{\tau_c}}\;\eta(t),
        \qquad \avg{\eta(t)\,\eta(t')} = \delta(t-t').
  \label{eq:cm-bottomline}
\end{align}

Read this as a coupled pair:

\begin{description}
  \item[Top line, \eqref{eq:cm-toptline}] is \emph{your} equation, except the
        colored force $\xi$ is no longer ``random by fiat'' --- it is now a
        genuine dynamical field.
  \item[Bottom line, \eqref{eq:cm-bottomline}] is the \emph{auxiliary} OU
        equation: $\xi$ relaxes toward zero with rate $1/\tau_c$, and is itself
        kicked by \emph{white} noise $\eta$ of unit strength, scaled by the
        amplitude $\sqrt{2c/\tau_c}$.
\end{description}

\begin{defn}
The \term{Ornstein--Uhlenbeck (OU) process} is the simplest mean-reverting linear
SDE, $\dd\xi/\dd t = -\xi/\tau_c + (\text{white})$. The \term{Markovian
embedding} is the act of enlarging the state space (adding auxiliary fields) so
that a non-Markovian, memory-carrying, colored process becomes Markovian
(memoryless, white-driven). The observable output field is unchanged; only the
representation grows.
\end{defn}

\subsection{Why the OU output is exactly the Lorentzian we wanted}

The reason the whole construction works is a single claim: the \emph{stationary}
autocovariance of the OU output $\xi$ defined by \eqref{eq:cm-bottomline} is
exactly the Lorentzian we set out to reproduce,
\begin{equation}
  \avg{\xi(t)\,\xi(t')} \;=\; c\,\ee^{-|t-t'|/\tau_c}.
  \label{eq:cm-ouvariance}
\end{equation}

\emph{Derivation.} The OU SDE $\dd\xi/\dd t = -\xi/\tau_c + \sigma\eta$, with
white $\eta$ ($\avg{\eta(t)\eta(t')} = \delta(t-t')$) and
$\sigma = \sqrt{2c/\tau_c}$, is linear, so its stationary two-point function is
known in closed form. The stationary \emph{variance} balances the injection rate
$\sigma^2$ against the dissipation rate $2/\tau_c$, giving $\sigma^2 \tau_c/2$;
the decay in the lag is governed by the relaxation rate $1/\tau_c$. Hence
\begin{equation}
  \avg{\xi\,\xi}(\tau)
    \;=\; \frac{\sigma^2 \tau_c}{2}\,\ee^{-|\tau|/\tau_c}
    \;=\; \frac{(2c/\tau_c)\,\tau_c}{2}\,\ee^{-|\tau|/\tau_c}
    \;=\; c\,\ee^{-|\tau|/\tau_c}.
  \label{eq:cm-ouvariance-derived}
\end{equation}
That last equality is exactly \eqref{eq:cm-ouvariance}. \hfill$\checkmark$

The derivation also pins the white drive amplitude. In \emph{variance} terms, the
strength of the white kick is $\sigma^2 = 2c/\tau_c$; \Daedalus{} declares white
noise as a CGF row with kernel $\delta(\tau)$, and the coefficient of such a row
is exactly the variance strength. So the auxiliary white-noise coefficient must be
\begin{equation}
  \text{aux white-noise coefficient} \;=\; \frac{2c}{\tau_c}.
  \label{eq:cm-auxdrive}
\end{equation}
This $2c/\tau_c$ is precisely what the detector computes and stores as the
\code{aux\_drive\_text}. The code builds it symbolically
(\file{pipeline/colored\_to\_markovian.py:183}):

\begin{lstlisting}
aux_drive_expr = (SR(2) * c_expr / tauc_expr).simplify_full()
\end{lstlisting}

For the worked example $c = 2D/\tau_c$, this gives
$2\cdot(2D/\tau_c)/\tau_c = 4D/\tau_c^2$, which the test suite pins down
(\file{tests/test\_markovianize.py:73}):
\code{assert info['aux\_drive\_text'] == '4*D/tauc\^{}2'}.

\section{From SDE to MSR-JD action: the three structural pieces}

\Daedalus{} works in the \msrjd{} (Martin--Siggia--Rose / Janssen--De~Dominicis)
path integral. Each physical field $x$ gets a conjugate \term{response field}
$\tilde x$ (written \code{xt} in code), and a linear equation of motion
$(\Dt + \mu)x = (\text{forcing})$ appears in the action as the bilinear term
$\tilde x\,[(\Dt + \mu)x - (\text{forcing})]$. The preprocessor's edits are
nothing more than the \msrjd{} encoding of the two SDE lines
\eqref{eq:cm-toptline}--\eqref{eq:cm-bottomline}. There are exactly three
structural pieces, plus two text edits to the action.

\paragraph{Piece 1 --- the auxiliary field.} For each noisy field, inject one new
physical field $\xi$ (code name \code{xi}). \Daedalus's authoring API auto-creates
its conjugate response \code{xit} and its saddle parameter \code{xistar}, just as
for any user-declared field.

\paragraph{Piece 2 --- the auxiliary OU drift equation.} The OU line
\eqref{eq:cm-bottomline}, stripped of its white drive, is the deterministic
operator $(\Dt + 1/\tau_c)\xi = 0$. The code injects this equation directly. In
the action it becomes a kinetic term contracted with the aux response field
\code{xit} (\file{pipeline/colored\_to\_markovian.py:431--435}):
\begin{lstlisting}[style=console]
    + xit*((Dt + 1/tauc)*xi)
\end{lstlisting}

\paragraph{Piece 3 --- the auxiliary white CGF row.} The white drive
$\sqrt{2c/\tau_c}\,\eta$ does \emph{not} appear in the action text. It is
represented by a \emph{new white CGF row} on the aux response legs --- the
$2c/\tau_c$ coefficient of \eqref{eq:cm-auxdrive} with kernel $\delta(\tau)$ ---
which is how \Daedalus{} injects Gaussian white noise everywhere. This is the
\code{new\_cgf\_rows} construction at
\file{pipeline/colored\_to\_markovian.py:397--405}.

That leaves one text edit to wire the auxiliary output into your dynamics:

\paragraph{Edit --- couple the user's field to the aux output.} The colored
force $\xi$ enters the SDE as $\dd x/\dd t = f(x) + \xi$, i.e. moving everything to
one side, $(\Dt + \mu)x + \varepsilon x^3 - \xi = 0$. So the code inserts
\code{- xi} inside your parenthesised dynamics block
(\file{pipeline/colored\_to\_markovian.py:472}):
\begin{lstlisting}[style=console]
    xt*((Dt+mu)*x + eps*x^3)   ->   xt*((Dt+mu)*x + eps*x^3 - xi)
\end{lstlisting}

Putting the three pieces together, the auxiliary field, the aux OU drift
equation, and the aux white CGF row jointly reproduce the bottom SDE line
$\dd\xi/\dd t = -\xi/\tau_c + \sqrt{2c/\tau_c}\,\eta$; the \code{- xi} coupling
plus the appended $\tilde\xi\,(\Dt + 1/\tau_c)\,\xi$ kinetic term wire $\xi$ into
the top line. The model that emerges is a pure white-noise theory.

\begin{note}
Note the asymmetry between the two text edits. The \code{- xi} \emph{coupling} is
inserted \emph{inside} an existing parenthesised block (it must land next to the
right field's dynamics), whereas the aux \emph{kinetic} term is simply
\emph{appended} to the end of the action string with a leading \code{+}. Appending
is safe because the kinetic term is self-contained; the coupling insertion is the
delicate one, and \S\ref{sec:cm-injectcoupling} is where its fragility lives.
\end{note}

\section{The cross-correlated multi-field case}

The embedding generalises cleanly to several fields whose colored forces are
\emph{cross}-correlated. For two fields $x, y$ with
\begin{align}
  \avg{\xi_x\,\xi_x}(\tau) &= c_{xx}\,\ee^{-|\tau|/\tau_c}, &
  \avg{\xi_y\,\xi_y}(\tau) &= c_{yy}\,\ee^{-|\tau|/\tau_c}, &
  \avg{\xi_x\,\xi_y}(\tau) &= c_{xy}\,\ee^{-|\tau|/\tau_c},
  \label{eq:cm-cross}
\end{align}
the embedding gives $x$ its own auxiliary OU process $\xi_x$ (code \code{xi}) and
$y$ its own $\xi_y$ (code \code{yi}), each driven by white noise, plus a single
\emph{cross} white CGF row between the two aux response legs. The shared
relaxation rate $1/\tau_c$ is what makes this work: because all three covariances
in \eqref{eq:cm-cross} decay at the same rate, the aux white covariance matrix
between $(\eta_x, \eta_y)$ factorises cleanly and reproduces the desired colored
cross-covariance.

The shipped two-field example
(\file{theories/ou\_quartic\_two\_dim\_color\_corr.theory.py:25--27}) declares
exactly this:
\begin{lstlisting}
.declare_cgf_term('Cxx', response_legs=['xt','xt'], order=2,
                  coefficient='D1/tauc',           kernel='exp(-abs(tau)/tauc)')
.declare_cgf_term('Cyy', response_legs=['yt','yt'], order=2,
                  coefficient='D2/tauc',           kernel='exp(-abs(tau)/tauc)')
.declare_cgf_term('Cxy', response_legs=['xt','yt'], order=2,
                  coefficient='rho*sqrt(D1*D2)/tauc', kernel='exp(-abs(tau)/tauc)')
\end{lstlisting}
The three colored rows \code{Cxx}, \code{Cyy}, \code{Cxy} become three white
aux rows: two auto rows on \code{[xit,xit]} and \code{[yit,yit]}, plus one cross
row on \code{[xit,yit]}. The docstring lays this out at
\file{pipeline/colored\_to\_markovian.py:33--39}.

\section{The kernel matcher: \texttt{detect\_lorentzian}}\label{sec:cm-detect}

Everything above is the physics. The rest of the chapter is the code that
implements it, starting with the function that decides whether a given CGF row
\emph{is} a single Lorentzian and, if so, extracts $\tau_c$, the amplitude $c$,
and the aux drive $2c/\tau_c$. That function is \code{detect\_lorentzian}
(\file{pipeline/colored\_to\_markovian.py:79}):

\begin{lstlisting}
def detect_lorentzian(kernel_text: str,
                      coefficient_text: str,
                      declared_params: list[str]) -> Optional[dict]:
\end{lstlisting}

It takes the row's \code{kernel} string (e.g. \code{'exp(-abs(tau)/tauc)'}), its
\code{coefficient} string (e.g. \code{'2*D/tauc'}), and the list of declared
parameter names (so it can build a namespace in which \code{tauc}, \code{D}, etc.
are known symbols). It returns a dict describing the embedding on a match, or
\code{None} if the kernel does not cleanly match the single-Lorentzian template.

This is the one part of the subsystem that does real symbolic algebra, so before
reading the steps we pause on the single library it leans on.

\begin{defn}
\term{SageMath} is a large open-source mathematics system built on top of Python;
it bundles many math libraries behind one Python-friendly API. The piece this
module uses is Sage's \term{Symbolic Ring}, abbreviated \code{SR} --- a
``calculator'' that works with algebraic expressions containing \emph{named
variables} (like \code{tauc}, \code{D}, \code{tau}) rather than just numbers. An
element of \code{SR} is a symbolic expression, e.g. \code{exp(-abs(tau)/tauc)},
that you can differentiate, simplify, substitute into, and pretty-print
symbolically. This is the same \emph{idea} as \code{sympy} (which \Daedalus{}
uses elsewhere), but Sage's \code{SR} is a separate, more powerful engine; the two
are \emph{not} interchangeable objects.
\end{defn}

The import line is (\file{pipeline/colored\_to\_markovian.py:70--73}):

\begin{lstlisting}
from sage.all import (
    SR, sage_eval, exp, abs_symbolic, simplify,
    sqrt as sage_sqrt,
)
\end{lstlisting}

What each name is, and exactly how this file uses it:

\begin{description}
  \item[\code{SR}] the Symbolic Ring object. \code{SR.var('name')} creates (and
        globally registers) a symbolic variable; \code{SR(x)} \emph{coerces} an
        arbitrary object $x$ into a symbolic expression so that symbolic methods
        can be called on it. Both forms appear throughout the file.
  \item[\code{sage\_eval}] Sage's \emph{safe, namespaced} string-to-expression
        evaluator. You hand it a string of Sage syntax plus a dictionary mapping
        names to Sage objects, and it parses and evaluates the string \emph{in
        that namespace only}. It is like Python's built-in \code{eval()} but tuned
        to Sage syntax (\code{\^{}} means power) and scoped to the dict you pass.
  \item[\code{exp}] Sage's symbolic exponential. It plays two roles: it is put
        \emph{into} the evaluation namespace so the kernel string's \code{exp(...)}
        resolves to it, and it is \emph{compared against} to recognise the kernel's
        top-level operation.
  \item[\code{abs\_symbolic}] Sage's symbolic absolute value. Plain Python
        \code{abs} would try to evaluate numerically; \code{abs\_symbolic} keeps
        $|\tau|$ as an \emph{unevaluated} symbolic object so the matcher can reason
        about it.
  \item[\code{simplify}, \code{sqrt as sage\_sqrt}] imported but \emph{not} called
        directly in this file. The code uses the more aggressive \emph{method}
        \code{.simplify\_full()} instead of the free \code{simplify}, and the
        \code{sqrt} appearing in coefficient strings is resolved by
        \code{sage\_eval} against Sage's default namespace rather than via this
        alias.
\end{description}

The matcher also calls a handful of methods on \code{SR} expressions:
\code{.operator()} returns the outermost operation node (for a Lorentzian that
must be \code{exp}); \code{.operands()} returns the arguments of that operation
($\exp(z)$ has exactly one operand, $z$); \code{.simplify\_full()} returns a
maximally-simplified copy; and \code{.has(sub)} is a boolean test for whether an
expression contains a given sub-expression anywhere in its tree.

\subsection{Step by step}

\begin{enumerate}
  \item \textbf{Empty guard} (\file{pipeline/colored\_to\_markovian.py:110}). If
        \code{kernel\_text} is falsy, return \code{None}.

  \item \textbf{Build the evaluation namespace}
        (\file{pipeline/colored\_to\_markovian.py:128--135}). Create a
        \emph{fresh local} symbol \code{tau\_sym = SR.var('\_lorentz\_tau')} but
        bind it to the \emph{name} \code{'tau'} in the namespace, so the kernel
        string's \code{tau} resolves to it. Add \code{exp} and \code{abs ->
        abs\_symbolic}, then add every declared parameter as a generic
        \code{SR.var(pname)} via \code{setdefault} (so a parameter the user
        already declared wins, and any other name is still defined):
\begin{lstlisting}
tau_sym = SR.var('_lorentz_tau')
eval_ns: dict = {
    'tau': tau_sym,
    'exp': exp,
    'abs': abs_symbolic,
}
for pname in declared_params:
    eval_ns.setdefault(pname, SR.var(pname))
\end{lstlisting}

  \item \textbf{Parse the kernel}
        (\file{pipeline/colored\_to\_markovian.py:137--141}).
        \code{K = sage\_eval(kernel\_text, eval\_ns)} inside a \code{try/except}
        that returns \code{None} on any parse error, then \code{K = SR(K)} to
        guarantee a symbolic object.

  \item \textbf{Top-level must be \texttt{exp}}
        (\file{pipeline/colored\_to\_markovian.py:143--148}). If
        \code{K.operator() != exp}, return \code{None}. Then read
        \code{operands = K.operands()}, require exactly one, take
        \code{arg = SR(operands[0])} (the exponent), and build
        \code{abs\_tau = abs\_symbolic(tau\_sym)}.

  \item \textbf{Decompose} $\mathtt{arg} = -|\tau|/\tau_c$
        (\file{pipeline/colored\_to\_markovian.py:154--161}). Compute
        \code{r = (arg / abs\_tau).simplify\_full()}. For a true Lorentzian, $r$
        should be $-1/\tau_c$ --- which is to say \emph{abs-free} and
        \emph{tau-free}. The code rejects (\code{return None}) if
        \code{r.has(abs\_tau)} or \code{r.has(tau\_sym)}. This single test is what
        distinguishes the three cases below.

  \item \textbf{Compute and validate} $\tau_c$
        (\file{pipeline/colored\_to\_markovian.py:163--170}).
        \code{tauc\_expr = (SR(-1) / r).simplify\_full()}, then reject if it still
        references \code{tau\_sym} or \code{abs\_tau}.

  \item \textbf{Parse the amplitude}
        (\file{pipeline/colored\_to\_markovian.py:174--180}).
        \code{c\_expr = sage\_eval(coefficient\_text or '0', eval\_ns)}, coerced
        and \code{.simplify\_full()}-ed; \code{return None} on failure. The
        \code{or '0'} guards an empty coefficient string. This is the $c$ in
        $c\,\ee^{-|\tau|/\tau_c}$.

  \item \textbf{Compute the aux drive}
        (\file{pipeline/colored\_to\_markovian.py:183}). The
        \code{aux\_drive\_expr = (SR(2) * c\_expr / tauc\_expr).simplify\_full()}
        from \eqref{eq:cm-auxdrive}.

  \item \textbf{Return the embedding dict}
        (\file{pipeline/colored\_to\_markovian.py:185--192}) with both the SR
        expressions and their string renderings (via \code{\_sr\_to\_text},
        \S\ref{sec:cm-srtotext}): \code{tauc\_expr}, \code{tauc\_text},
        \code{amplitude\_expr}, \code{amplitude\_text}, \code{aux\_drive\_expr},
        \code{aux\_drive\_text}.
\end{enumerate}

The decomposition in step~5 is the whole matcher in miniature. Three kernels and
what $r = \mathtt{arg}/|\tau|$ becomes for each:

\begin{center}
\begin{tabular}{@{}lll@{}}
\toprule
kernel & $r = \mathtt{arg}/|\tau|$ & verdict\\
\midrule
$\ee^{-|\tau|/\tau_c}$  & $-1/\tau_c$ (abs-free, tau-free)        & \textbf{match} \\
$\ee^{-\tau/\tau_c}$    & $-\tau/(\tau_c\,|\tau|)$ (still has $\tau$) & reject \\
$\ee^{-\tau^2/(\cdots)}$ & still has $\tau$                       & reject \\
\bottomrule
\end{tabular}
\end{center}

The middle row is the physically important rejection: a bare exponential
$\ee^{-\tau/\tau_c}$ without the absolute value is \emph{asymmetric} in $\tau$,
and a true autocovariance must be symmetric (it depends on $\tau = t-t'$ only
through $|\tau|$). The matcher correctly refuses it. The test suite checks all
three verdicts directly, including the Gaussian and the bare-exponential
rejections (\file{tests/test\_markovianize.py}, the
\code{test\_lorentzian\_kernel\_detection} case).

\begin{gotcha}
\textbf{The \texttt{tau} pollution guard.} Step~2 deliberately mints a fresh
\code{SR.var('\_lorentz\_tau')} instead of reusing the global \code{tau}. The long
comment at \file{pipeline/colored\_to\_markovian.py:116--127} explains why. Other
parts of the pipeline routinely register a global \code{tau} symbol carrying a
\code{domain='positive'} assumption (for instance an exponential-synapse time
constant). If $\tau > 0$ is assumed globally, then \code{abs(tau)} collapses to
\code{tau}, and the matcher would \emph{wrongly accept} an asymmetric kernel like
$\ee^{-\tau/\tau_c}$ --- precisely the case it is supposed to reject. Using a name
the caller never declares positive keeps the matcher robust against
test-ordering pollution. This is subtle and order-dependent: a regression would
only surface if a positive \code{tau} had been registered earlier in the same
Sage session.
\end{gotcha}

\begin{gotcha}
\textbf{Positivity of $\tau_c$ is documented but not enforced.} The docstring
says $\tau_c$ ``must be POSITIVE for a physical Lorentzian'' and that the matcher
requires the decay rate to ``evaluate to a strictly positive expression''
(\file{pipeline/colored\_to\_markovian.py:91--93, 163--164}), but the code only
ever checks that $\tau_c$ is abs-free and tau-free --- never that it is
numerically positive. A symbolic $\tau_c$ that happens to be negative for some
parameter values would still be accepted. Keep your $\tau_c$ a genuinely positive
parameter.
\end{gotcha}

\section{The builder transform: \texttt{markovianize\_spec}}\label{sec:cm-transform}

The detector is called from one place: \code{markovianize\_spec}
(\file{pipeline/colored\_to\_markovian.py:198}), the entry point that actually
rewrites the builder. Its signature is deceptively simple:

\begin{lstlisting}
def markovianize_spec(builder) -> None:
\end{lstlisting}

It takes a \code{TheoryBuilder} (in practice a \code{TemporalTheoryBuilder}) and
returns \emph{nothing}. All of its effects are in-place mutations: it reads
\code{builder.\_cgf\_terms}, \code{builder.\_markovianize\_default},
\code{builder.parameters}, \code{builder.physical\_fields}, and
\code{builder.\_action\_text}, and it writes them back --- crucially, by calling
the very same public authoring methods a user would call
(\code{builder.physical\_field(...)}, \code{builder.equation(...)}) plus direct
overwrites of \code{\_cgf\_terms} and \code{\_action\_text}.

\begin{note}
\textbf{Idempotency.} Re-running \code{markovianize\_spec} on an
already-markovianized builder is a no-op: no remaining row matches, because the
injected aux rows carry \code{markovianize=False} and the original colored rows
are gone. This safety property rests entirely on that one flag
(\file{pipeline/colored\_to\_markovian.py:404}); if it were dropped, a second
\code{build()} pass could try to re-embed the white aux rows.
\end{note}

\subsection{Step by step}

\begin{enumerate}
  \item \textbf{Snapshot \& empty guard}
        (\file{pipeline/colored\_to\_markovian.py:223--225}). Copy the rows;
        return immediately if there are none.

  \item \textbf{Builder-level gating}
        (\file{pipeline/colored\_to\_markovian.py:227--230}). Read
        \code{builder\_default = \_markovianize\_default} (default \code{True}).
        Compute \code{any\_explicit\_on = any(t.get('markovianize') is True ...)}.
        If the default is OFF \emph{and} no row explicitly opts in, return. This
        is the \code{.markovianize(False)} opt-out path.

  \item \textbf{Collect declared parameter names}
        (\file{pipeline/colored\_to\_markovian.py:232}) for the matcher.

  \item \textbf{Classify each row}
        (\file{pipeline/colored\_to\_markovian.py:240--276}) into
        \code{matched\_rows} / \code{passthrough\_rows}:
        \begin{itemize}
          \item \code{markovianize is False} $\to$ passthrough.
          \item \code{order != 2} $\to$ passthrough (the colored embedding is a
                \emph{second}-cumulant story).
          \item otherwise run \code{detect\_lorentzian(...)}.
          \item on \textbf{no match}: if the row had \code{markovianize is True}
                (an \emph{explicit} assertion), \textbf{raise \code{ValueError}}
                telling the user their kernel doesn't match the template;
                otherwise pass through silently.
          \item on \textbf{match}: append \code{\{'term': term, 'info': info\}} to
                \code{matched\_rows}.
        \end{itemize}

  \item \textbf{Early out}
        (\file{pipeline/colored\_to\_markovian.py:278--279}). If nothing matched,
        return.

  \item \textbf{Gather response-leg fields}
        (\file{pipeline/colored\_to\_markovian.py:292--298}). Read each matched
        row's \code{response\_legs}; if absent, broadcast its
        \code{response\_field} to a 2-list; collect all leg names into
        \code{source\_field\_names}.

  \item \textbf{Invert the response$\to$natural map}
        (\file{pipeline/colored\_to\_markovian.py:303--317}). For each response
        field name \code{rname} (e.g. \code{'xt'}), find the physical field whose
        auto-response equals it by testing \code{f'\{nat\}t' == rname} (recall
        \code{physical\_field('x')} auto-creates the response \code{xt}); store
        \code{response\_to\_natural[rname] = nat} (e.g. \code{'x'}). If none is
        found, store \code{None} --- that row will fall back to scipy because it
        cannot be embedded.

  \item \textbf{Single-$\tau_c$ enforcement}
        (\file{pipeline/colored\_to\_markovian.py:323--334}). Collect the set of
        distinct \code{tauc\_text} strings across matched rows. If there is more
        than one, \textbf{raise \code{NotImplementedError}} --- v1 supports a
        single shared $\tau_c$. Else take the one $\tau_c$.

  \item \textbf{Assign aux field names}
        (\file{pipeline/colored\_to\_markovian.py:339--352}). For each source
        natural name, in sorted order for determinism, call
        \code{\_pick\_aux\_field\_name(...)} (\S\ref{sec:cm-pickname}) to get a
        collision-safe aux name (\code{x -> xi}), record it, and \emph{reserve} it
        against later collisions in the same pass. If the map ends up empty,
        return.

  \item \textbf{Inject aux physical fields}
        (\file{pipeline/colored\_to\_markovian.py:355--360}). For each
        \code{(source\_nat, aux\_nat)}, call
        \code{builder.physical\_field(aux\_nat, indexed=True, description=...)}.
        This is the standard authoring method --- it auto-creates the conjugate
        response \code{<aux\_nat>t} and the saddle parameter
        \code{<aux\_nat>star}.

  \item \textbf{Inject aux OU drift equations}
        (\file{pipeline/colored\_to\_markovian.py:363--367}). For each aux field,
        \code{builder.equation(lhs=f'(Dt + 1/\{tauc\_text\}) * \{aux\_nat\}',
        rhs='0')} --- the deterministic part of the OU process,
        $(\Dt + 1/\tau_c)\xi = 0$.

  \item \textbf{Build replacement white CGF rows}
        (\file{pipeline/colored\_to\_markovian.py:370--405}). For each matched
        row, recover its two response legs, map them to natural names
        \code{nat\_a}, \code{nat\_b}, derive the aux response names
        \code{aux\_resp\_a}, \code{aux\_resp\_b}, and emit one white row:
\begin{lstlisting}
{
    'name':           f'{term["name"]}_markov_aux',
    'response_field': None,
    'response_legs':  [aux_resp_a, aux_resp_b],
    'order':          order,
    'coefficient':    info['aux_drive_text'],   # 2c / tauc
    'kernel':         'dirac_delta(tau)',        # white
    'markovianize':   False,                     # don't re-process
}
\end{lstlisting}
        For an auto-cumulant \code{[xt, xt]} both aux legs are \code{xit} (one
        row). For a cross-cumulant \code{[xt, yt]} the legs are \code{[xit, yit]}
        (the cross row); the auto rows \code{Cxx}, \code{Cyy} separately produce
        \code{[xit,xit]} and \code{[yit,yit]}.

  \item \textbf{Swap the CGF list}
        (\file{pipeline/colored\_to\_markovian.py:408}).
        \code{builder.\_cgf\_terms = passthrough\_rows + new\_cgf\_rows} --- the
        colored rows are gone, replaced by white aux rows; untouched rows are
        preserved.

  \item \textbf{Augment the action text}
        (\file{pipeline/colored\_to\_markovian.py:424--436}). For each source
        field, call \code{\_inject\_aux\_coupling(action\_text, resp\_field,
        aux\_nat)} with \code{resp\_field = f'\{source\_nat\}t'} to insert
        \code{- xi} inside the user's \code{xt*(...)} block. Then for each aux
        field, append \code{+ \{aux\_resp\}*((Dt + 1/\{tauc\_text\})*\{aux\_nat\})}
        --- the OU kinetic term. Write back to \code{builder.\_action\_text}.

  \item \textbf{Record diagnostics}
        (\file{pipeline/colored\_to\_markovian.py:439--443}). Set
        \code{builder.\_markovianize\_applied} to a small dict recording the aux
        field names, the matched row names, and the shared \code{tauc\_text}.
\end{enumerate}

\subsection{Where it is wired into \texttt{build()}}

The transform is invoked from exactly one place,
\file{pipeline/theory.py:1735--1737}, inside \code{build()}:

\begin{lstlisting}
if self._cgf_terms:
    from pipeline.colored_to_markovian import markovianize_spec
    markovianize_spec(self)
\end{lstlisting}

This call is deliberately placed \emph{before} \code{self.\_compile\_text\_declarations()}
(\file{pipeline/theory.py:1742}). Ordering is the whole point: by the time the
downstream compiler turns your text into lambdas, it only ever sees the augmented,
white-noise spec. The colored row never reaches the compiler.

\begin{note}
The two gating knobs you can reach as a user are: \code{.markovianize(False)} on
the builder (turns the whole transform OFF), and \code{markovianize=False} (or
\code{=True}) as a keyword on an individual \code{declare\_cgf\_term} call.
\code{False} skips a row; \code{True} asserts ``this IS a Lorentzian --- fail if
you disagree.'' The default (\code{None}/\code{'auto'}) means ``match if you can,
otherwise pass through silently.'' The opt-out path is exercised directly by the
test suite, which checks that \code{.markovianize(False)} leaves the original
colored row in place and injects no auxiliary field
(\file{tests/test\_markovianize.py}, the
\code{test\_markovianize\_opt\_out\_keeps\_colored\_row} case).
\end{note}

\section{The two helpers}

\subsection{\texttt{\_pick\_aux\_field\_name} --- collision-safe naming}
\label{sec:cm-pickname}

\code{\_pick\_aux\_field\_name(source\_nat, taken, additional\_taken=())}
(\file{pipeline/colored\_to\_markovian.py:449}) picks a name for the auxiliary
field that will not clash with anything already in the model. It unions
\code{taken} with \code{additional\_taken}, then tries the natural candidate
\code{f'\{source\_nat\}i'} (so \code{x -> xi}, \code{y -> yi}); if that is taken it
tries \code{f'xi\_\{source\_nat\}'}; and failing that it loops
\code{f'xi\_\{source\_nat\}\_\{n\}'} for $n = 2, 3, \dots$, which always terminates
because the integer suffix grows without bound.

\subsection{\texttt{\_inject\_aux\_coupling} --- action-text surgery}
\label{sec:cm-injectcoupling}

\code{\_inject\_aux\_coupling(action\_text, resp\_field, aux\_natural)}
(\file{pipeline/colored\_to\_markovian.py:472}) is the one piece of literal
string editing. It inserts \code{- <aux>} inside the parenthesised block that
follows \code{<resp\_field>*(} in the action text, or returns the text unchanged
if that pattern is not found. It is the only place the standard-library
\code{re} (regular expressions) module is used.

\begin{defn}
\code{re} is Python's built-in regular-expression module: it matches text
patterns. A \term{regular expression} is a small pattern language --- \code{\textbackslash{}b}
means ``word boundary,'' \code{\textbackslash{}s*} means ``zero or more whitespace characters,''
and \code{re.escape(s)} quotes any special characters in \code{s} so they match
literally.
\end{defn}

The function works in three moves
(\file{pipeline/colored\_to\_markovian.py:486--508}):
\begin{enumerate}
  \item Compile the pattern (word boundary, the field token, optional whitespace,
        \code{*}, optional whitespace, \code{(}), with \code{re.escape} quoting the
        field name:
\begin{lstlisting}
import re
pattern = re.compile(
    rf'\b{re.escape(resp_field)}\s*\*\s*\(', flags=0)
m = pattern.search(action_text)
\end{lstlisting}
  \item \code{.search(...)} for the first occurrence; if there is none, return the
        text unchanged.
  \item From the match end, \code{start = m.end() - 1} is the index of the
        \code{(}. Walk forward with a \code{depth} counter --- \code{+1} on
        \code{(}, \code{-1} on \code{)} --- until \code{depth} returns to $0$; that
        is the \emph{matching} close paren. Insert \code{f' - \{aux\_natural\}'}
        just before it.
\end{enumerate}
Regex only \emph{locates} the insertion point; the brace matching is a manual
depth counter, because regular expressions cannot reliably match balanced
parentheses. The function edits only the \emph{first} \code{<resp\_field>*(}
occurrence per call; \code{markovianize\_spec} calls it once per source field.

\begin{gotcha}
\textbf{Action-text surgery is fragile by design.} \code{\_inject\_aux\_coupling}
only handles the canonical \code{<rt>*(...)} layout. If your action is written in
a more baroque form, the insertion silently no-ops and the \code{- xi} coupling is
simply \emph{missing} --- you get a silently \emph{wrong} theory, not an error.
There is no validation that the coupling actually landed. This is the most
surprising failure mode in the subsystem. The remedy the docstring prescribes
(\file{pipeline/colored\_to\_markovian.py:418--423, 481--485}) is to opt out with
\code{.markovianize(False)} and hand-roll the embedding yourself.
\end{gotcha}

\subsection{\texttt{\_sr\_to\_text} --- symbolic back to re-parseable string}
\label{sec:cm-srtotext}

\code{\_sr\_to\_text(expr)} (\file{pipeline/colored\_to\_markovian.py:511}) is a
one-liner: \code{return str(SR(expr))}. The default Sage string form uses
\code{\^{}} for power and function names like \code{sqrt}, which is exactly the
syntax that \code{pipeline.theory\_compiler.\_CGFKernelCallable} re-evaluates
later (\file{pipeline/colored\_to\_markovian.py:517--519}). So the round-trip
symbolic~$\to$~text~$\to$~symbolic is lossless: every coefficient the detector
computes can be written into a new CGF row and re-parsed downstream without loss.

\section{Worked example: \texttt{ou\_quartic\_colored} before and after}

Concretely, take the shipped \code{ou\_quartic\_colored} theory and watch the
builder state change across \code{markovianize\_spec}.

\paragraph{Before.} Just before the call, the builder holds:
\begin{itemize}
  \item \textbf{physical fields:} \code{x} (internal \code{dx}, response
        \code{xt}).
  \item \textbf{one CGF row:}
        \code{\{'name': 'CXX', 'response\_legs': ['xt','xt'], 'order': 2,
        'coefficient': '2*D/tauc', 'kernel': 'exp(-abs(tau)/tauc)',
        'markovianize': None\}}.
  \item \textbf{action text:} \code{'xt*((Dt+mu)*x + eps*x\^{}3)'}.
  \item \code{\_markovianize\_default}: \code{True}.
\end{itemize}

\paragraph{Detection.} \code{detect\_lorentzian('exp(-abs(tau)/tauc)', '2*D/tauc',
[...])} returns \code{\{'tauc\_text': 'tauc', 'amplitude\_text': '2*D/tauc',
'aux\_drive\_text': '4*D/tauc\^{}2', ...\}}. The test suite pins both
\code{tauc\_text == 'tauc'} and \code{aux\_drive\_text == '4*D/tauc\^{}2'}
(\file{tests/test\_markovianize.py:71--73}).

\paragraph{After.} The builder now holds:
\begin{itemize}
  \item \textbf{physical fields:} \code{['dx', 'dxi']} --- the auxiliary \code{xi}
        (internal \code{dxi}, auto-response \code{xit}, auto-saddle \code{xistar})
        was added. The test asserts \code{phys\_names == ['dx', 'dxi']}
        (\file{tests/test\_markovianize.py:116}).
  \item \textbf{equations:} the user's $(\Dt+\mu)x = -\varepsilon x^3$ \emph{plus}
        the new aux drift $(\Dt + 1/\tau_c)\xi = 0$ --- two DAE equations.
  \item \textbf{one CGF row}, the colored \code{CXX} replaced by
        \code{\{'name': 'CXX\_markov\_aux', 'response\_legs': ['xit','xit'],
        'order': 2, 'coefficient': '4*D/tauc\^{}2', 'kernel': 'dirac\_delta(tau)',
        'markovianize': False\}}. The test confirms the row name ends in
        \code{\_markov\_aux} and the resolved legs are \code{['xit','xit']}
        (\file{tests/test\_markovianize.py:130, 132}).
  \item \textbf{action text:}
        \code{'xt*((Dt+mu)*x + eps*x\^{}3 - xi) + xit*((Dt + 1/tauc)*xi)'}.
\end{itemize}

\paragraph{Downstream.} \code{make\_correlated\_noises\_block} turns the white aux
row into a \code{model['correlated\_noises']} entry whose order-2 kernel callable
is \code{\_CGFKernelCallable('4*D/tauc\^{}2', 'dirac\_delta(tau)', 2)}. The diagram
engine attaches it to the \code{xit}--\code{xit} noise vertex. Because the kernel
is now $\delta(\tau)$, the fast white-noise analytic integrators run --- which was
the entire goal.

\subsection{The tree-level analytic cross-check}

The payoff is checkable in closed form. With $\varepsilon = 0$ (the free theory),
\code{compute\_cumulants(model, k=2, max\_ell=0, ...)} on the markovianized
\code{ou\_quartic\_colored} must return the analytic colored-OU equal-time
variance
\begin{equation}
  C(0) \;=\; \frac{2D}{\mu\,(\mu\,\tau_c + 1)},
  \label{eq:cm-treevariance}
\end{equation}
the simple-pole result the test pins down
(\file{tests/test\_markovianize.py}, the
\code{test\_colored\_ou\_tree\_level\_matches\_analytic} case, which runs at
$\mu = 0.1,\ \tau_c = 1,\ D = 1$). That the embedded white-noise model reproduces
\eqref{eq:cm-treevariance} is direct confirmation that the OU filter recovers the
original colored statistics.

\section{Limitations and how to live with them}

The v1 embedding is deliberately the minimum that unblocks the working
OU-colored theories. Its boundaries are worth knowing precisely.

\begin{description}
  \item[Single Lorentzian only.] Only $c\,\ee^{-|\tau|/\tau_c}$ matches. Gaussian
        kernels, the bare-exponential $\ee^{-\tau/\tau_c}$ (asymmetric),
        double-exponential, oscillatory, and polynomially-modulated kernels all
        fail the matcher and pass through to the slow scipy fallback. The future
        templates are enumerated in the docstring
        (\file{pipeline/colored\_to\_markovian.py:51--64}).

  \item[Order-2 only.] Only second cumulants are embedded; an \code{order != 2}
        row is passed through silently
        (\file{pipeline/colored\_to\_markovian.py:250--252}). A colored
        \emph{third} cumulant would hit the slow path.

  \item[Single $\tau_c$.] If matched rows carry \emph{different} $\tau_c$
        expressions, v1 raises \code{NotImplementedError}
        (\file{pipeline/colored\_to\_markovian.py:323--333}). All shipped examples
        use one shared $\tau_c$.

  \item[The \texorpdfstring{$\mu = 1/\tau_c$}{mu = 1/tauc} double pole.] This is
        the constraint most likely to bite. See the dedicated box below.
\end{description}

\begin{gotcha}
\textbf{Keep $\mu \neq 1/\tau_c$.} The free ($\ell = 0$) two-point function of the
embedded \emph{user} field is the convolution of the user-field propagator
$1/(\ii\omega + \mu)$ with the OU output spectrum $\propto 1/(\omega^2 +
1/\tau_c^2)$. In the frequency-domain residue calculus this product has poles at
$\omega = \ii\mu$ and $\omega = \ii/\tau_c$, and the analytic mode-sum integrator
computes residues \emph{assuming these poles are simple} (multiplicity~1). When
$\mu = 1/\tau_c$ exactly, the two poles \emph{collide} into a single double pole,
and the simple-pole residue infrastructure cannot form the multiplicity-2
residue. The closed-form variance \eqref{eq:cm-treevariance} is itself the
simple-pole result; at $\mu = 1/\tau_c$ it is the limit of a $0/0$. This is a
\emph{downstream} limitation --- it lives in the residue calculus, not in this
module --- and the embedding code does \emph{not} guard against it. The test
deliberately picks $\mu = 0.1,\ \tau_c = 1$ to avoid the degeneracy, with an
explicit inline note (\file{tests/test\_markovianize.py:165--168}). When you
parameterise a markovianized theory, keep $\mu$ away from $1/\tau_c$.
\end{gotcha}

\begin{note}
\textbf{Roadmap (v2).} The docstring enumerates the natural extensions, each of
which simply enlarges the auxiliary state space
(\file{pipeline/colored\_to\_markovian.py:51--64}): \emph{underdamped oscillatory}
kernels $\ee^{-|\tau|/\tau_d}\cos(\omega\tau)$ as a 2-state OU process with an
imaginary eigenvalue; \emph{double-exponential} kernels as a sum of two
independent OU processes; and \emph{polynomially-modulated} kernels as a chain of
OU processes (a higher-order linear filter). The single-Lorentzian case in this
chapter is the lowest-cost member of that family.
\end{note}

\section{Summary}

Colored noise is the one input that white-noise field theory cannot integrate
analytically. The fix is structural rather than numerical: rewrite the model so
the memory lives in an extra dynamical field. This chapter showed the full chain
--- a single-Lorentzian autocovariance \eqref{eq:cm-secondcumulant}, the OU filter
\eqref{eq:cm-toptline}--\eqref{eq:cm-bottomline} whose stationary covariance is
exactly that Lorentzian \eqref{eq:cm-ouvariance}, the three \msrjd{} structural
pieces and two action edits that encode it, the symbolic kernel matcher
\code{detect\_lorentzian} (\S\ref{sec:cm-detect}), and the in-place builder
rewrite \code{markovianize\_spec} (\S\ref{sec:cm-transform}) wired into
\code{build()} before compilation. The transformation is invisible to everything
downstream: by the time the field-theory core in the next part of the manual sees
your model, the colored noise has become an auxiliary OU field driven by white
noise, and the fast analytic integrators apply again.


% =====================================================================
%  PART III  --  BUILDING THE PERTURBATION THEORY
% =====================================================================
\part{Building the Perturbation Theory}
% ---------------------------------------------------------------------
%  Chapter 6 : The Field-Theory Core --- Action, Taylor Expansion, Vertices
% ---------------------------------------------------------------------
\chapter{The Field-Theory Core: Action, Taylor Expansion, Vertices}\label{ch:fieldtheory-core}

This is the first of the machinery chapters, and it opens the very first of the
seven phases you watched scroll past in the quickstart trace: \code{[1/7]
FieldTheory.expand}. Everything downstream --- the propagator, the mean-field
solve, the diagram enumeration, the loop integrals --- consumes the output of
this one stage. If you understand nothing else about the internals, understand
this: the field-theory core is the machine that turns a \emph{declarative model}
(a Python dictionary of strings, lists, and small functions) into a
\emph{classified symbolic action} --- a dictionary of polynomial pieces sorted by
how many response legs and how many physical legs each piece carries. The free
propagator, the noise sources, and the interaction vertices all fall out of that
single classification.

The subsystem lives in four files, all under \file{msrjd/core/}:

\begin{itemize}
  \item \file{field\_theory.py} --- the \code{FieldTheory} expander itself: it
        builds the symbolic namespace, evaluates the action, Taylor-expands every
        nonlinearity, classifies the result by ``bigrade,'' and verifies the
        mean-field sector vanishes.
  \item \file{vertices.py} --- decomposes each classified sector into individual
        typed records (\code{VertexType}, \code{SourceType}, and their
        kernel-carrying subclasses) that the diagram machinery can build with.
  \item \file{convolution.py} --- the formal \code{Conv(kernel, field)} operator
        and the rules that resolve it, so a synaptic kernel times a field means
        \emph{convolution in time}, not ordinary multiplication.
  \item \file{serialize.py} --- saves an expanded theory to disk (so the
        expensive Taylor step is not repeated) and re-imports a model on demand.
\end{itemize}

We will work two examples in parallel. The running ``toy'' is the quartic
Ornstein--Uhlenbeck model from Chapter~\ref{ch:quickstart} --- one physical field,
white noise, a single cubic vertex --- because it is small enough that you can
hold the entire expanded action in your head. The running ``real'' example is the
quadratic Hawkes model shipped in \file{models/hawkes\_quad\_expg.py} --- two
populations, a synaptic convolution kernel, a formal transfer function --- because
it exercises every feature the subsystem has. Reading that model file alongside
this chapter is recommended; it is the canonical shape of the ``model dict'' the
core consumes.

\begin{note}
A vocabulary reminder from the physics, since the code leans on it constantly. In
the \msrjd{} construction every \term{physical field} (the thing that actually
fluctuates --- a voltage $v$, a spike rate $n$, our $x$) is paired with a
conjugate \term{response field}, written with a tilde ($\tilde v$, $\tilde n$,
$\tilde x$) and called the ``tilde'' or ``hat'' field. In diagram language a
physical leg is an ``in-leg'' and a response leg is an ``out-leg.'' The entire
classification scheme of this chapter is built on counting these two kinds of
legs separately.
\end{note}

\section{The MSR-JD action as data}

The single most important design choice in this subsystem --- the one that makes
it model-independent --- is that \emph{a model is data, not code}. You do not
write a new expander for each physics problem. You hand the one expander a Python
dictionary describing your model, and it does the rest. So before any algebra, we
need to understand the shape of that dictionary and the object the framework
builds from it.

\subsection{The model dictionary}

A model is a \code{dict}. Its values are a mix of plain data (strings, numbers,
lists of small spec dictionaries) and \emph{lambdas} --- small functions that,
given a namespace object, return a symbolic expression. The keys the core reads
are summarized in Table~\ref{tab:modeldict}.

\begin{table}[h]
\centering
\caption{The keys of a model dict that the field-theory core reads. ``lambda
\code{ns}$\to$\,\dots'' means a function taking the namespace object and returning
the indicated thing. Optional keys are noted.}
\label{tab:modeldict}
\small
\begin{tabular}{@{}lll@{}}
\toprule
key & type & meaning \\
\midrule
\code{name}              & str  & human-readable model name \\
\code{index\_sets}       & dict & named index ranges; first is the primary index \\
\code{populations}       & list & per-population \code{\{name, size\}} (optional) \\
\code{response\_fields}  & list & specs for the tilde fields \\
\code{physical\_fields}  & list & specs for the fluctuation fields \\
\code{parameters}        & list & \code{\{name, indexed\_by, domain, default, \dots\}} \\
\code{kernels}           & list & \code{\{name, sage\_name, indexed\_by, \dots\}} \\
\code{operators}         & list & \code{\{name, sage\_name\}} (e.g.\ \code{Dt}) \\
\code{functions}         & list & nonlinear fns: \code{expression} or \code{taylor\_coeffs} \\
\code{mf\_substitutions} & list & \code{\{name, value: lambda ns$\to$SR\}} \\
\code{mf\_bg\_conditions}& lambda & saddle substitutions for sector verification \\
\code{specializations}   & lambda & rename abstract derivative symbols to values \\
\code{correlated\_noises}& dict & colored-noise cumulant declarations (optional) \\
\code{kernel\_ft\_image} & lambda & $\{$kernel\,$\mapsto\hat g(\omega)\}$ (Fourier image) \\
\code{kernel\_td\_image} & lambda & $\{$kernel\,$\mapsto g(\tau)\}$ (time-domain image) \\
\code{action}            & lambda & \textbf{the action} \\
\bottomrule
\end{tabular}
\end{table}

You will rarely author this dict by hand. The theory builders of
Chapter~\ref{ch:theory-spec} generate it for you from a much friendlier fluent
API (recall
\code{TemporalTheoryBuilder('OU Quartic')\allowbreak.physical\_field('x')\dots}).
But the core does not know about the builders; it only sees the dict. So
everything in this chapter is phrased in terms of these keys.

The one key that does the real work is \code{action}: a function that, handed the
namespace, returns one big symbolic expression for $S$. For the quartic OU model
the builder produces an action equivalent to
\[
  S \;=\; \sum_{i}\Bigl[\tilde x_i\,\bigl((\partial_t + \mu)\,x_i + \varepsilon\,x_i^{3}\bigr)
          \;-\; D\,\tilde x_i^{2}\Bigr],
\]
which is exactly the action text \code{xt[i]*((Dt+mu)*x[i] + eps*x[i]\^{}3) - D*xt[i]\^{}2}
you saw in the quickstart theory file. For the Hawkes model the action authored at
\file{hawkes\_quad\_expg.py:236} is the richer
\[
  S \;=\; \sum_i\Bigl[\tilde n_i\,(\dot{n}^{*}_i + \delta n_i)
        \;-\;(\ee^{\tilde n_i}-1)\,\phi_i(\delta v_i)
        \;+\;\tilde v_i\bigl[(\tau\partial_t+1)\delta v_i + v^{*}_i - E_i
              - \textstyle\sum_j w_{ij}\,g\!\cdot\!(n^{*}_j+\delta n_j)\bigr]\Bigr],
\]
with response fields $\tilde n_i,\tilde v_i$, fluctuation fields
$\delta n_i,\delta v_i$, a synaptic kernel $g$ (a convolution), and a nonlinear
transfer function $\phi_i(v)=a\,v^2$.

\subsection{The namespace object}

The action lambda takes one argument, conventionally called \code{ns}, and writes
the action in terms of attributes of \code{ns}: \code{ns.xt[i]}, \code{ns.x[i]},
\code{ns.mu}, \code{ns.Dt}, and so on. That \code{ns} object is an instance of a
deliberately empty class:

\begin{lstlisting}
# msrjd/core/field_theory.py:185
class _Namespace:
    pass
\end{lstlisting}

It is empty because it is a blank slate that \code{FieldTheory.\_build\_namespace}
fills, one attribute at a time, with a symbolic variable for every field, every
parameter, every kernel, every operator, and every function the model declares
(we walk that construction in detail in \S\ref{sec:buildns}). The action lambda
never has to import anything model-specific; it just reaches into the namespace it
was handed. This indirection is the whole trick that lets one expander serve every
model.

\begin{gotcha}
\code{\_Namespace} also carries a pile of private bookkeeping attributes whose
names start with an underscore: \code{ns.\_all\_field\_sr\_vars} (the list of every
field symbol), \code{ns.\_ring\_var\_names} (the ring generator names, tilde
first), \code{ns.\_deriv\_rename\_subs} (the formal-derivative rename map), and
more. These are not for the model author to touch; they are how the expander
passes state from \code{\_build\_namespace} to the rest of \code{expand}. The
namespace is also \emph{unpicklable} --- it holds inner closures --- which becomes
important in \S\ref{sec:pickle}.
\end{gotcha}

\section{SageMath in one page}

Every field, parameter, kernel, and operator in this subsystem is a SageMath
symbolic object. If you have never used Sage, this section is the minimum you need
to read the code; everything here recurs throughout the manual.

\term{SageMath} is a large open-source mathematics system built on Python. It
bundles many computer-algebra engines (Maxima, Pari, Singular, FLINT, GMP, SymPy,
\dots) behind one Python API. This subsystem uses exactly two pieces of it.

\paragraph{The Symbolic Ring \code{SR}.} \code{SR} is Sage's universal ring of
symbolic expressions --- think ``anything you could write on a blackboard'':
\code{a*x\^{}2 + sin(t) + exp(nt)}. \code{SR.var('a')} creates a symbolic variable
named \code{a}. Wrapping a value in \code{SR(...)} coerces it into the symbolic
ring, so \code{SR(0)} is the symbolic zero. Crucially, an \code{SR} variable can
carry an \emph{assumption} (e.g.\ \code{domain='positive'}, which tells Sage the
symbol is a positive real) and a \emph{LaTeX display name} (\code{latex\_name=...},
so a notebook renders $\mu$ instead of \code{mu}).

\paragraph{\code{PolynomialRing(SR, names)}.} This builds a polynomial ring
\emph{whose coefficients live in \code{SR}} and whose \emph{generators} are the
named field variables. So \code{PolynomialRing(SR, ['xt','dx'])} represents the
action as a genuine polynomial in the two field generators \code{xt} and
\code{dx}, while the parameters ($\mu$, $\varepsilon$, $D$) live symbolically
\emph{in the coefficients}. This separation is the linchpin of the whole chapter:
a polynomial-ring element exposes its monomials through \code{.dict()}, a map
\code{\{exponent\_vector: coefficient\}}, and the exponent vector immediately tells
you how many of each field a monomial contains.

The exact imports the expander uses make the surface concrete:

\begin{lstlisting}
# field_theory.py:26-29
from sage.all import (
    SR, PolynomialRing, factorial, QQ, latex, LatexExpr,
    diff, function, exp, dirac_delta, heaviside, integrate, oo, I, pi, taylor
)
\end{lstlisting}

A quick gloss of the less obvious names: \code{factorial} and \code{QQ} (Sage's
field of \emph{exact} rationals) build Taylor coefficients $1/k!$ without ever
touching floating point; \code{diff} is symbolic differentiation; \code{function}
builds a \emph{formal} (unevaluated) symbolic function (\S\ref{sec:formal});
\code{dirac\_delta} and \code{heaviside} are $\delta(t)$ and the step function;
\code{integrate} is symbolic integration; \code{oo} is $\infty$, \code{I} is
$\sqrt{-1}$, \code{pi} is $\pi$; and \code{taylor} is the multivariate Taylor
expander (\S\ref{sec:taylor}).

Two Sage behaviors trip up newcomers so hard that the code comments call them out
by name. They are worth a callout box of their own, because getting either wrong
\emph{silently corrupts the answer}.

\begin{gotcha}
\textbf{A Sage sum is \code{add\_vararg}, not \code{operator.add}.} When you ask a
Sage expression for its top-level operator with \code{expr.operator()}, a sum does
\emph{not} report Python's \code{operator.add}; it reports Sage's own internal
n-ary callable \code{add\_vararg}. So the code never compares operators by
identity. Everywhere it needs to ask ``is this a sum?'' it does a name-based test:
\begin{lstlisting}
# field_theory.py:266-267  (inside _sr_to_ring)
_op_obj = expanded.operator()
if _op_obj is not None and 'add' in getattr(_op_obj, '__name__', '').lower():
    summands = expanded.operands()
\end{lstlisting}
The same pattern appears in \file{convolution.py:254}. Getting this wrong does not
raise an error --- it just fails to recognize the sum, and terms are silently
dropped.
\end{gotcha}

\begin{gotcha}
\textbf{\code{expr != 0} is unreliable for symbolic zero.} Under Sage, when an
expression contains an unbound real parameter with no positivity assumption (e.g.\
a symbol declared \code{domain='real'}), the comparison \code{expr != 0} can return
\code{False} --- Sage cannot decide the sign and refuses to commit. The correct
\emph{syntactic} zero test is \code{.is\_trivial\_zero()}, which is \code{True}
exactly when the expression is literally zero, regardless of any parameter
assumptions. The cumulant injector relies on it (\file{field\_theory.py:490}),
with the reason spelled out in the comment right above the call.
\end{gotcha}

\section{From action to polynomial: the \texttt{expand()} pipeline}

With the vocabulary in place we can read the heart of the subsystem:
\code{FieldTheory.expand} at \file{field\_theory.py:814}. It is a single method
that performs an eight-step transform. We list the steps first, then spend the
following sections opening up the ones that carry real subtlety (Taylor expansion,
convolution, bigrade classification, mean-field verification). Here is the
skeleton, lightly trimmed:

\begin{lstlisting}
# field_theory.py:814  (FieldTheory.expand, abridged)
def expand(self) -> None:
    ns, R, n_tilde = self._build_namespace()           # 1. build namespace
    self._ns, self._R, self._n_tilde = ns, R, n_tilde

    S_sr = SR(self.model['action'](ns))                # 2. evaluate the action

    # 3. resolve Conv(kernel, field) atoms
    S_sr = reduce_conv_in_action(S_sr, fluct_vars,
               taylor_order=self.taylor_order, attachments_out=attachments)
    ns._kernel_attachments = attachments

    # 4. inject correlated-noise cumulant series
    S_sr = S_sr + _build_cumulant_action(ns, self.model)

    # 5. multivariate Taylor to total degree taylor_order
    if ns._all_field_sr_vars:
        S_sr = taylor(S_sr, *[(v, 0) for v in ns._all_field_sr_vars],
                      self.taylor_order)

    # 6. rename formal-derivative symbols  D[0](phi_1)(0) -> phi1_1
    S_sr = S_sr.subs(ns._deriv_rename_subs)

    # 7. apply model specializations  phi1_i -> 2 a vstar_i, ...
    S_sr = S_sr.subs(self.model['specializations'](ns))

    # 8. coerce to the polynomial ring and classify BEFORE MF subs
    S_poly = _sr_to_ring(S_sr.expand(), R, ns._ring_var_names)
    by_tp  = _collect_bigrade(S_poly, n_tilde)
    # ... then verify+zero the mean-field sector, rebuild S_raw ...
    self._by_tp = by_tp
\end{lstlisting}

Read top to bottom:

\begin{enumerate}
  \item \textbf{Build the namespace.} \code{\_build\_namespace} returns the triple
        \code{(ns, R, n\_tilde)}: the populated namespace, the polynomial ring, and
        the number of response-field generators. Detailed in \S\ref{sec:buildns}.
  \item \textbf{Evaluate the action.} \code{self.model['action'](ns)} calls the
        model's action lambda with the namespace; the result is one large \code{SR}
        expression. Wrapping it in \code{SR(...)} guarantees it is in the symbolic
        ring even if the lambda returned a bare integer.
  \item \textbf{Resolve convolutions.} Every \code{Conv(g, field)} atom is rewritten
        by \code{reduce\_conv\_in\_action} (\S\ref{sec:conv}). The whole step is
        wrapped in a \code{try/except (AttributeError, TypeError)} so a model with
        no convolutions passes through untouched.
  \item \textbf{Inject correlated-noise cumulants.} If the model declares colored
        noise, \code{\_build\_cumulant\_action} appends a cumulant series to the
        action; otherwise it returns \code{SR(0)}, a true no-op (\S\ref{sec:cumulant}).
  \item \textbf{Taylor-expand.} Every nonlinear function of the fields becomes a
        finite polynomial, truncated at total degree \code{taylor\_order}
        (\S\ref{sec:taylor}).
  \item \textbf{Rename derivative symbols.} Sage produces ugly atoms like
        \code{D[0](phi\_1)(0)} when it differentiates a formal function; these are
        renamed to clean polynomial-friendly names like \code{phi1\_1}
        (\S\ref{sec:formal}).
  \item \textbf{Apply specializations.} The model substitutes concrete values for
        those abstract derivative symbols (for $\phi=a v^2$, $\phi''\to 2a$). These
        are pure renames, safe to apply to every sector.
  \item \textbf{Coerce and classify.} The \code{SR} expression is converted into a
        polynomial-ring element by \code{\_sr\_to\_ring}, then split into bigrade
        sectors by \code{\_collect\_bigrade} (\S\ref{sec:bigrade}). Only \emph{after}
        this does the mean-field sector get verified and zeroed (\S\ref{sec:mf}).
\end{enumerate}

\begin{gotcha}
The docstring of \code{expand} (\file{field\_theory.py:818}) still describes step~5
as ``sequential \code{taylor()}.'' The code at \file{field\_theory.py:894} actually
does a single one-shot multivariate \code{taylor}. Trust the code: the expansion is
multivariate, for the reason given in \S\ref{sec:taylor}.
\end{gotcha}

The order of these steps is not arbitrary. The most important ordering constraint
--- classify \emph{before} applying mean-field substitutions --- is the subject of
\S\ref{sec:mf}, and the comment at \file{field\_theory.py:913--922} is worth
reading in the source.

\subsection{Building the namespace}\label{sec:buildns}

\code{FieldTheory.\_build\_namespace} (\file{field\_theory.py:1041}) is the longest
method in the file. You do not need every line, but you should know the order in
which it populates \code{ns}, because that order fixes the ring generator ordering
that the entire classification depends on.

\begin{enumerate}
  \item \textbf{Index sets and populations.} Each entry of \code{model['index\_sets']}
        is bound onto \code{ns}; the first is the ``primary'' index. For each
        \code{model['populations']} entry it builds a $0\!\dots\!(\text{size}-1)$
        index list, so action text can write \code{for i in pop}. Legacy theories
        with no populations fall back to a single flat \code{pop}.
  \item \textbf{Field variables.} For each \code{response\_fields} and
        \code{physical\_fields} spec it creates the \code{SR} symbols: an indexed
        field becomes a list \code{[SR.var(f"\{fname\}\{i+1\}", latex\_name=...)]}
        (so field \code{x} in a population of one becomes the symbol \code{x1}); a
        scalar field becomes a single symbol. It accumulates the tilde symbols and
        the physical symbols separately and stashes them as
        \code{ns.\_tilde\_sr\_vars}, \code{ns.\_phys\_sr\_vars}, and the union
        \code{ns.\_all\_field\_sr\_vars}.
  \item \textbf{The polynomial ring.} It forms
        \code{ring\_var\_names = tilde\_names + phys\_names} --- \emph{tilde names
        first} --- sets \code{n\_tilde = len(tilde\_names)}, and builds
        \code{R = PolynomialRing(SR, ring\_var\_names)}. This is why every later
        exponent-vector split at index \code{n\_tilde} cleanly separates response
        from physical legs.
  \item \textbf{Parameters.} Scalar parameters get one symbol (with a \code{domain}
        assumption if declared); vector/matrix parameters get a flat list of symbols
        sized by population.
  \item \textbf{Kernels and operators.} A scalar kernel becomes one symbol named by
        its \code{sage\_name}; a vector kernel becomes \code{z\_<name>\_<i+1>}; a
        matrix kernel becomes a list-of-lists \code{z\_<name>\_<i+1>\_<j+1>}.
        Operators (like the time derivative \code{Dt}) become one symbol each. It
        also defines \code{ns.delta\_D = SR.var('z\_delta', \dots)} and a primed
        variant. The leading \code{z} in these internal names is deliberate: \code{z}
        sorts \emph{after} \code{phi}/\code{tau} in any Sage product, giving a
        canonical, reproducible factor ordering.
  \item \textbf{Nonlinear functions.} Each \code{model['functions']} spec installs
        a callable on \code{ns} and registers its derivative-rename entries. There
        are two authoring paths --- an \code{'expression'} path and a legacy
        \code{'taylor\_coeffs'} path --- detailed in \S\ref{sec:formal}.
  \item \textbf{Mean-field substitutions.} For each \code{model['mf\_substitutions']}
        it does \code{setattr(ns, sub['name'], sub['value'](ns))}, computing the
        value once at build time (this is how a saddle quantity like \code{vstar}
        gets onto the namespace so the action lambda can reference it).
\end{enumerate}

The method returns \code{(ns, R, n\_tilde)}. The take-away to carry forward: the
ring generators are ordered \emph{all tilde fields, then all physical fields}, and
\code{n\_tilde} is the boundary between the two halves.

\section{Taylor expansion and formal-function renaming}\label{sec:taylor}

A field theory's interactions live in the nonlinear terms of its action: the
$\varepsilon x^3$ of the OU model, the $\ee^{\tilde n}-1$ and the $\phi(v)$ of the
Hawkes model. To compute Feynman diagrams these must be replaced by finite
polynomials. That is what step~5 of \code{expand} does.

\subsection{Why total-degree truncation}

For a single-argument function expanded around the saddle ($\delta\phi=0$), the
Taylor series truncated at order $N$ is the familiar
\[
  f(\delta\phi) \;=\; \sum_{k=0}^{N} \frac{f^{(k)}(0)}{k!}\,\delta\phi^{\,k}.
\]
For a multi-argument function the framework uses the multivariate Taylor series
truncated at \emph{total} degree $N$:
\[
  f(x) \;=\; \sum_{|\alpha|\le N} \frac{1}{\alpha!}\,(\partial^{\alpha} f)(0)\,x^{\alpha},
\]
where $\alpha=(k_1,\dots,k_n)$ is a multi-index, $|\alpha|=\sum_j k_j$,
$\alpha!=\prod_j k_j!$, and $x^{\alpha}=\prod_j x_j^{k_j}$.

Truncating at \emph{total} degree (not per-variable degree) is the physically
correct choice, and it is worth being explicit about why. The total degree of a
monomial in the fields equals the number of legs of the diagrammatic vertex it
produces. So ``Taylor order $N$'' means exactly ``keep every vertex with up to $N$
legs.'' This is the single knob \code{taylor\_order} (default $4$), and it is the
same knob the quickstart's \code{[1/7]} trace printed as \code{taylor\_order=...}.

The code realizes this with one Sage call:

\begin{lstlisting}
# field_theory.py:894
S_sr = taylor(
    S_sr,
    *[(v, 0) for v in ns._all_field_sr_vars],
    self.taylor_order,
)
\end{lstlisting}

Sage's \code{taylor(expr, (v1,0), (v2,0), ..., N)} expands \code{expr} jointly in
all listed variables, each about $0$, to total order $N$. The \code{*[(v, 0) for
...]} splat hands it one \code{(variable, 0)} pair per field.

\begin{gotcha}
This is a \emph{one-shot multivariate} Taylor, and it must be. An earlier version
of the code chained single-variable \code{.taylor()} calls in a loop. That is wrong
for multi-argument formal functions: chained \code{.taylor()} does not compose at
non-zero expansion points, because the outer call freezes the inner-argument
substructure of the formal function. The comment at \file{field\_theory.py:884--892}
documents the switch. For single-argument functions the one-shot result is identical
to the old loop, so the generalization is safe.
\end{gotcha}

\subsection{Formal functions and the rename dance}\label{sec:formal}

A model author often does not want to commit to a closed form for a transfer
function. Instead they declare $\phi$ as a \term{formal function}: a Sage symbolic
function left unevaluated, built with \code{function('phi\_1')(v)}. A formal
function stays inert through symbolic manipulation --- you can differentiate it,
multiply it, Taylor-expand it, all abstractly. That is exactly what we want: we
expand $\phi$ without ever saying what $\phi$ is.

The catch is what Sage produces when it Taylor-expands a formal function. The
$k$-th derivative at $0$ comes out as an atom like \code{D[0](phi\_1)(0)} (first
derivative of \code{phi\_1} at $0$) or \code{D[0,0](phi\_1)(0)} (second
derivative). These are not valid polynomial-ring coefficients. So step~6 renames
them to clean symbols:
\[
  \texttt{phi\_1(0)} \;\to\; \texttt{phi0\_1}, \qquad
  \texttt{D[0](phi\_1)(0)} \;\to\; \texttt{phi1\_1}, \qquad
  \texttt{D[0,0](phi\_1)(0)} \;\to\; \texttt{phi2\_1}.
\]
For a single-argument function the suffix is just the derivative order $k$; for a
multi-argument function it is the concatenated multi-index, so $f^{(2,1)}$ becomes
\code{f21\_<i>}. That suffix encoding is a one-liner:

\begin{lstlisting}
# field_theory.py:225
def _multi_index_suffix(multi_idx):
    return ''.join(str(k) for k in multi_idx)
\end{lstlisting}

The rename map itself, \code{ns.\_deriv\_rename\_subs}, is built during
\code{\_build\_namespace}. To know \emph{which} derivative atoms to expect, the
framework enumerates every multi-index up to the Taylor order with a small
generator:

\begin{lstlisting}
# field_theory.py:202
def _iter_multi_indices(n_args, max_total):
    if n_args == 0:
        yield ()
        return
    from itertools import product
    for combo in product(range(max_total + 1), repeat=n_args):
        if sum(combo) <= max_total:
            yield combo
\end{lstlisting}

\begin{note}
Why \code{itertools.product} rather than a tidy self-recursive generator? Because
this code runs in Jupyter under \code{\%autoreload}. A self-recursive generator
holds a reference to itself through the module's global namespace; when
\code{\%autoreload} hot-swaps the module, that self-reference is broken and the
generator fails mid-iteration. \code{itertools.product} carries no such
self-reference, so it survives the reload. The comment at
\file{field\_theory.py:210--215} records this hard-won lesson. You will see the same
defensive pattern in a few other places in the codebase.
\end{note}

Once the abstract symbols are in place, a model gives them concrete values through
its \code{specializations} lambda (step~7). For the Hawkes transfer function
$\phi=a v^2$, the specialization (at \file{hawkes\_quad\_expg.py:186}) sets
\code{phi1\_i}$\to 2a\,v^{*}_i$, \code{phi2\_i}$\to 2a$, and every higher
derivative to $0$ --- and just like that the formal $\phi$ collapses to its
quadratic reality. The \code{'expression'} authoring path in
\code{\_build\_namespace} automates even this: the model supplies the closed-form
expression, and the framework differentiates it symbolically at the all-zero point
with \code{diff} to compute each $\partial^{\alpha}f(0)$ itself, skipping any
derivative that comes out purely numeric. The legacy \code{'taylor\_coeffs'} path
instead takes a literal list of derivative coefficients and feeds them to
\code{\_poly\_taylor} (\file{field\_theory.py:193}), which assembles
$\sum_n (c_n/n!)\,x^n$ using the exact rational $1/n!$ from \code{QQ}.

\section{The convolution operator}\label{sec:conv}

This section explains step~3 of \code{expand}, and it is the one piece of the
pipeline whose \emph{whole reason to exist} is a subtle physics correctness issue.

\subsection{Why \texorpdfstring{$g \ast n$}{g*n} is not multiplication}

When a synaptic kernel $g$ multiplies a field $n$ in the action, the \code{*} in
the source text \emph{reads} as scalar multiplication but \emph{means} convolution
in time:
\[
  (g \ast n)(t) \;=\; \int g(t-s)\,n(s)\,\dd s.
\]
The framework represents this with a formal two-argument Sage function:

\begin{lstlisting}
# convolution.py:70
Conv = _sr_function('Conv', nargs=2)
\end{lstlisting}

\begin{defn}
\code{Conv(kernel, field)} is a formal time-convolution operator: by convention the
\emph{kernel} is the first argument and the \emph{field} is the second. It never
auto-evaluates --- it stays a formal symbol through all symbolic manipulation, so
each downstream consumer (the stationary mean field, the Fourier propagator, the
time-domain integrator) can apply its own reduction. The \code{nargs=2} declaration
makes Sage reject malformed calls like \code{Conv(g)} or \code{Conv(g, n, extra)}
immediately rather than producing a silently wrong expression.
\end{defn}

The asymmetry between the two arguments is the entire point. Only the second
argument (the field) carries time-domain content; the first (the kernel) is a
fixed translation-invariant filter. Consider a conductance-style term
$v\cdot\mathrm{Conv}(g,n)$ and expand both fields around the saddle,
$v=v^{*}+\delta v$, $n=n^{*}+\delta n$:
\[
  v\cdot\mathrm{Conv}(g,n)
  \;=\; v^{*}n^{*} \;+\; v^{*}\,g\,\delta n \;+\; \delta v\,n^{*} \;+\; \delta v\,g\,\delta n.
\]
Notice $\delta v$ multiplies the bare $n^{*}$ --- \emph{not} $g\cdot n^{*}$. A naive
``flatten $\mathrm{Conv}(g,n)\to g\cdot n$'' would instead give
$v\cdot g\cdot n = v^{*}g\,n^{*}+\dots$, spuriously attaching the kernel to the
$\delta v$ side of the bilinear. That single wrong coefficient would poison the free
propagator. The \code{Conv} operator exists precisely so this expansion comes out
right.

\subsection{The four reduction rules}

\code{reduce\_conv\_in\_action} (\file{convolution.py:110}) resolves every
\code{Conv} atom by handing Sage a Python callback through
\code{.substitute\_function(Conv, \_reduce\_arg)} --- Sage's mechanism for replacing
every occurrence of a formal function with the result of a callback that receives
the function's arguments. The callback applies, in order:

\begin{description}
  \item[Rule 0 --- pre-Taylor.] When a \code{taylor\_order} is supplied (it always
        is, from \code{expand}), Taylor-expand the \emph{argument} first, so a
        nonlinear inner field like $\phi(v)$ becomes a polynomial that the linearity
        rule can distribute over. A no-op for an already-polynomial argument such as
        a bare \code{n[j]}.
  \item[Rule 1 --- linearity.] $\mathrm{Conv}(g, a+b)\to\mathrm{Conv}(g,a)+\mathrm{Conv}(g,b)$,
        recursing over the summands of a sum.
  \item[Rule 2 --- pull constants.] $\mathrm{Conv}(g, c\cdot X)\to c\cdot\mathrm{Conv}(g,X)$
        for any factor $c$ with no fluctuation-field dependence (time-constant
        factors come out of the time convolution).
  \item[Rule 3 --- DC reduction.] $\mathrm{Conv}(g, c)\to\hat g(0)\cdot c = c$ for a
        time-constant leaf, under the normalized-kernel assumption $\hat g(0)=1$.
  \item[Rule 4 --- defer to the Fourier transform.] $\mathrm{Conv}(g, X)\to g\cdot X$
        for a genuinely fluctuating $X$. The kernel \emph{symbol} $g$ is left in
        place; downstream it is replaced by $\hat g(\omega)$ via
        \code{model['kernel\_ft\_image']}, which is just the convolution theorem
        $(g\ast n)\hat{\phantom{x}}(\omega)=\hat g(\omega)\,\hat n(\omega)$.
\end{description}

\begin{gotcha}
The \code{normalized=True} assumption ($\hat g(0)=1$, i.e.\ every kernel integrates
to one) is hard-wired in rule~3. Non-normalized kernels would need per-kernel DC
gains, which are not yet implemented; see \file{convolution.py:182--185}. For the
neural-style kernels this framework targets --- a synaptic filter normalized to unit
area --- the assumption holds by construction.
\end{gotcha}

\subsection{Kernel attachments}

Rule~4 has a side effect that matters two phases downstream. When it emits
$g\cdot X$ it records, into the \code{attachments\_out} dictionary, which
fluctuation variables the surviving kernel symbol $g$ attached to (set-union
accumulating across reuse). \code{expand} stashes this on the namespace as
\code{ns.\_kernel\_attachments}. Later, vertex extraction reads it to recover which
\emph{physical leg} each leftover kernel symbol belongs to, so the time-domain
integrator can place that leg at the correct convolved time. We pick this thread
back up in \S\ref{sec:typed} when we meet \code{ConvVertexType}.

\section{Bigrade classification}\label{sec:bigrade}

We have reached the conceptual core. After Taylor expansion, renaming, and
specialization, the action is one big polynomial in the field generators. Step~8
turns it into something the rest of the pipeline can navigate: a dictionary keyed
by \emph{bigrade}.

\begin{defn}
The \term{bigrade} of a monomial is the pair
$(n_{\tilde{}}, n_{\mathrm{phys}})$ where $n_{\tilde{}}$ is the total exponent over
the \emph{response} (tilde) generators and $n_{\mathrm{phys}}$ is the total exponent
over the \emph{physical} (fluctuation) generators. It is the number of response legs
and the number of physical legs the corresponding diagrammatic object carries.
\end{defn}

The physical meaning of each sector is the whole reason to classify this way:

\begin{table}[h]
\centering
\caption{The bigrade sectors and their physical roles. The first three rows are
collectively the mean-field (saddle-equation) sector and must vanish at the saddle.}
\label{tab:bigrade}
\small
\begin{tabular}{@{}lll@{}}
\toprule
bigrade $(n_{\tilde{}},n_p)$ & name & meaning \\
\midrule
$(0,0)$        & constant      & must vanish (no constant piece at the saddle) \\
$(1,0)$        & tadpole       & response-linear; must vanish at the saddle \\
$(0,1)$        & EOM residual  & physical-linear; must vanish at the background \\
$(1,1)$        & \textbf{free action}   & the bilinear kernel; its inverse is the free propagator \\
$(\ge 2,\,0)$  & \textbf{noise kernel}  & response-only, $\ge$ quadratic; the noise sources \\
total $\ge 3$  & \textbf{interaction vertex} & every higher-order monomial is a vertex \\
\bottomrule
\end{tabular}
\end{table}

The mechanics are remarkably short because the polynomial ring does the heavy
lifting. First, \code{\_sr\_to\_ring} (\file{field\_theory.py:241}) converts the
\code{SR} expression into a ring element. It expands the expression, splits it into
summands (using the name-based \code{add} test from the earlier gotcha), and for
each summand reads off the degree in each field generator, stripping it away and
accumulating whatever symbolic coefficient remains:

\begin{lstlisting}
# field_theory.py:272-281  (inside _sr_to_ring)
for term in summands:
    exponents = []
    coeff = SR(term)
    for v_sr in gen_sr:
        deg = int(coeff.degree(v_sr))
        exponents.append(deg)
        if deg > 0:
            coeff = coeff.coefficient(v_sr, deg)
    result += SR(coeff) * R.monomial(*exponents)
\end{lstlisting}

\begin{gotcha}
The summand/single-term distinction in \code{\_sr\_to\_ring} is explicit for a
reason. If the expanded action is a single product --- say $\tau\cdot\mathrm{Dt}\cdot v$
--- then calling \code{.operands()} on it would iterate its \emph{factors}
($\tau$, $\mathrm{Dt}$, $v$), not its summands, and the coefficient $\tau\cdot\mathrm{Dt}$
would be torn apart and lost. So the code wraps any non-sum in a one-element list
\code{[expanded]} before iterating (\file{field\_theory.py:268--271}). This is the
same \code{add\_vararg} hazard wearing a different hat.
\end{gotcha}

Then \code{\_collect\_bigrade} (\file{field\_theory.py:570}) walks the ring
element's \code{.dict()} and buckets each monomial by its bigrade. Because the ring
generators are ordered tilde-first, the split index is exactly \code{n\_tilde}:

\begin{lstlisting}
# field_theory.py:570
def _collect_bigrade(poly, n_tilde: int) -> dict:
    R = poly.parent()
    sectors: dict = {}
    for exp_vec, coeff in poly.dict().items():
        n_t = int(sum(exp_vec[:n_tilde]))
        n_p = int(sum(exp_vec[n_tilde:]))
        key = (n_t, n_p)
        sectors[key] = sectors.get(key, R.zero()) + SR(coeff) * R.monomial(*exp_vec)
    return sectors
\end{lstlisting}

That five-line loop is the payoff of every design decision so far. The result is
\code{ft.\_by\_tp}, a dictionary \code{\{(n\_tilde, n\_phys): polynomial\}}. The
accessors on \code{FieldTheory} read directly off it:

\begin{itemize}
  \item \code{free\_action()} (\file{field\_theory.py:1010}) returns the $(1,1)$
        sector --- the bilinear kernel whose inverse is the free propagator
        (Chapter~\ref{ch:propagator}).
  \item \code{noise\_kernel()} (\file{field\_theory.py:1015}) returns all
        $(\ge 2,0)$ sectors --- the noise sources.
  \item \code{vertices()} (\file{field\_theory.py:1021}) returns all
        total-degree-$\ge 3$ sectors --- the interaction vertices.
  \item \code{sectors()} (\file{field\_theory.py:1027}) returns the full nonzero
        bigrade dictionary; \code{summary()} pretty-prints each labeled sector.
\end{itemize}

The human-readable labels come from the static helper \code{\_sector\_label}
(\file{field\_theory.py:1380}), which is itself a compact transcription of
Table~\ref{tab:bigrade}.

\begin{example}
\textbf{The OU quartic, classified.} With $\phi=\varepsilon x^3$ and white noise
$-D\tilde x^2$, expanding the action and classifying gives a tiny \code{\_by\_tp}.
The $(1,1)$ sector is $\tilde x\,(\partial_t+\mu)x$ --- the free OU kernel. The
$(2,0)$ sector is $-D\tilde x^2$ --- the single noise source. And the
total-degree-4 $(1,3)$ sector is $\varepsilon\,\tilde x\,x^3$ --- the lone
interaction vertex, with one response leg and three physical legs. That one vertex
is the source of every loop diagram in the quickstart's $C_{xx}(\tau)$ calculation.
\end{example}

\section{Mean-field sector verification}\label{sec:mf}

Three of the bigrade sectors --- $(0,0)$, $(1,0)$, $(0,1)$ --- are special. They are
the constant, the response-linear ``tadpole,'' and the physical-linear
``equation-of-motion residual.'' Collectively they are the \term{mean-field (saddle)
sector}, and at the correct saddle point they must \emph{vanish}. The framework does
not solve the saddle equations itself --- that is a downstream module --- but it does
two things here: it \emph{verifies} this sector vanishes under the model's declared
saddle substitutions, and then it forcibly \emph{zeroes} it.

\subsection{Why classify before substituting}

This is the single most important ordering constraint in \code{expand}, and it is
easy to get backwards. The mean-field substitutions are rewrites like
$v^{*}\to (E+\sum_j w_{ij}g\cdot n^{*}_j)/(1+\tau\,n^{*})$. If you applied them to the
\emph{whole} action before classifying, that rewrite would fire \emph{everywhere}
$v^{*}$ appears --- including inside the $(1,1)$ propagator kernel and the $\ge 2$
interaction vertices, where $v^{*}$ and $n^{*}$ must remain \emph{free symbolic
parameters} to be bound numerically much later. You would inject saddle-equation
algebra into the propagator. So the code coerces and classifies first
(\file{field\_theory.py:921--922}), and applies the saddle substitutions \emph{only}
to the three mean-field sectors, purely as a test. The long comment at
\file{field\_theory.py:913--922} states this explicitly.

\subsection{The two-tier check}

\code{\_verify\_and\_zero\_mf\_sector} (\file{field\_theory.py:585}) applies the
saddle substitutions to each monomial coefficient in the three sectors and checks it
vanishes, in two tiers:

\begin{enumerate}
  \item \textbf{Symbolic.} Substitute the mean-field conditions (then the
        specializations), and if \code{.simplify\_full() == 0}, the term passes.
  \item \textbf{Numerical fallback.} Only if the symbolic check fails, evaluate the
        residual numerically at a representative parameter point via
        \code{\_mf\_numerical\_residual}; pass if its magnitude is below a tolerance
        (default $10^{-9}$), recording a \emph{soft pass} (warned, not failed).
\end{enumerate}

Any term that passes neither tier raises an \code{AssertionError} naming the
offending bigrade, its post-substitution symbolic form, and its numerical residual.
The error message is blunt: \emph{``Either the MF solver is wrong or the action's
bigrade-$\le 1$ sector is not the saddle-eq sector.''} This assertion is a genuine
structural safety net --- it catches a miswired mean-field solver or a mis-specified
action before any diagrams are computed. After the checks, each mean-field sector
key is set to \code{R.zero()}.

\begin{note}
The numerical fallback is non-trivial machinery in its own right.
\code{\_mf\_numerical\_residual} (\file{field\_theory.py:662}) builds a representative
parameter point (preferring each parameter's declared \code{default}, since theories
ship defaults that admit a well-behaved saddle), then prefers the differential-algebraic
mean-field solver when the model declares \code{equations} --- the legacy iterative
solver cannot handle self-referential implicit saddles like
$x^{*}=-\varepsilon (x^{*})^3$ --- and falls back to the simple solver otherwise. The
companion \code{\_default\_fundamental\_point} (\file{field\_theory.py:730}) warns in
its docstring that an all-ones fallback is a poor universal default for a nonlinear
closure (for $\phi=av^2$ with $E=a=w=1$ the saddle equation $2v^2-v+1=0$ has no real
root), which is exactly why declared defaults are preferred.
\end{note}

\begin{gotcha}
Specializations run \emph{twice}. Once after the Taylor (step~7, on every sector),
and again \emph{inside} the verifier (\file{field\_theory.py:949}). The reason: the
\code{mf\_bg} lambda builds its right-hand sides \emph{before} specializations have
run, so a substitution like $v^{*}\to(\dots\phi_0\dots)$ reintroduces raw
$\phi$-Taylor tokens that need a second specialization pass to resolve.
\end{gotcha}

After verification the action is rebuilt from the now-zeroed \code{by\_tp} into
\code{self.\_S\_raw}, and \code{self.\_by\_tp} is the authoritative output. The quick
public sanity check \code{sanity\_check()} (\file{field\_theory.py:966}) simply
confirms those three sectors are \code{R.zero()} and prints a PASS/FAIL table --- and
it always prints a FAIL, even when called with \code{verbose=False}, so a failure is
never silent.

\section{Correlated-noise cumulants}\label{sec:cumulant}

This section covers step~4, which is a no-op for most models and so can be skipped on
a first read. White noise (the OU model) contributes a single local term
$-D\tilde x^2$ that lives entirely in the $(2,0)$ sector; nothing special is needed.
But \emph{colored} or \emph{correlated} external noise contributes a whole series to
the action: the cumulant-generating functional $-W_m[\tilde m]$,
\[
  S_{\mathrm{cum}} \;=\; -\sum_n \frac{1}{n!}\int \dd t_1\cdots\dd t_n
     \sum_{i_1\cdots i_n} \kappa^{(n)}_{i_1\cdots i_n}(\tau\text{'s})\,
     \tilde m_{i_1}(t_1)\cdots\tilde m_{i_n}(t_n),
\]
where $\kappa^{(n)}$ is the $n$-th cumulant of the noise and the $\tilde m$ are the
response legs the noise couples to. \code{\_build\_cumulant\_action}
(\file{field\_theory.py:293}) injects this series symbolically from the declarative
\code{model['correlated\_noises']} block, returning \code{SR(0)} when there is no such
block.

The key idea is that each per-leg kernel $K(\tau)$ is split into a
\emph{local (delta-correlated)} part and a \emph{smooth non-local} part:
\[
  K(\tau) \;=\; c_{\mathrm{local}}\,\delta(\tau) \;+\; K_{\mathrm{smooth}}(\tau).
\]
The framework peels off the local part by repeatedly taking the coefficient of
$\delta(\tau_k)$:

\begin{lstlisting}
# field_theory.py:460-463  (inside _build_cumulant_action)
c_local = K
for tau_k in tau_syms:
    try:
        c_local = c_local.coefficient(dirac_delta(tau_k))
\end{lstlisting}

The local part collapses every time integral and is injected as an ordinary monomial
(using the \code{is\_trivial\_zero()} test from the earlier gotcha, at
\file{field\_theory.py:490}). The smooth residual, by contrast, cannot collapse: it
introduces a placeholder symbol \code{z\_kappa\_...} (the \code{z} prefix again for
canonical ordering) and registers the actual kernel function in
\code{ns.\_cumulant\_kernels} for the downstream integrator to pick up. The injector
also handles \emph{cross-field} cumulants by enumerating distinct permutations of the
leg-field tuple --- so a noise correlating two different fields $[x_t,y_t]$ sums over
both orderings.

\begin{gotcha}
The smooth (non-$\delta$) residual is only handled at cumulant order $n=2$. At order
$n\ge 3$, a smooth residual is \emph{silently dropped} with a \code{warnings.warn}
(the local $\delta$ part of the same kernel, if any, is still injected correctly).
The reason is purely a limitation of the time-domain integrator downstream, which
currently allocates one relative-time variable per noise vertex and would need an
$n$-leg time map to handle higher non-local cumulants (\file{field\_theory.py:550--565}).
For the Gaussian-and-shifts noise models this machinery was built for, all
order-$\ge 3$ cumulants happen to be fully local, so the dropped branch never fires
in practice --- but it is a real ceiling for future non-local higher-order kernels.
\end{gotcha}

\section{Typed vertices and sources}\label{sec:typed}

We now cross from \file{field\_theory.py} into \file{vertices.py}. The bigrade
sectors are polynomials; the diagram machinery wants \emph{individual typed records},
one per monomial, each tagged with its legs. That conversion is what
\file{vertices.py} does.

\subsection{Leg tuples and the record types}

A \term{leg} is one attachment point of a vertex or source, tagged
\code{(field\_base\_name, population\_index)} --- for example \code{('dv', 1)} for the
first population's $\delta v$ field. The population index is 1-based, or $0$ when the
field is not indexed. A monomial's exponent in a given generator becomes that many
repeated leg entries.

The records themselves are lightweight classes that use \code{\_\_slots\_\_} (a
Python optimization that removes the per-instance \code{\_\_dict\_\_}, saving memory
and forbidding stray attributes). There are four:

\begin{description}
  \item[\code{VertexType}] (\file{vertices.py:80}) --- one interaction monomial
        (total degree $\ge 3$). Carries \code{coefficient} (the SR coupling times
        combinatorial prefactor), \code{response\_legs}, \code{physical\_legs}, and
        \code{bigrade}. Its properties \code{in\_degree} (number of physical legs),
        \code{out\_degree} (number of response legs), and \code{total\_degree} are
        the diagram-facing vocabulary.
  \item[\code{ConvVertexType}] (\file{vertices.py:134}) --- a \code{VertexType}
        subclass for a conductance-style vertex, where one or more physical legs sit
        at independent times linked to the vertex by a synaptic kernel $g(\tau)$. It
        adds \code{kernel\_attachments}: a list of dicts, each recording the kernel
        \code{'symbol'}, the \code{'leg'} it attaches to, that leg's 0-based
        \code{'leg\_index'}, and a \code{'kernel\_td\_fn'} callable
        $\tau\mapsto g(\tau)$. This is where the kernel attachments from
        \S\ref{sec:conv} are finally consumed.
  \item[\code{SourceType}] (\file{vertices.py:216}) --- one noise-kernel monomial
        ($n_{\tilde{}}\ge 2$, $n_{\mathrm{phys}}=0$). Carries \code{coefficient},
        \code{response\_legs}, \code{bigrade}, and an \code{out\_degree} property.
  \item[\code{NoiseSourceType}] (\file{vertices.py:256}) --- a \code{SourceType}
        subclass backed by a non-local cumulant kernel $\kappa^{(n)}$, whose response
        legs sit at independent times. It adds \code{cumulant\_specs}, the records
        that tie each placeholder symbol back to its kernel function. Locally
        correlated ($\delta$) cumulants stay plain \code{SourceType}.
\end{description}

The bridge from polynomial to records is \code{decompose\_sector}
(\file{vertices.py:349}). For every \code{(exponent\_vector, coefficient)} in a
sector's \code{.dict()}, it parses each nonzero generator's name into a leg tuple and
appends it (to the response legs if the generator index is below \code{n\_tilde}, else
to the physical legs), repeated by the exponent. If there are no physical legs it
emits a \code{SourceType}; otherwise a \code{VertexType}:

\begin{lstlisting}
# vertices.py:386-393  (inside decompose_sector)
n_t = len(resp_legs)
n_p = len(phys_legs)
bigrade = (n_t, n_p)
if n_p == 0:
    results.append(SourceType(SR(coeff), resp_legs, bigrade))
else:
    results.append(VertexType(SR(coeff), resp_legs, phys_legs, bigrade))
\end{lstlisting}

The generator-name parsing is done by \code{\_parse\_field\_name}
(\file{vertices.py:327}): it walks backward over trailing digits to split a name like
\code{dn2} into base \code{dn} and population index $2$ (returning index $0$ when
there are no trailing digits).

The two top-level extractors \code{extract\_vertex\_types} (\file{vertices.py:500})
and \code{extract\_source\_types} (\file{vertices.py:629}) wrap
\code{decompose\_sector} and then do the kernel/cumulant upgrade: scanning each
coefficient for kernel symbols, resolving which leg each attaches to, and promoting
the plain record to a \code{ConvVertexType} or \code{NoiseSourceType} when an
attachment resolves.

\begin{gotcha}
If a kernel symbol cannot be resolved to a specific leg, \code{extract\_vertex\_types}
deliberately \emph{leaves it in the coefficient} and emits a plain \code{VertexType}
(\file{vertices.py:589--595}). The point is fail-loud: a surviving kernel symbol in a
coefficient will trip a downstream diagnostic, which is far better than silently
treating the unresolved attachment as zero.
\end{gotcha}

\subsection{What the enumerator actually asks for}

The diagram enumerator does not consume the full vertex objects --- it only needs to
know which vertex \emph{shapes} exist. That is the entire job of
\code{available\_degrees} (\file{vertices.py:777}):

\begin{lstlisting}
# vertices.py:777
def available_degrees(vertex_types, source_types):
    interaction_degrees = {(vt.in_degree, vt.out_degree) for vt in vertex_types}
    source_degrees      = {st.out_degree for st in source_types}
    return interaction_degrees, source_degrees
\end{lstlisting}

It returns a set of \code{(in\_degree, out\_degree)} pairs from the vertices and a set
of \code{out\_degree} integers from the sources. The diagram enumerator
(\file{msrjd/enumeration/degree\_scan.py}) reads exactly this to know which
``shapes'' of vertex and source it may attach when it builds graphs --- the topic of
Chapter~\ref{ch:enumeration}.

\section{Picklability and the process boundary}\label{sec:pickle}

A short but load-bearing section, because it explains design choices in
\file{vertices.py} that otherwise look arbitrary. The downstream time-domain
integrator (Phase~J) can farm diagram evaluation across a
\code{multiprocessing.Pool}, which means the vertex and source objects may be
\emph{pickled} and shipped to worker processes. This subsystem never imports
\code{multiprocessing}, but it is engineered around that possibility. Three
consequences:

\begin{itemize}
  \item \code{\_NamespaceBoundKernel} (\file{vertices.py:25}) is a \emph{module-level}
        class, not a closure, so \code{pickle} can locate it by its fully-qualified
        name. It \emph{pre-evaluates} the user's kernel function once against the
        namespace at construction time and caches the resulting \code{SR} expression
        (\file{vertices.py:56}) --- because the full namespace is unpicklable (it
        holds inner closures), but a plain \code{SR} expression pickles cleanly.
  \item The record classes implement explicit
        \code{\_\_getstate\_\_}/\code{\_\_setstate\_\_} because \code{\_\_slots\_\_}
        classes have no \code{\_\_dict\_\_} for pickle to grab automatically.
  \item In a \code{\_\_slots\_\_} inheritance chain, the base class's pickle methods
        must reference the base \code{\_\_slots\_\_} \emph{by name}.
        \code{SourceType.\_\_getstate\_\_} therefore writes
        \code{SourceType.\_\_slots\_\_} explicitly, not \code{self.\_\_slots\_\_} ---
        on a \code{NoiseSourceType} instance the latter resolves only to the
        subclass's own \code{('cumulant\_specs',)}, which would drop the base fields
        from the pickle (\file{vertices.py:235--238}).
\end{itemize}

\begin{gotcha}
Because \code{\_NamespaceBoundKernel} pre-evaluates \emph{eagerly} at extraction time,
a kernel function that errors will throw during \code{extract\_source\_types}, not
later during integration. That is usually a feature (you find out early), but it can
surprise you if you expected lazy evaluation.

A project-wide caveat, recorded in the framework's memory: fork-based multiprocessing
has crashed this user's macOS machine. The spatial path is serial-only and the
temporal path is fork-guarded. The picklability machinery here is what \emph{allows}
these objects to cross a (guarded) process boundary at all, but the guard itself lives
elsewhere in the pipeline.
\end{gotcha}

\section{Serialization}

The last file, \file{serialize.py}, persists an expanded theory so the expensive
multivariate Taylor step is not repeated every run. \code{save\_theory}
(\file{serialize.py:52}) writes a directory with two artifacts:

\begin{itemize}
  \item \file{metadata.json} --- plain-Python metadata: format and Sage version, a
        timestamp, the model identity (\code{model\_name}, \code{model\_file},
        \code{model\_var\_name}), the expansion info (\code{taylor\_order},
        \code{n\_tilde}, \code{ring\_var\_names}), the field/parameter/kernel specs
        with callables stripped, and a sorted list of nonzero bigrade keys for quick
        inspection.
  \item \file{symbolic\_data.sobj} --- a Sage \emph{binary} object dump (via
        \code{sage\_save}) of the polynomial ring, the raw action, the bigrade
        dictionary, \code{n\_tilde}, and the propagator matrices. Sage's \code{.sobj}
        format can pickle arbitrary Sage objects --- polynomial rings, symbolic
        expressions, matrices --- that plain \code{pickle}/\code{json} cannot.
\end{itemize}

\begin{gotcha}
Two deliberate omissions. First, the model dict itself (with its lambdas) is
\emph{not} serialized; instead the metadata records the model's file path and
variable name so it can be re-imported. Second, the \code{VertexType}/\code{SourceType}
objects are \emph{not} saved either --- they are cheap to re-extract from the bigrade
dictionary on load, and they carry namespace-derived callables that are better
reconstructed than persisted.
\end{gotcha}

\code{load\_theory} (\file{serialize.py:156}) returns the \code{(metadata, data)} pair,
and \code{reload\_model} (\file{serialize.py:192}) re-imports the model dict from the
stored path. The diagram enumerator uses exactly this round-trip to re-expand a saved
theory at a \emph{higher} Taylor order on demand, when a particular diagram needs more
vertex legs than the cached expansion provides.

\begin{gotcha}
A trap for the Sage-literate. The comment at \file{serialize.py:232} says ``Use
SageMath's \code{load()} which handles .py files,'' but the body actually uses a plain
Python \code{exec(compile(code, full\_path, 'exec'), ns)}. The comment is stale. The
practical consequence: this \code{exec} does \emph{not} run Sage's preparser, so a
model file relying on Sage-only syntax (\code{\^{}} for power, a bare \code{2/3} meaning
an exact rational) would not be preparsed. Every model in this repository imports from
\code{sage.all} and uses plain Python syntax, so it works --- but do not author a model
file in Sage-preparser dialect and expect \code{reload\_model} to cope.
\end{gotcha}

\section{Where this sits in the pipeline}

Step back and place the whole subsystem. It is the front end --- phase \code{[1/7]}
--- and its job is one transform:
\[
  \underbrace{\text{model dict}}_{\text{declarative}}
  \;\xrightarrow{\ \text{\code{FieldTheory(...).expand()}}\ }\;
  \underbrace{\code{ft.\_by\_tp}}_{\{(n_{\tilde{}},n_p)\,:\,\text{polynomial}\}}.
\]

What \emph{feeds} it is the model dict, authored via the theory builders of
Chapter~\ref{ch:theory-spec}. The main driver \file{pipeline/compute.py} makes the
entry call \code{FieldTheory(model, taylor\_order).expand()} as its first phase.

What \emph{consumes} it splits three ways, all reading off \code{ft.\_by\_tp}:

\begin{itemize}
  \item \code{ft.free\_action()} (the $(1,1)$ sector) feeds the \textbf{propagator
        extractor} of Chapter~\ref{ch:propagator}, which inverts the bilinear kernel
        to get the free propagator $\Gret$.
  \item \code{extract\_source\_types(ft)} and \code{extract\_vertex\_types(ft)} feed
        \code{available\_degrees}, which feeds the \textbf{diagram enumerator} of
        Chapter~\ref{ch:enumeration}.
  \item The full \code{VertexType}/\code{SourceType}/\code{ConvVertexType}/\code{NoiseSourceType}
        records feed the \textbf{integrators} (temporal and spatial) of Parts~IV
        and~V, which finally turn the diagrams into numbers.
\end{itemize}

Everything else in this manual builds on the classified action you have just watched
get constructed. With the free action in hand, the next chapter inverts it to obtain
the propagator.

% ---------------------------------------------------------------------
%  Chapter 7 : The Propagator
% ---------------------------------------------------------------------
\chapter{The Propagator}\label{ch:propagator}

Every Feynman diagram in this framework is a graph whose \emph{edges} are
copies of one object: the free two-point function, the \term{propagator}
$\Gret$. Wire vertices together with copies of $\Gret$, integrate over the
internal times, and you have a diagram. Nothing downstream --- not the
enumeration, not the type assignment, not a single loop integral --- can run
until $\Gret$ exists. This chapter is about how the engine builds it.

The story has a clean shape. The Gaussian (free) part of the \msrjd{} action is
a bilinear form coupling each response field to each physical field; collect its
coefficients into a matrix $K$ (the ``kinetic'' or ``kernel'' matrix), and the
propagator is simply its inverse $G = K^{-1}$. Because $K$ contains time
derivatives, the natural place to invert it is \emph{frequency space}: Fourier
transform in time to get $K_{\mathrm{ft}}(\omega)$, invert the small matrix once,
and read off a rational function of $\omega$. The physics is then entirely in
that rational function's \emph{poles} (which set relaxation rates), its
\emph{residues} (which set amplitudes), and one \emph{instantaneous} piece that
becomes a $\delta(t)$ in the time domain. The closed-form inverse Fourier
transform reassembles those three ingredients into the time-domain retarded
propagator that the diagram integrators actually call.

This is phases \texttt{[2/7]} (\emph{Build propagator}) and \texttt{[4/7]}
(\emph{Numerical poles + residues}) of the seven-phase trace you watched in
Chapter~\ref{ch:quickstart}. We take them together because they are two halves
of one object: phase~2 builds the \emph{symbolic} frequency-domain propagator
(parameter-independent, cached to disk); phase~4 substitutes the numeric
parameters and extracts the poles and residues. We will build the mathematics
from the ground up, then walk the two source files that implement it ---
\file{pipeline/\_propagator.py} (the engine) and
\file{msrjd/integration/time\_domain/propagator\_td.py} (the time-domain
assembler) --- quoting real code with file and line numbers throughout.

\begin{gotcha}
There is a file named \file{msrjd/core/propagator.py}. It is a \emph{stub} --- a
ten-line docstring saying the logic ``currently lives in the demo notebook.''
Nothing in the working pipeline imports it. The real propagator code is
\file{pipeline/\_propagator.py}. Do not go looking for \code{build\_propagator}
in \file{msrjd/core/propagator.py}; it is not there.
\end{gotcha}

\section{Where the propagator sits}

It helps to fix the data flow before the mathematics. Two functions in
\file{pipeline/\_propagator.py} do the work, and they run at two different points
in \file{pipeline/compute.py}:

\begin{description}
  \item[\code{build\_propagator(ft, model, \ldots)}] (\file{pipeline/\_propagator.py:524})
        Called at phase~2 (\file{pipeline/compute.py:357}). Takes the
        \emph{expanded} \code{FieldTheory} \code{ft} and the \code{model} dict;
        builds the symbolic chain $K_{\mathrm{ker}} \to K_{\mathrm{ft}} \to
        G_{\mathrm{ft}}$, plus the adjugate, the determinant, and the
        instantaneous matrix \code{D\_delta}. It returns a \code{prop} dict with
        the poles and residues still set to \code{None} --- those are
        parameter-dependent and come later. The symbolic stages \emph{are}
        parameter-independent, so this dict is cached to disk.
  \item[\code{compute\_poles\_and\_residues(prop, num\_params, \ldots)}] (\file{pipeline/\_propagator.py:851})
        Called at phase~4 (\file{pipeline/compute.py:632}). Takes the \code{prop}
        dict and \code{num\_params} --- the numeric parameter substitution
        produced by the mean-field solve (phase~3). It substitutes those numbers,
        finds the retarded poles, computes the residue matrices, and
        \emph{mutates} \code{prop} in place to fill \code{prop['pole\_vals']} and
        \code{prop['C\_mats']}.
\end{description}

A third function lives in the other file. \code{build\_G\_t\_matrix}
(\file{msrjd/integration/time\_domain/propagator\_td.py:135}) turns the
frequency-domain data (poles, residues, \code{D\_delta}) into a \emph{time-domain}
matrix $G(t)$, and the per-edge lookups \code{G\_t\_entry} /
\code{G\_t\_delta\_coeff} hand single propagator entries to the diagram
integrators. Schematically:

\begin{lstlisting}[style=console]
  FieldTheory (expanded action)        model dict (params, kernel images)
           |                                       |
           +------------------+--------------------+
                              v
  build_propagator():   K_ker -> K_ft -> G_ft / adj_ft / D_omega / D_delta
                              |   (cached to saved_theories/<tag>/propagator.sobj)
                              v
  compute_poles_and_residues(prop, num_params):
                              fills prop['pole_vals'], prop['C_mats'], prop['D_delta']
                              v
  build_G_t_matrix(prop, t):  {'smooth': sum_k C_k exp(i p_k t),  'delta': D_delta,  't_var'}
                              v
  per-edge   G_t_entry()  /  G_t_delta_coeff()
                              v
  diagram tree integrators (final_integral.py, grouped_integral.py)
\end{lstlisting}

\begin{note}
There is also a precompute entry point at \file{pipeline/\_precompute.py:168}
that calls \code{build\_propagator} at Taylor order~2 purely to warm the cache.
This works because the propagator depends \emph{only} on the bilinear sector of
the action, which is fully captured at Taylor order $\ge 2$ --- so a single
order-2 precompute fills the cache once and every higher-order run reuses it.
The propagator is, in this precise sense, Taylor-order-independent.
\end{note}

\section{From the action to the kernel matrix \texorpdfstring{$K$}{K}}

A stochastic field theory written in \msrjd{} form has an action $S[\varphi,
\tilde\varphi]$. The fields come in two halves of equal count $n_f$:

\begin{itemize}
  \item \term{physical fields} $\varphi_j$, $j = 1,\dots,n_f$ (a density, a
        membrane potential, \ldots),
  \item \term{response fields} $\tilde\varphi_i$, $i = 1,\dots,n_f$, also called
        \emph{auxiliary} or \emph{tilde} fields --- the conjugate variables that
        the \msrjd{} construction introduces, one per physical field.
\end{itemize}

The \term{free action} is the part of $S$ that is bilinear with exactly one
response leg and one physical leg. In the engine's bookkeeping this is
\emph{bigrade $(1,1)$} --- one tilde factor, one physical factor --- and it is
exactly what \code{ft.free\_action()} returns. Schematically,
\begin{equation}
  S_{\mathrm{free}} \;=\; \sum_{i,j}\int\!\dd t\;
  \tilde\varphi_i(t)\,\hat K_{ij}(\Dt)\,\varphi_j(t),
  \label{eq:freeaction}
\end{equation}
where each $\hat K_{ij}$ is a differential/kernel operator in time. For a plain
Ornstein--Uhlenbeck field $\hat K = \Dt + \mu$, so $S_{\mathrm{free}} \supset
\int \tilde n\,(\Dt n + \mu n)$.

\begin{defn}
The \term{kernel matrix} $K$ is the matrix of bilinear coefficients of
\eqref{eq:freeaction}, with layout \emph{[response row, physical column]}:
\[
  K[i, j] \;=\; \text{coefficient of }(\tilde\varphi_i\,\varphi_j)
  \text{ in } S_{\mathrm{free}}.
\]
The propagator is its inverse, $G = K^{-1}$.
\end{defn}

The code extracts $K$ by reading the coefficient of each $(\tilde\varphi_i\,
\varphi_j)$ monomial out of the free action. \code{build\_propagator} first asks
the \code{FieldTheory} for the free action and the polynomial-ring generator
names, then splits the names into the response half and the physical half
(\file{pipeline/\_propagator.py:586}):

\begin{lstlisting}
S_free = ft.free_action()
ring_gen_names = [str(g) for g in R.gens()]

resp_names = ring_gen_names[:ft._n_tilde]
phys_names = ring_gen_names[ft._n_tilde:]
pos_to_row = {ring_gen_names.index(nm): i for i, nm in enumerate(resp_names)}
pos_to_col = {ring_gen_names.index(nm): j for j, nm in enumerate(phys_names)}
\end{lstlisting}

The first \code{ft.\_n\_tilde} generators are the response fields and the rest are
the physical fields; \code{pos\_to\_row} and \code{pos\_to\_col} map a generator's
position in the ring to its row (if it is a response field) or column (if it is
physical). The extraction loop then walks the \emph{monomials} of the free action
--- \code{S\_free.dict()} returns a mapping from exponent vector to coefficient ---
and for each monomial finds which generators carry a positive exponent
(\file{pipeline/\_propagator.py:595}):

\begin{lstlisting}
nf = len(resp_names)
K_data = [[SR(0)] * nf for _ in range(nf)]
for exp_vec, coeff in S_free.dict().items():
    row = col = None
    for idx in range(len(ring_gen_names)):
        if exp_vec[idx] > 0:
            if idx in pos_to_row:
                row = pos_to_row[idx]
            if idx in pos_to_col:
                col = pos_to_col[idx]
    if row is not None and col is not None:
        K_data[row][col] += SR(coeff)
K_mat = matrix(SR, K_data)
\end{lstlisting}

A bilinear $(1,1)$ monomial has exactly one response generator and one physical
generator with positive exponent, so it pins down exactly one $(row, col)$, and
its coefficient is accumulated into $K[\,row, col\,]$. The result \code{K\_mat} is
a Sage symbolic matrix of size $n_f \times n_f$.

\begin{note}
\code{SR} is SageMath's \term{Symbolic Ring} --- the general symbolic-expression
type, a Python wrapper around the GiNaC/Pynac C++ engine. \code{SR(x)} coerces
anything (an integer, a coefficient, a string) to a symbolic expression, and
\code{matrix(SR, data)} builds a matrix whose entries live in that ring. Almost
every matrix in this chapter lives over \code{SR}. We meet several more Sage
rings below; each is introduced where it first appears.
\end{note}

\section{Kernel form: the \texorpdfstring{$\delta$}{delta} and \texorpdfstring{$\delta'$}{delta-prime} symbols}

We want to Fourier transform $K$ in time, but a raw coefficient like $\Dt n$ (a
time derivative) does not Fourier-transform cleanly as written. The engine first
rewrites every entry of $K$ into \term{kernel form}, a sum of a Dirac delta and
its derivative:
\begin{equation}
  c \;\longmapsto\; c_0\,\delta(t) \;+\; c_1\,\delta'(t),
  \qquad c_0 = c\big|_{\Dt \to 0}.
  \label{eq:kernelform}
\end{equation}
There are three abstract namespace symbols involved, all carried on the
\code{FieldTheory}'s namespace object \code{ns}:

\begin{itemize}
  \item \code{ns.Dt} --- a marker for the time-derivative operator $\Dt$,
  \item \code{ns.delta\_D} --- the Dirac delta $\delta(t)$ (displayed as $\delta$),
  \item \code{ns.delta\_Dp} --- its derivative $\delta'(t)$ (displayed as $\delta'$).
\end{itemize}

The helper \code{\_to\_kernel} (\file{pipeline/\_propagator.py:507}) performs the
rewrite \eqref{eq:kernelform}:

\begin{lstlisting}
def _to_kernel(c, Dt, delta_D, delta_Dp):
    c = SR(c)
    if c.has(delta_D):
        return c
    p0   = c.subs({Dt: 0})
    rest = (c - p0).subs({Dt: delta_Dp})
    return p0 * delta_D + rest
\end{lstlisting}

It splits the coefficient into a constant part \code{p0} (set $\Dt \to 0$) and a
derivative part \code{rest} (the remainder, with $\Dt$ replaced by the symbol
$\delta'$), and returns \code{p0*delta\_D + rest}. So a plain constant $\mu$
becomes $\mu\,\delta(t)$ --- an instantaneous coupling --- and a derivative $\Dt$
becomes $\delta'(t)$.

Why wrap the \emph{constant} part in $\delta(t)$ at all? Because of what the
Fourier transform does to a bare constant. The transform of $\delta(t)$ is the
constant $1$, but the transform of a bare constant $\mu$ is the distribution
$2\pi\,\mu\,\delta(\omega)$ --- not what we want. Writing $\mu$ as $\mu\,\delta(t)$
ensures the transform returns the clean constant $\mu$ back (the docstring at
\file{pipeline/\_propagator.py:514} spells this out). Abstract \emph{kernel}
symbols (for example a synaptic kernel \code{ns.g} in a coloured-noise model) are
not constants and are left untouched here; their frequency images are applied
later, after the transform.

\code{build\_propagator} applies \code{\_to\_kernel} entrywise to build
\code{K\_ker} (\file{pipeline/\_propagator.py:615}):

\begin{lstlisting}
Dt       = ns.Dt
delta_D  = ns.delta_D
delta_Dp = ns.delta_Dp

K_ker = matrix(
    SR,
    [[_to_kernel(K_mat[i, j], Dt, delta_D, delta_Dp)
      for j in range(nf)] for i in range(nf)],
)
\end{lstlisting}

\section{The Fourier convention --- and why poles live in the upper half-plane}

The engine uses a single, fixed Fourier convention pipeline-wide, defined in
\code{fourier\_transform} (\file{msrjd/core/field\_theory.py:33}). It is the
angular-frequency convention with no $2\pi$ in the exponent:
\begin{equation}
  F(\omega) \;=\; \int_{-\infty}^{\infty} f(t)\,\ee^{-\ii\omega t}\,\dd t.
  \label{eq:ftconv}
\end{equation}
Under \eqref{eq:ftconv} the kernel symbols transform as
\[
  \delta(t) \;\longmapsto\; 1, \qquad
  \delta'(t) \;\longmapsto\; \ii\omega,
\]
so a kernel-form entry $c_0\,\delta + c_1\,\delta'$ transforms to $c_0 + c_1\,
(\ii\omega)$. For the OU example $\hat K = \Dt + \mu$ becomes $K_{\mathrm{ft}} =
\ii\omega + \mu$.

\code{build\_propagator} performs the transform by substituting the abstract
$\delta$ symbols for their concrete Sage forms --- \code{dirac\_delta(t)} for
$\delta$ and \code{diff(dirac\_delta(t), t)} for $\delta'$ --- and calling
\code{fourier\_transform} on each nonzero entry
(\file{pipeline/\_propagator.py:624}):

\begin{lstlisting}
t_var = SR.var('t')
omega = SR.var('omega', latex_name=r'\omega')
time_subs = {
    delta_D:  dirac_delta(t_var),
    delta_Dp: diff(dirac_delta(t_var), t_var),
}
K_ft_data = [[SR(0)] * nf for _ in range(nf)]
for i in range(nf):
    for j in range(nf):
        c = K_ker[i, j]
        if not c.is_zero():
            K_ft_data[i][j] = fourier_transform(
                SR(c).subs(time_subs), t_var, omega
            )
K_ft = matrix(SR, K_ft_data)
\end{lstlisting}

After the transform, any abstract kernel symbol still present (the
coloured-noise case) is replaced by its frequency image $\hat g(\omega)$ via a
model-supplied hook (\file{pipeline/\_propagator.py:642}):

\begin{lstlisting}
kft_hook = model.get('kernel_ft_image')
if kft_hook is not None:
    kft_subs = kft_hook(ns, omega)
    K_ft = K_ft.apply_map(lambda e: SR(e).subs(kft_subs))
\end{lstlisting}

The \code{kernel\_ft\_image} hook is a callable the theory file supplies; it maps
abstract kernel symbols to their frequency images (for a white-noise model there
are none, and the hook is absent).

\begin{defn}
\code{dirac\_delta(t)} is Sage's symbolic Dirac delta; \code{diff(expr, var)} is
symbolic differentiation, so \code{diff(dirac\_delta(t), t)} is the symbolic
$\delta'(t)$. \code{fourier\_transform} is the engine's own wrapper
(\file{msrjd/core/field\_theory.py}); internally it dispatches the integral to
SymPy first (which handles \code{positive=True} assumptions on time constants
without stopping to ask the sign of $\omega$) and falls back to Maxima. SymPy is
therefore on the hot path even though no \code{import sympy} appears in the
propagator files.
\end{defn}

\begin{gotcha}
\textbf{Retarded $=$ $\mathrm{Im}(\omega) > 0$ is the single most important sign
convention in this subsystem, and it is born directly from the $\ee^{-\ii\omega
t}$ in \eqref{eq:ftconv}.} Because the forward transform carries
$\ee^{-\ii\omega t}$, the \emph{inverse} transform carries $\ee^{+\ii\omega t}$
and closes its $\omega$-contour in the \emph{upper} half-plane for $t > 0$. So the
retarded (causal) propagator is parameterised by poles with $\mathrm{Im}(\omega) >
0$: each $\ee^{\ii\omega_k t}$ then decays for $t > 0$ and grows for $t < 0$.
Every pole filter in the code keeps roots with \code{imag > 1e-9}. If you ever
flipped \code{fourier\_transform} to $\ee^{+\ii\omega t}$, every pole-side sign
would flip and the entire downstream causality machinery would break. This is
restated in the header of \file{propagator\_td.py:12}.
\end{gotcha}

\section{The frequency-domain propagator \texorpdfstring{$G_{\mathrm{ft}} = K_{\mathrm{ft}}^{-1}$}{G = K-inverse}}

With $K_{\mathrm{ft}}(\omega)$ in hand, the propagator is the matrix inverse
\begin{equation}
  G_{\mathrm{ft}}(\omega) \;=\; K_{\mathrm{ft}}(\omega)^{-1},
  \label{eq:Ginv}
\end{equation}
a matrix of \emph{rational functions of $\omega$}. Two companion objects come for
free and matter later:

\begin{itemize}
  \item the \term{determinant} $D_\omega = \det K_{\mathrm{ft}}$. Its numerator is
        the \emph{characteristic polynomial}, whose roots are the system poles.
  \item the \term{adjugate} $\mathrm{adj}(K_{\mathrm{ft}})$, the transpose of the
        cofactor matrix. By Cramer's rule $K^{-1} = \mathrm{adj}(K)/\det(K)$, so
        $\mathrm{adj}(K_{\mathrm{ft}}) = G_{\mathrm{ft}}\cdot\det(K_{\mathrm{ft}})$.
\end{itemize}

For a diagonally-dominant $K_{\mathrm{ft}}$ every entry of $G_{\mathrm{ft}}$
shares one common denominator $Q(\omega)$, the polynomial numerator of
$\det(K_{\mathrm{ft}})$.

\subsection{The complexity gate: when symbolic inversion is skipped}

Sage's symbolic \code{K\_ft.inverse()} is \emph{cofactor-based}: it expands all
$n_f^2$ cofactor sub-determinants symbolically, an $O(n_f!)$ operation in the
matrix dimension. For a small matrix (the OU model is $1\times1$; a two-field
system is $2\times2$) this is instantaneous. For a large matrix --- a spike-reset
model can reach $n_f = 8$ --- it can run for \emph{hours}, even when the symbol set
is lean, simply because $8! = 40\,320$ minor determinants must each be simplified
as rational symbolic expressions.

So \code{build\_propagator} gates the symbolic inversion on matrix size
(\file{pipeline/\_propagator.py:681}):

\begin{lstlisting}
free_syms = set()
for i in range(nf):
    for j in range(nf):
        free_syms.update(SR(K_ft[i, j]).variables())
free_syms.discard(omega)
rich = nf >= 6 or len(free_syms) > 20
\end{lstlisting}

A matrix is \term{rich} if $n_f \ge 6$ (the dominant trigger) \emph{or} if it
carries more than twenty free symbols beyond $\omega$ (a secondary trigger). The
comment at \file{pipeline/\_propagator.py:672} records the lesson that motivated
the matrix-size metric: spike-reset at $n_f = 8$ with only $\sim18$ free symbols
ran for hours under the old symbol-count gate, because the cost is $O(n_f!)$ in
the \emph{dimension}, not in the symbol count.

\begin{itemize}
  \item If the matrix is \emph{not} rich, the engine attempts the symbolic inverse
        under a wall-clock budget (next subsection).
  \item If it \emph{is} rich, symbolic inversion is skipped entirely:
        \code{G\_ft}, \code{adj\_ft}, \code{D\_omega}, and \code{D\_delta} are all
        left \code{None}, and \emph{everything} is deferred to phase~4, which
        works numerically after the parameters are substituted (at which point
        every entry is a univariate rational in $\omega$, far cheaper to handle).
\end{itemize}

\subsection{The budgeted symbolic-inverse block}

When the matrix is lean, the symbolic inverse, determinant, adjugate, and
\code{D\_delta} are all computed in one block, wrapped in a $120$-second SIGALRM
budget (\file{pipeline/\_propagator.py:706}):

\begin{lstlisting}
def _do_symbolic_inverse_block():
    t0 = _time.perf_counter()
    G    = K_ft.inverse()
    D_om = K_ft.det()
    # Adjugate via the identity  adj(K) = K^(-1) . det(K).
    adj  = G * D_om
    ...
    return G, D_om, adj, matrix(SR, Dd_data)
\end{lstlisting}

Two design choices in this tiny block are worth calling out, because each is a
hard-won performance fix.

\begin{gotcha}
\textbf{The adjugate is computed as \code{G * D\_om}, never via Sage's
\code{.adjugate()}.} Sage's \code{K\_ft.adjugate()} re-expands all $n_f^2$
cofactor sub-determinants from scratch, throwing away the work the inverse
already did. The comment at \file{pipeline/\_propagator.py:710} records the
measurement: $734$~s for \code{.adjugate()} versus $\sim11$~s for the inverse
itself on a $4\times4$ spike-reset matrix. Reusing the inverse via
$\mathrm{adj} = G\cdot\det$ is the whole point.
\end{gotcha}

If the block exceeds its budget --- or raises any exception, or aborts deep inside
a native routine with a \code{SignalError} --- the four symbolic stages are set to
\code{None} and phase~4 picks up the slack numerically
(\file{pipeline/\_propagator.py:731}):

\begin{lstlisting}
try:
    G_ft, D_omega, adj_ft, D_delta = _run_with_timeout(
        _do_symbolic_inverse_block, _SYMBOLIC_INVERSE_BUDGET_SEC)
except _PolynomialPathTimeout:
    ...
    G_ft = adj_ft = D_omega = D_delta = None
\end{lstlisting}

\code{\_run\_with\_timeout} (\file{pipeline/\_propagator.py:57}) is a small helper
that installs a POSIX \code{SIGALRM} handler raising \code{\_PolynomialPathTimeout}
after the budget expires, arms \code{signal.alarm}, runs the function, and cancels
the alarm in a \code{finally}. On a platform without \code{SIGALRM} (Windows) or
when the budget is \code{None}, it simply calls the function directly --- a no-op
timeout. This is the lower-level cousin of the \code{cysignals} alarm we meet
later for \code{factor()}.

The resulting symbolic stages (or their \code{None} placeholders) are assembled
into the \code{prop} dict, together with the as-yet-unfilled poles and residues
(\file{pipeline/\_propagator.py:765}):

\begin{lstlisting}
prop = {
    'K_ker':         K_ker,
    'K_ft':          K_ft,
    'G_ft':          G_ft,
    'adj_ft':        adj_ft,
    'D_omega':       D_omega,
    'D_delta':       D_delta,
    't_var':         t_var,
    'omega':         omega,
    'nf':            nf,
    'ring_gen_names': ring_gen_names,
    # Filled in later by compute_poles_and_residues():
    'pole_vals':     None,
    'C_mats':        None,
}
\end{lstlisting}

\section{Poles as zeros of \texorpdfstring{$\det K_{\mathrm{ft}}$}{det K}}

Phase~4 begins by substituting the numeric parameters. Once that is done, every
entry of $K_{\mathrm{ft}}$ is a univariate rational function in $\omega$ with
numeric coefficients, and the propagator's poles are the roots of the
characteristic polynomial $Q(\omega) = \operatorname{num}(\det K_{\mathrm{ft}})$
that lie in the upper half-plane.

\begin{defn}
A \term{pole} $\omega_k$ is a (retarded) zero of $\det K_{\mathrm{ft}}$ with
$\mathrm{Im}(\omega_k) > 0$. It sets a relaxation or oscillation rate of the free
theory. A \term{simple} pole is a single root; a higher-multiplicity (repeated)
root is a \emph{Jordan-block} pole and is handled differently --- see the
limitations section.
\end{defn}

Let $\omega_k$ be a simple zero of $Q(\omega)$. The residue of $G_{\mathrm{ft}}$
there is, entrywise,
\begin{equation}
  \operatorname*{Res}_{\omega=\omega_k} G_{\mathrm{ft}}(\omega)
  \;=\; \frac{P_{ij}(\omega_k)}{Q'(\omega_k)},
  \qquad
  G_{\mathrm{ft}}[i,j] = \frac{P_{ij}(\omega)}{Q(\omega)}.
  \label{eq:residue}
\end{equation}
The code stores the residue matrix \emph{scaled by $\ii$} --- the ``$C_k$
convention'' --- because that factor makes the inverse Fourier transform come out
clean:
\begin{equation}
  \texttt{C\_mats}[k] \;=\; \ii\cdot
  \operatorname*{Res}_{\omega=\omega_k} G_{\mathrm{ft}}(\omega).
  \label{eq:Ck}
\end{equation}

The residues are \emph{always} computed numerically, even when a symbolic
adjugate is available. This is not laziness; it is precision. The reason gets its
own gotcha below. First, the primary path.

\section{Residues, tier 1: exact arithmetic over \texorpdfstring{$\mathbb{Q}(i)[\omega]$}{Q(i)}}

\code{compute\_poles\_and\_residues} tries three strategies in order. The first,
\code{\_compute\_residues\_via\_polynomial\_fracfield}
(\file{pipeline/\_propagator.py:81}), is the one that runs for almost every
model. It does \emph{exact} arithmetic --- no floating point until the very last
root-find --- over the field $\mathbb{Q}(i)$.

\subsection{Why an exact field}

The trick is the choice of coefficient ring. Sage can invert a matrix over the
field of rational functions $\mathrm{Frac}(R[\omega])$, and when $R$ is an
\emph{exact} field, polynomial GCDs are reliable: Sage returns every entry of the
inverse in genuinely \emph{canonical irreducible} form $P_{ij}/Q$, with no
spurious cancellable factors and no floating-point GCD noise. The denominator is
then \emph{exactly} the characteristic polynomial.

The exact field used is \code{CyclotomicField(4, 'I\_')}, which is
$\mathbb{Q}(\zeta_4) = \mathbb{Q}(i)$ --- the rationals extended by a primitive
fourth root of unity, which is just the imaginary unit $i$.

\begin{defn}
\term{\code{CyclotomicField(4)}} $= \mathbb{Q}(i)$ is the field of Gaussian
rationals $a + b\,i$ with $a, b \in \mathbb{Q}$. \term{\code{PolynomialRing(R, 'x')}}
builds the ring $R[x]$ of univariate polynomials over $R$; calling
\code{.fraction\_field()} on it gives $\mathrm{Frac}(R[x])$, the rational
functions $P/Q$. The point of doing the matrix inversion over an exact field is
that GCD-over-a-field is exact, so each inverse entry is reduced to true
irreducible $P/Q$ form. A floating-coefficient (\code{CDF}$[\omega]$) path, by
contrast, once left ``degree-18 polynomials with $\sim13$ spurious roots''
(\file{pipeline/\_propagator.py:104}).
\end{defn}

\subsection{Walking the exact path}

The function builds the exact rings, substitutes the parameters, converts every
entry into the exact fraction field, and inverts under the same SIGALRM watchdog
(\file{pipeline/\_propagator.py:115}):

\begin{lstlisting}
CF = CyclotomicField(4, 'I_')
i_CF = CF.gen()
PR_exact = PolynomialRing(CF, 'om_exact')
F_exact  = PR_exact.fraction_field()

K_ft_num = K_ft.apply_map(lambda e: SR(e).subs(num_params))
\end{lstlisting}

The two converters are the heart of it. \code{\_sr\_complex\_to\_CF}
(\file{pipeline/\_propagator.py:122}) turns a symbolic complex \emph{constant}
into an exact $\mathbb{Q}(i)$ element by rationalising its real and imaginary
parts, and \code{\_sr\_to\_F\_exact} (\file{pipeline/\_propagator.py:128}) lifts a
whole rational entry into \code{F\_exact} by converting its numerator and
denominator coefficient lists:

\begin{lstlisting}
def _sr_complex_to_CF(c):
    c = SR(c)
    re = QQ(c.real_part())
    im = QQ(c.imag_part())
    return CF(re) + CF(im) * i_CF
\end{lstlisting}

Here \code{QQ} is Sage's field of rational numbers $\mathbb{Q}$; \code{QQ(x)}
rationalises a float. Once every entry is in \code{F\_exact}, the matrix is built
and inverted (\file{pipeline/\_propagator.py:140}):

\begin{lstlisting}
K_ft_F = matrix(F_exact, K_ft_F_data)
G_ft_F = _run_with_timeout(
    K_ft_F.inverse, _POLYNOMIAL_PATH_BUDGET_SEC)
\end{lstlisting}

\begin{gotcha}
The exact path \emph{requires} parameters that \code{QQ} can rationalise. The line
\code{re = QQ(c.real\_part())} throws on a non-rationalisable (for example,
irrational symbolic) parameter, the whole tier returns a triple of \code{None},
and the caller falls through to the numerical cofactor path
(\file{pipeline/\_propagator.py:122}). In practice parameters are floats, which
\code{QQ} rationalises happily, so the exact path is what runs.
\end{gotcha}

\subsection{The universal denominator must be an LCM}

After the inverse, the code needs the characteristic polynomial $Q(\omega)$. It
is tempting to read it off the first nonzero entry's denominator --- but that is
wrong for a structured $K_{\mathrm{ft}}$. The fix is to take the \emph{least
common multiple} of all the per-entry denominators
(\file{pipeline/\_propagator.py:172}):

\begin{lstlisting}
Q_poly = None
for i in range(nf):
    for j in range(nf):
        if G_ft_F[i, j] == 0:
            continue
        d_ij = G_ft_F[i, j].denominator()
        if Q_poly is None:
            Q_poly = d_ij
        else:
            Q_poly = Q_poly.lcm(d_ij)
\end{lstlisting}

\begin{gotcha}
\textbf{The universal denominator is an LCM, not the first entry's denominator.}
An upper-triangular $K_{\mathrm{ft}}$ --- the shape produced by the
Markovian-embedding preprocessor for coloured noise --- gives entries with
\emph{different} canonical denominators: one entry might carry a single pole
factor while another carries the product of two. Picking the first nonzero entry
would under-count the system poles. The LCM is the universal denominator whose
roots include every system pole; a pole that does not divide a given entry's own
denominator simply yields a vanishing residue there
(\file{pipeline/\_propagator.py:157}).
\end{gotcha}

\subsection{The squarefree guard}

Before root-finding, the code checks that $Q(\omega)$ is \emph{squarefree} --- no
repeated roots (\file{pipeline/\_propagator.py:196}):

\begin{lstlisting}
if not Q_poly.is_squarefree():
    if verbose:
        print(...)
    return None, None, None
\end{lstlisting}

A non-squarefree $Q$ signals a higher-multiplicity (Jordan-block) pole, for
example $Q = (\ii\omega + \mu)^2$ from a Markovianised OU process at $\mu =
1/\tau_c$. The single-pole residue formula \eqref{eq:residue} cannot represent the
$\tau\,\ee^{-pt}$ term such a pole produces, so the exact path bails and lets the
fallback tier take over (the fallback shares the limitation but is more robust to
the numerical near-degeneracy). This is the most significant \emph{known
limitation} of the subsystem; the limitations section returns to it.

\subsection{Roots and residues}

With a squarefree $Q$, the polynomial is converted from the exact field to
\code{CDF}$[\omega]$ (floating-point) for numerical root-finding, the retarded roots
are kept and sorted, and the residues are read off entry by entry.

\begin{defn}
\term{\code{CDF}} is Sage's \term{Complex Double Field}: complex numbers backed by
the machine \code{double} type ($\approx 10^{-16}$ precision). It is used for
root-finding (\code{Q\_cdf.roots(CDF)}) and to store the residue matrices as
\code{matrix(CDF, \ldots)}. \code{complex(value)} converts a \code{CDF} number back
to a plain Python \code{complex}.
\end{defn}

The retarded filter and sort (\file{pipeline/\_propagator.py:214}):

\begin{lstlisting}
roots = Q_cdf.roots(CDF)
pole_vals = [complex(r) for r, _ in roots
             if complex(r).imag > 1e-9]
pole_vals.sort(key=lambda p: (p.imag, p.real))
\end{lstlisting}

The residue loop then implements \eqref{eq:Ck} per entry, $C_k[i,j] = \ii\,
P_{ij}(\omega_k)/Q'_{ij}(\omega_k)$, skipping any entry for which $\omega_k$ is
not a pole of that entry's own denominator (\file{pipeline/\_propagator.py:246}):

\begin{lstlisting}
for pk in pole_vals:
    C_entries = [[0j] * nf for _ in range(nf)]
    for i in range(nf):
        for j in range(nf):
            ...
            qij_at = complex(Q_entry_cdf[i][j](pk))
            if abs(qij_at) > 1e-9:
                continue           # pk is not a pole of this entry
            qijp = complex(Q_entry_prime_cdf[i][j](pk))
            if abs(qijp) < 1e-30:
                continue           # higher-multiplicity: leave 0
            P_at = complex(P_cdf[i][j](pk))
            C_entries[i][j] = 1j * P_at / qijp
\end{lstlisting}

The function returns the triple \code{(pole\_vals, C\_mats, D\_delta)}
(\file{pipeline/\_propagator.py:301}); the third element is the instantaneous
matrix, computed in the same pass and discussed shortly. On success,
\code{compute\_poles\_and\_residues} stores the poles and residues and
\emph{returns early}, never touching the fallback tiers
(\file{pipeline/\_propagator.py:892}):

\begin{lstlisting}
pole_vals_tier1, C_mats_tier1, D_delta_tier1 = (
    _compute_residues_via_polynomial_fracfield(
        K_ft, omega, num_params, nf, verbose=verbose))
if pole_vals_tier1 is not None:
    prop['pole_vals'] = pole_vals_tier1
    prop['C_mats']    = C_mats_tier1
    if prop.get('D_delta') is None and D_delta_tier1 is not None:
        prop['D_delta'] = D_delta_tier1
    ...
    return prop
\end{lstlisting}

Note the third clause: if \code{build\_propagator}'s symbolic block left
\code{D\_delta} as \code{None} (because it busted its budget or the matrix was
rich), the exact path's \code{D\_delta} fills it in.

\begin{note}
The function's docstring (\file{pipeline/\_propagator.py:109}) says it returns
\code{(None, None)} on failure, but every actual \code{return} statement in the
function --- success and failure alike --- is a \emph{three}-tuple
\code{(pole\_vals, C\_mats, D\_delta)}. The caller unpacks three values, so the
code is self-consistent; the docstring is simply stale. When code and prose
disagree, the code is authoritative.
\end{note}

\section{Residues, tiers 2 and 3: the numerical cofactor fallback}

When the exact path returns \code{None} --- non-rationalisable parameters, a
non-squarefree denominator, or a timeout --- \code{compute\_poles\_and\_residues}
substitutes the parameters into $K_{\mathrm{ft}}$ and proceeds numerically
(\file{pipeline/\_propagator.py:926} onward). Which numerical branch runs depends
on whether \code{build\_propagator} left a symbolic adjugate behind:

\begin{itemize}
  \item the \emph{rich / heterogeneous} branch fires when \code{adj\_ft} or
        \code{G\_ft} is \code{None} (\file{pipeline/\_propagator.py:939}), and
  \item the \emph{legacy symbolic} branch fires when both are present
        (\file{pipeline/\_propagator.py:1167}).
\end{itemize}

Both branches find poles the same way --- they take the symbolic
$\det(K_{\mathrm{ft,num}})$, extract its numerator's $\omega$-coefficients, build a
characteristic polynomial, and root it. A sanity check guards against a common
silent failure (\file{pipeline/\_propagator.py:971}):

\begin{lstlisting}
free_syms = set()
for i in range(nf):
    for j in range(nf):
        free_syms.update(SR(K_ft_num[i, j]).variables())
free_syms.discard(omega)
if free_syms and verbose:
    print(f'      WARNING: K_ft_num has free symbols beyond omega: '
          f'{sorted(str(s) for s in free_syms)}.  Pole-finding '
          f'will likely return 0 roots.  ...')
\end{lstlisting}

\begin{gotcha}
\textbf{Pole-finding silently returns zero roots if any free symbol other than
$\omega$ survives.} If a kernel symbol was not substituted (a
\code{kernel\_ft\_image} hook missing an entry) or a saddle value is absent from
\code{num\_params}, then $\det(K_{\mathrm{ft,num}})$ has non-numeric coefficients
and the root-finder finds nothing. The rich branch \emph{warns} about this
explicitly (\file{pipeline/\_propagator.py:976}); the empty pole list then
surfaces downstream as a warning inside \code{build\_G\_t\_matrix}. If you ever
see ``$0$ retarded poles'' on a model you expect to have some, an
unsubstituted symbol is the first thing to check.
\end{gotcha}

Both branches cache $K$'s entrywise numerator/denominator coefficients once
(\file{pipeline/\_propagator.py:1035}) so that the kernel matrix can be evaluated
at any $\omega$ in microseconds via Horner's rule (\code{\_horner},
\file{pipeline/\_propagator.py:841}), candidates are deduplicated, and each is
\term{Newton-refined} on the numerical $\det(K(\omega))$ --- Newton's method
polishing each root to machine precision, with acceptance gated on
$|\det(K(\omega))|/\text{scale} < 10^{-9}$, staying in the upper half-plane, and
deduplication.

The residue evaluation is the common, important part. In \emph{both} branches the
residue is computed via the numerical \term{cofactor adjugate}, never the
symbolic \code{adj\_ft} (\file{pipeline/\_propagator.py:1470}):

\begin{lstlisting}
def _cofactor_adj(M):
    """adj(M) = transpose of cofactor matrix.
    adj[i,j] = (-1)^(i+j) . det(M with row j and col i removed)."""
    n = M.shape[0]
    A = np.zeros_like(M)
    for i in range(n):
        for j in range(n):
            minor = np.delete(np.delete(M, j, axis=0), i, axis=1)
            A[i, j] = ((-1) ** (i + j)) * np.linalg.det(minor)
    return A
\end{lstlisting}

At a pole $K(\omega_k)$ is singular, so $K^{-1}$ does not exist --- but the
adjugate does, by direct cofactor expansion. The determinant derivative
$Q'(\omega_k)$ is taken by central difference, and the residue follows
(\file{pipeline/\_propagator.py:1482}):

\begin{lstlisting}
for pk in pole_vals:
    K_at_pk = _K_at(pk)
    adj_at_pk = _cofactor_adj(K_at_pk)
    h = 1e-6 * (1 + abs(pk))
    d_prime = (_det_at(pk + h) - _det_at(pk - h)) / (2 * h)
    ...
    C_np = 1j * adj_at_pk / d_prime
\end{lstlisting}

\begin{gotcha}
\textbf{The cofactor index convention is transposed.} \code{\_cofactor\_adj} sets
\code{adj[i,j] = (-1)\^{}(i+j) * det(M with row j and col i removed)} --- note
row~$j$, column~$i$. The adjugate is the transpose of the cofactor matrix, and
getting this backwards would transpose every residue matrix
(\file{pipeline/\_propagator.py:1470}).
\end{gotcha}

\begin{gotcha}
\textbf{Residues are always numerical, even when a symbolic \code{adj\_ft}
exists.} Sage stores the symbolic adjugate as a sum of rationals whose per-term
denominators diverge by factors of $10^3$ to $10^4$ as $\omega$ approaches a
kernel pole. Evaluating that sum at the pole causes \term{catastrophic
cancellation} --- nearly-equal large terms subtracting to lose $3$ to $5$
significant digits (a measured $11\%$ to $42\%$ relative error on a quartic
model's residues). The numpy cofactor path builds the minor determinants directly
and has no such leak, so both numerical branches use it
(\file{pipeline/\_propagator.py:1503}). The ``lean symbolic'' branch keeps the
symbolic \code{adj\_ft} around for display, but computes residues numerically all
the same.
\end{gotcha}

\begin{note}
The upshot is that the three tiers are \emph{numerically equivalent} in current
code: the legacy symbolic path was patched to use the numpy cofactor evaluation
too, so tiers~2 and~3 differ only in how they \emph{find the poles}, not in how
they evaluate the residues. One cosmetic redundancy is worth knowing: for the
rich branch there are two nearly-identical Newton-refine blocks
(\file{pipeline/\_propagator.py:1297} and \file{:1389}), so a candidate can be
refined twice --- harmless, since the second refine of an already-converged root
is a near no-op.
\end{note}

\section{The instantaneous limit \texorpdfstring{$D_\delta$}{D-delta}}

If $G_{\mathrm{ft}}(\omega)$ does not vanish as $\omega \to \infty$ --- a
\emph{non-proper} rational function --- then its inverse Fourier transform
produces a $\delta(t)$ term. Decompose
\begin{equation}
  G_{\mathrm{ft}}(\omega) \;=\; Q_{\mathrm{poly}}(\ii\omega)
  \;+\; G_{\mathrm{proper}}(\omega),
  \qquad
  G_{\mathrm{proper}} \xrightarrow{\;\omega\to\infty\;} 0,
\end{equation}
and the polynomial part transforms to derivatives of $\delta(t)$. For the common
case where that polynomial part is a constant matrix,
\begin{equation}
  D_\delta[i,j] \;=\; \lim_{\omega\to\infty} G_{\mathrm{ft}}[i,j](\omega).
  \label{eq:Ddelta}
\end{equation}

\begin{defn}
A rational function is \term{strictly proper} if $\deg(\text{numerator}) <
\deg(\text{denominator})$; it vanishes at infinity and contributes no $\delta(t)$.
A \term{non-proper} rational has a polynomial part and does contribute one. The
matrix $D_\delta$ collects the $\delta(t)$ coefficients of the time-domain
propagator. Physically, a nonzero entry is an \emph{immediate, same-time}
response --- in the \msrjd{} action, a response-field source $\tilde n$ at time
$t$ produces a $\delta n$ at the \emph{same} $t$.
\end{defn}

Concretely, per entry $G_{\mathrm{ft}}[i,j] = P_{ij}/Q$:
\[
  \deg(P_{ij}) < \deg(Q) \;\Rightarrow\; D_\delta[i,j] = 0,
  \qquad
  \deg(P_{ij}) = \deg(Q) \;\Rightarrow\; D_\delta[i,j] =
  \frac{\operatorname{lc}(P_{ij})}{\operatorname{lc}(Q)},
\]
the ratio of leading coefficients. The engine computes this three different ways
depending on which path produced the propagator:

\begin{itemize}
  \item In the \emph{exact path}, it is the leading-coefficient ratio over the
        exact polynomials (\file{pipeline/\_propagator.py:283}).
  \item In the \emph{lean symbolic} block, it is \code{\_omega\_inf\_limit\_fast}
        applied per entry (\file{pipeline/\_propagator.py:719}).
  \item In the \emph{rich} numerical branch, it is a \term{two-probe scaling test}
        (\file{pipeline/\_propagator.py:1150}): evaluate $G(\omega)$ at $\omega =
        10^8$ and $\omega = 10^{10}$; an entry whose value is the same at both
        probes is a constant (so it becomes a $\delta$), while one that scales like
        $1/\omega$ has decayed to zero.
\end{itemize}

\code{\_omega\_inf\_limit\_fast} (\file{pipeline/\_propagator.py:460}) deserves a
look, because it exists purely to avoid a slow tool:

\begin{lstlisting}
def _omega_inf_limit_fast(expr, omega_var):
    expr = SR(expr)
    try:
        num = expr.numerator()
        den = expr.denominator()
        num_deg = int(num.degree(omega_var))
        den_deg = int(den.degree(omega_var))
    except (AttributeError, TypeError, ValueError):
        return None
    if num_deg < den_deg:
        return SR(0)
    if num_deg > den_deg:
        return None
    # Same degree: leading-coefficient ratio.
    ...
\end{lstlisting}

\begin{gotcha}
\textbf{Sage's \code{limit} is avoided for the $\omega \to \infty$ limit.} It
routes through Maxima (a Lisp computer-algebra system bundled with Sage) and, on
raw cofactor or determinant entries, triggers thousands of FLINT divide-by-zero
warnings and Singular thrashing. \code{\_omega\_inf\_limit\_fast} does the
textbook rational-limit recipe --- compare numerator and denominator degrees,
return the leading-coefficient ratio --- entirely by hand, with no Maxima
(\file{pipeline/\_propagator.py:464}).
\end{gotcha}

\section{The time-domain propagator \texorpdfstring{$G_R(t)$}{G\_R(t)}}

Putting the pieces together --- the inverse Fourier transform of a sum of simple
poles plus a constant --- gives the time-domain \term{retarded propagator}:
\begin{equation}
  G_R[i,j](t) \;=\; D_\delta[i,j]\,\delta(t)
  \;+\; \Theta(t)\sum_k C_k[i,j]\,\ee^{\ii\omega_k t}.
  \label{eq:GR}
\end{equation}
The Heaviside step $\Theta(t)$ enforces \emph{retardation} (response only after
the cause); with $\mathrm{Im}(\omega_k) > 0$ each mode decays for $t > 0$, so the
product is well-defined.

This is exactly what \code{build\_G\_t\_matrix}
(\file{msrjd/integration/time\_domain/propagator\_td.py:135}) assembles. It
returns a dict with a \code{smooth} matrix (the pole-residue sum, \emph{without}
the $\Theta$ --- the caller applies that), a \code{delta} matrix (the $\delta(t)$
coefficients), and the time variable:

\begin{lstlisting}
smooth = matrix(SR, n_propagator, n_propagator, 0)
for k in range(len(pole_vals)):
    smooth = smooth + C_mats[k] * exp(I * pole_vals[k] * t_var)
\end{lstlisting}

Note the accumulator is initialised as a proper $n_f \times n_f$ zero matrix
\emph{before} the loop, so that even when \code{pole\_vals} is empty the result is
still a matrix (a bare \code{sum} over an empty generator would collapse to the
Python integer \code{0} and break every downstream \code{.dimensions()} call ---
the comment at \file{propagator\_td.py:252} explains the trap). When
\code{pole\_vals} \emph{is} empty, the function also issues a warning
(\file{propagator\_td.py:271}), because then only the $\delta$-part survives and
that usually means the characteristic-polynomial solve found no retarded roots.

\subsection{The complex-comparison pitfall}

Before substituting numbers, \code{build\_G\_t\_matrix} normalises the poles and
residues with a small helper, \code{\_to\_sr\_ab}
(\file{propagator\_td.py:225}):

\begin{lstlisting}
def _to_sr_ab(value):
    """Normalise only raw complex / CDF scalars; leave SR alone."""
    if isinstance(value, complex):
        return SR(float(value.real)) + SR(float(value.imag)) * I
    if isinstance(value, ComplexDoubleElement):
        c = complex(value)
        return SR(float(c.real)) + SR(float(c.imag)) * I
    # Symbolic SR expression, Sage number, etc. -- preserve exactly.
    return SR(value)
\end{lstlisting}

\begin{gotcha}
\textbf{\code{\_to\_sr\_ab} must run before \code{.subs(num\_params)}, and it must
\emph{not} cast genuine SR expressions through \code{complex()}.} The first half
of that rule: a raw Python \code{complex} fed to \code{SR(\ldots)} becomes an
opaque GiNaC node, and when that node later propagates into a product or
expansion, GiNaC's internal term-sort tries to compare two \code{complex} objects
with \code{<} and raises \code{TypeError: '<' not supported between instances of
'complex' and 'complex'}. Rewriting it as $a + b\,\ii$ with real \code{float}
components sidesteps that. The second half: an \emph{exact} closed-form pole such
as $\ii(1 - \sqrt6/10)$ must keep its algebraic content; casting it through
\code{complex()} would silently collapse it to double precision and bias
correlators that compound the error across many $\tau$ evaluations
(\file{propagator\_td.py:204}). So \code{\_to\_sr\_ab} normalises only raw complex
scalars and leaves true SR expressions untouched.
\end{gotcha}

\subsection{The delta coefficients}

\code{build\_G\_t\_matrix} prefers the precomputed \code{D\_delta}, applying
\code{num\_params} if given, and only falls back to Sage's symbolic \code{limit}
when no \code{D\_delta} was precomputed (\file{propagator\_td.py:311}):

\begin{lstlisting}
D_precomputed = propagator_data.get('D_delta')
if D_precomputed is not None:
    if num_params:
        delta_coeffs = D_precomputed.apply_map(
            lambda e: SR(e).subs(num_params) if not SR(e).is_zero() else SR(0)
        )
    else:
        delta_coeffs = D_precomputed
else:
    # Compute from G_ft via symbolic limit.
    ...
\end{lstlisting}

Since every path in phase~4 fills \code{D\_delta}, the symbolic-limit fallback is
rarely exercised.

\subsection{The per-edge lookups and the transpose}

The diagram integrators do not consume the matrix directly; they ask for one
entry at a time via \code{G\_t\_entry} (\file{propagator\_td.py:379}) for the
smooth part and \code{G\_t\_delta\_coeff} (\file{propagator\_td.py:433}) for the
$\delta$ coefficient. The smooth lookup is where the index transpose is applied:

\begin{lstlisting}
if isinstance(G_t_obj, dict):
    G_t_matrix = G_t_obj['smooth']
else:
    G_t_matrix = G_t_obj

t_var = _infer_time_variable(G_t_matrix)
entry = SR(G_t_matrix[phys_idx, resp_idx])
if t_var is not None:
    entry = entry.subs({t_var: t_expr})
if include_heaviside:
    entry = entry * heaviside(t_expr)
return entry
\end{lstlisting}

\begin{gotcha}
\textbf{The kernel layout is [response, physical], but the physical retarded
propagator is read with the indices transposed.} The physical object is ``the
response of physical field $j$ to a response-field source $i$,'' written $G^R_{j
\leftarrow i}$, which is the matrix entry $G[j, i]$ --- physical row, response
column. So \code{G\_t\_entry(phys\_idx=j, resp\_idx=i, \ldots)} reads
\code{smooth[phys\_idx, resp\_idx]} to apply the transpose. Every consumer must
use the same transpose; this one matches \code{\_get\_propagator\_entry} in
\file{msrjd/integration/symbolic.py} (\file{propagator\_td.py:44}). For an edge
$u \to v$ the time argument is $t_{\mathrm v} - t_{\mathrm u}$.
\end{gotcha}

The \code{include\_heaviside} flag multiplies in $\Theta(t)$ to enforce
retardation; only the \emph{smooth} part is returned, because the $\delta(t)$ part
is handled separately by the tree integrator (where a $\delta$ edge collapses one
integration variable). The companion \code{G\_t\_delta\_coeff}
(\file{propagator\_td.py:433}) returns that $\delta$ coefficient as a Python
\code{complex}, demoting it to a real \code{float} when the imaginary part is
negligible, and returning \code{0.0+0.0j} when there is no $\delta$ component.

\begin{gotcha}
\textbf{$\Theta(0)$ semantics differ by path.} The numerical integrator treats
$\Theta(0) = 0$ --- a zero time difference is strictly excluded from the retarded
support, enforced by a Heaviside-filter integrand wrapper in
\file{final\_integral.py} and by strict polytope inequalities. But Sage's
\emph{symbolic} \code{heaviside(0)} returns $1/2$ by default. That $1/2$ is used
only in the symbolic-display path, never in the JIT-compiled numerical integrand
(\file{propagator\_td.py:30}). The two coexist precisely because they never meet.
\end{gotcha}

\subsection{Inferring the time variable without triggering Maxima}

A small but instructive helper, \code{\_infer\_time\_variable}
(\file{propagator\_td.py:459}), recovers the symbolic $t$ from the smooth matrix.
Once parameters are substituted, $t$ is the only free variable left in each
entry. The naive way to find a nonzero entry would be \code{entry == 0} --- but
that comparison routes to Maxima, which fails on entries carrying embedded Python
\code{complex} coefficients. So the helper checks \code{entry.variables()}
directly and only uses \code{is\_trivial\_zero()} for scalar entries. It is a tiny
function, but it embodies a recurring theme in this subsystem: \emph{avoid the
slow, fragile symbolic tool when a cheap structural check will do.}

\section{A concrete example: the OU propagator}

It is worth grounding all of this in the simplest model, the white-noise OU
process whose free action is $S_{\mathrm{free}} \supset \tilde n(\Dt n + \mu n)$.
There is one field, so $n_f = 1$ and every matrix is $1\times1$:

\begin{itemize}
  \item The single kernel entry is $K[\tilde n, n] = \Dt + \mu$, which
        \code{\_to\_kernel} rewrites as $\delta' + \mu\,\delta$.
  \item Under the Fourier convention \eqref{eq:ftconv}, $K_{\mathrm{ft}} = \ii\omega
        + \mu$.
  \item The propagator is $G_{\mathrm{ft}} = 1/(\ii\omega + \mu)$, with a single
        pole at $\omega = \ii\mu$ --- which has $\mathrm{Im} > 0$ for $\mu > 0$, so
        it is retarded.
  \item The residue of $G_{\mathrm{ft}}$ at the pole is $1/\ii = -\ii$, so by the
        $C_k$ convention \eqref{eq:Ck} the residue matrix is $C_0 = \ii\cdot(-\ii)
        = 1$.
  \item $G_{\mathrm{ft}}$ is strictly proper ($\deg P = 0 < 1 = \deg Q$), so
        $D_\delta = 0$.
\end{itemize}

Assembling \eqref{eq:GR}: $G_R(t) = \Theta(t)\cdot 1\cdot\ee^{\ii(\ii\mu)t} =
\Theta(t)\,\ee^{-\mu t}$ --- the textbook retarded relaxation. A two-field system
(a density coupled to a potential, say) produces a $2\times2$ $K_{\mathrm{ft}}$
with two poles, and a $\tilde n \leftrightarrow \delta n$ coupling that yields a
nonzero $D_\delta$ entry --- an instantaneous $\delta(t)$ piece in $G_R$.

\section{Robustness machinery and external tools}

This subsystem leans almost entirely on \term{SageMath} and, through it, several
computer-algebra back-ends. None of nauty, networkx, numba, or sympy is called
\emph{directly} here --- though, as noted, Sage's \code{fourier\_transform} routes
through SymPy internally. A consolidated tour of the tools, since this is where
they all first appear:

\begin{description}
  \item[SageMath / \code{SR}] The Symbolic Ring (GiNaC/Pynac-backed) and the glue
        that routes each operation to the right back-end. \code{matrix(ring, data)}
        builds matrices over \code{SR}, \code{CDF}, or an exact field; matrix
        methods used include \code{.inverse()}, \code{.det()}, \code{.apply\_map(fn)}
        (entrywise function application), and \code{.dimensions()}.
  \item[\code{QQ}, \code{CDF}, \code{CyclotomicField(4)}] The rationals
        $\mathbb{Q}$, the machine-precision complex field, and $\mathbb{Q}(i)$ ---
        the three number fields the residue paths move between.
  \item[\code{PolynomialRing} / fraction field] $R[\omega]$ and
        $\mathrm{Frac}(R[\omega])$; the latter reduces each inverse entry to
        canonical irreducible $P/Q$ form when $R$ is an exact field.
  \item[NumPy] Imported lazily inside the numerical branches; \code{np.linalg.inv}
        / \code{np.linalg.det} for the matrix evaluations, \code{np.delete} to form
        the minors of the cofactor adjugate, and \code{np.polyfit} for the
        fallback pole-finder. NumPy is preferred for residues because its LU-based
        linear algebra avoids the sum-of-rationals cancellation the symbolic path
        suffers near poles.
  \item[\code{cysignals}] A small Cython library (bundled with Sage) that can
        interrupt or recover from signals raised inside native C/Cython routines.
        Plain Python \code{try/except} cannot catch a hang or crash deep inside
        Singular; \code{cysignals} can.
  \item[Python \code{signal}] The POSIX signal API. \code{\_run\_with\_timeout}
        uses \code{SIGALRM} + \code{signal.alarm} to wall-clock-bound the exact
        inverse and the symbolic-inverse block, degrading to a plain call where
        \code{SIGALRM} is absent.
\end{description}

The most colourful piece of robustness machinery is the \code{factor()} hazard.
\code{build\_propagator} deliberately returns the \emph{raw}, unfactored inverse,
because \code{factor()} can hang or segfault Singular (Sage's polynomial CAS). A
separate, purely cosmetic helper \code{factor\_propagator}
(\file{pipeline/\_propagator.py:366}) applies factoring only for report display,
and it does so with \emph{two} layers of protection (\code{\_safe\_factor},
\file{pipeline/\_propagator.py:422}):

\begin{lstlisting}
from cysignals.alarm import alarm, cancel_alarm, AlarmInterrupt
try:
    alarm(timeout)
    try:
        return e.factor()
    finally:
        cancel_alarm()
except (RuntimeError, ValueError, TypeError, ArithmeticError,
        AlarmInterrupt):
    return e
except BaseException as exc:
    if exc.__class__.__name__ == 'SignalError':
        return e
    raise
\end{lstlisting}

The \code{alarm(timeout)} arms a timer that raises \code{AlarmInterrupt} to stop a
runaway \code{factor()} loop, and on any failure the helper falls back to the
unfactored entry. \code{factor\_propagator} adds a one-shot ``bail'' flag that
permanently skips factoring after any single entry exceeds a threshold, because
the alarm does not always interrupt Singular's native loop.

\begin{gotcha}
\textbf{The \code{cysignals} \code{SignalError} is caught by class name, not by
\code{isinstance}.} In some Sage builds \code{SignalError} (Sage's wrapper around
a native segfault) is \emph{not} a subclass of \code{Exception}. So both
\code{\_safe\_factor} and \code{build\_propagator} catch \code{BaseException} and
check \code{exc.\_\_class\_\_.\_\_name\_\_ == 'SignalError'}, re-raising everything
else --- which is what keeps \code{KeyboardInterrupt} and \code{SystemExit} from
being swallowed (\file{pipeline/\_propagator.py:451}, \file{:746}).
\end{gotcha}

\section{Caching}

The symbolic stages of the propagator are parameter-independent, so they are
cached to disk. \code{build\_propagator} slugifies the model name into a tag and
opens a \code{PipelineCache} on that directory
(\file{pipeline/\_propagator.py:547}):

\begin{lstlisting}
prop_tag = re.sub(r'[^A-Za-z0-9]+', '_', model['name']).strip('_').lower()
cache_dir = f"{cache_dir_root}/{prop_tag}"
cache = PipelineCache(cache_dir)
\end{lstlisting}

\code{PipelineCache} (\file{msrjd/core/cache.py}) wraps Sage's pickling ---
\code{save} writes a \code{.sobj} (Sage object) file, \code{load} reads it back ---
so the cache lands at \file{saved\_theories/<tag>/propagator.sobj}. The cache has
two staleness guards (\file{pipeline/\_propagator.py:551}): it rebuilds if the
cached \code{nf} differs from \code{ft.\_n\_tilde} (catching a field-list edit that
did not change the model name), and it rebuilds if the model is spatial but the
cached propagator predates the spatial block (no \code{G\_tx\_sym}).

\begin{gotcha}
\textbf{Closures do not pickle, so the spatial \code{G\_tx} callables are attached
\emph{after} the cache save.} For a spatial model, \code{build\_propagator} builds
the heat-kernel propagator block and merges it into \code{prop}, but the runtime
callables \code{G\_tx} are functions, which Sage cannot pickle. So the cache is
saved \emph{first} (\file{pipeline/\_propagator.py:817}), with only the picklable
symbolic data, and the \code{G\_tx} callables are rebuilt and attached afterwards
via \code{make\_g\_tx\_callables} (\file{pipeline/\_propagator.py:830}). On a cache
hit the same rebuild step runs. The spatial heat-kernel propagator is the subject
of its own chapter; here it appears only at its caching seam.
\end{gotcha}

\section{Known limitations and failure modes}

A short field guide to the ways this subsystem can refuse or mislead, gathered in
one place:

\begin{description}
  \item[Higher-multiplicity (Jordan-block) poles are not handled.] A
        multiplicity-$\ge 2$ pole produces a $\tau^k\,\ee^{-pt}$ mode that the
        single-pole residue formula \eqref{eq:residue} cannot represent. The exact
        path \emph{bails} on a non-squarefree denominator
        (\file{pipeline/\_propagator.py:196}); the numerical path \emph{silently
        leaves a zero residue} where $|Q'| < 10^{-30}$. A future version would need
        the $m$-th-derivative residue formula \emph{and} a downstream integrator
        that can absorb the polynomial-times-exponential modes. This is the most
        significant known limitation.
  \item[Non-rationalisable parameters knock out the exact path.] An irrational
        symbolic parameter throws inside \code{\_sr\_complex\_to\_CF}; the exact
        tier returns \code{None} and the numerical cofactor path takes over.
  \item[Surviving free symbols give zero poles.] If a kernel symbol is
        unsubstituted or a saddle value is missing from \code{num\_params}, the
        characteristic polynomial has non-numeric coefficients and the root-finder
        returns nothing --- surfaced as a warning, then an empty pole list.
  \item[An empty pole list makes the smooth propagator identically zero.] Only
        $D_\delta$ then contributes. \code{build\_G\_t\_matrix} warns about this
        (\file{propagator\_td.py:271}) precisely so an unexpected
        no-retarded-roots solve does not slip through silently.
\end{description}

\section{What you just saw}

The propagator is two phases of the trace --- \texttt{[2/7]} and \texttt{[4/7]}
--- and one object, $\Gret$, built in three movements:

\begin{enumerate}
  \item \textbf{Action to frequency domain.} Read the kernel matrix $K$ out of the
        bilinear free action, rewrite each entry in $\delta/\delta'$ kernel form,
        Fourier transform under $\ee^{-\ii\omega t}$, and invert to get
        $G_{\mathrm{ft}}(\omega)$ --- with a matrix-size gate that skips symbolic
        inversion for large matrices and defers them to the numerical stage. All of
        this is parameter-independent and cached.
  \item \textbf{Poles and residues.} Substitute the numeric parameters; find the
        retarded poles ($\mathrm{Im}(\omega) > 0$) and the residue matrices $C_k =
        \ii\,\mathrm{Res}\,G_{\mathrm{ft}}$, preferring an exact $\mathbb{Q}(i)$
        path and falling back to a numerical cofactor path, with residues always
        evaluated numerically to dodge catastrophic cancellation.
  \item \textbf{Back to the time domain.} Assemble $G_R(t) = D_\delta\,\delta(t) +
        \Theta(t)\sum_k C_k\,\ee^{\ii\omega_k t}$, and serve it one edge at a time
        --- transposed to [physical, response] --- to the diagram integrators.
\end{enumerate}

With $\Gret$ in hand, the engine has the building block for every Feynman edge.
The next chapters open up what is wired together with those edges: the mean-field
saddle the parameters are evaluated at --- which is where \code{num\_params} comes
from --- the diagram enumeration that decides \emph{which} graphs to build, and
the loop integration that finally turns a graph of $\Gret$ edges into a number.

% ---------------------------------------------------------------------
%  Chapter 8 : The Mean-Field Saddle and the DAE Solver
% ---------------------------------------------------------------------
\chapter{The Mean-Field Saddle and the DAE Solver}\label{ch:mean-field}

Every loop expansion needs a point to expand \emph{about}. In the \msrjd{}
formalism that point is the \term{mean-field saddle}: the deterministic steady
state of the system with all fluctuations switched off. It is the firing rate
$n^{*}$, the membrane voltage $v^{*}$, the variance-suppressed equilibrium value
$x^{*}$ --- whatever the deterministic field settles to when the noise is muted.
The whole perturbative series (the diagrams) is a systematic correction
\emph{around} that background, so before a single diagram can be evaluated the
engine must find the saddle and turn it into numbers.

This chapter is about the subsystem that does exactly that. Its single job is
narrow and entirely concrete: \textbf{given a fully specified model and a
dictionary of numeric parameter values, find the saddle and package it into a
substitution dictionary that the rest of the pipeline plugs into its symbolic
expressions}. That dictionary is called \code{num\_params}, and it is the only
thing this subsystem hands downstream. The chapter walks the two solvers that
produce it --- one for theories that declare their steady state as a
self-consistency closure, one for theories that declare it as a
differential-algebraic system --- and then the linear-stability machinery that
decides whether a given saddle is even a legitimate place to expand.

The code lives in two files. The legacy \emph{iteration} solver is
\file{pipeline/\_mean\_field.py} (553 lines); the \emph{DAE multi-root} solver
plus linear stability is \file{pipeline/\_mean\_field\_dae.py} (977 lines). They
are, by design, completely independent of each other, and which one runs is
decided by a single line in \file{pipeline/compute.py}.

This is phase \code{[3/7]} of the seven-phase trace you met in
Chapter~\ref{ch:quickstart} (\code{[3/7] Solve MF self-consistency...}). It sits
after the propagator has been built symbolically and before the propagator's
poles and residues are computed numerically --- because those poles need the
saddle's numbers to exist.

\section{Why a saddle at all}

\subsection{The background expansion in MSR-JD}

Recall the \msrjd{} picture. A stochastic dynamical system --- a Langevin
equation, a master equation, a spiking network --- is rewritten as a generating
functional
\begin{equation}
  Z[J] \;=\; \int \mathcal{D}[\phi]\,\mathcal{D}[\tilde\phi]\;
    \exp\!\big(-S[\phi,\tilde\phi] + \text{sources}\big),
  \label{eq:mf-genfunc}
\end{equation}
where $\phi$ is the physical field (a rate, a voltage, a density) and
$\tilde\phi$ is the conjugate \term{response field}. The action $S$ is the object
the entire pipeline manipulates.

Perturbation theory expands the physical field about a background,
$\phi = \phi^{*} + \delta\phi$, and the background $\phi^{*}$ is chosen so that
the action has \emph{no linear term} in the fluctuation $\delta\phi$. That
``vanishing linear term'' condition is not an arbitrary convenience: it is
exactly the classical equation of motion, i.e.\ the deterministic steady state.
A linear term in $\delta\phi$ would be a tadpole --- a one-point function that
should vanish on a correctly-chosen background --- and leaving it in would make
the diagrammatic bookkeeping wrong. So the saddle is the \emph{tree-level}
solution, and the loop corrections (the diagrams) are the systematic corrections
to it.

\begin{defn}
The \textbf{mean-field saddle} $\phi^{*}$ is the field configuration that makes
the action's first variation vanish, $\delta S/\delta\phi|_{\phi^{*}} = 0$. It is
the deterministic steady state of the dynamics. The loop expansion is an
expansion in fluctuations $\delta\phi = \phi - \phi^{*}$ \emph{around} this point.
\end{defn}

\subsection{The canonical self-consistency problem}

For a rate network the saddle is a self-consistency problem. The textbook neural
example, which the legacy solver is built around, is
\begin{align}
  v^{*}_i &= E_i + \sum_j w_{ij}\, g_{ij}\!\star n^{*}_j,
    && \text{(voltage $=$ drive $+$ recurrent input)} \label{eq:mf-volt}\\
  n^{*}_i &= \phi_i(v^{*}_i),
    && \text{(rate $=$ transfer function of voltage).} \label{eq:mf-rate}
\end{align}
Here $g_{ij}$ is a synaptic kernel, $E_i$ an external drive, $w_{ij}$ a coupling
weight, and $\phi_i$ a (possibly nonlinear) \term{transfer function} from input
to output rate. The symbol $\star$ is convolution in time.

One feature of \eqref{eq:mf-volt} is load-bearing throughout the code: at the
saddle every kernel \emph{integrates to one}. Kernels are normalized,
$\int g(t)\,\dd t = 1$, so the convolution $g\star n^{*}$ of a \emph{constant}
$n^{*}_j$ is just $n^{*}_j$ --- a stationary mean field sees only a kernel's total
weight, never its shape. This is precisely why both solvers substitute every
kernel symbol $\to 1$ before evaluating anything (you will see this as
\code{g\_to\_one} and \code{kernel\_to\_one} below).

Substituting the voltage equation \eqref{eq:mf-volt} into the rate equation
\eqref{eq:mf-rate} closes the problem in the rates alone:
\begin{equation}
  n^{*}_i \;=\; \phi_i\!\Big(E_i + \sum_j w_{ij}\, n^{*}_j\Big)
    \;\equiv\; T_i(n^{*}).
  \label{eq:mf-closure}
\end{equation}
The legacy solver finds $n^{*}$ by driving the residual $n^{*}_i - T_i(n^{*})$ to
zero. This is the split that pervades \file{\_mean\_field.py}: the rates $n^{*}$
are the \term{iteration saddle} (the variables the root finder actually
iterates), and the voltages $v^{*}$ are the \term{compound saddle} (computed from
the rates on the fly each step).

\subsection{Where it sits in the pipeline}

The subsystem is fed by \code{TheoryBuilder.build()} (the \code{model} dict plus a
compiled \code{FieldTheory}) and the user's numeric parameters, and it hands its
\code{num\_params} dict to the propagator's pole finder. Schematically:

\begin{lstlisting}[style=console]
   TheoryBuilder.build()              (pipeline/theory.py)
        |  model dict + compiled FieldTheory ft
        v
   compute_cumulants(...)             (pipeline/compute.py)
        |  [1] expand action -> vertices / sources
        |  [2] build_propagator -> K_ft, G_ft   (symbolic, in omega)
        |
        +--[3] MEAN-FIELD SOLVE  <----------  THIS SUBSYSTEM
        |        if model['equations']: solve_mean_field_dae_compat(...)
        |        else:                   solve_mean_field(...)
        |        -> mf['num_params']   {SR symbol -> float}
        |
        v
   compute_poles_and_residues(prop, num_params)   (pipeline/_propagator.py)
        |  substitute num_params into K_ft, find retarded poles, residues
        v
   Phase J / time-domain machinery -> connected cumulant C(tau)
\end{lstlisting}

What feeds it:
\begin{itemize}
  \item \code{ft} --- a compiled \code{FieldTheory}. The solver touches only
        \code{ft.\_ns} (the Sage symbol namespace) and \code{ft.taylor\_order}
        (how many transfer-function derivatives to precompute). In the DAE path
        \code{ft} is used \emph{only} by the compatibility wrapper, to map saddle
        names back onto the \code{ft.\_ns} symbol arrays.
  \item \code{model} --- the declaration dict from \code{build()}. It carries
        \code{parameters}, \code{kernels}, \code{populations},
        \code{physical\_fields}, \code{functions}, \code{phi\_concrete},
        \code{mf\_bg\_conditions}, \code{iteration\_saddles}, \code{equations},
        and \code{stability\_analysis}.
  \item \code{fundamental} --- the user's concrete numeric parameter values, keyed
        by parameter \emph{name} (e.g.\ \code{\{'E': [...], 'w': [[...]], 'tau': 1.0\}}).
\end{itemize}

What consumes its output: the \emph{primary} consumer is
\code{compute\_poles\_and\_residues(prop, num\_params)}
(\file{pipeline/\_propagator.py:851}), which substitutes \code{num\_params} into
the symbolic kernel $K_{\mathrm{ft}}(\omega)$ to get a numeric matrix-valued
rational function of $\omega$, then finds its retarded poles and residue
matrices. A second consumer, \code{\_solve\_mf\_at\_saddle}
(\file{pipeline/\_precompute.py:200}), re-runs the solver to verify the action's
linear term vanishes at the saddle. The handoff from the saddle to the
perturbative expansion is, in both cases, \emph{purely substitutional}.

\section{Two solvers, one routing decision}

A model declares its steady state in one of two styles, and the framework picks a
solver to match. The decision is a single truthiness test on
\code{model['equations']}:

\begin{lstlisting}
# pipeline/compute.py:364
if model.get('equations'):
    mf = solve_mean_field_dae_compat(
        ft, model, fundamental,
        fixed_point_index=fixed_point_index,
        n_starts=mf_dae_n_starts,
        seed_box=mf_dae_seed_box,
        verbose=verbose,
    )
else:
    mf = solve_mean_field(ft, model, fundamental, verbose=verbose)
num_params = mf['num_params']
\end{lstlisting}

The two declaration styles, and the solvers they route to:

\begin{description}
  \item[Self-consistency closure $\to$ legacy iteration solver.] A theory whose
        saddle is described as a closure ($n^{*}_i = \phi(v^{*}_i)$ with $v^{*}_i$
        an algebraic function of the $n^{*}_j$) goes to
        \code{solve\_mean\_field} in \file{\_mean\_field.py}. It iterates only the
        rate-like saddle variables and treats the rest (voltages) as compound
        expressions evaluated on the fly. It has a single-population branch and a
        heterogeneous (\code{E}/\code{I}) branch.
  \item[Differential-algebraic system $\to$ DAE multi-root solver.] A theory that
        declares its dynamics directly, one equation per field, via
        \code{TheoryBuilder.equation(lhs, rhs, population)}
        (\file{pipeline/theory.py:1087}), goes to
        \code{solve\_mean\_field\_dae\_compat} in \file{\_mean\_field\_dae.py}. It
        sets the time-derivative operator $\code{Dt}\to 0$ on every equation,
        solves the resulting algebraic root problem with many random restarts,
        deduplicates the roots, optionally classifies each root's linear
        stability, and selects one root by index.
\end{description}

\begin{defn}
\textbf{DAE} stands for \emph{differential-algebraic equation}: a system that
mixes genuine differential equations (those carrying the time-derivative operator
\code{Dt}, such as a voltage relaxation $(\tau\,\code{Dt}+1)v=\dots$) with purely
algebraic constraints (those without \code{Dt}, such as a rate readout
$n=\phi(v)$). At the saddle the differential equations collapse to algebraic ones
($\code{Dt}\to 0$), but the differential structure is \emph{remembered} and
reused for stability analysis.
\end{defn}

\begin{gotcha}
The routing is \emph{purely} the truthiness of \code{model.get('equations')}
(\file{compute.py:364}, and identically at \file{\_precompute.py:204}). A theory
that declares both \code{.equation(...)} \emph{and} \code{mf\_bg\_conditions} will
silently take the DAE path. The two modules are, in the words of the source,
``INTENTIONALLY decoupled'' (\file{\_mean\_field\_dae.py:16}) --- there is no
shared code and no negotiation between them.
\end{gotcha}

\section{The legacy iteration solver}

\code{solve\_mean\_field(ft, model, fundamental, verbose=True)}
(\file{\_mean\_field.py:14}) is the entry point. It immediately routes to the
heterogeneous branch if the model declares populations:

\begin{lstlisting}
# pipeline/_mean_field.py:40
if model.get('populations'):
    return _solve_mean_field_hetero(ft, model, fundamental, verbose)
\end{lstlisting}

We will follow the single-population path first, then note what the
heterogeneous branch adds. The solver returns a five-key dict:
\code{nstar\_vals} (list of floats), \code{vstar\_vals} (list),
\code{phi\_deriv\_vals} (a \code{\{(dk, i): float\}} map of saddle derivatives),
\code{num\_params} (the deliverable), and \code{param\_subs} (the raw user
parameters expanded to indexed symbols).

\subsection{Step 1 --- expand the parameters into Sage symbols}

The first job is to turn the user's parameter values, keyed by name, into a dict
keyed by the \emph{indexed Sage symbols} that actually appear inside the
propagator. A matrix parameter \code{w} becomes \code{w11}, \code{w12}, \dots; a
vector becomes \code{E1}, \code{E2}, \dots; a scalar stays bare. The axis sizes
come from the population-size map \code{ns.\_pop\_size}, but the \emph{shape of
the user's value} is the source of truth for whether something is scalar, vector,
or matrix:

\begin{lstlisting}
# pipeline/_mean_field.py:62-87  (abridged)
ib = pspec.get('indexed_by')
if ib:
    if len(ib) == 2:
        n_rows = pop_size_map.get(ib[0], len(ns.pop))
        n_cols = pop_size_map.get(ib[1], len(ns.pop))
        for i in range(n_rows):
            for j in range(n_cols):
                param_subs[SR.var(f'{pname}{i+1}{j+1}')] = val[i][j]
    elif len(ib) == 1:
        ...
        param_subs[SR.var(f'{pname}{i+1}')] = val[i]
    ...
else:
    param_subs[SR.var(pname)] = val
\end{lstlisting}

\code{SR.var('w12')} returns (and globally registers) the symbolic variable named
\code{w12}; calling it again anywhere with the same name returns the \emph{same}
symbol object. That identity is exactly why a substitution dict keyed by
\code{SR.var('w12')} matches the \code{w12} that the propagator built
independently. (Section~\ref{sec:mf-tools} explains Sage's symbolic ring from the
ground up.)

\subsection{Step 2 --- the transfer-function derivatives, symbolically}

The diagrammatic vertices are obtained by Taylor-expanding the action about the
saddle. A transfer function $\phi$ that is nonlinear in the field contributes
vertices whose couplings are the \emph{derivatives of $\phi$ evaluated at the
saddle}:
\begin{equation}
  \phi(v^{*} + \delta v) = \phi(v^{*}) + \phi'(v^{*})\,\delta v
    + \tfrac{1}{2}\phi''(v^{*})\,\delta v^{2}
    + \tfrac{1}{6}\phi'''(v^{*})\,\delta v^{3} + \cdots
  \label{eq:mf-taylor}
\end{equation}
The pipeline names these $\code{phi0\_i} = \phi(v^{*}_i)$ (which equals the rate
$n^{*}_i$), $\code{phi1\_i} = \phi'(v^{*}_i)$, $\code{phi2\_i} = \phi''(v^{*}_i)$,
and so on up to \code{ft.taylor\_order}. The solver builds them by symbolic
differentiation of the \emph{concrete} $\phi$ expression and then evaluating at
the numeric $v^{*}_i$:

\begin{lstlisting}
# pipeline/_mean_field.py:90-102
v_sym = SR.var('_v_mf_')
phi_derivs = {}
for i in ns.pop:
    phi_expr = model['phi_concrete'](ns, i, v_sym)
    phi_derivs[i] = {}
    for dk in range(taylor_order + 1):
        if dk == 0:
            phi_derivs[i][dk] = phi_expr
        else:
            phi_derivs[i][dk] = diff(phi_expr, v_sym, dk)

def phi_num(i, v_val):
    return float(phi_derivs[i][0].subs(param_subs).subs({v_sym: v_val}))
\end{lstlisting}

Here \code{model['phi\_concrete'](ns, i, v\_sym)} returns the concrete symbolic
expression for $\phi$ on a dummy variable \code{\_v\_mf\_}, and
\code{diff(phi\_expr, v\_sym, dk)} takes its $dk$-th derivative. \code{phi\_num}
substitutes the parameters and the numeric voltage, then coerces to a Python
\code{float}.

\begin{defn}
The symbol $\code{phik\_i}$ denotes the $k$-th derivative of the transfer
function at the saddle of population $i$: $\code{phik\_i} = \phi^{(k)}(v^{*}_i)$.
The zeroth one is special, $\code{phi0\_i} = \phi(v^{*}_i) = n^{*}_i$. These are
the numeric vertex couplings the diagrams need.
\end{defn}

\subsection{Step 3 --- the symbolic voltage and the iteration target}

The voltage saddle is read \emph{symbolically} from the model rather than
hardcoded, so that any model whose voltage carries extra terms (feedforward
drives, additional couplings) is picked up automatically. The kernel symbol is
forced to $1$ and the parameters baked in:

\begin{lstlisting}
# pipeline/_mean_field.py:105-118
vstar_subs_dict = model.get('mf_bg_conditions', lambda ns: {})(ns)
vstar_sym = {i: SR(vstar_subs_dict.get(ns.vstar[i], ns.E[i])) for i in ns.pop}
g_to_one = {ns.g: SR(1)}
vstar_baked = {i: SR(vstar_sym[i]).subs(g_to_one).subs(param_subs) for i in ns.pop}

def vstar_num(i, nstar_vec):
    sub = {ns.nstar[j]: float(nstar_vec[j]) for j in ns.pop}
    return float(vstar_baked[i].subs(sub))
\end{lstlisting}

\code{mf\_bg\_conditions} is a callable on the namespace that returns a dict
mapping each saddle LHS symbol to its RHS expression. The line
\code{vstar\_subs\_dict.get(ns.vstar[i], ns.E[i])} reads the declared RHS for
$v^{*}_i$, defaulting to the bare drive \code{ns.E[i]} if the model declares
none. \code{g\_to\_one} is the kernel-normalization substitution discussed above
(\code{ns.g} $\to 1$).

The iteration target for $n^{*}$ is likewise taken from the model when declared,
falling back to the legacy hardcode $n^{*} = \phi(v^{*})$:

\begin{lstlisting}
# pipeline/_mean_field.py:124-141  (abridged)
nstar_rhs_baked = None
if hasattr(ns, 'nstar') and ns.nstar[0] in vstar_subs_dict:
    nstar_rhs_baked = {i: SR(vstar_subs_dict[ns.nstar[i]])
                            .subs(g_to_one).subs(param_subs) for i in ns.pop}

def nstar_target(i, nstar_vec):
    if nstar_rhs_baked is None:
        return phi_num(i, vstar_num(i, nstar_vec))   # legacy: n* = phi(v*)
    ...
\end{lstlisting}

\subsection{Step 4 --- multi-start \code{fsolve} with a singularity guard}

The self-consistency residual is $n^{*}_i - T_i(n^{*})$, and it is wrapped so that
symbolic singularities (a spike-reset RHS like $1/(1-n^{*})$ evaluated near
$n^{*}=1$) become a large \emph{finite} penalty rather than raising --- which lets
the root finder back away and keep searching:

\begin{lstlisting}
# pipeline/_mean_field.py:149-179  (abridged)
def mf_residual(nstar_vec):
    try:
        return [nstar_vec[i] - nstar_target(i, nstar_vec) for i in ns.pop]
    except (ValueError, ZeroDivisionError):
        return [1e10] * len(ns.pop)

initial_guesses = [[0.1]*npop, [0.5]*npop, [1.0]*npop, [0.01]*npop]
sol = None
for x0 in initial_guesses:
    try:
        attempt = fsolve(mf_residual, x0, full_output=True)
    except (ValueError, ZeroDivisionError):
        continue
    sol = attempt
    if attempt[2] == 1:        # ier == 1  =>  converged
        break
\end{lstlisting}

\code{fsolve} (SciPy's wrapper around MINPACK's Powell hybrid method, see
Section~\ref{sec:mf-tools}) returns a 4-tuple under \code{full\_output=True};
element \code{[2]} is the integer convergence flag \code{ier}, and \code{ier == 1}
means clean convergence. The solver tries the guesses $0.1, 0.5, 1.0, 0.01$ in
order and keeps the first that converges. The small leading guesses are
deliberate: they are safe for theories whose saddle sits below a pole at larger
$n^{*}$.

After the solve, explicit finiteness checks guard against a root finder that dove
into a near-pole and returned garbage:

\begin{lstlisting}
# pipeline/_mean_field.py:193-197
if not all(__import__('math').isfinite(x) for x in nstar_vals):
    raise RuntimeError(
        f'solve_mean_field: fsolve returned non-finite n* = '
        f'{nstar_vals!r}.  The mf_eq probably has a singularity '
        f'within the iteration basin; check parameter values.')
\end{lstlisting}

\begin{gotcha}
The legacy solver has \emph{no} multi-root machinery. It tries the four guesses
in order and returns whatever the first converging one lands on. A theory with a
saddle far outside those basins may converge to the \emph{wrong} fixed point or
fail outright. If you need to discover multiple fixed points or pick a specific
branch, that is precisely what the DAE solver is for.
\end{gotcha}

\subsection{Step 5 --- assemble \code{num\_params} and the \code{phi0} rescue}

With the saddle in hand, the voltages and $\phi$-derivatives are evaluated at the
fixed point and everything is poured into \code{num\_params}: the raw parameters,
the saddle values \code{ns.nstar[i]}, \code{ns.vstar[i]}, the GTaS rate
\code{ns.mstar[i]} (guarded behind a presence check), and the numeric derivative
symbols \code{phi1\_i}, \code{phi2\_i}, \dots

\begin{lstlisting}
# pipeline/_mean_field.py:234-258  (abridged)
num_params = dict(param_subs)
for i in ns.pop:
    num_params[ns.nstar[i]] = float(nstar_vals[i])
    num_params[ns.vstar[i]] = float(vstar_vals[i])
...
for i in ns.pop:
    for dk in range(1, taylor_order + 1):
        sym = SR.var(f'phi{dk}_{i+1}')
        if (dk, i) in phi_deriv_vals:
            num_params[sym] = phi_deriv_vals[(dk, i)]
\end{lstlisting}

The last act is the most easily overlooked, and the source comments it heavily.
Template-built theories keep $\phi$ as a \emph{formal} function symbol, so the
$(2,0)$ Poisson-noise sector of the action carries the formal symbol
\code{phi0\_i} (the mean rate $n^{*}_i$) rather than a number. The per-derivative
loop above only resolves \code{phi1\_i} upward. If \code{phi0\_i} is left
symbolic, the noise-source coefficient $-\tfrac12\,\code{phi0\_i}$ never becomes
numeric --- and \emph{the entire correlator collapses to zero}. The fix is to
evaluate the remaining \code{mf\_bg\_conditions} entries (notably the identity
$\code{phi0\_i}\to n^{*}_i$) at the solved saddle:

\begin{lstlisting}
# pipeline/_mean_field.py:272-279
for _lhs, _rhs in vstar_subs_dict.items():
    if _lhs in num_params:
        continue          # saddle vars already resolved above
    try:
        num_params[_lhs] = float(SR(_rhs).subs(g_to_one).subs(num_params))
    except (TypeError, ValueError):
        pass              # still symbolic -- leave it out
\end{lstlisting}

\begin{gotcha}
If a new template forgets to map \code{phi0\_i} $\to$ \code{nstar\_i} in its
\code{mf\_bg\_conditions}, the symbol stays formal, the noise vertex
$-\tfrac12\,\code{phi0\_i}$ never goes numeric, and the correlator silently
evaluates to zero. This is a quiet failure mode, not a crash --- watch for an
all-zero $C(\tau)$ when bringing up a new template-built theory.
\end{gotcha}

\subsection{The heterogeneous branch}

\code{\_solve\_mean\_field\_hetero} (\file{\_mean\_field.py:291}) handles multiple
populations (e.g.\ excitatory and inhibitory), each with its own $\phi$ and its
own saddle arrays. It generalizes the single-pop path along three axes. First, it
identifies which parameters are saddles (those with \code{mean\_field=True}) and
records the population each ranges over (\file{\_mean\_field.py:338-351}). Second,
it classifies saddles as iteration vs.\ compound, preferring the authoritative
\code{model['iteration\_saddles']} list computed at build time and falling back to
``all are iteration'' only if it is missing (\file{\_mean\_field.py:367-375}).
Third, it builds a \emph{flat} fsolve vector that concatenates all
iteration-saddle elements, with helpers \code{\_unflatten} and
\code{\_eval\_compound} to convert between the flat vector and per-saddle arrays
and to evaluate the compound (voltage) saddles from the iteration (rate) values
each step. The kernel-to-one normalization here maps \emph{every} kernel symbol
on the namespace --- scalar, vector, and matrix --- to $1$
(\file{\_mean\_field.py:387-401}).

\begin{note}
The heterogeneous path returns \code{phi\_deriv\_vals: \{\}} (empty). In the
multi-population case the propagator builds the $\phi$ derivatives symbolically
rather than relying on this dict. The same is true of the entire DAE path. If a
consumer ever needs explicit saddle derivatives down those routes, it would have
to populate this dict --- the code does not do so today.
\end{note}

\section{The DAE formulation}

The DAE path declares the dynamics directly, one equation per field, through
\code{TheoryBuilder.equation(lhs, rhs, population)}. The neural example becomes
\begin{align}
  (\tau_i\,\code{Dt} + 1)\,v_i &= E_i + \textstyle\sum_j w_{ij}\,n_j
    && \text{(differential --- carries \code{Dt})}, \\
  n_i &= \phi_i(v_i)
    && \text{(algebraic --- no \code{Dt})}.
\end{align}
The symbol \code{Dt} is a formal stand-in for the time-derivative operator
$\dd/\dd t$. At the stationary mean-field saddle everything is time-independent,
so $\code{Dt}\to 0$ and the system becomes purely algebraic:
\begin{equation}
  v_i = E_i + \sum_j w_{ij}\,n_j,
  \qquad
  n_i = \phi_i(v_i).
\end{equation}
The solver hard-substitutes $\code{Dt}=0$ (and, for spatial theories,
$\code{Laplacian}=0$, since a spatially-uniform background has zero Laplacian)
into every equation and finds the joint root $(v^{*}, n^{*})$. The residual vector
is $F_k = \text{LHS}_k - \text{RHS}_k$, and the root satisfies $F(x^{*}) = 0$.

\subsection{Evaluating the residual via sandboxed \code{eval}}

This is the design choice that defines the DAE solver: it does \emph{not} parse
the user's equation text into a syntax tree. It \code{eval}s the text directly,
in a carefully constructed namespace. The residual builder
\code{\_build\_residual} (\file{\_mean\_field\_dae.py:236}) first expands each
declared equation over its population's index range into a list of single-scalar
residuals, then returns a closure \code{R(x)} that packs the flat state vector
into per-variable NumPy arrays and evaluates each residual:

\begin{lstlisting}
# pipeline/_mean_field_dae.py:268-307  (abridged)
def R(x: np.ndarray) -> np.ndarray:
    state = {var: np.asarray(x[start:end], dtype=float)
             for var, (start, end) in slices.items()}
    ns = {
        **params_np, **state, **phi_callables, **index_sets, **_MATH_NS,
        'Dt':           0,
        'Laplacian':    0,
        'sum':          sum,
        '__builtins__': {},
    }
    out = np.empty(len(expanded), dtype=float)
    for k, (lhs_text, rhs_text, i, pop) in enumerate(expanded):
        ns['i'] = i
        try:
            lhs = eval(lhs_text, ns)
            rhs = eval(rhs_text, ns)
            out[k] = float(lhs - rhs)
        except (ValueError, ZeroDivisionError, OverflowError, FloatingPointError):
            out[k] = 1e10
    return out
\end{lstlisting}

Several details here are load-bearing. The math functions in \code{\_MATH\_NS}
(\file{\_mean\_field\_dae.py:33-38}) are drawn from NumPy (\code{np.tanh},
\code{np.exp}, \dots) so the expression evaluates correctly over both scalars and
arrays. \code{Dt} and \code{Laplacian} are bound to $0$.
\code{'\_\_builtins\_\_': \{\}} sandboxes the evaluation so that only the
explicitly-bound names are visible. Exceptions from singular regions become the
$10^{10}$ penalty that steers the solver away.

\begin{gotcha}
The namespace is passed to \code{eval} as the \emph{globals} argument, not
locals. Python's comprehension-scoping rule (PEP 3104) runs a generator
expression such as \code{sum(w[i,j]*n[j] for j in E)} in its own function scope,
which sees the \code{eval}-globals but \emph{not} the \code{eval}-locals. Passing
the namespace as locals would make the genexpr raise \code{NameError} on
\code{w}, \code{n}, \code{E}. This subtlety is documented in the source at
\file{\_mean\_field\_dae.py:273-279}.
\end{gotcha}

\begin{gotcha}
Because \code{'\_\_builtins\_\_'} is empty, \emph{only} the names the solver
explicitly binds are available (parameters, state variables, $\phi$ callables,
index sets, the \code{\_MATH\_NS} functions, \code{Dt}, \code{Laplacian},
\code{sum}). A user equation that reaches for any other Python builtin --- say
\code{min}, \code{max}, or \code{len} --- will raise, and that exception is
swallowed into the $10^{10}$ penalty, which can masquerade as a failure to
converge. If a DAE theory mysteriously finds no root, an unsupported builtin in
the equation text is a prime suspect.
\end{gotcha}

\begin{note}
\code{\_sage\_to\_python} (\file{\_mean\_field\_dae.py:41}) translates the power
operator \code{\^{}} to \code{**} before any \code{eval}. Sage preparses
\code{\^{}} as exponentiation, but in plain Python \code{\^{}} is bitwise XOR, so
\code{eps*x[i]\^{}3} would raise \code{TypeError: ufunc 'bitwise\_xor' not
supported} on a NumPy float. The replacement is unconditional and safe, since XOR
is never meaningful in a mean-field equation over reals.
\end{note}

\subsection{The transfer functions as numeric callables}

The $\phi$ functions referenced in the equation RHSs are turned into numeric
callables by \code{\_make\_phi\_callables} (\file{\_mean\_field\_dae.py:147}).
Each declared single-argument function becomes a per-population closure that
\code{eval}s the function's expression text in the same sandboxed style. They are
wrapped in a small helper class, \code{\_PhiCallableList}
(\file{\_mean\_field\_dae.py:114}), that supports both indexed access
\code{phi[i](v)} and --- when there is exactly one population position --- a bare
scalar call \code{phi(v)}:

\begin{lstlisting}
# pipeline/_mean_field_dae.py:134-144
def __call__(self, *args, **kwargs):
    if len(self._callables) != 1:
        raise TypeError(
            f'bare call ``f(...)`` requires exactly one population '
            f'position; got {len(self._callables)}.  Use ``f[i](...)`` '
            f'with an explicit index instead.')
    return self._callables[0](*args, **kwargs)
\end{lstlisting}

This mirrors the action-side \code{\_IndexedFormalFunction} convention, so a user
can write \code{phi(v)} in a single-population theory and \code{phi[i](v)} in a
multi-population one.

\begin{gotcha}
\code{\_make\_phi\_callables} only handles \emph{single-argument} functions
(\file{\_mean\_field\_dae.py:174-175}: \code{if len(args\_text) != 1: continue}).
Multi-argument functions --- e.g.\ CGF-style cumulant generating functions --- are
silently skipped here. The DAE solver only needs single-argument transfer
functions for the saddle, so this is by design, but a multi-argument function
referenced in an equation RHS will surface as an unbound name (and hence the
swallowed penalty).
\end{gotcha}

\section{Finding all fixed points}

Nonlinear self-consistency equations generically have \emph{several} solutions:
bistability, up/down states, multiple firing-rate fixed points. The DAE solver
embraces this with a multi-start strategy. The entry point is
\code{solve\_mean\_field\_dae} (\file{\_mean\_field\_dae.py:390}).

\subsection{Seeds, the seed box, and domains}

It draws \code{n\_starts} random initial guesses (default $64$) from a
per-variable box. The box is domain-aware: a variable declared
\code{domain='positive'} is sampled from $[0,\,5s]$, anything else from
$[-3s,\,3s]$, where the scale $s$ is the largest absolute parameter value
(default $1.0$). The default domain, when a field declares none, is
\code{'positive'} --- the Hawkes-like assumption that rates are non-negative:

\begin{lstlisting}
# pipeline/_mean_field_dae.py:336-342
domain = (field_specs.get(var, {}).get('domain') or 'positive')
if domain == 'positive':
    boxes[var] = (0.0, 5.0 * scale)
else:
    boxes[var] = (-3.0 * scale, 3.0 * scale)
\end{lstlisting}

The RNG is the modern \code{np.random.default\_rng(rng\_seed)}
(\file{\_mean\_field\_dae.py:475}), seeded for reproducibility, so the same theory
and parameters always sample the same seeds. A caller can override the box for
any variable via the \code{seed\_box} argument.

\subsection{Multi-start Newton with a re-checked residual}

Each seed is handed to SciPy's \code{root} with the \code{hybr} method (the same
Powell hybrid as \code{fsolve}, but returning a richer result object). The solver
keeps only roots that survive three filters: SciPy's own success flag, a
\emph{re-evaluated} residual below $10^{-7}$, and finiteness:

\begin{lstlisting}
# pipeline/_mean_field_dae.py:481-499  (abridged)
for s in range(n_starts):
    x0 = seeds[s]
    try:
        sol = _scipy_root(R, x0, method='hybr')
    except (ValueError, ZeroDivisionError, OverflowError, FloatingPointError):
        continue
    if not sol.success:
        continue
    resid = R(sol.x)
    if np.max(np.abs(resid)) > 1e-7:
        continue
    if not np.all(np.isfinite(sol.x)):
        continue
    n_converged += 1
    raw_roots.append(np.asarray(sol.x, dtype=float))
\end{lstlisting}

\begin{gotcha}
SciPy's \code{sol.success} is not trusted. The solver re-evaluates \code{R(sol.x)}
and rejects any root whose residual exceeds $10^{-7}$ even when \code{success} is
\code{True} (\file{\_mean\_field\_dae.py:492-494}) --- scipy sometimes reports
success with a residual above its own internal tolerance.
\end{gotcha}

\subsection{Dedup and the deterministic sort}

Many seeds converge to the same root, so the raw list is clustered:
\code{\_dedup\_roots} (\file{\_mean\_field\_dae.py:368}) merges two roots when
$\max|r - \text{kept}| < \code{atol} + \code{rtol}\cdot\text{scale}$, keeping one
representative per cluster. The deduplicated roots are then sorted ascending by
the \emph{first declared physical field's first population index}:

\begin{lstlisting}
# pipeline/_mean_field_dae.py:506-508
first_var = state_vars[0][0]
first_start = slices[first_var][0]
deduped.sort(key=lambda r: float(r[first_start]))
\end{lstlisting}

This gives a deterministic, reproducible labeling: \code{fixed\_point\_index=0}
always means ``the root with the smallest first-field value.'' The list of state
variables comes from \code{\_state\_variables} (\file{\_mean\_field\_dae.py:75}),
which reads them in declaration order and prefers each field's user-facing
\code{natural\_name} (\code{'n'}, \code{'v'}) over its internal MSR name
(\code{'dn'}, \code{'dv'}).

\begin{gotcha}
The sort key is fragile by construction. Reordering the \code{physical\_field}
declarations in the theory changes which field is ``first,'' and therefore which
root \code{fixed\_point\_index=0} selects. If a previously-correct theory starts
picking the wrong branch after an edit, check whether the field declaration order
moved.
\end{gotcha}

\section{Linear stability}

Not every fixed point is a legitimate place to expand. The loop expansion is only
well-defined about a \emph{linearly stable} saddle: perturbations must decay,
otherwise the free propagator carries a growing mode and the perturbative series
diverges. Stability is read from the linearization of the DAE about the root,
implemented in \code{linear\_stability} (\file{\_mean\_field\_dae.py:631}).

\subsection{The matrix pencil}

Linearize $F_k(x,\code{Dt}) = 0$ about $x = x^{*}$, writing a perturbation
$\delta x \propto \ee^{\sigma t}$ so that \code{Dt} acting on it becomes
multiplication by $\sigma$. To first order,
\begin{equation}
  \sum_j \Big[\,\frac{\partial F_k}{\partial x_j}\Big|_{x^{*},\,\code{Dt}=0}
    + \sigma\,\frac{\partial^{2} F_k}{\partial \code{Dt}\,\partial x_j}\Big|_{x^{*}}
    \Big]\,\delta x_j = 0.
\end{equation}
Defining two real $M\times M$ matrices ($M$ = total number of scalar state
variables),
\begin{align}
  B_{kj} &= \frac{\partial F_k}{\partial x_j}\bigg|_{x^{*},\,\code{Dt}=0}
    && \text{(the algebraic Jacobian)}, \\
  A_{kj} &= \frac{\partial^{2} F_k}{\partial \code{Dt}\,\partial x_j}\bigg|_{x^{*}}
    && \text{(the ``mass'' matrix, the coefficient of $\sigma$)},
\end{align}
the linearized condition is $(\sigma A + B)\,\delta x = 0$ --- a
\term{generalized eigenvalue problem}. Rearranged,
$(-B)\,\delta x = \sigma A\,\delta x$, so the stability exponents $\sigma$ are the
eigenvalues of the matrix pencil $(-B, A)$.

\begin{defn}
A \textbf{generalized eigenvalue problem} is $C\,v = \lambda D\,v$ for two
matrices $C, D$. It reduces to the ordinary eigenproblem when $D$ is the identity.
It is the \emph{correct} tool here because the mass matrix $A$ is \emph{singular}:
algebraic equations (those without \code{Dt}) have zero rows in $A$, so $A$ cannot
be inverted and one cannot simply take the eigenvalues of $A^{-1}B$.
\end{defn}

The matrices are built by symbolic differentiation at the root. Note that the
\code{Dt} derivative for $A$ is taken \emph{first}, which drops the polynomial
degree before differentiating with respect to $x_j$:

\begin{lstlisting}
# pipeline/_mean_field_dae.py:798-811
for k, F_k in enumerate(residuals_sym):
    # B[k, j] = dF_k/dx_j at Dt=0, x=x*.
    F_k_alg = F_k.subs({Dt_sym: 0})
    for j, x_j in enumerate(flat_syms):
        d = diff(F_k_alg, x_j)
        B_mat[k, j] = float(d.subs(root_subs))

    # A[k, j] = d^2 F_k / dDt dx_j at x*.  dDt first to drop the degree.
    F_k_dDt = diff(F_k, Dt_sym)
    for j, x_j in enumerate(flat_syms):
        d = diff(F_k_dDt, x_j)
        A_mat[k, j] = float(d.subs(root_subs).subs({Dt_sym: 0}))
\end{lstlisting}

The symbolic residuals \code{residuals\_sym} are themselves built by \code{eval}ing
the equation text in a \emph{Sage} namespace (Section~\ref{sec:mf-tools} contrasts
this with the NumPy namespace the residual path uses), wrapping single-position
state-variable symbol lists in \code{\_FieldScalar}
(\file{theory\_compiler.py:268}) so that bare-name arithmetic such as \code{xt*x}
works alongside indexed access \code{xt[0]*x[0]}.

\subsection{Solving the pencil and the stability verdict}

The generalized eigenproblem is solved with \code{scipy.linalg.eig} called with
two matrix arguments. The infinite/NaN eigenvalues --- the algebraic-constraint
modes from the zero rows of $A$ --- are filtered out, and the verdict is read from
the surviving finite eigenvalues:

\begin{lstlisting}
# pipeline/_mean_field_dae.py:817-832  (abridged)
sigmas, eigvecs = scipy.linalg.eig(-B_mat, A_mat)
finite_mask = np.isfinite(sigmas)
sigmas_finite = sigmas[finite_mask]
EIG_TOL = 1e-9
stable = bool(sigmas_finite.size > 0 and np.all(np.real(sigmas_finite) < -EIG_TOL))
unstable = [complex(s) for s in sigmas_finite if np.real(s) >= -EIG_TOL]
\end{lstlisting}

The convention is: the fixed point is stable iff there is at least one finite
eigenvalue \emph{and every} finite eigenvalue has $\mathrm{Re}(\sigma) < 0$. The
small tolerance \code{EIG\_TOL} $= 10^{-9}$ keeps an eigenvalue that is zero to
roundoff from being misclassified. A growing mode ($\mathrm{Re}\,\sigma \ge 0$)
marks the root unstable.

\begin{note}
The lazy import of Sage at \file{\_mean\_field\_dae.py:676} is deliberate. The DAE
\emph{root-finding} path is pure NumPy/SciPy and never imports Sage; \emph{only}
the stability classification needs symbolic differentiation, so Sage (a heavy
import) is loaded only when \code{linear\_stability} actually runs.
\end{note}

\subsection{The \code{stability\_analysis} toggle}

Whether the eigenvalue machinery runs at all is gated by
\code{model['stability\_analysis']} (default \code{False}), set on the builder via
\code{.stability\_analysis(...)} (\file{pipeline/theory.py:1187}):

\begin{description}
  \item[OFF (the default).] Skip stability entirely. Every converged root is
        selectable, and \code{fixed\_point\_index} ranges over all of them. Each
        record is tagged \code{stable=None} (to distinguish ``not classified''
        from ``classified unstable''). This is the \emph{correct} choice when
        every equation is algebraic --- voltages integrated out, no \code{Dt}
        anywhere --- because then $A \equiv 0$, every eigenvalue is infinite, and
        ``linear stability'' is meaningless. A Hawkes-type rate model is the
        canonical example.
  \item[ON.] Classify every root, filter to the linearly-stable subset, expose
        \code{fixed\_point\_index} over \emph{that} subset, and \emph{raise} if no
        stable root exists. Required for bistable or multi-saddle theories where
        you must pick the stable branch.
\end{description}

\begin{lstlisting}
# pipeline/_mean_field_dae.py:569-586  (abridged)
if stab_on:
    selectable_records = [r for r in all_root_records if r['stable']]
    if not selectable_records:
        raise ValueError(
            "... found N fixed point(s) but NONE were classified linearly "
            "stable.  ... (If your equations are all algebraic -- voltages "
            "integrated out, no Dt anywhere -- set ``.stability_analysis(False)`` "
            "on the TheoryBuilder; ``fixed_point_index`` will then range over "
            "every converged root.)")
\end{lstlisting}

\begin{gotcha}
Turning \code{stability\_analysis} ON for an all-algebraic theory is a guaranteed
failure: with no \code{Dt} anywhere, $A \equiv 0$, every eigenvalue is infinite,
the stable subset is empty, and the solver raises ``NONE were classified linearly
stable.'' The error message itself tells you the fix --- call
\code{.stability\_analysis(False)}.
\end{gotcha}

\begin{gotcha}
Per-root \code{linear\_stability} calls are wrapped in \code{except Exception}
(\file{\_mean\_field\_dae.py:542-552}): a stability failure marks that root
\emph{unstable} (with an \code{'error'} annotation) rather than aborting the whole
solve. The consequence is that a bug in the stability code can make a genuinely
stable root silently vanish from the selectable set.
\end{gotcha}

\subsection{Index selection and clamping}

Finally, \code{fixed\_point\_index} selects from the selectable subset. An
out-of-range index is \emph{clamped} with a warning rather than raising:

\begin{lstlisting}
# pipeline/_mean_field_dae.py:594-603
if fixed_point_index < 0 or fixed_point_index >= n_selectable:
    clamped = max(0, min(fixed_point_index, n_selectable - 1))
    warnings.warn(
        f"... requested fixed_point_index={fixed_point_index} but only "
        f"{n_selectable} {index_label} found ...; using index {clamped}.",
        stacklevel=2)
    fixed_point_index = clamped
\end{lstlisting}

\begin{gotcha}
An out-of-range \code{fixed\_point\_index} never errors --- it clamps to the valid
range, and a single \code{warnings.warn} is the only signal. In a noisy notebook
that warning is easy to miss, so if you asked for the third fixed point and got
the first, check for a clamp warning.
\end{gotcha}

\section{From saddle to diagrams}

The handoff to the perturbative expansion is \emph{purely substitutional}. The
propagator kernel $K_{\mathrm{ft}}(\omega)$ and the vertices are symbolic Sage
expressions containing parameter symbols and saddle symbols. The mean-field
solver's deliverable, \code{num\_params}, is a dict
$\{\text{SR symbol}\to\text{float}\}$ covering every numeric parameter (\code{E1},
\code{w12}, \code{tau}, \dots), every saddle value (\code{nstar[i]},
\code{vstar[i]}, \code{mstar[i]}), every transfer-function derivative at the
saddle (\code{phi1\_i}, \code{phi2\_i}, \dots, on the legacy path), and the
\code{phi0\_i} $\to$ \code{nstar\_i} identities.

Downstream, \code{compute\_poles\_and\_residues(prop, num\_params)}
(\file{pipeline/\_propagator.py:851}) substitutes \code{num\_params} into
$K_{\mathrm{ft}}$, obtaining a numeric matrix-valued rational function of
$\omega$. It then finds the retarded poles $\omega_k$ (those with
$\mathrm{Im}\,\omega > 0$) and the residue matrices that assemble the free
propagator $G(\tau)$ (the sibling chapter, Chapter~\ref{ch:propagator}, covers
that). From there the loop corrections build the connected cumulant $C(\tau)$. The
saddle enters as \emph{the numbers that turn symbolic diagrams into numeric ones}
--- there is no further re-solve; the saddle is frozen and the loops are
corrections about it.

\begin{note}
The symbolic-derivative renaming that \emph{creates} the \code{phi1\_i},
\code{phi2\_i} symbols lives upstream in the compiler, in
\code{ns.\_deriv\_rename\_subs} (\file{pipeline/theory\_compiler.py:1048}): the
compiler differentiates each formal function at its saddle and registers a rename
$\langle\text{derivative}\rangle \to \code{SR.var('phik\_<i+1>')}$. The legacy
solver then supplies the \emph{numeric} value of each \code{phik\_i}. The DAE
compatibility wrapper does \emph{not} populate \code{phi\_deriv\_vals} (it returns
\code{\{\}}); the propagator builds those derivatives symbolically in that path.
\end{note}

\subsection{The compatibility wrapper}

\code{solve\_mean\_field\_dae\_compat} (\file{\_mean\_field\_dae.py:853}) is what
\code{compute.py} actually calls on the DAE path. It is a drop-in replacement for
the legacy \code{solve\_mean\_field}: it runs \code{solve\_mean\_field\_dae},
builds \code{param\_subs} in the legacy Sage-symbol convention, and bakes the
selected root's saddle values into \code{num\_params} by looking up the Sage
arrays on \code{ft.\_ns}:

\begin{lstlisting}
# pipeline/_mean_field_dae.py:943-950
num_params = dict(param_subs)
for saddle_name, vals in mf_values.items():
    sr_array = getattr(ft._ns, saddle_name, None)
    if sr_array is None:
        continue
    for i, v in enumerate(vals):
        num_params[sr_array[i]] = float(v)
\end{lstlisting}

It returns the legacy-shape keys (\code{nstar\_vals}, \code{vstar\_vals},
\code{phi\_deriv\_vals=\{\}}, \code{num\_params}, \code{param\_subs},
\code{saddle\_values}) \emph{plus} the DAE extras (\code{mf\_all\_roots},
\code{mf\_stable\_roots}, \code{mf\_unstable\_roots}, \code{mf\_index\_used},
\code{state\_var\_order}, \code{n\_seeds\_converged}). \code{compute\_cumulants}
surfaces those extras onto its returned result dict, so a notebook can introspect
every fixed point and its stability --- which makes them the natural first stop
when a DAE saddle behaves unexpectedly.

\section{External tooling primer}\label{sec:mf-tools}

This subsystem leans on four external pieces. Since the manual assumes you know
the physics but not necessarily the tooling, here is each from scratch, with the
exact import lines.

\subsection{SageMath and the symbolic ring \code{SR}}

\term{SageMath} is a large open-source mathematics system built on top of Python;
it bundles many computer-algebra packages behind one Python API. The piece this
subsystem uses is its \emph{symbolic ring} \code{SR} (the ``Symbolic Ring''),
which lets you build, differentiate, and substitute into expressions made of
symbols (\code{x}, \code{E1}, \code{nstar}) rather than numbers. Think of it as a
calculator that does algebra with letters.

\begin{itemize}
  \item \code{SR.var('name')} returns (and globally registers) a symbolic
        variable; calling it twice with the same name returns the \emph{same}
        symbol. This is why a substitution dict keyed by \code{SR.var('E1')}
        matches the \code{E1} the propagator built separately.
  \item \code{SR(value)} coerces a Python number or another expression into the
        ring.
  \item \code{diff(expr, var, k)} computes the $k$-th symbolic derivative.
  \item \code{.subs(dict)} substitutes symbols $\to$ values (or expressions);
        chained \code{.subs(...).subs(...)} applies left to right. After all
        symbols are numbers, \code{float(expr)} evaluates to a Python float.
\end{itemize}

The legacy solver imports Sage at module top (\file{\_mean\_field.py:10}:
\code{from sage.all import SR, diff}). The DAE solver imports it \emph{lazily},
inside \code{linear\_stability} (\file{\_mean\_field\_dae.py:676-681}), because
only stability needs symbolic differentiation.

\begin{note}
The contrast between the two math namespaces is the key to the DAE design. The
\emph{residual} path binds \code{tanh} $\to$ \code{np.tanh} (\code{\_MATH\_NS}),
so an equation evaluates to a \emph{float}. The \emph{stability} path binds
\code{tanh} $\to$ Sage's \code{tanh} (\code{sage\_math\_ns}), so the same equation
stays \emph{symbolic} and can be differentiated. The very same text is evaluated
two different ways for two different purposes.
\end{note}

\subsection{SciPy --- numerical solvers}

\term{SciPy} is the standard scientific-computing library, built on NumPy. It
provides numerical (not symbolic) algorithms; nothing in it knows about symbols.

\begin{itemize}
  \item \code{scipy.optimize.fsolve} (legacy solver, \file{\_mean\_field.py:11})
        finds a root of $f(x)=0$ from a starting guess using MINPACK's modified
        Powell hybrid method. With \code{full\_output=True} it returns a 4-tuple
        \code{(x, infodict, ier, mesg)}; \code{ier == 1} signals clean
        convergence.
  \item \code{scipy.optimize.root} (DAE solver, imported as \code{\_scipy\_root}
        at \file{\_mean\_field\_dae.py:27}) is the newer, more general
        root-finder. With \code{method='hybr'} it is the same Powell hybrid, but
        it returns a richer \code{OptimizeResult} with \code{.x} and
        \code{.success}.
  \item \code{scipy.linalg.eig} (stability, \file{\_mean\_field\_dae.py:682})
        solves eigenvalue problems. Called with \emph{two} matrices,
        \code{eig(C, D)}, it solves the generalized problem $C v = \lambda D v$.
        The code calls \code{eig(-B, A)}, so the returned eigenvalues \emph{are}
        the stability exponents $\sigma$.
\end{itemize}

\subsection{NumPy --- numerical arrays}

\term{NumPy} is the foundational array library: the \code{ndarray}, elementwise
math (\code{np.tanh}, \code{np.exp}), and reductions (\code{np.max}, \code{np.abs},
\code{np.all}, \code{np.isfinite}). It is the substrate SciPy stands on. The DAE
solver uses it to coerce parameters into arrays so that \code{w[i,j]} (2-D) and
\code{Em[i]} (1-D) work under \code{eval} (\code{\_numpy\_params},
\file{\_mean\_field\_dae.py:198}), to draw the random seeds via the modern
generator \code{np.random.default\_rng} (\file{\_mean\_field\_dae.py:475}), and to
pack and filter the residual and root vectors.

\subsection{\code{eval} and the namespace conventions}

\code{eval(source, globals\_dict)} executes a Python expression string in a
supplied namespace and returns its value. The DAE solver \code{eval}s the
(Sage-to-Python translated) equation text directly. To keep that safe and
predictable it passes a \emph{closed} namespace with
\code{'\_\_builtins\_\_': \{\}} (so no builtins leak in) and binds exactly the
names an expression may reference. The two non-obvious rules --- pass the namespace
as \emph{globals} (the PEP-3104 comprehension-scoping fix) and the empty-builtins
sandbox --- were covered in the gotchas above; they are the difference between an
equation that evaluates and one that raises a swallowed \code{NameError}.

\section{Data structures}

For reference, the dictionaries that cross this subsystem's boundaries.

\subsection{The legacy solver return dict}

\begin{lstlisting}
{
  'nstar_vals':     [float, ...],       # rate saddles (per pop / concatenated)
  'vstar_vals':     [float, ...],       # voltage saddles
  'phi_deriv_vals': {(dk, i): float},   # phi^(dk)(v*_i); {} on the hetero path
  'num_params':     {SR var: float},    # THE handoff dict
  'param_subs':     {SR var: float},    # raw user params only
  # hetero path also adds:
  'saddle_values':  {saddle_name: [float]},
  'saddle_info':    {saddle_name: {'pop', 'size', 'sr_array'}},
}
\end{lstlisting}

\subsection{The DAE solver return dict}

\begin{lstlisting}
{
  'mf_values':         {'<var>star': [floats]},   # the selected root
  'mf_all_roots':      [ {'values': {...},         # every root, sorted, annotated
                          'stable': bool|None,
                          'eigenvalues_finite': np.ndarray[complex]}, ... ],
  'mf_stable_roots':   [ {'<var>star': [...]}, ... ],  # selectable subset
  'mf_unstable_roots': [ {...}, ... ],             # only meaningful when stab ON
  'mf_index_used':     int,                        # index after clamping
  'state_var_order':   [var_name, ...],            # declaration order (sort key)
  'n_seeds_converged': int,                        # pre-dedup convergence count
  'stability_analysis': bool,                      # was classification run?
}
\end{lstlisting}

\subsection{The \code{linear\_stability} return dict}

\begin{lstlisting}
{
  'stable':               bool,
  'eigenvalues_finite':   np.ndarray[complex],   # dynamical modes sigma
  'eigenvalues_all':      np.ndarray[complex],   # includes +/-inf (constraints)
  'A':                    np.ndarray[float] (MxM),  # d^2 F / dDt dx   (pencil)
  'B':                    np.ndarray[float] (MxM),  # dF / dx at Dt=0  (Jacobian)
  'unstable_eigenvalues': [complex, ...],        # finite sigma with Re >= 0
}
\end{lstlisting}

\section{Pitfalls and diagnostics}

A consolidated checklist of the failure modes scattered through this chapter, for
when a saddle misbehaves:

\begin{itemize}
  \item \textbf{Two solvers, one switch.} Which solver runs is purely
        \code{model.get('equations')} truthiness. A theory carrying both
        \code{.equation(...)} and \code{mf\_bg\_conditions} silently takes the DAE
        path.
  \item \textbf{Singularities are penalized, not raised.} Both solvers turn
        symbolic singularities (spike-reset $1/(1-n^{*})$) into a $10^{10}$
        residual to let the solver back away. But a root that \emph{converges}
        into a near-pole is caught by explicit finiteness checks, which raise a
        \code{RuntimeError} with guidance (\file{\_mean\_field.py:193-217}).
  \item \textbf{\code{phi0\_i} must be resolved or the correlator is zero.}
        Template-built theories keep $\phi$ formal; the noise coefficient
        $-\tfrac12\,\code{phi0\_i}$ stays symbolic unless \code{phi0\_i} $\to$
        \code{nstar\_i} is baked in (\file{\_mean\_field.py:260-279}).
  \item \textbf{DAE \code{eval} namespace must be globals, with empty builtins.}
        Get the globals/locals wrong and comprehensions raise \code{NameError};
        reference an unbound builtin and the exception is swallowed into a
        $10^{10}$ penalty that looks like non-convergence.
  \item \textbf{\code{stability\_analysis} must be OFF for all-algebraic
        theories.} Otherwise $A\equiv 0$, no eigenvalue is finite, and the solver
        raises ``NONE were classified linearly stable.''
  \item \textbf{\code{fixed\_point\_index} clamps with a warning, never errors.}
        An out-of-range index is silently clamped.
  \item \textbf{Sort-key fragility.} Roots are ordered by the first declared
        field's first index; reordering field declarations changes which root is
        index $0$.
  \item \textbf{Swallowed per-root stability failures.} A bug in
        \code{linear\_stability} can make a genuinely stable root vanish from the
        selectable set, because the per-root call is wrapped in
        \code{except Exception}.
  \item \textbf{Reading \code{mf\_all\_roots}.} When a DAE theory picks an
        unexpected saddle, inspect \code{res['mf\_all\_roots']} (every root with
        its stability flag and finite eigenvalues) and call
        \code{linear\_stability(model, fundamental, root)} directly on a root to
        see its $A$, $B$, and eigenvalues.
\end{itemize}

\begin{note}
\code{Laplacian} and \code{Dt} are both hard-substituted to $0$ on the
homogeneous mean field. Stability to \emph{spatial} ($k\neq 0$) perturbations ---
the pattern-formation question --- is \emph{not} handled here; it lives in the
propagator's momentum dependence, as the source notes at
\file{\_mean\_field\_dae.py:769-773}.
\end{note}

% ---------------------------------------------------------------------
%  Chapter 9 : Diagram Enumeration -- Prediagrams, nauty, and the
%              loop-order completeness bound
% ---------------------------------------------------------------------
\chapter{Diagram Enumeration}\label{ch:enumeration}

Every loop calculation in this manual ultimately reduces to a sum over Feynman
diagrams. Before the engine can evaluate that sum, it must first \emph{know which
diagrams are in it} --- and it must know that the list is \emph{complete}, with
nothing missing and nothing double-counted. That bookkeeping is the job of the
subsystem dissected in this chapter. Given just two integers --- $k$, the number
of external legs of the correlation function we want, and $\ell$, the loop order
of the correction we are computing --- it produces the complete, deduplicated
list of every bare graph topology that could possibly contribute, \emph{before
any physical field labels are attached}. Those bare graphs are called
\term{prediagrams}.

The whole subsystem lives in two files and we read both in full:
\file{msrjd/enumeration/loop\_diagram\_enumeration.py} (the topological
generator) and \file{msrjd/enumeration/degree\_scan.py} (a small feedback bridge
back to the algebra). The generator is \emph{theory-agnostic}: it never looks at
the action, the fields, or the parameters. It answers a purely combinatorial
question --- \emph{what graph shapes are even possible at $(k,\ell)$?} --- and
leaves the physical question --- \emph{which of those shapes the theory actually
populates, and with what propagators and vertices} --- to the type-assignment
subsystem in the next chapter.

The reader is assumed to know \msrjd{} field theory but \emph{not} the tooling.
The heavy lifting here is done by \term{SageMath} and, inside it, by the graph
canonical-labeling engine \term{nauty}; both are introduced from first
principles, as are the handful of Python standard-library helpers
(\code{itertools}, \code{collections}, \code{concurrent.futures}) the module
leans on. When in doubt we over-explain.

\section{What a prediagram is, and why we strip the physics first}

In \msrjd{} (Martin--Siggia--Rose--Janssen--De~Dominicis) field theory, a genuine
Feynman diagram is a graph in which \emph{every edge is a specific propagator}
and \emph{every internal vertex is a specific interaction monomial} from the
action. The doubled field content of \msrjd{} means there are two flavors of
propagator: a \emph{response} propagator $\avg{\phi\,\tilde\phi}$ (causal,
directed --- it carries a definite time direction, from cause to effect) and a
\emph{correlation} propagator $\avg{\phi\,\phi}$. The vertices come from the
Taylor-expanded interaction part of the action --- a cubic vertex $\tilde\phi\,
\phi^2$, a noise vertex $\tilde\phi^2$, and so on.

A \term{prediagram} strips all of that physics away and keeps only the
combinatorial skeleton. Concretely it retains exactly three things:

\begin{itemize}
  \item the \textbf{shape} of the graph --- which vertices connect to which;
  \item the \textbf{leaf-vs-internal} distinction --- which vertices are
        \emph{external legs} of the correlation function, and which are
        \emph{internal} interaction points; and
  \item an \textbf{orientation} on every edge (an arrow $u \to v$), because
        \msrjd{} propagators are causal: a response propagator flows in a
        definite time direction, so the edges of a valid diagram must carry
        arrows.
\end{itemize}

So a prediagram is, in one phrase, ``a Feynman diagram with the physics not yet
filled in.''

\begin{defn}
A \textbf{prediagram} is a bare, oriented graph topology: a partition of vertices
into \emph{leaves} (external legs) and \emph{internal} vertices, together with an
arrow on every edge, but with \emph{no} propagator types and \emph{no} vertex
monomials assigned yet. It is the output object of the enumeration subsystem,
carried as the 4-tuple \code{(D, G, leaves, internal)}.
\end{defn}

\subsection{Why separate the topology from the physics?}

Two reasons, one practical and one conceptual.

The conceptual reason is that the set of admissible \emph{shapes} depends only on
$k$ and $\ell$ --- it is the same for an Ornstein--Uhlenbeck particle, the
Kardar--Parisi--Zhang equation, or a coupled multi-field reaction--diffusion
model. Solving the topology problem once, model-independently, means the
expensive graph-isomorphism work is never repeated per theory.

The practical reason is that it cleanly factors a hard combinatorial problem into
two tractable stages. Stage one (this chapter) enumerates shapes. Stage two
(Chapter~\ref{ch:type-assignment}) walks each shape and tries every consistent
way to label its edges with propagators and its vertices with monomials, throwing
away the labelings the theory does not support. Keeping the two apart means each
stage has a single, checkable correctness criterion: \emph{completeness} for the
topology stage (we miss no shape), and \emph{validity} for the typing stage (we
keep only labelings the action provides).

\subsection{Where it sits in the pipeline}

The project is organized into lettered ``build phases'' (documented in
\file{docs/BUILD\_PHASE\_OUTLINES.md}). This subsystem is the topological
generator --- it is invoked as phase~\code{[5/7]} of the temporal trace you saw
in Chapter~\ref{ch:quickstart} (``Enumerate prediagrams''). Schematically:

\begin{lstlisting}[style=console]
  Phase 1 (FieldTheory.expand)     Phase 2 (THIS SUBSYSTEM)        Phase D/E
  ----------------------------     -------------------------       ------------
  Action -> Taylor-expanded        enumerate_prediagrams(k,ell)    type_assignment
  vertices + noise kernel,     --> trees -> topologies ->      --> .enumerate_typed
  saved as a "theory"              prediagrams                     _diagrams
                                        |
                                        v
                               Phase C: degree_scan
                               .ensure_taylor_order
                               (feedback to Phase 1)
\end{lstlisting}

\begin{itemize}
  \item \textbf{What feeds it:} essentially nothing but the two integers $k$ and
        $\ell$. The generation is theory-agnostic; the action enters only later.
  \item \textbf{What consumes it:} two downstream readers.
        \code{msrjd/diagrams/type\_assignment.py} is the direct consumer --- its
        \code{enumerate\_typed\_diagrams(prediagram, ...)} unpacks one prediagram
        tuple \code{(D, G, leaves, internal)} and produces every valid fully-typed
        diagram for it. And \code{msrjd/enumeration/degree\_scan.py} (build
        phase~C) scans the same tuples for their maximum vertex degree, to decide
        whether the saved theory was Taylor-expanded far enough --- the feedback
        bridge of \S\ref{sec:degree-scan}.
\end{itemize}

\subsection{The three-stage refinement}

The generator works in three successive refinements. Each stage is a
\emph{filter-and-deduplicate} pass on the previous one:

\begin{enumerate}
  \item \textbf{Trees} --- connected acyclic graphs with the right number of
        leaves. A tree is the loop-free ``spanning skeleton'' of every possible
        diagram. Trees have zero loops, so to reach $\ell$ loops we will later add
        $\ell$ extra edges.
  \item \textbf{Topologies} --- undirected multigraphs obtained by adding exactly
        $\ell$ extra edges to a tree (creating exactly $\ell$ independent
        cycles), keeping only those that satisfy \msrjd's structural constraints,
        then removing isomorphic duplicates.
  \item \textbf{Prediagrams} --- topologies with every edge given an arrow,
        keeping only physically admissible (causal, acyclic) orientations, then
        removing isomorphic duplicates again.
\end{enumerate}

The public entry points (\code{enumerate\_all}, \code{enumerate\_prediagrams},
and friends) wrap these three stages; \S\ref{sec:api} is the reference.

\section{A graph-theory primer for field theorists}

This subsystem is pure graph theory, so we build the vocabulary from the ground
up. A reader who knows \msrjd{} but not the combinatorics should be able to
follow every later step.

\subsection{Vertices, edges, degree, leaves}

A \term{graph} is a set of \emph{vertices} (points) and \emph{edges}
(connections between pairs of points). The \term{degree} of a vertex $v$, written
$\deg(v)$, is the number of edge-ends touching it. If two vertices are joined by
a \emph{double edge} --- a multi-edge --- that contributes $2$ to each of their
degrees.

We partition the vertices by degree:

\begin{itemize}
  \item \term{Leaves} are the degree-1 vertices. In our setting these are the
        \emph{external legs} of the correlation function; a $k$-point function has
        exactly $k$ of them in the final diagram.
  \item \term{Internal vertices} are the degree-$\ge 2$ vertices. These are the
        interaction points.
  \item Internal vertices split further into \emph{degree-2} vertices ($V_2$) and
        \emph{degree-$\ge 3$} vertices ($V_3$, written \code{degree\_3plus} in the
        code).
\end{itemize}

In code this partition is computed by \code{classify\_vertices\_sage}
(\file{loop\_diagram\_enumeration.py:32}), which returns the four lists
\code{(leaves, internal, degree\_2, degree\_3plus)} by one filtered
list-comprehension per category over \code{G.vertices()}:

\begin{lstlisting}
def classify_vertices_sage(G):
    leaves       = [v for v in G.vertices() if G.degree(v) == 1]
    internal     = [v for v in G.vertices() if G.degree(v) > 1]
    degree_2     = [v for v in G.vertices() if G.degree(v) == 2]
    degree_3plus = [v for v in G.vertices() if G.degree(v) >= 3]
    return leaves, internal, degree_2, degree_3plus
\end{lstlisting}

Note \code{internal} is the union of \code{degree\_2} and \code{degree\_3plus}.
This is the single most-called helper in the module --- nearly every constraint
check begins by classifying the vertices.

\subsection{Trees and the loop number}

A \term{tree} is a connected graph with no cycles. The defining counting
identity of a tree is
\begin{equation}
  |E| \;=\; |V| - 1
  \qquad\text{(a tree on $|V|$ vertices has exactly $|V|-1$ edges).}
  \label{eq:tree-edges}
\end{equation}
For \emph{any} connected graph $G$, the number of \term{independent cycles} ---
equivalently the \term{loop number} or first Betti number --- is
\begin{equation}
  L(G) \;=\; |E| - |V| + 1.
  \label{eq:loop-number}
\end{equation}
This is exactly \code{count\_cycles\_sage} (\file{loop\_diagram\_enumeration.py:68}):

\begin{lstlisting}
def count_cycles_sage(G):
    if not G.is_connected():
        return -1
    return G.size() - G.order() + 1
\end{lstlisting}

In SageMath, \code{G.size()} is the number of edges $|E|$ and \code{G.order()} is
the number of vertices $|V|$. The function returns $-1$ as a sentinel when the
graph is disconnected, since \eqref{eq:loop-number} is only meaningful for a
connected graph. (That $-1$ is load-bearing: it is never equal to any $\ell\ge0$,
so a disconnected candidate is automatically rejected by the
``\code{== ell}'' loop test downstream --- see Gotcha~\ref{gotcha:sentinel}.)

\begin{note}
The consequence we use \emph{everywhere}: starting from a tree ($L=0$) and adding
$\ell$ extra edges \emph{without adding new vertices} raises $|E|$ by $\ell$ and
leaves $|V|$ unchanged, so by \eqref{eq:loop-number} it raises the loop number by
exactly $\ell$. This is the engine of the whole topology stage: \emph{take a
tree, add $\ell$ edges, demand $L(G)=\ell$.}
\end{note}

\section{The leaf surplus \texorpdfstring{$j$}{j} and the loop budget}

The final diagram must have exactly $k$ external legs. But the \emph{intermediate
trees} may temporarily carry \emph{more} than $k$ leaves, because some of those
extra leaves will be ``eaten'' when we add the loop edges: an added edge can
attach to a former leaf, raising its degree above 1 so it stops being a leaf.

The code calls this surplus $j$. A tree is generated with $k+j$ leaves, where
$j\ge 0$ counts the leaves that the $\ell$ added edges will consume.

\subsection{How large can the surplus be?}

Each added edge has two endpoints, so $\ell$ added edges provide $2\ell$
endpoint-attachments. In the most leaf-consuming case, attaching those endpoints
to leaves is what converts leaves into internal vertices. The code caps the
surplus at
\begin{equation}
  j_{\max} \;=\; \ell + \left\lfloor \tfrac{\ell}{2} \right\rfloor
  \;=\; \left\lfloor \tfrac{3\ell}{2} \right\rfloor,
  \label{eq:jmax}
\end{equation}
implemented verbatim as \code{j\_max = ell + (ell // 2)}
(\file{loop\_diagram\_enumeration.py:124}; \code{//} is Python's integer
floor-division). Intuitively this is the maximum number of tree-leaves the
$\ell$ loop-edges can absorb across all valid decompositions; the $3\ell/2$ form
is the integer ceiling on how many degree-1 nodes the loop structure can use up.
For $\ell=1$ this gives $j_{\max}=1$; for $\ell=2$, $j_{\max}=3$; for $\ell=3$,
$j_{\max}=4$.

\section{The degree bounds and the completeness theorem}
\label{sec:bounds}

Not every tree with $k+j$ leaves can be the skeleton of a valid \msrjd{} diagram.
Two structural bounds prune the tree search. Both are computed at the top of the
$j$-loop in \code{generate\_trees\_with\_constraints}
(\file{loop\_diagram\_enumeration.py:118}).

\subsection{Bound on degree-\texorpdfstring{$\ge 3$}{>=3} vertices}

The first bound limits the number of high-degree (branching) vertices:
\begin{equation}
  |V_3| \;\le\; k + j - 2 \;=\; \code{v3\_max}
  \qquad (\file{loop\_diagram\_enumeration.py:128}).
  \label{eq:v3max}
\end{equation}
This is the standard fact about trees: a tree with $m$ leaves has at most $m-2$
vertices of degree $\ge 3$ (the extreme case being a binary tree, where every
internal branching is exactly three-way). Here $m = k+j$, so $\code{v3\_max} =
k+j-2$.

\subsection{Bound on degree-2 vertices: the completeness bound}

The second bound is the subtle one --- the \term{enumeration completeness bound}
referenced throughout the project. The code
(\file{loop\_diagram\_enumeration.py:139}) sets
\begin{lstlisting}
v2_max = 2 * v3_max + 3 * ell - k - 3 * j + 3
\end{lstlisting}
and the in-code comment records that this algebraically equals the \emph{proven}
upper bound on the number of degree-2 vertices of the final topology $G$:
\begin{equation}
  |V_2^{T}| \;\le\; k + 3\ell - j - 1.
  \label{eq:v2max}
\end{equation}
Let us verify the algebra by substituting $\code{v3\_max}=k+j-2$:
\begin{align*}
  2\,\code{v3\_max} + 3\ell - k - 3j + 3
   &= 2(k+j-2) + 3\ell - k - 3j + 3 \\
   &= 2k + 2j - 4 + 3\ell - k - 3j + 3 \\
   &= k - j - 1 + 3\ell \;=\; k + 3\ell - j - 1. \qquad\checkmark
\end{align*}
So the code's \code{v2\_max} is exactly the bound \eqref{eq:v2max}.

\subsection{Why \texorpdfstring{$k+3\ell-j-1$}{k+3l-j-1} is the right bound}

The full proof lives in the paper appendix; the in-code comment
(\file{loop\_diagram\_enumeration.py:129--138}) and the project note
\file{project\_enumeration\_bound\_fix.md} summarize the three ingredients:

\begin{enumerate}
  \item a \textbf{degree-partition identity} --- the handshake/incidence count
        relating leaves, $V_2$, $V_3$, edges, and the loop number;
  \item the \textbf{orientability constraint} $|V_2^{G}| \le k + \ell - 1$ --- a
        degree-2 internal vertex of an \msrjd{} diagram must be orientable as
        one-in/one-out, and there is a hard ceiling on how many such vertices a
        valid causal orientation can admit; and
  \item an \textbf{incidence-counting} bound $\theta \le 2\ell - j$, where
        $\theta$ counts certain doubled-edge / multi-edge incidences.
\end{enumerate}
Combining the three yields \eqref{eq:v2max}.

\begin{defn}
The \textbf{completeness bound} $|V_2^{T}| \le k+3\ell-j-1$ is the proven upper
bound on the number of degree-2 vertices a valid topology can have at fixed
$(k,\ell,j)$. Pruning the tree search by this bound is what \emph{guarantees}
the search misses no contributing topology --- this is the sense in which the
enumeration is provably complete.
\end{defn}

\subsection{The bound-correction story: a false lemma at \texorpdfstring{$\ell\ge 3$}{l>=3}}

This bound has a cautionary history worth telling, because it encodes a principle
the whole subsystem now lives by. An \emph{earlier} version of the code used a
\emph{tighter} bound carrying an extra $-\lfloor j/3\rfloor$ term:
\begin{equation}
  \code{v2\_max\_OLD} \;=\; k + 3\ell - j - 1 - \left\lfloor \tfrac{j}{3}
  \right\rfloor .
  \label{eq:v2old}
\end{equation}
That extra tightening came from a claimed lemma $\theta \le 2\ell - j -
\lfloor j/3\rfloor$, which is \emph{false at $\ell\ge 3$}. The counterexample,
recorded in the comment at \file{loop\_diagram\_enumeration.py:133} (and in
\file{scratch/bound\_check.py}), is \emph{three doubled-edge bubbles on a hub},
$|V|=8$, where every decomposition has $j=3$ and $\theta = 3 > 2\ell - j -
\lfloor j/3\rfloor$.

The insidious part: the false lemma never actually \emph{bound} at the small
orders that had been tested. Exhaustive enumeration at $(k,\ell)\in\{(2,1),(3,1),
(2,2),(3,2),(4,1),(2,3)\}$ showed a slack of $\ge 2$, so the tighter
\eqref{eq:v2old} and the true \eqref{eq:v2max} produced \emph{identical output} at
every tested order. Only the proven bound \eqref{eq:v2max}, however,
\emph{guarantees} completeness at $\ell\ge 3$.

\begin{gotcha}
\label{gotcha:proof-not-slack}
The headline lesson, recorded in the project memory: \textbf{a prune must follow
only a proven bound; empirical slack at small orders is not a proof.} If anyone
re-tightens \code{v2\_max} (or \code{j\_max}), the new expression must come from a
\emph{proven} statement. The June~2026 fix (commit \code{40454e7}) removed the
false $-\lfloor j/3\rfloor$ term; it changed only the \emph{proof}, not the
observed output at the tested orders. The deduplicated \emph{topology} counts at the tested orders ---
$(k,\ell)=(2,1)\to 9$, $(3,1)\to 67$, $(2,2)\to 289$ --- are the same before and
after the fix. (The comment at \file{loop\_diagram\_enumeration.py:138} records
the \emph{orders} that were exhaustively slack-verified,
$\{(2,1),(3,1),(2,2),(3,2),(4,1),(2,3)\}$, not these counts.)
\end{gotcha}

\section{Stage 1: generating the trees}

With the bounds in hand we can read the tree generator. The function
\code{generate\_trees\_with\_constraints(k, ell, max\_vertices\_search=50)}
(\file{loop\_diagram\_enumeration.py:118}) returns a list of triples
\code{(tree, j, num\_leaves)} --- every non-isomorphic tree satisfying the degree
constraints, tagged with its surplus $j$ and its leaf count $\code{num\_leaves} =
k+j$.

\begin{lstlisting}
def generate_trees_with_constraints(k, ell, max_vertices_search=50):
    valid_trees = []
    j_max = ell + (ell // 2)

    for j in range(0, j_max + 1):
        num_leaves = k + j
        v3_max = k + j - 2
        v2_max = 2 * v3_max + 3 * ell - k - 3 * j + 3

        min_n = num_leaves if num_leaves == 1 else num_leaves + 1
        max_n = min(num_leaves + v2_max + v3_max, max_vertices_search)
        max_n = min(max_n, num_leaves + min(10, v2_max + v3_max))

        if max_n < min_n:
            continue

        for n in range(min_n, max_n + 1):
            for tree in graphs.trees(n):
                leaves, internal, degree_2, degree_3plus = classify_vertices_sage(tree)
                if len(leaves) != num_leaves:
                    continue
                if len(degree_3plus) > v3_max:
                    continue
                if len(degree_2) > v2_max:
                    continue
                valid_trees.append((tree, j, num_leaves))

    return valid_trees
\end{lstlisting}

Reading it in order:

\begin{enumerate}
  \item $j_{\max}$ is set by \eqref{eq:jmax}, and we loop $j$ from $0$ to
        $j_{\max}$. For each $j$ we fix $\code{num\_leaves}=k+j$, then compute
        \code{v3\_max} \eqref{eq:v3max} and \code{v2\_max} \eqref{eq:v2max}.
  \item The vertex-count search window $[\code{min\_n}, \code{max\_n}]$ is set:
        \code{min\_n} is \code{num\_leaves} if there is a single leaf (the
        degenerate one-vertex tree), else $\code{num\_leaves}+1$ (you need at
        least one internal vertex). The upper end is $\code{num\_leaves}+
        \code{v2\_max}+\code{v3\_max}$ (the most vertices the bounds permit),
        capped by \code{max\_vertices\_search}, and then capped \emph{again} by
        the \code{+10} guard discussed below.
  \item If the window is empty (\code{max\_n < min\_n}), this $j$ is skipped.
  \item For each $n$ in the window, \code{graphs.trees(n)} yields every
        non-isomorphic tree on $n$ vertices exactly once (\S\ref{sec:sage}). A
        tree is \emph{kept} iff its leaf count equals \code{num\_leaves}, its
        $V_3$ count is $\le\code{v3\_max}$, and its $V_2$ count is
        $\le\code{v2\_max}$.
\end{enumerate}

\begin{gotcha}
\label{gotcha:plus10}
\textbf{The \code{+10} vertex-search cap can in principle undercut completeness at
high order.} The line
\code{max\_n = min(max\_n, num\_leaves + min(10, v2\_max + v3\_max))}
(\file{loop\_diagram\_enumeration.py:143}) clamps the \emph{internal}-vertex
budget to at most $10$. The proven bound only guarantees that all needed
topologies fit within $\code{num\_leaves}+\code{v2\_max}+\code{v3\_max}$
vertices. For small $(k,\ell)$ the proven budget is well under $10$, so the cap is
inert. But at large enough orders $\code{v2\_max}+\code{v3\_max} > 10$, and then
the cap could \emph{skip trees the proven bound says we need}, silently dropping
topologies. This is a performance guard, not a proven-safe bound --- flag it for
review at high order.
\end{gotcha}

\section{Stage 2: adding edges to make topologies}

A tree has $L=0$. To reach $\ell$ loops we add exactly $\ell$ edges, then keep
only the results that have exactly $\ell$ loops, satisfy the structural
constraints, and are not isomorphic duplicates of something already kept.

\subsection{Enumerating the edges to add}

The candidate edges are produced by \code{generate\_edge\_multisets}
(\file{loop\_diagram\_enumeration.py:166}):

\begin{lstlisting}
def generate_edge_multisets(all_vertices, ell):
    if len(all_vertices) < 2:
        return []
    possible_edges = list(combinations(all_vertices, 2))
    return list(combinations_with_replacement(possible_edges, ell))
\end{lstlisting}

The two \code{itertools} generators here matter (\S\ref{sec:itertools}):
\code{combinations(all\_vertices, 2)} lists every \emph{possible} edge --- every
unordered pair of distinct vertices. Then
\code{combinations\_with\_replacement(possible\_edges, ell)} lists every
\emph{multiset} of $\ell$ edges to add. ``With replacement'' is essential: it
permits choosing the \emph{same} edge twice, which creates a double-edge between
that pair of vertices --- a \emph{bubble}, the simplest one-loop topology. A
plain \code{combinations} would forbid bubbles and the enumeration would be
incomplete.

The chosen edges are spliced onto a copy of the tree by \code{add\_edges\_to\_tree}
(\file{loop\_diagram\_enumeration.py:173}), which builds the copy as a Sage
\code{Graph(tree, multiedges=True, loops=False)} --- multi-edges \emph{allowed}
(we just argued we need them), self-loops \emph{forbidden} (an edge from a vertex
to itself is never physical here).

\subsection{The structural constraints on a topology}

Beyond having the right loop number, a topology must satisfy local structural
rules so that it reduces to a unique canonical diagram. These are enforced by
\code{check\_topology\_constraints(G, k, ell=None)}
(\file{loop\_diagram\_enumeration.py:180}):

\begin{lstlisting}
def check_topology_constraints(G, k, ell=None):
    leaves, internal, degree_2, degree_3plus = classify_vertices_sage(G)
    if len(leaves) != k:
        return False
    if not G.is_connected():
        return False
    if ell is None or ell > 0:
        if has_adjacent_degree2_sage(G, degree_2):
            return False
        if not check_deg3_has_non_deg2_neighbor(G, degree_3plus, degree_2):
            return False
        if not check_leaf_neighbors_not_all_deg2(G, leaves, degree_2):
            return False
    return True
\end{lstlisting}

The unconditional checks are the obvious ones: the final topology must have
exactly $k$ external legs (\code{len(leaves) == k}), and it must be
\code{is\_connected} (a diagram contributing to a \emph{connected} cumulant must
be connected). The three conditional checks are all ``degree-2 pruning'' rules
that remove redundant subdivisions of a single propagator line:

\begin{description}
  \item[No two adjacent degree-2 vertices] (\code{has\_adjacent\_degree2\_sage},
        \file{loop\_diagram\_enumeration.py:41}). A chain of two degree-2 vertices
        in a row is a redundant subdivision of one propagator line; only one
        canonical representative is kept.
  \item[Every degree-$\ge3$ vertex has a non-degree-2 neighbor]
        (\code{check\_deg3\_has\_non\_deg2\_neighbor},
        \file{loop\_diagram\_enumeration.py:50}). This returns \code{False} if
        some branching vertex is surrounded \emph{only} by subdivision chains ---
        an ``isolated hub.''
  \item[Leaf neighbors are not all degree-2]
        (\code{check\_leaf\_neighbors\_not\_all\_deg2},
        \file{loop\_diagram\_enumeration.py:58}). An external leg should not
        attach \emph{only} through a chain of degree-2 subdivisions.
\end{description}

\begin{gotcha}
\label{gotcha:ell0-skip}
The three degree-2 pruning checks are \emph{skipped at tree level} ($\ell=0$). The
gate is \code{if ell is None or ell > 0:} (\file{loop\_diagram\_enumeration.py:189}),
and the comment explains why: at $\ell=0$ ``the tree IS the topology --- no
contraction ambiguity'' (\file{loop\_diagram\_enumeration.py:188}). The constraint
set is therefore \emph{different} at tree level vs.\ loop level. Passing the wrong
\code{ell} into this function would change which topologies survive --- a subtle
correctness trap.
\end{gotcha}

\subsection{Processing one tree, then deduplicating}

The unit of work is \code{process\_tree\_parallel(args)}
(\file{loop\_diagram\_enumeration.py:203}), where \code{args = (tree, j,
num\_leaves, k, ell)}. For each edge-multiset it adds the edges, demands
\code{count\_cycles\_sage(G) == ell}, demands \code{check\_topology\_constraints},
then relabels the survivor leaves-first and re-classifies:

\begin{lstlisting}
def process_tree_parallel(args):
    tree, j, num_leaves, k, ell = args
    local_candidates = []
    all_verts = list(tree.vertices())
    for edge_multiset in generate_edge_multisets(all_verts, ell):
        G = add_edges_to_tree(tree, edge_multiset)
        if count_cycles_sage(G) != ell:
            continue
        if not check_topology_constraints(G, k, ell=ell):
            continue
        G_relabeled, _ = relabel_leaves_first(G)
        leaves_final, internal_final, _, _ = classify_vertices_sage(G_relabeled)
        local_candidates.append((G_relabeled, leaves_final, internal_final))
    return local_candidates
\end{lstlisting}

The name says ``parallel'' because this is the function submitted to the thread
pool (\S\ref{sec:parallel}); it runs perfectly well serially too. The relabeling
step \code{relabel\_leaves\_first} (\file{loop\_diagram\_enumeration.py:74})
renumbers the graph so that \emph{leaves are $0\ldots|L|-1$ and internal vertices
are $|L|\ldots|V|-1$}; it calls \code{G.relabel(dict, inplace=False)}, which
(because of \code{inplace=False}) returns a \emph{new} relabeled graph rather than
mutating the original. This canonical leaf-first labeling is the invariant the
downstream code relies on (Gotcha~\ref{gotcha:leaffirst}).

The candidates from all trees are then collapsed by
\code{\_remove\_isomorphic\_undirected} (\file{loop\_diagram\_enumeration.py:219}):

\begin{lstlisting}
def _remove_isomorphic_undirected(candidates):
    unique = []
    for G, leaves, internal in candidates:
        if not any(graphs_isomorphic_with_labels(G, leaves, G2, l2)
                   for G2, l2, _ in unique):
            unique.append((G, leaves, internal))
    return unique
\end{lstlisting}

This is a classic $O(n^2)$ dedup: walk the candidates, keep one only if it is
\emph{not} isomorphic to any already-kept representative. The isomorphism
predicate \code{graphs\_isomorphic\_with\_labels} is the subject of
\S\ref{sec:iso}; each comparison is cheap because it fast-rejects on cheap
invariants before invoking the expensive test.

\section{Stage 3: orienting the edges into prediagrams}

An \msrjd{} edge is a \emph{directed} propagator, so each undirected topology must
be oriented. With $|E|$ edges there are $2^{|E|}$ raw orientations; only the
physically admissible ones survive.

\subsection{Building one orientation}

\code{orient\_edges(G, orientation\_bits)}
(\file{loop\_diagram\_enumeration.py:264}) turns a bit-vector into a
\code{DiGraph}:

\begin{lstlisting}
def orient_edges(G, orientation_bits):
    D = DiGraph(multiedges=True, loops=False)
    D.add_vertices(G.vertices())
    for i, (u, v) in enumerate(G.edges(labels=False)):
        if orientation_bits[i] == 0:
            D.add_edge(u, v, i)
        else:
            D.add_edge(v, u, i)
    return D
\end{lstlisting}

The $i$-th edge keeps its direction $u\to v$ when $\code{orientation\_bits}[i]=0$,
and flips to $v\to u$ otherwise. Crucially \textbf{the edge index $i$ is used as
the edge label} (\code{D.add\_edge(u, v, i)}). This preserves multi-edge identity:
when a topology has a double edge between the same pair, the two parallel arrows
remain distinguishable by their labels, and downstream code
(\code{type\_assignment} iterates \code{D.edges()} as 3-tuples \code{(u, v,
label)}) depends on those labels being present and distinct
(Gotcha~\ref{gotcha:edgelabels}). The per-edge order comes from
\code{G.edges(labels=False)}, which is deterministic for a fixed graph object.

\subsection{The MSR-JD causality constraints}

The admissible orientations are selected by \code{check\_orientation\_constraints}
(\file{loop\_diagram\_enumeration.py:275}):

\begin{lstlisting}
def check_orientation_constraints(D, leaves):
    leaves_set = set(leaves)
    if not D.is_directed_acyclic():
        return False
    for v in D.vertices():
        in_deg  = D.in_degree(v)
        out_deg = D.out_degree(v)
        if (in_deg + out_deg) == 1 and in_deg != 1:
            return False
        if out_deg == 0 and in_deg >= 2:
            return False
        if in_deg == 1 and out_deg == 1:
            return False
    sources = [v for v in D.vertices() if D.in_degree(v) == 0]
    for v in sources:
        if any(u in sources for u in D.neighbors_out(v)):
            return False
    return True
\end{lstlisting}

Each clause has a direct physical reading:

\begin{description}
  \item[Acyclic.] \code{D.is\_directed\_acyclic()} --- \msrjd{} response
        propagators are causal, so a valid diagram has no directed cycle: you
        cannot return to your own past.
  \item[A leaf must be incoming.] If $\code{in\_deg}+\code{out\_deg}=1$ (i.e.\ the
        vertex is a degree-1 leaf) we require $\code{in\_deg}=1$: the external leg
        carries the field \emph{in}; the arrow points \emph{into} the leaf.
  \item[No multi-line pure sink.] $\code{out\_deg}=0$ with $\code{in\_deg}\ge 2$ is
        forbidden --- a pure sink absorbing two or more lines is not an admissible
        \msrjd{} vertex.
  \item[No pass-through degree-2 vertex.] $\code{in\_deg}=1$ and
        $\code{out\_deg}=1$ is forbidden --- a one-in/one-out degree-2 vertex is
        again a trivial line-subdivision.
  \item[No two adjacent sources.] A \emph{source} is a vertex with
        $\code{in\_deg}=0$ (and not a leaf). Two adjacent sources are forbidden:
        sources represent noise-injection points, and two of them cannot be wired
        directly together.
\end{description}

This is the \term{source / interaction / leaf} trichotomy that drives all the
downstream physics: \textbf{sources} (in-degree 0, not a leaf) are
\emph{noise-kernel} insertion points; \textbf{interaction vertices} (both in- and
out-edges) are \emph{action-vertex} insertion points; \textbf{leaves} are
\emph{external legs}. The next chapter and \code{degree\_scan.py} re-derive
exactly this trichotomy.

\subsection{Enumerating and deduplicating orientations}

The brute-force orientation sweep is \code{enumerate\_orientations(G, leaves)}
(\file{loop\_diagram\_enumeration.py:295}):

\begin{lstlisting}
def enumerate_orientations(G, leaves):
    edges = list(G.edges(labels=False))
    return [orient_edges(G, [(bits >> i) & 1 for i in range(len(edges))])
            for bits in range(2 ** len(edges))
            if check_orientation_constraints(
                orient_edges(G, [(bits >> i) & 1 for i in range(len(edges))]),
                leaves)]
\end{lstlisting}

It iterates \code{bits} over $0 \ldots 2^{|E|}-1$, decoding each integer into a
length-$|E|$ orientation vector via the bit-extraction \code{(bits >> i) \& 1}
(shift right by $i$, mask the low bit), builds the \code{DiGraph}, and keeps it
iff the constraints pass.

\begin{gotcha}
\label{gotcha:double-build}
\textbf{\code{enumerate\_orientations} builds each candidate \code{DiGraph}
twice} (\file{loop\_diagram\_enumeration.py:297--301}) --- once \emph{inside} the
\code{check\_orientation\_constraints(...)} call, and once again to materialize
the kept result in the list comprehension. This is functionally correct but
doubles the orientation-construction cost; a clean rewrite would build once and
reuse.
\end{gotcha}

The oriented candidates are deduplicated by \code{remove\_isomorphic\_directed}
(\file{loop\_diagram\_enumeration.py:304}) --- the same $O(n^2)$ pattern as the
undirected case, but using the \emph{directed} isomorphism predicate
\code{directed\_graphs\_isomorphic\_with\_labels}, in which arrows must match too.
Each result is carried as the final 4-tuple \code{(D, G, leaves, internal)}.

\section{Graph isomorphism and nauty, from scratch}
\label{sec:iso}

The dedup steps in both stages rest on \emph{graph isomorphism}, which deserves a
proper introduction since it is where almost all the compute time goes at high
order.

\subsection{Why dedup is needed, and what isomorphism means}

After adding edges, and again after orienting, the same \emph{abstract} graph is
produced many times with different vertex \emph{labels} --- ``vertex $5$ here is
vertex $7$ there.'' Two graphs are \term{isomorphic} if one can be turned into
the other purely by \emph{relabeling vertices}: a bijection of vertices that maps
edges to edges. For Feynman-diagram bookkeeping, isomorphic graphs are the
\emph{same diagram} and must be counted exactly once.

Deciding isomorphism naively means trying all $n!$ relabelings --- hopeless even
for modest $n$. The classic engine that does it efficiently is \term{nauty}
(``No AUTomorphisms, Yes?''), by Brendan~McKay. Nauty computes a \emph{canonical
form}: a deterministic relabeling such that two graphs are isomorphic \emph{iff}
their canonical forms are \emph{literally identical}. It does this by iteratively
\emph{refining a coloring} of the vertices --- partitioning them by degree, then
by their neighbors' degrees, and so on --- and backtracking through the residual
symmetries, pruning the search with the automorphisms it discovers along the way.

\begin{note}
You will not see the string \code{nauty} anywhere in this code --- it lives
\emph{inside} SageMath. Every \code{is\_isomorphic} call routes to nauty's
canonical-form machinery. That is the \emph{only} graph-isomorphism technology
this subsystem relies on, and it is what makes deduplication tractable at all.
\end{note}

\subsection{Colored isomorphism: leaves are not interchangeable}

A plain isomorphism may map \emph{any} vertex to any other. But in our diagrams
that is too permissive: \textbf{a leaf must not be mapped onto an internal
vertex}. An external leg and an interaction point play different physical roles,
so a relabeling that swaps the two is \emph{not} a legal symmetry.

The fix is \term{vertex coloring}. Assign color $0$ to leaves and color $1$ to
internal vertices, and ask for a \emph{color-preserving} isomorphism: nauty/Sage
honors the coloring as an initial vertex partition and only searches over
relabelings that keep colors fixed. The undirected predicate
\code{graphs\_isomorphic\_with\_labels} (\file{loop\_diagram\_enumeration.py:84})
implements exactly this:

\begin{lstlisting}
def graphs_isomorphic_with_labels(G1, leaves1, G2, leaves2):
    if G1.order() != G2.order() or G1.size() != G2.size():
        return False
    if len(leaves1) != len(leaves2):
        return False
    leaves1_set = set(leaves1)
    leaves2_set = set(leaves2)
    G1c = G1.copy(); G2c = G2.copy()
    for v in G1.vertices():
        G1c.set_vertex(v, 0 if v in leaves1_set else 1)
    for v in G2.vertices():
        G2c.set_vertex(v, 0 if v in leaves2_set else 1)
    return G1c.is_isomorphic(G2c)
\end{lstlisting}

Three things to read here. First, the \emph{fast rejects}: different vertex count
(\code{order}), different edge count (\code{size}), or different leaf count
returns \code{False} immediately, without invoking nauty. These cheap invariants
are what keep the $O(n^2)$ dedup affordable at small orders. Second, the coloring:
\code{set\_vertex(v, 0\_or\_1)} attaches a payload to each vertex that Sage
interprets as its color. Third, the test itself: \code{G1c.is\_isomorphic(G2c)} is
now a color-preserving comparison, so leaves can only map to leaves.

The directed analog \code{directed\_graphs\_isomorphic\_with\_labels}
(\file{loop\_diagram\_enumeration.py:99}) is structurally identical; the only
difference is that the final \code{is\_isomorphic} on \code{DiGraph}s also
requires the arrows to match.

\section{External tools, from first principles}

This subsystem touches SageMath heavily, three Python standard-library modules,
and --- only transitively, through \code{degree\_scan.py} --- the project's own
serialization and \code{FieldTheory}. It does \emph{not} itself import
\code{sympy}, \code{scipy}, \code{numba}, or \code{networkx}. Each tool is
explained from scratch below; the imports at the top of
\file{loop\_diagram\_enumeration.py:22--25} are:

\begin{lstlisting}
from sage.all import *
from itertools import combinations, combinations_with_replacement
from collections import Counter, defaultdict
from concurrent.futures import ThreadPoolExecutor, as_completed
\end{lstlisting}

\subsection{SageMath}
\label{sec:sage}

\term{SageMath} is a large open-source mathematics system built on top of Python.
It bundles dozens of specialized math libraries (Pari, GAP, Singular, nauty,
NetworkX, NumPy, $\ldots$) behind a uniform Python API, and adds its own types
for exact arithmetic, symbolic expressions (the ``Symbolic Ring'' \code{SR}),
and --- most relevant here --- \emph{graph theory}: the \code{Graph}, \code{DiGraph},
and \code{graphs.*} generator classes.

Because Sage replaces some Python built-ins (e.g.\ integer literals become Sage
\code{Integer}s) and needs its bundled libraries on the path, \emph{this module
can only run inside a SageMath kernel/interpreter}. The module docstring says so
plainly --- \code{"Requires a SageMath kernel."}
(\file{loop\_diagram\_enumeration.py:7}) --- and the tests are run via \code{sage
-python -m pytest}. The whole API is pulled in with the conventional wildcard
import \code{from sage.all import *} (\file{loop\_diagram\_enumeration.py:22}),
which brings \code{Graph}, \code{DiGraph}, \code{graphs}, \code{Integer}, and the
rest into scope.

The specific Sage features this code uses:

\begin{itemize}
  \item \code{graphs.trees(n)} (\file{loop\_diagram\_enumeration.py:149}) --- a
        \emph{generator} yielding every non-isomorphic tree on $n$ vertices,
        exactly once each. This is the brute-force engine of the tree stage; Sage
        generates these efficiently so the user never enumerates trees by hand.
  \item \code{Graph(tree, multiedges=True, loops=False)}
        (\file{loop\_diagram\_enumeration.py:174}) --- an undirected graph that
        \emph{allows} multiple edges between the same pair (\code{multiedges=True},
        needed because a loop-edge can double an existing edge to make a bubble)
        but \emph{forbids} self-loops (\code{loops=False}).
  \item \code{DiGraph(multiedges=True, loops=False)}
        (\file{loop\_diagram\_enumeration.py:265}) --- the directed analog, used to
        hold an oriented topology (a prediagram).
  \item \textbf{Inspection methods:} \code{G.vertices()}, \code{G.edges(labels=...)},
        \code{G.degree(v)}, \code{G.neighbors(v)}, \code{G.is\_connected()},
        \code{G.order()} ($=|V|$), \code{G.size()} ($=|E|$), \code{G.add\_edge(u, v)},
        \code{G.copy()}, \code{G.relabel(dict, inplace=False)}.
  \item \code{G.set\_vertex(v, value)} (\file{loop\_diagram\_enumeration.py:93})
        --- attaches an arbitrary payload (here the color $0$/$1$) to a vertex;
        Sage's isomorphism test treats these payloads as a vertex coloring.
  \item \code{G.is\_isomorphic(H)} (\file{loop\_diagram\_enumeration.py:96}) ---
        the \textbf{nauty-backed} isomorphism test; with vertex colors set first
        it becomes the color-preserving test of \S\ref{sec:iso}.
  \item \textbf{DiGraph-specific:} \code{D.add\_vertices(...)}, \code{D.add\_edge(u,
        v, label)}, \code{D.in\_degree(v)}, \code{D.out\_degree(v)},
        \code{D.is\_directed\_acyclic()}, \code{D.neighbors\_out(v)}.
  \item \textbf{Plotting:} \code{graphics\_array}, \code{G.plot(...)},
        \code{.show()}, \code{.save()} --- used only by the \code{\_plot\_*}
        helpers and the \code{show\_*} API for inline notebook display, not by the
        data pipeline.
\end{itemize}

\subsection{\texttt{itertools}}
\label{sec:itertools}

A Python standard-library module of memory-efficient combinatorial iterators. Two
functions are used:

\begin{itemize}
  \item \code{combinations(iterable, r)} yields every $r$-subset (no repeats,
        order ignored). In \code{generate\_edge\_multisets}
        (\file{loop\_diagram\_enumeration.py:169}),
        \code{combinations(all\_vertices, 2)} enumerates every possible edge ---
        every unordered pair of distinct vertices.
  \item \code{combinations\_with\_replacement(iterable, r)} is like
        \code{combinations} but \emph{allows the same element more than once}. At
        \file{loop\_diagram\_enumeration.py:170},
        \code{combinations\_with\_replacement(possible\_edges, ell)} enumerates
        every multiset of $\ell$ edges to add. ``With replacement'' is what permits
        adding the \emph{same} edge twice --- the double-edge that makes a bubble.
\end{itemize}

\subsection{\texttt{collections}}

\begin{itemize}
  \item \code{defaultdict(list)} --- a dict that auto-creates an empty list on
        first access to a missing key. Used in \code{\_plot\_trees}
        (\file{loop\_diagram\_enumeration.py:349}) to bucket trees by their $j$
        value for grouped display.
  \item \code{Counter} is imported (\file{loop\_diagram\_enumeration.py:24}) but
        \textbf{not actually referenced} anywhere in the file --- a harmless dead
        import (see Gotcha~\ref{gotcha:deadimport}).
\end{itemize}

\subsection{\texttt{concurrent.futures}}
\label{sec:parallel}

\begin{itemize}
  \item \code{ThreadPoolExecutor(max\_workers=n)} runs callables on a pool of OS
        \emph{threads} (not processes). Used in \code{\_enumerate\_topologies\_raw}
        (\file{loop\_diagram\_enumeration.py:248}) and
        \code{\_enumerate\_prediagrams\_raw}
        (\file{loop\_diagram\_enumeration.py:332}) to fan the per-tree /
        per-topology work across threads when \code{n\_threads > 1}.
  \item \code{as\_completed(futures)} yields each submitted future \emph{as it
        finishes} (not in submission order), so results are collected as soon as
        they are ready.
\end{itemize}

\begin{gotcha}
\label{gotcha:threads}
\textbf{Thread-only parallelism is intentional and load-bearing for safety.} The
choice of a thread pool over a process/fork pool is deliberate. Per the project's
macOS fork-safety history, \emph{fork-based multiprocessing crashes the user's
machine} when a kernel forks after BLAS/Cocoa/matplotlib initialization. Threads
avoid \code{fork} entirely. The trade-off is the Python GIL --- but Sage's graph
routines spend most of their time in compiled C (nauty, the graph backend), which
\emph{releases} the GIL, so threads still parallelize the heavy isomorphism work.
The default is \code{n\_threads=1} (serial) everywhere; the pool only activates
when the caller opts in. Do \textbf{not} ``optimize'' this to a
\code{ProcessPoolExecutor} or fork --- it can crash the machine.
\end{gotcha}

\section{The public API}
\label{sec:api}

The module exposes a small, regular surface. The \code{enumerate\_*} functions
return data; the \code{show\_*} variants additionally render the graphs inline in
a notebook and return just the count(s). The colors quoted below are documented in
the corresponding \code{show\_*} docstrings.

\begin{itemize}
  \item \code{enumerate\_all(k, ell, n\_threads=1, max\_vertices\_search=50,
        verbose=True)} (\file{loop\_diagram\_enumeration.py:434}) runs the full
        pipeline \emph{once}, sharing the tree list across stages: it passes
        \code{trees\_with\_j=trees} into \code{\_enumerate\_topologies\_raw} so the
        trees are not regenerated. Returns \code{(trees, topologies, prediagrams,
        counts)} where \code{counts = \{n\_trees, n\_topologies, n\_prediagrams\}}.
  \item \code{enumerate\_trees(k, ell, ...)}
        (\file{loop\_diagram\_enumeration.py:478}) --- a thin wrapper over
        \code{generate\_trees\_with\_constraints}; returns \code{(trees, count)}.
  \item \code{enumerate\_topologies(k, ell, ...)}
        (\file{loop\_diagram\_enumeration.py:498}) --- returns \code{(topologies,
        count)}.
  \item \code{enumerate\_prediagrams(k, ell, ...)}
        (\file{loop\_diagram\_enumeration.py:515}) --- runs topologies then
        orientations; returns \code{(prediagrams, count)}. \textbf{This is the
        canonical entry point the rest of the pipeline calls.}
  \item \code{show\_all} / \code{show\_trees} / \code{show\_topologies} /
        \code{show\_prediagrams} (\file{loop\_diagram\_enumeration.py:462},
        \file{:535}, \file{:552}, \file{:569}) --- display each non-empty stage and
        return the count(s). The prediagram plot colors \emph{leaves} black,
        \emph{sources} red, and other internal vertices light blue, recomputing
        sources on the fly as the vertices with \code{in\_degree == 0}.
\end{itemize}

\begin{note}
The shared-tree optimization in \code{enumerate\_all} matters because tree
generation walks \code{graphs.trees(n)} for every $n$ in the search window, which
is not free. Running the three stages through \code{enumerate\_all} (rather than
calling \code{enumerate\_trees}, \code{enumerate\_topologies}, and
\code{enumerate\_prediagrams} separately) generates the trees exactly once and
threads them through.
\end{note}

\section{The degree-scan feedback bridge (build phase~C)}
\label{sec:degree-scan}

The second file, \file{msrjd/enumeration/degree\_scan.py}, is small but
important: it closes the loop between the \emph{topological} world (this chapter)
and the \emph{algebraic} world (the Taylor-expanded action of Phase~1). The
generator does not look at the action --- but the prediagrams it produces may
\emph{demand} interaction monomials of a higher degree than the action was
expanded to. Phase~C detects that and triggers a re-expansion.

Unlike the generator, this file imports no third-party libraries --- only project
internals (\file{degree\_scan.py:13--14}):

\begin{lstlisting}
from msrjd.core.serialize import load_theory, save_theory, reload_model
from msrjd.core.field_theory import FieldTheory
\end{lstlisting}

\subsection{Reading the degree a theory must support}

\code{max\_vertex\_degree(prediagrams)} (\file{degree\_scan.py:17}) returns the
maximum total degree $\code{in\_degree}(v)+\code{out\_degree}(v)$ over all
\emph{non-leaf} vertices across all prediagrams (or $0$ if none):

\begin{lstlisting}
def max_vertex_degree(prediagrams):
    max_deg = 0
    for (D, G, leaves, internal) in prediagrams:
        leaf_set = set(leaves)
        for v in D.vertices():
            if v in leaf_set:
                continue
            deg = D.in_degree(v) + D.out_degree(v)
            if deg > max_deg:
                max_deg = deg
    return max_deg
\end{lstlisting}

This is the \emph{required Taylor order}: a degree-$d$ internal vertex needs a
$d$-leg interaction monomial, which exists only if the action was expanded to
order $d$. Its sibling \code{scan\_source\_vertices(prediagrams)}
(\file{degree\_scan.py:45}) returns the \emph{set of out-degrees} exhibited by
\emph{source} vertices (non-leaf, \code{in\_degree == 0}) --- these out-degrees
tell the noise-kernel machinery which noise arities are needed.

\subsection{Comparing to the saved order, and re-expanding}

\code{check\_taylor\_order(meta, max\_degree)} (\file{degree\_scan.py:70}) is a
one-liner comparison returning \code{(sufficient, current\_order,
required\_order)} with $\code{sufficient} = (\code{meta['taylor\_order']} \ge
\code{max\_degree})$. The driver is \code{ensure\_taylor\_order(theory\_path,
prediagrams, project\_root=None)} (\file{degree\_scan.py:91}):

\begin{lstlisting}
def ensure_taylor_order(theory_path, prediagrams, project_root=None):
    meta, data = load_theory(theory_path)
    max_deg = max_vertex_degree(prediagrams)
    sufficient, current, required = check_taylor_order(meta, max_deg)
    if sufficient:
        return meta, data
    # Re-expand at the required order
    print(f'Re-expanding theory from order {current} to {required}...')
    model = reload_model(meta, project_root=project_root)
    ft = FieldTheory(model, taylor_order=required)
    ft.expand()
    save_theory(
        theory_path, ft,
        propagator_data=None,  # propagator must be recomputed separately
        stationarity=meta.get('stationarity', True),
        model_file=meta.get('model_file'),
        model_var_name=meta.get('model_var_name'),
    )
    print(f'Theory re-expanded and saved. Note: propagator data must be recomputed.')
    return load_theory(theory_path)
\end{lstlisting}

In the common case (\code{sufficient == True}) the function returns the
already-loaded theory immediately and never touches \code{FieldTheory}. Only on
the re-expansion branch does it reload the model definition
(\code{reload\_model}), rebuild a \code{FieldTheory(model,
taylor\_order=required)}, call \code{.expand()}, re-save, and re-read from disk.
The supporting serialization calls are: \code{load\_theory(path)} returning
\code{(meta, data)} and \code{save\_theory(path, ft, ...)} (read/write a
previously-expanded theory, where \code{meta} carries \code{meta['taylor\_order']}
and, on the re-expansion path, \code{meta['stationarity']},
\code{meta['model\_file']}, \code{meta['model\_var\_name']}); and
\code{reload\_model(meta, project\_root=...)} (re-import the model definition from
\code{meta['model\_file']} / \code{meta['model\_var\_name']}).

\begin{gotcha}
\label{gotcha:stale-prop}
\textbf{\code{ensure\_taylor\_order} leaves the propagator stale after
re-expansion.} It passes \code{propagator\_data=None} to \code{save\_theory} and
prints ``propagator data must be recomputed''
(\file{degree\_scan.py:128--133}). A caller that re-expands and then immediately
uses the propagator without recomputing it will get a wrong or missing
propagator. This is by design --- propagator computation is a separate phase ---
but it is an easy footgun.
\end{gotcha}

\section{Data structures and the downstream contract}

The subsystem passes around a small number of tuple shapes; no dataclasses are
defined here (the \code{TypedDiagram} dataclass lives downstream in
\code{type\_assignment.py}).

\begin{description}
  \item[Tree record] \code{(tree, j, num\_leaves)} --- \code{tree} a Sage
        \code{Graph} (acyclic); \code{j} the leaf surplus; \code{num\_leaves =
        k + j} the leaf count.
  \item[Topology record] \code{(G, leaves, internal)} --- \code{G} a Sage
        \code{Graph} (undirected, multi-edges allowed, exactly $\ell$ loops),
        relabeled leaves-first; \code{leaves} the vertex IDs $0\ldots k-1$;
        \code{internal} the remaining IDs.
  \item[Prediagram record] \code{(D, G, leaves, internal)} --- \emph{the} output
        object. \code{D} a Sage \code{DiGraph} (the oriented graph, edge labels
        carrying the original edge index for multi-edge identity); \code{G} the
        underlying undirected topology; \code{leaves} / \code{internal} as above.
        This 4-tuple is exactly what
        \code{type\_assignment.enumerate\_typed\_diagrams} unpacks and what
        \code{degree\_scan.*} iterate over.
  \item[Counts dict] \code{\{'n\_trees', 'n\_topologies', 'n\_prediagrams'\}},
        built in \code{enumerate\_all} (\file{loop\_diagram\_enumeration.py:454}).
  \item[Edge-multiset] a tuple of $\ell$ vertex-pairs, e.g.\ \code{((0,3),(0,3))}
        for a double-edge bubble between vertices $0$ and $3$.
  \item[Orientation vector] a length-$|E|$ list of bits ($0$ = keep $(u,v)$,
        $1$ = flip to $(v,u)$), produced by \code{(bits >> i) \& 1}.
\end{description}

The end-to-end flow of a single \code{enumerate\_prediagrams(k, ell)} call is then:

\begin{lstlisting}[style=console]
  (k, ell)
     |  generate_trees_with_constraints
  [ (tree, j, num_leaves), ... ]                      trees passing degree bounds
     |  process_tree_parallel  (add ell edges; demand L==ell; constraints;
     |                          relabel leaves-first)
  [ (G, leaves, internal), ... ]  (with isomorphic duplicates)
     |  _remove_isomorphic_undirected   (nauty, leaf-colored)
  [ (G, leaves, internal), ... ]  unique TOPOLOGIES
     |  enumerate_orientations  (2^|E| orientations, keep admissible)
  [ (D, G, leaves, internal), ... ]  (with isomorphic duplicates)
     |  remove_isomorphic_directed   (nauty, leaf-colored, directed)
  [ (D, G, leaves, internal), ... ]  unique PREDIAGRAMS  --->  downstream
\end{lstlisting}

\section{A worked walk-through: \texorpdfstring{$(k,\ell)=(2,1)$}{(k,l)=(2,1)}}

It helps to trace the smallest non-trivial case --- the one-loop correction to a
two-point function. We have $k=2$, $\ell=1$, so by \eqref{eq:jmax} the surplus
runs $j\in\{0,1\}$.

\begin{itemize}
  \item For $j=0$: $\code{num\_leaves}=2$, $\code{v3\_max}=k+j-2=0$ (no branching
        vertices allowed), $\code{v2\_max}=k+3\ell-j-1=4$. The trees here are
        simple paths with two leaves and a short chain of degree-2 vertices
        between them.
  \item For $j=1$: $\code{num\_leaves}=3$, $\code{v3\_max}=1$ (one branching
        vertex allowed --- a ``Y'' joint), $\code{v2\_max}=3$. The third leaf is
        the one the single loop-edge will consume.
\end{itemize}

Each surviving tree then has $\ell=1$ extra edge added in every
\code{combinations\_with\_replacement} way (including the double-edge that makes a
bubble), and only the results with loop number exactly $1$ that pass the
structural constraints are kept. After leaf-colored deduplication this yields
\textbf{9 topologies} for $(k,\ell)=(2,1)$. Orienting those 9 topologies under the
causality constraints of \S\ref{sec:iso} and deduplicating again gives the final
prediagram list that flows downstream to type assignment. (The corresponding
deduplicated counts at the next orders are $(3,1)\to 67$ and $(2,2)\to 289$.)

\section{Performance notes and known limitations}

We collect the sharp edges that were flagged inline, plus the remaining
performance caveats, so they live in one place.

\begin{gotcha}
\label{gotcha:deadimport}
\code{Counter} is imported but unused (\file{loop\_diagram\_enumeration.py:24}) ---
only \code{defaultdict} is referenced. Harmless dead import.
\end{gotcha}

\begin{gotcha}
\label{gotcha:on2}
\textbf{$O(n^2)$ isomorphism dedup.} Both \code{\_remove\_isomorphic\_undirected}
and \code{remove\_isomorphic\_directed} compare each new candidate against
\emph{every} kept representative. The cheap invariant fast-rejects (order, size,
leaf-count) keep this affordable at small orders, but the comparison count is
quadratic in the candidate count and will dominate at high order. A canonical-form
bucket (Sage's \code{Graph.canonical\_label}, hashing the canonical form and
bucketing) would make it near-linear; that is not done here.
\end{gotcha}

\begin{gotcha}
\label{gotcha:edgelabels}
\textbf{Edge labels carry the original edge index.} \code{orient\_edges} adds each
arrow as \code{D.add\_edge(u, v, i)} with \code{i} the edge index
(\file{loop\_diagram\_enumeration.py:269--271}). This is what preserves multi-edge
identity through orientation and into the prediagram \code{D}. Downstream
multi-edge handling (\code{type\_assignment} iterating \code{D.edges()} as 3-tuples
\code{(u, v, label)}) depends on these labels being present and distinct.
\end{gotcha}

\begin{gotcha}
\label{gotcha:leaffirst}
\textbf{The leaf-first labeling is assumed downstream.} \code{relabel\_leaves\_first}
is applied to every accepted topology inside \code{process\_tree\_parallel}
(\file{loop\_diagram\_enumeration.py:213}), and downstream code
(\code{type\_assignment}, \code{degree\_scan}) treats \code{leaves} as the
canonical external-leg set, relying on the relabeling having been done. Any code
path that builds prediagrams \emph{without} going through
\code{process\_tree\_parallel} (e.g.\ hand-built mocks in tests) will not have
leaf-first labels.
\end{gotcha}

\begin{gotcha}
\label{gotcha:sentinel}
\code{count\_cycles\_sage} returns $-1$ for a disconnected graph
(\file{loop\_diagram\_enumeration.py:69--71}), which is not $\ell$ for any
$\ell\ge 0$; the disconnected graph is therefore rejected automatically by the
\code{== ell} test in \code{process\_tree\_parallel}. The $-1$ is a deliberate
sentinel, not an error.
\end{gotcha}

The remaining structural limitations were stated in context and are recalled here
for completeness: the \code{+10} vertex-search cap (Gotcha~\ref{gotcha:plus10})
can in principle undercut the proven completeness bound at high order; the
double-build in \code{enumerate\_orientations} (Gotcha~\ref{gotcha:double-build})
doubles the orientation-construction cost; the $\ell=0$ constraint skip
(Gotcha~\ref{gotcha:ell0-skip}) means the topology constraint set differs at tree
level; and the whole module requires a SageMath kernel (\S\ref{sec:sage}). None of
these affect correctness at the small orders the pipeline routinely runs, but all
are worth keeping in view when pushing $\ell$ to three loops and beyond.

\section{Where to go next}

The prediagram list this chapter produces is consumed in
Chapter~\ref{ch:type-assignment}, where each bare \code{(D, G, leaves, internal)}
tuple is walked and every consistent assignment of propagators to edges and
interaction monomials to vertices is tried, the invalid ones discarded, and a
fully-typed \code{TypedDiagram} emitted for each survivor. That is the step where
the physics --- the response-versus-correlation propagator distinction, the
specific vertices the action provides, and the symmetry factor $\Sym(\Gamma)$ ---
finally re-enters. The completeness guaranteed here (every shape) and the validity
guaranteed there (only labelings the action supports) together ensure the diagram
sum the engine evaluates is exactly the right one.

% ---------------------------------------------------------------------
%  Chapter 10 : From Bare Graphs to Labelled Diagrams
%               (Type assignment, causality, symmetry factor, filter)
% ---------------------------------------------------------------------
\chapter{From Bare Graphs to Labelled Diagrams}\label{ch:type-assignment}

The previous chapter handed us a pile of \term{prediagrams}: bare directed
multigraphs with $k$ external leaves and $\ell$ loops, deduplicated up to graph
isomorphism, but with \emph{no physics on them yet}. A prediagram knows that
``two arrows leave this vertex and one enters it''; it does not yet know which
monomial of the action sits there, which field rides each leg, or which
propagator each edge carries. This chapter is the bridge from those bare
topologies to fully labelled Feynman diagrams that an integrator can actually
evaluate.

The bridge is four pipeline phases, run in order, and they are the subject of
this chapter:

\begin{description}
  \item[Phase D --- Filter] (\file{msrjd/diagrams/filter.py}). A cheap
        set-membership pre-screen that throws away any prediagram whose vertex
        \emph{degrees} no vertex or source of the theory could ever supply ---
        before we pay for the expensive labelling.
  \item[Phase E --- Type assignment] (\file{msrjd/diagrams/type\_assignment.py}).
        The heart of the chapter. Given one bare graph, produce \emph{every}
        valid way to decorate it: a monomial at each internal vertex, an external
        field on each leaf, and a propagator $G_{ij}$ on each edge. This is a
        constraint-satisfaction search; its output is a stream of
        \code{TypedDiagram} objects.
  \item[Phase F --- Causality] (\file{msrjd/diagrams/causality.py}). Discard typed
        diagrams that violate the retarded boundary condition of \msrjd{}.
        Structurally, that means the diagram must be a directed acyclic graph.
  \item[Phase G --- Symmetry factor \& dedup]
        (\file{msrjd/diagrams/symmetry.py}). Deduplicate the typed diagrams (Phase
        E emits the same diagram many times, relabelled) and compute the
        combinatorial \emph{symmetry factor} $\Sym(\Gamma)$ that every Feynman
        rule carries.
\end{description}

A fifth file, \file{msrjd/fork\_safety.py}, is a small shared safety guard that
the parallel branch of Phase E consults; we treat it at the end of the chapter
because it is shared with the temporal Phase~J path and is one of the most
load-bearing six lines in the whole repository.

\begin{note}
A quick orientation on \emph{why these phases are separate}. Filtering is
$O(\text{vertices})$ and uses no algebra; type assignment is a backtracking
search; deduplication and the symmetry factor both need a graph-isomorphism
engine (\code{nauty}, via Sage). Keeping them apart lets the cheap check run
first, lets the search be tested in isolation, and confines every call into the
isomorphism machinery to one file. The single function that wires all four
together with disk caching is
\code{enumerate\_unique\_diagrams} in \file{pipeline/\_diagrams.py}; we read it in
\S\ref{sec:ta-orchestrator}.
\end{note}

% =====================================================================
\section{The two field species and the edge convention}\label{sec:ta-convention}

Everything in this chapter rests on one convention, so we state it before any
code. A reader fluent in \msrjd{} field theory will recognise all of it; we
restate it from the ground up so the code's index choices are unambiguous.

\subsection{Physical fields and response fields}

A Langevin equation $\Dt\varphi = F[\varphi] + \xi$ with noise $\xi$ is rewritten,
\`a la Martin--Siggia--Rose / Janssen--De~Dominicis (\msrjd{}), as a path integral
over \emph{two} fields per physical field. The generating functional is
\begin{equation}
  Z \;=\; \int \mathcal{D}[\varphi]\,\mathcal{D}[\tilde\varphi]\;
  \ee^{-S[\varphi,\,\tilde\varphi]} .
  \label{eq:ta-Z}
\end{equation}
The two species are:

\begin{itemize}
  \item the \term{physical field} $\varphi$ --- the actual dynamical variable. In
        the code's naming its generator base names are letters like \code{vt},
        \code{nt}, \code{xt} (the trailing ``t'' is part of the base name, not a
        tilde);
  \item the \term{response field} $\tilde\varphi$ (``tilde'') --- the auxiliary
        conjugate field of \msrjd{}. In the code its base names carry a \code{d}
        prefix: \code{dv}, \code{dn}, \code{dx}.
\end{itemize}

The quadratic (Gaussian) part of $S$ defines the \emph{propagators}; the
higher-order part defines the \emph{interaction vertices}; and the noise enters
as \emph{source vertices} --- pure response-field monomials drawn from the noise
cumulant-generating functional.

\subsection{What a directed edge means}

Every internal line of a diagram is a propagator, and \msrjd{} has essentially
one forward-in-time propagator family: the retarded response $\avg{\varphi\,
\tilde\varphi}$. The code encodes this with a strict directionality convention,
stated at the top of \file{type\_assignment.py:11--13}:

\begin{quote}\itshape
each directed edge $u\to v$ carries a propagator $G_{ij}$ where $i$ is the
response-field index contributed by vertex $u$ (the tail) and $j$ is the
physical-field index contributed by vertex $v$ (the head).
\end{quote}

So a directed edge is a \emph{Wick contraction} between a response leg
$\tilde\varphi_i$ at its tail and a physical leg $\varphi_j$ at its head. The
value carried is an entry of the propagator matrix
$G_{\mathrm{ft}} = K_{\mathrm{ft}}^{-1}$, the inverse of the quadratic-form matrix
$K$. Now the part that is genuinely easy to get backwards, and which the code
warns about in two places (\file{type\_assignment.py:293--296} and
\file{:490--502}):

\begin{defn}
\textbf{The $(pi, ri)$ index order.}
$G_{\mathrm{ft}} = K_{\mathrm{ft}}^{-1}$ has \emph{rows indexed by physical fields
and columns by response fields}, so
\[
  G_{\mathrm{ft}}[i,\,j] \;=\; \avg{\varphi_i\,\tilde\varphi_j}.
\]
The propagator on an edge whose \emph{tail} contributes response index $ri$ and
whose \emph{head} contributes physical index $pi$ is therefore
\[
  \avg{\varphi_{\mathrm{phys}}\,\tilde\varphi_{\mathrm{resp}}}
  \;=\; G_{\mathrm{ft}}[\,pi,\; ri\,].
\]
You index $G_{\mathrm{ft}}$ with \textbf{(physical-row, response-col)}, in that
order.
\end{defn}

This ordering is load-bearing. The code stores the pair the ``natural'' way
round --- \code{propagator\_indices[edge] = (ri, pi)} --- but every \emph{lookup}
into the matrix or its zero-mask uses \code{(pi, ri)}. We will see exactly that
asymmetry in the backtracker, and \S\ref{sec:ta-gotchas} records the historical
bug that motivated the comment.

\begin{figure}[ht]
\centering
\begin{tikzpicture}[
    >={Stealth[length=2.6mm]},
    vert/.style={circle,draw,fill=notebg,minimum size=7mm,inner sep=0pt},
    every node/.style={font=\small},
  ]
  \node[vert] (u) at (0,0) {$u$};
  \node[vert] (v) at (4.4,0) {$v$};
  \draw[->,thick] (u) -- (v)
    node[midway,above=3pt] {$G_{\mathrm{ft}}[\,pi,\,ri\,]
        = \avg{\varphi_j\,\tilde\varphi_i}$};
  \node[below=2pt,align=center] at (u.south)
    {tail: response\\leg $\tilde\varphi_i$ ($ri$)};
  \node[below=2pt,align=center] at (v.south)
    {head: physical\\leg $\varphi_j$ ($pi$)};
\end{tikzpicture}
\caption{The edge convention. A directed edge $u\to v$ is a Wick contraction of a
response leg at the tail with a physical leg at the head; its value is
$G_{\mathrm{ft}}$ indexed (physical-row, response-col).}
\label{fig:ta-edge}
\end{figure}

\subsection{Retarded causality, stated two ways}

The retarded Green's function $\Gret(t)$ is nonzero only for $t>0$. In Fourier
space that is the statement that \emph{all poles of $\det\hat K(\omega)$ lie in
the upper half of the complex $\omega$ plane}, $\operatorname{Im}\omega_k > 0$, so
that closing the contour the ``correct'' way reproduces the retarded boundary
condition. There are two equivalent manifestations, and Phase F uses the first:

\begin{itemize}
  \item \textbf{Structural.} A diagram built only from retarded responses and
        noise vertices is a \term{directed acyclic graph} (DAG): there is no way
        to return to a vertex by following time-forward arrows. A directed cycle
        would close a loop of strictly-retarded propagators, which integrates to
        zero (or signals an acausal construction). This is the operative check,
        \code{D.is\_directed\_acyclic()}.
  \item \textbf{Analytic.} If explicit pole locations are supplied, every pole
        must have $\operatorname{Im} > 0$. This lives in
        \code{check\_pole\_structure}, but in the current pipeline the typed-diagram
        path passes \emph{no} poles, so only the DAG check runs.
\end{itemize}

% =====================================================================
\section{Phase D --- the degree filter}\label{sec:ta-filter}

We begin with the cheapest phase. Phase D
(\file{msrjd/diagrams/filter.py}) never looks at a single field type; it asks
only ``could the \emph{degrees} of this graph's vertices conceivably be supplied
by the theory's monomials?'' and throws away the hopeless graphs in
$O(\text{vertices})$ time.

\subsection{Classifying a prediagram's vertices}

The first helper sorts the non-leaf vertices of a prediagram into the two kinds
of internal vertex. A \term{source} is a noise vertex: it has \emph{no} incoming
edges (in-degree $0$), because a pure response-field monomial offers only
response legs, which feed \emph{outgoing} edges. Anything else is an
\term{interaction} vertex.

\begin{lstlisting}
def classify_prediagram_vertices(D, leaves):  # filter.py:12-40
    leaf_set = set(leaves)
    source = []
    interaction = []
    for v in D.vertices():
        if v in leaf_set:
            continue
        if D.in_degree(v) == 0:
            source.append(v)
        else:
            interaction.append(v)
    return source, interaction
\end{lstlisting}

\code{D} is a Sage \code{DiGraph} (we meet Sage properly in
\S\ref{sec:ta-tools}); \code{D.in\_degree(v)} is the number of edges pointing
into $v$. The exact same source-vs-interaction split, by the same in-degree-zero
test, reappears verbatim inside Phase E
(\file{type\_assignment.py:167--177}) --- the two phases agree on what counts as a
noise vertex.

\subsection{The filter itself}

\begin{lstlisting}
def filter_prediagrams(prediagrams, vertex_types, source_types):  # filter.py:43-81
    from msrjd.core.vertices import available_degrees
    int_degs, src_degs = available_degrees(vertex_types, source_types)
    kept = []
    for pd in prediagrams:
        D, G, leaves, internal = pd
        source_verts, interaction_verts = classify_prediagram_vertices(D, leaves)
        valid = True
        for v in source_verts:
            if D.out_degree(v) not in src_degs:
                valid = False
                break
        if valid:
            for v in interaction_verts:
                if (D.in_degree(v), D.out_degree(v)) not in int_degs:
                    valid = False
                    break
        if valid:
            kept.append(pd)
    return kept, len(prediagrams) - len(kept)
\end{lstlisting}

The pivot is \code{available\_degrees(vertex\_types, source\_types)}, defined in
\file{msrjd/core/vertices.py:777}. It returns two sets:

\begin{itemize}
  \item \code{int\_degs} --- the set of $(\text{in\_degree},
        \text{out\_degree})$ \emph{pairs} that the theory's interaction vertices
        can supply, built as
        \code{\{(vt.in\_degree, vt.out\_degree) for vt in vertex\_types\}}
        (\file{vertices.py:793}). Here the \code{in\_degree} of a \code{VertexType}
        is the count of its physical legs and \code{out\_degree} the count of its
        response legs;
  \item \code{src\_degs} --- the set of \code{out\_degree} \emph{values} the
        sources can supply, \code{\{st.out\_degree for st in source\_types\}}
        (\file{vertices.py:794}).
\end{itemize}

Each prediagram passes only if \emph{every} source vertex has an out-degree in
\code{src\_degs} \emph{and} every interaction vertex has a $(\text{in},
\text{out})$ pair in \code{int\_degs}. The first failing vertex flips
\code{valid} to \code{False} and short-circuits the loop --- there is no point
checking the rest.

\begin{example}[Why this saves work]
Suppose the theory has only a cubic interaction $\varphi^2\tilde\varphi$ (so
\code{int\_degs} contains only the bidegree $(2,1)$) and a single quadratic noise
$\tilde\varphi^2$ (so \code{src\_degs}$=\{2\}$). A prediagram containing a
degree-4 internal vertex --- a quartic --- is rejected here in a handful of
integer comparisons, instead of after a fruitless backtracking search in Phase E
that would have found no candidate type for that vertex and bailed out anyway.
The filter is purely a performance gate; it never changes the final answer.
\end{example}

\begin{note}
Phase D is \emph{not} in the production orchestrator. The loop enumerator
(Phase~C, Chapter~\ref{ch:enumeration}) already emits only degree-valid
prediagrams for the chosen $(k,\ell)$, so \code{enumerate\_unique\_diagrams}
skips the explicit filter. It remains the canonical cheap pre-screen --- it is
unit-tested in \file{tests/test\_filter.py} and is the right tool for any caller
that enumerates prediagrams more loosely --- but a reader tracing the live
pipeline will not see it called.
\end{note}

% =====================================================================
\section{Phase E --- type assignment as constraint satisfaction}\label{sec:ta-typeassign}

This is the engine room. Given one bare prediagram and the theory's menu of
vertex and source types, Phase E
(\file{msrjd/diagrams/type\_assignment.py}) enumerates \emph{every} valid way to
turn the bare graph into a labelled Feynman diagram. ``Valid'' means three things
must all hold simultaneously: vertex degrees must match a real monomial, every
leg must point the right way (response legs feed outgoing edges, physical legs
feed incoming edges), and the propagator the labelling demands on each edge must
not be identically zero. It is a backtracking constraint solver, and its output
is a stream of \code{TypedDiagram} objects.

\subsection{The output object: \texttt{TypedDiagram}}

Before the algorithm, the thing it produces. A \code{TypedDiagram}
(\file{type\_assignment.py:29--71}) is a \code{\_\_slots\_\_} class --- meaning
its instances carry only the five named attributes and cannot accidentally grow
others, which keeps them memory-lean since we may make thousands. The five
fields are:

\begin{description}
  \item[\code{prediagram}] the \code{(D, G, leaves, internal)} 4-tuple this was
        built from.
  \item[\code{vertex\_assignments}] a dict
        \code{\{vertex\_id: VertexType | SourceType\}} for the internal vertices
        (empty in the no-internal case).
  \item[\code{edge\_types}] a dict
        \code{\{(u, v, label): (resp\_leg, phys\_leg)\}}. The 3-tuple edge key
        keeps parallel edges distinct (the \code{label} disambiguates them); each
        value is a pair of \code{(field\_base, pop\_idx)} legs.
  \item[\code{external\_legs}] a dict
        \code{\{leaf\_vertex: (field\_base, pop\_idx)\}}.
  \item[\code{propagator\_indices}] a dict
        \code{\{(u, v, label): (ri, pi)\}} --- the row/col into
        $G_{\mathrm{ft}}$, stored as $(ri, pi)$ (recall the lookup flips this to
        $(pi, ri)$).
\end{description}

It also defines \code{\_\_getstate\_\_}/\code{\_\_setstate\_\_}
(\file{type\_assignment.py:60--65}) so that a \code{\_\_slots\_\_} object can be
pickled across the fork pool (\S\ref{sec:ta-fork}), and a \code{\_\_repr\_\_}
that summarises vertex / edge / assignment counts.

\subsection{A field-leg is a \texttt{(base, pop\_idx)} tuple}

Throughout, a single field-leg identity is a tuple \code{(field\_base,
pop\_idx)}, e.g. \code{('nt', 1)} or \code{('dn', 2)}. The \code{field\_base} is
the generator's base name (letters); the \code{pop\_idx} is a 1-based
\term{population index} parsed off the trailing digits (the \code{2} in
\code{nt2}), distinguishing populations or components. Index \code{0} means ``no
trailing digits.'' Whether a leg is a response or physical leg is \emph{not}
encoded in the tuple --- it is decided by which generator list the base name came
from, which is exactly what the next helper records.

\subsection{The bridge from legs to matrix indices}

A leg \code{(base, pop\_idx)} must be turned into a row or column of
$G_{\mathrm{ft}}$. \code{build\_field\_index\_map}
(\file{type\_assignment.py:76--110}) builds that bridge once:

\begin{lstlisting}
def build_field_index_map(ring_var_names, n_tilde):  # type_assignment.py:76-110
    from msrjd.core.vertices import _parse_field_name
    resp_index = {}
    phys_index = {}
    resp_counter = 0
    phys_counter = 0
    for i, name in enumerate(ring_var_names):
        base, pop_idx = _parse_field_name(name)
        if i < n_tilde:
            resp_index[(base, pop_idx)] = resp_counter
            resp_counter += 1
        else:
            phys_index[(base, pop_idx)] = phys_counter
            phys_counter += 1
    return resp_index, phys_index
\end{lstlisting}

\code{ring\_var\_names} is the ordered list of polynomial-ring generator names,
e.g. \code{['vt1','vt2','nt1','nt2','dv1','dv2','dn1','dn2']}, and \code{n\_tilde}
is the number of response generators --- crucially, the response generators come
\emph{first}. So a generator at position \code{i < n\_tilde} is a response field
(it gets the next \code{resp\_counter}); otherwise it is a physical field. The
helper \code{\_parse\_field\_name} (from \file{core/vertices.py:327}) splits a
name like \code{'nt1'} into \code{('nt', 1)}. The two returned dicts,
\code{resp\_index} and \code{phys\_index}, are the maps every later phase uses to
turn a leg into a matrix index.

\subsection{The top-level enumerator}

\code{enumerate\_typed\_diagrams} (\file{type\_assignment.py:115--266}) is the
top-level generator for \emph{one} prediagram. Reading its body in order is the
clearest way to understand the algorithm; it proceeds in seven steps.

\paragraph{Step 1 --- the propagator zero mask
(\file{:158--165}).} Before searching, it precomputes which entries of
$G_{\mathrm{ft}}$ are identically zero, so the inner loop can prune any labelling
that would place a zero propagator:

\begin{lstlisting}
g_zero_mask = {}
for _i in range(G_ft.nrows()):
    for _j in range(G_ft.ncols()):
        g_zero_mask[(_i, _j)] = (str(G_ft[_i, _j]) == '0')   # type_assignment.py:165
\end{lstlisting}

The test is \code{str(G\_ft[i,j]) == '0'} --- a deliberately \emph{structural,
string-based} test, \emph{not} the semantically-correct \code{is\_zero()}. The
reason is performance, and it matters enough that \S\ref{sec:ta-zeromask} is
devoted to it. For now: this catches only entries that are the literal
\code{SR(0)}, which is exactly what the propagator builder writes into empty
cells.

\paragraph{Step 2 --- classify the non-leaf vertices
(\file{:167--177}).} The same in-degree-zero split as Phase D, producing
\code{source\_verts} and \code{interaction\_verts}; their concatenation
\code{ordered\_internal = source\_verts + interaction\_verts} fixes a
deterministic processing order for the backtracker.

\paragraph{Step 3 --- per-vertex incident-edge lists (\file{:179--184}).} For
every vertex, cache \code{D.outgoing\_edges(v)} and \code{D.incoming\_edges(v)}.
These come back as Sage 3-tuples \code{(u, v, label)} so that parallel edges stay
distinct.

\paragraph{Step 4 --- candidate types per vertex (\file{:186--201}).} A source
vertex of out-degree \code{od} gets every \code{SourceType} with
\code{st.out\_degree == od}; an interaction vertex of bidegree \code{(ind, od)}
gets every \code{VertexType} whose \code{(in\_degree, out\_degree)} matches. If
\emph{any} vertex has \emph{no} candidate, the whole prediagram is impossible and
the generator simply \code{return}s, yielding nothing.

\paragraph{Step 5 --- leaf direction constraints (\file{:204--213}).} Each
external leaf is classified by the edges touching it: an only-outgoing leaf can
carry only a response field (\code{'resp'}); only-incoming, only a physical field
(\code{'phys'}); a leaf with both is \code{'both'}.

\paragraph{Step 6 --- external-field permutations (\file{:223--242}).} This is
the step a reader is most likely to find surprising, so it gets its own
subsection below. The enumerator loops over \emph{every distinct ordering} of
\code{external\_fields} across the leaves, keeping only orderings in which each
leaf's field is compatible with its direction.

\paragraph{Step 7 --- dispatch (\file{:244--266}).} If there are no internal
vertices, the edges merely connect leaves, and the work is delegated to
\code{\_try\_build\_diagram\_no\_internal}. Otherwise the enumerator takes the
Cartesian product \code{product(*candidate\_lists)} over the internal vertices'
candidate type lists, and for each combination builds a \code{vert\_assignment}
and delegates to \code{\_try\_build\_diagram}.

\subsection{Why enumerate external-field permutations at all?}

Step 6 deserves its own explanation because it looks like double work. The
relevant code is:

\begin{lstlisting}
for ext_perm in _distinct_permutations(tuple(external_fields)):  # type_assignment.py:223
    ext_assignment = {}
    valid_ext = True
    for leaf_idx in range(len(ext_perm)):
        lf = leaves[leaf_idx]
        field = ext_perm[leaf_idx]
        direction = leaf_directions[lf]
        if direction == 'resp' and field not in resp_index:
            valid_ext = False; break
        if direction == 'phys' and field not in phys_index:
            valid_ext = False; break
        if direction == 'both':
            if field not in resp_index and field not in phys_index:
                valid_ext = False; break
        ext_assignment[lf] = field
    if not valid_ext:
        continue
    ...
\end{lstlisting}

The upstream prediagram dedup (Phase~C) treats all leaves as interchangeable
stubs, so a \emph{single} prediagram stands in for every way of distributing the
external fields across its leaves. To recover the diagrams that differ in
\emph{which leaf carries which field} --- for example \code{dn2} sitting on a
source-leaf versus on an interaction-leaf --- we must materialise those
distributions here. The helper \code{\_distinct\_permutations}
(\file{type\_assignment.py:22--24}) is a one-liner that returns
\code{set(permutations(seq))}, so two identical external fields do not produce two
``different'' orderings. Any orderings that nevertheless turn out to be truly
identical diagrams are merged later by Phase~G's deduplication.

\subsection{The degenerate case: a bare propagator}

When a prediagram has \emph{no} internal vertices --- e.g. a single edge directly
joining two external leaves, which is just a bare propagator ---
\code{\_try\_build\_diagram\_no\_internal} (\file{type\_assignment.py:269--305})
handles it directly. For each edge $u\to v$: the tail $u$ contributes the
response leg (its external field), the head $v$ the physical leg; both must be
direction-compatible; we look up \code{ri = resp\_index[resp\_field]} and
\code{pi = phys\_index[phys\_field]}; and we reject the edge if
\code{g\_zero\_mask[(pi, ri)]} is true --- note again the \textbf{(pi, ri)} order.
On success it records the edge's types and propagator indices and yields a
\code{TypedDiagram} with empty \code{vertex\_assignments}.

\begin{example}[A single edge, end to end]
The micro-example from \file{tests/test\_type\_assignment.py} makes the data
shapes concrete. With a prediagram that is one directed edge $1\to2$ between two
leaves, and external fields \code{[('nt', 1), ('dn', 1)]}:
\begin{lstlisting}
pd = (DiGraph([(1, 2)]), None, [1, 2], [])
external_fields = [('nt', 1), ('dn', 1)]
results = list(enumerate_typed_diagrams(pd, external_fields, [], [],
                                        G_ft, resp_idx, phys_idx))
# -> one TypedDiagram with
#    edge_types         = {(1, 2, None): (('nt',1), ('dn',1))}
#    external_legs      = {1: ('nt',1), 2: ('dn',1)}
#    propagator_indices = {(1, 2, None): (ri, pi)}
#    vertex_assignments = {}
\end{lstlisting}
The response field \code{nt} sits at the tail, the physical field \code{dn} at
the head, and the edge carries $G_{\mathrm{ft}}[pi, ri]$ --- one bare propagator,
exactly as Fig.~\ref{fig:ta-edge} drew it.
\end{example}

\subsection{Leg matchings and the backtracker}

For a prediagram \emph{with} internal vertices, \code{\_try\_build\_diagram}
(\file{type\_assignment.py:308--337}) sets up the search. For each internal
vertex it computes the list of ways to map that vertex type's legs onto its
incident edges, via \code{\_leg\_matchings}; if any vertex admits zero such
matchings the type combination is dead and it \code{return}s. Otherwise it
launches the recursive backtracker.

\code{\_leg\_matchings(vertex\_type, out\_edges, in\_edges)}
(\file{type\_assignment.py:340--416}) enumerates the distinct bijections from a
vertex's legs to its edges: each outgoing edge must receive a \emph{response}
leg, each incoming edge a \emph{physical} leg. Its core is short:

\begin{lstlisting}
resp_legs = vertex_type.response_legs
has_phys = hasattr(vertex_type, 'physical_legs')
phys_legs = vertex_type.physical_legs if has_phys else []     # sources have none
if len(resp_legs) != len(out_edges) or len(phys_legs) != len(in_edges):
    return
from sympy.utilities.iterables import multiset_permutations    # type_assignment.py:406
resp_perms = (list(multiset_permutations(list(resp_legs))) if resp_legs else [[]])
phys_perms = (list(multiset_permutations(list(phys_legs))) if phys_legs else [[]])
for rp in resp_perms:
    resp_map = dict(zip(out_edges, rp))
    for pp in phys_perms:
        phys_map = dict(zip(in_edges, pp))
        yield resp_map, phys_map
\end{lstlisting}

The \code{has\_phys} check is how a source is told apart from an interaction
vertex: a \code{SourceType} has no \code{physical\_legs} attribute at all, so
\code{phys\_legs} defaults to the empty list and the source contributes only
response legs. The choice of \code{multiset\_permutations} over a full
\code{N!} enumeration is the single most important optimisation in the file, and
\S\ref{sec:ta-canonical} is devoted to proving it is correct.

The recursive core \code{\_backtrack} (\file{type\_assignment.py:419--538})
carries three pieces of state: \code{vertex\_idx} (which internal vertex we are
assigning), and two partial maps \code{assigned\_resp} and \code{assigned\_phys}
(\code{\{edge: leg\}} so far). It has two parts.

\paragraph{The base case (\file{:431--472}).} All internal vertices are assigned.
For each edge, the code resolves its response leg --- from \code{assigned\_resp}
if an internal tail set it, else from the external field if the tail is a leaf
--- and, symmetrically, its physical leg. It looks up \code{ri}/\code{pi}; if
either is missing or the propagator is zero (\code{g\_zero\_mask[(pi, ri)]}), it
\code{return}s. On full success it yields the finished \code{TypedDiagram}.

\paragraph{The recursive step (\file{:474--538}).} For each leg-matching option at
the current vertex, it merges that option into fresh \code{new\_resp}/
\code{new\_phys} copies and runs an \emph{early propagator-consistency check} on
every edge that is now fully determined --- both ends known. It checks both
internal-internal edges and edges where one end is a leaf, and prunes immediately
if any propagator is zero. Only if the partial assignment is still
\code{consistent} does it recurse to \code{vertex\_idx + 1}. This early pruning
is what keeps the search tractable on large diagrams: a dead branch is abandoned
the moment a single edge is forced onto a zero propagator, rather than after a
full assignment.

\begin{gotcha}
The convention note at \file{type\_assignment.py:490--502} is worth reading in
the source. It states explicitly that the correct lookup order is
\code{g\_zero\_mask[(pi, ri)]} (physical-row, response-col) and records a real
historical bug: a previous version used \code{[ri, pi]}, which happened to read
the right entry for \emph{diagonal} couplings (where the response and physical
indices coincide) but silently filtered out valid diagrams for \emph{off-diagonal}
ones. The cited example is a GTaS coupling $\tilde m \to \delta n$ with physical
index $0$ and response index $4$: there \code{[ri, pi]} and \code{[pi, ri]} point
at different cells, so the wrong order rejected diagrams whose propagator was in
fact nonzero. Any new propagator-touching code must index $(pi, ri)$.
\end{gotcha}

% ---------------------------------------------------------------------
\subsection{The canonical-leg-matching optimisation}\label{sec:ta-canonical}

Why does \code{\_leg\_matchings} use SymPy's \code{multiset\_permutations}
instead of \code{itertools.permutations}? The difference is a factor that can be
$15$--$20\times$ on the size of the intermediate diagram set, and getting it
right without changing the final number is a small theorem the code's docstring
(\file{type\_assignment.py:356--387}) carefully proves.

\begin{defn}
\textbf{Multiset permutations.}
For a multiset of length $N$ with multiplicities $n_r$ (i.e. $n_r$ copies of the
$r$-th distinct element), the number of \emph{distinct} orderings is the
multinomial coefficient
\[
  \frac{N!}{\prod_r n_r!},
\]
not the full $N!$. SymPy's \code{multiset\_permutations(seq)} yields each distinct
ordering exactly once. Because the leg values are hashable \code{(field\_base,
pop\_idx)} tuples, the multiset comparison just works.
\end{defn}

The argument for correctness is this. Earlier code enumerated all $N!$ index
permutations and \emph{relied on} the downstream deduplication to discard the
overcount. That overcount factor is \emph{exactly} the symmetry factor
$\Sym(\Gamma)$. So the old pipeline generated $\Sym(\Gamma)$ identical copies of
each diagram, each copy carrying its full weight; dedup kept one; and the
integrator then multiplied by $\Sym(\Gamma)$. The new code skips the inflation
entirely: it generates the single \emph{canonical} ordering directly, and the
integrator multiplies by the \emph{same} $\Sym(\Gamma)$. The output is
bit-identical; only the intermediate count (and the cold-start enumeration wall
time) shrinks. This equivalence is pinned by the regression tests
\code{test\_leg\_matchings\_canonical} and
\code{test\_enumerate\_typed\_signatures\_match\_pre\_change} in
\file{tests/test\_type\_assignment.py}.

\begin{gotcha}
This optimisation is correct \emph{only because} the integrator multiplies by
$\Sym(\Gamma)$ afterward. The two are a matched pair. If you ever change how the
symmetry factor is applied downstream, you must revisit \code{\_leg\_matchings} ---
the regression tests above exist precisely to catch a drift between the two.
\end{gotcha}

% ---------------------------------------------------------------------
\subsection{The propagator zero mask, and why a string test}\label{sec:ta-zeromask}

We deferred one design choice from Step 1: why is the zero test
\code{str(G\_ft[i,j]) == '0'} rather than the obviously-correct
\code{G\_ft[i,j].is\_zero()}? The answer is entirely about cost, and the code's
comment (\file{type\_assignment.py:144--157}) spells it out.

$G_{\mathrm{ft}}$ entries can be long rational functions of the symbolic
parameters. Calling \code{is\_zero()} (or, equivalently, \code{e == SR(0)}) on
such an entry triggers Sage's full symbolic canonicalisation engine, which can
cost \emph{several seconds per entry} on a rich theory --- and the backtracker
would otherwise hit that cost on every leg assignment. The string test sidesteps
canonicalisation completely: it returns \code{True} only for entries that are
\emph{literally} the symbol \code{SR(0)}, which is what the propagator builder
writes into empty cells of $K^{-1}$.

The trade-off is a deliberate, harmless over-generation. An entry that is zero
only \emph{after cancellation} --- a nontrivial expression that happens to
simplify to $0$ --- is \emph{not} caught by the string test, so it slips through
and a diagram or two carrying it gets generated. But those diagrams simply
integrate to $0$ downstream. The string test can therefore over-generate, never
under-generate: it never rejects a valid diagram, so it can never produce a wrong
number, only a few extra ``vanishing'' diagrams.

\begin{note}
This is a recurring pattern in the codebase: when an exact symbolic decision is
expensive and a cheap structural proxy errs only on the safe side (here,
over-generation that integrates to zero), the cheap proxy wins. Be aware, though,
that the typed-diagram set you inspect may contain a handful of diagrams that
contribute nothing --- this is by design, not a bug.
\end{note}

% =====================================================================
\section{Phase F --- causality}\label{sec:ta-causality}

Phase F (\file{msrjd/diagrams/causality.py}) prunes typed diagrams that violate
the retarded boundary condition. It is short, because in the current pipeline
only the structural DAG check is live.

\subsection{The per-diagram check}

\begin{lstlisting}
def check_causality(typed_diagram, pole_vals=None):  # causality.py:97-123
    D = typed_diagram.prediagram[0]
    # Structural check: the prediagram must be a DAG
    if not D.is_directed_acyclic():
        return False, 'Prediagram contains a directed cycle - acausal.', []
    # Pole check (if available)
    if pole_vals is not None and len(pole_vals) > 0:
        return check_pole_structure(pole_vals)
    return True, 'Structural DAG check passed. No poles to verify.', []
\end{lstlisting}

It pulls the directed graph \code{D = typed\_diagram.prediagram[0]} and asks
\code{D.is\_directed\_acyclic()}. A directed cycle means you can return to a
vertex by following retarded arrows --- a closed loop of strictly-retarded
propagators --- which is acausal, so the diagram is rejected. Only if
\code{pole\_vals} are supplied does it fall through to the analytic pole check;
in the live pipeline none are passed, so \textbf{the DAG check is the operative
gate}.

\code{filter\_causal(typed\_diagrams)} (\file{causality.py:126--152}) is the
list-level wrapper the orchestrator calls (\file{pipeline/\_diagrams.py:185}): it
runs \code{check\_causality} on each diagram and partitions them into a
\code{kept} list and a list of one human-readable rejection reason per discarded
diagram, returning \code{(kept, n\_discarded, discarded\_details)}.

\subsection{The latent pole machinery}

\code{check\_pole\_structure(pole\_vals, omega=None)}
(\file{causality.py:19--94}) is the analytic counterpart, present but dormant.
Given a list of pole locations it coerces each with \code{SR(pole)}, takes its
imaginary part with Sage's \code{imag\_part}, attempts \code{simplify\_full()},
and then tries to decide the sign:

\begin{itemize}
  \item \code{bool(im\_simplified > 0)} --- the pole passes (upper half-plane);
  \item \code{bool(im\_simplified <= 0)} --- the pole fails;
  \item \code{bool(im\_simplified == 0)} --- a pole exactly on the real axis is
        treated as a \emph{failure} (it is not retarded).
\end{itemize}

Each \code{bool(...)} is wrapped in \code{try/except (TypeError, ValueError)}
because Sage raises when it cannot decide a symbolic inequality; an undecidable
pole is recorded as \emph{conditional}. The final verdict: any failure
$\Rightarrow$ \code{False}; else any conditional $\Rightarrow$ \code{True}
\emph{with} the symbolic $\operatorname{Im}(\text{pole}) > 0$ constraints returned
as \code{conditions}; else everything passed. So an inconclusive symbolic pole
returns \code{passed=True} (carrying its constraint), while a real-axis pole is a
hard fail.

% =====================================================================
\section{Phase G --- the symmetry factor and deduplication}\label{sec:ta-symmetry}

Phase G (\file{msrjd/diagrams/symmetry.py}) does the two jobs whose
correctness matters most, because a mistake here produces a numerically
\emph{plausible} but \emph{wrong} answer rather than an obvious crash. First, it
deduplicates the typed diagrams. Second, it computes the symmetry factor
$\Sym(\Gamma)$. Both rest on a single construction --- a coloured incidence
digraph fed to \code{nauty} --- so we build that first.

\subsection{The canonical Feynman rule}

The code implements what its comments call ``Path A'': the canonical Feynman
rule (\file{symmetry.py:403--446}),
\begin{equation}
  \Sym(\Gamma) \;=\;
  \frac{\displaystyle\prod_v \prod_\ell n_{v,\ell}!}
       {\bigl|\Aut_{\text{fixed ext}}(\Gamma)\bigr|}.
  \label{eq:ta-sym}
\end{equation}
There are two pieces, a numerator and a denominator, and the per-diagram weight
the integrator ultimately builds (\file{symmetry.py:421--423}) is
\begin{equation}
  \mathrm{weight}(\Gamma) \;=\;
  \Sym(\Gamma)\;\times\;\prod_v c_{\text{user},v}\;\times\;\prod_e P_e
  \;\times\;\int \dd t_v,
  \label{eq:ta-weight}
\end{equation}
where $c_{\text{user},v}$ is the \emph{literal} coefficient written in the action
(no implicit $1/n!$), $P_e$ is the propagator on edge $e$, and the integral runs
over the internal vertex times.

\paragraph{The numerator --- the Wick leg factor.} For each vertex $v$ and each
leg-type $\ell$ (a leg-type lumps \emph{direction} --- response vs physical ---
with the \emph{field}), let $n_{v,\ell}$ be the number of legs of type $\ell$ at
$v$. The numerator is the product of $n_{v,\ell}!$ over all of them --- ``the Wick
combinatorial you would get if every leg at every vertex were distinguishable.''
A degree-3 vertex $\varphi^2\tilde\varphi$ contributes $2!$ from its two
identical $\varphi$ legs. This is \code{\_wick\_leg\_factor}
(\file{symmetry.py:115--146}):

\begin{lstlisting}
def _wick_leg_factor(typed_diagram):  # symmetry.py:115-146
    n = 1
    for v, vtype in typed_diagram.vertex_assignments.items():
        resp_counts = Counter(vtype.response_legs)      # outgoing/response legs
        for c in resp_counts.values():
            n *= factorial(c)
        if hasattr(vtype, 'physical_legs'):             # incoming/physical legs
            phys_counts = Counter(vtype.physical_legs)
            for c in phys_counts.values():
                n *= factorial(c)
    return n
\end{lstlisting}

\code{Counter} (from the standard library's \code{collections}) tallies the
multiplicity of each distinct leg, and \code{factorial} (from \code{math}) is
$n!$.

\paragraph{The denominator --- the automorphism order.}
$|\Aut_{\text{fixed ext}}(\Gamma)|$ is the order of the automorphism group of the
\emph{fully coloured} diagram, with the external legs \emph{pinned} at their
canonical positions. An \term{automorphism} is a relabelling of internal vertices
and edges that maps the diagram to itself while preserving every label physics
cares about: vertex type, vertex coefficient, the propagator on each edge, and
the direction of each edge. It captures three families of redundancy ---
same-type internal vertex swaps, parallel-edge swaps between the same vertex
pair, and self-loop / tadpole leg-pair swaps --- and dividing the Wick numerator
by it removes exactly the over-counting those redundancies introduce.

\subsection{The coloured incidence digraph}

The obstacle is that \code{nauty} (via Sage) colours \emph{vertices}, not edges
--- yet an automorphism must respect the colour (the propagator type) of each
\emph{edge}. The resolution, in \code{\_colored\_incidence\_digraph}
(\file{symmetry.py:149--244}), is to build a fresh bipartite ``vertex-nodes plus
edge-nodes'' graph:

\begin{itemize}
  \item each Feynman vertex becomes a node \code{('V', v)};
  \item \emph{each propagator edge becomes its own node} \code{('E', ek)};
  \item an original edge $u\to v$ is encoded as \emph{two} bipartite edges,
        \code{('V', u) -> ('E', ek)} and \code{('E', ek) -> ('V', v)}.
\end{itemize}

Now an edge's colour lives on its edge-node, where \code{nauty} can see it. The
node colours (the \emph{partition}, in nauty's language) are:

\begin{description}
  \item[Internal vertex node] colour
        \code{('vertex', type(vt).\_\_name\_\_, str(vt.coefficient), vt.bigrade)}
        --- so only same-type, same-coefficient, same-bigrade vertices may be
        swapped.
  \item[External leaf node] if \code{fix\_external=True}, colour
        \code{('leaf', v, field)} --- each leaf gets its own colour, so leaves can
        never move; this gives $\Aut_{\text{fixed ext}}$. If
        \code{fix\_external=False}, colour \code{('leaf', field)} --- same-field
        leaves share a colour and so \emph{may} permute; this gives
        $\Aut_{\text{free ext}}$.
  \item[Edge node] colour \code{('edge', resp\_leg, phys\_leg, prop)}, where
        \code{prop} is the $(ri, pi)$ propagator-index tuple --- so an
        automorphism must preserve the propagator type \emph{and} direction on
        every edge.
\end{description}

Parallel edges between the same vertex pair with identical colour become
genuinely swappable edge-nodes, which is exactly the symmetry the Feynman rules
want counted. The function returns the digraph \code{D}, the \code{partition}
(the list of colour classes, ready to hand to nauty), and bookkeeping maps.

\subsection{Counting automorphisms with nauty}

With the coloured digraph in hand, \code{\_automorphism\_order}
(\file{symmetry.py:328--340}) is the single place the automorphism count is
invoked:

\begin{lstlisting}
def _automorphism_order(typed_diagram, fix_external=True):  # symmetry.py:328-340
    D, partition, _, _ = _colored_incidence_digraph(
        typed_diagram, fix_external=fix_external)
    grp = D.automorphism_group(partition=partition)
    return int(grp.order())
\end{lstlisting}

\code{D.automorphism\_group(partition=partition)} asks Sage --- and underneath,
\code{nauty} --- for the group of colour-preserving relabellings; \code{grp} is a
Sage permutation group, and we keep only its \code{order()} (an integer),
discarding the group itself. Passing \code{fix\_external=True} yields
$|\Aut_{\text{fixed ext}}|$; \code{False} yields $|\Aut_{\text{free ext}}|$.

The production symmetry factor \code{combinatorial\_factor}
(\file{symmetry.py:403--446}) then divides numerator by denominator:

\begin{lstlisting}
def combinatorial_factor(typed_diagram):  # symmetry.py:403-446
    numer = _wick_leg_factor(typed_diagram)
    aut = _automorphism_order(typed_diagram, fix_external=True)
    if aut <= 0 or numer % aut != 0:
        return numer // max(aut, 1)     # defensive floor-division fallback
    return numer // aut
\end{lstlisting}

The integer division should always be exact --- a finite group acting freely on a
finite set --- and the floor-division branch is purely defensive against a Sage
edge case. \code{compute\_all\_combinatorial\_factors}
(\file{symmetry.py:449--462}) is the obvious list comprehension over this.

\begin{note}
\file{symmetry.py} also carries an older, vertex-local way to count Wick
pairings: \code{\_vertex\_combinatorial\_factor} (\file{symmetry.py:51--112}),
which computes a per-vertex factor $M_v$ and would give
$\Sym(\Gamma) = \prod_v M_v$ in the ``Path B'' formulation. The production
\code{combinatorial\_factor} does \emph{not} use it --- it uses the Aut-based
Eq.~\eqref{eq:ta-sym}. The $M_v$ function is retained because it is the
per-vertex semantic that \S\ref{sec:ta-canonical}'s leg-matching optimisation
appeals to (the overcount factor it avoids generating is exactly this $M_v$
product), and as a documentation cross-check.
\end{note}

% ---------------------------------------------------------------------
\subsection{Deduplication and the complete invariant}\label{sec:ta-dedup}

Phase E over-generates: it emits the same diagram many times, relabelled and
leaf-permuted. Deduplication collapses those copies, and the key it uses must be
a \emph{complete} graph-isomorphism invariant --- two typed diagrams must hash
equal \emph{if and only if} they are isomorphic coloured graphs. A weaker
invariant that occasionally collides non-isomorphic diagrams would silently
\emph{drop} one class's integral, which is the worst kind of bug: a wrong number
with no error.

\code{diagram\_signature(td)} (\file{symmetry.py:467--524}) provides exactly such
a complete invariant, built from \code{nauty}'s canonical labelling:

\begin{lstlisting}
def diagram_signature(td):  # symmetry.py:467-524
    D, _, _, color_groups = _colored_incidence_digraph(td, fix_external=False)
    keys = sorted(color_groups.keys(), key=str)         # label-independent order
    partition = [color_groups[k] for k in keys]
    C, cert = D.canonical_label(partition=partition, certificate=True)
    cells = tuple(tuple(sorted(cert[v] for v in color_groups[k])) for k in keys)
    edges = tuple(sorted(C.edges(labels=False)))
    return (tuple(str(k) for k in keys), cells, edges)
\end{lstlisting}

\begin{defn}
\textbf{Canonical label and certificate.}
\code{nauty} computes a \term{canonical form}: a relabelling of a graph's
vertices into a normal order such that two graphs are isomorphic \emph{iff} their
canonical forms are identical. Sage exposes it as
\code{D.canonical\_label(partition=..., certificate=True)}, which returns a pair
\code{(C, cert)}: \code{C} is the canonically relabelled copy of \code{D}, and the
\term{certificate} \code{cert} is the dict mapping each original vertex id to its
canonical-label id.
\end{defn}

The signature is built from three pieces: the colour keys (as strings, for
hashability), \emph{which canonical ids each colour class occupies}
(\code{cells}), and the canonical edge list (\code{edges}). The \code{cells} piece
is essential --- the canonical edge list \emph{alone} would confuse two graphs of
the same shape but different colourings, so we must also record where each colour
landed. It is built with \code{fix\_external=False} so that same-field external
leaves are allowed to permute; leaf-permuted variants are thereby merged here and
re-expanded downstream by the integrator's mapping sum (\S\ref{sec:ta-comp}).

\code{deduplicate\_with\_multiplicities(typed\_diagrams)}
(\file{symmetry.py:553--593}) is the function the orchestrator actually calls
(\file{pipeline/\_diagrams.py:186}). It keeps one representative per signature in
first-seen order, alongside each equivalence class's size:

\begin{lstlisting}
sig_to_idx = {}
unique = []
multiplicities = []
for td in typed_diagrams:
    sig = diagram_signature(td)
    idx = sig_to_idx.get(sig)
    if idx is None:
        sig_to_idx[sig] = len(unique)
        unique.append(td)
        multiplicities.append(1)
    else:
        multiplicities[idx] += 1
return unique, multiplicities
\end{lstlisting}

(\code{deduplicate\_typed\_diagrams} (\file{symmetry.py:527--550}) is a thin
wrapper returning just the \code{unique} list.)

\begin{gotcha}
The \code{multiplicities} are \textbf{diagnostic only} --- never multiply them
back into the weight. Under Path A the entire weight is carried by
\code{combinatorial\_factor}; multiplying by the class size would double-count,
because the merged entries are isomorphic copies whose pairings $\Sym(\Gamma)$
already includes (\file{symmetry.py:556--568} is explicit). This subtlety has
teeth: the on-disk cache stage tag was bumped (to \code{unique\_typed\_mult\_v3})
specifically to invalidate caches predating the complete-invariant signature,
since loading them would resurrect the old class-collision bug.
\end{gotcha}

\begin{example}[The collision bug the complete invariant closed]
The history note at \file{symmetry.py:484--498} records the concrete failure. At
$k=3$, one loop, the OU$+ax^2+bx^3$ theory's $a^3$ sector has $11$ distinct
diagram classes. A previous, hand-rolled signature --- a depth-1 isomorphism
\emph{invariant} rather than a complete one --- collided $4$ of those $11$ into
$2$ representatives. Because the dedup multiplicity is correctly \emph{not}
multiplied back, the collided classes' integrals were silently dropped, and the
$\kappa_3$ one-loop coefficient came out $-\tfrac{68}{3}a^3$ instead of the
correct $-32\,a^3$. A canonical-form invariant \emph{cannot} collide
non-isomorphic graphs, so this failure mode is closed for every $(k, \ell,
\text{theory})$. The fix was validated against the exact Boltzmann series and a
brute-force labelled Wick enumeration per class.
\end{example}

% ---------------------------------------------------------------------
\subsection{External Wick compensation and the integrator hand-off}\label{sec:ta-comp}

There are \emph{two} automorphism orders, differing only in whether external
leaves may permute, and their ratio is the bridge to the integrator. With leaves
pinned we get $|\Aut_{\text{fixed ext}}|$ (the denominator of $\Sym(\Gamma)$);
with same-field leaves free we get $|\Aut_{\text{free ext}}|$. Since
$\Aut_{\text{fixed ext}}$ is a subgroup of $\Aut_{\text{free ext}}$, the latter
order is always an integer multiple of the former, and that integer is the
\term{external Wick compensation}:
\begin{equation}
  \mathrm{comp} \;=\;
  \frac{|\Aut_{\text{free ext}}(\Gamma)|}{|\Aut_{\text{fixed ext}}(\Gamma)|}.
  \label{eq:ta-comp}
\end{equation}
This is \code{external\_wick\_compensation} (\file{symmetry.py:343--400}):

\begin{lstlisting}
def external_wick_compensation(typed_diagram):  # symmetry.py:343-400
    aut_free = _automorphism_order(typed_diagram, fix_external=False)
    aut_fixed = _automorphism_order(typed_diagram, fix_external=True)
    if aut_fixed <= 0 or aut_free % aut_fixed != 0:
        raise RuntimeError(
            f'external_wick_compensation: |Aut_free|={aut_free} not '
            f'divisible by |Aut_fixed|={aut_fixed}')
    return aut_free // aut_fixed
\end{lstlisting}

Why it exists: the integrators enumerate \emph{every} field-respecting assignment
of canonical external positions to a diagram's leaves and sum the integrand over
all of them. By the \term{orbit--stabilizer theorem} ($|\text{orbit}| =
|\text{group}| / |\text{stabilizer}|$), two assignments give the same
pinned-external diagram iff they differ by a leaf-permuting automorphism, and the
stabilizer of any assignment is exactly $\Aut_{\text{fixed ext}}$. So each
distinct pinned diagram is counted $|\text{free}|/|\text{fixed}|$ times by the
mapping sum; dividing by \code{comp} makes the sum equal the sum over distinct
pinned diagrams --- which is what the Feynman rules require, since each pinned
diagram already carries its full weight through $\Sym(\Gamma)$, whose denominator
is the \emph{same} $\Aut_{\text{fixed ext}}$.

\begin{gotcha}
\code{external\_wick\_compensation} fails \emph{loud}: a non-divisor raises
\code{RuntimeError} rather than guessing (\file{symmetry.py:391--399}). This is a
deliberate asymmetry with \code{combinatorial\_factor}, which falls back to
floor-division (\file{symmetry.py:440--445}). The reasoning: a wrong compensation
silently \emph{mis-weights} diagrams --- the exact bug class this function exists
to fix --- whereas the Wick factor's divisibility is a theorem that should always
hold. So the compensation refuses to guess, while the Wick factor is allowed a
defensive fallback.
\end{gotcha}

\begin{gotcha}
A nearby function, \code{vertex\_role\_signature} (\file{symmetry.py:247--325}),
is a hashable per-vertex signature invariant under relabelling. It once served as
a Phase~J external divisor (a $\prod N!$ over leaves grouped by this signature),
but that depth-1 heuristic \emph{over-divided} when two structurally distinct
vertices happened to share a signature without any automorphism relating them ---
breaking every $k\ge3$ one-loop cumulant. It has been \emph{superseded} by the
exact \code{external\_wick\_compensation} of Eq.~\eqref{eq:ta-comp}. The function
survives in the file but the production weight path no longer uses it; do not
reintroduce the heuristic divisor.
\end{gotcha}

\subsection{Coefficient classification}\label{sec:ta-coeff}

One more Phase G function is called directly by the integrators:
\code{classify\_coefficient\_factors} (\file{symmetry.py:618--698}). Its job is to
split each vertex coefficient into a \emph{scalar prefactor} that can be pulled
outside the integral versus \emph{factors that must stay inside} (time-dependent
couplings, non-white noise amplitudes). It returns a dict with keys \code{'Scal'}
(the symmetry factor $\Sym(\Gamma)$ itself), \code{'scalar\_prefactor'} (the
product of all pull-out-able pieces, including $\Sym(\Gamma)$),
\code{'vertex\_time\_factors'}, \code{'source\_time\_info'}, and
\code{'is\_stationary'}.

The mechanics, at a high level: it seeds \code{scalar\_parts} with \code{SR(Scal)}
and, for each vertex, takes \code{coeff = -SR(vtype.coefficient)} --- the minus
sign coming from $Z = \int \ee^{-S}$ in Eq.~\eqref{eq:ta-Z}, so each interaction
factor from $S$ acquires a sign. For a \emph{source}, white noise with a
time-independent amplitude is pulled out; coloured or time-dependent amplitudes
stay inside and are recorded in \code{source\_time\_info}. For an
\emph{interaction vertex}, the constant part (obtained by substituting every
time-dependent symbol to $1$) joins \code{scalar\_parts}, and the residual
\code{td\_part = coeff / const\_part} (tidied with \code{simplify\_rational})
goes into \code{vertex\_time\_factors}. Finally
\code{scalar\_prefactor = reduce(mul, scalar\_parts, SR(1))} multiplies the
pull-out pieces into one product, and \code{is\_stationary} is \code{True} when
there are no vertex time factors and no time-dependent or non-standard noise
amplitudes.

% =====================================================================
\section{The orchestrator: wiring the four phases}\label{sec:ta-orchestrator}

All four phases are wired together --- with disk caching --- by
\code{enumerate\_unique\_diagrams} in \file{pipeline/\_diagrams.py}. The heart of
it, after a cache miss, is four lines
(\file{pipeline/\_diagrams.py:176--186}):

\begin{lstlisting}
_, _, prediagrams, _ = enumerate_prediagrams_all(k=k, ell=ell, verbose=False)
typed = enumerate_all_typed(prediagrams, external_fields, vtypes, stypes,
                            G_ft=G_ft, resp_index=resp_idx, phys_index=phys_idx,
                            parallel=parallel, n_workers=n_workers)
causal, n_disc, _ = filter_causal(typed)
unique, multiplicities = deduplicate_with_multiplicities(causal)
\end{lstlisting}

Read top to bottom this is the whole chapter in miniature: Phase~C produces the
prediagrams; \code{enumerate\_all\_typed} (the cross-prediagram driver of Phase~E,
which we meet by its full name \code{enumerate\_all} in \S\ref{sec:ta-fork})
labels them; \code{filter\_causal} (Phase~F) keeps the DAGs; and
\code{deduplicate\_with\_multiplicities} (Phase~G) collapses isomorphic copies.
The flow of intermediate objects is:

\begin{center}
\begin{tikzpicture}[
    >={Stealth[length=2.4mm]},
    box/.style={draw,rounded corners,fill=codebg,align=center,
                font=\footnotesize,inner sep=4pt,minimum height=9mm},
    lbl/.style={font=\scriptsize\itshape,midway,above},
  ]
  \node[box] (pc)  at (0,0)    {prediagrams\\\scriptsize $(D,G,\text{leaves},\text{internal})$};
  \node[box] (pe)  at (3.9,0)  {typed\\\scriptsize list[TypedDiagram]};
  \node[box] (pf)  at (7.5,0)  {causal\\\scriptsize list[TypedDiagram]};
  \node[box] (pg)  at (11.1,0) {unique\\\scriptsize $+$ multiplicities};
  \draw[->] (pc) -- (pe) node[lbl] {E};
  \draw[->] (pe) -- (pf) node[lbl] {F};
  \draw[->] (pf) -- (pg) node[lbl] {G};
\end{tikzpicture}
\end{center}

Note that Phase~D does not appear: as \S\ref{sec:ta-filter} explained, the loop
enumerator already emits degree-valid prediagrams, so the explicit filter is
skipped on the live path.

Downstream, the integrators (Phase~J) take each unique \code{TypedDiagram} and,
per diagram: read \code{vertex\_assignments} for coefficients (split via
\code{classify\_coefficient\_factors}), read \code{edge\_types} and
\code{propagator\_indices} to place propagators $G_{\mathrm{ft}}[pi, ri]$ on
edges, read \code{external\_legs} to line the leaves up with
\code{external\_fields}, multiply by $\Sym(\Gamma) =
\code{combinatorial\_factor(td)}$, sum over external-leaf mappings while dividing
by \code{external\_wick\_compensation(td)}, and finally perform the time /
frequency / momentum integral.

% =====================================================================
\section{Parallelism and fork-safety}\label{sec:ta-fork}

The cross-prediagram driver of Phase~E, \code{enumerate\_all}
(\file{type\_assignment.py:582--684}), is the only function here with a
performance story worth telling, and it comes with a genuine landmine. Labelling
one prediagram is independent of labelling any other --- the work is
\emph{embarrassingly parallel} across prediagrams --- so \code{enumerate\_all}
offers an optional fork-based multiprocessing path. The serial and parallel paths
produce \textbf{bit-identical} output, in the same order, pinned by
\code{test\_enumerate\_all\_parallel\_matches\_serial}.

\subsection{Why fork, and why it is lethal in a notebook}

The path uses POSIX \code{fork} specifically because Sage \code{DiGraph} objects,
\code{VertexType}/\code{SourceType} instances, and the propagator matrix do not
all pickle cleanly through stdlib \code{pickle} --- so the ``spawn'' start method,
which passes inputs to workers by pickling, would fail. With \code{fork}, the
workers simply \emph{inherit} the parent's memory and already have those objects;
only the \emph{return} value (a list of \code{TypedDiagram}, which does pickle via
its \code{\_\_getstate\_\_}) crosses back. The inputs are stashed in a module
global \code{\_ENUM\_WORKER\_STATE} \emph{before} the fork so the children inherit
them rather than receiving them by pickling.

The catastrophe is platform-specific. On \textbf{macOS}, Apple's runtime forbids
using many Objective-C / Cocoa APIs after a \code{fork()} that is not followed by
an \code{exec()}. Once Cocoa, BLAS, or matplotlib has initialised in the parent
--- which \emph{always} happens inside a Jupyter kernel --- forking and then
touching those libraries in the child can \textbf{hard-crash the kernel and the
entire OS}. (This crashed the project's dev machine more than once.) The
``fix'' people reach for, setting the environment variable
\code{OBJC\_DISABLE\_INITIALIZE\_FORK\_SAFETY=YES}, makes it \emph{worse}: it
converts a controlled abort into a hard crash.

\subsection{The guard}

The shared guard lives in \file{msrjd/fork\_safety.py} and is consulted by every
fork-based path in the repository (Phase~E here, and the temporal Phase~J
batch driver). Its core predicate returns \code{True} only in the precise lethal
situation:

\begin{lstlisting}
def fork_unsafe_in_notebook(start_method='fork'):  # fork_safety.py:27-44
    if start_method != 'fork':
        return False
    if sys.platform != 'darwin':
        return False
    try:
        from IPython import get_ipython
        ip = get_ipython()
    except Exception:
        return False
    return ip is not None and ip.__class__.__name__ == 'ZMQInteractiveShell'
\end{lstlisting}

The Jupyter fingerprint is the class-name check \code{ZMQInteractiveShell} --- the
IPython shell class that \emph{only} a ZMQ/Jupyter kernel uses. Terminal IPython
uses \code{TerminalInteractiveShell}; plain scripts and pytest return \code{None}
from \code{get\_ipython()}. So fork stays available --- and fast --- everywhere it
is safe (scripts run as \code{sage -python run.py}, pytest, Linux, terminal
IPython), and is suppressed only inside a macOS Jupyter kernel.

When the guard fires, \code{enumerate\_all} warns once and degrades to serial
(\file{type\_assignment.py:630--634}):

\begin{lstlisting}
from msrjd.fork_safety import fork_unsafe_in_notebook, warn_fork_fallback_once
if (parallel and len(prediagrams) > 1
        and fork_unsafe_in_notebook(start_method)):
    warn_fork_fallback_once('diagram type-assignment enumeration')
    parallel = False
\end{lstlisting}

\code{warn\_fork\_fallback\_once(where)} (\file{fork\_safety.py:47--61}) emits a
single \code{RuntimeWarning} per call-site label per process, deduped via the
module-level \code{\_WARNED} set, so a notebook user sees the explanation once
rather than on every call.

\begin{gotcha}
The fork path is \emph{off by default} (\code{parallel=False},
\file{type\_assignment.py:584}). And note the subtle reason the bit-identity
tests do \emph{not} certify notebook safety: pytest is not a ZMQ kernel, so
\code{fork\_unsafe\_in\_notebook} returns \code{False} under pytest and those
tests run on the \emph{fork} path. They prove serial and parallel agree
numerically; they say nothing about whether forking a notebook is safe. Only the
guard protects the notebook. Relatedly, \code{\_ENUM\_WORKER\_STATE} is a module
global populated before the fork and is \emph{not} thread-safe --- do not call
\code{enumerate\_all(parallel=True)} from two threads at once
(\file{type\_assignment.py:657--660}).
\end{gotcha}

% =====================================================================
\section{External tools used in this chapter}\label{sec:ta-tools}

This subsystem leans on three external pieces beyond plain Python. None of NumPy,
SciPy, numba, or networkx appears here. Because the manual assumes you know the
physics but not necessarily the tooling, here is each from the ground up.

\paragraph{SageMath and its symbolic ring \code{SR}.} \term{SageMath} (``Sage'')
is a large open-source mathematics system that bundles many specialised libraries
behind one uniform Python API; the project runs inside a Sage-provided Python
(the \code{MSRJD\_diagrams} conda environment). Its symbolic ring \code{SR} is a
ring of computer-algebra expression trees. \code{SR(x)} coerces a number or
expression into a symbolic expression; methods like \code{.is\_zero()},
\code{.simplify\_full()}, \code{.variables()}, \code{.subs(\{...\})}, and
\code{.simplify\_rational()} operate on it. It is imported in all three core
files (\file{type\_assignment.py:19}, \file{causality.py:16},
\file{symmetry.py:46}). Notably, in \code{type\_assignment.py} the
\emph{semantic} \code{SR(...).is\_zero()} path is \emph{deliberately bypassed} in
favour of the structural string test of \S\ref{sec:ta-zeromask}; in
\code{causality.py}, \code{imag\_part(...)} is Sage's symbolic ``take the
imaginary part'' function.

\paragraph{Sage's graph theory, powered by \code{nauty}.} Sage's \code{DiGraph}
is a directed-graph object; crucially its isomorphism, automorphism, and
canonical-form algorithms are powered by \term{nauty} (``No AUTomorphisms,
Yes?''), Brendan McKay's gold-standard C library for graph canonical labelling
and automorphism-group computation. Given a graph --- optionally vertex-coloured
via a \term{partition} --- nauty computes a canonical form (two graphs are
isomorphic iff their canonical forms match) and the automorphism group (all
relabellings mapping the graph to itself). The partition is how you pass a
colouring: nauty may only map a vertex to another vertex in the \emph{same}
partition cell, which is exactly how the code forbids, say, mapping a noise source
onto an interaction vertex. This subsystem calls into nauty at two points only:
\code{D.automorphism\_group(partition=...)} for the symmetry factor and
compensation (\file{symmetry.py:339}), and
\code{D.canonical\_label(partition=..., certificate=True)} for the dedup signature
(\file{symmetry.py:517}).

\paragraph{SymPy's \code{multiset\_permutations}.} \term{SymPy} is a pure-Python
computer-algebra library, independent of Sage. This subsystem uses exactly one
helper from it, \code{multiset\_permutations} (\file{type\_assignment.py:406}),
the engine of the canonical-leg-matching optimisation of \S\ref{sec:ta-canonical}.

\paragraph{Python standard library.} The rest is stdlib:
\code{itertools.permutations} and \code{itertools.product}
(\file{type\_assignment.py:18}) for orderings and Cartesian products;
\code{collections.Counter} and \code{defaultdict} (\file{symmetry.py:41},
\file{:301}) for tallying multiplicities and grouping edges;
\code{functools.reduce}, \code{operator.mul}, and \code{math.factorial}
(\file{symmetry.py:42--44}) for forming products and factorials;
\code{multiprocessing} (imported lazily, \file{type\_assignment.py:648}) for the
optional fork pool; and \code{IPython.get\_ipython}, \code{sys}, \code{warnings}
(\file{fork\_safety.py}) for the notebook detection and the fallback warning.

% =====================================================================
\section{Gotchas, consolidated}\label{sec:ta-gotchas}

A single reference list of the traps in this subsystem; each is documented in
context above and in the code at the cited lines.

\begin{enumerate}
  \item \textbf{Propagator index order is $(pi, ri)$, not $(ri, pi)$.}
        $G_{\mathrm{ft}}$ is physical-row, response-col, so an edge propagator is
        $G_{\mathrm{ft}}[pi, ri]$. The code \emph{stores}
        \code{propagator\_indices = (ri, pi)} but \emph{looks up} the mask/matrix
        with \code{(pi, ri)}. The historical \code{[ri, pi]} bug
        (\file{type\_assignment.py:490--502}) worked for diagonal couplings and
        silently dropped off-diagonal ones.
  \item \textbf{The zero-check is a string test, not \code{is\_zero()}.}
        \code{str(G\_ft[i,j]) == '0'} (\file{type\_assignment.py:165}) catches
        only literal \code{SR(0)} to avoid multi-second symbolic simplification;
        it can over-generate vanishing diagrams but never produces a wrong number.
  \item \textbf{\code{\_leg\_matchings} emits only the canonical leg ordering.}
        Correct only because the integrator multiplies by $\Sym(\Gamma)$
        afterward; pinned by \code{test\_leg\_matchings\_canonical}.
  \item \textbf{Dedup multiplicity is diagnostic only.} Never multiply it back;
        under Path A the weight is entirely in \code{combinatorial\_factor}
        (\file{symmetry.py:556--568}).
  \item \textbf{\code{diagram\_signature} must stay a complete invariant.} A
        shallow heuristic collided non-isomorphic classes and silently dropped
        integrals; the canonical-form signature cannot
        (\file{symmetry.py:484--498}).
  \item \textbf{\code{external\_wick\_compensation} fails loud on a non-divisor.}
        It raises rather than guess (\file{symmetry.py:391--399}), unlike
        \code{combinatorial\_factor}'s defensive floor-division --- a deliberate
        asymmetry.
  \item \textbf{\code{vertex\_role\_signature} is superseded for weighting.} Its
        depth-1 $\prod N!$ over-divided; the exact compensation replaced it. Kept
        but unused on the production weight path.
  \item \textbf{The fork path is off by default and unsafe in notebooks.}
        \code{parallel=False} (\file{type\_assignment.py:584});
        \code{fork\_unsafe\_in\_notebook} forces serial inside a macOS Jupyter
        kernel. The bit-identity tests run under pytest (not a ZMQ kernel) and
        never certified notebook safety.
  \item \textbf{\code{check\_pole\_structure} is latent.} No \code{pole\_vals} are
        passed in the live pipeline, so only the structural DAG check gates
        diagrams. A pole exactly on the real axis is a hard fail; an inconclusive
        symbolic pole returns \code{passed=True} with conditions.
  \item \textbf{Phase D is not on the live path.} The loop enumerator already
        emits degree-valid prediagrams, so \code{enumerate\_unique\_diagrams}
        skips \code{filter\_prediagrams}; it is exercised by tests and looser
        callers.
\end{enumerate}

% =====================================================================
\section{What you just learned}

This chapter took a bare directed graph and turned it into a set of unique,
causally valid, fully labelled Feynman diagrams, each tagged with its symmetry
factor --- the form the integrators consume. Along the way you saw:

\begin{itemize}
  \item the response/physical edge convention and the load-bearing
        $G_{\mathrm{ft}}[pi, ri]$ index order;
  \item Phase~D's cheap degree pre-screen, and why it is absent from the live
        path;
  \item Phase~E as a backtracking constraint solver, with candidate types,
        external-field permutations, leg matchings, early propagator pruning, and
        the canonical-leg-matching optimisation that keeps the intermediate set
        $15$--$20\times$ smaller with bit-identical output;
  \item the structural string-based zero mask and the performance reasoning
        behind it;
  \item Phase~F's DAG causality check (and its dormant analytic sibling);
  \item Phase~G's symmetry factor via a coloured incidence digraph fed to
        \code{nauty}, the complete-invariant deduplication that closed a real
        silent-wrong-answer bug, and the orbit--stabilizer external Wick
        compensation that hands off to the integrator;
  \item the embarrassingly-parallel fork path and the fork-safety guard that
        makes it survivable on a Mac.
\end{itemize}

The next chapter follows the unique diagrams into the caching layer, which is
what lets a cold enumeration like this one be paid for exactly once.

% ---------------------------------------------------------------------
%  Chapter 11 : Caching --- the enumeration and expand caches
% ---------------------------------------------------------------------
\chapter{Caching: The Enumeration and Expand Caches}\label{ch:caching}

Two of the steps on the road from a model spec to a number are
\emph{brutally expensive} and \emph{highly reusable}, and this chapter is
about how \Daedalus{} avoids paying for them twice. The first is expanding
the symbolic action: a multivariate Taylor expansion that, for a four-field,
two-population compartmental theory at \code{taylor\_order=4}, runs
\emph{roughly ninety minutes single-threaded} inside SageMath. The second is
enumerating Feynman diagrams: a four-stage combinatorial pipeline whose cost
grows fast with the external-leg count $k$ and the loop order $\ell$.

The key observation that makes caching possible --- and the organising idea
of the whole chapter --- is that \emph{neither result depends on the numerical
parameter values a user later plugs in}. A higher value of the friction
$\mu$, a different noise strength $D$: these change the numbers that come out
the far end of the pipeline, but they do not change \emph{which} diagrams
exist or \emph{what} the expanded action looks like symbolically. Those
artefacts depend only on the \emph{structure} of the theory --- its field
list, its interaction vertices, the requested $k$ and $\ell$, and which fields
sit on the external legs. So both are computed once, written to disk keyed by
that structure, and reloaded across notebook restarts and across runs that
vary only the numbers.

This chapter documents the caching layer for those two artefacts. It sits
squarely between the model definition and the numerical integration, and it is
built on four files:

\begin{description}
  \item[\file{msrjd/core/cache.py}] the low-level, stage-keyed \code{.sobj}
        store, \code{PipelineCache}, that both higher-level caches build on.
  \item[\file{pipeline/\_expand\_cache.py}] the disk cache for
        \code{FieldTheory.expand()} results --- the expanded action.
  \item[\file{pipeline/\_diagrams.py}] the disk cache wrapping the four-stage
        diagram pipeline --- \code{enumerate\_unique\_diagrams}.
  \item[\file{pipeline/\_precompute.py}] the cheap one-time structural pass,
        \code{precompute}, that warms the durable caches up front.
\end{description}

By the end you will understand exactly which file lands where on disk, why a
cached order-6 expansion satisfies an order-2 request for free, why the
enumeration cache embeds a version string that you \emph{must} bump when you
change the diagram semantics, and which single corner of this subsystem is the
most fragile.

\section{Why cache at all}

\Daedalus{} turns a model spec --- a Python \code{dict} describing an \msrjd{}
field theory --- into the connected cumulants of the fields, computed
perturbatively from Feynman diagrams. Two stages on that road dominate the
cost and are perfectly reusable.

\paragraph{Expanding the action.}
To do perturbation theory you need the action written as a power series in the
\emph{fluctuations} about the mean-field saddle. \Daedalus{} hands the
symbolic \msrjd{} action $S[\tilde\varphi,\varphi]$ to SageMath and runs a
multivariate Taylor expansion to a chosen order. The docstring of
\file{pipeline/\_expand\_cache.py:6--9} is blunt about the cost:

\begin{quote}\itshape
``\code{FieldTheory.expand()} is the dominant cost for non-trivial theories
--- for a 4-field $\times$ 2-population compartmental Bernoulli theory at
\code{taylor\_order=4} it runs $\sim$90 minutes single-threaded inside Sage's
multivariate \code{taylor()}.''
\end{quote}

\paragraph{Enumerating diagrams.}
For each loop order $\ell = 0, 1, \dots, \texttt{max\_ell}$ the pipeline
enumerates every topologically distinct, type-consistent, causal,
\emph{non-isomorphic} Feynman diagram with $k$ external legs. This is a
four-stage combinatorial pipeline --- prediagrams, then typing, then a
causality filter, then deduplication --- whose cost grows fast with $k$ and
$\ell$.

\paragraph{The separation of structure from numbers.}
The thing that makes both cacheable is that they are functions of the theory's
\emph{structure} alone. Concretely, the two artefacts depend on the field
list, the interaction vertices, the requested $k$ and $\ell$, and the
external-field labels --- and on \emph{nothing} numerical. So \Daedalus{}
caches both on disk, keyed by structure, and reloads them across notebook
restarts and across parameter sweeps. This is the single most important fact
in the chapter: \emph{the cache key is structural, never numerical}.

The data path looks like this. \code{TheoryBuilder.build()} produces the model
\code{dict}; \code{precompute(model)} runs a one-time structural pass that
warms the durable caches; and \code{compute\_cumulants(\dots)} consumes them
on its way to the integrator:

\begin{lstlisting}[style=console]
TheoryBuilder.build()  ->  model dict
        |
        v
   precompute(model)  -- one-time structural pass --+
        |   (expand@order2 + propagator + MF check) |  writes
        v                                            v
  compute_cumulants(...)                  saved_theories/<slug>/
        |                                    +-- propagator.sobj
        +-[1] FieldTheory.expand            +-- expand_taylor2.sobj
        |     <- EXPAND CACHE               +-- expand_taylor4.sobj
        |       (pipeline/_expand_cache.py) +-- unique_typed_..._k<k>_l<l>_taylor<N>.sobj
        +-[5] enumerate_unique_diagrams     +-- manifest.json
        |     <- ENUMERATION CACHE
        |       (pipeline/_diagrams.py)
        v
   diagrams -> integration -> cumulants
\end{lstlisting}

The two caches have deliberately \emph{different} keys, and that difference is
the central design decision:

\begin{itemize}
  \item The \term{expand cache} is keyed by \code{taylor\_order} \emph{only}
        --- one file per order. A higher-order expansion is a strict superset
        of a lower one (we prove this below), so a high-order cache serves any
        lower-order request for free.
  \item The \term{enumeration cache} is keyed by the finer tuple
        $(\text{model}, \texttt{taylor\_order}, k, \ell,
        \text{external-fields})$, because the diagram \emph{set} does depend on
        which interaction vertices are in scope, and that in turn depends on
        \code{taylor\_order}.
\end{itemize}

Both caches share one physical store, \code{PipelineCache}, which writes Sage
\code{.sobj} files into a single per-theory directory and maintains a
human-readable \code{manifest.json} index. (As we will see, the expand cache
actually bypasses \code{PipelineCache} and calls Sage's \code{save}/\code{load}
directly --- but it writes into the \emph{same} directory.)

\section{The on-disk layout, and \texttt{.sobj} files from scratch}

Everything for one theory lives in a single directory under
\file{saved\_theories/}. The directory name is the theory's \term{slug}: a
filesystem-safe version of \code{model['name']}. All three caches compute the
same slug the same way, which is why they share one directory. The rule is
(\file{pipeline/\_expand\_cache.py:104--106}):

\begin{lstlisting}
def _slug(model: dict) -> str:
    """Filesystem-safe slug for the cache directory."""
    return re.sub(r'[^A-Za-z0-9]+', '_', model['name']).strip('_').lower()
\end{lstlisting}

The regular expression \code{[\^{}A-Za-z0-9]+} matches any run of characters
that are \emph{not} a letter or digit, and \code{re.sub} replaces each such run
with a single underscore; \code{.strip('\_')} trims leading and trailing
underscores, and \code{.lower()} normalises case. So a model named
\code{'1d KPZ (per-leg gradient vertex)'} becomes the directory
\file{1d\_kpz\_per\_leg\_gradient\_vertex}. A live such directory looks like:

\begin{lstlisting}[style=console]
saved_theories/1d_kpz_per_leg_gradient_vertex/
    expand_taylor2.sobj        propagator.sobj
    expand_taylor4.sobj        manifest.json
    unique_typed_mult_v3_<ext>_k<k>_l<l>_taylor<N>.sobj   (one per (k, l, ext, order))
\end{lstlisting}

Read the filenames left to right. \file{expand\_taylor2.sobj} and
\file{expand\_taylor4.sobj} are two different orders of the same expanded
action --- siblings in one directory. \file{propagator.sobj} is the symbolic
propagator (built by a different subsystem, but landing in the same place
because it shares the slug). \file{manifest.json} is the human-readable index
of the \code{PipelineCache} entries. And the long
\file{unique\_typed\_mult\_v3\_\dots} files are the per-$(k,\ell)$ enumeration
results; each filename carries the external-field tag, the version stamp, the
$k$ and $\ell$, and the taylor order. We dissect every piece of that name in
Section~\ref{sec:enumcache}.

\begin{defn}
A \term{\code{.sobj}} (``Sage object'') file is SageMath's serialization
format: a gzip-compressed pickle written by \code{sage.all.save(obj, path)}
and read back by \code{sage.all.load(path)}. It is the on-disk format for
\emph{every} cached artefact in this subsystem.
\end{defn}

\subsection{Why Sage's serializer, and not plain \texttt{pickle}}

The objects this subsystem writes to disk are \emph{Sage objects}: the
symbolic action, the polynomial-ring elements stored in the expanded action,
the propagator matrix, and the diagram objects (which wrap Sage graphs). Plain
Python \code{pickle} does not round-trip these reliably, because Sage objects
carry references to \emph{parent structures} --- rings, fields --- whose
identity matters. After an unpickle, the reconstructed parent is a
\emph{different} Python object than a freshly built one, even when the two have
identical generators, and arithmetic that mixes elements from the two parents
can break. SageMath's own \code{save}/\code{load} understand this and handle it
correctly, so \Daedalus{} uses them everywhere it touches disk.

\begin{note}
SageMath is a large open-source mathematics system built on top of Python. It
bundles symbolic algebra (its symbolic ring \code{SR}), exact and
arbitrary-precision arithmetic, polynomial rings, and graph theory under one
Python API. When you import from \code{sage.all} you get its global namespace,
including the two object-serialization helpers used throughout this chapter:
\code{save} and \code{load}. If you have never used Sage, the only two facts
you need for this chapter are: (i) \code{sage.all.save(obj, path)} pickles
\code{obj} to \code{path.sobj}, and (ii) it \emph{appends} the \code{.sobj}
extension if you do not supply one --- which is why every call here strips it
off first.
\end{note}

There is one subtlety in how \code{save} is called that recurs in every file.
Because Sage's \code{save} appends \code{.sobj}, the caller passes a path with
the extension already \emph{removed}, so that auto-append produces exactly one
\code{.sobj} and not \code{....sobj.sobj}:

\begin{lstlisting}
sage_save(data, path.removesuffix('.sobj'))  # sage_save adds .sobj
\end{lstlisting}

That exact line appears at \file{msrjd/core/cache.py:108},
\file{pipeline/\_expand\_cache.py:345}, and inside \code{save\_expand}. The
matching \code{sage\_load(path)} accepts \emph{either} form (with or without
the extension), so the load side passes the full \code{.sobj} path.

\section{\texttt{PipelineCache}: the low-level store}

\file{msrjd/core/cache.py} is the foundation both higher-level caches sit on.
It is a stage-based disk cache: one directory per model, with files keyed by
the triple \code{(stage\_name, k, loop\_order)}. The class is small enough to
walk method by method.

\subsection{Construction and the stage key}

\code{PipelineCache.\_\_init\_\_} (\file{cache.py:57}) takes a \code{root}
directory and stores \code{self.root = os.path.expanduser(root)} --- the
\code{expanduser} call turns a leading \code{\~{}} into the user's home
directory. It does \emph{not} create the directory yet; that happens lazily on
the first \code{save}.

A cache file's name is built from the stage name plus the optional $k$ and
$\ell$. Two helpers do this (\file{cache.py:62--83}):

\begin{lstlisting}
@staticmethod
def _to_int(val):
    """Coerce SageMath Integer / numpy int to plain Python int."""
    if val is None:
        return None
    return int(val)

@classmethod
def _stage_key(cls, stage, k=None, loop_order=None):
    """Build a filename stem from stage name and optional (k, l)."""
    k = cls._to_int(k)
    loop_order = cls._to_int(loop_order)
    parts = [stage]
    if k is not None:
        parts.append(f'k{k}')
    if loop_order is not None:
        parts.append(f'l{loop_order}')
    return '_'.join(parts)
\end{lstlisting}

\code{\_to\_int} earns its place because $k$ and $\ell$ frequently arrive as
SageMath \code{Integer} objects, not plain Python \code{int}s. Coercing them
keeps filenames and manifest JSON clean --- you get \code{k2}, not the
\code{repr} of a Sage \code{Integer}. So \code{('unique\_typed', 2, 1)} becomes
the stem \code{'unique\_typed\_k2\_l1'}, and \code{\_sobj\_path}
(\file{cache.py:81}) joins that to the root with a \code{.sobj} suffix.

\subsection{Save, load, exists, and the memoise helper}

The core API is four methods. \code{exists} (\file{cache.py:87}) is a one-liner
--- \code{os.path.isfile(self.\_sobj\_path(\dots))} --- and it is the
\emph{filesystem} that is consulted, not the manifest. \code{save}
(\file{cache.py:91}) creates the directory if needed, writes via Sage, and
then updates the manifest:

\begin{lstlisting}
def save(self, stage, data, k=None, loop_order=None):
    os.makedirs(self.root, exist_ok=True)
    path = self._sobj_path(stage, k, loop_order)
    sage_save(data, path.removesuffix('.sobj'))  # sage_save adds .sobj
    self._update_manifest(stage, k, loop_order)
\end{lstlisting}

The \code{exist\_ok=True} flag on \code{os.makedirs} means ``do not raise if
the directory already exists'' --- the idiomatic way to ensure a directory is
present without a prior existence check. \code{load} (\file{cache.py:111})
reads the file back through \code{sage\_load}, and \emph{raises}
\code{FileNotFoundError} with a descriptive message if the file is absent. That
exception is the documented contract, but --- as we will see --- the
enumeration cache calls \code{exists} first and wraps \code{load} in a
\code{try}/\code{except}, so it never actually relies on the raise.

\code{get\_or\_compute} (\file{cache.py:125}) is the classic memoise pattern:

\begin{lstlisting}
def get_or_compute(self, stage, compute_fn, k=None, loop_order=None):
    if self.exists(stage, k, loop_order):
        return self.load(stage, k, loop_order)
    data = compute_fn()
    self.save(stage, data, k, loop_order)
    return data
\end{lstlisting}

\begin{note}
The enumeration cache does \emph{not} use \code{get\_or\_compute}. It does its
own manual \code{exists}/\code{load}/\code{save} dance so that it can print
per-stage progress messages and degrade gracefully on a load failure. The
memoise helper is here for callers who want the simple path; the production
diagram cache wants the verbose, fault-tolerant one.
\end{note}

\subsection{The manifest --- a best-effort index}

\code{PipelineCache} maintains a \file{manifest.json} alongside the
\code{.sobj} files: a human-readable index of what has been cached. Its shape
is \code{\{'entries': [\{key, stage, k, loop\_order, saved\_at\}], 'created':
iso, 'updated': iso\}}. \code{\_update\_manifest} (\file{cache.py:168}) loads
the current manifest, removes any stale entry with the same key, appends a
fresh one with an ISO-8601 timestamp, and writes it back with
\code{json.dump(\dots, indent=2)} for readability.

The crucial design point is that the manifest is \emph{not} correctness-bearing.
\code{\_load\_manifest} (\file{cache.py:153}) is deliberately robust: if the
file is missing \emph{or corrupt} it silently returns a fresh empty structure
rather than raising:

\begin{lstlisting}
try:
    with open(mp) as f:
        manifest = json.load(f)
    if not isinstance(manifest, dict):
        raise ValueError
    return manifest
except (json.JSONDecodeError, ValueError):
    # Corrupt manifest -- start fresh.
    return {'entries': [], 'created': datetime.now().isoformat()}
\end{lstlisting}

The \code{.sobj} files are the source of truth; the manifest is a convenience.
Because \code{exists} checks the filesystem and not the manifest, a missing or
truncated manifest entry can never cause a spurious cache miss.

Finally, \code{clear} (\file{cache.py:194}) removes cached data: with all three
arguments \code{None} it \code{shutil.rmtree}s the whole directory; otherwise
it deletes just the one matching \code{.sobj} and prunes its manifest entry.
And \code{\_\_repr\_\_} reports \code{PipelineCache('<root>', <N> entries)}.

\section{The expand cache I: bigrades and the superset theorem}

We now turn to the larger of the two caches, the one that stores the
ninety-minute artefact. To understand what it stores you first need to
understand the object \code{FieldTheory.expand()} produces: the
bigrade-classified action.

\subsection{The \msrjd{} action and its Taylor expansion}

\Daedalus{} works in the Martin--Siggia--Rose--Janssen--De~Dominicis
representation of a stochastic dynamical system. A Langevin equation
$\partial_t\varphi = F[\varphi] + \text{noise}$ becomes a path integral with an
action over a \emph{physical} field $\varphi$ and a \emph{response} (or
``tilde'') field $\tilde\varphi$:
\begin{equation}
  S[\tilde\varphi,\varphi]
  = \int \dd t\,
    \Big\{
      \tilde\varphi\,(\Dt\varphi - F[\varphi])
      - \tfrac12\,\tilde\varphi\,B\,\tilde\varphi + \cdots
    \Big\}.
  \label{eq:msrjd-action}
\end{equation}
Cumulants are read off from a perturbative expansion of $\ee^{-S}$ around the
mean-field (saddle) point. To do that expansion you need the action written in
powers of the fluctuations about the saddle. Write
$\varphi = \varphi^{*} + \delta\varphi$ and
$\tilde\varphi = \tilde\varphi^{*} + \delta\tilde\varphi$, and Taylor-expand $S$
as a multivariate series in $(\delta\tilde\varphi, \delta\varphi)$.

The central move is then to \emph{classify each monomial by its bigrade}.

\begin{defn}
The \term{bigrade} of an action monomial is the pair
\[
  (n_{\tilde{}},\, n_{\mathrm{phys}})
  = (\#\,\text{response-field factors},\ \#\,\text{physical-field factors}).
\]
Its total degree is $n_{\tilde{}} + n_{\mathrm{phys}}$.
\end{defn}

The expanded action is stored as a dictionary, \code{ft.\_by\_tp}, that maps
each bigrade to the \emph{sum} of all action monomials of that bigrade:
\[
  \texttt{\_by\_tp}[(n_{\tilde{}},\, n_{\mathrm{phys}})]
  \;=\; \sum (\text{action monomials of that bigrade}),
\]
each value being a live Sage polynomial-ring element. A Taylor truncation at
order $N$ keeps exactly the bigrades with
$n_{\tilde{}} + n_{\mathrm{phys}} \le N$.

The bigrades carry direct physical meaning, and knowing it tells you why
order~2 is special (Section~\ref{sec:precompute}):

\begin{itemize}
  \item \textbf{$(0,0)$, $(1,0)$, $(0,1)$} --- the \term{mean-field / saddle
        sectors}. These are order zero in the fluctuations: the classical
        equation of motion. The saddle condition is precisely that they vanish
        at $(\varphi^{*}, \tilde\varphi^{*})$, which is what
        \code{sanity\_check} verifies.
  \item \textbf{$(1,1)$} --- the \term{propagator kernel}, the bilinear sector.
        The free (Gaussian) propagator $\Gret$ is the matrix inverse of the
        $(1,1)$ coefficients. The propagator builder reads exactly this sector.
  \item \textbf{Total degree $\ge 3$} --- the \term{interaction vertices}. A
        monomial of bigrade $(a,b)$ with $a+b=d$ is a degree-$d$ vertex with
        $a$ response legs and $b$ physical legs. These are the Feynman rules
        the diagram pipeline consumes.
\end{itemize}

\subsection{The superset theorem and why downgrade is exact}

Here is the mathematical fact that justifies keying the expand cache by
\code{taylor\_order} alone. The docstring states it at
\file{pipeline/\_expand\_cache.py:14--18}:

\begin{quote}\itshape
``\code{\_by\_tp} at \code{taylor\_order=N} is a STRICT SUPERSET of
\code{\_by\_tp} at any \code{taylor\_order=M} with \code{M <= N}: the bigrade
entries at total degree \code{<= M} are byte-identical (the Taylor truncation
only adds higher-order entries, never modifies lower ones).''
\end{quote}

Stated as a theorem:

\begin{defn}
\textbf{(Superset theorem.)} For $M \le N$,
\[
  \texttt{by\_tp}@N \;\supseteq\; \texttt{by\_tp}@M,
  \qquad
  \texttt{by\_tp}@N\big|_{\{(a,b)\,:\,a+b\le M\}}
  \;=\; \texttt{by\_tp}@M .
\]
\end{defn}

The reason is elementary: Taylor truncation at a higher order adds
higher-degree terms but cannot alter the lower-degree ones, because the
order-$d$ Taylor coefficient is fixed by the $d$-th derivative at the saddle and
is independent of where you truncate. Each bigrade entry is therefore
\emph{byte-identical} between the two expansions.

The practical consequence is the \term{downgrade}: to satisfy an order-$M$
request from a cached order-$N$ file, you simply \emph{discard} every bigrade of
total degree above $M$. This is exact --- it is not a re-expansion and not an
approximation. The implementation is one comprehension
(\file{pipeline/\_expand\_cache.py:181--190}):

\begin{lstlisting}
def downgrade_by_tp_dict(by_tp_dict: dict, target_order: int) -> dict:
    """Keep only bigrades whose total degree (a + b) is <= target_order."""
    return {key: poly_dict
            for key, poly_dict in by_tp_dict.items()
            if sum(key) <= target_order}
\end{lstlisting}

Note \code{sum(key)}: \code{key} is the bigrade tuple $(a,b)$, so \code{sum}
gives the total degree $a+b$. This is why a cached order-6 file fully satisfies
an order-2 request, and it is why the cache picks the \emph{smallest} cached
order that suffices: every bigrade kept must later be coerced back into the
freshly built ring, so loading the minimal sufficient file is the cheapest. The
chooser is \code{find\_best\_cached\_order}
(\file{pipeline/\_expand\_cache.py:136--146}):

\begin{lstlisting}
def find_best_cached_order(model, target_order, cache_dir_root='saved_theories'):
    """Return the smallest cached order >= target_order, or None."""
    cached = list_cached_orders(model, cache_dir_root)
    candidates = [o for o in cached if o >= target_order]
    return min(candidates) if candidates else None
\end{lstlisting}

\code{list\_cached\_orders} (\file{pipeline/\_expand\_cache.py:121--133}) finds
all available orders by listing the directory and matching each filename
against \code{re.fullmatch(r'expand\_taylor(\textbackslash d+)\textbackslash
.sobj', fname)} --- the parenthesised \code{(\textbackslash d+)} captures the
integer order, which \code{int(m.group(1))} extracts. It returns \code{[]} if
the directory does not exist.

\begin{gotcha}
Downgrade is exact \emph{only because of} the superset theorem, which in turn
relies on Taylor truncation never modifying lower-order entries. If a future
change made the expansion non-truncating --- a resummation, say --- then
\code{downgrade\_by\_tp\_dict} would silently be \emph{wrong}: it merely drops
high-degree bigrades; it does not recompute the low-degree ones.
\end{gotcha}

\section{The expand cache II: dict-form serialization}

We said Sage objects carry parent-identity baggage across pickle. Polynomial
objects are the worst offenders: an unpickled polynomial ring is a different
Python object than the freshly built one even with identical generators, so a
polynomial unpickled against the old ring will not cleanly combine with
elements of the new ring. The expand cache sidesteps this entirely by
\emph{not storing polynomial objects at all}. Instead it stores their
\term{dict form}.

\begin{defn}
The \term{dict form} of a polynomial-ring element is what Sage's
\code{poly.dict()} method returns: a mapping
\code{\{exponent\_tuple: coefficient\}}. The exponents are plain integers and
the coefficients are \code{SR} (symbolic-ring) elements, both of which survive
pickle and unpickle cleanly.
\end{defn}

The conversions are a handful of one-liners
(\file{pipeline/\_expand\_cache.py:152--175}):

\begin{lstlisting}
def _poly_to_dict_form(poly) -> dict:
    """Returns {exp_tuple: SR_coeff}."""
    return dict(poly.dict())

def _dict_form_to_poly(R, exp_to_coeff: dict):
    """Build a polynomial in ring R from its exp_tuple -> coeff dict."""
    return R(exp_to_coeff) if exp_to_coeff else R.zero()
\end{lstlisting}

The reconstruction calls \code{R(exp\_to\_coeff)} --- the ring's \emph{call
operator}, which builds a polynomial from such a dict. The all-important detail
is that \code{R} here is the \emph{freshly built} ring on load, so the
reconstructed polynomial lives in the live ring and the parent-identity problem
never arises. The whole \code{\_by\_tp} dictionary is mapped through these in
both directions by \code{\_by\_tp\_to\_dict\_form} (save side) and
\code{\_by\_tp\_from\_dict\_form(R, \dots)} (load side, with the rehydrated
ring).

\begin{note}
Because the bundle's values are stored as plain Python dicts of integer
exponents and \code{SR} coefficients, most of the bundle is \emph{pickle-trivial}
--- the only piece where Sage's serializer truly earns its keep is the \code{SR}
coefficients themselves. The dict form thus does most of the
parent-identity work for free, and leaves only the symbolic coefficients to
Sage.
\end{note}

\subsection{The bundle written to disk}

\code{save\_expand} (\file{pipeline/\_expand\_cache.py:322--348}) assembles the
on-disk bundle. It is a plain Python dict:

\begin{lstlisting}
bundle = {
    'by_tp':           _by_tp_to_dict_form(ft._by_tp),
    'S_raw_dict':      (_poly_to_dict_form(ft._S_raw)
                        if ft._S_raw is not None else {}),
    'mf_sector_raw':   _by_tp_to_dict_form(
                          getattr(ft, '_mf_sector_raw', {}) or {}),
    'ring_var_names':  _get_ring_var_names(ft),
    'taylor_order':    int(target),
    'n_tilde':         int(ft._n_tilde),
    'vertex_signature': vertex_form_factor_signature(ft._ns),
    'cache_version':   2,
}
sage_save(bundle, path.removesuffix('.sobj'))
\end{lstlisting}

The fields are summarised in Table~\ref{tab:bundle}.

\begin{table}[t]
\centering\small
\begin{tabular}{@{}lll@{}}
\toprule
key & type & meaning \\
\midrule
\code{'by\_tp'}            & \code{\{(a,b): \{exp: coeff\}\}} & the bigrade-classified action, dict form \\
\code{'S\_raw\_dict'}      & \code{\{exp: coeff\}}            & dict form of the full polynomial (written, \emph{not read on load}) \\
\code{'mf\_sector\_raw'}   & \code{\{(a,b): \{exp: coeff\}\}} & the pre-zero $(0,0)/(1,0)/(0,1)$ sectors (diagnostic) \\
\code{'ring\_var\_names'}  & \code{list[str]}                 & generator names; structural sanity check \\
\code{'taylor\_order'}     & \code{int}                       & the order this was computed at (may exceed the request) \\
\code{'n\_tilde'}          & \code{int}                       & number of response variables; MF-sector filter check \\
\code{'vertex\_signature'} & \code{str} or \code{None}        & operator-IR form-factor fingerprint \\
\code{'cache\_version'}    & \code{int} (=2)                  & bundle-format version \\
\bottomrule
\end{tabular}
\caption{The expand bundle stored in \file{expand\_taylor<N>.sobj}, produced by
\code{save\_expand} and consumed by \code{load\_expand}.}
\label{tab:bundle}
\end{table}

\begin{gotcha}
\code{'S\_raw\_dict'} is written by \code{save\_expand} but \emph{never read by}
\code{load\_expand}. As we will see, the load path reconstructs \code{\_S\_raw}
by summing the (possibly filtered) \code{by\_tp} polynomials instead, which is
the correct value for the requested order. The stored field is therefore
dead-on-load --- harmless, but redundant.
\end{gotcha}

\section{The expand cache III: hidden state and its reconstruction}
\label{sec:hidden-state}

This is the single most fragile corner of the subsystem, and it is worth slowing
down for. The problem is stated plainly in the module docstring
(\file{pipeline/\_expand\_cache.py:67--73}): the on-disk slug is only
\code{model['name']} plus the taylor order, and that key \emph{under-determines}
the theory. Two pieces of state that a fresh \code{expand()} produces as
\emph{side effects} are not captured by \code{\_by\_tp}, hence not by the
bundle, and must be reconstructed when the cache is loaded.

\subsection{The operator-IR form-factor table}

Spatial theories with derivative interaction vertices --- Model~B, KPZ, Burgers
--- represent those vertices through an \term{operator IR} (intermediate
representation) of binding nodes for the Laplacian $\Lap$, the time derivative
$\Dt$, and the spatial gradient $\partial_x$. The per-vertex \emph{form-factor
table} that records which momentum form factor each interaction node carries is
populated as a \emph{side effect of evaluating the action lambda}: it is built
inside \code{model['action'](ns)}, which runs during
\code{FieldTheory.expand()}, \emph{not} during \code{\_build\_namespace}. A
cache load deliberately skips \code{expand()} --- that is the whole point --- so
the loaded namespace would carry \emph{no} table, and every derivative vertex
would silently collapse to form factor $1$. The docstring at
\file{pipeline/\_expand\_cache.py:263--264} spells out the failure mode: a
Model-B$\oplus$KPZ one-loop variance would ``load as the bare $\tilde\varphi
\varphi^2$ number instead of the correct composite\,+\,per-leg one.''

The cure is to re-run the action lambda on load. It is cheap (it builds the
symbolic action and lowers the IR --- no Taylor expansion) and reproduces the
table bit-identically. The helper is \code{\_reconstruct\_operator\_ir\_table}
(\file{pipeline/\_expand\_cache.py:253--282}):

\begin{lstlisting}
ns = getattr(ft, '_ns', None)
if ns is None or not getattr(ns, '_operator_ir', False):
    return
if hasattr(ns, '_operator_ir_vertex_terms'):
    return
try:
    ft.model['action'](ns)
except Exception:
    # Leave the table absent; the signature check below (or a
    # downstream fresh expand) catches the resulting inconsistency.
    pass
\end{lstlisting}

It is \code{hasattr}-guarded (a no-op when the table is already present), and it
returns immediately for theories that do not use the operator IR. On an
exception it swallows the error (\code{pass}) and leaves the table absent,
trusting the signature check to catch the inconsistency.

\subsection{The form-factor signature}

Because the slug cannot encode this per-vertex state, the bundle stores a
\emph{fingerprint} of it, and the load path \emph{validates} that fingerprint.
This is \code{vertex\_form\_factor\_signature}
(\file{pipeline/\_expand\_cache.py:208--250}). For a theory that does not use
the operator IR it returns \code{None} --- so every pre-existing plain or
temporal theory keeps a \code{None}-versus-\code{None} match and loads
unchanged. For an operator-IR theory it returns
\code{'operator\_ir:'} followed by the first sixteen hex digits of a SHA-256
hash, built from each vertex term's four identifying pieces: its \code{mode},
its physical-leg count \code{n\_phys}, its operator \code{chain}, and its
coupling-mixing \code{weight}.

\begin{lstlisting}
canon: list = []
for t in (terms or []):
    n_phys = t.get('n_phys')
    canon.append([
        str(t.get('mode')),
        int(n_phys) if n_phys is not None else None,
        _canon_chain(t.get('chain')),
        str(t.get('weight')),
    ])
# Order-independent: the table-build order is not contractual.
canon.sort(key=lambda e: json.dumps(e, sort_keys=True))
blob = json.dumps(canon, sort_keys=True)
return 'operator_ir:' + hashlib.sha256(blob.encode('utf-8')).hexdigest()[:16]
\end{lstlisting}

Three implementation choices deserve a word. First, every piece is stringified
(via \code{\_canon\_chain} for the operator chain, \code{str()} for mode and
weight) so it round-trips identically through pickle and across Sage sessions.
Second, the list of terms is \emph{sorted} before hashing --- using
\code{json.dumps(e, sort\_keys=True)} as the sort key --- because the
table-build order is not contractual, so the signature must be
order-independent. Third, \code{hashlib} is the Python standard-library hashing
module; \code{sha256(blob).hexdigest()[:16]} produces a short, stable
fingerprint. Adding the KPZ $(\partial_x\varphi)^2$ vertex to a Model-B-only
theory, dropping it, or re-weighting the mix all change this signature, so a
stale on-disk bundle whose signature does not match the freshly built model is
\emph{rejected} by \code{load\_expand}.

\begin{note}
The docstring at \file{pipeline/\_expand\_cache.py:229--233} records a stability
guarantee: the \code{weight} is serialised via \code{str()} of an
already-\code{simplify\_full}'d \code{SR} expression. A purely cosmetic
re-ordering of that expression across Sage versions would, at worst, force one
extra fresh \code{expand()} (and then a re-save with the new signature). It can
never cause a \emph{wrong} load --- the failure mode is conservative.
\end{note}

\subsection{The structural-integrity checks in \texttt{load\_expand}}

\code{load\_expand} (\file{pipeline/\_expand\_cache.py:351--482}) ties it all
together. Its contract is strict: it returns \code{True} on success (after which
\code{ft} is equivalent to a fresh \code{expand()} at \code{target\_order}), and
\code{False} on \emph{any} problem --- no usable cache, a load exception, or a
failed sanity check --- leaving \code{ft} unchanged so the caller can fall back
to a fresh expand. It runs a sequence of guards:

\begin{enumerate}
  \item \textbf{Discover the file.} If \code{cached\_order} was not supplied,
        re-discover it via \code{find\_best\_cached\_order}; if still
        \code{None}, return \code{False}. Then \code{sage\_load} the bundle
        inside a \code{try}/\code{except} (\file{:379--384}).
  \item \textbf{Ring-name check} (\file{:388--394}). The freshly built ring's
        generator names must equal \code{bundle['ring\_var\_names']},
        \emph{order included}. A mismatch means a model edit changed the field
        list and invalidated the cache; return \code{False}.
  \item \textbf{\code{n\_tilde} check} (\file{:396--399}).
        \code{bundle['n\_tilde']} must match \code{ft.\_n\_tilde}.
  \item \textbf{Form-factor signature check} (\file{:411--418}). Reconstruct the
        live table, then require \code{vertex\_form\_factor\_signature(ft.\_ns)}
        to equal \code{bundle['vertex\_signature']}. This catches both a changed
        model that re-used the slug \emph{and} any pre-signature bundle for a
        derivative-vertex theory.
  \item \textbf{Downgrade filter} (\file{:421--423}). If
        \code{cached\_order > target\_order}, filter the bigrades via
        \code{downgrade\_by\_tp\_dict}.
  \item \textbf{Rehydrate} (\file{:425--435}). Set
        \code{ft.\_by\_tp = \_by\_tp\_from\_dict\_form(ft.\_R, by\_tp\_dict)};
        rebuild \code{\_S\_raw} as the \emph{sum} of the (filtered) \code{by\_tp}
        polynomials; and rebuild \code{\_mf\_sector\_raw} from its dict form.
\end{enumerate}

There is one more reconstruction, and missing it would be catastrophic for
colored-noise theories. A fresh \code{expand()} calls
\code{\_build\_cumulant\_action}, which has two effects: it bakes a symbolic
noise-cumulant term into the action (captured by the cached \code{\_by\_tp}),
\emph{and} it populates \code{ns.\_cumulant\_kernels} with kernel-function
closures that the downstream type extractor reads to promote a plain noise
source into a colored one. Those closures do not round-trip through pickle, so
the load path re-executes \code{\_build\_cumulant\_action} purely for that side
effect (\file{:464--472}):

\begin{lstlisting}
try:
    from msrjd.core.field_theory import _build_cumulant_action
    _build_cumulant_action(ft._ns, model)
except Exception as e:
    if verbose:
        print(f'[expand-cache] _cumulant_kernels rebuild failed ...')
    return False
\end{lstlisting}

A failure here returns \code{False} --- forcing a fresh expand --- because a
missing \code{\_cumulant\_kernels} would silently collapse every colored-noise
diagram to zero. This, with the form-factor signature, is the protection that
keeps the slug's under-determination from producing wrong numbers.

\code{load\_expand} has one precondition the caller must honour:
\code{ft.\_ns}, \code{ft.\_R}, and \code{ft.\_n\_tilde} must already be
populated, because the dict-form rehydration needs a live ring to coerce into.
The helper \code{prepare\_for\_load} (\file{pipeline/\_expand\_cache.py:300--319})
does exactly that --- it runs \code{ft.\_build\_namespace()} (but \emph{not}
\code{expand()}) to set the namespace, ring, and \code{n\_tilde}, then calls
\code{\_reconstruct\_operator\_ir\_table}. Both production callers,
\code{precompute} and \code{compute\_cumulants}, call \code{prepare\_for\_load}
immediately before \code{load\_expand}.

\begin{gotcha}
The expand cache's on-disk slug under-determines the bundle. The filename is
only \code{model['name']} plus the taylor order; it captures neither the
operator-IR form-factor table nor the \code{\_cumulant\_kernels} closures. Both
are reconstructed on load, and the form-factor one is validated by the
signature check. If either reconstruction silently no-ops, a derivative-vertex
or colored-noise theory loads a \emph{wrong} (collapsed) result. Three separate
docstring blocks in \file{\_expand\_cache.py} are devoted to this corner ---
that is how fragile it is.
\end{gotcha}

\section{\texttt{precompute}: the one-time structural pass}
\label{sec:precompute}

Recall that bigrade order~2 already contains the $(0,0)/(1,0)/(0,1)$ mean-field
sectors \emph{and} the $(1,1)$ propagator kernel --- everything the structural
validation and propagator construction need, regardless of how high the
downstream cumulant calculation reaches. \code{precompute}
(\file{pipeline/\_precompute.py:55}) exploits this: it is a cheap (a few
seconds for any theory) one-time pass that expands at order~2, runs the sanity
check, solves the mean-field saddle, and builds the propagator --- writing
\file{expand\_taylor2.sobj} and \file{propagator.sobj} so that all later runs
get the propagator and MF check for free.

\code{precompute(model, *, force=False, verbose=True)} returns a status
\code{dict} and proceeds in four stages:

\begin{enumerate}
  \item \textbf{Expand at order~2} (\file{\_precompute.py:106--139}). Build
        \code{ft = FieldTheory(model, taylor\_order=2)}. Unless \code{force},
        try the expand cache --- \code{find\_best\_cached\_order(model, 2)}, and
        on a hit \code{prepare\_for\_load} then \code{load\_expand}. On a miss,
        \code{ft.expand()} fresh (catching exceptions into an
        \code{'EXPAND\_RAISED'} status) and \code{save\_expand} (a save failure
        is non-fatal).
  \item \textbf{Sanity check} (\file{\_precompute.py:142--152}).
        \code{ft.sanity\_check()} confirms the MF sector vanishes at the saddle;
        sets \code{mf\_check} to \code{'PASS'} or \code{'FAIL'} and returns
        early on failure.
  \item \textbf{Solve the mean field} (\file{\_precompute.py:155--164}). Build
        the parameter defaults and solve the saddle. Non-fatal on failure --- the
        propagator can still be built.
  \item \textbf{Build the propagator} (\file{\_precompute.py:167--176}).
        \code{build\_propagator(ft, model, use\_cache=True, force=force)} writes
        \file{propagator.sobj} and sets \code{propagator\_built}.
\end{enumerate}

The status dict carries \code{mf\_check} (\code{'PASS'}/\code{'FAIL'}/error
string), \code{sanity\_ok}, \code{mf\_values}, \code{taylor\_order} (always
\code{2}, since this is the pre-compute order), \code{cache\_dir},
\code{propagator\_built}, \code{wall\_seconds}, and a \code{log} list of the
status messages.

\begin{note}
The imports of \code{FieldTheory}, \code{build\_propagator}, and
\code{\_expand\_cache} sit \emph{inside} \code{precompute}
(\file{\_precompute.py:91--100}), not at module scope. This keeps importing the
module itself cheap, and lets a missing heavy dependency degrade to an
\code{'IMPORT\_FAILED'} status string rather than blowing up at import time.
\end{note}

The two internal helpers are small. \code{\_build\_fundamental\_defaults}
(\file{\_precompute.py:186}) assembles a parameter dict from each model
parameter's \code{default=} declaration, \emph{skipping} parameters marked
\code{mean\_field} (those are saddles, solved not configured) and any whose
default is absent or \code{None}. \code{\_solve\_mf\_at\_saddle}
(\file{\_precompute.py:200}) runs the DAE solver if the model declares
\code{equations}, otherwise the legacy iteration solver, and returns the saddle
values.

\section{The enumeration cache}
\label{sec:enumcache}

The second cache wraps the four-stage diagram pipeline. It lives in
\file{pipeline/\_diagrams.py} and exposes one headline function,
\code{enumerate\_unique\_diagrams}.

\subsection{The four stages it wraps}

For each loop order $\ell$, the pipeline produces the list of
\emph{non-isomorphic typed causal diagrams} with $k$ external legs, plus a
per-diagram multiplicity. The four stages (imported at
\file{\_diagrams.py:28--34} and run at \file{\_diagrams.py:176--186}):

\begin{enumerate}
  \item \textbf{Prediagrams.} \code{enumerate\_prediagrams\_all(k=k, ell=ell)}
        generates all undirected topologies --- trees, then topologies, then
        prediagrams --- with $k$ external legs and $\ell$ loops, \emph{before}
        any field assignment. (Only the prediagrams of its four-tuple return are
        kept.)
  \item \textbf{Typing.} \code{enumerate\_all\_typed} assigns each vertex an
        interaction type from the theory's vertex list (\code{vtypes}) or a
        source type (\code{stypes}), and each internal edge a propagator type
        consistent with the zero/nonzero pattern of the propagator matrix
        \code{G\_ft}. Many prediagrams admit several typings; many admit none.
  \item \textbf{Causality.} \code{filter\_causal} drops diagrams that violate
        the retarded (It\^o/causal) structure of the \msrjd{} propagator --- for
        instance an acausal response-to-response loop. It returns
        \code{(causal, n\_discarded, \_)}.
  \item \textbf{Deduplication.} \code{deduplicate\_with\_multiplicities} merges
        diagrams that are isomorphic as coloured graphs, returning the surviving
        representatives and the size of each equivalence class.
\end{enumerate}

\subsection{The equivalence relation, and what the multiplicity is for}

The deduplication step is where the genuinely subtle physics lives, and it pays
to be precise. Two typed diagrams are treated as duplicates exactly when they
have the same \term{diagram signature}: the canonical form of their coloured
incidence digraph, with leaves coloured by field only
(\file{msrjd/diagrams/symmetry.py:467} onward, called with
\code{fix\_external=False}). The canonical form is computed by Sage's
\code{DiGraph.canonical\_label(\dots, certificate=True)}, which delegates to
\code{nauty}.

\begin{note}
\code{nauty} (``No AUTomorphisms, Yes?'') is a C library by Brendan McKay for
computing graph automorphism groups and canonical labellings. Given a coloured
graph it returns a canonical form such that two graphs have the \emph{same}
canonical form if and only if they are isomorphic. SageMath embeds \code{nauty}
and exposes it through \code{Graph}/\code{DiGraph} methods like
\code{canonical\_label}. This caching subsystem never imports \code{nauty}
directly --- it inherits the dependency transitively through the
\code{deduplicate\_with\_multiplicities} call.
\end{note}

Now the subtle part. The \term{symmetry factor} $\Sym(\Gamma)$ of a diagram ---
its Feynman-rule combinatorial weight --- is carried entirely by a separate
function, \code{combinatorial\_factor}, as the orbit--stabiliser count
\[
  \Sym(\Gamma)
  \;=\; \frac{\prod_v n_{\mathrm{leg}}!}{\big|\Aut_{\text{fixed-ext}}(\Gamma)\big|},
\]
computed (\file{msrjd/diagrams/symmetry.py:438--439}) as
\code{\_wick\_leg\_factor(td)} divided by the automorphism order with external
legs held fixed. The dedup \emph{multiplicity} --- the size of each equivalence
class --- is therefore \emph{diagnostic only}. Multiplying by it would
\emph{double-count}, because the merged copies are isomorphic and their pairings
are already inside $\Sym(\Gamma)$. The multiplicity is retained in the cache
purely for diagnostics and historical compatibility (it once compensated for an
\emph{incomplete} signature --- see the Gotchas).

\subsection{The headline function and its cache key}

\code{enumerate\_unique\_diagrams} (\file{\_diagrams.py:54}) takes the expanded
\code{ft} (only \code{ft.taylor\_order} is read here, for the key), the
\code{model} (for its name), $k$ and \code{max\_ell}, the
\code{external\_fields}, the propagator \code{G\_ft}, the field-index maps
\code{resp\_idx}/\code{phys\_idx}, and the type lists \code{vtypes}/\code{stypes}.
It returns the triple \code{(unique\_by\_ell, multiplicity\_by\_ell,
all\_unique)}: a dict from $\ell$ to the surviving diagrams, a parallel dict of
their equivalence-class sizes, and a flat $\ell$-ordered concatenation.

The cache directory is per-theory, computed by \code{\_model\_cache\_dir}
(\file{\_diagrams.py:42}), which uses the very same slug rule as the expand
cache so all artefacts share one directory. The cache key, though, is finer
than the expand cache's --- it must distinguish $k$, $\ell$, external fields,
\emph{and} taylor order. The stage name (\file{\_diagrams.py:150}) is:

\begin{lstlisting}
stage_name = f'unique_typed_mult_v3_{ext_tag}_taylor{ft.taylor_order}'
\end{lstlisting}

with $k$ and $\ell$ appended by \code{PipelineCache} as \code{\_k<k>\_l<l>}.
Read the pieces:

\begin{description}
  \item[\code{unique\_typed}] the stage --- deduplicated typed diagrams.
  \item[\code{\_mult\_}] records that the file carries multiplicities.
  \item[\code{\_v3\_}] the cache-\emph{version} tag (see below).
  \item[\code{\_<ext\_tag>}] the external-field tag from
        \code{\_ext\_fields\_tag} (\file{\_diagrams.py:37}), e.g.
        \code{[('v',1),('v',2)]} becomes \code{v1\_v2}. This is what makes
        different external-field permutations land in distinct cache files
        automatically.
  \item[\code{\_taylor<N>}] the taylor-order dependence, so sibling files for
        different orders coexist in one directory.
\end{description}

The loop body (\file{\_diagrams.py:156--209}) is the manual memoise pattern. On
each $\ell$: if \code{use\_cache and cache.exists(\dots)}, try to \code{load}
the \code{\{'unique', 'multiplicities'\}} dict, populate the return structures,
and \code{continue}; any exception prints a \code{cache load failed \dots
rebuilding} message and falls through to recompute. Otherwise build the four
stages, and if \code{use\_cache} \code{save} the result --- wrapped in
\code{try}/\code{except} so a save failure only warns.

\subsection{Version tags are load-bearing}

The \code{\_v3\_} in the stage name is not decoration. It is the \emph{only}
invalidation mechanism for the enumeration cache: old files are never read
because their filename stem differs. The long comment block at
\file{\_diagrams.py:123--149} records the history. The tag was bumped
\code{v1}$\to$\code{v2} (2026-05-26) when the source-type extractor began
promoting cross-/auto-cumulant monomials to a colored \code{NoiseSourceType}; a
\code{v1} cache would encode the plain type and silently collapse every
colored-noise diagram to the white-noise answer. It was bumped
\code{v2}$\to$\code{v3} (2026-06-10) when \code{diagram\_signature} became a
\emph{complete} isomorphism invariant; a \code{v2} cache had been deduplicated
with the old incomplete signature, which collided non-isomorphic diagram classes
at $k\ge 3$ and silently dropped their integrals.

\begin{gotcha}
If you change the dedup or typing semantics, you \emph{must} bump the version
string in \code{stage\_name}. A stale cache with the old stem will load and
silently produce \emph{wrong numbers}. The \code{v2}$\to$\code{v3} bump is the
concrete cautionary tale: with the old incomplete signature, a particular
three-point cumulant came out as $-\tfrac{68}{3}\,a^3$ instead of the correct
$-32\,a^3$, because four of the eleven one-loop classes at $k=3$ had been merged
into the wrong representatives.
\end{gotcha}

\section{The bridge between the two caches: the taylor-order formula}

The two caches are linked by one formula. By default
(\file{pipeline/compute.py:279}):

\begin{lstlisting}
taylor_order = max(k + 2 * max_ell, 2)
\end{lstlisting}

The intuition: a connected diagram with $k$ external legs and $\ell$ loops needs
interaction vertices whose total leg-count covers the $k$ external legs plus the
internal propagator legs. Each loop adds (heuristically) two leg-degrees of
internal structure, giving the $2\,\texttt{max\_ell}$ term; the floor of $2$
guarantees the propagator sector $(1,1)$ is always present, even at
$k=2,\ \texttt{max\_ell}=0$. This is the taylor order needed so that
\emph{every} diagram at $(k,\ell)$ has its vertices in scope.

This formula is the bridge: it tells you which \code{taylor\_order} the expand
cache must serve to satisfy a given diagram request, and it determines the
\code{\_taylor<N>} suffix in the enumeration-cache filename. A historical floor
of $4$ was dropped to $2$ once the cache layout could side different orders
cleanly --- worth $\sim$90 minutes on a heavy Bernoulli theory at
$k=2,\ \texttt{max\_ell}=0$, because that run now needs only the cheap order-2
expansion.

\section{Data-flow walk-through}

Put the pieces together with one full
\code{compute\_cumulants(model, k=2, max\_ell=1, \dots)} run, end to end through
both caches.

\begin{enumerate}
  \item \textbf{Taylor order chosen} (\file{compute.py:279}):
        $\texttt{taylor\_order} = \max(2 + 2\cdot 1,\, 2) = 4$.
  \item \textbf{Expand cache} (\file{compute.py:319--332}). Build
        \code{ft = FieldTheory(model, taylor\_order=4)}, then
        \code{cached\_order = find\_best\_cached\_order(model, 4)} --- which is
        \code{4} if an order-4 file exists, \code{6} if only a higher order is
        cached (then downgrade-filtered), or \code{None} on a miss. On a hit,
        \code{prepare\_for\_load(ft)} then
        \code{load\_expand(model, ft, target\_order=4, cached\_order=\dots)}
        populates \code{ft.\_by\_tp}. On a miss, \code{ft.expand()} then
        \code{save\_expand(model, ft)} writes \file{expand\_taylor4.sobj}.
  \item \textbf{Downstream of expand.} The vertex/source type extractors, the
        propagator builder, and the field-index-map builder all read
        \code{ft.\_by\_tp} to produce \code{vtypes}, \code{stypes},
        \code{G\_ft}, and \code{resp\_idx}/\code{phys\_idx}.
  \item \textbf{Enumeration cache} (\file{compute.py:649}).
        \code{enumerate\_unique\_diagrams(\dots, use\_cache=True)} with
        \code{external\_fields=[('v',1),('v',2)]} forms
        \code{stage\_name = 'unique\_typed\_mult\_v3\_v1\_v2\_taylor4'}. For
        $\ell=0$ it consults
        \file{unique\_typed\_mult\_v3\_v1\_v2\_taylor4\_k2\_l0.sobj} (hit:
        load; miss: build four stages, save); for $\ell=1$ the matching
        \code{\_k2\_l1} file. It returns
        \code{unique\_by\_ell=\{0:[\dots], 1:[\dots]\}}, the parallel
        multiplicities, and the flat \code{all\_unique}.
  \item \textbf{Consumers} (\file{compute.py:680+}). Walk \code{unique\_by\_ell}
        per $\ell$, classify each diagram's coefficient factors, group by
        kernel, and integrate.
\end{enumerate}

The \code{precompute} flow is the same but truncated: it exercises only the
expand cache at order~2 and the propagator cache, writing
\file{expand\_taylor2.sobj} and \file{propagator.sobj}. A later
\code{compute\_cumulants} at order~2 then hits both immediately; at order~4 it
still pays the order-4 \code{expand()} once.

\begin{note}
The \term{spatial bridge} (\file{msrjd/integration/spatial/pipeline\_bridge.py:242})
calls \code{enumerate\_unique\_diagrams} with \code{use\_cache=False}: it
re-enumerates every time, because the diagram topology is cheap relative to the
spatial integral and the bridge does not want to manage cache keys. Note that
\code{use\_cache=False} also means \emph{never write} --- such a run does not
pollute the cache, but gets no persistence either.
\end{note}

\section{Gotchas and caveats}

A consolidated list of the sharp edges, most of which we have met in passing.

\begin{gotcha}
\textbf{Fork-based \code{multiprocessing} is dangerous on macOS notebooks.}
\code{enumerate\_unique\_diagrams} exposes \code{parallel=True}, which fans the
type-assignment stage across a \emph{fork-based} pool. Forking a Jupyter kernel
after the GUI/BLAS libraries have initialised can crash the kernel \emph{and}
the OS; the analogous temporal path was guarded for exactly this reason, and the
spatial path was made serial-only. The default here is \code{parallel=False}. Do
not flip \code{parallel=True} from a notebook on macOS --- the cache wrapper
does not re-apply the notebook-safety guard; it trusts the caller. (Cache hits
skip the parallel path entirely, since they never reach the typing stage.)
\end{gotcha}

\begin{gotcha}
\textbf{A model edit that does not change the name silently re-uses the cache.}
The slug derives from \code{model['name']} only. Editing the field list or a
vertex without bumping the name leaves a stale \file{expand\_taylor<N>.sobj} and
stale enumeration files. The partial defences are the \code{ring\_var\_names}
and \code{n\_tilde} structural checks (which catch field-list changes that alter
the ring) and the form-factor signature (which catches operator-IR vertex
changes) --- but a change that alters a coupling \emph{value} baked into a
coefficient while keeping the same ring would \emph{not} be caught. The
docstring advice (\file{\_diagrams.py:18--22}): bump \code{model['name']} or
delete the cache directory on model edits.
\end{gotcha}

\begin{gotcha}
\textbf{\code{\_model\_cache\_dir} ignores its \code{taylor\_order} argument.}
The directory is per-theory; the parameter is in the signature but never read ---
a vestige of an older per-taylor-order directory layout. Harmless, but confusing
the first time you read the call.
\end{gotcha}

\begin{gotcha}
\textbf{Graceful degradation differs slightly between the two caches.} On a load
problem the expand cache returns \code{False} (a clean miss, triggering a fresh
expand), whereas the enumeration cache \emph{rebuilds the single $\ell$ slot}
(\file{\_diagrams.py:170--173}). Both degrade gracefully and never raise out of
a cache miss --- but the granularity is per-bundle for one and per-loop-order
for the other.
\end{gotcha}

In short: the cache key is structural, the expand cache trades disk for a
ninety-minute expansion, the enumeration cache trades disk for a combinatorial
blow-up, the superset theorem makes downgrade exact, and the version string and
form-factor signature are the two guards standing between a stale file and a
wrong number. Mind those two, and the rest takes care of itself.


% =====================================================================
%  PART IV  --  EVALUATING DIAGRAMS : TEMPORAL
% =====================================================================
\part{Evaluating Diagrams: Temporal Models}
% ---------------------------------------------------------------------
%  Chapter 12 : Phase J --- The Temporal Loop-Integral Engine
% ---------------------------------------------------------------------
\chapter{Phase J: The Temporal Loop-Integral Engine}\label{ch:phasej-temporal}

Every chapter so far has been preparation for this one. The theory file gave us
an action; the field-theory core expanded it into vertices; the propagator
chapter built the free two-point function; the mean-field solve found the
saddle; the enumeration and type-assignment chapters produced a list of typed
Feynman diagrams, each tagged with a scalar prefactor. What we do \emph{not} yet
have is a single number. A diagram is still an abstract graph plus a
prefactor; turning it into a value --- a function of the external times we can
plot against a simulation --- is the job of \term{Phase~J}, the final,
time-domain evaluation layer of the temporal pipeline.

This is the numerical heart of the whole temporal branch. It is also the most
intricate piece of code in the repository: \file{final\_integral.py} is a single
file of roughly 5{,}300 lines. The intricacy is not gratuitous. The integral a
diagram represents is, in general, a multi-dimensional integral over a region
carved out by causality constraints, and the engine goes to considerable
lengths to evaluate that integral \emph{analytically} --- in closed form ---
rather than handing it to a black-box numerical quadrature routine. This chapter
explains, from the ground up, why that integral has a closed form, what the
closed form is in each of the cases the engine distinguishes, and how the code
realises it.

\begin{note}
The name ``Phase~J'' is a historical artifact. The temporal pipeline was
originally built as a hybrid: the early phases (``I'' and before) work in
frequency space, and ``J'' was the last phase, the one that returns to the time
domain to finish the job. It began life as a tree-level ``minimum viable
product,'' but the core integrator (\code{integrate\_diagram}) has since been
generalised to handle \emph{any} loop order. The label stuck even though the
algorithm outgrew its original scope.
\end{note}

\section{What Phase J computes}

In the Martin--Siggia--Rose--Janssen--De~Dominicis (\msrjd) formalism the
quantity we ultimately want is a time-domain correlation or response function of
a stochastic dynamical system --- for example the autocorrelation
$C(\tau) = \avg{\delta n(t)\,\delta n(t+\tau)}$ of a Hawkes process, or a
higher-order cumulant of any of the temporal models in the library. Perturbation
theory expresses that correlator as a sum over Feynman diagrams. Each diagram is
a graph whose

\begin{itemize}
  \item \textbf{edges} are \emph{retarded propagators} $\Gret(t_v - t_u)$ --- the
        linear response of one field to another, carrying a Heaviside step
        $\Theta$ that enforces causality (a response can only happen \emph{after}
        its cause), and whose
  \item \textbf{internal vertices} are interaction points whose \emph{times}
        must be integrated over the whole real line, subject to the causal
        orderings the Heavisides impose.
\end{itemize}

So the value of one diagram is a multi-dimensional integral of the schematic
form
\begin{equation}
  C_\Gamma(\text{external times})
   \;=\; \text{prefactor} \cdot
   \int \dd s_1 \cdots \dd s_m \;\prod_{e \in E} \Gret\big(t_{\text{head}(e)} - t_{\text{tail}(e)}\big),
  \label{eq:diagram-integral}
\end{equation}
where $s_1, \dots, s_m$ are the internal-vertex times. \textbf{Phase~J is the
machine that performs that integral}, for every diagram, at every loop order,
returning a Python callable \code{f(*external\_times) -> complex} that you can
evaluate on any grid of external times.

\begin{defn}
A \textbf{retarded propagator} $\Gret(\Delta t)$ is the linear response of a
physical field to a perturbation in its conjugate response field a time
$\Delta t$ earlier. ``Retarded'' means it vanishes for $\Delta t < 0$: the
effect cannot precede the cause. In the diagrams this causality is carried by an
explicit Heaviside factor $\Theta(\Delta t)$ on every edge, and those Heavisides
are precisely what turn \eqref{eq:diagram-integral} from an integral over all of
$\mathbb{R}^m$ into an integral over a bounded \emph{causal region}.
\end{defn}

The cleverness of the engine is that it almost never falls back to brute-force
numerical quadrature. Because every retarded propagator turns out to be a
\emph{finite sum of decaying complex exponentials} (we derive this in
\S\ref{sec:modesum}), the whole integrand is a sum of terms of the form
$A \cdot \exp(\text{linear combination of the times})$, and the integral of such
a term over the causal region --- which is a \emph{polytope}, a region bounded by
flat faces --- has a closed form. The engine therefore carries three dedicated
\emph{analytic} integrators, one each for $1$, $2$, and $\ge 3$ internal-vertex
integration variables, and only routes to \code{scipy}'s numerical quadrature
when the analytic path hits a genuine degeneracy or a floating-point overflow.

\subsection{Where it sits, and what feeds it}

Phase~J consumes four things, all produced by earlier phases:

\begin{enumerate}
  \item a list of \term{typed diagrams} (\code{TypedDiagram} objects) from the
        enumeration, typing, and symmetry-deduplication pipeline;
  \item a \emph{scalar prefactor} per diagram --- the combinatorial and coupling
        weight, from \code{classify\_coefficient\_factors};
  \item the retarded propagator in \term{pole--residue form}: a dictionary
        \code{propagator\_data} with keys \code{pole\_vals} (the poles $p_\alpha$),
        \code{C\_mats} (one residue matrix per pole), and \code{D\_delta} (the
        instantaneous $\delta$-function part);
  \item the user's external-field list and the numerical parameter
        substitutions \code{num\_params}.
\end{enumerate}

\noindent The public entry point is \code{compute\_correction\_td} in
\file{pipeline.py}, which loops over the diagrams calling \code{integrate\_diagram}
(in \file{final\_integral.py}) on each. Downstream, the result feeds the
notebook-facing callables \code{total\_C} and \code{total\_C\_batch} that produce
the theory curves you overlay on simulation data.

The three source files this chapter dissects are:

\begin{description}
  \item[\file{final\_integral.py}] the $\sim$5{,}300-line core: the per-diagram
        integrator and all of the analytic primitives.
  \item[\file{pipeline.py}] the orchestration and batching layer: it loops over
        diagrams and fans the $\tau$-grid evaluation out over worker processes.
  \item[\file{subgraph.py}] a near-stub for a frequency-domain loop-reduction
        route that was designed but never finished --- the time-domain integrator
        superseded it. We cover it briefly at the end so you know it is dead
        weight, not a live path.
\end{description}

\section{The retarded propagator as a mode sum}\label{sec:modesum}

Everything downstream rests on one fact: the time-domain retarded propagator is
a finite sum of complex exponentials. Let us derive it, because the code is a
literal transcription of the result.

\subsection{From frequency to time by residues}

For a \emph{linear} response system the frequency-domain propagator
$\widehat{G}(\omega)$ is a rational matrix function of $\omega$ --- a ratio of
polynomials, entry by entry. Its inverse Fourier transform, under the project's
fixed convention
\begin{equation}
  G(t) \;=\; \frac{1}{2\pi}\int \dd\omega \; \ee^{\ii \omega t}\,\widehat{G}(\omega),
  \label{eq:fourier-convention}
\end{equation}
is computed by closing the $\omega$ contour and summing residues. \emph{Why
upward?} For $t>0$ the kernel $\ee^{\ii\omega t}$ decays as
$\operatorname{Im}\omega\to+\infty$, so by Jordan's lemma the contour may be
closed by a large semicircle in the \emph{upper} half plane at no added cost,
enclosing exactly the upper-half-plane poles (and for $t<0$ one closes downward,
enclosing none). The causality
filter built into the propagator construction guarantees that every pole
$p_\alpha$ of $\widehat{G}$ lies in the \emph{upper} half plane,
$\operatorname{Im} p_\alpha > 0$. Each residue term therefore carries a factor
$\ee^{\ii p_\alpha t}$, which \emph{decays} for $t > 0$ (its exponent has real
part $-\operatorname{Im} p_\alpha < 0$) and \emph{grows} for $t < 0$. The
physical retarded propagator is this analytic sum multiplied by a Heaviside that
kills the growing $t<0$ half:
\begin{equation}
  \Gret[p, r](\Delta t)
   \;=\; \underbrace{\texttt{delta\_coeff}[p,r]\cdot\delta(\Delta t)}_{\text{instantaneous part}}
   \;+\; \Theta(\Delta t)\cdot\sum_\alpha C_\alpha[p,r]\,\ee^{\ii p_\alpha \Delta t}.
  \label{eq:GR-decomp}
\end{equation}

Three pieces of nomenclature, each of which has a direct counterpart in the
code:

\begin{itemize}
  \item The indices $[p, r]$ are the \textbf{(physical-row, response-column)}
        matrix indices. The convention is $G[p_i, r_i] = \avg{\phi_{p_i}\,
        \tilde n_{r_i}}$ --- the response of physical field $p_i$ to a
        response-field (hatted) source $r_i$. There is a transpose to watch
        here: the kernel matrix $K$ from which the propagator is built is laid
        out \code{[resp, phys]}, so to read off a retarded entry you index
        $G$\code{[phys, resp]}. The propagator chapter covers this; Phase~J just
        respects it.
  \item $\texttt{delta\_coeff}[p,r] = \lim_{\omega\to\infty}\widehat{G}[p,r](\omega)$
        is the \textbf{instantaneous response} --- the polynomial (non-proper)
        part of the rational $\widehat{G}$. It is nonzero exactly when there is
        an instantaneous coupling: for example the $\tilde n \times \delta n$
        term in an \msrjd{} action, where a $\tilde n$ source at time $t$
        produces an \emph{immediate} $\delta n$ at the same $t$. It is captured
        in \code{propagator\_data['D\_delta']}.
  \item $C_\alpha[p,r]$ is the residue of $\widehat{G}$ at pole $p_\alpha$,
        position $[p,r]$.
\end{itemize}

To keep the inner loops in plain Python arithmetic, the code stores not the pole
$p_\alpha$ but the combination $\lambda_\alpha := \ii\, p_\alpha$, so that a mode
reads simply $C_\alpha\,\ee^{\lambda_\alpha \Delta t}$ with no factor of $\ii$ to
re-apply on every evaluation.

\subsection{The \code{EdgeModeSum} record}

This entire object --- one edge's worth of \eqref{eq:GR-decomp} --- is packaged
into a single frozen dataclass, \code{EdgeModeSum}
(\file{final\_integral.py:116}). Its documented fields are the residue and pole
indices, the $\delta$-coefficient, and the list of modes:

\begin{lstlisting}
@dataclass(frozen=True)
class EdgeModeSum:
    ri: int
    pi: int
    delta_coeff: complex
    modes: tuple                 # tuple[tuple[complex, complex], ...]  i.e. (C_alpha, lambda_alpha)
    dt_c0: float = 0.0
    dt_int_pairs: tuple = ()     # sparse coefficients on integration vars
    dt_ext_pairs: tuple = ()     # sparse coefficients on free external times
\end{lstlisting}

\begin{note}
\code{@dataclass} is a standard-library decorator (from the \code{dataclasses}
module) that auto-generates a class's constructor, \code{\_\_repr\_\_}, and
equality from a list of typed field declarations --- you write the fields, Python
writes the boilerplate. \code{frozen=True} makes instances immutable (assigning
to a field raises), which is what lets these records be safely cached and shared
across the evaluation without fear of one code path mutating another's data.
\end{note}

The \code{modes} tuple holds the $(C_\alpha, \lambda_\alpha)$ pairs. The three
\code{dt\_*} fields are left empty at construction; they hold the linear form for
$\Delta t$ in terms of the integration variables and external times, and they
are filled in later, once we know which subset of edges we are in (see
\S\ref{sec:polytope}). Carrying them on the same object means the one
\code{EdgeModeSum} threads through the whole evaluation chain.

The builder \code{\_build\_edge\_mode\_sums(edge\_info, propagator\_data)}
(\file{final\_integral.py:150}) does the extraction \emph{once} per edge. It
reads the poles and residue matrices out of \code{propagator\_data} and converts
them to machine-precision Python \code{complex} numbers:

\begin{lstlisting}
modes_lambdas = tuple(complex(CDF(SR(p))) * 1j for p in pole_vals)
# ... then per edge, per pole:
residues = tuple(
    complex(CDF(SR(C_mats[k][pi, ri])))
    for k in range(n_poles)
)
modes = tuple(zip(residues, modes_lambdas))
\end{lstlisting}

\noindent The idiom \code{complex(CDF(SR(expr)))} is how this codebase turns a
fully-substituted symbolic scalar into a Python \code{complex}: \code{SR} coerces
to a symbolic expression, \code{CDF} (the Complex Double Field) evaluates it to
machine precision, and \code{complex(...)} hands back a native Python number. We
explain \code{SR} and \code{CDF} in \S\ref{sec:tools}. If the propagator data is
incomplete --- missing \code{pole\_vals} or \code{C\_mats} --- the builder returns
\code{None}, which is the signal for the caller to fall back to the slow
symbolic path.

\begin{gotcha}
Extracting the residues once, here, rather than re-extracting them inside the
per-call evaluator, is a deliberate and important optimisation. The original
code re-derived the smooth-part data on every \code{(diagram, subset)} call even
though it depends only on the $(p_i, r_i)$ of the edge. Lifting the extraction to
this single per-edge step is what keeps the hot numerical loop in plain
\code{cmath}, never touching Sage --- Sage is far too slow to call per
evaluation.
\end{gotcha}

\section{From a diagram to a polytope integral}\label{sec:polytope}

We now assemble \eqref{eq:diagram-integral} concretely and watch it collapse to
the analytic form. There are three moving parts: the $\delta$-subset expansion,
the elimination of variables by $\delta$-edge equations, and the half-space
constraints that build the polytope.

\subsection{The $\delta$-subset expansion}

A diagram with edge set $E$ has integrand $\prod_{e\in E}\Gret(\Delta t_e)$.
Substitute the two-piece decomposition \eqref{eq:GR-decomp} into each factor and
\emph{expand the product}. Each edge is taken either in its $\delta$-form or in
its smooth (exponential) form, so the product becomes a sum over subsets
$S\subseteq E$ of ``which edges are $\delta$'':
\begin{equation}
  C_\Gamma = \sum_{S\subseteq E}\int \dd s_1\cdots \dd s_m\;
    \underbrace{\Big(\prod_{e\in S}\texttt{delta\_coeff}[e]\,\delta(\Delta t_e)\Big)}_{\delta\text{ edges}}
    \cdot
    \underbrace{\Big(\prod_{e\notin S}\Theta(\Delta t_e)\sum_\alpha C_\alpha[e]\,\ee^{\lambda_\alpha \Delta t_e}\Big)}_{\text{smooth edges}}
    \cdot \text{prefactor}.
  \label{eq:subset-sum}
\end{equation}

This is the central organising idea of the engine. Within each subset:

\begin{itemize}
  \item A \textbf{$\delta$-edge equation} $\Delta t_e = 0$ (for $e\in S$) is
        linear in the vertex times. Each one is solved to \emph{eliminate one
        integration variable} by substitution --- it pins one internal vertex's
        time to another's. This \emph{reduces} the dimension $m$ of the remaining
        integral.
  \item If a $\delta$-edge equation has \emph{no integration variable left to
        solve for}, it becomes a constraint among the external times alone --- a
        \textbf{shot-noise $\delta(\tau)$ spike}. The correlator then contains a
        term $\propto \delta(a\cdot\tau + c)$. These are distributional and are
        \emph{not} added to the smooth callable; they are emitted separately as
        structured \code{delta\_contributions} (covered in
        \S\ref{sec:deltaspike}).
  \item A \textbf{smooth edge} contributes its mode sum
        $\sum_\alpha C_\alpha\,\ee^{\lambda_\alpha \Delta t_e}$ and, via its
        Heaviside, a \textbf{half-space constraint} $\Delta t_e > 0$. Together,
        the smooth edges' constraints carve a convex \term{polytope} out of the
        space of remaining integration variables.
\end{itemize}

\begin{defn}
A \textbf{polytope} is the intersection of finitely many half-spaces --- a
convex region with flat faces (a polygon in 2D, a polyhedron in 3D, and so on).
In Phase~J the polytope is the set of internal-vertex times that respect all the
causal orderings $\{\Delta t_e > 0\}$ a subset imposes. The integral
\eqref{eq:subset-sum} is, per subset, an exponential integrated over this
polytope.
\end{defn}

\begin{gotcha}
Naively the subset sum has $2^{|E|}$ terms. The engine shrinks this
dramatically: edges whose $\delta$-coefficient is numerically zero
(\code{|delta\_coeff| < 1e-15}) are \emph{forced smooth} --- they have no
$\delta$-form to branch on --- so the sum runs only over the edges that actually
carry an instantaneous part, turning $2^{|E|}$ into $2^{|\text{branch}|}$. For
most temporal models only a handful of edges (the noise vertices) have a
$\delta$-part at all, so $|\text{branch}|$ is small.
\end{gotcha}

\subsection{Affine $\Delta t$ and the pole-tuple expansion}

After $\delta$-elimination, every surviving $\Delta t_e$ is an \emph{affine}
(linear-plus-constant) function of the integration variables
$s = (s_0, \dots, s_{m-1})$, the free external times $t_{\text{free}}$, and a
constant:
\begin{equation}
  \Delta t_e = c^{(e)}_0 + \sum_v a^{(e)}_{\text{int}}[v]\, s_v + \sum_j a^{(e)}_{\text{ext}}[j]\, t_{\text{free}}[j].
  \label{eq:dt-affine}
\end{equation}

Now choose one pole $\alpha_e$ per smooth edge --- a \term{pole tuple}. With that
choice fixed, the product of the chosen modes is a \emph{single} exponential
$C_{\text{prod}}\cdot\exp\!\big(\sum_e \lambda_{\alpha_e}\Delta t_e\big)$. Group
the exponent by what it multiplies:
\begin{align}
  \alpha_s[v]  &= \sum_e \lambda_{\alpha_e}\, a^{(e)}_{\text{int}}[v]   &&\text{(coefficient of } s_v\text{)}, \label{eq:alpha-s}\\
  \gamma       &= \sum_e \lambda_{\alpha_e}\Big(c^{(e)}_0 + \sum_j a^{(e)}_{\text{ext}}[j]\,t_{\text{free}}[j]\Big) &&\text{(the } s\text{-independent part)}. \label{eq:gamma}
\end{align}
So the integrand for one pole tuple is
$\text{prefactor}\cdot C_{\text{prod}}\cdot \ee^{\gamma}\cdot \exp\!\big(\sum_v \alpha_s[v]\,s_v\big)$,
and the whole subset integral is the sum, over all
$\prod_e (\#\text{poles}_e)$ pole tuples, of
\begin{equation}
  \int_{\text{polytope}} \exp\!\Big(\sum_v \alpha_s[v]\,s_v\Big)\,\dd s.
  \label{eq:exp-over-polytope}
\end{equation}
Each of these is a closed-form \emph{exponential-over-a-polytope}. This is the
integral the three analytic backends evaluate.

\begin{defn}
A \textbf{pole tuple} is a choice of exactly one pole $\alpha_e$ for each smooth
edge $e$ of a subset. Because the propagator's time-domain form is a sum over
poles, the product over edges of those sums expands into a sum over pole tuples;
each tuple contributes one pure exponential, and the polytope integral of a pure
exponential is closed-form. The number of pole tuples is the product of the
per-edge pole counts.
\end{defn}

\subsection{The plan cache}

The coefficients $\alpha_s[v]$ depend only on the pole tuple and the diagram
geometry --- \emph{not} on the value of the external times. Only $\gamma$ depends
on $t_{\text{free}}$, and it does so \emph{linearly}. The engine exploits this by
pre-computing, once per subset, everything that is $\tau$-invariant. That is what
\code{\_build\_modesum\_plan(...)} (\file{final\_integral.py:595}) does: it
returns a dict whose keys are \code{pole\_tuples}, \code{alphas\_per\_tuple} (the
$\alpha_s[v]$ of \eqref{eq:alpha-s}), \code{gamma\_const\_per\_tuple} ($\gamma$
evaluated at $t_{\text{free}}=0$), and \code{gamma\_slope\_per\_tuple\_per\_ext}
($\partial\gamma/\partial t_{\text{free}}[j]$). The per-$\tau$ work then collapses
to a single linear combination
\[
  \gamma = \gamma_{\text{const}} + \sum_j \text{slope}[j]\cdot t_{\text{free}}[j],
\]
followed by the closed-form polytope integral. The expensive part --- enumerating
pole tuples and building the $\alpha_s$ vectors --- is paid once, not once per
grid point.

The pole tuples themselves are produced by \code{\_enumerate\_pole\_tuples}
(\file{final\_integral.py:529}), a generator yielding $(C_{\text{product}},
\lambda\text{s})$ over the Cartesian product of the per-edge modes. It advances a
stack of indices by hand rather than calling \code{itertools.product}, to avoid
the per-tuple object-allocation overhead; an empty edge set yields the single
trivial tuple $(1.0+0j, ())$.

\section{Translation invariance and external-Wick sums}\label{sec:wick}

Before we reach the integrators, two correctness subtleties have to be handled in
the construction of the callable: pinning the origin, and summing over Wick
contractions of identical external legs.

\subsection{Origin pinning}

A stationary correlator depends only on \emph{time differences}. There is no
absolute clock, so we are free to fix one external leaf at $t=0$. The engine pins
the leaf at canonical position \code{origin\_leaf\_idx} (default $0$) to
\code{SR(0)} during integrand construction. The callable it returns therefore
computes a function of $(t_j - t_{\text{origin}})$ --- exactly the physical
content, with the redundant overall time shift removed. For a $k=2$
autocorrelator this means \code{total\_C(t\_1, t\_2)} returns $C$ at
$\tau = t_2 - t_1$.

\subsection{Summing Wick contractions, and the compensation divisor}

When several external legs carry the \emph{same} field type --- say two $\delta
n_1$ legs --- the symmetry-deduplication upstream has merged diagrams that differ
only in which same-type leaf attaches where. Phase~J must compensate by summing
over all such leaf permutations. Concretely, it enumerates every
field-respecting assignment of canonical positions to leaves (the
\code{\_all\_mappings} list), sums the integrand over them, and divides by the
\textbf{external-Wick compensation factor}
\begin{equation}
  \texttt{comp} \;=\; \frac{|\Aut(\Gamma,\ \text{leaves free})|}{|\Aut(\Gamma,\ \text{leaves fixed})|}.
  \label{eq:wick-comp}
\end{equation}
This is computed by \code{external\_wick\_compensation} in
\file{symmetry.py:343}. The justification is orbit--stabilizer: two assignments
yield the same pinned-external diagram iff they differ by an automorphism that
permutes leaves, and the stabilizer of any assignment is exactly
$\Aut_{\text{fixed}}$, so the mapping sum reproduces each distinct pinned diagram
$|\Aut_{\text{free}}|/|\Aut_{\text{fixed}}|$ times. Dividing by that index makes
the sum equal the sum over \emph{distinct} pinned diagrams --- which is what the
Feynman rules require, since each pinned diagram already carries its full weight
through its combinatorial factor (whose denominator is the same
$\Aut_{\text{fixed}}$).

\begin{gotcha}
The compensation divisor must be the \emph{exact} orbit--stabilizer index of
\eqref{eq:wick-comp}, not a cheaper proxy. A previous implementation approximated
it as $\prod N!$ over leaves grouped by their attachment vertex's
\emph{role signature} (a shallow depth-1 invariant). That over-divides whenever
two attachment vertices share a role signature but are \emph{not} related by any
automorphism --- e.g. the $k=4$ OU$+\varepsilon x^3$ one-loop cascades, where
three noise vertices all read ``noise feeding a cubic'' yet hang off
structurally distinct cubics. The resulting $\times\frac13$ / $\times\frac12$
deficits broke every $k\ge 3$ one-loop cumulant while coincidentally leaving
$k=2$ exact. The fix (validated against the exact Boltzmann series
$\kappa_4 = -6\varepsilon + 126\varepsilon^2$) is the true automorphism index.
When \code{external\_fields} is not supplied and a field-count mismatch forces
the identity-mapping fallback, the divisor must be exactly $1$, never the old
heuristic.
\end{gotcha}

\subsection{Re-pinning the origin per permutation}

There is a second, subtler point in the final callable
(\file{final\_integral.py:3759}). The integrand was \emph{built} with the origin
leaf pinned to $0$. When a non-identity permutation relabels which leaf occupies
which canonical position, the origin leaf may land at a non-zero time in the
user's input frame. The callable therefore time-translates each permuted
configuration so the origin returns to $0$ --- it subtracts the permuted origin's
time from all the others:

\begin{lstlisting}
total = 0.0 + 0.0j
for perm in _perms:
    permuted = [ext_time_values[perm[j]] for j in range(_k)]
    if origin_leaf_idx is not None:
        t_origin = permuted[origin_leaf_idx]
        # ... subtract t_origin from the free legs ...
\end{lstlisting}

\begin{gotcha}
This per-permutation re-pinning is load-bearing. For an \emph{asymmetric}
integrand --- which is every one-loop-and-higher diagram with a distinguished leaf,
such as a cubic-vertex tadpole --- the different permutations are physically
distinct evaluations (for a $k=2$ swap, $V(\tau)$ and $V(-\tau)$) and must be
genuinely summed, not merely counted. Before this fix the swap permutation was
fed the pinned origin's actual time, $\code{free\_val}=0$, producing spurious
asymmetry in $C(\tau)$ for every non-tree $k=2$ identical-externals case.
\end{gotcha}

\section{The $m=1$ interval integrator}\label{sec:m1}

We now turn to the three analytic backends, in order of increasing dimension.
The dispatch on $m$ (the number of integration variables surviving
$\delta$-elimination) happens inside each subset's contribution closure.

For $m=1$ the polytope is an interval $[L, U]$ on the single integration variable
$s$. With one variable, \eqref{eq:exp-over-polytope} is the elementary
\begin{equation}
  \int_L^U \ee^{\alpha_s\, s + \gamma}\,\dd s
   = \ee^{\gamma}\cdot\frac{\ee^{\alpha_s U} - \ee^{\alpha_s L}}{\alpha_s},
  \label{eq:m1-closed}
\end{equation}
with the removable-singularity limit $(U-L)\,\ee^{\gamma}$ as $\alpha_s\to 0$. The
implementation is \code{\_integrate\_1d\_polytope\_modesum}
(\file{final\_integral.py:1976}). The bounds $L$ and $U$ come from intersecting
the smooth-edge half-line constraints $a^{(e)}_{\text{int}}[0]\,s + c_{\text{ext}}^{(e)} > 0$.

\begin{note}
\textbf{Unbounded endpoints are allowed.} If a side of the interval is $\pm\infty$,
the corresponding boundary term in \eqref{eq:m1-closed} evaluates to $0$
\emph{provided} the integrand decays in that direction --- that is, provided
$\operatorname{sign}(\operatorname{Re}\alpha_s)$ matches the unbounded side. If it
does not, the integral genuinely diverges along the analytic path, and the
integrator returns \code{None}, which tells the caller to fall back to
\code{scipy.quad} (which handles $\pm\infty$ bounds natively via a coordinate
transform). Returning \code{None} is the universal ``I cannot do this one
analytically, try numerics'' signal throughout the engine.
\end{note}

This same function is also called directly by the \emph{grouped} Phase~J path
(Chapter~\ref{ch:grouped-phasej}) through its optional \code{pole\_tuples}
argument, which lets that path inject pre-summed residues. We mention it so the
shared interface is on your radar; the grouped path is the subject of the next
chapter.

\section{The $m=2$ polygon integrator}\label{sec:m2}

For $m=2$ the polytope is a convex \emph{polygon} in the plane $(s_0, s_1)$ ---
the intersection of the half-planes from the retardation constraints. The
integral \eqref{eq:exp-over-polytope} becomes $\iint_{\text{polygon}}
\exp(\alpha s_0 + \beta s_1)\,\dd A$. The engine evaluates it in four steps,
implemented across \code{\_integrate\_2d\_polygon\_modesum}
(\file{final\_integral.py:2175}) and its helpers.

\subsection{Step 1 --- build the polygon by clipping}

Start from a counter-clockwise bounding box of half-width \code{POLYGON\_BBOX\_CAP}
($=200$) and \emph{clip} it against each constraint half-plane in turn. The
clipping uses the \term{Sutherland--Hodgman} algorithm,
\code{\_clip\_polygon\_to\_halfplane} (\file{final\_integral.py:454}).

\begin{defn}
\textbf{Sutherland--Hodgman polygon clipping} is a classic computer-graphics
algorithm. It walks the edges of a polygon and, for each clipping half-plane,
keeps the vertices on the ``inside,'' inserts the intersection point wherever an
edge crosses the boundary, and drops the outside vertices. Applied once per
constraint, it produces the convex intersection of the box with all the
half-planes. Vertices lying exactly on a boundary are counted as inside --- a
measure-zero choice that does not affect the area integral.
\end{defn}

The driver \code{\_polygon\_from\_2d\_constraints} (\file{final\_integral.py:499})
substitutes the current external-time values into each constraint's affine offset
and clips:

\begin{lstlisting}
polygon = [(-bbox_cap, -bbox_cap), (bbox_cap, -bbox_cap),
           (bbox_cap, bbox_cap), (-bbox_cap, bbox_cap)]
for (a_int, a_ext, c0) in constraint_data:
    c_eff = float(c0) + sum(float(a_ext[j]) * float(free_ext_vals[j])
                            for j in range(len(a_ext)))
    polygon = _clip_polygon_to_halfplane(polygon, float(a_int[0]), float(a_int[1]), c_eff)
    if not polygon:
        return []
\end{lstlisting}

\noindent An empty result (the constraints are jointly infeasible at this
$\tau$) integrates to $0$.

\subsection{Steps 2--4 --- triangulate and map to the unit triangle}

Fan-triangulate the polygon from vertex $0$ (every triangle shares vertex $0$).
For each triangle, affine-map it to the \emph{unit triangle}
$\{0\le u, w;\ u+w\le 1\}$; the exponential integrand becomes
$\exp(\alpha_0 + p\,u + q\,w)$ for some complex $p, q$ read off the vertices. The
unit-triangle integral then has the closed form
\begin{equation}
  J(p, q) = \int_0^1\!\dd u\int_0^{1-u}\!\dd w\;\ee^{p u + q w}
   = \frac{1}{q}\left[\frac{\ee^{p} - \ee^{q}}{p - q} - \frac{\ee^{p} - 1}{p}\right],
  \label{eq:Jpq}
\end{equation}
implemented in \code{\_exp\_over\_unit\_triangle} (\file{final\_integral.py:356}).
The per-triangle wrapper \code{\_exp\_over\_triangle} (\file{final\_integral.py:406})
maps the vertices, calls $J$, and scales by the Jacobian
$|\det|\cdot\ee^{v_0\text{-term}}$.

\begin{gotcha}
\textbf{Equation \eqref{eq:Jpq} is a minefield of removable singularities.} The
denominators $q$, $p$, and $p-q$ each vanish in a different degenerate regime,
and the formula then takes the indeterminate $0/0$ form --- numerically, that is
catastrophic cancellation. The code guards \emph{all four} regimes with
4th-order Taylor expansions, triggered when any of $|p|, |q|, |p-q|$ drops below
\code{\_J\_TAYLOR\_EPS} $= 10^{-6}$. For instance the doubly-degenerate
$|p|,|q|\to 0$ case returns the moment expansion of the triangle directly:
\begin{lstlisting}
return (0.5
        + (p + q) / 6.0
        + (p*p)/24.0 + (q*q)/24.0 + (p*q)/24.0
        + (p**3 + q**3)/120.0
        + (p*p*q + p*q*q)/120.0)
\end{lstlisting}
the coefficients being the moments $\iint 1 = \tfrac12$, $\iint u = \iint w =
\tfrac16$, and so on.
\end{gotcha}

\subsection{The bilateral overflow guard}

\code{\_exp\_over\_triangle} returns \code{None} --- abort, fall back to scipy ---
whenever any of $|\operatorname{Re} p|$, $|\operatorname{Re} q|$,
$|\operatorname{Re}(v_0\text{-term})|$ exceeds \code{EXP\_REAL\_LIMIT} $= 600$.
Note this guard is \emph{bilateral}: it checks the magnitude on \emph{both} signs,
not just the overflow ($+$) side.

\begin{gotcha}
The bilateral guard is not an over-caution; it is a hard-won correctness fix. The
structural cancellation in \eqref{eq:Jpq}, $(\ee^p-\ee^q)/(p-q) -
(\ee^p-1)/p$, becomes dominated by floating-point noise when \emph{one} of
$\ee^p$, $\ee^q$ underflows to zero and the other does not. A 2026 optimisation
(commit 7f0bf05) relaxed the guard to one-sided, reasoning that the underflow
direction is harmless; in fact it introduced wrong-direction loop corrections in
the spike-reset $k=2$, $\ell=1$ case (which exercises the $m=2$ path heavily) and
was reverted. The lesson, recorded in the code comment at line~436: keep the
guard bilateral so the fragile-cancellation cases route to \code{scipy.nquad},
which handles the well-conditioned integral natively.
\end{gotcha}

\section{The $m\ge 3$ causal-poset integrator}\label{sec:m3}

For three or more integration variables the polygon picture no longer scales, and
the engine switches to a combinatorial decomposition built on the
\emph{structure} of the causal constraints. This is the most involved of the
three backends, \code{\_integrate\_nd\_polytope\_poset\_modesum}
(\file{final\_integral.py:1726}).

\subsection{The constraint set as a directed acyclic graph}

Read the constraint set $\{\Delta t_e > 0\}$ as a directed graph on the $m$
integration variables. Each inter-axis constraint $s_v - s_u > 0$ becomes a
directed edge $u\to v$ (read ``$u$ precedes $v$''); each single-variable
constraint becomes a scalar bound $s_v > c$ or $s_v < c$. Because all the
constraints come from causal orderings, this graph is acyclic --- a
\term{directed acyclic graph} (DAG), equivalently a \term{partial order} (poset).
This is extracted into the frozen \code{\_CausalPoset} dataclass
(\file{final\_integral.py:691}) by \code{\_extract\_causal\_poset}
(\file{final\_integral.py:720}), which classifies each constraint as an inter-axis
edge, a scalar lower bound, a scalar upper bound, or --- if it is none of these
clean forms (a mixed or shifted constraint) --- returns \code{None} to force the
scipy fallback.

\subsection{Linear extensions and the chain simplex}

The key geometric fact: the full polytope is the \emph{disjoint union of simplex
regions, one per linear extension of the poset}.

\begin{defn}
A \textbf{linear extension} of a partial order is a total order consistent with
it --- a way of laying all the variables in a single line $s_{\sigma(0)} \le
s_{\sigma(1)} \le \cdots$ that never violates a ``precedes'' relation. It is
exactly a topological sort of the DAG. Each linear extension corresponds to one
\textbf{chain simplex}
\[
  L \le s_{\sigma(0)} \le s_{\sigma(1)} \le \cdots \le s_{\sigma(m-1)} \le U_\sigma,
\]
and these simplices tile the polytope without overlap.
\end{defn}

The linear extensions are enumerated by \code{\_enumerate\_linear\_extensions}
(\file{final\_integral.py:806}), a recursive Kahn-style topological-sort generator
that yields every extension in deterministic order. The common scalar lower bound
$L$ is resolved by \code{\_causal\_poset\_consistent\_scalar\_lower}
(\file{final\_integral.py:857}); when no variable carries an explicit lower bound,
a physical-margin fallback sets $L = (\text{earliest external time}) -
\code{POSET\_PHYSICAL\_MARGIN}$ with \code{POSET\_PHYSICAL\_MARGIN} $= 50$.

\begin{gotcha}
Why the physical margin $50$ rather than the polygon's $200$? Because
\code{POLYGON\_BBOX\_CAP}$=200$ is too loose for the chain simplex: combined with
retarded poles ($|\operatorname{Re}\beta| > 3$ when three or four poles sum), the
factor $\ee^{\beta\cdot L}$ at $L=-200$ overflows the $600$ guard and bails to
scipy. The poset path uses the tighter margin $50$ --- several correlation times
beyond the earliest external time, where the retarded integrand is already
negligible --- so the analytic path actually completes. The comment block at
\file{final\_integral.py:290} records this trade-off in detail.
\end{gotcha}

\subsection{The chain-simplex closed form}

The nested exponential integral over one chain
$\{L\le s_1\le\cdots\le s_N\le U\}$ has a closed form obtained by integrating
inside-out. The reference implementation is \code{\_exp\_over\_chain\_simplex}
(\file{final\_integral.py:910}). Its derivation, quoted from the docstring:

\begin{quote}\small
Integrate inside-out. After each step the running integrand is a sum of terms
$C\cdot\exp(\beta\cdot s_{\text{outer}} + \text{const})$. Each inner integration
produces two new terms --- one ``upper-bound'' piece (whose $\beta$ merges with the
next-outer variable's coefficient) and one ``lower-bound'' piece (a numerical
prefactor $\ee^{\beta\cdot L}$ falls out). After $N$ steps the sum has $2^N$
constant terms; their sum is the integral.
\end{quote}

\noindent Two boundary behaviours are essential:

\begin{itemize}
  \item An \textbf{empty chain} ($U \le L$) returns $0$ --- the domain has zero
        measure. This is precisely the $\tau < 0$ situation where the chain-bottom
        lower bound (from a leaf at $t=0$) and the chain-top upper bound (from a
        leaf at $t=\tau<0$) cross. Causality makes the diagram vanish there, and
        the code returns exactly $0$.
  \item A \textbf{degenerate cumulative coefficient} ($|\beta| < \epsilon$) makes
        the closed form's $1/\beta$ singular; the function returns \code{None} so
        the caller switches to the polynomial-prefactor variant below. There is
        also the same bilateral $600$ overflow guard as in the polygon path.
\end{itemize}

Two variants extend this base case:

\begin{description}
  \item[\code{\_exp\_over\_chain\_simplex\_polynomial}
        (\file{final\_integral.py:1339})] When a level's cumulative $\beta$
        vanishes, the antiderivative is a \emph{polynomial} in $s$ rather than an
        exponential. This variant carries a polynomial (in the basis
        $(s-\text{lower})^k$) through the subsequent levels, using the closed form
        $\int (s-\text{lower})^k\,\ee^{\beta s}\,\dd s$. It returns \code{None}
        only on overflow, never on degenerate $\beta$ --- it is the engine's answer
        to the $1/\beta$ singularity.
  \item[\code{\_chain\_with\_intermediate\_uppers}
        (\file{final\_integral.py:1489})] The base chain assumes the only scalar
        upper bound sits at the top of the chain. When scalar uppers appear at
        \emph{intermediate} positions, this function computes a non-decreasing
        ``effective upper'' per position, groups them into levels, and enumerates
        monotone ``cut tuples'' marking where the chain crosses each lower upper
        bound, multiplying the chain-simplex value of each non-empty piece. It
        subsumes an older, cruder ``maximality bail.''
\end{description}

\subsection{The numba fast path}

For $m\ge 3$ the $2^N$ term expansion is the dominant cost, so it is also
reimplemented for speed. \code{\_exp\_over\_chain\_simplex\_numba\_core}
(\file{final\_integral.py:1219}) is a \code{numba}-compiled translation operating
on pre-allocated complex buffers.

\begin{defn}
\textbf{Numba} is a just-in-time (JIT) compiler for a subset of Python and NumPy.
You decorate a function with \code{@numba.njit} (``no-python JIT'') and the first
time it runs, numba compiles it to native machine code through LLVM --- typically
$10$--$100\times$ faster than the interpreted version for tight numeric loops. The
\code{cache=True} option persists the compiled code to disk so the compilation
cost is paid only once across runs.
\end{defn}

The numba core returns a \code{(status, value)} pair --- status $0$ for ok, $1$
for degenerate $\beta$, $2$ for overflow --- so the pure-Python caller can take the
same fallback branches it would for the reference implementation. The import is
guarded, and there is a master switch:

\begin{lstlisting}
try:
    import numba as _numba
    _HAVE_NUMBA = True
except ImportError:
    _HAVE_NUMBA = False
    _numba = None

USE_NUMBA_CHAIN_SIMPLEX = True
\end{lstlisting}

\noindent The dispatcher \code{\_exp\_over\_chain\_simplex\_fast}
(\file{final\_integral.py:1303}) prefers the numba core when it is available and
enabled, and otherwise uses the pure-Python reference --- which is guaranteed
bit-compatible on converged results. If numba is not installed, nothing breaks;
the engine is simply slower on $m\ge 3$ diagrams.

\section{The scipy fallback layer}\label{sec:scipy}

Whenever an analytic backend returns \code{None} --- degenerate, overflowing,
mixed-constraint, or unbounded-divergent --- the subset closure routes to a
numerical-quadrature fallback built on SciPy. The dispatcher
\code{\_integrate\_polytope} (\file{final\_integral.py:4372}) branches on $m$ and
forwards to \code{\_integrate\_1d\_polytope}, \code{\_integrate\_2d\_polytope}, or
\code{\_integrate\_nd\_polytope}, all built on \code{scipy.integrate.quad} (1D)
and \code{nquad} (nested multi-D). Real and imaginary parts are integrated
separately by \code{\_complex\_quad} (\file{final\_integral.py:4520}), since
\code{quad} works on real-valued integrands. The accuracy is tuned by the
module-level \code{QUAD\_OPTS = \{'limit': 200\}}.

This fallback layer has three subtle correctness requirements, each of which has
its own regression test:

\begin{description}
  \item[The Heaviside filter.] The polytope \emph{bounds} passed to \code{nquad}
        can overshoot the true causal region (the cap fallbacks and deferred
        cross-axis constraints are loose). For retarded poles the smooth
        propagator \emph{grows} for $\Delta t < 0$, so an overshooting box would
        let the integrand add a large spurious positive contribution.
        \code{\_make\_heaviside\_filtered\_integrand}
        (\file{final\_integral.py:4411}) wraps the integrand to return $0$
        whenever any $\Delta t_e \le 0$, enforcing the $\Theta(0)=0$ convention.
  \item[Deferred constraints.] In \code{\_integrate\_nd\_polytope}
        (\file{final\_integral.py:4805}) each axis's bound function must
        \emph{skip} any constraint that still couples to a more-inner axis ---
        otherwise the bound would depend on a variable not yet integrated. This is
        the deferred-constraint correctness handled by \code{\_make\_bound\_fn},
        guarded by the test
        \code{test\_phase\_J\_nd\_polytope\_preserves\_deferred\_constraints}.
  \item[The degenerate-empty sentinel.] \code{\_resolve\_1d\_bounds}
        (\file{final\_integral.py:4550}) returns a degenerate $(0.0, 0.0)$ --- not
        $(\infty, -\infty)$ --- when the constraints are infeasible. The reason is
        sharp: \code{scipy.quad(f, +inf, -inf)} returns the \emph{sign-flipped}
        full-line integral, which would silently poison any outer \code{nquad}.
        A zero-width interval correctly integrates to $0$.
\end{description}

\begin{gotcha}
A specific trap in \code{\_integrate\_2d\_polytope}: a constraint with both
$a_{\text{int}}[0]\approx 0$ and $a_{\text{int}}[1]\approx 0$ is a
\emph{pure-external} inequality (a condition on $\tau$ alone), not a bound on
$s_1$. Treating it as an $s_1$ bound sets a \code{pure\_s1\_found} flag that
skips the $\pm 200$ cap fallback and oversamples an unbounded axis --- the source
of roughly $12\%$ of the $k=4$ theory-versus-simulation overshoot. The code is
careful to distinguish the two.
\end{gotcha}

\section{Non-local kernels: noise sources and conductance vertices}\label{sec:kernels}

Two vertex types break the pure rational-propagator structure and need extra
integration variables. Both are recognised by importing the vertex classes
\code{NoiseSourceType} and \code{ConvVertexType} from \code{msrjd.core.vertices}.

\paragraph{\code{NoiseSourceType} (coloured-noise cumulants).} The source's legs
sit at \emph{independent} times coupled by a noise kernel $\kappa^{(n)}(\tau)$.
The engine adds one $\tau$ per noise vertex, routes one leg to
$\text{anchor} - \tau$, and substitutes the kernel $\kappa(\tau)$ into the
prefactor. Because $\kappa$ is generally \emph{not} a sum of exponentials --- a
Gaussian $\ee^{-\tau^2/2\sigma^2}$, say --- these diagrams cannot stay on the
analytic path and are forced through the symbolic-plus-\code{scipy.nquad} route.

\paragraph{\code{ConvVertexType} (conductance/synaptic kernels $g(\tau)$).} Same
scaffold, but with a rescue. If $g(\tau)$ \emph{single-exponential-extracts} ---
i.e. it is of the form $C\,\ee^{\lambda\tau}$ --- the engine can absorb it. The
extraction is \code{\_extract\_exp\_mode} (\file{final\_integral.py:203}), which
uses the log-derivative trick: $\lambda = \tfrac{\dd}{\dd\tau}\log g\big|_{\tau=0}$
and $C = g(0)$. A successfully extracted kernel becomes a synthetic
\textbf{pseudo-edge} $(C, \lambda)$ with constraint $\Delta t = +\tau > 0$, and
the diagram stays on the fast analytic path.

\begin{gotcha}
\code{\_extract\_exp\_mode} returns \code{None} for \emph{polynomial-prefactor}
kernels --- the classic example being the ``alpha'' synaptic kernel
$\tau/\tau_g^2\cdot\ee^{-\tau/\tau_g}$, for which $g(0)=0$. Such kernels need a
multi-mode decomposition the single-exponential extractor cannot provide, so they
fall back to the slow path. The single-exponential conductance kernel
$\tau_g\,\ee^{-\tau/\tau_g}$ is fine; the alpha kernel is not.
\end{gotcha}

\begin{gotcha}
The extra $\tau$ variables are capped to $\pm\code{TAU\_KERNEL\_CAP}$ ($=50$). This
is a quadrature band-aid for the slow path: without the cap, \code{scipy.quad}
over $(\text{retard}_L, +\infty)$ lets the standard $\tan$-substitution compress a
Gaussian kernel's peak against the boundary, where the adaptive sampler
intermittently misses it and produces spurious spikes. Conductance $\tau$ uses a
lower bound of $0$ (causality); noise $\tau$ uses the symmetric $\pm$ cap.
\end{gotcha}

\section{Batching and parallelism}\label{sec:parallel}

The orchestration layer (\file{pipeline.py}) builds, on top of the per-diagram
callables, two batched evaluators for sweeping a $\tau$ grid:

\begin{itemize}
  \item \code{total\_C\_batch(tau\_points, parallel=True, ...)}
        (\file{pipeline.py:374}, a closure inside \code{compute\_correction\_td})
        --- evaluates the full summed correlator over a grid.
  \item \code{eval\_per\_diagram\_batch(...)} (\file{pipeline.py:502}) --- returns
        the full \code{(n\_tau, n\_diagram)} grid of per-diagram contributions, for
        diagnostics.
\end{itemize}

Both can fan out over a pool of worker processes. The parallelism uses
\code{multiprocessing.Pool} with the \code{fork} start method --- and that choice
is forced, for a reason that also makes it dangerous.

\begin{defn}
\textbf{Fork vs.\ spawn.} These are the two ways the \code{multiprocessing} module
creates worker processes. \code{fork} (POSIX only) copies the parent's entire
memory image into each child --- the children inherit every live object for free.
\code{spawn} starts a fresh Python interpreter and \emph{pickles} the arguments
across the process boundary.
\end{defn}

Fork is \emph{mandatory} here because the per-diagram \code{contribution}
callables are nested closures (defined inside
\code{integrate\_diagram.<locals>.contribution}), and the standard-library
\code{pickle} cannot serialise a nested closure. With \code{fork}, the workers
inherit those closures through the copied memory image; a module-level
\code{\_WORKER\_STATE} dict carries them in. The module-level worker entry points
\code{\_worker\_eval\_total\_C} (\file{pipeline.py:174}) and
\code{\_worker\_eval\_one\_diagram} (\file{pipeline.py:186}) are themselves
picklable (they are top-level functions) and simply look up the inherited
closures.

There are two parallel strategies, auto-chosen by batch shape: per-$\tau$
parallelism (each worker evaluates the full \code{total\_C} at one $\tau$) when
there are at least \code{n\_workers} grid points, and per-$(\tau, \text{diagram})$
nested parallelism otherwise --- the latter aggregates results in ascending
diagram order so that floating-point accumulation is bit-identical to the serial
path.

\begin{gotcha}
\textbf{Fork can crash the OS inside a notebook.} Forking a process \emph{after}
the Cocoa/BLAS libraries have initialised --- which is the state of any macOS
Jupyter kernel --- has hard-crashed the development machine twice. The environment
variable \code{OBJC\_DISABLE\_INITIALIZE\_FORK\_SAFETY=YES} makes it \emph{worse}:
it removes the guard rail, not the hazard. Phase~J therefore consults
\code{\_fork\_unsafe\_in\_notebook} (\file{pipeline.py:146}) --- which is true iff
the start method is \code{fork} \emph{and} the platform is macOS \emph{and} the
process is running inside a ZMQ/Jupyter kernel --- and silently degrades to serial
evaluation with a one-time warning. The bit-identity tests
\code{test\_phase\_J\_total\_C\_batch\_*} pass under \emph{pytest}, which is a plain
interpreter and never a ZMQ kernel, so they certify numerical determinism but
\emph{not} notebook crash-safety. The lesson, learned twice: in a notebook on a
Mac, keep \code{parallel=False}.
\end{gotcha}

\begin{gotcha}
Windows has no \code{fork}, and the per-diagram closures are unpicklable, so
\code{total\_C\_batch(parallel=True)} raises \code{ValueError: cannot find context
for 'fork'} on Windows. The parallel batch path is POSIX-only; on Windows, run
serially.
\end{gotcha}

\section{$\delta$-spike discretisation}\label{sec:deltaspike}

Recall from \S\ref{sec:polytope} that a $\delta$-edge equation with no
integration variable to solve for becomes a constraint among external times --- a
distributional shot-noise spike $\delta(a\cdot\tau + c)$. These are collected
into the result's \code{delta\_contributions} list rather than folded into the
smooth callable, because a $\delta$-function has no value to evaluate pointwise;
it must be deposited onto a grid.

\code{eval\_delta\_contributions\_on\_tau\_grid} (\file{final\_integral.py:3828})
does this for a 1D $\tau$ grid. For each contribution with equality
$a\cdot x + c = 0$ it solves for the firing lag $\tau_{\text{fire}} =
-c'/a_{\text{vary}}$, evaluates the (compiled) coefficient callable there, checks
the retardation half-spaces, and deposits $\text{coeff}/|a_{\text{vary}}|/\dd\tau$
into the nearest bin --- the $1/\dd\tau$ being the discrete stand-in for the
$\delta$'s unit area.

\begin{gotcha}
There is a degenerate case the code handles explicitly: when $a_{\text{vary}}
\approx 0$ \emph{and} $c'\approx 0$, the $\delta$ fires along the \emph{whole}
slice rather than at a single point. The code then deposits the
\emph{continuous} coefficient $\text{coeff}_{\text{fc}}(\tau)$ with \emph{no}
$1/\dd\tau$ divisor --- treating it as a smooth contribution along the slice, not a
point spike.
\end{gotcha}

The 2D analogue \code{eval\_delta\_contributions\_on\_2d\_grid}
(\file{final\_integral.py:4004}) sweeps the $\delta$-line
$a_1\tau_1 + a_2\tau_2 + c = 0$ along whichever axis carries the larger
coefficient, depositing into the nearest cell.

\section{External tools used}\label{sec:tools}

This subsystem straddles SageMath and the scientific-Python stack. Each library
appears for a specific, narrow reason; here is the inventory.

\subsection{SageMath (\code{sage.all})}

SageMath is a large open-source mathematics system that wraps dozens of
specialised libraries (Maxima for symbolic algebra, PARI, GMP, and more) behind a
unified Python API. Phase~J uses it only as the \emph{authoring and extraction}
layer --- the diagram's propagator entries, the $\delta$-coefficients, and the
prefactor all arrive as Sage symbolic expressions. The relevant objects:

\begin{description}
  \item[\code{SR} --- the Symbolic Ring.] \code{SR.var('x')} makes a symbolic
        variable; \code{SR(expr)} coerces a value into a symbolic expression.
        Symbolic expressions support \code{.subs(dict)}, \code{.expand()},
        \code{.diff()}, \code{.coefficient(v)}, and \code{.variables()}.
  \item[\code{CDF} --- the Complex Double Field.] Machine-precision complex
        numbers. \code{complex(CDF(SR(expr)))} is the standard idiom for turning a
        fully-substituted symbolic scalar into a Python \code{complex}.
  \item[\code{fast\_callable(expr, vars, domain=CDF)}.] Sage's JIT compiler
        \emph{for symbolic expressions}. It walks the expression tree once and
        produces a fast callable (backed by a bytecode interpreter) that evaluates
        \code{expr} numerically. Phase~J uses it to compile the residual symbolic
        integrand for the \code{scipy.nquad} fallback path only.
  \item[\code{solve} (imported as \code{sage\_solve}).] The symbolic equation
        solver, used on the $\delta$-edge equations $\Delta t_e == 0$ to solve for
        an integration variable.
\end{description}

\noindent The exact import (\file{final\_integral.py:83}):

\begin{lstlisting}
from sage.all import SR, fast_callable, CDF, solve as sage_solve
\end{lstlisting}

The discipline is: extract the numerical pole/residue data \emph{once}, solve the
$\delta$ equations, read linear coefficients off the constraint expressions --- and
then never touch Sage again. The hot numerical loops are pure Python
\code{complex} and \code{cmath}, because Sage is too slow per call.

\subsection{\code{cmath} / \code{math} (standard library)}

\code{cmath.exp} is the complex exponential at the core of every analytic
integrator's inner loop; \code{math.inf}, \code{math.isinf}, and
\code{math.factorial} appear in the bound resolution and the polynomial-prefactor
chain. These are imported \emph{locally} inside the hot functions (e.g.
\code{import cmath}) so the module-level namespace stays clean.

\subsection{NumPy and Numba}

\term{NumPy} (\code{import numpy as np}, \file{final\_integral.py:1041}) plays two
roles: it backs the numba chain-simplex kernel, which works on pre-allocated
\code{np.complex128} buffers, and it provides the output arrays (and
\code{np.argmin}/\code{np.abs}) for the $\delta$-spike grid discretisers.
\term{Numba} is the JIT compiler for the chain-simplex inner loop covered in
\S\ref{sec:m3}; its import is guarded so the engine runs without it.

\subsection{mpmath}

\term{mpmath} is a pure-Python arbitrary-precision floating-point library;
\code{mpmath.mpc} is a complex number with a configurable number of decimal
digits (\code{mp.dps}). Phase~J uses it in two \emph{experimental,
off-by-default} precision-rescue paths for the $m\ge 3$ chain simplex, where the
$2^N$ cancellation can lose precision on close-paired poles. Both
\code{USE\_CHAIN\_SIMPLEX\_PRECISION\_FIX} and
\code{USE\_POSET\_MPMATH\_ACCUMULATION} default to \code{False} --- they did not
actually cure the targeted bug (see \S\ref{sec:knownbugs}).

\subsection{SciPy}

\term{SciPy}'s \code{integrate} submodule provides adaptive numerical quadrature:
\code{quad(f, a, b)} for adaptive 1D integration (Gauss--Kronrod / QUADPACK,
handling $\pm\infty$ bounds natively), and \code{nquad(f, [bounds...])} for nested
multi-dimensional quadrature where each axis's bounds may be a callable of the
outer variables. It is the engine's universal fallback, reached whenever an
analytic integrator returns \code{None}.

\subsection{What is \emph{not} used}

NetworkX, nauty, and SymPy do \emph{not} appear in this subsystem. Graph
structure is handled through SageMath's own graph objects
(\code{D.vertices()}, \code{D.edges()}, \code{D.num\_edges()}); nauty (graph
canonicalisation) lives in the upstream enumeration subsystem; and all symbolic
work here is Sage's \code{SR}, never SymPy.

\section{Diagnostics and tuning knobs}\label{sec:diagnostics}

The engine exposes a handful of module-level knobs and a runtime-counter
mechanism, all so you can see \emph{which} path actually fired and tune the
trade-offs.

\begin{description}
  \item[\code{\_RUNTIME\_COUNTERS} (\file{final\_integral.py:315}).] A diagnostic
        dict counting, per $m$, how often each analytic path was \emph{attempted},
        how often it \emph{returned \code{None}}, and how often the run fell
        through to \code{scipy.nquad}. Call \code{\_reset\_runtime\_counters()}
        before a timed run, then read the dict after. These are \emph{runtime}
        observations; the \emph{intent} of each subset (which path it was
        \emph{set up} to take) is recorded separately in
        \code{subset\_diagnostics['evaluator']}.
  \item[\code{QUAD\_OPTS = \{'limit': 200\}}.] The accuracy knob for the scipy
        fallback's quadrature.
  \item[\code{TAU\_KERNEL\_CAP = 50.0}.] The cap on the extra non-local-kernel
        $\tau$ variables (\S\ref{sec:kernels}).
  \item[\code{POLYGON\_BBOX\_CAP = 200.0} vs.\ \code{POSET\_PHYSICAL\_MARGIN = 50.0}.]
        The bounding-box half-widths for the $m=2$ and $m\ge 3$ paths respectively
        (\S\ref{sec:m2}, \S\ref{sec:m3}).
  \item[The \code{USE\_*} flags.] \code{USE\_NUMBA\_CHAIN\_SIMPLEX} (default
        \code{True}) toggles the numba fast path; \code{USE\_POLYGON\_M2\_INTEGRATOR}
        and \code{USE\_1D\_INTEGRATOR} toggle the analytic $m=2$ and $m=1$ paths;
        and the precision-rescue flags
        \code{USE\_CHAIN\_SIMPLEX\_PRECISION\_FIX},
        \code{USE\_POSET\_MPMATH\_ACCUMULATION}, and
        \code{USE\_POSET\_CAP\_MATCH\_SCIPY} all default \code{False}.
\end{description}

\section{Known precision caveats and open bugs}\label{sec:knownbugs}

Honesty about the limits of the engine matters more than a clean narrative. There
is one genuinely open bug and a cluster of near-misses around it.

\begin{gotcha}
\textbf{The $m\ge 3$ chain-simplex close-pole overestimate (OPEN).} On
close-paired poles --- e.g.\ the spike-reset poles $0.329\ii$ and $0.351\ii$, a
difference of only $0.022\ii$ --- the $2^N$ closed form suffers
cancellation-driven precision loss. Empirically this shows up as roughly a
$4\times$ aggregate overestimate of the one-loop value at spike-reset $k=2$,
$\ell=1$. Two attempted fixes are \emph{both off by default} because neither
actually cured it: \code{USE\_CHAIN\_SIMPLEX\_PRECISION\_FIX} (an mpmath dispatch
on close poles) and \code{USE\_POSET\_MPMATH\_ACCUMULATION} (mpmath outer
accumulation, which improved the cancellation factor only $\sim 21\times$). The
audit (\file{docs/m\_ge3\_precision\_bug\_audit.md}) now points at a
\emph{bookkeeping difference in pole-tuple construction between the grouped and
per-diagram paths} --- by linearity they should agree, so the $4\times$ discrepancy
points to a construction bug, not an arithmetic-precision one. This is
unresolved.
\end{gotcha}

The related near-miss flag \code{USE\_POSET\_CAP\_MATCH\_SCIPY} was meant to align
the analytic poset's unbounded-below fallback ($L = \text{earliest\_ext} - 50$)
with the scipy \code{OUTER\_CAP} of $200$, for marginally-stable theories where
the integral is genuinely cap-dependent. The code comment records that aligning
the caps does \emph{not} bring the per-diagram path into agreement with the
grouped path --- confirming the dominant error lies elsewhere. It too defaults
\code{False}.

\section{The loop-subgraph stub}\label{sec:subgraph}

For completeness: \file{subgraph.py} is a near-stub for a route that was designed
but never built. Its \code{identify\_loop\_subgraphs} (\file{subgraph.py:66})
returns \code{[]} for tree-level diagrams (empty \code{free\_freqs}) and
\emph{raises \code{NotImplementedError}} otherwise. The frequency-domain
loop-reduction algorithm this file was meant to enable was superseded: the
current \code{integrate\_diagram} handles loops directly in the time domain,
treating every internal vertex as an integration variable regardless of loop
order. The module is never exercised on the live path; its docstring documents
only the \emph{planned} algorithm. If you go looking for ``where the loop
reduction happens,'' the answer is that there is no separate loop reduction ---
the DAG-and-polytope machinery of this chapter \emph{is} the loop handling.

\begin{note}
The same point bears on a stale docstring you may notice while reading the code:
\file{pipeline.py}'s module header and the \code{compute\_correction\_td}
docstring both say loop diagrams are \emph{skipped}. They are not. The live code
evaluates every loop order through \code{integrate\_diagram}; the inline comment at
\file{pipeline.py:291} explains why (the enumerator produces DAGs, so tree and
loop share the same vertex-time integration algorithm). Trust the code, not the
header.
\end{note}

\section{What you just learned}

Phase~J turns each typed diagram into a number by recognising that the retarded
propagator is a finite sum of decaying exponentials \eqref{eq:GR-decomp}, so the
diagram integral \eqref{eq:diagram-integral} is a sum --- over $\delta$-subsets
and pole tuples --- of pure exponentials integrated over a causal polytope
\eqref{eq:exp-over-polytope}. Three analytic backends evaluate that polytope
integral in closed form, dispatched by the surviving dimension $m$: an interval
formula for $m=1$ (\S\ref{sec:m1}), Sutherland--Hodgman clipping plus a
unit-triangle formula for $m=2$ (\S\ref{sec:m2}), and a causal-poset
linear-extension decomposition into chain simplices for $m\ge 3$
(\S\ref{sec:m3}). A \code{scipy} quadrature layer catches every degeneracy
(\S\ref{sec:scipy}), non-local noise and conductance kernels add extra
$\tau$ variables (\S\ref{sec:kernels}), distributional spikes are deposited onto
the grid separately (\S\ref{sec:deltaspike}), and a fork-based batch layer ---
with a hard-won notebook safety guard --- sweeps the $\tau$ grid
(\S\ref{sec:parallel}).

The natural next step is Chapter~\ref{ch:grouped-phasej}, the \emph{grouped}
Phase~J path. It reuses the very same analytic primitives ---
\code{\_integrate\_1d\_polytope\_modesum}, \code{\_integrate\_2d\_polygon\_modesum},
and the poset integrator --- but sums residues \emph{across} diagrams before
integrating, trading per-diagram independence for a reduction in the pole-tuple
count. The discrepancy between that path and this one is, as
\S\ref{sec:knownbugs} noted, the current prime suspect for the $m\ge 3$ open bug.

% ---------------------------------------------------------------------
%  Chapter 13 : Grouped Phase J --- the fully-analytic mode-sum
% ---------------------------------------------------------------------
\chapter{Grouped Phase J: the Fully-Analytic Mode-Sum}\label{ch:grouped-phasej}

The previous chapter walked the \emph{per-diagram} time-domain integrator: for
each typed Feynman diagram, Phase~J takes that diagram's product of retarded
propagators and vertex coefficients and integrates over the internal vertex
times to produce a number $C(t_1, \dots, t_k)$. This chapter is about an
\emph{optimisation} of that loop --- one that is mathematically exact, not an
approximation --- together with the analytic machinery (``Stage~4a'') that
replaces adaptive numerical quadrature with closed-form mode-sums and so lifts
the precision floor from scipy's $\sim\!10^{-8}$ all the way to machine epsilon.

The motivation is brute arithmetic. The enumerator does not emit one diagram per
topology; it emits one diagram per \emph{typed} topology, and a single bare
topology fans out into hundreds of typed copies that differ only in bookkeeping
the integrator can factor out. The grouped path recognises that redundancy, does
the expensive per-topology setup \emph{once}, and sums the integrands of the
whole family before integrating. The result is bit-for-bit the same numbers the
per-diagram path produces (a contract the regression suite enforces to
round-off), produced far faster and, on the analytic path, far more precisely.

Two source files carry the whole story, and we will quote both with
\code{file:line} throughout:
\begin{itemize}
  \item \file{pipeline/\_grouped\_phase\_j.py} --- the \emph{grouping wrapper}
        \code{compute\_correction\_td\_grouped}, which buckets typed diagrams
        into groups and dispatches each group, with a three-tier fallback to the
        per-diagram path when grouping cannot be trusted.
  \item \file{msrjd/integration/time\_domain/grouped\_integral.py} --- the
        \emph{grouped evaluator} \code{integrate\_grouped\_diagram}, which is
        where the sum-then-integrate and the merged-residue mode-sum actually
        happen.
\end{itemize}

\begin{note}
This is a \term{prototype}, opt-in by a single flag. It is \emph{not} wired into
the engine by default: \code{compute\_cumulants} takes a keyword
\code{use\_grouped\_phase\_j} that defaults to \code{False}
(\file{pipeline/compute.py:129}). Pass \code{use\_grouped\_phase\_j=True} to
route Phase~J through the grouped path. Everything downstream --- the $\tau$-grid
evaluation, the plotting, the saving --- is untouched, because the wrapper
returns the \emph{same dict shape} as the per-diagram path. The two paths are
interchangeable behind that one flag.
\end{note}

\section{Why group? The per-diagram redundancy}

To see the redundancy clearly we need two definitions that recur throughout the
chapter. They are the same objects the enumeration and type-assignment chapters
introduced, viewed now from the integrator's side.

\begin{defn}
A \term{prediagram} is the bare graph topology of a Feynman diagram: its
vertices, its \emph{leaves} (the external legs), its edges, and which internal
vertices are integration variables. It carries \emph{no} field-component
information on its edges.

A \term{typed diagram} is a prediagram with two further things fixed: per edge,
\emph{which} propagator matrix entry $(r_i, p_i)$ flows along that edge (i.e.\
which response/physical field-component pair), and the per-vertex coefficients.
\end{defn}

Because a typed diagram only \emph{adds} edge indices and coefficients to a
prediagram, all typed diagrams sharing one prediagram have \emph{identical graph
topology, identical leaves, identical internal vertices, and identical
integration variables}. They differ only in those per-edge $(r_i, p_i)$ indices
and per-vertex coefficients. The module docstring states exactly this
(\file{grouped\_integral.py:9--13}):

\begin{lstlisting}[style=console]
Typed diagrams sharing the same prediagram have identical graph
topology, leaves, internal vertices, and integration variables.
They differ only in per-edge propagator field-type indices (ri, pi)
and per-vertex coefficients.
\end{lstlisting}

The per-diagram integrator does not exploit this. It pays the \emph{entire}
per-diagram setup cost --- the graph walk, the allocation of vertex-time symbols,
the $2^{|E|}$ $\delta$-edge subset bookkeeping (more on that below), the
\code{fast\_callable} compile, the polytope and pole extraction --- once for
\emph{every} typed diagram. The docstring records the measurement that motivated
the whole effort (\file{grouped\_integral.py:18--21}):

\begin{lstlisting}[style=console]
For (k=2, ell=1) multipop the user saw 7736 typed diagrams across
9 prediagrams - ~860 diagrams per prediagram, all paying full setup
cost.
\end{lstlisting}

Eight hundred and sixty typed diagrams per prediagram, each repeating the same
graph walk and the same JIT compile. The setup is identical across all of them;
only a handful of complex numbers differ. That is the redundancy grouping
removes.

\section{Linearity of integration: why grouping is exact}

The mathematical justification is a single line, and it is worth dwelling on
because it is what makes grouping a \emph{drop-in} rather than a tradeoff. The
contribution of one typed diagram is an integral over the internal vertex times;
the contribution of the whole group is the sum of those integrals. Because the
integration \emph{domain} is identical across every typed diagram of one
prediagram --- the domain comes from the shared graph topology, which is the same
for all of them --- you may sum the integrands first and integrate once
(\file{grouped\_integral.py:33--41}):
\begin{equation}
  \sum_{\Gamma}\ \int \dd s_1 \cdots \dd s_m\ \mathcal{I}_\Gamma
  \;=\;
  \int \dd s_1 \cdots \dd s_m\ \sum_{\Gamma}\ \mathcal{I}_\Gamma ,
  \label{eq:linearity}
\end{equation}
where $\mathcal{I}_\Gamma$ is the time-domain integrand of typed diagram
$\Gamma$ and the $s_i$ are the surviving internal-vertex times. This is the
linearity of the integral: it is \emph{exact}, holds term by term, and is not an
approximation in any limit.

So the grouped path builds the prediagram-level scaffolding once, forms the
right-hand-side sum $\sum_\Gamma \mathcal{I}_\Gamma$, and runs a single
integration. The expensive parts (graph walk, symbol allocation, the
$2^{|E|}$ structure, the compile, the polytope) are shared; only the cheap parts
(reading off each diagram's edge residues and coefficients) repeat per diagram.

\begin{note}
\eqref{eq:linearity} is the entire conceptual content of the chapter. Everything
that follows --- the subset expansion, the polytope, the mode-sum, the merged
residue --- is just \emph{how} the right-hand side is built and integrated
efficiently. Keep the picture ``sum the integrands of a family that share a
domain, then integrate once'' in mind and the code reads straightforwardly.
\end{note}

\section{From retarded propagators to a polytope integral}

Before the grouping can buy anything, we need the shape of the integral it is
grouping. This section rebuilds it from the field theory down to the convex
polytope the closed forms integrate over. The construction is shared with the
per-diagram path; the grouped path differs only in the very last step.

\subsection{The retarded propagator in the time domain}

Phase~I (the frequency-domain assembly of the previous part) hands Phase~J a
rational propagator matrix $\hat G(\omega)$ already decomposed into poles and
residues. Concretely it passes a \code{propagator\_data} dict whose keys are:
\begin{description}
  \item[\code{pole\_vals}] the poles $p_k$ of $\hat G(\omega)$. The causality
        filter guarantees $\operatorname{Im}(p_k) > 0$ for every one of them.
  \item[\code{C\_mats}] one residue matrix $C_k$ per pole, so that
        $C_k[i,j]$ is the residue of $\hat G[i,j](\omega)$ at the pole $p_k$.
  \item[\code{D\_delta}] (optional) the $\omega \to \infty$ limits of
        $\hat G$ --- the instantaneous response weights.
\end{description}
The helper \code{build\_G\_t\_matrix} (\file{propagator\_td.py:135}) turns this
into the \term{time-domain retarded propagator}. Its docstring states the
decomposition (\file{propagator\_td.py:140--141}):
\begin{equation}
  \Gret[i,j](t)
  \;=\;
  \underbrace{\text{delta\_coeffs}[i,j]\cdot \delta(t)}_{\text{instantaneous}}
  \;+\;
  \underbrace{\Theta(t)\cdot\Bigl(\sum_k C_k[i,j]\,\ee^{\ii\,p_k\,t}\Bigr)}_{\text{smooth}} .
  \label{eq:GR}
\end{equation}
The first term is the \emph{instantaneous} $\delta$-response (it appears for any
entry whose frequency-domain propagator has a nonzero $\omega\to\infty$ limit,
e.g.\ a response source producing an immediate response at the same time); the
second is the \emph{smooth} pole--residue sum. With the pipeline-wide Fourier
convention, each summand $C_k\,\ee^{\ii p_k t}$ \emph{decays} for $t>0$ and grows
for $t<0$, and the Heaviside $\Theta(t)$ is exactly what makes the retarded
propagator well-defined on the real line (\file{propagator\_td.py:157--161}).

It is convenient to fold the $\ii$ into the rate. Define the \term{mode}
convention the analytic code uses throughout:
\begin{equation}
  \lambda_k \;=\; \ii\,p_k ,
  \qquad
  C_\alpha \;=\; C_k[p_i, r_i] ,
  \label{eq:mode}
\end{equation}
so $\operatorname{Re}\lambda_k < 0$ for a decaying retarded mode, and the smooth
time-domain propagator on an edge reads $\sum_\alpha C_\alpha\,
\ee^{\lambda_\alpha\,\Delta t}$. Here $\Delta t$ is the time difference across the
edge: for an oriented edge $u \to v$, $\Delta t = t_v - t_u$. (The dataclass
\code{EdgeModeSum} records this convention in its own docstring,
\file{final\_integral.py:128--133}, with the index ordering
$G[p_i, r_i] = \langle \varphi_{p_i}\,\tilde n_{r_i}\rangle$ --- physical row,
response column.)

\subsection{One diagram's integral and the $2^{|E|}$ subset expansion}

Give each internal vertex a time --- the integration variables $s_v$ --- and each
leaf its external time, with one leaf (the \term{origin leaf}) pinned to $t=0$.
The contribution of one typed diagram is then
\begin{equation}
  C_\Gamma(t_1,\dots,t_k)
  \;=\;
  (\text{prefactor})\cdot
  \int \prod_v \dd s_v\ \prod_{\text{edges } e} \Gret[p_i, r_i](\Delta t_e) .
  \label{eq:onediagram}
\end{equation}
Each retarded factor is the sum of a $\delta$-part and a $\Theta$-times-smooth
part, \eqref{eq:GR}. Expand the product of $|E|$ such factors over the binary
choice ``$\delta$-piece versus smooth-piece'' on each edge: that gives a sum over
$2^{|E|}$ \term{subsets}. In a given subset, the edges chosen as
``$\delta$-edges'' each carry a $\delta(\Delta t_e)$, which \emph{collapses one
integration variable} via the linear equality $\Delta t_e = 0$; the remaining
``smooth-edges'' each carry $\sum_\alpha C_\alpha \ee^{\lambda_\alpha \Delta t_e}$
together with the retardation constraint $\Delta t_e > 0$ coming from the
Heaviside.

After $\delta$-elimination, suppose $m$ integration variables $s_1,\dots,s_m$
survive. The smooth retardation constraints $\Delta t_e > 0$ are \emph{linear
inequalities} in the $s_i$ and the free external times --- they cut out a
\term{convex polytope} $P$. The subset's contribution is
\begin{equation}
  I_{\text{subset}}
  \;=\;
  (\text{prefactor})\cdot
  \int_P \dd s_1 \cdots \dd s_m\
  \prod_{e\,\in\,\text{smooth}}
  \Bigl[\,\sum_\alpha C^{(e)}_\alpha\, \ee^{\lambda_\alpha \Delta t_e}\,\Bigr] .
  \label{eq:subset}
\end{equation}

\begin{defn}
The integer $m$ --- the number of integration variables surviving
$\delta$-elimination in a subset --- is the single most important quantity in the
evaluator. It selects which closed form is used: $m=0$ (no integral, just a
constraint check), $m=1$ (a 1-D interval), $m=2$ (a polygon), or $m\ge 3$ (a
general polytope). We return to each case in \S\ref{sec:closedforms}.
\end{defn}

\subsection{Pole expansion into mode-sums}

Expand the product of mode-sums in \eqref{eq:subset} into a sum over \term{pole
tuples} $\alpha = (\alpha_1, \dots, \alpha_{n_{\text{smooth}}})$ --- one pole
index per smooth edge. Each tuple contributes a single exponential whose exponent
is \emph{linear} in the $s_i$:
\begin{equation}
  \prod_e \Bigl[\sum_\alpha C_\alpha \ee^{\lambda_\alpha \Delta t_e}\Bigr]
  \;=\;
  \sum_{\bm\alpha}\ \Bigl(\prod_e C_{\alpha_e}\Bigr)\,
  \exp\!\Bigl(\sum_e \lambda_{\alpha_e}\,\Delta t_e\Bigr) .
  \label{eq:poletuple}
\end{equation}
Writing each edge's time difference as a linear form
$\Delta t_e = c_{0,e} + \sum_i a^{\text{int}}_{e,i}\,s_i + \sum_j a^{\text{ext}}_{e,j}\,t^{\text{free}}_j$,
the exponent of \eqref{eq:poletuple} becomes $\gamma_{\bm\alpha} + \sum_i
\alpha^{(i)}_s\,s_i$ with
\begin{equation}
  \alpha^{(i)}_s \;=\; \sum_e \lambda_{\alpha_e}\,a^{\text{int}}_{e,i},
  \qquad
  \gamma_{\bm\alpha} \;=\; \sum_e \lambda_{\alpha_e}\Bigl(c_{0,e} + \sum_j a^{\text{ext}}_{e,j}\,t^{\text{free}}_j\Bigr).
  \label{eq:slopes}
\end{equation}
So the subset contribution becomes a sum over pole tuples, each a constant times
a geometric integral of an exponential-of-linear over the polytope:
\begin{equation}
  I_{\text{subset}}
  \;=\;
  \sum_{\bm\alpha}\Bigl(\prod_e C_{\alpha_e}\Bigr)\,\ee^{\gamma_{\bm\alpha}}
  \int_P \exp\!\Bigl(\sum_i \alpha^{(i)}_s\,s_i\Bigr)\,\dd s .
  \label{eq:subsetfinal}
\end{equation}
The integral $\int_P \exp(\text{linear})\,\dd s$ has a closed form that depends
only on $m$ and the polytope geometry --- it does \emph{not} know which pole tuple
or which diagram produced it. That separation is what the grouped path exploits.

\section{The merged-residue tensor \texorpdfstring{$B_{\bm\alpha}$}{B-alpha}}

Now apply linearity, \eqref{eq:linearity}, to \eqref{eq:subsetfinal}. Take one
prediagram with $N$ typed diagrams indexed by $j$. All $N$ share the same pole
spectrum $\{\lambda_\alpha\}$ (they share \code{propagator\_data}) and the same
polytope (the constraints come from the shared $\Delta t_e$ symbols). Only the
residues $C_\alpha$ and the prefactor $\mathrm{cp}_j$ differ between typed
diagrams --- through their per-edge $(r_i, p_i)$ indices. Summing
\eqref{eq:subsetfinal} over $j$:
\begin{equation}
  I_{\text{subset}}^{\text{group}}
  \;=\;
  \sum_{\bm\alpha}
  \underbrace{\Bigl[\sum_j \mathrm{cp}_j
    \Bigl(\textstyle\prod_{\delta\text{-edge}} \delta_{\mathrm{coeff}\,j,e}\Bigr)
    \prod_{\text{smooth }e} C^{(j)}_{\alpha_e,e}\Bigr]}_{\displaystyle B_{\bm\alpha}}
  \ee^{\gamma_{\bm\alpha}}
  \int_P \exp\!\Bigl(\sum_i \alpha^{(i)}_s s_i\Bigr)\dd s .
  \label{eq:merged}
\end{equation}
Crucially, the $\ee^{\gamma_{\bm\alpha}} \int_P(\cdots)$ factor does \emph{not}
depend on $j$. So the code pre-sums the residue weights over $j$ \emph{inside}
the pole tuple, defining the \term{merged-residue tensor}
(\file{grouped\_integral.py:71--75}):
\begin{equation}
  B_{\bm\alpha}
  \;=\;
  \sum_j\ \Bigl(\mathrm{cp}_j \cdot
    \textstyle\prod_{\delta\text{-edge}} \delta_{\mathrm{coeff}\,j,e}\Bigr)
  \cdot \prod_{\text{smooth }e} C^{(j)}_{\alpha_e,e} .
  \label{eq:Bdef}
\end{equation}

\begin{defn}
The merged-residue tensor $B_{\bm\alpha}$ is the grouped path's per-tuple
coefficient. For the \emph{per-diagram} path the per-tuple coefficient is the
Cartesian product $\prod_e C_{\alpha_e}$ produced by \code{\_enumerate\_pole\_tuples}
(\file{final\_integral.py:529}). The grouped path replaces that single-diagram
product with the sum $B_{\bm\alpha}$ over the whole group. The underlying analytic
integration --- the $\ee^\gamma \int_P$ factor --- is byte-for-byte identical; only
the residue accumulation moved inside the tuple sum
(\file{grouped\_integral.py:79--86}).
\end{defn}

This is the precise sense in which grouping is a drop-in. The per-diagram analytic
evaluators already accept the per-tuple coefficient through an iterator; the
grouped path feeds them the merged $B_{\bm\alpha}$ instead of the Cartesian
product, and the same closed form runs unchanged.

\begin{gotcha}
There is one subtlety, the subject of an open precision investigation, and it
lives entirely in this reordering. The grouped path sums $B_{\bm\alpha}$
\emph{before} the integration, so intra-group cancellations happen at the level
of (possibly large) individual residue products. The per-diagram path sums
\emph{after} integration. By linearity these are equal in exact arithmetic, but
in IEEE float64 they can disagree when individual pole-tuple terms are huge ---
close-paired poles give $O(1/\Delta p)$ residues, products of several such give
$O(10^8)$, and catastrophic cancellation drags them down to an $O(1)$ integral.
The disagreement is documented at \file{final\_integral.py:1690--1700} (the
\code{USE\_POSET\_CAP\_MATCH\_SCIPY} comment block) and affects the $m\ge 3$
case; see \S\ref{sec:precision}.
\end{gotcha}

\section{Closed forms by \texorpdfstring{$m$}{m}}\label{sec:closedforms}

With $B_{\bm\alpha}$ in hand, the only thing left is the geometric integral
$\int_P \exp(\text{linear})\,\dd s$ for each subset. Its closed form is selected
by $m$. The grouped path implements $m=0$ and $m=1$ as \emph{stand-alone twins}
inside \file{grouped\_integral.py}, and routes $m=2$ and $m\ge 3$ through the
shared per-diagram evaluators in \file{final\_integral.py} via a
\code{pole\_tuples} override. All four read the merged residue
$(B_{\bm\alpha}, \text{lambdas})$ pairs the same way.

\subsection{$m=0$: a constraint check}

No integration variable survived $\delta$-elimination, so there is nothing to
integrate --- only the constraints to check and the exponentials to sum. The
helper \code{\_evaluate\_grouped\_m0\_modesum} (\file{grouped\_integral.py:138})
does exactly that:
\begin{lstlisting}
def _evaluate_grouped_m0_modesum(pole_tuples, subset_constraint_data, free_ext_vals):
    import cmath
    for (a_int, a_ext, c0) in subset_constraint_data:
        c_eff = float(c0) + sum(float(a_ext[i]) * float(free_ext_vals[i])
                                for i in range(len(a_ext)))
        if c_eff <= 0:
            return 0.0 + 0.0j
    ...
    total = 0.0 + 0.0j
    for (B_alpha, lambdas) in pole_tuples:
        gamma = 0.0 + 0.0j
        for e in range(len(lambdas)):
            gamma += lambdas[e] * dt_per_edge[e]
        if abs(gamma.real) > _GROUPED_EXP_REAL_LIMIT:
            return None
        try:
            total += B_alpha * cmath.exp(gamma)
        except (OverflowError, ValueError):
            return None
    return total
\end{lstlisting}
It returns $\sum_{\bm\alpha} B_{\bm\alpha}\,\exp(\sum_e \lambda_{\alpha_e}\Delta t_e)$
as a Python \code{complex}; it returns $0$ if any constraint $\Delta t_e > 0$ is
violated (the test \code{c\_eff <= 0} at \file{grouped\_integral.py:155}, the
$\Theta(0)=0$ convention --- the boundary is \emph{outside}); and it returns
\code{None} on overflow, which signals the caller to fall back to numerical
quadrature.

\subsection{$m=1$: an interval, with $\pm\infty$ handled exactly}

One integration variable $s$ survives, on an interval $[L,U]$. The closed form is
elementary,
\begin{equation}
  \int_L^U \ee^{\alpha_s s + \gamma}\,\dd s
  \;=\;
  \frac{\ee^{\alpha_s U + \gamma} - \ee^{\alpha_s L + \gamma}}{\alpha_s},
  \label{eq:m1}
\end{equation}
with the limit $\alpha_s \to 0$ (a constant integrand, value $(U-L)\,\ee^\gamma$)
handled separately. The interesting part is the unbounded endpoints.
\code{\_integrate\_grouped\_m1\_modesum} (\file{grouped\_integral.py:179}) first
resolves the feasible interval from the linear constraints, tracking $\pm\infty$
exactly:
\begin{lstlisting}
    L = -math.inf
    U = +math.inf
    for (a_int, a_ext, c0) in subset_constraint_data:
        a = float(a_int[0]) if a_int else 0.0
        c_eff = float(c0) + sum(float(a_ext[i]) * float(free_ext_vals[i])
                                for i in range(len(a_ext)))
        if abs(a) < 1e-15:
            if c_eff <= 0:
                return 0.0 + 0.0j
            continue
        bound = -c_eff / a
        if a > 0:
            if bound > L: L = bound
        else:
            if bound < U: U = bound
    if L >= U:
        return 0.0 + 0.0j
\end{lstlisting}
A constraint with $|a| < 10^{-15}$ is $s$-independent: if it is violated the
subset is infeasible (return $0$), otherwise it is skipped. A constraint with
$a>0$ tightens the lower bound, $a<0$ the upper bound. When a side is left
unbounded, the corresponding boundary term in \eqref{eq:m1} is \emph{zeroed} only
if the integrand decays in that direction (the sign of
$\operatorname{Re}\alpha_s$ matches); if it would instead diverge, the function
returns \code{None} to force the scipy fallback, which handles the unbounded
endpoint with an adaptive substitution.

\begin{gotcha}
Why $\pm\infty$ rather than clipping at some large finite cap? The docstring
(\file{grouped\_integral.py:196--202}) is explicit, and it is load-bearing for
precision: clipping the interval at $\pm$\code{bbox\_cap} would leave a residual
$\ee^{\alpha_s\cdot\text{bbox\_cap}} \approx 10^{-9}$ boundary term for typical
Hawkes timescales --- measured on the \code{quad\_exp\_k2\_ell0} fixture before
this fix. Treating the genuine $\pm\infty$ exactly removes that error entirely.
The \code{bbox\_cap} argument is still passed in, but on this path it is used only
to \emph{not} clip.
\end{gotcha}

\subsection{$m=2$: a polygon via fan-triangulation}

For two surviving variables the polytope is a convex polygon. The shared
evaluator \code{\_integrate\_2d\_polygon\_modesum} (\file{final\_integral.py:2175})
clips the half-planes, fan-triangulates the polygon from vertex~$0$, and
integrates $\exp(\alpha x + \beta y)$ over each triangle using the unit-triangle
closed form $J(p,q)$ (\file{final\_integral.py:356}):
\begin{equation}
  J(p,q) \;=\; \int_0^1\!\dd u \int_0^{1-u}\!\dd w\ \ee^{p\,u + q\,w} .
  \label{eq:J}
\end{equation}
Its docstring gives the derivation (do the $w$-integral first, then split and
integrate in $u$) and notes a 4th-order Taylor fallback when any of $|p|$, $|q|$,
$|p-q|$ is tiny --- the same catastrophic-cancellation guard that recurs across
the analytic code. The grouped path calls this evaluator with the merged residues
supplied as \code{pole\_tuples=grouped\_pole\_tuples}, a dummy
\code{smooth\_edge\_modes=list(dummy\_modes)} (explained in
\S\ref{sec:dummymodes}), and \code{prefactor\_complex=1.0+0.0j} --- because the
prefactor is already folded into $B_{\bm\alpha}$.

\subsection{$m\ge 3$: a causal-poset chain simplex}

For three or more variables the polytope is handled by a \term{causal-poset
chain-simplex} decomposition in \code{\_integrate\_nd\_polytope\_poset\_modesum}
(\file{final\_integral.py:1726}). The retardation constraints define a partial
order $s_u < s_v$ on the vertex times; each linear extension of that poset is a
nested simplex $L \le s_{\sigma(0)} \le \cdots \le s_{\sigma(m-1)} \le U$, and the
integrator sums $\int \exp(\sum \lambda s)$ over every nested simplex. The grouped
path uses it with the same \code{pole\_tuples} override. This is the path with the
known float64 cancellation issue (\S\ref{sec:precision}).

\begin{note}
The two flags \code{USE\_POLYGON\_M2\_INTEGRATOR}
(\file{final\_integral.py:287}) and \code{USE\_POSET\_INTEGRATOR}
(\file{final\_integral.py:1661}) gate the $m=2$ and $m\ge 3$ analytic paths
respectively. The grouped file imports the \emph{module}, not the flag values
(\code{import \ldots final\_integral as \_fi\_mod}, \file{grouped\_integral.py:119}),
so each per-call read picks up a notebook override live. Flip
\code{\_fi\_mod.USE\_POSET\_INTEGRATOR = False} in a notebook and the grouped
$m\ge 3$ analytic path silently turns off and routes to scipy.
\end{note}

\section{The grouping wrapper}

We now walk the two files top to bottom, starting with the wrapper in
\file{pipeline/\_grouped\_phase\_j.py}. Its job is to bucket typed diagrams into
groups, dispatch each group to the evaluator, and --- the part that earns its
keep --- fall back to the per-diagram path whenever the grouped result cannot be
trusted.

\subsection{The skip-status guard}

Before any logic, the module declares the set of subset-level statuses that mean
``the grouped evaluator silently \emph{dropped} this subset rather than
integrating it'' (\file{\_grouped\_phase\_j.py:49}):
\begin{lstlisting}
_SUBSET_SKIP_STATUSES = frozenset({
    'shotnoise_skipped_in_prototype',
    'dt_sym_mismatch_skipped',
})
\end{lstlisting}
This frozenset is the heart of the wrapper's correctness story. A group can return
an overall \code{status='ok'} and yet have a \emph{subset} diagnostic carrying one
of these codes --- which means a piece of the integrand was dropped and the group
total is numerically \emph{wrong}. The comment (\file{\_grouped\_phase\_j.py:42--48})
calls this ``the production bug we're fixing'' and lists the statuses explicitly
so that a \emph{new} skip status added later cannot silently regress the guard.

\begin{gotcha}
\textbf{If you add a new skip status to the evaluator, you MUST add it here.}
Otherwise the wrapper will accept an incomplete integrand as if it were
complete. The isolation test cited in the comment found
``5/7 multi-diagram groups dropped to 0'' at spike-reset $k=1$, $\ell=2$ ---
exactly the failure mode this guard exists to catch.
\end{gotcha}

\subsection{The entry point and the group key}

The entry point is \code{compute\_correction\_td\_grouped}
(\file{\_grouped\_phase\_j.py:55}):
\begin{lstlisting}
def compute_correction_td_grouped(
    typed_diagrams, prefactors, propagator_data, k,
    num_params=None, ext_time_vars=None,
    origin_leaf_idx=0, external_fields=None,
):
\end{lstlisting}
It takes a flat list of \code{typed\_diagrams} at one loop order, a parallel list
of \code{prefactors} (the \code{combined\_prefactor} per diagram), the Phase~I
\code{propagator\_data}, the leg count \code{k}, optional numeric substitutions
\code{num\_params}, the external time symbols \code{ext\_time\_vars} (auto-built as
$t_1,\dots,t_k$ if \code{None}, \file{\_grouped\_phase\_j.py:72--76}), the pinned
\code{origin\_leaf\_idx}, and the canonical \code{external\_fields} list. It
returns the same dict shape as \code{compute\_correction\_td}.

The grouping itself is a \code{defaultdict} keyed on a pair
(\file{\_grouped\_phase\_j.py:96--100}):
\begin{lstlisting}
groups = defaultdict(lambda: {'tds': [], 'cps': []})
for td, pf in zip(typed_diagrams, prefactors):
    key = (id(td.prediagram), _ext_legs_key(td))
    groups[key]['tds'].append(td)
    groups[key]['cps'].append(pf)
\end{lstlisting}

\begin{defn}
\code{collections.defaultdict} is a \code{dict} subclass that auto-creates a
default value for any missing key, by calling a zero-argument factory you supply.
Here the factory \code{lambda: \{'tds': [], 'cps': []\}} makes a fresh
two-list bucket the first time a key is seen, so
\code{groups[key]['tds'].append(td)} works without a membership check. It is the
idiomatic Python way to bucket items into groups.
\end{defn}

The key has two components. The first, \code{id(td.prediagram)}, is the prediagram
identity --- Python's built-in \code{id()} returns a unique integer per object, and
the enumerator emits the \emph{same} prediagram object for all typed diagrams of one
topology, so identity is a valid grouping key. The second, the nested
\code{\_ext\_legs\_key(td)} (\file{\_grouped\_phase\_j.py:87}), returns a hashable
sorted tuple of the external-legs mapping:
\begin{lstlisting}
def _ext_legs_key(td):
    return tuple(sorted((lf, fld) for lf, fld in td.external_legs.items()))
\end{lstlisting}
Subdividing by external-legs is \emph{required}, not cosmetic. Two typed diagrams
that differ only in which leaf carries which external field have different Wick
mappings and different compensation factors inside the evaluator, so they cannot
share a sum (\file{\_grouped\_phase\_j.py:78--86}). The per-diagram path handles
each independently; the grouped path subdivides at this finer key to stay safe.

\subsection{The three-tier dispatch}

For each group the wrapper tries the grouped evaluator and falls back in tiers
(\file{\_grouped\_phase\_j.py:125--228}):

\begin{description}
  \item[Tier 1 --- accept the grouped result.] Call
        \code{integrate\_grouped\_diagram}. If it returns \code{status='ok'}
        \emph{and} no subset diagnostic is in \code{\_SUBSET\_SKIP\_STATUSES}
        (the \code{has\_skipped\_subset} scan at
        \file{\_grouped\_phase\_j.py:147--153}), accept it: append
        \code{result['contribution']} to the running list of group callables and
        record metadata with \code{'handled\_by': 'grouped\_evaluator'}.
  \item[Tier 2 --- per-diagram fallback on a silent subset skip.] If
        \code{status='ok'} but a subset \emph{was} silently dropped, re-compute
        the \emph{entire group} via the per-diagram path
        (\code{\_perdiag\_fallback}, \file{\_grouped\_phase\_j.py:107}), append
        that callable, and record \code{'handled\_by': 'perdiag\_fallback'} with
        the skip reasons. This trades the grouped speedup on that group for the
        correct numbers --- exactly what \code{use\_grouped\_phase\_j=False} would
        have produced.
  \item[Tier 3 --- per-diagram fallback on a group-level failure.] If the grouped
        evaluator failed outright (e.g.\ early validation), also fall back to
        per-diagram. If \emph{that} throws too, record a \emph{loud}
        \code{'skipped'} entry (\file{\_grouped\_phase\_j.py:213--228}) rather than
        silently dropping the whole group.
\end{description}

The nested \code{\_perdiag\_fallback(tds, cps)} (\file{\_grouped\_phase\_j.py:107})
simply calls the per-diagram \code{compute\_correction\_td} on the group's diagrams
and returns \code{(pd\_result['total\_C'], len(tds))}, isolating the fallback so its
returned callable has the same \code{(*ext\_time\_values) -> complex} signature as
the grouped \code{contribution}.

\subsection{Assembling the result}

Finally the wrapper builds the two callables the downstream code expects
(\file{\_grouped\_phase\_j.py:231--241}):
\begin{lstlisting}
def total_C(*ext_time_values):
    total = 0.0 + 0.0j
    for fn in group_callables:
        total = total + complex(fn(*ext_time_values))
    return total

def total_C_batch(tau_points, parallel=False, n_workers=None):
    return [complex(total_C(*pt)) for pt in tau_points]
\end{lstlisting}
\code{total\_C} sums every group callable; \code{total\_C\_batch} serially
evaluates \code{total\_C} at each $\tau$-point. The returned dict carries
\code{total\_C}, \code{total\_C\_batch}, the grouped bookkeeping (\code{groups},
\code{skipped\_kernel\_ids}, \code{fallback\_to\_perdiag}), and \code{ext\_time\_vars};
it sets \code{eval\_per\_diagram\_batch} to \code{None} and
\code{delta\_contributions} to \code{[]} --- both unsupported in the prototype
(\file{\_grouped\_phase\_j.py:243--252}).

\begin{gotcha}
\code{total\_C\_batch} ignores its \code{parallel} and \code{n\_workers}
arguments --- the prototype is serial. It keeps the arguments only to match the
per-diagram API signature. This is deliberate and safe: as the codebase memory
records, fork-based parallelism crashes a Jupyter kernel on the user's macOS
machine, so serial-only is the correct default here.
\end{gotcha}

\section{The grouped evaluator, walked}

The evaluator \code{integrate\_grouped\_diagram} (\file{grouped\_integral.py:283})
is where \eqref{eq:linearity} and \eqref{eq:merged} are realised. We follow its
control flow.

\subsection{Validation and shared scaffolding}

It first validates the group (\file{grouped\_integral.py:328--358}): an empty list
fails; a length mismatch between diagrams and prefactors fails; and --- the
defining invariant --- \emph{all} typed diagrams must share the same prediagram
object by identity (\file{grouped\_integral.py:349--350}):
\begin{lstlisting}
pd0 = typed_diagrams[0].prediagram
if not all(td.prediagram is pd0 for td in typed_diagrams):
    return {'status': 'failed', ...}
\end{lstlisting}
Then it builds the shared scaffolding once (\file{grouped\_integral.py:360--379}):
the loop number $L = |E| - |V| + 1$ from \code{\_loop\_number\_from\_graph}
(\file{final\_integral.py:2401}); the Sage graph \code{D = td0.prediagram[0]} and
its leaves \code{td0.prediagram[2]}; and the time-domain propagator matrix
\code{G\_t\_obj = build\_G\_t\_matrix(propagator\_data, t\_sym, num\_params)} built
\emph{once} for the whole group.

\subsection{Wick enumeration and compensation}

Next it enumerates the external-Wick mappings and computes the compensation
divisor (\file{grouped\_integral.py:381--452}). If \code{external\_fields} is
provided and matches the leaf count, it buckets canonical external positions and
leaves by field and takes \code{itertools.permutations} within each field,
building the Cartesian product of per-field permutations into \code{\_all\_mappings}.
If any field's counts mismatch it falls back to the identity mapping with
\code{\_compensation = 1}. Otherwise it computes the \emph{exact} compensation via
\code{external\_wick\_compensation} on every group member and asserts they all
agree (\file{grouped\_integral.py:434--449}); a mismatch returns
\code{status='failed'}, because a shared divisor would mis-weight the group. We
unpack the compensation in \S\ref{sec:compensation}.

\subsection{Vertex times and per-diagram edge info}

For the first mapping it assigns each leaf its external time, pinning the origin
leaf to \code{SR(0)}, and allocates a fresh symbol \code{s\_v...} per internal
vertex into \code{integration\_vars} (\file{grouped\_integral.py:454--469}). A
per-diagram scan handles \code{NoiseSourceType} vertices that introduce extra
$\tau$ integration variables (\file{grouped\_integral.py:471--520}); for plain
groups this loop is a no-op and the analytic path stays enabled.

Then comes \emph{the only genuinely per-diagram work}
(\file{grouped\_integral.py:522--591}). For each typed diagram and each edge
$(u,v,\text{lbl})$ of the shared prediagram, it looks up $(r_i, p_i)$ via
\code{\_lookup\_prop\_indices} (\file{final\_integral.py:5230}), resolves the head
and tail times, and builds the per-edge data:
\begin{lstlisting}
dt = SR(t_v - t_u)
delta_c = G_t_delta_coeff(G_t_obj, pi, ri)
smooth_factor = G_t_entry(G_t_obj, pi, ri, dt, include_heaviside=False)
edge_info.append({'u': u, 'v': v, 'lbl': lbl, 'ri': ri, 'pi': pi,
                  'dt_sym': dt, 'delta_coeff': delta_c,
                  'smooth_factor': smooth_factor})
\end{lstlisting}
Note \code{include\_heaviside=False}: causality is enforced by the polytope
constraints, not by an explicit $\Theta$. (\code{G\_t\_entry} reads the smooth
$[p_i, r_i]$ entry --- the transpose of the kernel's $(r_i, p_i)$ convention, which
both Phase~I and Phase~J must agree on.) It then coerces the diagram's
\code{combined\_prefactor} to \code{SR} and substitutes \code{num\_params}.

\subsection{The analytic mode-sum precomputation}

The analytic path is enabled only if \code{USE\_GROUPED\_ANALYTIC\_MODESUM} is set,
\code{pole\_vals} and \code{C\_mats} are present, \emph{and} no diagram has a
non-empty \code{vertex\_leg\_time} (\file{grouped\_integral.py:602--607}):
\begin{lstlisting}
grouped_analytic_enabled = (
    USE_GROUPED_ANALYTIC_MODESUM
    and pole_vals is not None
    and C_mats is not None
    and not any(vlt for vlt in per_td_vertex_leg_time)
)
\end{lstlisting}
When enabled it precomputes, once, the per-pole rates
$\lambda_k = \texttt{complex(CDF(SR(p)))}\cdot 1j$
(\file{grouped\_integral.py:612}) and the per-diagram/per-edge/per-pole residue
tensor \code{residues\_per\_td\_edge} (\file{grouped\_integral.py:619--628}). These
are reused across all $2^{|E|}$ subsets.

\subsection{The subset loop and merged-residue construction}

The core is the loop over \code{subset\_bits} in \code{range(2**n\_edges)}
(\file{grouped\_integral.py:654--1033}). For each subset it: splits edges into
$\delta$-edges (bit set) and smooth-edges (bit clear); drops any diagram with a
near-zero $\delta$-coefficient on a $\delta$-edge; checks that all contributing
diagrams agree on the symbolic $\Delta t$ of each $\delta$-edge (a mismatch
records \code{dt\_sym\_mismatch\_skipped} and skips --- one of the two
\code{\_SUBSET\_SKIP\_STATUSES}); performs $\delta$-elimination with
\code{sage\_solve}; and skips residual shot-noise $\delta$-spikes with status
\code{shotnoise\_skipped\_in\_prototype} (the other guarded status).

It then builds the \emph{summed} subset integrand --- the right-hand side of
\eqref{eq:linearity} --- as a single \code{SR} expression
(\file{grouped\_integral.py:759--770}):
\begin{lstlisting}
subset_factor_group = SR(0)
for td_idx in contributing_td:
    ei = per_td_edge_info[td_idx]
    cp = per_td_cp[td_idx]
    term = cp
    for i in delta_edges:
        term = term * SR(ei[i]['delta_coeff'])
    for i in smooth_edges:
        term = term * ei[i]['smooth_factor']
    term = term.subs(substitutions)
    subset_factor_group = subset_factor_group + term
\end{lstlisting}
After extracting the linear constraint data $(a^{\text{int}}, a^{\text{ext}}, c_0)$
for each smooth edge (\file{grouped\_integral.py:830--847}) and compiling a single
\code{fast\_callable} over the summed integrand (used only on the scipy fallback),
it constructs the merged-residue pole tuples --- the realisation of
\eqref{eq:Bdef} (\file{grouped\_integral.py:925--948}):
\begin{lstlisting}
grouped_pole_tuples = []
idx_alpha = [0] * n_smooth
while True:
    B_alpha = 0.0 + 0.0j
    for j in range(len(merged_pf_per_j)):
        prod = merged_pf_per_j[j]
        for ee in range(n_smooth):
            prod *= smooth_residues_per_j[j][ee][idx_alpha[ee]]
        B_alpha += prod
    lambdas = tuple(lambdas_per_pole[idx_alpha[ee]] for ee in range(n_smooth))
    grouped_pole_tuples.append((B_alpha, lambdas))
    # advance the multi-index alpha ...
\end{lstlisting}
The outer \code{while} enumerates the multi-index $\bm\alpha$ over per-edge poles;
the inner sum over $j$ accumulates $B_{\bm\alpha} = \sum_j (\mathrm{cp}_j
\prod_\delta \delta) \prod_{\text{smooth}} C^{(j)}_{\alpha_e}$, with
\code{merged\_pf\_per\_j} the per-diagram prefactor-times-$\delta$-coefficients
built at \file{grouped\_integral.py:904--913}. Each $(B_{\bm\alpha},
\text{lambdas})$ pair is exactly what the closed-form evaluators consume.

\subsection{The dummy modes}\label{sec:dummymodes}

The per-diagram $m=2$ and $m\ge 3$ evaluators expect a list of \code{EdgeModeSum}
objects of length $n_{\text{smooth}}$ --- a length check. On the analytic path
those modes are never read (the real weights are in \code{pole\_tuples}), so the
grouped code builds throwaway \code{EdgeModeSum} instances purely to satisfy the
check (\file{grouped\_integral.py:949--961}). Their \code{modes} come from the
first diagram's residues.

\begin{gotcha}
Do not mistake \code{grouped\_dummy\_modes} for the actual weights. They carry the
first diagram's residues so the \code{n\_smooth} length check passes, but when
\code{pole\_tuples} is supplied the evaluator reads its coefficients from there,
not from the modes. The real merged weights live in \code{grouped\_pole\_tuples}.
\end{gotcha}

\subsection{The Wick-permutation sum and the divisor}

The final \code{contribution(*ext\_time\_values)} closure
(\file{grouped\_integral.py:1049--1074}) loops over the Wick permutations: for
each it permutes the external times, time-shifts so the permuted origin leaf
returns to $0$ (the integrand is a function of time \emph{differences}, so this
symmetrises asymmetric integrands and is a no-op for symmetric ones), extracts the
free-time values, and sums every subset closure. Finally it divides by the
compensation:
\begin{lstlisting}
        for cfn in subset_contributions:
            total = total + complex(cfn(free_vals))
    return total / _comp
\end{lstlisting}
That \code{/ \_comp} at \file{grouped\_integral.py:1074} is the external-Wick
compensation, the subject of the next section.

\section{External-Wick compensation}\label{sec:compensation}

The integrators enumerate \emph{every} field-respecting assignment of canonical
external positions to diagram leaves and sum the integrand over all of them. Some
of those assignments are redundant --- they map the diagram to itself via a graph
automorphism --- so the raw sum over-counts each \emph{distinct} pinned-external
diagram. The exact divisor that corrects this is the \term{orbit--stabilizer
index}
\begin{equation}
  \mathrm{comp}
  \;=\;
  \frac{|\Aut(\Gamma,\ \text{leaves free})|}{|\Aut(\Gamma,\ \text{leaves fixed})|},
  \label{eq:comp}
\end{equation}
computed by \code{external\_wick\_compensation} (\file{symmetry.py:343}). By
orbit--stabilizer, two assignments yield the same pinned-external diagram iff they
differ by a leaf-permuting automorphism, and the stabilizer of any assignment is
exactly $\Aut(\Gamma, \text{leaves fixed})$; so the mapping sum counts each
distinct class \eqref{eq:comp} times, and dividing by it recovers the sum over
distinct diagrams the Feynman rules require. The grouped path applies this divisor
once, at \file{grouped\_integral.py:1074}.

\begin{note}
\code{external\_wick\_compensation} computes \eqref{eq:comp} via two calls to
\code{\_automorphism\_order}, which builds a coloured incidence digraph and asks
SageMath for the order of its automorphism group. If $|\Aut_{\text{free}}|$ is not
divisible by $|\Aut_{\text{fixed}}|$ --- impossible in exact group theory, since the
fixed group is a subgroup of the free one, but reachable as a Sage edge case --- it
raises a \code{RuntimeError} rather than return a wrong divisor
(\file{symmetry.py:368--381}). Failing loudly is the whole point of a function
that exists to prevent silent mis-weighting.
\end{note}

\subsection{Why the old heuristic broke at \texorpdfstring{$k\ge 3$}{k>=3}}

The history here is instructive, because it explains a function name that still
appears in the codebase. The previous implementation \emph{approximated} the index
\eqref{eq:comp} as a product of factorials, $\prod N_{\text{sig,field}}!$, over
leaves grouped by the \code{vertex\_role\_signature} (\file{symmetry.py:247}) of
their attachment vertex. The \code{vertex\_role\_signature} is a \emph{depth-1}
graph invariant --- a vertex's type, coefficient, bigrade, external-leaf attachment
pattern, and internal-edge incidence pattern.

That heuristic \emph{over-divides} whenever two attachment vertices share a role
signature (a shallow, depth-1 invariant) but are \emph{not} related by any
automorphism. The documented example is the $k=4$ OU$+\varepsilon x^3$ one-loop
cascades, where three noise vertices all read ``noise feeding a cubic'' at depth~1
yet hang off structurally distinct cubics. The resulting deficits broke
\emph{every} $k\ge 3$ one-loop cumulant while leaving $k=2$ exact --- at $k=2$ the
heuristic happens to coincide with the true index. It was replaced by the exact
Sage automorphism ratio (\file{symmetry.py:368--381}), validated against the exact
Boltzmann series $\kappa_4 = -6\varepsilon + 126\varepsilon^2$.

\begin{gotcha}
\code{vertex\_role\_signature} still exists and is still referenced elsewhere
(\file{final\_integral.py:2573}), but it is \emph{not} the current compensation
mechanism and is not on the grouped path's numerical route. Do not reach for it
when reasoning about the divisor; the grouped path uses the exact
\code{external\_wick\_compensation}. The signature survives only as the historical
basis of the superseded heuristic.
\end{gotcha}

\section{The tooling, for the physicist}

This subsystem leans almost entirely on \term{SageMath}, with a fallback into
\code{scipy}. There is no direct use of nauty, sympy, numba, or networkx in the
two primary files --- though Sage embeds several of those internally. Here is each
piece, from the ground up, with the actual import and call sites.

\subsection{SageMath: \texttt{SR}, \texttt{fast\_callable}, \texttt{CDF}, \texttt{solve}}

\term{SageMath} is a large open-source mathematics system built on Python. It
bundles dozens of specialised libraries (a symbolic-algebra engine, a numerical
stack, graph-automorphism engines, number-theory backends) behind one Python API.
When the code writes \code{from sage.all import ...} it is importing Sage's
``everything'' namespace. The grouped file imports four names
(\file{grouped\_integral.py:101}):
\begin{lstlisting}
from sage.all import SR, fast_callable, CDF, solve as sage_solve
\end{lstlisting}

\begin{description}
  \item[\code{SR} --- the Symbolic Ring.] Sage's ring of symbolic expressions
        (its analogue of SymPy's \code{Symbol}/\code{Expr}, but a different
        engine). \code{SR.var('x')} makes a symbolic variable; \code{SR(3)} wraps
        a number; arithmetic builds expression trees you can \code{.subs()},
        \code{.expand()}, \code{.coefficient()}, \code{.variables()}, and so on.
        The grouped path uses it everywhere: to allocate the internal time symbol
        \code{t\_sym = SR.var('\_t\_td\_')} (\file{grouped\_integral.py:378}), the
        per-vertex time symbols (\file{grouped\_integral.py:467}), the symbolic
        per-edge $\Delta t$ via \code{dt = SR(t\_v - t\_u)}
        (\file{grouped\_integral.py:549}), and the summed subset integrand
        accumulated with \code{+} and \code{*}
        (\file{grouped\_integral.py:760--770}).
  \item[\code{fast\_callable} --- a JIT-style expression compiler.] A symbolic
        \code{SR} expression evaluated naively re-walks its expression tree on
        every call, which is slow. \code{fast\_callable(expr, vars=[...],
        domain=CDF)} compiles it into a fast closure taking the listed variables
        as positional arguments. It is Sage's analogue of SymPy's
        \code{lambdify}. The grouped path builds it once per subset over the
        \emph{summed} integrand (\file{grouped\_integral.py:815--818}) --- one
        compile instead of $N$ --- but only the scipy fallback ever calls the
        result; the analytic path never touches it.
  \item[\code{CDF} --- the Complex Double Field.] Sage's complex numbers backed
        by hardware \code{double} precision. \code{complex(CDF(SR(expr)))} is the
        idiom for ``evaluate this symbolic expression to a fast machine-precision
        Python \code{complex}.'' It is the \code{fast\_callable} domain and the
        cast target for poles and residues, e.g.\ \code{complex(\_CDF(SR(p)))*1j}
        for a rate (\file{grouped\_integral.py:612}).
  \item[\code{solve} (as \code{sage\_solve}) --- symbolic equation solver.] Sage's
        \code{solve(eqn, var, solution\_dict=True)} solves an equation
        symbolically and returns a list of solution dicts. The grouped path uses
        it to \emph{eliminate} an integration variable on each $\delta$-edge: the
        $\delta$-edge imposes the linear equation $\Delta t_e = 0$, and solving it
        for one surviving variable removes that variable from the integral
        (\file{grouped\_integral.py:719--722}).
\end{description}

\subsection{nauty / bliss, reached through Sage}

The compensation factor needs the \emph{order} of the automorphism group of a
coloured graph. Sage computes this with its built-in graph-isomorphism engine
(historically nauty/bliss-class algorithms). The grouped path reaches it
indirectly, through \code{external\_wick\_compensation} $\to$
\code{\_automorphism\_order} $\to$ Sage's \code{automorphism\_group().order()} on
a coloured digraph. The novice takeaway: \emph{nauty/bliss answer ``how many ways
can I relabel this graph's vertices and leave it unchanged, respecting the
colours?'', and that count is the divisor \eqref{eq:comp} that prevents the
external-leaf permutation sum from over-counting.}

\subsection{scipy, the numerical fallback}

\code{scipy} is the standard Python scientific-computing library;
\code{scipy.integrate.nquad} performs \term{adaptive multidimensional numerical
quadrature} --- it samples the integrand on an adaptively-refined grid until a
relative tolerance (default $\sim\!10^{-8}$) is met. The grouped path does not
import scipy directly, but its fallback routes through \code{\_integrate\_polytope}
(\file{final\_integral.py:4372}), which is scipy-backed. The grouped path only
reaches scipy when an analytic evaluator returns \code{None} (overflow, an
unbounded-divergent endpoint, a non-rational kernel). This is the slower,
lower-precision route; Stage~4a exists precisely to avoid it.

\begin{note}
What is \emph{not} used here: \code{numba} (the analytic mode-sum is pure
Python-level complex arithmetic, fast enough without it; numba lives in the
spatial path); \code{networkx} (all graph work goes through Sage \code{Graph}
objects carried on the prediagram); and \code{sympy} (Daedalus's CAS is Sage's
\code{SR}, not SymPy). The standard-library helpers are \code{defaultdict} (for
the group buckets), \code{itertools.permutations} (for the Wick mappings), and
\code{cmath}/\code{math} (the complex exponential and exact $\pm\infty$ tracking
in the hot helpers).
\end{note}

\section{Precision: analytic versus scipy, grouped versus per-diagram}\label{sec:precision}

There are two precision stories here, and they should not be confused.

The first is \term{analytic versus scipy}, and the grouped path wins it cleanly.
On the analytic merged-residue path every subset integral is a closed-form
mode-sum, evaluated in machine-precision complex arithmetic, so the precision
floor is machine $\varepsilon$ ($\sim\!10^{-15}$). The scipy fallback is adaptive
quadrature at $\sim\!10^{-8}$ relative. Stage~4a (the master switch
\code{USE\_GROUPED\_ANALYTIC\_MODESUM = True} at \file{grouped\_integral.py:130})
is what moved the grouped path onto the analytic route by default.

The second is \term{grouped versus per-diagram}, and here lies the open issue. By
\eqref{eq:linearity} the two paths are equal in exact arithmetic, and the codebase
treats agreement to floating-point round-off as a \emph{correctness contract}: the
regression suite \file{tests/test\_grouped\_vs\_perdiag.py} computes the sum of $N$
per-diagram results and the single grouped result and asserts they match to
$\sim\!10^{-14}$ relative on the analytic path. The working fixture for this is
\code{quad\_exp\_k2\_ell0}.

\begin{gotcha}
\textbf{Open precision bug at $m\ge 3$ with close-paired poles.} The grouped path
sums $B_{\bm\alpha}$ \emph{before} the chain-simplex integration; the per-diagram
path sums \emph{after}. When close-paired poles produce $O(1/\Delta p)$ residues
and a tuple multiplies several of them into $O(10^8)$, the two summation orders
disagree in float64 even though they agree exactly. The spike-reset $k=2$,
$\ell=1$ case shows a $\sim\!4\times$ discrepancy (per-diagram $+2.66\times10^{-3}$
versus grouped $+6.38\times10^{-4}$). The \code{USE\_POSET\_CAP\_MATCH\_SCIPY}
analysis (\file{final\_integral.py:1690--1700}) concludes the bounding-box cap is
\emph{not} the dominant cause and points to a bookkeeping difference in pole-tuple
construction between the two paths. An attempted fix using extended-precision
accumulation (\code{USE\_POSET\_MPMATH\_ACCUMULATION}, off by default) ``did NOT
address the spike-reset bug.'' The authoritative writeup is
\file{docs/m\_ge3\_precision\_bug\_audit.md}.
\end{gotcha}

\begin{note}
A separate, easily-confused failure is \emph{not} an integrator bug at all. Four
\code{spike\_reset\_*} cases in \file{tests/test\_grouped\_vs\_perdiag.py} fail
because a fixture still passes a \code{taug} parameter that the theory file
dropped when it switched to a delta synaptic kernel; the failure happens in the
pole/residue computation \emph{before} any Phase~J integration runs. That is
fixture/theory drift, not a Phase~J defect. Use \code{quad\_exp\_k2\_ell0} as the
trustworthy regression fixture.
\end{note}

\section{Limitations and the fallback matrix}

The grouped path is a prototype, and it is honest about what it does not handle.
In each of the following cases the grouped evaluator either disables its analytic
route (falling to scipy within the grouped path) or signals the wrapper to fall
back to the per-diagram path entirely. Knowing the matrix tells you when you must
use \code{use\_grouped\_phase\_j=False}.

\begin{description}
  \item[Noise-source kernels disable the analytic path.] If \emph{any} diagram in
        the group has a non-empty \code{vertex\_leg\_time} --- a
        \code{NoiseSourceType} cumulant kernel --- the integrand contains
        non-rational factors and the merged-residue closed form cannot fire. The
        subset falls through to \code{fast\_callable} $+$ scipy
        (\file{grouped\_integral.py:602--607}). The numbers are still correct,
        just at scipy precision.
  \item[Shot-noise $\delta$-spikes are skipped.] A residual external-time equality
        after $\delta$-elimination ($\tau=0$ same-population legs in heterogeneous
        Hawkes) is recorded as \code{shotnoise\_skipped\_in\_prototype} and the
        subset is dropped (\file{grouped\_integral.py:736--757}). The per-diagram
        path accounts for these in \code{delta\_contributions}, which the grouped
        path does not rebuild --- so the wrapper's Tier~2 falls the whole group back
        to per-diagram.
  \item[\code{delta\_contributions} and \code{eval\_per\_diagram\_batch} are
        unsupported.] The grouped return sets them to \code{[]} and \code{None}
        (\file{\_grouped\_phase\_j.py:246--248}). Any caller that needs a
        per-diagram breakdown or the $\delta$-spike accounting must use the
        per-diagram path.
  \item[The group must be homogeneous in compensation.] If
        \code{external\_wick\_compensation} differs across group members the
        evaluator returns \code{status='failed'}
        (\file{grouped\_integral.py:437--448}); the wrapper's Tier~3 then falls
        back to per-diagram. (The wrapper already subdivides by external-legs to
        make this rare.)
  \item[The batch is serial.] As noted, \code{total\_C\_batch} ignores
        \code{parallel}/\code{n\_workers} --- the safe default on the user's
        machine.
\end{description}

\begin{note}
The convention details that must match exactly between the grouped and
per-diagram paths, or boundary contributions would shift: $\Theta(0) = 0$ (a
constraint $\Delta t > 0$ is \emph{strict}, the boundary is outside --- enforced by
the \code{c\_eff <= 0} rejection in both
\code{\_evaluate\_grouped\_m0\_modesum} at \file{grouped\_integral.py:155} and
\code{\_integrate\_grouped\_m1\_modesum} at \file{grouped\_integral.py:217}); and
the overflow guard at $|\gamma_{\text{real}}| > 600$
(\code{\_GROUPED\_EXP\_REAL\_LIMIT}, \file{grouped\_integral.py:135}), a deliberate
margin below the $709$ hard ceiling of \code{cmath.exp} on IEEE doubles, which
returns \code{None} (to fall back) rather than raising.
\end{note}

\section{Where this sits in the pipeline}

To close, the data flow in one picture. The dispatch site is
\code{compute\_cumulants} in \file{pipeline/compute.py}. For each loop order
$\ell$ it gathers the typed diagrams and their prefactors, and if
\code{use\_grouped\_phase\_j} is set it calls the grouped wrapper instead of the
per-diagram path (\file{pipeline/compute.py:757--774}):
\begin{lstlisting}
if use_grouped_phase_j:
    from pipeline._grouped_phase_j import compute_correction_td_grouped
    td_result_ell = compute_correction_td_grouped(
        typed_diagrams = [r['typed_diagram'] for r in records_ell],
        prefactors     = [r['combined_prefactor'] for r in records_ell],
        k              = k,
        propagator_data= propagator_data,
        external_fields= external_fields,
        num_params     = num_params,
        origin_leaf_idx= origin_leaf_idx,
    )
else:
    td_result_ell = compute_correction_td(...)   # the per-diagram path
\end{lstlisting}
Because both paths return the same dict shape, the lines that follow ---
\code{total\_C\_by\_ell[ell] = td\_result\_ell['total\_C']} and the $\tau$-grid
evaluation \code{td\_result\_ell['total\_C\_batch'](tau\_points, ...)}
(\file{pipeline/compute.py:787--794}) --- never learn which path produced the
numbers. The master \code{total\_C} sums across loop orders and $C(\tau) =
\sum_\ell C_\ell(\tau)$ exactly as before. The grouped bookkeeping keys are
available for diagnostics but are not required downstream.

That is the whole subsystem: recognise that a prediagram's typed diagrams share a
domain, sum their integrands once (\eqref{eq:linearity}), pre-sum the residues
into $B_{\bm\alpha}$ (\eqref{eq:Bdef}), integrate the shared exponential-of-linear
in closed form by $m$, divide by the orbit--stabilizer compensation, and fall back
to the per-diagram path the moment any of that cannot be trusted. The payoff is
one JIT compile instead of hundreds and a precision floor at machine $\varepsilon$,
for numbers that match the per-diagram path bit-for-bit wherever the contract
holds.


% =====================================================================
%  PART V  --  EVALUATING DIAGRAMS : SPATIAL
% =====================================================================
\part{Evaluating Diagrams: Spatial Models}
% ---------------------------------------------------------------------
%  Chapter 14 : The Spatial Integrator --- Symanzik Polynomials and
%               Causal Chambers
% ---------------------------------------------------------------------
\chapter{The Spatial Integrator: Symanzik Polynomials and Causal Chambers}\label{ch:spatial-core}

The previous part of this manual evaluated diagrams \emph{in the time domain}:
each diagram became a product of exponential propagators, and the loop integrals
were done over the internal vertex times. That route works beautifully for
temporal (time-only) models, but it carries a structural wound --- the
\term{close-pair pathology}, a $1/\beta$ factor that blows up whenever two poles
of the time integrand approach each other. For a \emph{spatial} model --- a
stochastic \emph{partial} differential equation, where the field lives on a
$d$-dimensional space as well as in time --- there is a much cleaner road, and
this chapter is about the subsystem that paves it.

The idea, in one sentence, is: \emph{do the momentum integral first, and do it
analytically}. Every enumerated diagram --- a tree, a one-loop bubble, a one-loop
tadpole, a two-loop sunset, for any number of external legs $k$, any loop order
$\ell$, and any spatial dimension $d$ --- is reduced to the \emph{same} integral
and evaluated by the \emph{same} code. There is no per-topology branch. The
phrase the codebase repeats, in the very first lines of the production module, is
that ``ONE genuine integral evaluates \emph{every} enumerated diagram''
(\file{msrjd/integration/spatial/full\_integrator.py:4}). This chapter explains
why that is possible and walks the machinery that does it.

\begin{note}
This is the first of three spatial chapters. This one builds the core ---
momentum routing, the Symanzik polynomials, the causal chambers, and the kinematic
integral. Chapter~\ref{ch:spatial-heatkernel} covers the analytic
inverse Fourier transform and the derivative (form-factor) vertices in detail.
Chapter~\ref{ch:spatial-coupled} covers the coupled multi-field (unequal-$D$)
extension. Everything here lives under
\file{msrjd/integration/spatial/}, and none of it imports SageMath or
\code{nauty} --- those are upstream, in the enumeration. The only heavy
dependencies here are NumPy, SciPy, and (in one small place) SymPy.
\end{note}

\section{Why momentum-first}

\subsection{The wound in the temporal route}

Recall how the temporal pipeline evaluates a loop. After enumerating and typing a
diagram it has a product of propagators in the $(\,$frequency, or time$\,)$
representation, and it integrates the internal vertex times by carving the time
domain into ordering regions and integrating a product of exponentials over each.
The integrand of one such region looks schematically like a sum of terms
$\ee^{-\beta_1 t_1}\ee^{-\beta_2 t_2}\cdots$, and doing the nested integrals
produces denominators built from \emph{differences} of the decay rates $\beta_i$.
When two of those rates coincide --- a degenerate or near-degenerate pole pair ---
the denominator $1/(\beta_i-\beta_j)$ diverges. The true integral is finite (the
apparent pole is removable, a confluent limit), but the naive expression is not,
and resolving it costs precision and special-case code. That is the close-pair
pathology, and Part~IV spends real effort taming it.

\subsection{The cure}

For a spatial model the field also carries a momentum. The crucial observation is
that in the heat-kernel representation the edge \emph{times are already Schwinger
parameters} (we make this precise in \S\ref{sec:heatkernel} below), so the
loop-\emph{momentum} integral $\int \dd^d\ell$ is a \emph{pure Gaussian}. A
Gaussian integral can be done in closed form, exactly, at any loop order and in
any dimension --- it collapses to the two \term{Symanzik polynomials}. And here is
the payoff: the close-pair pathology lived entirely in the \emph{momentum}
integral. Once that integral is done analytically, what remains is a smooth
integral over the internal vertex times and a few auxiliary parameters, with no
poles to collide. Ordinary quadrature converges fast, and close-pair
\emph{cannot arise}. The comment block in \file{causal\_chambers.py:17-24} states
this explicitly: the temporal pipeline ``integrates products of exponentials per
chamber $\dots$ whose $1/\beta$ factors are the close-pair pathology. Here the
loop momenta are already integrated $\dots$ so the per-chamber integrand is
\emph{smooth} $\dots$ and close-pair cannot arise.''

\begin{defn}
The \term{one mental model} for this whole subsystem: a heat-kernel propagator in
the momentum--time representation is $G(k,t)=\theta(t)\,\ee^{-(\mu+Dk^2)t}$. The
edge times act as Schwinger parameters, so $\int\dd^d\ell$ is a pure Gaussian and
collapses \emph{exactly} to the Symanzik polynomials $\Usym$ and $\Fsym$ at any
loop order $L$ and any dimension $d$. After that collapse the only thing left is a
smooth integral over the internal vertex times and the noise-source Schwinger
parameters --- no pole, no close pair. That is the entire design.
\end{defn}

\subsection{Where this sits in the pipeline}

The subsystem consumes a fully enumerated and typed diagram and produces that
diagram's contribution to the connected $k$-point cumulant. Upstream, the
enumeration has (a)~expanded the action into vertices, (b)~enumerated all
topologically distinct Feynman diagrams up to the chosen loop order, and
(c)~``typed'' each diagram --- decided which lines are response/retarded ($R$)
lines and which are noise/correlation ($C$) lines, and attached a symbolic
prefactor carrying the couplings, noise amplitudes, and symmetry factor
$\Sym(\Gamma)$. What is left is a hard multidimensional integral \emph{per
diagram}, and this subsystem performs it. The flow is:

\begin{lstlisting}[style=console]
   theory (.theory.py)  ->  action -> vertices
        |
   diagram enumeration  (msrjd/diagrams/...)   -> prediagrams
        |  field-type assignment
        v
   TypedDiagram          (msrjd/diagrams/type_assignment.py)
        |  .prediagram, .vertex_assignments, .propagator_indices, prefactor
        v
 +-----------------------  THIS SUBSYSTEM  --------------------------+
 |  diagram_to_cstack(td)              -> CStackDiagram (descriptor) |
 |  route_momenta(td) -> edge_coeffs() -> per-edge (a_e, b_e)        |
 |  _symanzik_kernel_batch / _momentum_factor_batch -> U, F          |
 |  causal_chambers(...)               -> orderings of {t_v}         |
 |  diagram_kinematic(...)             -> int dt int dsigma          |
 |  _heat_kernel_x_general / IFT       -> dC(x,tau)  directly        |
 +------------------------------------------------------------------+
        |  sum over diagrams, x 2^{-n_C} . S(Gamma) . prefactor
        v
   C(q,tau) or dC(x,tau)  ->  pipeline_bridge.py  ->  compute_cumulants(...)
\end{lstlisting}

The six files of the subsystem, and their roles:

\begin{center}
\begin{tabular}{@{}>{\ttfamily\footnotesize}l l@{}}
\toprule
\normalfont\textbf{file (under \file{msrjd/integration/spatial/})} & \textbf{role} \\
\midrule
diagram\_descriptor.py & typed diagram $\to$ flat ``C-stack'' edge list \\
momentum\_routing.py   & momentum conservation $\to$ routing coefficients $(a_e,b_e)$ \\
spatial\_reduce.py     & the textbook Symanzik reduction (single-sample reference) \\
causal\_chambers.py    & the causal time-ordering chambers \\
full\_integrator.py    & the production integrator (batched + analytic IFT) \\
generic\_evaluator.py  & \normalfont\emph{oracle only} --- an older cross-check \\
\bottomrule
\end{tabular}
\end{center}

\section{The heat-kernel representation}\label{sec:heatkernel}

This section builds the relevant theory from the ground up. A reader who knows
\msrjd{} field theory but not this particular representation should be able to
follow every step.

\subsection{The two propagators}

Start from a linearized (Gaussian, ``tree'') theory with one field. In the
momentum--time representation the \term{retarded propagator} --- the response
function, the $R$ line --- is
\begin{equation}
  \Gret(k,t) \;=\; \theta(t)\,\ee^{-m(k)\,t},
  \qquad
  m(k) \;=\; \mu + D\,k^2,
  \label{eq:GR}
\end{equation}
where $\mu$ is a mass / relaxation rate, $D$ a diffusion constant,
$k^2=|k|^2$ the squared momentum, and $\theta$ the Heaviside step. The step
function is causality made manifest: a response only to the past. The
\term{correlation function} --- the $C$ line, $\avg{\phi\phi}$ --- of the
\emph{tree} theory is
\begin{equation}
  C_0(k,t) \;=\; \frac{T}{m(k)}\;\ee^{-m(k)\,|t|},
  \label{eq:C0}
\end{equation}
with $T$ the noise strength (the action carries a noise vertex $-T\,\tilde\phi^2$,
or $2T\,\tilde\phi\,\tilde\phi$ in the symmetric normalization). Real-space
correlators are inverse Fourier transforms of these over $k$.

\subsection{A correlation line is two retarded segments glued at a noise source}

The identity the whole subsystem rests on is that a correlation line can be
written as two retarded segments joined at a noise source, with the source time
integrated out. Schematically,
\begin{equation}
  C_0(\Delta t) \;=\; \frac{T}{m}\,\ee^{-m|\Delta t|}
  \;=\; T\!\int_0^{\infty}\!\dd\sigma\;\ee^{-m(|\Delta t|+\sigma)}\cdot(\text{const}),
  \label{eq:Cglue}
\end{equation}
i.e.\ the $C$ line is an $R$--noise--$R$ object. Integrating the noise-source
time produces a single \term{Schwinger parameter} $\sigma\in[0,\infty)$ and the
weight $\ee^{-m(|\Delta t|+\sigma)}$; the convention factor of $2$ is bookkept by
$2^{-n_C}$ (\S\ref{sec:norm}). This is exactly why the enumerator can work
entirely with all-$\Gret$ propagators plus explicit noise sources, and why the
descriptor re-contracts a 2-point noise source back into one $C$ edge.

\begin{defn}
A \term{Schwinger parameter} is an auxiliary integration variable that linearizes
a propagator denominator into an exponential. Here there are two kinds: an
$R$-edge time $w$, and a noise-source parameter $\sigma\in[0,\infty)$. Their job
is to turn the momentum dependence of every propagator into a single Gaussian
exponent $\exp(-D\sum_e w_e k_e^2)$, so that $\int\dd^d\ell$ is exactly doable.
\end{defn}

\subsection{Edge weights and the loop-momentum integral}

For a diagram with internal vertices at times $\{t_v\}$ (to be integrated),
external leaves at fixed times $\{\tau_j\}$, and edges $e$, each edge carries a
\term{Schwinger weight} $w_e$ (\file{full\_integrator.py:16-19}):
\begin{align}
  R \text{ edge } (t_{\text{head}},t_{\text{tail}}):&\quad
     w_e = t_{\text{head}} - t_{\text{tail}} \;\ge\; 0 \ \text{(by causality)},\\
  C \text{ edge between } t_a,t_b:&\quad
     w_e = |t_a - t_b| + \sigma_e,\quad \sigma_e\in[0,\infty).
\end{align}
Each edge also carries a routed momentum $k_e$ (next section). The diagram's
loop-momentum integral, in the heat-kernel representation, is the pure Gaussian
\begin{equation}
  I_{\text{mom}} \;=\; \int \prod_{i=1}^{L}\frac{\dd^d\ell_i}{(2\pi)^d}\;
     \exp\!\Big[-D\sum_e w_e\,k_e^2\Big],
  \label{eq:Imom}
\end{equation}
with $L$ the number of loops. This is the central object
(\file{spatial\_reduce.py:12}). The next two sections compute it.

\section{Momentum routing}\label{sec:routing}

\subsection{Conservation at a spatial vertex}

A spatial interaction vertex is \emph{local in space}: the action contains
$\int\dd x_v$ of the product of the fields meeting at the vertex. Fourier
transforming, that integral enforces \term{momentum conservation}
$\sum_e k_e = 0$ at every internal vertex --- the spatial analogue of the
time-domain vertex, which carries one integrated time $t_v$. Solving this linear
conservation system writes each edge momentum as a linear form
\begin{equation}
  k_e \;=\; \sum_{i=1}^{L} a_{ei}\,\ell_i \;+\; \sum_{j=1}^{n_{\text{ext}}} b_{ej}\,q_j,
  \label{eq:route}
\end{equation}
in the \term{loop momenta} $\ell_i$ (the $L$ free parameters of the conservation
system) and the \term{external momenta} $q_j$ (one per external leg). Overall
conservation $\sum_j q_j = 0$ is imposed by eliminating the last leaf, so
$n_{\text{ext}} = k-1$. The coefficients $(a_e,b_e)$ are the per-edge
\term{routing coefficients}; for a translation-invariant routing they are
integers, $\pm1$ or $0$.

Two cases make this concrete (\file{momentum\_routing.py:29-34}):

\begin{itemize}
  \item \textbf{Tree ($L=0$).} Every edge carries $\pm q$ only, so $k_e^2 = q^2$
        uniformly. A single-mode tree therefore needs only a
        $\Lap \to -q^2$ substitution.
  \item \textbf{One-loop bubble.} The two internal lines carry $\ell$ and
        $q-\ell$ respectively --- genuinely different momenta --- so the
        substitution becomes per-edge.
\end{itemize}

\subsection{The symbolic solve}

The routing is done with SymPy, the pure-Python computer-algebra system, because
momentum conservation is a small \emph{exact} linear-algebra problem and a CAS
gives integer $\pm1/0$ coefficients with no floating-point ambiguity.

\begin{note}
\textbf{What SymPy is.} SymPy (\code{import sympy as sp}) represents mathematical
\emph{symbols} (\code{sp.Symbol}), does exact rational/integer arithmetic, and
--- crucially here --- solves linear systems symbolically with
\code{sp.linsolve}. It is the only place in this entire subsystem that SymPy
appears; the heavy numerics are all NumPy/SciPy.
\end{note}

The function is \code{route\_momenta} (\file{momentum\_routing.py:99}). Reading
its body, the steps are:

\begin{enumerate}
  \item Create the external-momentum symbols $q_0\dots q_{k-1}$ and one unknown
        momentum symbol per edge:
        \begin{lstlisting}
q = sp.symbols(f'q0:{k}') if k > 0 else ()
kk = {e: sp.Symbol(f'k{i}') for i, e in enumerate(edges)}   # one per edge
        \end{lstlisting}
  \item Build the conservation \term{residual} at each vertex --- the signed sum
        that must vanish. The convention (\file{momentum\_routing.py:103-107}) is
        that edge $(u,v,\text{lbl})$ is oriented $u\to v$, so it contributes
        $+k_e$ at the vertex it flows \emph{into} and $-k_e$ at the vertex it
        flows \emph{out of}; at a leaf the injected external momentum $q_j$ is
        subtracted:
        \begin{lstlisting}
for v in vertices:
    s = sp.Integer(0)
    for e in edges:
        u, w, _lbl = e
        if w == v: s += kk[e]          # edge flows INTO v
        if u == v: s -= kk[e]          # edge flows OUT of v
    if v in leaves:
        s -= q[leaves.index(v)]        # external momentum injected at leaf
    residuals.append(sp.expand(s))
        \end{lstlisting}
  \item Impose overall conservation by substituting
        $q_{k-1} = -(q_0+\dots+q_{k-2})$ into every residual.
  \item Solve the system for the edge momenta with SymPy's linear solver:
        \begin{lstlisting}
sol_set = sp.linsolve(residuals, list(kk.values()))
        \end{lstlisting}
        An empty solution set means the system is inconsistent --- which never
        happens for a valid \msrjd{} diagram, so the code raises a clear error if
        it does (\file{momentum\_routing.py:139-142}).
  \item The \emph{free symbols} SymPy leaves unpinned in the solution \emph{are}
        the loop momenta. The code collects them
        (\code{expr.free\_symbols}), subtracts the external symbols, sorts by
        name, and relabels them to $\ell_0\dots\ell_{L-1}$
        (\file{momentum\_routing.py:149-158}).
\end{enumerate}

\subsection{The routing result}

The result is a small dataclass (\file{momentum\_routing.py:50}):

\begin{lstlisting}
@dataclass
class RoutingResult:
    edge_momenta: dict   # (u,v,lbl) -> sympy expr in q_syms + loop_syms
    q_syms: tuple        # external-momentum symbols q_0 ... q_{k-1}
    loop_syms: tuple     # loop-momentum symbols l_0 ... l_{L-1}
    n_loops: int
\end{lstlisting}

It exposes two extractors. \code{edge\_coeffs()}
(\file{momentum\_routing.py:62}) returns \code{\{edge: (a, b)\}} --- the signed
linear coefficients of \eqref{eq:route}, each coerced to a Python \code{int} by a
small helper \code{\_num} (\file{momentum\_routing.py:79}). \textbf{This is what
the Symanzik step consumes.} The other extractor, \code{edge\_k2()}
(\file{momentum\_routing.py:57}), returns the symbolic $k_e^2$ --- but it must
\emph{not} be used for the reduction.

\begin{gotcha}
\code{edge\_k2()} squares the momentum and loses the \emph{signed} cross-term
coefficients (the $a_{ei}b_{ej}$ products that build the loop/external coupling
block). Only \code{edge\_coeffs()} carries the $(a,b)$ the Symanzik step needs.
The code warns about this in the docstring at
\file{momentum\_routing.py:72-74}. If you ever extend the integrator, consume
\code{edge\_coeffs()}, never \code{edge\_k2()}.
\end{gotcha}

\section{Symanzik polynomials from scratch}\label{sec:symanzik}

Now do the Gaussian integral \eqref{eq:Imom}. Substitute the routing
\eqref{eq:route} and collect the loop momenta into a quadratic form. Define three
matrices by summing the routing coefficients weighted by the Schwinger weights
(\file{spatial\_reduce.py:18-21}):
\begin{align}
  \Lambda_{ii'} &= \sum_e w_e\,a_{ei}\,a_{ei'} &&(L\times L) &&\text{first-Symanzik kernel,}\\
  N_{ij} &= \sum_e w_e\,a_{ei}\,b_{ej} &&(L\times n_{\text{ext}}) &&\text{loop/external cross block,}\\
  Q_{jj'} &= \sum_e w_e\,b_{ej}\,b_{ej'} &&(n_{\text{ext}}\times n_{\text{ext}}) &&\text{pure external block.}
\end{align}
Then $\sum_e w_e k_e^2 = \ell^{\mathsf T}\Lambda\ell + 2q^{\mathsf T}N^{\mathsf T}\ell
+ q^{\mathsf T}Q\,q$. Completing the square in $\ell$ and doing the Gaussian
integral gives the closed form (\file{spatial\_reduce.py:24-26}):
\begin{equation}
  I_{\text{mom}} \;=\; (4\pi D)^{-Ld/2}\;\Usym^{-d/2}\;
     \exp\!\big[-D\,q^{\mathsf T} Q_{\text{eff}}\, q\big],
  \label{eq:collapse}
\end{equation}
with
\begin{equation}
  \Usym = \det\Lambda,
  \qquad
  Q_{\text{eff}} = Q - N^{\mathsf T}\Lambda^{-1}N,
  \qquad
  \Fsym \equiv \Usym\cdot Q_{\text{eff}}\ \text{(scalar case)}.
  \label{eq:UF}
\end{equation}

\subsection{What these polynomials are}

It is worth saying what $\Usym$ and $Q_{\text{eff}}$ \emph{are} classically,
because a reader may know the graph-theoretic definitions.

\begin{description}
  \item[$\Usym$ --- the first Symanzik (first graph) polynomial.] A homogeneous
        degree-$L$ polynomial in the edge weights. Graph-theoretically,
        $\Usym = \sum_{\text{spanning trees }T}\ \prod_{e\notin T} w_e$ --- a sum
        over spanning trees of the product of the weights of the edges \emph{not}
        in the tree. In this code it is computed \emph{numerically} as
        $\det\Lambda$, which equals that combinatorial sum (the matrix-tree
        theorem). Validated examples (\file{spatial\_reduce.py:36-38}): the
        one-loop bubble has $\Usym = w_1+w_2$; the two-loop sunset has
        $\Usym = w_1w_2 + w_2w_3 + w_3w_1$.
  \item[$Q_{\text{eff}}$ --- the second Symanzik / $\Fsym$.] Carries the
        external-momentum dependence. $Q_{\text{eff}} = Q - N^{\mathsf T}\Lambda^{-1}N$
        is the \term{Schur complement} of $\Lambda$ in the full quadratic form ---
        exactly the residual external quadratic form after integrating out the
        loops. For the bubble, $Q_{\text{eff}} = w_1w_2/(w_1+w_2)$; for the
        sunset, $Q_{\text{eff}} = w_1w_2w_3/\Usym$. For $n_{\text{ext}}=1$ (the
        2-point function, $k=2$) $Q_{\text{eff}}$ is a scalar and
        $\Fsym = \Usym\cdot Q_{\text{eff}}$.
\end{description}

\begin{defn}
The \term{Schur complement} of a block $\Lambda$ in a partitioned quadratic form
$\left(\begin{smallmatrix}\Lambda & N\\ N^{\mathsf T} & Q\end{smallmatrix}\right)$ is
$Q - N^{\mathsf T}\Lambda^{-1}N$. It is precisely what is left of the $q$-quadratic
form after the $\ell$ block has been integrated out (Gaussian-eliminated). That is
why $Q_{\text{eff}}$ is the residual external form: integrating the loops
\emph{is} taking a Schur complement.
\end{defn}

The exponents $-Ld/2$ and $-d/2$ are the \emph{only} places the spatial dimension
$d$ enters the momentum reduction. The reduction is otherwise exact and
$d$-agnostic --- the same algebra works in $d=1$ as in $d=7$.

\subsection{The reference implementation}

\file{spatial\_reduce.py} is the textbook, single-sample reference. Its three
functions are deliberately small and quotable. \code{symanzik\_matrices}
(\file{spatial\_reduce.py:58}) builds $(\Lambda,N,Q)$ for one set of edge weights
by exactly the contractions above:
\begin{lstlisting}
aw  = a * w[:, None]          # (E, L)
Lam = aw.T @ a                # (L, L)
N   = aw.T @ b                # (L, n_ext)
Q   = (b * w[:, None]).T @ b  # (n_ext, n_ext)
\end{lstlisting}
\code{symanzik\_polynomials} (\file{spatial\_reduce.py:82}) returns
$(\Usym, Q_{\text{eff}})$ with $\Usym = \det\Lambda$ and $Q_{\text{eff}}$
computed via \code{np.linalg.solve} (not an explicit inverse --- solving a linear
system is both faster and more accurate than forming $\Lambda^{-1}$). And
\code{momentum\_integral} (\file{spatial\_reduce.py:101}) evaluates
\eqref{eq:collapse} for one sample. The production code re-implements all three,
\emph{batched}, in \file{full\_integrator.py}; the reference exists so the batched
version can be checked against it.

\begin{note}
\textbf{Why \code{np.linalg.solve} and not \code{np.linalg.inv}.} To compute
$N^{\mathsf T}\Lambda^{-1}N$ you never need $\Lambda^{-1}$ itself --- you need the
product $\Lambda^{-1}N$, which is the solution $X$ of the linear system
$\Lambda X = N$. \code{np.linalg.solve(Lam, N)} computes exactly that, using an
LU factorization rather than a full matrix inversion. It is the standard numerical
hygiene: never invert when you can solve.
\end{note}

\section{The full diagram integral}

Putting the pieces together, one diagram's contribution to the connected
$k$-point cumulant is (\file{full\_integrator.py:15-20}):
\begin{equation}
  \Gamma(q,\{\tau\}) \;=\; 2^{-n_C}\,\Sym(\Gamma)\!
     \int \prod_v \dd t_v \prod_{C\text{ edges}}\!\dd\sigma_e\;
     \mathbf{1}(\theta\text{'s})\;
     \ee^{-\mu\sum_e w_e}\;\text{MomFactor}(w,q),
  \label{eq:diagram}
\end{equation}
\begin{equation}
  \text{MomFactor} \;=\; (4\pi D)^{-Ld/2}\,\Usym(w)^{-d/2}\,
     \ee^{-D\,q^{\mathsf T}Q_{\text{eff}}(w)\,q}.
\end{equation}
The factor $\ee^{-\mu\sum_e w_e}$ is the \term{mass damping}: every segment of
length $w$ decays as $\ee^{-\mu w}$, since the $\mu$ part of the mass $m(k)=\mu+Dk^2$
multiplies $w_e$ directly. The $Dk^2$ part of the mass already lives \emph{inside}
$\text{MomFactor}$, through the Symanzik exponent. The $\mathbf{1}(\theta\text{'s})$
are the causal step functions of the retarded edges. The residual integral is over
the internal vertex times $t_v$ and the correlation Schwinger parameters
$\sigma_e$ only --- and it is \emph{smooth}, because the loop integral, the source
of the close-pair pathology, was already done analytically.

\section{Normalization --- derived, not fitted}\label{sec:norm}

The integrator's normalization is \emph{derived}, not tuned to a fixture
(\file{full\_integrator.py:27-31}). The enumeration prefactor uses the $2T$
noise-vertex convention. A kinematic $C$ edge here, by contrast, is the
unit-amplitude Schwinger factor $\int_0^\infty\dd\sigma\,\ee^{-m\sigma}=1/m$. The
factor $2^{-n_C}$ ($n_C$ = number of $C$ edges) converts between the two. The
tree check makes this airtight: a tree 2-point function has no loops and one $C$
edge between the two leaves ($n_C=1$), prefactor $2T$, and the integral gives
\begin{equation}
  2^{-1}\cdot 2T\cdot\frac{1}{m}\,\ee^{-m\tau}
  \;=\; \frac{T}{m}\,\ee^{-m\tau}
  \;=\; C_0,
\end{equation}
exactly the tree correlator \eqref{eq:C0}. This $2^{-n_C}$ is applied in
\code{diagram\_value} / \code{diagram\_value\_x} (\file{full\_integrator.py:1153},
\file{full\_integrator.py:1210}).

\section{Causal chambers}\label{sec:chambers}

\subsection{Posets and linear extensions}

What remains after the momentum collapse is
$\int\prod_v\dd t_v\,\mathbf{1}(\theta\text{'s})\,(\text{smooth integrand})$. The
retarded edges impose orderings $t_v > t_u$ on the internal vertices --- a
\term{partial order} (a poset). The domain of a poset decomposes into
\term{ordering chambers}, one per \term{linear extension} (a total order
consistent with the partial order). The key fact: within each chamber every
$|\Delta t|$ has a \emph{fixed sign}, so the integrand is smooth (no cusps) and
ordinary quadrature converges fast (\file{causal\_chambers.py:18-24}).

\begin{defn}
A \term{linear extension} of a poset is a total order of its elements that
respects all the partial-order relations. Geometrically, each linear extension is
one simplicial chamber of the cube $[\,\text{lo},\text{hi}\,]^{n}$, and the
chambers \emph{partition} the constrained domain. Integrating over the whole
constrained domain therefore equals summing the integral over the chambers --- no
over- or under-counting.
\end{defn}

\subsection{The code, reused verbatim from the temporal pipeline}

The poset machinery is not reimplemented for the spatial path. The single public
function (\file{causal\_chambers.py:41}) wraps the temporal pipeline's
battle-tested \code{\_CausalPoset} and \code{\_enumerate\_linear\_extensions}:

\begin{lstlisting}
def causal_chambers(n_vertices, retarded_edges):
    poset = _CausalPoset(
        m=int(n_vertices),
        edges=tuple(sorted({(int(u), int(v)) for (u, v) in retarded_edges})),
        scalar_lowers=(), scalar_uppers=())
    return list(_enumerate_linear_extensions(poset))
\end{lstlisting}

Here \code{retarded\_edges} is an iterable of pairs $(u,v)$ meaning $t_v > t_u$.
The function returns the list of linear extensions, each a length-$n$ tuple giving
the vertex order (earliest first).

\begin{gotcha}
\textbf{The empty-edge case tiles the cube.} With \emph{no} retarded ordering
constraints, \code{causal\_chambers} returns all $n!$ orderings, which
\emph{partition} the time cube --- the integral over the cube equals the sum over
chambers. Forgetting this gives an $n!$ over- or under-count. This is why the
production integrator always loops over the returned chambers even when there are
no ordering constraints.
\end{gotcha}

The module also ships a reference integrator, \code{integrate\_over\_chambers}
(\file{causal\_chambers.py:84}), which does the chamber simplex integral by nested
\code{scipy.integrate.quad} (adaptive Gauss--Kronrod 1-D quadrature). That path is
an oracle; the \emph{production} integrator uses fixed Gauss--Legendre product
grids instead, for vectorization, which we turn to next.

\section{The descriptor: \texorpdfstring{\texttt{diagram\_to\_cstack}}{diagram\_to\_cstack}}\label{sec:descriptor}

Before the integrator can run, the typed diagram must be flattened into the
canonical \term{C-stack} form that every diagram shares. That is the job of
\code{diagram\_to\_cstack} (\file{diagram\_descriptor.py:105}). The two record
types it produces are frozen dataclasses.

\begin{lstlisting}
@dataclass(frozen=True)
class CEdge:
    a: tuple        # loop-momentum routing coefficients (over the L loops)
    b: tuple        # external-q routing coefficients (over the n_ext externals)
    kind: str       # 'R' (retarded segment) or 'C' (correlation line)
    u: int          # endpoint vertex id
    v: int          # endpoint vertex id  (u == v  =>  self-loop = tadpole)
    external: bool  # True iff the edge touches an external leaf
    fpairs: tuple = ()   # coupled-field propagator matrix indices (Dyson 3c)
\end{lstlisting}

\begin{lstlisting}
@dataclass(frozen=True)
class CStackDiagram:
    internal_vertices: tuple   # interaction-vertex ids (their times integrated)
    external_legs: tuple       # leaf vertex ids, in correlation-function order
    edges: tuple               # tuple of CEdge
    n_loops: int               # == diagram loop_number
\end{lstlisting}

The mapping itself, reading \file{diagram\_descriptor.py:105-193}:

\begin{enumerate}
  \item Unpack \code{td.prediagram} into the graph \code{D}, the leaves, and the
        internal vertices; read \code{td.vertex\_assignments}.
  \item \textbf{Classify each non-leaf vertex.} A \code{VertexType} is an
        interaction vertex (its time is integrated). A 2-point \code{SourceType}
        --- the Gaussian noise insertion $\avg{\tilde\phi\tilde\phi}$ that
        represents a $C$ line --- is a \emph{noise} vertex, to be contracted away.
        A \code{SourceType} with $\ge 3$ response legs (non-Gaussian noise) is
        \emph{kept} as an internal vertex with $n$ outgoing $R$ edges --- the
        paper's all-retarded formulation --- and its time is integrated like any
        vertex. A source with fewer than 2 response legs raises
        \code{NotImplementedError}.
  \item Call \code{route\_momenta(td)} and read
        \code{ec = rr.edge\_coeffs()} into a per-edge $(a,b)$ map.
  \item Build an incidence map from each vertex to its incident edges.
  \item \textbf{Contract each 2-point noise source} into a single $C$ edge: its
        two incident $\Gret$ half-edges become one $C$ edge between the two
        \emph{other} endpoints (the source time integrated out gives
        $C(\Delta t)$). A noise source whose two edges land on the \emph{same}
        vertex becomes a $C$ \emph{self-loop} ($u=v$) --- the tadpole. Degree
        $\ne 2$ raises.
  \item \textbf{Every remaining $\Gret$ edge becomes an $R$ edge.} Edges touching
        a leaf are flagged \code{external=True}.
  \item Assemble and return the \code{CStackDiagram}, internal vertices sorted,
        external legs in correlation-function order, with
        \code{n\_loops = rr.n\_loops}.
\end{enumerate}

\begin{note}
\textbf{Why the $C$ edge may take either half's routing coefficients.} When the
noise source is contracted, the new $C$ edge takes the $(a,b)$ of one of its two
halves (\file{diagram\_descriptor.py:168}). This is correct because flipping
$(a,b)\to(-a,-b)$ leaves every Symanzik sum $\sum_e w_e a a$, $\sum_e w_e a b$,
$\sum_e w_e b b$ invariant --- the routing sign is irrelevant to the quadratic
forms. The descriptor \emph{relies} on this; do not try to ``fix'' the sign.
\end{note}

\begin{gotcha}
\textbf{There is no bubble / tadpole / sunset branch in the descriptor.} Whether a
loop momentum couples to the external $q$ (a bubble) or not (a tadpole) is a
property of the Symanzik $\Fsym$ \emph{downstream}, not of this descriptor
(\file{diagram\_descriptor.py:10-13}). \code{CStackDiagram.is\_tadpole\_like()}
exists (\file{diagram\_descriptor.py:93}) but is \emph{diagnostic only} --- the
production evaluator never branches on it. This topology-blindness is the whole
point: one integral, every diagram.
\end{gotcha}

\section{The kinematic integral: \texorpdfstring{\texttt{diagram\_kinematic}}{diagram\_kinematic}}\label{sec:kinematic}

\code{diagram\_kinematic} (\file{full\_integrator.py:469}) is the heart of the
subsystem. It computes the kinematic (unit-amplitude, no couplings) full-diagram
integral
$\int\prod\dd t_v\prod\dd\sigma_e\,\mathbf{1}(\theta)\,\ee^{-\mu\sum w}\,\text{MomFactor}$
by causal-chamber quadrature. Its signature is:

\begin{lstlisting}
def diagram_kinematic(descr, q_vec, external_times, mu, D, spatial_dim=1,
                      W=None, n_t=22, n_s=24, formfactor=None, gh_order=6,
                      xs=None, method='grid', mc_n=1000000, mc_seed=0):
\end{lstlisting}

Walking the body:

\begin{enumerate}
  \item \textbf{Extract the structure} (\file{full\_integrator.py:487-493}). Pull
        \code{edges}, \code{internal} vertices, the $a$/$b$ coefficient matrices
        (one row per edge), and $n_C$ (the number of $C$ edges).
  \item \textbf{Set the time window} (\file{full\_integrator.py:495-499}).
        $[\,\text{lo},\text{hi}\,]$ comes from the external times and a window
        $W=22/\mu$ --- about twenty-two relaxation times, well past where the
        integrand has decayed.
  \item \textbf{Build the retarded poset on the internal vertices}
        (\file{full\_integrator.py:503-515}). An internal$\to$internal $R$ edge
        becomes a poset edge \code{internal\_R}; an $R$ edge to or from a leaf
        becomes a fixed-time scalar bound \code{s\_up}/\code{s\_lo} (the leaf time
        clamps the vertex from one side).
  \item \textbf{Build the correlation Schwinger grid}
        (\file{full\_integrator.py:517-536}). This step is subtle enough to
        deserve its own paragraph below.
  \item \textbf{Dispatch on \code{method}} (\file{full\_integrator.py:544-555}).
        \code{'bessel'} routes to \code{\_diagram\_bessel\_xs}; otherwise
        enumerate \code{chambers = causal\_chambers(n\_V, internal\_R)}.
  \item \textbf{Per chamber, build the internal-time grid.} \code{'mc'}
        importance-samples the chamber via nested $\mathrm{Exp}(\mu)$ time gaps
        (\file{full\_integrator.py:564-578}); \code{'grid'} (the default) builds
        the nested causal-chamber product grid latest$\to$earliest, each level
        bounded above by the next-later time and its external scalar bound, via
        \code{\_gl\_on} (\file{full\_integrator.py:586-603}).
  \item \textbf{Compute the edge weights} \code{w\_batch} and the residual mass
        \code{mu\_resid} ($R$: $t_v-t_u$; $C$: $|t_u-t_v|+\sigma$) and the
        per-sample amplitude \code{amp} (\file{full\_integrator.py:606-623}).
  \item \textbf{Accumulate the result}, in one of four sub-cases
        (\file{full\_integrator.py:625-709}): a plain analytic-IFT heat kernel; a
        multivariate ($k\ge3$) derivative-vertex IFT; a single-external
        ($k=2$) derivative-vertex IFT; or the $q$-evaluation path. The last three
        are the subject of Chapter~\ref{ch:spatial-heatkernel}.
  \item \textbf{Return} a real vector for the analytic-IFT path, complex for a
        derivative vertex, float otherwise
        (\file{full\_integrator.py:714-716}).
\end{enumerate}

\subsection{The retarded-time quadrature: \texorpdfstring{\texttt{\_gl\_on}}{\_gl\_on}}

The internal vertex times are integrated by \code{\_gl\_on}
(\file{full\_integrator.py:349}), Gauss--Legendre nodes on $[\text{lo},\text{hi}]$
with a $\sqrt{\,}$-concentrating substitution toward the upper bound, where the
retarded integrand peaks:

\begin{lstlisting}
def _gl_on(lower, upper, n):
    xg, wg = np.polynomial.legendre.leggauss(n)
    v = 0.5 * (xg + 1.0)
    span = (upper - lower)[..., None]
    t = upper[..., None] - span * (v * v)   # nodes cluster near `upper`
    w = span * (wg * v)                     # Jacobian 2.span.v . 1/2
    return t, w
\end{lstlisting}

The $t = \text{upper} - \text{span}\cdot v^2$ map clusters nodes near the upper
bound (because $v^2$ is dense near $v=0$) and carries the Jacobian
$\text{span}\cdot(w_g\,v)$. This concentration matters: the retarded propagator
$\ee^{-m\,w}$ is largest when the segment length $w$ is smallest, i.e.\ when the
vertex time sits just below its upper bound.

\subsection{The Schwinger \texorpdfstring{$\sigma$}{sigma} grid is load-bearing}

The correlation Schwinger parameters $\sigma_e\in[0,\infty)$ are integrated by the
\emph{same} kind of substitution. Reading
\file{full\_integrator.py:517-525}:

\begin{lstlisting}
s_cap = 32.0 / mu
xv, wv = np.polynomial.legendre.leggauss(n_s)
vv = 0.5 * (xv + 1.0)
s_nodes = s_cap * vv * vv                         # sigma = s_cap . v^2
s_w = wv * s_cap * vv * np.exp(-mu * s_nodes)     # w.s_cap.v.e^{-mu.sigma}
\end{lstlisting}

That is Gauss--Legendre on a finite window $[0,s_{\text{cap}}]$ with
$\sigma = s_{\text{cap}}\,v^2$ and the $\ee^{-\mu\sigma}$ weight folded into the
node weights, \emph{not} Gauss--Laguerre on $[0,\infty)$.

\begin{gotcha}
\textbf{The $\sigma$-node concentration is not optional.} The substitution
$\sigma = s_{\text{cap}}\,v^2$ concentrates nodes near $\sigma = 0$
specifically to resolve the integrable self-loop singularity
$\Usym^{-d/2}\sim\sigma^{-d/2}$ (a $C$ self-loop has $\Usym = \sigma$, so as
$\sigma\to0$ the integrand has an integrable spike). Plain Gauss--Laguerre
\emph{under-resolves} that spike. The code comment at
\file{full\_integrator.py:519} says exactly this. Changing this quadrature will
quietly degrade tadpole accuracy without raising any error. (Note that the
function's own docstring at \file{full\_integrator.py:482} still calls this
``Gauss--Laguerre''; the body is the Legendre-plus-substitution above, and the
body is what runs.)
\end{gotcha}

\subsection{A worked shape: the one-loop bubble}

Make the abstract concrete with a one-loop bubble, $k=2$, $d=1$. It has two
internal vertices ($n_V=2$), two correlation lines ($n_C=2$), and one loop
($L=1$). \code{causal\_chambers(2, internal\_R)} returns the orderings of the two
vertices. Per chamber the grid has $P \approx 22^2\cdot 24^2$ samples; $a$ is
$(E,1)$ and $b$ is $(E,1)$; $\Lambda$ is $(P,1,1)$ so $\Usym = w_1+w_2$ and
$Q_{\text{eff}} = w_1w_2/(w_1+w_2)$; $\mathcal{B} = D\cdot Q_{\text{eff}}$ is a
scalar; and the heat kernel $(4\pi\mathcal{B})^{-1/2}\ee^{-x^2/4\mathcal{B}}$ is
summed over the grid into $\delta C(x)$ (the analytic IFT of
Chapter~\ref{ch:spatial-heatkernel}).

\section{The batched Symanzik core}\label{sec:batched}

The production integrator does \emph{not} call the single-sample
\file{spatial\_reduce.py} functions in its inner loop. Instead it builds the
Symanzik matrices for an entire grid of Schwinger samples at once, using NumPy's
Einstein-summation engine \code{np.einsum}.

\begin{note}
\textbf{What \code{np.einsum} is.} \code{np.einsum} evaluates a tensor
contraction written in index notation. The string \code{'pe,ej,ek->pjk'} reads
``take an array indexed by $(p,e)$ and two indexed by $(e,j)$ and $(e,k)$, sum over
the repeated index $e$, and keep $(p,j,k)$''. It runs in compiled C, and for the
heavy linear algebra (\code{det}, \code{solve}) NumPy calls BLAS/LAPACK, which
releases the Python GIL. Here $p$ indexes the Schwinger sample, $e$ the edge, $l/m$
the loop momenta, and $j/k$ the external momenta.
\end{note}

The batched factory is \code{\_momentum\_factor\_batch}
(\file{full\_integrator.py:56}). Its core (\file{full\_integrator.py:73-90}):

\begin{lstlisting}
Q   = np.einsum('pe,ej,ek->pjk', w_batch, b, b)   # (P, n_ext, n_ext)
if L == 0:
    quad = np.einsum('j,pjk,k->p', qv, Q, qv)
    return np.exp(-D * quad)                       # tree: exp(-D q^T Q q)
Lam = np.einsum('pe,el,em->plm', w_batch, a, a)    # (P, L, L)
N   = np.einsum('pe,el,ej->plj', w_batch, a, b)    # (P, L, n_ext)
U   = np.linalg.det(Lam)                           # (P,)  the first Symanzik
ok  = U > u_floor                                  # degeneracy mask
LamiN = np.linalg.solve(Lam[ok], N[ok])            # Lam^{-1} N, batched
Qeff  = Q[ok] - np.einsum('plj,plk->pjk', N[ok], LamiN)
pref  = (4.0*math.pi*D)**(-0.5*spatial_dim*L) * np.power(U[ok], -0.5*spatial_dim)
out[ok] = pref * np.exp(-D * np.einsum('j,pjk,k->p', qv, Qeff, qv))
\end{lstlisting}

Every line corresponds to one symbol in \eqref{eq:collapse}: \code{Q}, \code{Lam},
and \code{N} are the three Symanzik matrices for the whole batch; \code{U} is
$\det\Lambda$ for every sample at once; \code{LamiN} is $\Lambda^{-1}N$ via the
batched solve; \code{Qeff} is the Schur complement; and \code{pref} is the
$(4\pi D)^{-Ld/2}\Usym^{-d/2}$ prefactor.

\begin{gotcha}
\textbf{Degeneracy floors are everywhere, and they silently zero samples.} The
mask \code{ok = U > u\_floor} (with \code{u\_floor = 1e-300}) selects the samples
where the loop matrix is non-degenerate; degenerate samples (a self-loop where
$\Usym\to0$) get \emph{zero}, with no warning. The same pattern appears for the
heat-kernel width ($\mathcal{B} > 10^{-300}$, $\det\mathcal{B} > 10^{-300}$). A
diagram whose \emph{entire} grid is degenerate returns $0$ silently. This is
correct (those configurations have measure zero), but it means a bug that
degenerates every sample looks like a clean zero, not an error.
\end{gotcha}

The sibling \code{\_symanzik\_kernel\_batch} (\file{full\_integrator.py:94}) is the
analytic-IFT version: it returns the $q$-Gaussian
$\text{pref}\cdot\ee^{-q^{\mathsf T}\mathcal{B}q}$ \emph{without} contracting in $q$,
so the $x$-dependence stays analytic. That function feeds the heat-kernel IFT of
Chapter~\ref{ch:spatial-heatkernel}.

\section{The three backends}\label{sec:backends}

\code{diagram\_kinematic} takes a \code{method} argument selecting one of three
quadrature backends. They trade memory against generality.

\begin{description}
  \item[\code{'grid'} (default).] The deterministic causal-chamber product grid
        described in \S\ref{sec:kinematic}. Exact (to quadrature order), validated,
        but its grid size is $P \approx n_t^{\,n_V}\cdot n_s^{\,n_C}$ \emph{per
        chamber} --- the curse of dimensionality.
  \item[\code{'mc'}.] Monte-Carlo importance sampling
        (\file{full\_integrator.py:564-578}): internal times via nested
        $\mathrm{Exp}(\mu)$ gaps, correlation $\sigma$'s as $\mathrm{Exp}(\mu)$,
        with $P=$\code{mc\_n} random points instead of the product grid. Bounded
        memory, error $O(1/\sqrt{N})$.
  \item[\code{'bessel'}.] The radial-Bessel $\times$ angular-Monte-Carlo backend
        (\file{full\_integrator.py:362}), covered in detail in
        Chapter~\ref{ch:spatial-heatkernel}. It reparametrizes the causal-region
        point as $\lambda\cdot\hat s$ with $\hat s$ on a simplex, does the radial
        $\lambda$-integral \emph{exactly} as a modified Bessel function $K_\nu$,
        and samples only the smooth angular simplex. The radial reduction handles
        the $\det\Lambda\to0$ degenerate-loop direction analytically.
\end{description}

\begin{gotcha}
\textbf{The memory wall at $\ell\ge2$.} The grid is $P=n_t^{\,n_V}\cdot n_s^{\,n_C}$
per chamber. A KPZ two-loop diagram ($n_V=4$, $n_C=3$) is $\approx 1.8\times10^8$
points per chamber --- tens of gigabytes. The bridge refuses up front (raising a
\code{SpatialPropagatorError}) and points you to the escape hatches: lower
\code{max\_ell}, set \code{SPATIAL\_INTEGRATOR=mc} or \code{=bessel}, coarsen
\code{SPATIAL\_GRID\_NT}/\code{NS}, or raise \code{SPATIAL\_MEM\_BUDGET\_GB}.
\end{gotcha}

\begin{gotcha}
\textbf{Monte-Carlo is biased for derivative vertices.} The \code{'mc'} backend is
validated to better than $0.1\%$ for \emph{plain} $\phi^n$ vertices but is
\emph{biased} for derivative (form-factor) vertices: the $\det\Lambda\to0$
singularity gives the Monte-Carlo estimator infinite variance. For derivative
vertices at $\ell\ge2$, use \code{SPATIAL\_INTEGRATOR=bessel} --- the radial
reduction cures the variance. This is stated repeatedly in the source
(\file{full\_integrator.py:560-563}).
\end{gotcha}

\section{Assembling a correlator}

The kinematic integral is unit-amplitude. The amplitude, the retarded$+$advanced
completion, and the diagram sum are layered on by a short chain of functions, all
near the bottom of \file{full\_integrator.py}.

\code{diagram\_value} (\file{full\_integrator.py:1142}) applies the amplitude:
\begin{lstlisting}
def diagram_value(descr, prefactor_val, q_vec, external_times, mu, D,
                  spatial_dim=1, **kw):
    n_C = sum(1 for e in descr.edges if e.kind == 'C')
    kin = diagram_kinematic(descr, q_vec, external_times, mu, D,
                            spatial_dim=spatial_dim, **kw)
    return (2.0 ** (-n_C)) * float(prefactor_val) * kin
\end{lstlisting}
That is exactly $2^{-n_C}\cdot\Sym(\Gamma)\cdot\text{kinematic}$, with the
$2^{-n_C}$ of \S\ref{sec:norm}.

\code{diagram\_correlator} (\file{full\_integrator.py:1166}) adds the
retarded$+$advanced completion. A \emph{retarded-type} insertion --- one whose two
external legs have \emph{different} propagator kinds, one $C$ and one $R$
(\code{\_is\_retarded\_type}, \file{full\_integrator.py:1156}) --- dresses both the
retarded and advanced sides of the line, so it appears as the pair
$\Sigma_R(\tau)+\Sigma_A(\tau) = \Gamma(\tau)+\Gamma(-\tau)$. A $\{R,R\}$ (Keldysh)
insertion is its own conjugate and contributes a single $\Gamma(\tau)$. At
$\tau=0$ the retarded-type case becomes $2\Gamma(0)$.

\code{correlator\_2pt} (\file{full\_integrator.py:1182}) sums over all enumerated
diagrams --- tree plus every loop, each via the \emph{same} full integral ---
skipping diagrams whose prefactor is numerically dead
($|\text{pre}| < 10^{-14}$, ``dead at the saddle''):
\begin{lstlisting}
def correlator_2pt(descrs_prefactors, q, tau, mu, D, spatial_dim=1, **kw):
    total = 0.0
    for descr, pre in descrs_prefactors:
        if abs(float(pre)) < 1e-14:
            continue
        total += diagram_correlator(descr, pre, q, tau, mu, D,
                                    spatial_dim=spatial_dim, **kw)
    return total
\end{lstlisting}

The three functions \code{diagram\_value\_x}, \code{diagram\_correlator\_x}, and
\code{correlator\_2pt\_x} (\file{full\_integrator.py:1203},
\file{full\_integrator.py:1213}, \file{full\_integrator.py:1228}) are the
analytic-IFT analogues: they return real-space $\delta C(x,\tau)$ as a vector over
the output positions \code{xs}, directly, with no $q$-grid. \textbf{These are the
production entry points} the bridge calls for real-space output, and they are the
bridge to Chapter~\ref{ch:spatial-heatkernel}.

\begin{gotcha}
\textbf{Complex per-diagram values are normal.} A derivative vertex with an odd
number of derivatives gives a complex per-diagram contribution, because
$\partial_x\to\ii k$ in momentum space. The \emph{physical} correlator is real:
the imaginary parts cancel in the diagram sum and are dropped at the real-space
output (\file{full\_integrator.py:710-716}). Do not panic at a complex
intermediate value --- it is a bookkeeping artifact of a single diagram, not a
sign of error.
\end{gotcha}

\section{The oracle evaluator}

The sixth file, \file{generic\_evaluator.py}, is \emph{not} on the production
path. Its banner says so in the first lines (\file{generic\_evaluator.py:2-4}):
``ORACLE-ONLY --- not on the production path. Superseded by
\code{full\_integrator.py} $\dots$ reached only by its own test(s). Kept as an
independent numerical cross-check --- \code{compute\_cumulants} does NOT use this
module.'' It evaluates a diagram by an older \emph{Dyson-convolution} route
(self-energy $\sigma(q,u)$ followed by a single-mode Dyson dressing).

\begin{gotcha}
\textbf{The oracle still branches on topology --- which is exactly what the
production design avoids.} \file{generic\_evaluator.py} branches on
\code{is\_tadpole\_like()} to split the bubble and tadpole cases (in its
\code{diagram\_delta\_C} dispatch). That is the \emph{only} place in the spatial
subsystem where a topology branch survives, and it is a relic of the older design.
It is kept solely so the topology-blind production integrator can be cross-checked
against an independent implementation. Never call it from production code.
\end{gotcha}

\section{What feeds and consumes this subsystem}

To close the loop on the data flow: the orchestrator that prepares the inputs and
consumes the outputs is \file{msrjd/integration/spatial/pipeline\_bridge.py}.

\textbf{In.} \code{build\_pipeline\_records} yields $(td,\text{pre})$ pairs per
loop order --- a \code{TypedDiagram} and a symbolic prefactor. The bridge
substitutes numeric couplings to get a float
$\text{pv} = \text{float}(\text{pre}\big|_{\text{saddle}})$, builds the descriptor
\code{dd = diagram\_to\_cstack(td)}, optionally a form factor \code{ff}, and tags
the loop order. Diagrams with $|\text{pv}| < 10^{-14}$ are dropped.

\textbf{Through.} For real-space output the bridge calls, per live diagram and per
$\tau$ (\file{pipeline\_bridge.py:1697}):
\begin{lstlisting}
dCx_by_ell[el][it, :] += diagram_correlator_x(
    dd, pv, xg, float(tau), mu0, D0, spatial_dim=d,
    n_t=nt, n_s=ns, formfactor=ff, ...)
\end{lstlisting}
which flows \code{dd, pv} $\to$ \code{diagram\_value\_x} $\to$
\code{diagram\_kinematic(..., xs=xg)} $\to$ the chamber loop $\to$
\code{\_symanzik\_kernel\_batch} $\to$ \code{\_heat\_kernel\_x\_general} $\to$ the
accumulated $\delta C(x)$.

\textbf{Out.} The bridge accumulates an array of shape $(n_\tau, n_x)$ per loop
order, which becomes \code{compute\_cumulants}'s $\delta C(x,\tau)$.

\begin{note}
\textbf{macOS fork crash --- a project-wide rule that touches this subsystem only
by mandate.} Fork-based multiprocessing crashes the user's Jupyter kernel and
operating system on this machine; the spatial path is therefore serial or
thread-only by mandate. \emph{None} of the six files in this subsystem fork --- but
if you extend the integrator, any parallelism you add must be thread-based (NumPy
releases the GIL for the heavy linear algebra) or $q$-grid batching, never
process- or fork-based.
\end{note}

\section{Recap}

You have now seen the spine of the spatial integrator:
\begin{enumerate}
  \item \textbf{Why momentum-first} --- the close-pair pathology lives in the
        momentum integral, so doing that integral analytically removes it
        (\S\ref{sec:heatkernel}).
  \item \textbf{Momentum routing} --- conservation at each spatial vertex, solved
        symbolically with SymPy, gives the routing coefficients $(a_e,b_e)$
        (\S\ref{sec:routing}).
  \item \textbf{Symanzik polynomials} --- the loop Gaussian collapses to
        $(4\pi D)^{-Ld/2}\Usym^{-d/2}\ee^{-Dq^{\mathsf T}Q_{\text{eff}}q}$, with
        $\Usym=\det\Lambda$ and $Q_{\text{eff}}$ the Schur complement, exact in $L$
        and $d$ (\S\ref{sec:symanzik}).
  \item \textbf{The descriptor} --- \code{diagram\_to\_cstack} flattens every
        typed diagram into one canonical form with no topology branch
        (\S\ref{sec:descriptor}).
  \item \textbf{Causal chambers} --- the retarded poset's linear extensions tile
        the time domain into smooth chambers (\S\ref{sec:chambers}).
  \item \textbf{The kinematic integral} --- \code{diagram\_kinematic} does the
        residual smooth integral by batched Gauss--Legendre product grids over the
        chambers (\S\ref{sec:kinematic}, \S\ref{sec:batched}), with three backends
        trading memory for generality (\S\ref{sec:backends}).
\end{enumerate}

What this chapter \emph{deferred} is the step from the per-sample $q$-Gaussian to
real-space $\delta C(x,\tau)$ --- the analytic inverse Fourier transform --- and
the handling of derivative (form-factor) vertices. Both are
Chapter~\ref{ch:spatial-heatkernel}. The coupled multi-field (unequal-$D$)
extension, which replaces the single mass $\mu$ with a table of per-segment masses,
is Chapter~\ref{ch:spatial-coupled}.

% ---------------------------------------------------------------------
%  Chapter 15 : The Spatial Heat-Kernel Propagator and the Analytic IFT
% ---------------------------------------------------------------------
\chapter{The Spatial Heat-Kernel Propagator and the Analytic IFT}\label{ch:spatial-heatkernel}

Every chapter before this one has worked, ultimately, in a \emph{mixed}
representation. The temporal pipeline lives in frequency $\omega$; the spatial
core that precedes this chapter keeps a formal Laplacian operator symbol
floating through the algebra, untouched, because the Laplacian cannot be folded
into the $\omega$-pole finder the way a mass or a damping rate can. At some
point that mixed representation has to be cashed out into a physical answer ---
a correlation function $C(x,\tau)$ you can plot against a separation $x$ and a
time lag $\tau$. This chapter is the layer that does the cashing out.

The whole subsystem rests on a single, very old idea: \emph{the inverse Fourier
transform of a Gaussian is a Gaussian}. For a diffusive (reaction--diffusion,
Allen--Cahn-like) field the momentum-space propagator is, after one contour
integral, a Gaussian in the wavevector $k$. Its inverse transform back to real
space is therefore another Gaussian --- the \term{heat kernel}, the Green's
function of the diffusion equation $\partial_t \phi = B\,\Lap\phi$. The payoff
of recognising this is enormous: the spatial transform is done \emph{with
pencil and paper}, in closed form, and never as a discrete numerical transform
on a grid of momenta. The slogan inside the codebase is that the analytic heat
kernel ``retired the $q$-grid.''

We will build this from the ground up. First the physics: the inverse
propagator of a diffusive field, the contour integral that turns it into
exponential decay in time, and the Gaussian integral that turns it into the
heat kernel in space (\S\ref{sec:hk-fromprop}--\S\ref{sec:hk-kernel}). Then the
two complications a real model can add --- a drift (advection) term
(\S\ref{sec:hk-drift}) and a periodic box (\S\ref{sec:hk-images}). Then the
code that reads the mass, diffusion, and drift symbolically off the inverse
propagator matrix using SageMath (\S\ref{sec:hk-extract}), and the code that
assembles and reconstructs the runtime propagator
(\S\ref{sec:hk-build}). With the propagator in hand we build the tree-level
correlator --- noise driving two retarded legs, collapsed to a single
error-function integral (\S\ref{sec:hk-tree})--- and read the noise strength
out of the action (\S\ref{sec:hk-noise}). Finally we cover what happens at
loop order: the batched analytic transform in the production integrator and the
gate that selects it (\S\ref{sec:hk-loops}), the numerical-transform
cross-check kept alive behind environment variables
(\S\ref{sec:hk-numerical}), the Symanzik-polynomial oracle modules that pin the
production numbers (\S\ref{sec:hk-symanzik}), and the guard rails that make the
whole thing refuse to lie to you (\S\ref{sec:hk-guards}).

\begin{note}
This chapter documents the \emph{real-space propagator layer} for spatial
\msrjd{} theories. The four files at its centre are
\file{msrjd/integration/spatial/heat\_kernel.py} (the closed-form kernels and
the symbolic extraction), \file{spatial\_correlator.py} (the tree-level
correlator and the output transforms), and two modules ---
\file{temporal\_integrate.py} and \file{loop\_parametric.py} --- that carry a
banner reading ``\code{ORACLE-ONLY --- not on the production path}.'' Those two
are independent numerical cross-checks used in tests to pin the production
loop integrator; they are how loops were \emph{validated}, not how loops are
\emph{computed today}. The production loop integrator is
\file{full\_integrator.py}, whose batched transform
\code{\_heat\_kernel\_x\_general} we cover in \S\ref{sec:hk-loops}. We are
careful throughout to keep that distinction visible.
\end{note}

\section{From the inverse propagator to real space}\label{sec:hk-fromprop}

\subsection{The \msrjd{} inverse propagator of a diffusive field}

In the Martin--Siggia--Rose--Janssen--De~Dominicis (\msrjd) formalism a
stochastic PDE is rewritten as a field theory in two fields: the \emph{physical}
field $\phi$ and a \emph{response} (or ``tilde'') field $\tilde\phi$. Take the
prototypical reaction--diffusion (Allen--Cahn-like) equation
\begin{equation}
  \Dt\phi \;=\; B\,\Lap\phi \;-\; A\,\phi \;-\;(\text{nonlinear terms})\;+\;\xi,
  \qquad
  \avg{\xi(x,t)\,\xi(x',t')} \;=\; 2\kappa\,\delta(x-x')\,\delta(t-t').
  \label{eq:hk-langevin}
\end{equation}
The quadratic (``Gaussian'', ``free'') part of the corresponding action gives an
\term{inverse propagator}: the matrix whose inverse is the two-point function.
After a Fourier transform in time and the replacement of the Laplacian operator
by its momentum symbol $\Lap \to -k^2$, the single-field diagonal entry reads
\begin{equation}
  K(\omega,k) \;=\; A \;+\; B\,k^{2} \;+\; \ii\,\omega.
  \label{eq:hk-Kft}
\end{equation}

Three coefficients, three meanings, and the rest of this chapter is in essence
the story of how each one is extracted and used:

\begin{description}
  \item[$A$ --- the \term{mass}.] The $k^0,\omega^0$ term: the relaxation rate
        of the slowest, longest-wavelength mode. For the free theory $A=\mu$.
        At a $\phi^4$ saddle the curvature of the potential shifts it to
        $A=\mu+3\lambda\,\phi_*^{2}$. The code tolerates a \emph{complex} $A$
        (a complex saddle, a Keldysh-rotated combination).
  \item[$B$ --- the \term{diffusion}.] The coefficient of $k^{2}$, i.e.\ the
        $D$ in $D\,\Lap$. Physically it must be positive: a negative $B$ is
        anti-diffusion, which amplifies short wavelengths without bound and is
        ill-posed. The code enforces this (\S\ref{sec:hk-guards}).
  \item[$\ii\,\omega$ --- the time derivative.] The standard \msrjd{}
        first-order-in-time term, and it must be \emph{unit-normalized}: the
        coefficient of $\omega$ has to be exactly $\ii$. The code checks this
        and rejects anything else as ``not a standard \msrjd{} kernel.''
\end{description}

\begin{note}
The textbook \msrjd{} convention sends $\Dt \to -\ii\omega$; this code uses the
$+\ii\omega$ convention, so the time-derivative term in \eqref{eq:hk-Kft}
appears as $+\ii\omega$ rather than $-\ii\omega$. Nothing physical depends on
the sign choice --- it only fixes which half-plane you close the contour in
below --- but it is why the normalization check looks for $\ii$ and not $-\ii$.
\end{note}

The propagator itself is the reciprocal of \eqref{eq:hk-Kft},
\begin{equation}
  G(\omega,k) \;=\; \frac{1}{A + B\,k^{2} + \ii\,\omega}.
  \label{eq:hk-Gft}
\end{equation}

\subsection{Closing the $\omega$-contour: exponential decay in time}

As a function of complex $\omega$, the propagator \eqref{eq:hk-Gft} has a single
simple pole, at $\omega = \ii\,(A + Bk^{2})$. The inverse transform in time,
$G_{\mathrm R}(k,t) = \int \frac{\dd\omega}{2\pi}\,\ee^{-\ii\omega t}\,G(\omega,k)$,
is done by closing the integration contour in the half-plane where the
exponential decays and picking up that one residue. The result is a pure
exponential gated by causality:
\begin{equation}
  \Gret(k,t) \;=\; \theta(t)\,\ee^{-(A + B\,k^{2})\,t}.
  \label{eq:hk-GRkt}
\end{equation}
Here $\theta(t)$ is the Heaviside step --- the retarded (causal) propagator is
identically zero for $t<0$. It is convenient to name the $k$-dependent decay
rate the \term{$k$-dependent mass},
\[
  m(k) \;\equiv\; A + B\,k^{2} \qquad(\text{also written } \mu + D k^{2}),
\]
so that $\Gret(k,t)=\theta(t)\,\ee^{-m(k)t}$.

\section{The heat kernel}\label{sec:hk-kernel}

\subsection{Inverse-transforming $k\to x$}

The remaining transform is from momentum $k$ back to position $x$. The
$k$-dependence of \eqref{eq:hk-GRkt} is the Gaussian $\ee^{-Bk^{2}t}$, and its
inverse spatial Fourier transform is the classic Gaussian-integral identity
\begin{equation}
  \int \frac{\dd^{d}k}{(2\pi)^{d}}\;\ee^{\ii\,k\cdot x}\,\ee^{-B\,k^{2}t}
  \;=\;
  (4\pi B t)^{-d/2}\,\ee^{-|x|^{2}/(4Bt)}.
  \label{eq:hk-gaussianFT}
\end{equation}
This is the \term{heat kernel} --- the Green's function of the diffusion
equation. Multiplying back the $k$-independent decay $\ee^{-At}$ gives the full
real-space retarded propagator:
\begin{equation}
  \boxed{\;
  \Gret(t,x) \;=\; \theta(t)\,(4\pi B t)^{-d/2}\,
  \exp\!\Big[-\frac{|x|^{2}}{4 B t} - A\,t\Big].
  \;}
  \label{eq:hk-GR}
\end{equation}
The entire dependence on the spatial dimension $d$ lives in the prefactor
$(4\pi B t)^{-d/2}$; everything else is dimension-agnostic. Both $A$ and $B$ may
be complex, in which case \eqref{eq:hk-GR} is read as an analytic continuation
of the real formula. This is the single most important equation in the chapter,
and the function that evaluates it is \code{gaussian\_heat\_kernel}.

\begin{defn}
The \term{heat kernel} $(4\pi B t)^{-d/2}\,\ee^{-|x|^{2}/4Bt}$ is the
probability density, at time $t$, of a diffusion that started as a point source
at the origin at $t=0$ with diffusivity $B$. It is the real-space image of a
momentum-space Gaussian. In this subsystem ``analytic IFT'' (inverse Fourier
transform) means exactly this: doing the $k\to x$ transform in closed form via
\eqref{eq:hk-gaussianFT}, rather than sampling the correlator on a grid of
momenta and transforming numerically.
\end{defn}

\subsection{The code: \code{gaussian\_heat\_kernel}}

The closed form \eqref{eq:hk-GR} is implemented essentially verbatim
(\file{heat\_kernel.py:58}). The function takes a scalar time \code{t} (which
may be complex), a scalar separation \code{x} (in $d=1$; for $d\ge 2$ you pass
the radial $|x|$), the mass \code{A} and diffusion \code{B} (both complex-ok),
the dimension \code{spatial\_dim}, and an optional drift \code{V} we return to
in \S\ref{sec:hk-drift}. It returns a Python \code{complex}.

\begin{lstlisting}
def gaussian_heat_kernel(t, x, A, B, spatial_dim: int = 1, V=0.0):
    t = float(t.real) if hasattr(t, 'real') and not isinstance(t, complex) else t
    if (t.real if isinstance(t, complex) else t) <= 0.0:
        return 0j
    A = complex(A); B = complex(B); V = complex(V)
    x_signed = float(x); xx = x_signed ** 2
    pref = (4.0 * math.pi * B * t) ** (-0.5 * spatial_dim)
    drift_exp = 0.0
    if V != 0:
        if spatial_dim != 1:
            raise SpatialPropagatorError(
                'drift heat kernel (V!=0) is implemented for d=1 only ...')
        drift_exp = -1j * V * x_signed / (2.0 * B) + V * V * t / (4.0 * B)
    return complex(pref * cmath.exp(-xx / (4.0 * B * t) - A * t + drift_exp))
\end{lstlisting}

Read it line by line. The first line coerces \code{t} to a real float when it
is a real-like object that is not literally a Python \code{complex} (Sage and
mpmath both produce ``real-like'' objects with a \code{.real} attribute). The
second line is the $\theta(t)$ \term{causality guard}: if the real part of
\code{t} is non-positive the kernel is zero and we return early. The prefactor
line is exactly $(4\pi B t)^{-d/2}$ with $d=$ \code{spatial\_dim}. The final
\code{return} assembles $\text{pref}\cdot\exp(-x^{2}/4Bt - A t + \text{drift})$,
using \code{cmath.exp} (the complex exponential, because $A$ and $B$ can be
complex). The library imports are deliberately minimal --- \code{cmath} and
\code{math} from the standard library (\file{heat\_kernel.py:42--43}) ---
because for the scalar kernel there is nothing exotic to do.

\begin{note}
\textbf{What is \code{cmath}?} It is Python's standard-library module for
complex-valued elementary functions: \code{cmath.exp}, \code{cmath.sqrt}, and so
on. The ordinary \code{math.exp} accepts only real arguments and would raise on
a complex one. Because the mass $A$ and diffusion $B$ here can be complex (a
complex saddle, a Keldysh combination), the kernel must use the complex
exponential. The prefactor $(4\pi B t)^{-d/2}$ is likewise computed with
Python's complex power, so $B<0$ or complex $B$ produces the correct branch
rather than a domain error.
\end{note}

\section{Drift and advection}\label{sec:hk-drift}

A model can carry a first-spatial-derivative term $\partial_x\phi$ ---
advection, a velocity. Under the substitution $\partial_x \to \ii k$ it appears
in the inverse propagator as a term \emph{linear} in $k$, so \eqref{eq:hk-Kft}
generalises to
\begin{equation}
  K(\omega,k) \;=\; A \;+\; V\,k \;+\; B\,k^{2} \;+\; \ii\,\omega,
  \label{eq:hk-Kdrift}
\end{equation}
with $V$ the \term{drift} (the $k^{1}$ coefficient). Completing the square in
the exponent $\ee^{-(A + Vk + Bk^{2})t}$ before doing the Gaussian integral
\eqref{eq:hk-gaussianFT} produces an extra factor in the real-space kernel,
\begin{equation}
  \exp\!\Big[-\frac{\ii V x}{2B} + \frac{V^{2} t}{4B}\Big].
  \label{eq:hk-driftfactor}
\end{equation}
For a drift of the form $V=\ii v$ (the physical advection velocity $v$) this
factor \emph{shifts the centre of the Gaussian to $x = v\,t$} --- the packet is
transported at velocity $v$ while it spreads. In \code{gaussian\_heat\_kernel}
this is the \code{drift\_exp} term, added inside the exponential
(\file{heat\_kernel.py:89}).

Two facts about drift matter for the rest of the manual:

\begin{itemize}
  \item \textbf{$V=0$ is bit-identical to the pure heat kernel.} When the drift
        is zero the \code{if V != 0} branch is skipped entirely and the function
        returns exactly \eqref{eq:hk-GR}. This is not an approximation; it is the
        same arithmetic.
  \item \textbf{Drift is $d=1$ only.} A velocity is a \emph{vector} in $d\ge 2$,
        which is outside the v1 scope, so \code{gaussian\_heat\_kernel} raises a
        \code{SpatialPropagatorError} if it is handed $V\ne 0$ with
        \code{spatial\_dim != 1} (\file{heat\_kernel.py:85--88}).
\end{itemize}

\begin{gotcha}
For a \emph{gradient nonlinearity} --- Burgers $\partial_x(\phi^{2})$, KPZ
$(\partial_x\phi)^{2}$ --- you might expect the propagator to carry a drift. It
usually does not, and that is correct. The only $\partial_x$ that reaches the
bilinear (propagator) sector is the saddle cross-term, which is proportional to
the saddle value $\phi_*$. At the homogeneous saddle $\phi_*=0$ that term
vanishes, so $V\to 0$ and the propagator is the \emph{pure} heat kernel even for
KPZ or Burgers. Do not be surprised to see a ``KPZ'' propagator print
$V(\text{drift})=0$; it is physically right at the trivial saddle.
\end{gotcha}

\section{Periodic boundaries: the method of images}\label{sec:hk-images}

On an infinite line the kernel is \eqref{eq:hk-GR}. On a ring --- a periodic box
of period $L$ --- the Green's function must itself be periodic, and the standard
construction is the \term{method of images}: sum the infinite-domain kernel over
all image sources displaced by integer multiples of $L$,
\begin{equation}
  G_{\text{periodic}}(t,x) \;=\; \sum_{n\in\mathbb{Z}} G_{\infty}(t,\,x + nL).
  \label{eq:hk-imagesum}
\end{equation}
Because the Gaussian tail decays \emph{super-exponentially} in $n$ (the exponent
scales like $-(nL)^{2}$), only a handful of terms matter and the sum truncates
almost immediately. This is the function \code{image\_sum}
(\file{heat\_kernel.py:132}):

\begin{lstlisting}
def image_sum(t, x, A, B, L, spatial_dim=1, eps=1e-12, n_max_cap=2000, V=0.0):
    if spatial_dim != 1:
        raise SpatialPropagatorError(
            'image_sum: periodic BC implemented for d=1 only in v1.')
    base = gaussian_heat_kernel(t, x, A, B, spatial_dim=1, V=V)
    total = base
    ref = abs(base) if abs(base) > 0 else 1.0
    n = 1
    while n <= n_max_cap:
        term_p = gaussian_heat_kernel(t, x + n * L, A, B, spatial_dim=1, V=V)
        term_m = gaussian_heat_kernel(t, x - n * L, A, B, spatial_dim=1, V=V)
        total += term_p + term_m
        if abs(term_p) < eps * ref and abs(term_m) < eps * ref and n >= 2:
            break
        n += 1
    return complex(total)
\end{lstlisting}

The $n=0$ term \code{base} sets a reference magnitude \code{ref}; the loop adds
the $\pm n$ images in pairs and stops once both have fallen below
\code{eps}~$\times$~\code{ref} (default $10^{-12}$) \emph{and} at least two
shells have been added (the \code{n >= 2} guard avoids a premature stop at a
node where a single term happens to be tiny). A hard cap \code{n\_max\_cap}
prevents a runaway. Like drift, periodic boundaries are $d=1$ only here:
$d\ge 2$ periodic transforms would need a lattice sum per axis, and the function
raises otherwise (\file{heat\_kernel.py:143--145}). The $d\ge 2$ periodic case
is handled instead by a separate lattice-sum transform, \code{periodic\_inverse\_ft}, in \S\ref{sec:hk-numerical}.

\section{Reading $(A,B,V)$ off the inverse propagator with Sage}\label{sec:hk-extract}

So far we have assumed the three coefficients $(A,B,V)$ are known. In practice
they arrive as \emph{symbolic} expressions buried inside the inverse-propagator
matrix \code{K\_ft} that the shared \code{build\_propagator} produces upstream:
a single diagonal entry might look like
\code{mu + 3*lambda*phistar1\^{}2 - D*Laplacian + I*omega}. Peeling
$(A,B,V)$ out of that is a small computer-algebra task, and the tool for it is
SageMath.

\begin{note}
\textbf{What is SageMath?} Sage is a large open-source mathematics system built
on top of Python; it bundles dozens of mathematical libraries behind one Python
API. The only piece this subsystem uses is Sage's \term{Symbolic Ring}, written
\code{SR} --- a computer-algebra engine for symbolic expressions, in the spirit
of Mathematica but callable from Python. A Sage symbolic expression can hold
unknowns (\code{mu}, \code{D}, \code{phistar1}, \code{k}, \code{omega}), be
expanded and substituted into, and have its polynomial coefficients read off.
The imports are \code{from sage.all import SR, I as SR\_I, matrix}
(\file{heat\_kernel.py:47}): \code{SR} is the ring itself, \code{SR\_I} is the
symbolic imaginary unit $\ii$ (renamed so it does not clash with loop counters
named \code{i} or \code{I}), and \code{matrix} builds a matrix whose entries
live in \code{SR}. The reason this layer stays in Sage rather than switching to
sympy is that \code{build\_propagator} already hands it a Sage matrix, so
staying in Sage avoids a translation round-trip.
\end{note}

\subsection{The single-entry extractor}

The workhorse is \code{extract\_mass\_diffusion} (\file{heat\_kernel.py:161}),
which takes one diagonal entry of \code{K\_ft} together with the relevant
symbols and returns the triple $(A,B,V)$. The substitution and guards are worth
quoting:

\begin{lstlisting}
subs = {lap_sym: -k_var**2}
if grad_sym is not None:
    subs[grad_sym] = SR_I * k_var          # d_x -> i k  (the drift symbol)
e = SR(kft_entry).subs(subs).expand()
if e.has(lap_sym) or (grad_sym is not None and e.has(grad_sym)):
    raise SpatialPropagatorError(...)      # a residual operator symbol survived

omega_coeff = e.coefficient(omega, 1)      # the coefficient of omega^1
if (omega_coeff - SR_I).simplify_full() != 0:
    raise SpatialPropagatorError(...)      # not +i*omega-normalized

e0 = e.coefficient(omega, 0)               # the omega-independent part
deg = e0.degree(k_var)
if deg > 2:
    raise SpatialPropagatorError(...)      # higher-derivative operator
B = e0.coefficient(k_var, 2)               # k^2 coefficient = diffusion
V = e0.coefficient(k_var, 1)               # k^1 coefficient = drift
A = e0.coefficient(k_var, 0)               # k^0 coefficient = mass
return A, B, V
\end{lstlisting}

Each Sage call earns its keep:

\begin{description}
  \item[\code{.subs(dict)}] substitutes one expression for another. Here the
        inert Laplacian symbol becomes $-k^{2}$ and, if a first-derivative symbol
        \code{grad\_sym} is present, it becomes $\ii k$. After the substitution
        the entry is an ordinary polynomial in $k$ and $\omega$.
  \item[\code{.expand()}] multiplies everything out so the coefficients can be
        read individually.
  \item[\code{.has(sym)}] is \code{True} iff the expression still contains
        \code{sym}. If a Laplacian (or gradient) symbol survives the
        substitution, the entry encoded a higher-derivative operator (a
        $(\Lap)^{2}$ would leave a $k^{4}$ and a residual symbol), and the function
        rejects it.
  \item[\code{.coefficient(var, n)}] returns the coefficient of $\code{var}^{n}$.
        This is how $A,B,V$ and the $\ii\omega$ coefficient are picked off.
  \item[\code{.simplify\_full()}] runs Sage's strongest simplifier so that the
        normalization test $(\text{omega\_coeff}-\ii)$ genuinely collapses to
        \code{0} when it should. A naive subtraction might leave an
        unsimplified expression that is nonzero only syntactically.
  \item[\code{.degree(k\_var)}] returns the highest power of $k$. The guard
        \code{deg > 2} is what enforces ``no higher-derivative operators'' (see
        \S\ref{sec:hk-guards}).
\end{description}

The mapping from these three coefficients to physics is exactly the dictionary
of \S\ref{sec:hk-fromprop}: \code{B} is the $k^{2}$ coefficient (diffusion),
\code{V} the $k^{1}$ coefficient (drift), \code{A} the $k^{0}$ coefficient
(mass). For the worked example above, $A=\code{mu + 3*lambda*phistar1\^{}2}$,
$B=\code{D}$, and $V=0$.

\begin{gotcha}
The gradient symbol is fetched as \code{grad\_sym = SR.var('GradX')}
(\file{heat\_kernel.py:360}). Sage \emph{caches symbolic variables by name}: any
two calls to \code{SR.var('GradX')} anywhere in the codebase return the
\emph{same} object. This is load-bearing --- the gradient symbol declared when
the operator IR is lowered and the one declared here must be identical, or the
\code{.subs} would silently fail to match. If you ever rename it in one place
you must rename it everywhere.
\end{gotcha}

\subsection{The coupled-field generalization}

\code{extract\_mass\_diffusion} reads \emph{one} diagonal entry. Its sibling
\code{reaction\_diffusion\_matrices} (\file{heat\_kernel.py:226}) does the same
for the \emph{full} matrix, returning $N\times N$ Sage matrices $(M,\mathcal{D},V)$
of masses, diffusions, and drifts. The extra requirement at the matrix level is
that the time derivative be \emph{diagonal and unit-normalized}: the $\ii\omega$
coefficient must be $\ii$ on the diagonal and exactly $0$ off it
(\file{heat\_kernel.py:262--268}). For a diagonal \code{K\_ft} this reduces, per
entry, to \code{extract\_mass\_diffusion} bit-identically. These matrices feed
the coupled-field spectral propagator (the subject of
Chapter~\ref{ch:spatial-coupled}); the diagonal Tier-1 path we are documenting
here does \emph{not} consume them.

\section{From propagator to runtime: build versus reconstruct}\label{sec:hk-build}

\subsection{Building the picklable spatial block}

\code{build\_spatial\_propagator} (\file{heat\_kernel.py:310}) is the assembler.
It is handed the inverse-propagator matrix \code{K\_ft}, the $\omega$ symbol, the
namespace \code{ns} that carries \code{ns.Laplacian}, the model dict, and the
field-name lists. It reads the spatial dimension, declares the $k$ and
\code{GradX} symbols, runs the diagonality and boundary guards, calls
\code{extract\_mass\_diffusion} on each diagonal entry, and returns a plain
dictionary --- the \emph{spatial block} of the propagator dict. The single most
important property of that dictionary is that it is \textbf{picklable}: it holds
only Sage expressions and plain data, no closures, so it caches cleanly through
the pipeline's \code{PipelineCache}.

The block it returns (\file{heat\_kernel.py:421--431}) has these keys:

\begin{center}
\begin{longtable}{@{}p{0.27\textwidth} p{0.66\textwidth}@{}}
\toprule
\textbf{key} & \textbf{meaning} \\
\midrule
\code{G\_tx\_sym} & \code{dict[(i,j)] -> (A\_expr, B\_expr)} or \code{None};
  per-entry symbolic mass/diffusion; \code{None} off-diagonal. \\
\code{k\_var} & the Sage momentum symbol $k$. \\
\code{spatial\_dim} & the dimension $d\in\{1,2,3\}$. \\
\code{bc\_mode} & \code{'infinite'} or \code{'periodic'}. \\
\code{bc\_params} & boundary parameters (e.g.\ \code{\{'length': 'L'\}}),
  with \code{'mode'} stripped out. \\
\code{initial\_mode} & \code{'stationary'} (default), etc. \\
\code{ac\_mass} & \code{dict[i] -> A\_expr}; per-field mass. \\
\code{ac\_diffusion} & \code{dict[i] -> B\_expr}; per-field diffusion. \\
\code{ac\_drift} & \code{dict[i] -> V\_expr}; per-field drift ($0$ means a pure
  heat kernel). \\
\bottomrule
\end{longtable}
\end{center}

Note that \code{G\_tx\_sym} keeps each diagonal entry as a 2-tuple $(A,B)$ for
backward compatibility; the drift $V$ travels separately in \code{ac\_drift}.
The off-diagonal entries of \code{G\_tx\_sym} are set to \code{None}, and the
diagonality guard (\file{heat\_kernel.py:369--377}) guarantees there is nothing
there to lose --- any nonzero off-diagonal \code{K\_ft[i,j]} raises a
\code{SpatialPropagatorError} before we ever get here, because the closed-form
Tier-1 kernel requires each field to be its own independent Allen--Cahn mode.

\subsection{Why the callables are not stored}

The one thing the block does \emph{not} contain is the runtime propagator
function itself --- the thing you actually call as $G(t,x)$. The reason is
purely mechanical: \textbf{closures do not pickle}. A function that has captured
\code{A\_num}, \code{B\_num}, and the boundary mode in its enclosing scope cannot
be serialized, so it cannot live in a cacheable dict. The fix is to store only
the symbolic data (which \emph{does} pickle) and rebuild the callables on demand.

That rebuild is \code{make\_g\_tx\_callables} (\file{heat\_kernel.py:434}). It is
idempotent and cheap, and it must be called after \code{build\_propagator}
returns --- whether that was a fresh build or a cache load. It walks
\code{G\_tx\_sym}, turns each symbolic $(A,B)$ (and the matching drift $V$) into
numeric callables via the helper \code{\_make\_numeric}, and wraps them in a
per-entry closure that dispatches on the boundary mode:

\begin{lstlisting}
def _g(t, x, **num_params):
    A = A_num(**num_params)
    B = B_num(**num_params)
    V = V_num(**num_params) if V_num is not None else 0.0
    if bc_mode == 'periodic':
        ...
        L = (num_params[L_name] if isinstance(L_name, str) else float(L_name))
        return image_sum(t, x, A, B, float(L), spatial_dim=d, V=V)
    return gaussian_heat_kernel(t, x, A, B, spatial_dim=d, V=V)
\end{lstlisting}

The closure evaluates the symbolic coefficients at the supplied numeric
parameters and then calls \code{image\_sum} (periodic) or
\code{gaussian\_heat\_kernel} (infinite). The helper \code{\_make\_numeric}
(\file{heat\_kernel.py:287}) is the bridge from symbol to number: it reads an
expression's free symbols with \code{expr.variables()}, sorts them by name, and
returns a closure that substitutes each by name from a \code{num\_params} dict
and coerces the result to a Python \code{complex} --- raising a \code{KeyError}
with a helpful message if a parameter or saddle value is missing. Entries whose
symbol is \code{None} (the off-diagonal ones) get a closure that simply returns
\code{0j}.

\begin{example}\label{ex:hk-gtx}
A concrete round trip. After building a $\phi^4$ propagator at the trivial
saddle and reconstructing the callables, calling the diagonal entry as
\begin{center}
\code{G\_tx[(0,0)](t=0.5, x=1.0, mu=1.0, lambda=0.0, phistar1=0.0, D=0.1)}
\end{center}
substitutes $A=\mu=1$, $B=D=0.1$, $V=0$ into \eqref{eq:hk-GR} and returns
$(4\pi\cdot 0.1\cdot 0.5)^{-1/2}\,\exp[-1/(4\cdot 0.1\cdot 0.5) - 0.5]$.
\end{example}

\section{The tree-level correlator}\label{sec:hk-tree}

\subsection{Noise driving two retarded legs}

The retarded propagator \eqref{eq:hk-GR} is the response to a kick. The
\emph{correlation} function $C(x,\tau)=\avg{\phi(x_1,t_1)\,\phi(x_2,t_2)}$ (with
$x=x_1-x_2$, $\tau=t_1-t_2$) is, at tree level, the \emph{noise sourcing two
retarded propagators}:
\begin{equation}
  C_{ij}(x,\tau) \;=\; 2\,D^{\text{noise}}_{ij}
  \int_{-\infty}^{\min(t_1,t_2)}\!\dd t_v \int \dd x_v\;
  {\Gret}_i(t_1-t_v,\,x-x_v)\,{\Gret}_j(t_2-t_v,\,-x_v).
  \label{eq:hk-Ctwoleg}
\end{equation}
The noise, injected at the vertex $(x_v,t_v)$ and white in space and time,
propagates causally to both external points. For a diagonal heat-kernel
propagator the spatial integral over $x_v$ collapses by the heat-kernel
\term{semigroup} property --- the convolution of two Gaussians is one Gaussian
--- and the remaining $t_v$ integral reduces to a single one-dimensional
integral of error-function type,
\begin{equation}
  C_{ii}(x,\tau) \;=\; \frac{D^{\text{noise}}_i}{\sqrt{4\pi B_i}}
  \int_{|\tau|}^{\infty} s^{-1/2}\,
  \exp\!\Big[-\frac{x^{2}}{4 B_i s} - A_i s\Big]\,\dd s.
  \label{eq:hk-Cerf}
\end{equation}

\subsection{The error-function closed form}

The integral $\int s^{-1/2}\,\ee^{-\beta/s - \alpha s}\,\dd s$ in
\eqref{eq:hk-Cerf} has a genuine closed form. Substituting $s=w^{2}$, the
antiderivative of $\ee^{-\alpha w^{2} - \beta/w^{2}}$ is the ``erf-split'' form
\begin{equation}
  F(w) \;=\; \frac{\sqrt\pi}{4\sqrt\alpha}\Big[
    \ee^{+2\sqrt{\alpha\beta}}\,\mathrm{erf}\!\big(\sqrt\alpha\,w + \tfrac{\sqrt\beta}{w}\big)
    + \ee^{-2\sqrt{\alpha\beta}}\,\mathrm{erf}\!\big(\sqrt\alpha\,w - \tfrac{\sqrt\beta}{w}\big)
  \Big],
  \label{eq:hk-Fw}
\end{equation}
so the definite integral equals $2\,[F(\sqrt{U})-F(\sqrt{L})]$, with
$\alpha=$ mass, $\beta=x^{2}/4B\ge 0$. The semi-infinite case $U\to\infty$
(valid when $\operatorname{Re}\sqrt\alpha>0$) has the simple limit
$F(\infty)=\tfrac{\sqrt\pi}{4\sqrt\alpha}\big(\ee^{2\sqrt{\alpha\beta}}+\ee^{-2\sqrt{\alpha\beta}}\big)$.
This is the function \code{erf\_time\_integral} (\file{heat\_kernel.py:93}).

\begin{gotcha}
The form \eqref{eq:hk-Fw} contains a \emph{near-cancellation of two large
exponentials}: when $\sqrt{\alpha\beta}$ is large, $\ee^{+2\sqrt{\alpha\beta}}$
and $\ee^{-2\sqrt{\alpha\beta}}$ differ by many orders of magnitude, and double
precision loses most of its significant digits in the subtraction. For that
reason \code{erf\_time\_integral} is evaluated at \textbf{30 decimal digits}
using \code{mpmath} and only then collapsed back to a \code{complex}. If you
ever port this to plain floats, expect it to silently lose precision.
\end{gotcha}

\begin{note}
\textbf{What is \code{mpmath}?} It is a pure-Python \emph{arbitrary-precision}
floating-point library: where ordinary floats give about 16 significant digits,
mpmath gives 30, 50, or 1000 on demand, and its special functions (\code{erf},
\code{sqrt}, \code{exp}) work on \emph{complex} arguments. Sage ships it. The
import is \code{import mpmath as mp} (\file{heat\_kernel.py:46}), and the whole
high-precision computation runs inside a \code{with mp.workdps(dps):} block ---
\code{workdps} is a context manager that temporarily sets the working precision
to \code{dps} decimal digits. The complex numbers are \code{mp.mpc(...)}, reals
are \code{mp.mpf(...)}.
\end{note}

The implementation also handles the $w\to 0^{+}$ endpoint carefully. The helper
\code{\_F(w)} (\file{heat\_kernel.py:115--122}) special-cases \code{w == 0}:
$\mathrm{erf}(\pm\infty)=\pm 1$ when $\beta>0$, and $\mathrm{erf}(0)=0$ when
$\beta=0$. This is the $s\to 0$ endpoint that appears whenever the lower limit
$L$ is zero (the equal-time, $x=0$ corner), and getting it wrong would produce a
spurious $0/0$ right at the most physically interesting point of the correlator.

\subsection{\code{free\_two\_point}: one field, both boundary conditions}

\code{free\_two\_point} (\file{spatial\_correlator.py:138}) wraps the erf
integral into a per-field tree correlator. Its signature is
\code{free\_two\_point(mu, D, kap, x, tau, bc\_mode='infinite', L=None, ...)},
where \code{mu} is the mass, \code{D} the diffusion, and \code{kap} the noise
spectral coefficient:

\begin{lstlisting}
def free_two_point(mu, D, kap, x, tau, bc_mode='infinite', L=None, n_images_cap=200):
    D = float(D); mu = complex(mu); kapf = float(kap)
    pref = kapf / math.sqrt(4.0 * math.pi * D)
    def _C_inf(xx):
        beta = (xx * xx) / (4.0 * D)
        integ = erf_time_integral(mu, beta, abs(tau), U_hi=None)
        return complex(pref * integ)
    if bc_mode == 'infinite':
        return _C_inf(x)
    if bc_mode == 'periodic':
        ...                                # sum images _C_inf(x +/- nL)
\end{lstlisting}

For the infinite domain it returns the single erf integral with
$\beta=x^{2}/4D$ and lower limit $|\tau|$ (the $U\to\infty$ semi-infinite form,
\code{U\_hi=None}). For a periodic box it sums the erf correlator over images
$x\to x\pm nL$ with the same super-exponential truncation as \code{image\_sum}.

\begin{gotcha}
The first two arguments of \code{free\_two\_point} are named \code{(mu, D)}, but
its callers pass them the mass and diffusion \emph{they} call $(A,B)$. In
\code{compute\_spatial\_correlator\_tree} the call is
\code{free\_two\_point(A, B, Dn, ...)} (\file{spatial\_correlator.py:291}). This
is semantically correct --- mass maps to \code{mu}, diffusion maps to \code{D}
--- but the name change can read as a bug at a glance. The rule to remember:
\emph{the first argument is always the mass, the second is always the
diffusion}, whatever the local letter.
\end{gotcha}

\subsection{Assembling the $(x,\tau)$ grid}

\code{compute\_spatial\_correlator\_tree} (\file{spatial\_correlator.py:178}) is
the driver that fills the output array. It takes the expanded FieldTheory, the
model, the propagator dict, the numeric parameters, the two external legs, and
the $(\tau,x)$ grids, and returns a complex array of shape
\code{(len(tau\_grid), len(spatial\_grid))} together with an \code{info} dict.
Its logic, in order:

\begin{enumerate}
  \item Require a spatial block (\code{prop['spatial\_dim']}); normalize the
        numeric parameters to string keys; resolve the periodic length $L$ by
        name or inline value.
  \item \textbf{Resolve the field index} of the external legs by matching the
        leg base name to the physical column names (stripping trailing
        population digits).
  \item \textbf{Tier-1 guard:} if \code{prop['G\_tx\_sym'] is None} the theory
        is not the diagonal heat-kernel case, and the function raises a clear
        \code{NotImplementedError} explaining that off-diagonal coupling and
        higher-derivative operators are v2 features
        (\file{spatial\_correlator.py:239--248}).
  \item Read $(A,B)$ for the resolved diagonal entry, substitute parameters, and
        take $A$ complex, $B$ real.
  \item \textbf{Diffusion-sign guard:} raise on $B<0$ (anti-diffusive, the heat
        kernel diverges) and on $B=0$ (no Laplacian: a time-only field, for
        which a spatial correlator is undefined). See \S\ref{sec:hk-guards}.
  \item Read the noise coefficient $D_{\text{noise}}$ for this field
        (\S\ref{sec:hk-noise}); raise if absent.
  \item Fill \code{C[it, ix] = free\_two\_point(A, B, Dn, x, tau, ...)} over the
        grid.
\end{enumerate}

The static ($\tau=0$) limit of \eqref{eq:hk-Cerf} has its own closed form that
the erf result must reproduce --- $C(x,0)=\frac{D^{\text{noise}}}{2\sqrt{AB}}\,
\ee^{-|x|\sqrt{A/B}}$ in $d=1$ --- and the function
\code{free\_correlator\_static\_closed\_form} (\file{spatial\_correlator.py:398})
supplies it in all three dimensions: the $d=1$ exponential, the $d=2$ modified
Bessel $K_0$, and the $d=3$ Yukawa $\ee^{-\kappa r}/r$. It serves as an
independent oracle for the transforms of \S\ref{sec:hk-numerical}.

\section{Reading the noise strength out of the action}\label{sec:hk-noise}

The correlator \eqref{eq:hk-Cerf} needs the noise strength $D^{\text{noise}}$,
and that lives in the action, not in the propagator. In the \msrjd{} action a
white-noise covariance $\avg{\xi\xi}=2\kappa\,\delta$ appears as a term
$-D^{\text{noise}}\,\tilde\phi^{2}$ --- a term with \emph{two response fields and
zero physical fields}. The framework classifies every action term by its
``bigrade'' $(\text{number of response factors},\,\text{number of physical
factors})$, so the noise sector is precisely the $(2,0)$ bigrade.

\code{extract\_noise\_coefficients} (\file{spatial\_correlator.py:46}) reads it.
It fetches the $(2,0)$ sector from the expanded FieldTheory
(\code{ft.\_by\_tp[(2,0)]}), iterates the monomials of that polynomial, and
identifies a \emph{pure} $\tilde\phi_i^{2}$ monomial as one with exponent $2$ on
ring position $i$ (response fields are the first \code{n\_tilde} generators of
the ring) and $0$ everywhere else. For such a monomial the white-noise
coefficient is $D^{\text{noise}}_i = -\operatorname{Re}(\text{coeff})$ --- the
minus sign because the action term is $-D^{\text{noise}}\tilde\phi^{2}$. The
return value is a dict \code{\{field\_index: D\_noise\}}. This convention is
exactly the one that reproduces the scalar OU result $\avg{x^{2}}=D/\mu$ from an
action term $-D\,\tilde x^{2}$.

Its coupled-field sibling \code{extract\_noise\_matrix}
(\file{spatial\_correlator.py:93}) returns the full $n_{\tilde\phi}\times
n_{\tilde\phi}$ noise-covariance matrix $N$ with the convention
$N_{ii}=-2\,\mathrm{coeff}(\tilde\phi_i^{2})$ and
$N_{ij}=-\mathrm{coeff}(\tilde\phi_i\tilde\phi_j)$ for $i\ne j$, so that
$N_{ii}=2\kappa_i$ with $\kappa_i$ the value the diagonal extractor returns.
Both extractors read only the $q^{0}$ (white-noise) part; $q^{2}$-dependent
(conserved, Model-B) noise is skipped and left at zero, matching the diagonal
heat-kernel scope.

\section{The analytic IFT at loop order}\label{sec:hk-loops}

Everything to here has been tree level. At loop order the correlator picks up a
self-energy correction $\delta C$, and the central question becomes how to do the
spatial transform of a \emph{loop} diagram. The answer is the same trick at a
larger scale: do the loop momentum integral analytically (it is Gaussian), do
the external $q\to x$ transform analytically (it is again Gaussian), and never
build a $q$-grid.

\subsection{The per-chamber Gaussian}

A loop diagram, after its loop momentum integral is done in Schwinger
parameters (the machinery of \S\ref{sec:hk-symanzik} and the spatial-core
chapter), leaves a residual Gaussian in the \emph{external} momentum $q$ of the
form $\ee^{-q^{\mathsf T}\mathcal{B}\,q}$, where for each Schwinger sample
$\mathcal{B}=D\cdot Q_{\text{eff}}$ (a scalar for a two-point function, an
$n\times n$ matrix for $k\ge 3$ external legs). Inverse-transforming that
Gaussian to real space is once more the heat-kernel identity
\eqref{eq:hk-gaussianFT}. The production code does this for a whole \emph{batch}
of Schwinger samples at once, vectorized, in \code{\_heat\_kernel\_x\_general}
(\file{full\_integrator.py:143}):

\begin{lstlisting}
def _heat_kernel_x_general(Bcal_g, xs_arr, spatial_dim):
    d = float(spatial_dim)
    xs_arr = np.asarray(xs_arr, dtype=float)
    if Bcal_g.ndim == 1:                                     # n_ext = 1 (k=2)
        if xs_arr.ndim == 1:
            x2 = xs_arr ** 2
        else:
            x2 = np.sum(xs_arr ** 2, axis=-1)
        return ((4.0 * math.pi * Bcal_g)[:, None] ** (-0.5 * d)
                * np.exp(-x2[None, :] / (4.0 * Bcal_g[:, None])))
    n = Bcal_g.shape[1]
    ...
    detB = np.linalg.det(Bcal_g)                             # (P',)
    Binv = np.linalg.inv(Bcal_g)                             # (P', n, n)
    if xs_arr.ndim == 2:                                     # d=1: (n_x, n)
        quad = np.einsum('xj,pjk,xk->px', xs_arr, Binv, xs_arr)
    else:                                                    # d>=2: (n_x, n, d)
        quad = np.einsum('xjc,pjk,xkc->px', xs_arr, Binv, xs_arr)
    return ((4.0 * math.pi) ** (-0.5 * n * d)
            * detB[:, None] ** (-0.5 * d)
            * np.exp(-0.25 * quad))
\end{lstlisting}

For a two-point function (\code{n\_ext = 1}) this is exactly the body of
\code{gaussian\_heat\_kernel} --- $(4\pi B)^{-d/2}\,\ee^{-|x|^{2}/4B}$ --- with
the $\theta(t)$, drift, and time pieces already folded into $\mathcal{B}$
upstream, and with the batch axis $P'$ over Schwinger samples vectorized. For
$k\ge 3$ external legs (\code{n\_ext >= 2}) it is the \emph{multivariate}
Gaussian inverse transform
\[
  \text{hk} \;=\; (4\pi)^{-nd/2}\,\det(\mathcal{B})^{-d/2}\,
  \exp\!\Big[-\tfrac14\sum_c \vec X_c^{\mathsf T}\,\mathcal{B}^{-1}\,\vec X_c\Big],
\]
one factor per spatial component $c$ (the components share the same
$\mathcal{B}$ by isotropy). The determinant and inverse are batched with
\code{np.linalg.det} / \code{np.linalg.inv}, and the quadratic form is contracted
with \code{np.einsum}.

\begin{note}
\textbf{What is \code{np.einsum}?} It is NumPy's Einstein-summation routine:
you write the index pattern of a tensor contraction as a string and it performs
exactly that sum. Here \code{'xj,pjk,xk->px'} reads ``contract the evaluation
point $x$'s component $j$ with $\mathcal{B}^{-1}_{jk}$ and the component $k$,
summing over $j$ and $k$, for every sample $p$,'' producing the array of
quadratic forms $\vec X^{\mathsf T}\mathcal{B}^{-1}\vec X$ indexed by sample $p$
and point $x$. The \code{[:, None]} index inserts a length-one axis so the
per-sample prefactor broadcasts against the per-point exponential.
\end{note}

\subsection{The gate: \code{\_use\_analytic}}

Whether the analytic transform is used at all is decided by one line in the
orchestrator \file{pipeline\_bridge.py}. The relevant flags
(\file{pipeline\_bridge.py:1667,1677}) are:

\begin{lstlisting}[style=console]
_all_plain    = all(rec[2] is None for rec in live_g)   # no derivative vertices
_force_num    = _os.environ.get('SPATIAL_FORCE_NUMERICAL_FT', '') == '1'
_use_analytic = (_all_plain or d == 1) and not _force_num
\end{lstlisting}

So the analytic heat-kernel transform covers two regimes: \textbf{(a)} plain
(derivative-free) vertices at any dimension $d$, and \textbf{(b)} any vertices at
$d=1$, where the derivative-vertex form factors are handled by a polynomial fit
plus closed-form heat-kernel $q$-moments. The one remaining gap is $d\ge 2$
\emph{with} derivative vertices, which still falls back to the numerical
transform of \S\ref{sec:hk-numerical}. When \code{\_use\_analytic} is true, each
diagram's per-chamber Schwinger samples produce $\mathcal{B}=D\cdot Q_{\text{eff}}$,
which is fed straight to \code{\_heat\_kernel\_x\_general} to get $\delta C(x,\tau)$
directly --- with no $q$-grid anywhere.

\section{The numerical-transform cross-check and the env knobs}\label{sec:hk-numerical}

The analytic transform is exact, but ``exact'' is a claim that has to be
\emph{checked}, and the check is a second, independent route: do the $q\to x$
transform numerically on a discretized momentum grid. That route lives in
\code{radial\_inverse\_ft} and \code{periodic\_inverse\_ft}
(\file{spatial\_correlator.py:299,341}).

\code{radial\_inverse\_ft} takes an isotropic momentum-space correlator
$C(|q|)$ sampled on a $q$-grid and transforms it to $C(|x|)$ via the radial
inverse transform appropriate to the dimension:
\begin{align*}
  d=1:&\quad C(x) = \tfrac{1}{\pi}\int_0^{\infty}\cos(qx)\,C(q)\,\dd q,\\
  d=2:&\quad C(x) = \tfrac{1}{2\pi}\int_0^{\infty} q\,J_0(qx)\,C(q)\,\dd q
       \quad(\text{Hankel order }0),\\
  d=3:&\quad C(x) = \tfrac{1}{2\pi^{2}x}\int_0^{\infty} q\,\sin(qx)\,C(q)\,\dd q,
\end{align*}
each evaluated as a \code{np.trapz} (trapezoidal numerical integration) over the
grid, with the $d=3$ $x\to 0$ limit ($\sin(qx)/x\to q$) handled explicitly. The
$d=2$ branch uses the Bessel function $J_0$ from \code{scipy.special.j0}.
\code{periodic\_inverse\_ft} is the discrete-momentum analogue for a periodic
cubic box: it builds the integer-mode lattice with \code{np.meshgrid},
interpolates the radial $C(q)$ onto the lattice magnitudes with \code{np.interp},
zeros out modes past the cutoff with \code{np.where}, and sums
$\cos(k_x x)\,C(|k_n|)$ with a $1/L^{d}$ prefactor. The two boundary conditions
share one momentum-space correlator and differ only in the transform.

\begin{gotcha}
Both numerical transforms \textbf{truncate at $k_{\max}=q\_grid[-1]$} --- they
evaluate the correlator \emph{at} a finite momentum cutoff, not in the
continuum. The continuum value is recovered only in the $k_{\max}\to\infty$,
fine-grid limit. This is the ``Regime 1'' (physical-cutoff) interpretation. The
\emph{analytic} heat-kernel transform has no such truncation --- that is the
whole point of ``retiring the $q$-grid'' --- so the two routes agree only as the
numerical grid is refined.
\end{gotcha}

Three environment variables, all read in \file{pipeline\_bridge.py}, control the
cross-check:

\begin{description}
  \item[\code{SPATIAL\_FORCE\_NUMERICAL\_FT}] (\file{pipeline\_bridge.py:610,1676})
        forces the numerical-transform path even where the analytic transform
        would apply (it is the \code{not \_force\_num} term in the gate above).
        This keeps the cosine/radial cross-check reachable. It is exact only in
        the $q\_cut\to\infty$, $n\_q\to\infty$ limit; use it for validation, not
        production.
  \item[\code{SPATIAL\_Q\_CUT}] (\file{pipeline\_bridge.py:1604}) overrides the
        $q$-grid upper cutoff for the numerical cross-check. The in-code comment
        warns that derivative-vertex form factors with a $q^{4}$ tail need a
        large cutoff for the truncated trapezoid to converge --- a limit the
        analytic path does not have.
  \item[\code{SPATIAL\_N\_Q}] (\file{pipeline\_bridge.py:1605}) overrides the
        number of $q$-grid points for the cross-check.
\end{description}

\begin{gotcha}
A separate guard on the same code path can refuse to run at all. Before
allocating a chamber's $n_t^{\,n_V}\cdot n_s^{\,n_C}\cdot n_x$ complex array,
\file{pipeline\_bridge.py:1641--1665} checks it against
\code{SPATIAL\_MEM\_BUDGET\_GB} (default 6\,GB) and aborts up front if it would
exceed the budget. This is a deliberate guard against an out-of-memory event
that can crash not just the kernel but the OS --- the same macOS hazard that
forced the spatial path to be thread-based rather than fork-based. It is not in
the four primary files of this chapter, but it governs whether the analytic
transform even gets to run.
\end{gotcha}

\section{The Symanzik momentum core (oracle path)}\label{sec:hk-symanzik}

We now reach the two modules that carry the \code{ORACLE-ONLY} banner. They
implement the loop \emph{momentum} integral in Schwinger parameters as an
independent numerical reference, and they exist to pin the production
\code{full\_integrator} numbers in tests. \emph{They are not on the production
path}: \code{compute\_cumulants} does not call them.

\subsection{Why momentum-first}

The motivation (stated in the module docstring,
\file{loop\_parametric.py:8--31}) is precision. Doing the loop integral
time-first reuses the temporal Phase~J machinery plus a numerical $\int\dd\ell$,
but spatial loops \emph{generically} trip the $m\ge 3$ close-pair precision slow
path, because the loop momentum sweeps the edge masses past one another. The cure
is to do the momentum integral \emph{analytically first}. With a Schwinger
parameter $w_e$ on each edge, every edge is $\ee^{-(\mu + D k_e^{2})w_e}$, and if
each edge momentum is routed as $k_e = a_e\,\ell + b_e\,q$ (one loop momentum
$\ell$, external $q$), the loop integral is a pure Gaussian:
\begin{equation}
  \int \frac{\dd^{d}\ell}{(2\pi)^{d}}\;\ee^{-D\sum_e w_e k_e^{2}}
  \;=\;
  (4\pi D U)^{-d/2}\,\ee^{-D q^{2}(W - V^{2}/U)},
  \label{eq:hk-symanzik1}
\end{equation}
with the \term{Symanzik forms}
\[
  U = \sum_e a_e^{2} w_e,\qquad
  V = \sum_e a_e b_e w_e,\qquad
  W = \sum_e b_e^{2} w_e,\qquad
  F_{\text{reduced}} = W - V^{2}/U.
\]
There are \emph{no momentum-dependent poles} here --- only a smooth
Schwinger/time integral remains, so close-pair singularities can never arise.
This is the single conceptual reason the spatial loop integrator works at all.

\subsection{\code{loop\_parametric.py} and \code{spatial\_reduce.py}}

\code{loop\_parametric.symanzik\_UF} (\file{loop\_parametric.py:41}) is the
$L=1$ core: it computes $U,V,W,F_{\text{reduced}}$ and the prefactor
$(4\pi D U)^{-d/2}$, raising on $U\le 0$ (which would mean no loop-momentum
damping). \code{gaussian\_momentum\_integral} (\file{loop\_parametric.py:73})
is the thin wrapper $\text{pref}\cdot\ee^{-Dq^{2}F_{\text{reduced}}}$. The
$L$-loop generalization lives one module over in \code{spatial\_reduce.py},
which replaces the scalars by matrices --- $\Lambda$ $(L\times L)$, $N$
$(L\times n_{\text{ext}})$, $Q$ $(n_{\text{ext}}\times n_{\text{ext}})$ --- with
$U=\det\Lambda$ and the reduced external form
$Q_{\text{eff}}=Q-N^{\mathsf T}\Lambda^{-1}N$. Its \code{momentum\_integral}
(\file{spatial\_reduce.py:101}) evaluates
\[
  \int \prod_i\frac{\dd^{d}\ell_i}{(2\pi)^{d}}\;\ee^{-D\sum_e w_e k_e^{2}}
  = (4\pi D)^{-Ld/2}\,U^{-d/2}\,\ee^{-D q^{\mathsf T} Q_{\text{eff}} q},
\]
using \code{np.linalg.det} for $U$ and \code{np.linalg.solve} to form
$\Lambda^{-1}N$ without explicitly inverting $\Lambda$. This is exactly the
$Q_{\text{eff}}$ that, multiplied by $D$, becomes the $\mathcal{B}$ fed to the
production transform of \S\ref{sec:hk-loops} --- the oracle and the production
path compute the \emph{same} object two different ways.

\subsection{\code{temporal\_integrate.py}: the causal time integral}

After the momentum integral, the residual integral is over the edges' time
(Schwinger) parameters with the causal structure. For a 2-vertex self-energy
where all internal edges span the same inter-vertex time $t$,
\code{sigma\_parametric} (\file{temporal\_integrate.py:56}) assembles
\[
  \Sigma(q,t) \;=\; T^{n_C}\int_{[|t|,\infty)^{n_C}}\!\prod\dd s_C\;
  \ee^{-\mu\sum_e w_e}\;I_{\text{mom}}\big(\{(a_e,b_e)\},\{w_e\},q\big),
\]
where a \emph{retarded} edge has fixed duration $w_e=t$ (from $\theta(t)$), and a
\emph{correlation} edge carries a Schwinger parameter $w_e=s_e$ integrated over
$[|t|,\infty)$, each contributing a factor $T$ (temperature). The combinatorial
factor $\Sym(\Gamma)$ is applied by the caller, not here. The dispatch on the
number of correlation edges $n_C$ is the interesting part: $n_C=0$ evaluates
directly, $n_C=1$ uses an adaptive \code{scipy.integrate.quad}, and $n_C\ge 2$
uses a \emph{Gauss--Laguerre tensor rule} over the correlation parameters.

\begin{note}
\textbf{Gauss--Laguerre quadrature} integrates $\int_0^{\infty}\ee^{-x}g(x)\,\dd x$
\emph{exactly} for polynomial $g$ up to degree $2\,\text{deg}-1$, using the nodes
and weights returned by \code{np.polynomial.laguerre.laggauss(deg)}. Because the
Schwinger integrand carries an $\ee^{-\mu s}$ weight, the substitution
$s_C = lo + x_k/\mu$ turns each correlation-edge integral into exactly that form,
and the $n_C\ge 2$ case takes the Cartesian product of node indices with
\code{itertools.product}. This is roughly 10--100$\times$ faster than the
adaptive \code{dblquad}/\code{nquad} it replaces --- the two-correlation-edge
\code{dblquad} in \code{loop\_parametric.sigma\_K\_kernel} was the bottleneck of
the older C stack.
\end{note}

The convenience edge-specs \code{bubble\_edges} and \code{sunset\_edges}
(\file{temporal\_integrate.py:127,135}) supply the routing triples for the
validated 2-vertex topologies (the 1-loop bubble and the 2-loop sunset), and
\code{bubble\_delta\_equal\_time\_via\_C} (\file{temporal\_integrate.py:144})
collapses the tabulated self-energies into $\delta C(q,0)$ via the \msrjd{}
Dyson equation. At $d=1$ this reproduces the backend-B golden reference; at
$d=2,3$ it gives the higher-dimensional bubble. Again: this is the
\emph{cross-check}, kept to certify that the production integrator produces the
same numbers.

\section{Guards, errors, and known limits}\label{sec:hk-guards}

The subsystem is built to \emph{refuse} cases it cannot handle correctly rather
than return a plausible-looking wrong number. Every refusal is a
\code{SpatialPropagatorError} (\file{heat\_kernel.py:50}) --- the subsystem's one
exception type, raised whenever the Tier-1 closed-form path does not apply ---
or, downstream, a clear \code{NotImplementedError}. Collected in one place, the
guards are:

\begin{description}
  \item[The diagonal (Tier-1) restriction.] The closed-form heat kernel requires
        a \emph{diagonal} inverse propagator (each field an independent
        Allen--Cahn mode). Any off-diagonal coupling makes
        \code{build\_spatial\_propagator} raise
        (\file{heat\_kernel.py:369--377}) and leaves \code{G\_tx\_sym = None};
        \code{compute\_spatial\_correlator\_tree} then raises a clear
        \code{NotImplementedError} (\file{spatial\_correlator.py:239--248}), and
        the bridge re-routes a coupled model to the spectral driver of
        Chapter~\ref{ch:spatial-coupled}.
  \item[Higher-derivative operators are rejected.]
        \code{extract\_mass\_diffusion} rejects any $k$-power $>2$
        (\file{heat\_kernel.py:213--216}) --- a $(\Lap)^{2}$ (a $k^{4}$, e.g.\ a
        conserved Model-B \emph{kinetic} operator) is a v2 feature. (Conserved,
        $q^{2}$-dependent \emph{noise} is separately skipped by the noise
        extractors, left at zero.)
  \item[The $\ii\omega$ normalization is strict.] The time-derivative coefficient
        must be exactly $\ii$ (\file{heat\_kernel.py:201}). A model that hides a
        factor on $\Dt$ is rejected as ``not a standard \msrjd{} kernel'' rather
        than silently rescaled.
  \item[Diffusion sign and zero.] \code{compute\_spatial\_correlator\_tree}
        raises on $B<0$ (anti-diffusive, ill-posed, the heat kernel diverges ---
        check the sign of $-D\,\Lap$ with $D>0$) and on $B=0$ (no Laplacian: a
        time-only field, for which a spatial correlator is undefined)
        (\file{spatial\_correlator.py:261--275}). The underlying reason is that
        the erf form divides by $\sqrt{4\pi B}$ and takes $\sqrt{}$ of the mass.
  \item[Periodic boundaries and drift are $d=1$ only.] \code{image\_sum} raises
        for $d\ne 1$ (\file{heat\_kernel.py:143--145});
        \code{gaussian\_heat\_kernel} raises for $V\ne 0$ with $d\ne 1$
        (\file{heat\_kernel.py:85--88}); and \code{build\_spatial\_propagator}
        raises for periodic boundaries with $d\ne 1$
        (\file{heat\_kernel.py:345--349}).
\end{description}

The v2 backlog these guards delimit is, in one list: off-diagonal (coupled
multi-field) propagators, higher-derivative operators ($k^{4}$+ kinetic terms),
and $d\ge 2$ derivative vertices (the one remaining numerical-transform gap of
\S\ref{sec:hk-loops}).

\section{What this chapter covered}

The spine of this chapter is one identity --- \emph{the inverse transform of a
Gaussian is a Gaussian} --- applied twice. Closing the $\omega$-contour on the
diffusive inverse propagator $A + Bk^{2} + \ii\omega$ gives exponential decay in
time, $\theta(t)\,\ee^{-(A+Bk^{2})t}$; inverse-transforming the surviving
momentum Gaussian gives the real-space heat kernel
$\theta(t)\,(4\pi Bt)^{-d/2}\,\ee^{-x^{2}/4Bt - At}$, the function
\code{gaussian\_heat\_kernel}. From there, drift shifts the Gaussian centre to
$x=vt$ (and is usually zero at the homogeneous saddle); a periodic box is the
super-exponentially convergent image sum \code{image\_sum}; the three coefficients
$(A,B,V)$ are read symbolically off \code{K\_ft} with Sage by
\code{extract\_mass\_diffusion}; the picklable propagator block is assembled by
\code{build\_spatial\_propagator} and its non-picklable callables rebuilt by
\code{make\_g\_tx\_callables}; the tree-level correlator is noise driving two
retarded legs, collapsed to the mpmath-precision erf integral
\code{erf\_time\_integral} inside \code{free\_two\_point}; and the noise strength
itself is read out of the action's $(2,0)$ sector by
\code{extract\_noise\_coefficients}.

At loop order the same Gaussian transform runs batched and vectorized in
\code{full\_integrator.\_heat\_kernel\_x\_general}, gated by \code{\_use\_analytic}
(plain vertices at any $d$, or any vertices at $d=1$), with a numerical-transform
cross-check (\code{radial\_inverse\_ft}, \code{periodic\_inverse\_ft}) kept alive
behind \code{SPATIAL\_FORCE\_NUMERICAL\_FT} and friends, and pinned by the
oracle-only Symanzik modules \code{loop\_parametric.py} and
\code{temporal\_integrate.py}. The next chapter,
Chapter~\ref{ch:spatial-coupled}, takes up the case this one explicitly refuses
--- the coupled, off-diagonal multi-field propagator --- where a single diagonal
heat kernel no longer suffices and the spectral construction takes over.

% ---------------------------------------------------------------------
%  Chapter 16 : Coupled Fields, Dyson Dressing, the Operator IR, and the Bridge
% ---------------------------------------------------------------------
\chapter{Coupled Fields, Dyson Dressing, and the Bridge}\label{ch:spatial-coupled}

The two preceding spatial chapters built the machinery for the \emph{easy}
spatial case: a single physical field whose linear dynamics is a plain heat
kernel, $\Gret(t,k)=\Theta(t)\,\ee^{-(\mu+Dk^{2})t}$, and whose only
complication is a derivative-free interaction vertex. That covers a surprising
amount of physics, but it leaves three genuinely harder situations untouched,
and each one breaks a different assumption of the diagonal path:

\begin{enumerate}
  \item \term{Coupled multi-field theories.} When several fields
        $\phi_{1},\dots,\phi_{N}$ have linear dynamics that \emph{mix} --- an
        off-diagonal reaction (``mass'') matrix $M$ --- the free propagator is a
        \emph{matrix} exponential $\ee^{-Mt}$, not a product of independent
        scalar modes. The single-pole-per-field picture is simply wrong.
  \item \term{Unequal diffusion.} When the diffusion matrix $\mathcal{D}$ is not
        a scalar multiple of the identity, $M$ and $\mathcal{D}$ do not commute,
        so $\ee^{-(M+\mathcal{D}k^{2})t}$ does not factor and has no simple
        $k$-dependence. There is no closed-form propagator to read off.
  \item \term{Derivative vertices.} A vertex like Model~B's conserved
        $\Lap(\phi^{2})$, Burgers' $\partial_{x}(\phi^{2})$, or KPZ's
        $(\partial_{x}\phi)^{2}$ carries spatial derivatives. These cannot be
        folded into a momentum-independent coupling; each one deposits a
        \emph{momentum form factor} that the loop integrator must average.
\end{enumerate}

This chapter is the part of the engine that handles all three. It is, in
practice, four cooperating modules:
\file{spectral\_propagator.py} (the coupled free propagator),
\file{dyson\_dressing.py} (the unequal-diffusion correction series),
\file{spatial\_operator\_ir.py} (the symbolic operator algebra for derivative
vertices), and \file{pipeline\_bridge.py} (the 2157-line dispatcher that ties
everything to the shared diagram pipeline). A fifth module,
\file{loop\_dyson.py}, is an off-the-production-path cross-check oracle we
describe at the very end.

The organising idea --- stated in the bridge's own module docstring --- is
worth fixing in your mind before any of the math: \textbf{route through the
shared pipeline, then certify.} Rather than re-implementing diagrammatics for
the spatial case, the bridge substitutes $\Lap\to-q^{2}$ and runs the
\emph{same} enumeration, classification, and time-domain integration that a
time-only theory uses; it then \emph{certifies} that the diagram answer
$C(q,\tau)$ equals the per-mode structure read off the propagator, before doing
the external $q\to x$ Fourier transform analytically.

\begin{note}
Everything in this chapter sits \emph{behind} the public
\code{compute\_cumulants} call you met in the quickstart. You do not invoke
\code{compute\_coupled\_tree\_correlator} or \code{dressed\_tree\_C} by hand;
\code{compute\_cumulants} dispatches to them based on the shape of your theory
(coupled or not, equal or unequal diffusion, tree or loop, $k=2$ or $k\ge3$).
This chapter explains \emph{what} it dispatches to and \emph{why}.
\end{note}

\section{Where this subsystem sits}

The end-to-end flow, from the text you type to the array that comes out, runs:

\begin{lstlisting}[style=console]
theory file (.theory.py)
  -> SpatialTheoryBuilder / theory_compiler   (operator IR lowering lives here)
  -> FieldTheory.expand()                      (multivariate Taylor -> vertices)
  -> build_propagator / build_spatial_propagator   (the `prop` dict)
  -> compute_cumulants   (the public entry; dispatches by spatial_dim/coupling)
       -> compute_spatial_correlator_via_pipeline    (tree, diagonal)
       -> compute_coupled_tree_correlator            (tree, coupled)
       -> compute_spatial_correlator_generic         (loops, single-field)
       -> compute_coupled_loop_correlator            (loops, coupled)
       -> compute_spatial_kpoint / compute_coupled_kpoint   (k>=3)
  -> C(x, tau) array  ->  notebook / sim comparison
\end{lstlisting}

The five drivers in the middle are all in \file{pipeline\_bridge.py}, and the
chapter walks through them roughly in that order. The bridge \emph{feeds on} the
expanded \code{FieldTheory} object (\code{ft}), the model dict (\code{model}),
the propagator dict (\code{prop}), and the numeric parameters (\code{num\_params},
which crucially includes the mean-field saddle). It \emph{calls into} the shared
diagram pipeline, the spatial \code{full\_integrator} (the one genuine loop
integral, Chapter~\ref{ch:spatial-core}), the coupled-field modules of this
chapter, and the analytic $q\to x$ transforms in \code{spatial\_correlator}. It
\emph{produces} a real-space correlator array \code{C[len(tau\_grid),
len(spatial\_grid)]} plus an \code{info} dict of diagnostics.

\section{Why coupled fields need a different propagator}

Start with the linear theory, because the interaction structure rides on top of
it. For an $N$-component field the linearized inverse propagator --- the
bilinear ``kernel'' in the \msrjd{} action --- is, in the mixed
frequency/momentum representation $(\omega,k)$,
\begin{equation}
  K(\omega,k) \;=\; -\ii\omega\,I \;+\; M \;+\; \mathcal{D}\,|k|^{2}.
  \label{eq:Kwk}
\end{equation}
Two matrices appear, and both can be non-trivial:
\begin{itemize}
  \item $M$ is the \term{reaction matrix} (or mass matrix),
        $M=\mathrm{diag}(\mu_{i})-A^{(0)}$. It is the $k^{0}$
        (momentum-independent) part. Its off-diagonal entries are linear
        couplings between fields; diagonal $M$ means decoupled fields.
  \item $\mathcal{D}$ is the \term{diffusion matrix},
        $\mathcal{D}=\mathrm{diag}(D_{i})+A^{(2)}$. It is the $k^{2}$ part.
        $\mathcal{D}\propto I$ means equal (scalar) diffusion.
\end{itemize}
This is exactly the structure recorded in the \file{spectral\_propagator.py}
module docstring (lines~10--16), which names \eqref{eq:Kwk} and the two
splittings verbatim.

The retarded free propagator is the matrix Green's function
\begin{equation}
  \Gret(t,k)\;=\;\Theta(t)\,\ee^{-(M+\mathcal{D}|k|^{2})t}.
  \label{eq:GRmatrix}
\end{equation}
Here is the trouble in one sentence: \emph{if $M$ and $\mathcal{D}$ do not
commute, the exponent $M+\mathcal{D}|k|^{2}$ cannot be split, so
\eqref{eq:GRmatrix} is not a product of two matrix exponentials and has no
simple $k$-dependence.} The diagonal pipeline's whole strategy --- one
$\omega$-pole per field, $k$ entering only through $\mu+Dk^{2}$ --- collapses.

\subsection{The reference-diffusion split}

The escape is to split off a \emph{scalar} reference diffusion and treat the
rest as a perturbation. Write
\begin{equation}
  \mathcal{D}\;=\;D_{0}\,I\;+\;\widehat{\mathcal{D}},
  \qquad D_{0}\in\mathbb{R},\qquad
  \widehat{\mathcal{D}}=\text{residual}\ (=0\ \text{iff}\ \mathcal{D}\propto I).
  \label{eq:Dsplit}
\end{equation}
The reference kernel $K_{0}=-\ii\omega\,I+M+D_{0}|k|^{2}\,I$ now has scalar
diffusion, which commutes with \emph{everything}. The only remaining
non-commuting object is $M$ itself, and a diagonalizable $M$ can always be
resolved into its spectral projectors:
\begin{equation}
  M\;=\;\sum_{\alpha} m_{\alpha}\,P_{\alpha},
  \qquad \sum_{\alpha}P_{\alpha}=I,
  \qquad P_{\alpha}P_{\beta}=\delta_{\alpha\beta}P_{\alpha}.
  \label{eq:spectral}
\end{equation}
The $m_{\alpha}$ are the eigenvalues of $M$ and $P_{\alpha}$ is the rank-one
projector onto the $\alpha$-th eigenspace, built as the outer product of the
right eigenvector with the corresponding row of the inverse eigenvector matrix.
With this, the \term{reference propagator} is closed form (the paper's
Appendix~B, eq.~B23):
\begin{equation}
  G_{0}(t,k)\;=\;\Theta(t)\sum_{\alpha}P_{\alpha}\,
     \ee^{-(m_{\alpha}+D_{0}|k|^{2})t}
   \;=\;\Theta(t)\,\ee^{-Mt}\,\ee^{-D_{0}|k|^{2}t}.
  \label{eq:G0}
\end{equation}

Two exactness facts make this more than a formal manoeuvre, and both are stated
in the module docstring (lines~33--37) and validated in
\file{tests/test\_spectral\_propagator.py}:
\begin{itemize}
  \item \textbf{If $\widehat{\mathcal{D}}=0$} (scalar diffusion, even with a
        coupled $M$), then $G_{0}$ is the \emph{exact} full propagator. No
        series is needed. This alone unlocks every coupled-reaction,
        equal-diffusion theory.
  \item \textbf{If $M$ and $\mathcal{D}$ are both diagonal}, $G_{0}$ reduces to
        the per-field scalar heat kernel $\ee^{-(\mu_{i}+D_{i}|k|^{2})t}$ that
        the diagonal pipeline (\code{heat\_kernel.py}) already builds. So the
        coupled code degenerates gracefully to the case you already know.
\end{itemize}

\begin{defn}
A matrix $M$ is \term{diagonalizable} if it has a full set of linearly
independent eigenvectors, i.e. $M=V\,\mathrm{diag}(m_{\alpha})\,V^{-1}$ for an
invertible $V$. The spectral resolution \eqref{eq:spectral} exists exactly when
$M$ is diagonalizable. A \term{defective} matrix (repeated eigenvalues with
non-trivial Jordan blocks) has no such resolution and would need the resolvent /
confluent form; that case is deliberately deferred.
\end{defn}

\section{Building the projectors: \texttt{spectral\_propagator.py}}

The numeric core lives in five functions plus a cached dataclass. We take them
in the order the data flows.

\subsection{Splitting the diffusion}

\code{split\_reference\_diffusion(D\_mat, D0=None)}
(\file{spectral\_propagator.py:55}) implements \eqref{eq:Dsplit}:

\begin{lstlisting}
def split_reference_diffusion(D_mat, D0=None):
    D_mat = np.asarray(D_mat, dtype=float)
    n = D_mat.shape[0]
    if D0 is None:
        D0 = float(np.trace(D_mat) / n)
    Dhat = D_mat - D0 * np.eye(n)
    return float(D0), Dhat
\end{lstlisting}

When you do not pass a $D_{0}$, it defaults to the isotropic part
$\mathrm{trace}(\mathcal{D})/N$ --- the mean eigenvalue. That is a sensible
reference because it minimizes a natural norm of $\widehat{\mathcal{D}}$, which
(as \S\ref{sec:dyson} shows) is exactly what governs the convergence of the
Dyson series. The override exists so a caller can retune $D_{0}$ to shrink the
ratio $\lVert\widehat{\mathcal{D}}\rVert/D_{0}$.

\subsection{The spectral projectors and the defectiveness guard}

\code{spectral\_projectors(M)} (\file{spectral\_propagator.py:70}) is the
literal translation of \eqref{eq:spectral}:

\begin{lstlisting}
def spectral_projectors(M):
    M = np.asarray(M, dtype=complex)
    w, V = np.linalg.eig(M)
    cond = np.linalg.cond(V)
    if not np.isfinite(cond) or cond > _COND_CAP:
        raise ValueError(...)   # near-defective: deferred
    Vinv = np.linalg.inv(V)
    proj = [np.outer(V[:, a], Vinv[a, :]) for a in range(len(w))]
    return w, proj
\end{lstlisting}

Two NumPy calls do the algebra. \code{np.linalg.eig(M)} returns the eigenvalues
\code{w} (the $m_{\alpha}$) and the matrix \code{V} whose columns are the right
eigenvectors. The projector $P_{\alpha}=$ \code{np.outer(V[:,a], Vinv[a,:])} is
the outer product of the $\alpha$-th eigenvector with the $\alpha$-th row of
$V^{-1}$ --- the standard construction of the rank-one spectral projector of a
diagonalizable matrix.

\begin{defn}
\code{np.linalg.cond(V)} is the \term{condition number} of \code{V}: the ratio
of its largest to smallest singular value. It measures how close \code{V} is to
singular. A nearly defective $M$ has nearly parallel eigenvectors, so \code{V}
is nearly singular and \code{cond} blows up. The guard
\code{\_COND\_CAP = 1e10} (\file{spectral\_propagator.py:52}) is the line in the
sand: above it, the projector construction is numerically unreliable and the
function \emph{raises} \code{ValueError} rather than return garbage.
\end{defn}

\begin{gotcha}
The guard catches a \emph{near}-defective $M$ before it produces silently wrong
numbers, but it is the only protection: there is no fallback to a confluent or
resolvent form. If your theory has a genuinely defective reaction matrix (a
repeated eigenvalue with a Jordan block), this subsystem will refuse it. In
practice this is rare --- it requires a fine-tuned degeneracy --- but it is a
hard boundary of the implementation, not a soft warning.
\end{gotcha}

\subsection{The cached reference}

\code{SpectralReference} (\file{spectral\_propagator.py:91}) is a
\code{@dataclass} that caches the spectral data so it is computed once per
$(M,\mathcal{D})$ and reused across the whole $q\times\tau$ grid. It carries
\code{M}, \code{D} (the full diffusion), \code{D0}, \code{Dhat}, \code{eigvals}
(the $m_{\alpha}$), and \code{projectors} (the $P_{\alpha}$). Three properties
encode the facts above:

\begin{itemize}
  \item \code{n\_fields} $=$ \code{M.shape[0]}.
  \item \code{is\_scalar\_diffusion} $=$ \code{bool(np.allclose(self.Dhat, 0.0))}.
        \textbf{When this is true, $G_{0}$ is the exact full propagator and no
        Dyson series is needed.} This single boolean is the most consequential
        branch point in the coupled tree path.
  \item \code{G0(ksq, t)} returns the matrix
        $\sum_{\alpha}P_{\alpha}\,\ee^{-(m_{\alpha}+D_{0}\cdot\code{ksq})t}$
        (the caller applies $\Theta(t)$).
\end{itemize}

\code{build\_reference(M, D, D0=None)} (\file{spectral\_propagator.py:116}) is
the single constructor every coupled path uses: it validates that $M$ and $D$
are square and the same shape, splits the diffusion, computes the projectors,
and assembles the dataclass.

\begin{note}
A \code{@dataclass} is a Python class whose boilerplate (the constructor, a
readable \code{repr}, equality) is generated automatically from a few typed
field declarations. You write the field list; Python writes
\code{\_\_init\_\_}. It is the idiomatic way to carry a small bundle of related
values, and you will see it again for the integrator's diagram descriptors.
\end{note}

\section{The tree-level coupled two-point}

The free two-point of a coupled \emph{linear} theory is \emph{not} a sum of
independent Ornstein--Uhlenbeck modes --- that diagonal form is precisely what
the single-field path builds and what fails here. The correct object comes from
two classical results of linear stochastic dynamics.

Let $A(q)=M+\mathcal{D}|q|^{2}$ be the relaxation matrix at momentum $q$, and
let $N$ be the symmetric noise covariance, read from the $(2,0)$ response-field
sector of the action (the statement $\avg{\xi\xi^{\mathsf T}}=N$). Then:

\begin{description}
  \item[Lyapunov / stationary covariance.] The equal-time covariance
        $\Sigma(q)$ solves the \term{continuous Lyapunov equation}
        \begin{equation}
          A(q)\,\Sigma(q)\;+\;\Sigma(q)\,A(q)^{\mathsf T}\;=\;N.
          \label{eq:lyap}
        \end{equation}
  \item[Fluctuation--regression.] The relaxation of an equilibrium fluctuation
        follows the same exponential law as a macroscopic perturbation
        (Onsager's regression hypothesis), so
        \begin{equation}
          C(q,\tau)=\ee^{-A(q)|\tau|}\,\Sigma(q)\ \ (\tau\ge0),
          \qquad
          C(q,-\tau)=\Sigma(q)\,\ee^{-A(q)^{\mathsf T}|\tau|},
          \label{eq:fdt}
        \end{equation}
        i.e. $C(q,\tau)=C(q,|\tau|)^{\mathsf T}$.
\end{description}

For scalar diffusion $A(q)=M+D_{0}|q|^{2}\,I$ shares $M$'s eigenprojectors, so
$\ee^{-A|\tau|}$ \emph{is} $G_{0}(|q|^{2},|\tau|)$ and \eqref{eq:fdt} is exact.
Unequal diffusion needs the Dyson series of the next section.

\subsection{\texttt{coupled\_two\_point} and the Lyapunov solver}

\code{lyapunov\_covariance(A, N)} (\file{spectral\_propagator.py:156}) is a
single call to SciPy:

\begin{lstlisting}
def lyapunov_covariance(A, N):
    A = np.asarray(A, dtype=float)
    N = np.asarray(N, dtype=float)
    return solve_continuous_lyapunov(A, N)
\end{lstlisting}

\code{coupled\_two\_point(ref, N, qsq, tau)}
(\file{spectral\_propagator.py:169}) assembles \eqref{eq:lyap}--\eqref{eq:fdt}:

\begin{lstlisting}
def coupled_two_point(ref, N, qsq, tau):
    if not ref.is_scalar_diffusion:
        raise ValueError(...)          # unequal diffusion needs Dyson
    n = ref.n_fields
    A = np.asarray(ref.M, dtype=float) + ref.D0 * float(qsq) * np.eye(n)
    Sigma = lyapunov_covariance(A, N)
    G = expm(-A * abs(float(tau)))     # e^{-A|tau|}  (= G0 for scalar D)
    return G @ Sigma if tau >= 0 else Sigma @ G.T
\end{lstlisting}

It forms $A=M+D_{0}q^{2}I$, solves for $\Sigma$, exponentiates, and returns
$G\Sigma$ (or its transpose for $\tau<0$). The returned $(N,N)$ matrix has
$C_{ij}$ equal to the cross-correlation of fields $i$ and $j$. The
\code{raise} on the first line is the guard: this routine is for scalar
diffusion only.

\begin{defn}
\code{scipy.linalg.solve\_continuous\_lyapunov(A, N)} solves
$A\,X+X\,A^{\mathsf H}=Q$ for $X$ (here $Q=N$ and $A^{\mathsf H}=A^{\mathsf T}$
since $A$ is real). It uses the \term{Bartels--Stewart algorithm} (a Schur
decomposition followed by back-substitution), which is far more stable than
forming and inverting the $N^{2}\times N^{2}$ Kronecker-product system by hand.
\code{scipy.linalg.expm(Z)} is the \term{matrix exponential} $\ee^{Z}$, computed
by a Pad\'e approximation with scaling-and-squaring --- robust even for
non-normal matrices, which is essential here because $A$ need not be symmetric.
\end{defn}

\section{The Dyson--Duhamel series for unequal diffusion}\label{sec:dyson}

When $\widehat{\mathcal{D}}\ne0$ there is no closed-form propagator, and we
expand $\Gret$ in powers of the insertion $\widehat{\mathcal{D}}|k|^{2}$ (paper
Appendix~B, \S B.24--B.30):
\begin{align}
  \Gret(t,k)&=\sum_{n\ge0}G_{n}(t,k),\notag\\
  G_{n}(t,k)&=(-|k|^{2})^{n}\,\ee^{-D_{0}|k|^{2}t}\,\mathcal{H}_{n}(t).
  \label{eq:B26}
\end{align}
The matrix factor $\mathcal{H}_{n}(t)$ comes from the $n$-fold \term{Duhamel
convolution} --- each $\widehat{\mathcal{D}}$ insertion sandwiched between
reference evolutions:
\begin{equation}
  G_{n}(t)=(-|k|^{2})^{n}\!\!\int\limits_{t\ge s_{1}\ge\cdots\ge s_{n}\ge0}
   \!\!\ee^{-M(t-s_{1})}\widehat{\mathcal{D}}\,
   \ee^{-M(s_{1}-s_{2})}\widehat{\mathcal{D}}\cdots
   \ee^{-Ms_{n}}\,\dd\mathbf{s}.
  \label{eq:duhamel}
\end{equation}
Inserting the spectral resolution \eqref{eq:spectral} for every $\ee^{-M(\cdot)}$
turns this nested time integral into a \term{divided difference}, and
shift-invariance factors out $\ee^{-m_{\alpha_{0}}t}$:
\begin{equation}
  \mathcal{H}_{n}(t)=\!\!\sum_{\alpha_{0}\dots\alpha_{n}}\!\!
   P_{\alpha_{0}}\widehat{\mathcal{D}}P_{\alpha_{1}}\cdots
   \widehat{\mathcal{D}}P_{\alpha_{n}}\,
   \ee^{-m_{\alpha_{0}}t}\,
   \Phi_{n}\!\big(t;\,m_{\alpha_{1}}-m_{\alpha_{0}},\dots,
   m_{\alpha_{n}}-m_{\alpha_{0}}\big).
  \label{eq:B27}
\end{equation}

\subsection{The one new primitive: \texorpdfstring{$\Phi_{n}$}{Phi\_n}}

Everything in \eqref{eq:B27} is projector bookkeeping except the scalar function
$\Phi_{n}$, the nested-simplex time integral
\begin{equation}
  \Phi_{n}(t;\nu_{1},\dots,\nu_{n})=\int_{\sigma_{n}}
    t^{n}\,\ee^{-t\sum_{i}u_{i}\nu_{i}}\,\dd\mathbf{u},
  \qquad \Phi_{0}(t)=1,
  \label{eq:phin}
\end{equation}
over the standard simplex $\sigma_{n}=\{u_{i}\ge0,\ \sum u_{i}\le1\}$. The
classical \term{Hermite--Genocchi formula} with $f(z)=\ee^{-tz}$ (so
$f^{(n)}(z)=(-t)^{n}\ee^{-tz}$) and nodes $\{0,\nu_{1},\dots,\nu_{n}\}$ states
that this simplex integral is, up to sign, the $n$-th divided difference of $f$:
\begin{equation}
  \Phi_{n}(t;\nu)=(-1)^{n}\,f[0,\nu_{1},\dots,\nu_{n}].
  \label{eq:hg}
\end{equation}

\begin{defn}
The \term{divided difference} $f[x_{0},\dots,x_{n}]$ is the leading coefficient
of the degree-$n$ polynomial interpolating $f$ at the nodes $x_{0},\dots,x_{n}$;
for distinct nodes it equals $\sum_{i}f(x_{i})/\prod_{j\ne i}(x_{i}-x_{j})$. The
\term{Hermite--Genocchi formula} writes that same object as a simplex integral
of $f^{(n)}$, which is what links the Duhamel integral \eqref{eq:duhamel} to the
divided difference.
\end{defn}

\subsection{Confluent-safe evaluation: the Opitz theorem}

The naive formula $\sum_{i}f(x_{i})/\prod_{j\ne i}(x_{i}-x_{j})$ is a numerical
trap: it is $0/0$ whenever two nodes coincide or nearly coincide. And the
eigenvalue differences $\nu=m_{\alpha_{i}}-m_{\alpha_{0}}$ \emph{routinely}
repeat --- whenever several $m_{\alpha}$ are equal, or come as the
complex-conjugate pairs that a real $M$ produces. The \term{Opitz theorem}
sidesteps division entirely:

\begin{quote}
For analytic $f$, build the $(n+1)\times(n+1)$ upper-bidiagonal matrix $Z$ with
the nodes $\{0,\nu_{1},\dots,\nu_{n}\}$ on the diagonal and ones on the
superdiagonal. Then $f(Z)[0,n]=f[0,\nu_{1},\dots,\nu_{n}]$ --- the top-right
corner of $f$ applied to $Z$.
\end{quote}

So $\Phi_{n}(t;\nu)=(-1)^{n}\,\code{expm}(-t\,Z)[0,n]$: a single matrix
exponential of a tiny $(n+1)\times(n+1)$ matrix, exact at coincident nodes and
stable at near-coincident ones. The implementation
(\file{spectral\_propagator.py:221}) is:

\begin{lstlisting}
def phi_n(t, nus):
    nus = np.asarray(nus, dtype=complex).ravel()
    n = nus.size
    if n == 0:
        return complex(1.0)
    Z = _opitz_bidiagonal(nus)
    return complex((-1) ** n * expm(-float(t) * Z)[0, n])
\end{lstlisting}

with \code{\_opitz\_bidiagonal} (\file{spectral\_propagator.py:211}) building
$Z$ via \code{np.diag(np.ones(n), k=1)} (ones on the superdiagonal) and then
writing the nodes onto the diagonal.

\begin{example}\label{ex:phi1}
The $n=1$ check from the docstring (\file{spectral\_propagator.py:240}). With
$Z=\begin{psmallmatrix}0&1\\0&\nu\end{psmallmatrix}$,
$\code{expm}(-tZ)[0,1]=-t\,(1-\ee^{-t\nu})/(t\nu)=-(1-\ee^{-t\nu})/\nu$; times
$(-1)$ gives $\Phi_{1}=(1-\ee^{-t\nu})/\nu$, which is exactly
$\int_{0}^{1}t\,\ee^{-tu\nu}\,\dd u$. The formula reduces to the textbook answer,
and --- crucially --- it has a finite limit $\Phi_{1}\to t$ as $\nu\to0$ with no
cancellation, because the matrix exponential never divides by $\nu$.
\end{example}

\begin{gotcha}
\code{phi\_n} always returns a Python \code{complex}, never a float, even when
the inputs look real. This is deliberate: a real $M$ can have a complex
spectrum, so the node differences $\nu$ come in conjugate pairs. The callers
decide when it is safe to take the real part. \code{phi\_n\_batch}
(\file{spectral\_propagator.py:259}) is the vectorized version --- one
\code{expm} per $t$ on the shared tiny $Z$ --- used by the dressing assembly.
\end{gotcha}

\subsection{Assembling the dressing: \texttt{dyson\_dressing.py}}

\code{hcal\_n(ts, M, Dhat, n)} (\file{dyson\_dressing.py:47}) is the direct
transcription of \eqref{eq:B27}:

\begin{lstlisting}
def hcal_n(ts, M, Dhat, n):
    eig, proj = spectral_projectors(M)
    nf = len(eig)
    out = np.zeros((ts.size, nf, nf), dtype=complex)
    for string in itertools.product(range(nf), repeat=n + 1):
        mat = proj[string[0]]
        for a_i in string[1:]:
            mat = mat @ Dhat @ proj[a_i]
        if not np.any(np.abs(mat) > 1e-300):
            continue                                   # skip a zero string
        m0 = eig[string[0]]
        nus = [eig[a_i] - m0 for a_i in string[1:]]
        scal = np.exp(-m0 * ts) * phi_n_batch(ts, nus)
        out += scal[:, None, None] * mat[None, :, :]
    return out
\end{lstlisting}

It loops over every projector string $(\alpha_{0},\dots,\alpha_{n})$ ---
\code{itertools.product(range(nf), repeat=n+1)} enumerates all $N^{n+1}$ of
them --- builds the outer-product chain
$P_{\alpha_{0}}\widehat{\mathcal{D}}P_{\alpha_{1}}\cdots\widehat{\mathcal{D}}P_{\alpha_{n}}$,
skips strings whose matrix is numerically zero, forms the scalar
$\ee^{-m_{\alpha_{0}}t}\Phi_{n}$ via \code{phi\_n\_batch}, and accumulates the
$\text{scalar}\otimes\text{matrix}$ product. The cost is $N^{n+1}$ strings ---
acceptable because the truncation order $n$ stays small ($\lesssim4$).

\begin{defn}
\code{itertools.product(range(nf), repeat=m)} is the Python standard-library
Cartesian-product generator: it yields every length-$m$ tuple of values drawn
from $0,\dots,nf-1$, i.e. exactly the index strings summed in \eqref{eq:B27}.
\code{np.einsum} (used below) is NumPy's Einstein-summation engine: you write the
index pattern of a tensor contraction as a string and it performs the sum in
compiled C.
\end{defn}

\code{dressed\_GR(ts, qsq, M, Dhat, D0, order)} (\file{dyson\_dressing.py:73})
assembles \eqref{eq:B26}: $\sum_{n=0}^{\text{order}}(-q^{2})^{n}\mathcal{H}_{n}$,
multiplied by $\ee^{-D_{0}q^{2}t}$. At \code{order=0} it is exactly the scalar
reference $G_{0}$.

\subsection{The dressed tree two-point and its convergence}

With the dressed retarded propagator in hand, the dressed tree two-point is the
same $\int\,\dd s\,\Gret N\Gret^{\mathsf T}$ structure as
\eqref{eq:fdt}, now evaluated by quadrature because $\Gret$ is a series rather
than a single exponential:
\begin{equation}
  C^{(N)}(q,\tau\ge0)=\int_{0}^{\infty}\dd s\;
    \Gret^{(N)}(\tau+s,q)\,N\,\Gret^{(N)}(s,q)^{\mathsf T}.
  \label{eq:dressedC}
\end{equation}
\code{dressed\_tree\_C(q, tau, M, Dhat, D0, N, order, n\_s=48,
s\_cap\_scale=32.0)} (\file{dyson\_dressing.py:87}) does this:

\begin{lstlisting}
    cap = s_cap_scale / (mu_min + float(D0) * float(q) ** 2)
    xv, wv = np.polynomial.legendre.leggauss(n_s)
    vv = 0.5 * (xv + 1.0)
    s_nodes = cap * vv * vv
    s_w = wv * cap * vv                                  # int ds = sum w*cap*v
    ...
    C = np.einsum('s,sij,jk,slk->il', s_w, G_late, N, G_early)
\end{lstlisting}

Two ideas are worth unpacking. First, it validates that
$\min\mathrm{Re}\,\mathrm{eig}(M)>0$ (a stable theory) and \emph{raises}
otherwise. Second, the $s$-integral uses a \term{$\sigma$-concentrated
substitution} $s=\code{cap}\cdot v^{2}$: \term{Gauss--Legendre quadrature} (fixed
nodes \code{xv} and weights \code{wv} on $[-1,1]$, from \code{leggauss}) is mapped
through $s=\code{cap}\,v^{2}$ so the nodes pack near $s=0$, where the integrand is
largest, with \code{cap} set by the slowest eigenvalue. This is the same trick
the chamber integrator uses, and it makes a modest \code{n\_s=48} accurate.

\begin{defn}
\term{Gauss--Legendre quadrature} approximates $\int_{-1}^{1}g\,\dd v$ by a
weighted sum $\sum_{i}w_{i}\,g(v_{i})$ at fixed nodes; an $n$-point rule is exact
for polynomials of degree $\le 2n-1$. \code{np.polynomial.legendre.leggauss(n)}
returns the nodes and weights. The substitution $s=\code{cap}\,v^{2}$ turns the
half-line $\int_{0}^{\infty}$ into $\int_{0}^{1}$ while concentrating sample
points near the origin.
\end{defn}

The convergence of the truncation is governed by the \term{residual ratio}
\begin{equation}
  \rho=\frac{\lVert\widehat{\mathcal{D}}\rVert}{D_{0}}
       =\frac{\max|\mathrm{eig}(\widehat{\mathcal{D}})|}{D_{0}}.
  \label{eq:rho}
\end{equation}
The series is geometric in the order at large $q$ and \emph{exact at $q=0$}
(where the insertion $\widehat{\mathcal{D}}|k|^{2}$ vanishes), so convergence
requires $\rho<1$. The bridge checks this (\S\ref{sec:dispatch}) and you can
shrink $\rho$ by retuning $D_{0}$.

\section{Loop-level coupled fields: spectral assignments}

So far this has all been tree-level (the linear theory). At loop order the
interaction vertices return, and a beautiful simplification appears for scalar
diffusion.

With $\widehat{\mathcal{D}}=0$, \emph{every} edge's momentum factor is the same
heat kernel $\ee^{-D_{0}k^{2}w}$. The Symanzik (parametric-momentum) machinery
of Chapter~\ref{ch:spatial-core} --- which only ever sees scalar diffusion --- is
therefore completely untouched. The coupling lives entirely in the time/matrix
factor: each retarded segment carries $\sum_{\alpha}P_{\alpha}\,\ee^{-m_{\alpha}w}$.
Expanding every segment (a retarded ``R'' edge is one segment; a correlation
``C'' edge is two glued half-segments) turns one coupled diagram into a
\emph{weighted sum of scalar diagrams}:
\begin{equation}
  \text{value}(\Gamma)=\code{pv}\cdot
   \sum_{\{\alpha_{r}\}}\prod_{r}[P_{\alpha_{r}}]_{p_{r}r_{r}}\cdot
   I_{\text{spec}}\big(\{m_{\alpha_{r}}\}\big).
  \label{eq:specassign}
\end{equation}
Here $(r_{r},p_{r})$ are the segment's (response, physical) matrix indices,
threaded through the edge descriptor's \code{fpairs}; \code{pv} is the
enumeration's $\Sym(\Gamma)\cdot\text{prefactor}$; and $I_{\text{spec}}$ is the
shared Symanzik/chamber quadrature evaluated for one mass assignment, with all
assignments batched against a single integral pass.

\code{compute\_coupled\_loop\_correlator} (\file{pipeline\_bridge.py:1090})
implements \eqref{eq:specassign}. Its docstring (lines~1104--1110) states the
formula and one subtle normalization point that is worth flagging:

\begin{gotcha}
The single-field generic path multiplies by a universal $2^{-n_{C}}$ factor (the
one-segment ``C'' convention). The coupled spectral-assignment paths do
\emph{not}: the two-glued-segment ``C'' representation natively produces the
Lyapunov denominator $1/(m_{\alpha}+m_{\beta}+2D_{0}k^{2})$, so the conversion
does not apply (\file{pipeline\_bridge.py:1108-1110}). Mixing the two conventions
is a classic source of a factor-of-$2^{n_{C}}$ error.
\end{gotcha}

Mechanically, the driver maps each enumerated loop diagram to the integrator's
``C-stack'' descriptor (\code{diagram\_to\_cstack}), expands it into rows
(\code{spectral\_rows}), reads the $(\text{resp},\text{phys})$ indices off each
edge's \code{fpairs}, and builds the full assignment grid with
\code{itertools.product(range(nf), repeat=n\_rows)}. Per-assignment weights are
the product $\prod_{r}[P_{\alpha_{r}}]_{p_{r}r_{r}}$, the live ones are kept, and
the mass table $\{m_{\alpha_{r}}\}$ feeds the one shared quadrature pass.

\begin{gotcha}
\textbf{Cross-correlator $\tau$-asymmetry (the $i\ne j$ fix).} For an
\emph{auto}-correlator ($i=j$) the two mirror-image records dedupe to one record
plus the $\Gamma(\tau)+\Gamma(-\tau)$ retarded completion. For a
\emph{cross}-correlator ($i\ne j$) \emph{both} mirror records survive and each is
evaluated once at the times matching its own leaf field order --- applying the
$\pm\tau$ completion there would double-count and symmetrize away the genuine
$\tau$-asymmetry of $C_{ij}$ (\file{pipeline\_bridge.py:1368-1400}). The
coupled-loop and coupled-$k$-point drivers handle this with the orbit-stabilizer
mapping sum (\S\ref{sec:kpoint}); the comment notes the two views coincide at
$k=2$.
\end{gotcha}

When the diffusion is unequal, the loop path layers the Dyson $\mathcal{H}_{n}$
insertions onto individual segments via a generalized partial-fraction expansion
of the $n$-fold Duhamel convolution (the \code{\_pf\_labels} / \code{\_hn\_labels}
machinery, with a log-derivative Bell expansion for repeated masses), and refuses
periodic boundaries or derivative vertices in this v1.

\section{The operator IR for derivative vertices}

We now switch from the matrix structure of coupled fields to the third hard
case: vertices that carry spatial derivatives. The home for these is
\file{spatial\_operator\_ir.py}, a small \term{intermediate representation} (IR)
that hosts the operators $\Lap$ ($\nabla^{2}$), $\partial_{t}$, and
$\partial_{x_{i}}$ as symbolic nodes and gives them an explicit algebra.

\subsection{Why an IR, and why inert nodes}

The module docstring (lines~10--20) argues the design from two failed
alternatives. A bare multiplicative symbol --- writing
\code{D*Laplacian*phi} --- loses \emph{which field the derivative acts on}
(\code{phi*Laplacian*phi} is ambiguous) and carries no operator algebra, so
linearity and the homogeneous-saddle annihilation would have to be hand-patched
downstream. Sage's own \code{laplacian()} (the Laplace--Beltrami operator from
SageManifolds) wants a manifold, a metric, and \code{DiffScalarField} objects,
and does not compose with the \msrjd{} field-doubling and multivariate-Taylor
machinery, which works on commutative polynomial-ring generators.

The solution is to host the operators as \term{inert Sage function
applications}:

\begin{lstlisting}
from sage.all import SR, I, function
_LAP = function('Lap')              # nabla^2 (.)
_DT  = function('Dt')               # d_t (.)
_DX  = function('Dx')               # d_{x_i} (., i)
\end{lstlisting}

(\file{spatial\_operator\_ir.py:59,68-70}).

\begin{defn}
Sage's \code{function('Lap')} creates an \term{inert symbolic function}: a
callable symbol with no defined evaluation rule. So \code{Lap(phi)} is just a
tree node \code{Lap(phi)} that survives substitution, printing, and tree-walking
unchanged --- Sage never tries to ``compute'' it. The entire IR design rests on
this: the operators are inert nodes, and \emph{all} semantics live in explicit
passes. Note also that Sage's \term{Symbolic Ring} \code{SR} caches symbols by
name, so \code{SR.var('Laplacian')} anywhere returns the same object --- a
property the bridge exploits when it substitutes $\Lap\to-q^{2}$.
\end{defn}

A companion bare symbol \code{GRADX\_SYM = SR.var('GradX', ...)}
(\file{spatial\_operator\_ir.py:78}) is what a \emph{bilinear} $\partial_{x}$
lowers to (the first-derivative analogue of \code{Laplacian}); the heat-kernel
propagator substitutes \code{GradX}$\to \ii k$ to read off a drift coefficient
$V$.

\subsection{The passes}

The IR's semantics are a sequence of tree-walking passes, each one a single
clearly-scoped transformation.

\paragraph{Linearity (always valid).}
\code{apply\_linearity(expr, fields, coords=('x','y','z'))}
(\file{spatial\_operator\_ir.py:147}) pushes every operator node through sums and
constant coefficients: $\Lap(a\,\delta\phi+b\,\delta\psi)\to a\,\Lap(\delta\phi)+b\,\Lap(\delta\psi)$
for field- and coordinate-independent $a,b$. The inner helper \code{\_lin\_node}
first \code{.expand()}s the argument so binomial powers become sums (so
$\Lap((\bar\phi+\delta\phi)^{3})$ distributes term by term, the Cahn--Hilliard
$\nabla^{2}\phi^{3}$ case), distributes over $+$, pulls out the constant factor,
and rebuilds the operator on the fluctuation part. An argument with no
fluctuation left --- $\Lap(\bar\phi)$, $\Lap(\bar\phi^{2})$, or a pure coordinate
coefficient --- stays \emph{atomic}.

\begin{gotcha}
A factor that depends on a \emph{coordinate} is deliberately \emph{not} pulled
out: $f(x)\,\nabla^{2}\phi$ is left atomic, because its Leibniz/product rule and
the resulting $\hat f(p)$ momentum injection are a separate, deferred layer
(\file{spatial\_operator\_ir.py:51-53}). \code{apply\_linearity} is the
\emph{constant}-coefficient rule only.
\end{gotcha}

\paragraph{Saddle expansion (mean kept).}
\code{expand\_about\_saddle(expr, replacements, ...)}
(\file{spatial\_operator\_ir.py:220}) substitutes
$\phi\to\bar\phi+\delta\phi$ per \code{replacements = \{field: (mean, fluct)\}}
and re-applies linearity, so $\Lap(\bar\phi+\delta\phi)\to\Lap(\bar\phi)+\Lap(\delta\phi)$.
The mean term is \emph{retained}.

\paragraph{Contingent annihilation.}
\code{kill\_means(expr, mean\_syms, ops=('Lap','Dt','Dx'))}
(\file{spatial\_operator\_ir.py:245}) sends $\mathrm{Op}(\text{arg})\to0$
whenever \code{arg} depends only on the mean symbols. This is the
\emph{contingent} fact that the saddle is spatially homogeneous (for
$\Lap$/$\partial_{x}$) and/or stationary (for $\partial_{t}$).

\begin{gotcha}
\code{apply\_linearity} (always-valid algebra) and \code{kill\_means} (contingent
on a homogeneous saddle) are \emph{two separate passes} on purpose. For an
\emph{inhomogeneous} saddle --- a front, a pattern --- you omit \code{Lap} from
\code{ops} so that $\Lap(\bar\phi)$ \emph{survives}; it cancels the rest of the
stationarity condition rather than being silently dropped
(\file{spatial\_operator\_ir.py:36-41}). Folding annihilation into linearity
would quietly destroy this.
\end{gotcha}

\paragraph{Lowering to ring generators.}
\code{to\_derived\_generators(expr, fluct\_syms, prefix='Dg')}
(\file{spatial\_operator\_ir.py:275}) is the trick that lets the existing
multivariate-Taylor \code{expand} machinery work unchanged. It replaces each
atomic $\mathrm{Op}(\delta\phi)$ --- e.g. $\Lap(\delta\phi)$,
$\Lap(\Lap(\delta\phi))$, $\partial_{x}(\delta\phi)$ --- with a fresh ring
generator symbol, so Taylor expansion treats $\nabla^{2}\delta\phi$ as if it
were just another field (the ``$u=\delta\phi,\ v=\nabla^{2}\delta\phi$'' trick).
It returns \code{(expr2, genmap)} where \code{genmap[gen] = (base, op\_chain)}
and \code{op\_chain} records the operators applied. It works bottom-up: an
operator wrapping an already-introduced generator extends that generator's chain,
so $\Lap(\Lap(\delta\phi))$ resolves to a single $\nabla^{4}$ generator with
chain \code{(('Lap',),('Lap',))}.

\begin{defn}
A \term{derived generator} is a fresh polynomial-ring symbol standing for a
derivative of a field. Treating $\nabla^{2}\delta\phi$ as an independent
generator is what lets the framework's ordinary commutative-ring Taylor
expansion handle a derivative vertex with \emph{zero} special-casing --- the
derivative information is parked in the \code{genmap} and reapplied later by the
Fourier rules.
\end{defn}

\paragraph{Composing the passes.}
\code{prepare\_action(S, fields, replacements=None, homogeneous=True, ...)}
(\file{spatial\_operator\_ir.py:346}) ties them together for two authoring
conventions: if \code{replacements=None} the action is already in fluctuation
fields, so it only applies linearity and lowers; if \code{replacements} is given
it expands about the saddle, kills means (for a homogeneous saddle), then lowers.

\paragraph{Classification.}
\code{classify\_generators(expr, genmap, fluct\_syms)}
(\file{spatial\_operator\_ir.py:388}) splits the generators into \emph{bilinear}
(appearing only in field-degree-$\le2$ terms, so they fold into the propagator
kernel $K(\omega,k)$) and \emph{derivative-vertex} (appearing in some
field-degree-$\ge3$ term, so they carry per-leg form factors). A generator's
degree contribution is the field-degree of its base. The worked distinction in
the docstring: KPZ's $\tilde\phi\,(\partial_{x}\delta\phi)^{2}$ is degree
$1+1+1=3$, a vertex; reaction-diffusion's $\tilde\phi\,\nabla^{2}\delta\phi$ is
degree $2$, bilinear.

\subsection{The Fourier rules and form factors}

\code{form\_factor(chain, k, omega=None)} (\file{spatial\_operator\_ir.py:420})
is the Fourier image of an operator chain on a leg of momentum $k$, frequency
$\omega$:
\begin{equation}
  \Lap\ \to\ -|k|^{2},\qquad
  \partial_{t}\ \to\ -\ii\omega,\qquad
  \partial_{x_{i}}\ \to\ \ii k_{i},
  \label{eq:fourierrules}
\end{equation}
composed multiplicatively along the chain (so $\nabla^{4}=\Lap\circ\Lap\to|k|^{4}$).
The code is a direct transcription:

\begin{lstlisting}
    for entry in chain:
        op = entry[0]
        if op == 'Lap':
            out *= -k2
        elif op == 'Dt':
            ...
            out *= -I * SR(omega)
        elif op == 'Dx':
            out *= I * kv[int(entry[1])]
        elif op == '__mean__':
            out *= SR(0)                 # an un-annihilated mean derivative
\end{lstlisting}

A derivative \emph{bilinear} term (degree $\le2$) folds into the propagator
kernel via these rules. A derivative \emph{vertex} (degree $\ge3$) instead
deposits a per-leg \term{form factor} $\mathcal{R}(q,\ell)$ on the diagram, a
product over the interaction vertices:
\begin{equation}
  \mathcal{R}(q,\ell)=\prod_{\text{vertices }v}\mathfrak{F}(v),
  \qquad
  \mathfrak{F}(v)=\sum_{t:\,n_{\text{phys}}(t)=\deg(v)}w_{t}\,\mathfrak{f}_{t}(v),
  \label{eq:Rcal}
\end{equation}
where $w_{t}=c_{t}/\sum c$ is the coupling weight and $\mathfrak{f}_{t}$ is the
type's Fourier kernel, evaluated either at the response-leg momentum
(\term{composite} mode: $\nabla^{2}(\phi^{2})$, $\partial_{x}(\phi^{2})$) or as a
product over physical-leg momenta (\term{perleg} mode: $(\partial_{x}\phi)^{2}$).

\begin{defn}
\term{Composite} vs \term{perleg} is the distinction between an operator acting
on the whole $\phi^{n}$ composite at a vertex (one factor, the response-leg
momentum) versus acting on each physical leg separately (a product of per-leg
factors). $\nabla^{2}(\phi^{2})$ is composite; $(\partial_{x}\phi)^{2}$ is
perleg. The vertex-terms table records each type's \code{mode}, and the diagram
form-factor assembler dispatches on it.
\end{defn}

The loop integrator must then average this polynomial form factor over the
loop-momentum Gaussian. There are two routes: a numerical Gauss--Hermite
quadrature, or --- the principled route --- a closed-form joint $(\ell,q)$
Gaussian moment, described next.

\section{The analytic spatial inverse Fourier transform}

A derivative vertex makes the loop-momentum integrand a Gaussian times a
polynomial form factor, and the polynomial part can be integrated \emph{exactly}
by Wick's theorem rather than sampled numerically. This retires the
$q$-grid entirely for $d=1$ derivative-vertex theories.

\code{\_build\_wick\_moment(Rcal, ls, qs)} (\file{pipeline\_bridge.py:815}) is
the construction (it is the $k=2$, $d=1$ case; \code{\_build\_wick\_moment\_multi}
at line~1893 generalizes to $k\ge3$). It splits the loop momentum
$\ell=a\cdot q+\xi$, turns the external $q$ into a complex Gaussian via the
Fourier-transform source, and writes the per-chamber inverse FT as
$\delta C(x)=A\cdot K(\mathcal{B},x)\cdot M_{F}$ with
\begin{equation}
  M_{F}=\mathbb{E}_{Q\sim\mathcal{N}(c,s)}\,
        \mathbb{E}_{\xi\sim\mathcal{N}(0,\Sigma)}
        \big[\mathcal{R}(a\,Q+\xi,\,Q)\big],
  \label{eq:MF}
\end{equation}
both expectations being closed-form Gaussian moments --- a non-central moment
$\mathbb{E}[Q^{m}]$ for the external part and an \term{Isserlis} pairing sum for
the loop part.

\begin{defn}
The \term{Isserlis (Wick) theorem} states that the expectation of a product of
jointly-Gaussian zero-mean variables equals the sum, over all perfect matchings
of the variables, of the products of pairwise covariances; an odd number of
factors gives zero. \code{\_isserlis(idx, Sget)} (\file{pipeline\_bridge.py:802})
is the recursive symbolic implementation: a sum over all perfect matchings of the
index multiset of $\prod_{\text{pairs}}\Sigma$. This is what makes the
loop-momentum average a finite \emph{exact} sum rather than a quadrature.
\end{defn}

The construction is built once per diagram and then \term{lambdified} into NumPy
callables. The key performance factorization, recorded in the docstring of
\code{moment\_x}: the polynomial-in-$X$ form factor has coefficients $g_{k}$ that
are \emph{$x$-independent}, so the expensive per-sample $g_{k}$ are evaluated once
and contracted with cheap powers of $X$ --- roughly an $n_{x}$ speedup over
recomputing the moment at every output point. A second variant,
\code{moment\_bessel}, grades the form factor by the radial scaling power for the
Bessel-K backend.

\begin{defn}
\code{sympy.lambdify(args, expr, 'numpy', cse=True)} \emph{compiles} a SymPy
symbolic expression into a fast Python function whose body is NumPy operations,
so a whole array of samples can be evaluated at once. \code{cse=True} runs
\term{common-subexpression elimination} first --- factoring out shared subterms
like $1/\mathcal{B}^{k}$ --- which is a real speedup for the Wick-moment
polynomials. This is the single most-used bridge from symbolic algebra to
numerics in the whole subsystem (\file{pipeline\_bridge.py:911,937,997,1043,1979}).
\end{defn}

\begin{note}
SymPy and Sage are \emph{both} present in this subsystem, deliberately. Sage's
Symbolic Ring \code{SR} is the canonical physics-facing algebra --- the action,
the inverse propagator $K_{\mathrm{ft}}$, the parameter substitutions. SymPy
(imported lazily as \code{import sympy as \_sp} inside functions) handles the
momentum-space \emph{polynomial} bookkeeping of diagrams --- routing momenta,
building form factors, taking Wick moments --- because that work is polynomial
algebra that does not need Sage's heavier algebraic-number machinery. They never
mix in a single expression.
\end{note}

\section{The bridge architecture and the dispatch tree}\label{sec:dispatch}

We can now assemble the pieces into the dispatcher,
\file{pipeline\_bridge.py}. Its defining trick is the \term{Laplacian
substitution}.

\subsection{The Laplacian substitution and certification}

\code{pipeline\_C\_q\_tau} (\file{pipeline\_bridge.py:259}) runs the shared,
time-only diagram pipeline at $\Lap=-q^{2}$:

\begin{lstlisting}
    Lap = SR.var('Laplacian')
    nps = dict(base_np_sr)
    nps[Lap] = -(qval ** 2)
    compute_poles_and_residues(prop, nps, verbose=False)
    ...
    res = compute_correction_td(typed_diagrams=[r[0] for r in records],
                                prefactors=[r[1] for r in records],
                                k=k, propagator_data=pdata, ...)
\end{lstlisting}

The line \code{nps[Lap] = -(qval**2)} injects $\Lap\to-q^{2}$ so the time-only
pipeline sees a rational propagator at effective mass $m(q)=\mu+Dq^{2}$, and
\code{compute\_poles\_and\_residues} re-solves the pole/residue cache per $q$,
exactly as \code{compute.py} does for a temporal model. The diagram answer
$C(q,\tau)$ is then handed to \code{certify\_modes}
(\file{pipeline\_bridge.py:286}), which computes the maximum relative error
between the pipeline's $C(q,\tau)$ and the per-mode reference
$\sum_{\alpha}\kappa_{\alpha}/(\mu_{\alpha}+D_{\alpha}q^{2})\,
\ee^{-(\mu_{\alpha}+D_{\alpha}q^{2})|\tau|}$ over a small sample grid.

\begin{note}
This is the operational meaning of ``route through the pipeline, then
certify.'' The Laplacian substitution lets the engine reuse the entire temporal
diagram machinery --- enumeration, classification, time-domain integration ---
with momentum carried as a parameter. Certification then confirms that the
diagrams the pipeline produced \emph{are} the heat-kernel modes the analytic
$q\to x$ transform is about to use. For a single-field reaction-diffusion theory
the relative error is around $10^{-15}$ at tree level.
\end{note}

\begin{gotcha}
Certification is a \emph{hard gate}, not a warning.
\code{compute\_spatial\_correlator\_via\_pipeline} \emph{raises}
\code{SpatialPropagatorError} if \code{certify\_max\_rel} exceeds
\code{certify\_tol} (default $10^{-8}$). The message says the $(\mu,D,\kappa)$
extraction or the diagram routing is wrong for the theory. A failed
certification is the engine refusing to return a number it cannot vouch for.
\end{gotcha}

\subsection{The five drivers and the branch points}

\code{compute\_cumulants} chooses among the bridge's drivers based on a short
sequence of predicates. The decisive ones are:

\begin{itemize}
  \item \code{prop['G\_tx\_sym'] is None} $\Rightarrow$ off-diagonal coupling
        $\Rightarrow$ the \emph{coupled} path. (The diagonal heat-kernel block is
        absent; the propagator's \code{K\_ft} is still present and the coupled
        path reads it.)
  \item \code{\_field\_index(leg0) != \_field\_index(leg1)} $\Rightarrow$ a
        \emph{cross}-correlator $\Rightarrow$ the coupled path (the diagonal path
        would wrongly reconstruct $\avg{\phi_{i}\phi_{i}}$).
  \item \code{ref.is\_scalar\_diffusion} (i.e. $\widehat{\mathcal{D}}=0$)
        $\Rightarrow$ the exact \code{coupled\_two\_point}; otherwise consume the
        Dyson policy and use \code{dressed\_tree\_C}.
  \item \code{max\_ell >= 1} $\Rightarrow$ a loop driver; if the tree was coupled,
        the coupled loop driver.
  \item \code{len(external\_fields) >= 3} $\Rightarrow$ a $k$-point driver.
\end{itemize}

\begin{figure}[ht]
\centering
\begin{tikzpicture}[
    >={Stealth[round]}, node distance=6mm,
    box/.style={draw=noteframe, fill=notebg, sharp corners, rounded corners=1pt,
                font=\footnotesize, align=center, inner sep=4pt, text width=33mm},
    leaf/.style={draw=defframe, fill=defbg, sharp corners, font=\footnotesize\ttfamily,
                 align=center, inner sep=4pt, text width=46mm},
    q/.style={font=\scriptsize\itshape, fill=white, inner sep=1pt}]
  \node[box] (start) {\code{compute\_cumulants}\\ (spatial)};
  \node[box, below=10mm of start] (coupled) {diagonal heat-kernel\\ block present?};
  \node[leaf, right=24mm of coupled] (treepipe)
        {compute\_spatial\_\\correlator\_via\_pipeline};
  \node[box, below=12mm of coupled] (scalar) {scalar diffusion\\ ($\widehat{\mathcal D}=0$)?};
  \node[leaf, right=24mm of scalar] (dyson) {dressed\_tree\_C\\ (Dyson series)};
  \node[leaf, below=10mm of scalar] (twopt) {coupled\_two\_point\\ (Lyapunov / FDT)};
  % edges
  \draw[->] (start) -- (coupled);
  \draw[->] (coupled) -- node[q,above]{yes (+ same field)} (treepipe);
  \draw[->] (coupled) -- node[q,right]{no / cross} (scalar);
  \draw[->] (scalar) -- node[q,above]{no} (dyson);
  \draw[->] (scalar) -- node[q,right]{yes} (twopt);
\end{tikzpicture}
\caption{The tree-level dispatch. A diagonal heat-kernel block and a same-field
correlator take the certified single-field path; off-diagonal coupling or a
cross-correlator routes to the coupled path, which then branches on scalar
vs.\ unequal diffusion. (\code{max\_ell}$\ge1$ and $k\ge3$ select loop and
$k$-point drivers on top of this.)}
\label{fig:dispatch}
\end{figure}

The coupled-tree driver \code{compute\_coupled\_tree\_correlator}
(\file{pipeline\_bridge.py:449}) reads \code{prop['K\_ft']}, extracts
$(M,\mathcal{D},V)$ via \code{reaction\_diffusion\_matrices}, numericizes them,
and \emph{requires $V=0$} (a genuine drift is not supported in the
scalar-diffusion v1 --- it raises \code{NotImplementedError} otherwise). It then
builds the \code{SpectralReference}, and:

\begin{itemize}
  \item if scalar diffusion, defines $C_{ij}(q,\tau)$ as
        \code{coupled\_two\_point(ref, N, q**2, tau)[i,j]};
  \item if unequal diffusion, consumes the Dyson policy
        (\code{model['spatial']['dyson']} or the \code{SPATIAL\_DYSON\_ORDER} env
        override), \emph{requires} a fixed truncation order, and checks the
        convergence ratio $\rho$ of \eqref{eq:rho}.
\end{itemize}

\begin{gotcha}
\textbf{The Dyson divergence guard.} \code{compute\_coupled\_tree\_correlator}
\emph{raises} \code{SpatialPropagatorError} if $\rho=\lVert\widehat{\mathcal{D}}\rVert/D_{0}\ge1$
(``the large-$|k|$ insertion outgrows the reference'') and warns if $\rho>0.6$
(slow convergence). The series is exact at $q=0$ regardless --- the insertion
vanishes there --- but for the whole grid to converge you need $\rho<1$. Retune
$D_{0}$ to shrink it (\file{pipeline\_bridge.py:519-523}).
\end{gotcha}

For $d=1$ infinite boundary with scalar diffusion the external $q\to x$ transform
is done by the \term{analytic spectral inverse FT}: each $(\alpha,\beta)$ term
$(P_{\alpha}NP_{\beta}^{\mathsf T})/(m_{\alpha}+m_{\beta}+2D_{0}q^{2})$ is a
single-mode correlator at denominator mass $\mu_{d}=(m_{\alpha}+m_{\beta})/2$,
transformed exactly via \code{free\_two\_point}$(\mu_{d},D_{0},\tfrac12;x,\tau)$,
with $C_{ij}(-\tau)=C_{ji}(\tau)$.

\section{The single-field generic loop path}

The single-field loop driver \code{compute\_spatial\_correlator\_generic}
(\file{pipeline\_bridge.py:1434}) is the one genuine momentum-integral path, and
it is where the form factors of the operator IR and the analytic IFT of the
previous sections actually get used. Its outline:

\begin{enumerate}
  \item Compute the tree $C_{0}$ via the certified
        \code{compute\_spatial\_correlator\_via\_pipeline}; if the tree turned
        out coupled, route to \code{compute\_coupled\_loop\_correlator}.
  \item Build the numeric derivative-vertex form-factor table from
        \code{ns.\_operator\_ir\_vertex\_terms} (substitute couplings into the
        symbolic weights $c_{t}/\sum c$).
  \item Enumerate records, map every diagram to a (descriptor, prefactor,
        form-factor callable, $\ell$) tuple, and filter to the live ones.
  \item Choose an integrator backend, check the memory guard, and do either the
        analytic or the numerical $q\to x$ transform.
\end{enumerate}

\begin{gotcha}
\textbf{Drift is refused on this path too.} For any field with a non-zero saddle
drift $V\ne0$ (a genuine advection $v\cdot\partial_{x}\phi$), the driver raises:
the drifting propagator is validated only at the heat-kernel oracle level, not
in the momentum integrator (which uses $m_{k}=\mu+Dk^{2}$). Gradient theories
like Burgers and KPZ, where $V\to0$ at $\phi^{*}=0$, run fine
(\file{pipeline\_bridge.py:1540-1554}).
\end{gotcha}

\subsection{Backends, the memory guard, and the analytic/numerical gate}

The \code{SPATIAL\_INTEGRATOR} environment variable selects the backend: the
default \code{grid} (deterministic product quadrature), \code{mc}
(importance-sampled Monte-Carlo, bounded memory but biased for derivative
vertices), or \code{bessel} (radial Bessel-K times angular Monte-Carlo, which
regularizes the $\det M\to0$ singularity and handles derivative vertices at two
loops).

\begin{gotcha}
\textbf{The memory guard / curse of dimensionality.} A chamber's quadrature is
$P=n_{t}^{\,n_{V}}\cdot n_{s}^{\,n_{C}}$ samples ($n_{V}$ internal vertices,
$n_{C}$ correlation edges). A KPZ two-loop diagram has $n_{V}=4$, $n_{C}=3$,
giving $P\approx1.8\times10^{8}$ at $(n_{t}{=}16,n_{s}{=}14)$, and the
$(P,n_{x})$ heat-kernel array alone is tens of gigabytes --- an out-of-memory
crash. The guard (\file{pipeline\_bridge.py:1633-1665}) refuses \emph{up front}
with the exact numbers and four mitigation knobs: lower \code{max\_ell}, use the
\code{mc} backend, coarsen \code{SPATIAL\_GRID\_NT}/\code{SPATIAL\_GRID\_NS}, or
raise \code{SPATIAL\_MEM\_BUDGET\_GB} (default 6). It is bypassed for the
bounded-memory \code{mc}/\code{bessel} backends.
\end{gotcha}

The final branch is the \term{analytic-vs-numerical FT gate}
(\file{pipeline\_bridge.py:1677}):
\begin{lstlisting}
    _all_plain = all(rec[2] is None for rec in live_g)
    _force_num = _os.environ.get('SPATIAL_FORCE_NUMERICAL_FT', '') == '1'
    _use_analytic = (_all_plain or d == 1) and not _force_num
\end{lstlisting}
The analytic path (Chapter~\ref{ch:spatial-heatkernel} for plain vertices; the
Wick-moment construction above for $d=1$ derivative vertices) computes
$\delta C(x,\tau)$ directly --- no $q$-grid, no ringing. The numerical path
computes $\delta C(q,\tau)$ on a grid and then Fourier-transforms; it is reachable
as a cross-check via \code{SPATIAL\_FORCE\_NUMERICAL\_FT=1} but is exact only in
the $q_{\text{cut}}\to\infty$, $n_{q}\to\infty$ limit.

\begin{gotcha}
\textbf{The \texttt{Rcal == 0} form factor must still evaluate to zero on every
path.} A diagram that vanishes by a conservation law (e.g. the
$\nabla^{2}(\phi^{2})$ tadpole with a forced-zero leg momentum) returns
\code{\_zero\_ff\_with\_moments}, which attaches zero-returning \code{moment\_x},
\code{moment\_x\_multi}, and \code{moment\_bessel} callables --- otherwise the
analytic IFT would hit a missing-moment error instead of contributing zero
(\file{pipeline\_bridge.py:972-990}).
\end{gotcha}

\subsection{The bubble-versus-tadpole test}

One small but load-bearing classifier deserves mention.
\code{\_diagram\_is\_bubble(td)} (\file{pipeline\_bridge.py:664}) decides whether
a one-loop diagram is a momentum-\emph{dependent} bubble:

\begin{lstlisting}
    for v in route_momenta(td).edge_k2().values():
        e = _sp.expand(v)
        syms = sorted(e.free_symbols, key=str)
        for ii in range(len(syms)):
            for jj in range(ii + 1, len(syms)):
                if _sp.expand(e.diff(syms[ii]).diff(syms[jj])) != 0:
                    return True
    return False
\end{lstlisting}

It routes the momenta, then for each edge's $k^{2}$ checks for a nonzero
\emph{mixed second partial} --- a $q\cdot\ell$ cross term. A bubble has such a
term (some edge carries $(q-\ell)^{2}$); a tadpole has every edge at pure
$q^{2}$, $\ell^{2}$, or $0$. The test is topology-agnostic.

\begin{gotcha}
A topological bubble whose scalar prefactor is $\propto\phi^{*2}$ is \emph{dead}
at $\phi^{*}=0$ and must not trigger the bubble route.
\code{\_prefactor\_is\_live} (\file{pipeline\_bridge.py:1058}) substitutes the
saddle-carrying \code{num\_params} and filters these out: only \emph{live}
bubbles count. (If free symbols remain it conservatively returns true.)
\end{gotcha}

\section{General-\texorpdfstring{$k$}{k} cumulants}\label{sec:kpoint}

For correlators with three or more external legs the bridge has two more drivers,
the $k$-point companions of the tree and loop paths.
\code{compute\_spatial\_kpoint} (\file{pipeline\_bridge.py:1788}) is the
single-field general-$k$ path; \code{compute\_coupled\_kpoint}
(\file{pipeline\_bridge.py:2000}) is the coupled scalar-diffusion companion. Both
take \code{points}, an array of $(x_{j},\tau_{j})$ offsets of slots
$1,\dots,k-1$ relative to slot~0.

The key structural difference from $k=2$ is how external legs are wired. The
$k$-point drivers use an \term{orbit-stabilizer external-Wick mapping sum}
(\code{field\_respecting\_mappings} plus \code{external\_wick\_compensation})
rather than the explicit orientation and $\pm\tau$ completion of the two-point
case. As the comment at \file{pipeline\_bridge.py:509} notes, the two views
coincide at $k=2$; at higher $k$ the mapping sum is the general machinery. Points
that share a $\tau$-configuration are batched, and the coupled $k$-point driver
is scoped to scalar diffusion, plain vertices, and infinite boundary.

\section{Operational guardrails}

Several of the boxes above are instances of a single design stance: this
subsystem prefers to \emph{refuse} a computation it cannot certify rather than
return a plausible-looking wrong number. Collected in one place, the guardrails
are:

\begin{itemize}
  \item \textbf{The macOS fork prohibition (the big one).} Fork-based parallelism
        was \emph{removed} from the spatial path. Forking a Jupyter kernel on
        macOS after the GUI/BLAS libraries have initialized crashes the kernel
        \emph{and} the OS, even with a single worker. The only safe speedup is
        thread-based (NumPy releases the GIL during heavy array ops) or
        $q$-grid NumPy batching. The numerical-FT path uses a
        \emph{smart-gated} \code{ThreadPoolExecutor}
        (\file{pipeline\_bridge.py:1713}): threads engage only when a
        loop-order-$\ge2$ diagram is present (about $2.5\times$ at two loops;
        \emph{slower} at one loop, where the tiny matrices are dispatch-bound, so
        one loop stays serial), with accumulation on the main thread for
        bit-identical results.
  \item \textbf{The memory budget.} As above: refuse up front when a chamber's
        quadrature would exceed \code{SPATIAL\_MEM\_BUDGET\_GB}.
  \item \textbf{Certification is a hard gate.} A failed mode certification raises
        \code{SpatialPropagatorError}.
  \item \textbf{Defective reaction matrix.} \code{spectral\_projectors} raises if
        the eigenvector condition number exceeds $10^{10}$.
  \item \textbf{Dyson divergence.} $\rho\ge1$ raises; $\rho>0.6$ warns.
  \item \textbf{Drift $V\ne0$ is refused} on both the single-field generic and the
        coupled-tree paths.
  \item \textbf{Never default a leg index silently.} Both
        \code{\_legs\_to\_phys\_idx} (\file{pipeline\_bridge.py:89}) and the
        per-leg index resolver \emph{raise} rather than fall back --- a wrong leg
        would quietly compute a \emph{different} correlator ($C_{aa}$ instead of
        $C_{ab}$). The point is emphasized twice in the source comments
        (\file{pipeline\_bridge.py:122-124,542-543}).
\end{itemize}

\begin{defn}
The \term{GIL} (Global Interpreter Lock) is the CPython mechanism that lets only
one thread execute Python bytecode at a time. The reason \emph{thread}
parallelism can still help here is that NumPy's heavy array operations release
the GIL while they run in compiled C --- so several threads can do NumPy work
genuinely in parallel. Process/fork parallelism would side-step the GIL too, but
it is forbidden here for the macOS-crash reason above.
\end{defn}

\begin{note}
\textbf{Environment escape hatches.} Several knobs exist for cross-checks and
debugging, not normal use: \code{SPATIAL\_FORCE\_NUMERICAL\_FT=1} keeps the
numerical-FT path reachable; \code{SPATIAL\_Q\_CUT} and \code{SPATIAL\_N\_Q} tune
it; \code{SPATIAL\_DYSON\_ORDER} overrides the model's Dyson policy;
\code{SPATIAL\_INTEGRATOR}, \code{SPATIAL\_GRID\_NT}/\code{NS},
\code{SPATIAL\_MC\_N}, and \code{SPATIAL\_MEM\_BUDGET\_GB} tune the loop
integrator. In ordinary use you set none of them.
\end{note}

\section{The oracle module: \texttt{loop\_dyson.py}}

The fifth module, \file{loop\_dyson.py}, carries a prominent banner
(\file{loop\_dyson.py:1-5}):

\begin{quote}\itshape
``$\triangle$ ORACLE-ONLY --- not on the production path. Superseded by
\code{full\_integrator.py}; reached only by its own test(s).
\code{compute\_cumulants} does NOT use this module.''
\end{quote}

It is an \emph{independent} implementation of the one-loop $\tilde\phi\phi^{2}$
reaction-diffusion \term{bubble} self-energy via the MSR Dyson equation --- both
the equal-time structure-factor correction $\delta C(q,0)$ and the full
$\tau$-dependent correction. It computes the retarded and Keldysh bubbles
($\Sigma_{R}=\int\,\dd\ell\,F(\ell)\,\Gret(\ell,t)\,C(q-\ell,t)$ and the
even-in-$t$ Keldysh analogue) by adaptive Gauss--Kronrod quadrature
(\code{scipy.integrate.quad}), and assembles the Dyson terms with the pinned
normalization \code{C\_R, C\_K = 4.0, 2.0} (\file{loop\_dyson.py:89}), validated
to give a fitted coefficient $B=0.99$ against a perturbative simulation.

Its value is precisely that it shares \emph{no code} with the production
\code{full\_integrator} path: a number that both produce independently is a
genuine cross-check. But the $C_{R}=4,C_{K}=2$ normalization and the $B=0.99$
validation are historical context, not the live computation.

\begin{defn}
\code{scipy.integrate.quad(f, a, b, limit=...)} is \term{adaptive Gauss--Kronrod
quadrature} (QUADPACK's \code{QAGS}/\code{QAGI}). Unlike the fixed-node
Gauss--Legendre rule, it recursively subdivides the interval until a local error
estimate is met; \code{limit} caps the number of subdivisions. The oracle uses it
for the self-energy momentum integrals and the Dyson time integrals.
\end{defn}

\section{What you just learned}

This chapter covered the part of the engine that handles every spatial case the
diagonal path cannot:

\begin{enumerate}
  \item \textbf{Coupled fields} need a matrix propagator. The reference-diffusion
        split $\mathcal{D}=D_{0}I+\widehat{\mathcal{D}}$ and the spectral
        resolution $M=\sum m_{\alpha}P_{\alpha}$ give a closed-form reference
        $G_{0}$ that is \emph{exact} when diffusion is scalar
        (\file{spectral\_propagator.py}).
  \item \textbf{The tree two-point} of a coupled linear theory is the
        Lyapunov/fluctuation--regression $C(q,\tau)=\ee^{-A|\tau|}\Sigma$, not a
        sum of independent modes (\code{coupled\_two\_point}).
  \item \textbf{Unequal diffusion} is handled by the Dyson--Duhamel series in
        powers of $\widehat{\mathcal{D}}|k|^{2}$, whose one new primitive
        $\Phi_{n}$ is evaluated confluent-safely by the Opitz matrix exponential
        (\file{dyson\_dressing.py}).
  \item \textbf{Derivative vertices} are hosted as inert operator nodes in an IR
        with explicit linearity, saddle-expansion, annihilation, and lowering
        passes; each deposits a momentum form factor whose loop average is taken
        either numerically or exactly by the Wick-moment construction
        (\file{spatial\_operator\_ir.py}, \code{\_build\_wick\_moment}).
  \item \textbf{The bridge} ties it all to the shared diagram pipeline via the
        $\Lap\to-q^{2}$ substitution, certifies the diagram answer against the
        per-mode reference, dispatches to the right driver, and guards
        aggressively against the cases it cannot certify
        (\file{pipeline\_bridge.py}).
\end{enumerate}

The recurring lesson is the one in the bridge's docstring and in nearly every
\emph{Gotcha} box above: the engine reuses the temporal machinery wherever it
can, and where it cannot prove the answer is right it refuses to return one. The
next part of the manual steps back out to the assembly layer ---
\code{compute\_cumulants} itself and the user-facing engine API --- which is what
selects every driver this chapter described.


% =====================================================================
%  PART VI  --  ASSEMBLY & THE USER SURFACE
% =====================================================================
\part{Assembly and the User Surface}
% ---------------------------------------------------------------------
%  Chapter 17 : The compute_cumulants Orchestrator
% ---------------------------------------------------------------------
\chapter{The \texttt{compute\_cumulants} Orchestrator}\label{ch:compute-orchestration}

Every chapter before this one has opened up a single subsystem in isolation: the
field-theory core that expands the action, the propagator builder, the
mean-field solver, the diagram enumerator, the symmetry-factor machinery, Phase~J
in the time domain, the spatial integrator. This chapter is where they are
\emph{wired together}. There is exactly one function in the codebase that calls
all of them in order, threads the data from each into the next, and hands you
back a finished correlation function: \code{compute\_cumulants}, defined in
\file{pipeline/compute.py:108}.

If the rest of the manual is the anatomy of the engine, this chapter is its
nervous system. \code{compute\_cumulants} is the function the front-desk module
\code{daedalus.dd.run} ultimately calls (you saw that call in
Chapter~\ref{ch:quickstart}, at \file{notebooks/daedalus.py:586}); it is the
function the simulation-comparison harnesses call; it is the function the example
notebooks narrate when they print their seven-phase trace. Understanding it is
understanding how a model declaration becomes a number.

The crucial thing to hold in mind from the first page is that
\code{compute\_cumulants} \emph{does almost no arithmetic of its own.} It is
\term{orchestration}, not computation. Its job is to sequence seven phases,
carry a handful of data dictionaries between them, decide whether the model is
temporal or spatial and branch accordingly, assemble the per-loop-order pieces
into a master correlator, and stamp wall-time onto each phase so you can see
where the seconds went. The heavy lifting all happens in the sibling modules
this chapter's code merely calls. So this chapter documents two things in
parallel: the \emph{glue} (what \code{compute\_cumulants} itself does, line by
line) and the \emph{contracts} (what each phase promises to return, so the next
phase can consume it). Where we name a function from another module we cite it
with its file and line so you can jump straight to its own chapter.

\section{Where this function sits in the pipeline}\label{sec:co-where}

Recall the spine from Chapter~\ref{ch:quickstart}. The user writes a
\term{theory file}, \code{dd.load\_theory} turns it into a \code{model} dict, and
\code{dd.run} packages the run choices and calls the engine. The engine is this
function. End to end:

\begin{lstlisting}[style=console]
  model dict  (from a TheoryBuilder / hand-written models/*.py)
       |
       |  + fundamental (numeric params)  + k, max_ell, external_fields
       v
  +-----------------------------------------------------------------+
  |          compute_cumulants   (pipeline/compute.py:108)          |
  |                                                                 |
  |  [1] expand      FieldTheory.expand()      -> bigraded action   |
  |  [2] propagator  build_propagator()        -> K_ft, G_ft, ...   |
  |  [3] mean_field  solve_mean_field[_dae]()  -> num_params, saddle|
  |  [3.5] spatial?  -> momentum integrator (early return)          |
  |  [4] poles       compute_poles_and_residues() -> pole_vals,Cmat |
  |  [5] diagrams    enumerate_unique_diagrams() -> typed diagrams  |
  |  [6] classify    classify_coefficient_factors() -> prefactors   |
  |  [7] phase_j     compute_correction_td() per ell -> C(tau)      |
  +-----------------------------------------------------------------+
       |
       v
  result dict  -> notebooks (plotting), pipeline.save (npz/csv),
                  sim-comparison harnesses
\end{lstlisting}

\textbf{What feeds it.} The \code{model} dict (the subject of the
\emph{theory-files} chapters), built either by a \code{TheoryBuilder} or written
by hand under \file{models/}. Alongside it: the numerical parameter point
\code{fundamental} (a dictionary mapping each coupling's name to its value), the
cumulant order \code{k} (the number of external legs), the maximum loop order
\code{max\_ell}, and \code{external\_fields} (which field sits on each external
leg).

\textbf{What consumes its output.} The returned \code{result} dict. Its
\code{total\_C} callable and its \code{C\_tau} / \code{C\_tau\_by\_ell} arrays
are what the notebooks plot and what the simulation harnesses compare against;
\code{pipeline.save.save\_npz} and \code{save\_csv} serialize it; the \code{mf}
and \code{params} accessors expose the saddle values and parameters by their
natural names. We dissect that dict in Section~\ref{sec:co-result}.

\begin{note}
The seven labels \code{[1/7]} through \code{[7/7]} that the engine prints in a
verbose run --- you saw them in Chapter~\ref{ch:quickstart} --- are emitted by
this file, one \code{print} per phase. The whole rest of this manual is, in a
sense, a deep dive into one or more of those seven labels. This chapter is the
table of contents that ties them together.
\end{note}

\section{The \protect\msrjd{} expansion in one page}\label{sec:co-msrjd}

To read the orchestrator you need only the skeleton of the physics; the
field-theory chapters fill it in. A stochastic dynamics for a field $\varphi$ ---
a firing rate, a density, a height field --- is recast as a path integral over
$\varphi$ \emph{and} a conjugate \term{response field} $\tilde\varphi$ (the
``tilde'' field). The generating functional is
\begin{equation}
  Z[j,\tilde\jmath]
  \;=\;
  \int \mathcal{D}\varphi\,\mathcal{D}\tilde\varphi \;
  \exp\!\Big(-S[\varphi,\tilde\varphi]
    + \!\int (\tilde\jmath\,\varphi + j\,\tilde\varphi)\Big).
  \label{eq:co-Z}
\end{equation}
This is the \term{Martin--Siggia--Rose / Janssen--De Dominicis} (\msrjd)
construction. The action $S$ splits into three pieces:

\begin{itemize}
  \item a \term{free (Gaussian, bilinear) part} $S_0$, quadratic in the fields;
        its inverse is the bare propagator;
  \item an \term{interaction part} $S_{\mathrm{int}}$, the cubic and higher
        vertices;
  \item a \term{mean-field / saddle part}, the linear and constant terms, which
        vanish exactly when the expansion is taken about the correct saddle
        point.
\end{itemize}

Connected correlation functions --- \term{cumulants} --- are obtained by
functionally differentiating $\ln Z$. The key fact for everything that follows is
that the connected $k$-point cumulant is precisely the sum of all
\emph{connected} Feynman diagrams with $k$ external legs. \code{compute\_cumulants}
automates that sum: it expands the action, builds the propagator, solves for the
saddle, enumerates every diagram up to \code{max\_ell} loops, computes each
diagram's symmetry factor, and integrates each one.

\begin{defn}
\textbf{Bigrade.} Each monomial in the expanded action carries a \term{bigrade}
$(n_t, n_p)$ = (number of response-field $\tilde\varphi$ factors, number of
physical-field $\varphi$ factors). The bigrade is how the expanded action is
sorted into sectors:
\begin{itemize}
  \item \textbf{Mean-field sector} --- bigrades $(0,0)$, $(1,0)$, $(0,1)$. These
        linear/constant terms vanish at the saddle.
  \item \textbf{Free action} --- bigrade $(1,1)$, the bilinear kernel. Its
        time-domain matrix is \code{K\_ker}; its Fourier image is \code{K\_ft};
        its inverse \code{G\_ft = K\_ft$^{-1}$} is the bare propagator.
  \item \textbf{Interaction sectors} --- total degree $\ge 2$ excluding $(1,1)$.
        These are the vertices of the diagrammatic expansion.
\end{itemize}
\end{defn}

The helper \code{\_print\_action\_sectors} (\file{pipeline/compute.py:50}) prints
exactly these three groups after the expand step, when \code{verbose=True}. It is
the ``here is your action'' diagnostic and has no effect on the computation; we
return to it in Section~\ref{sec:co-helpers}.

\section{The seven phases at a glance}\label{sec:co-glance}

The body of \code{compute\_cumulants} is, at heart, a straight-line sequence of
seven numbered phases with a single spatial-versus-temporal fork wedged between
phases [3] and [4]. Here is the whole tour in one table; the sections that follow
open each row.

\begin{center}
\small
\begin{longtable}{@{}p{0.10\textwidth}p{0.32\textwidth}p{0.50\textwidth}@{}}
\toprule
\textbf{Phase} & \textbf{Calls} & \textbf{Produces} \\
\midrule
\endhead
{}[1] expand     & \code{FieldTheory.expand()} & the bigraded action \code{ft} (free, interaction, MF sectors) \\
{}[2] propagator & \code{build\_propagator()}  & \code{prop} dict (\code{K\_ft}, \code{G\_ft}, \code{D\_delta}, $\omega$, \code{nf}) \\
{}[3] mean field & \code{solve\_mean\_field[\_dae]()} & \code{num\_params}: the $\{$\code{SR}$\to$float$\}$ substitution dict \\
{}[3.5] spatial? & spatial integrator & \emph{early return} for spatial models \\
{}[4] poles      & \code{compute\_poles\_and\_residues()} & fills \code{prop['pole\_vals']}, \code{prop['C\_mats']} \\
{}[5] diagrams   & \code{enumerate\_unique\_diagrams()} & \code{unique\_by\_ell}: $\{\ell:$ typed diagrams$\}$ \\
{}[6] classify   & \code{classify\_coefficient\_factors()} & \code{diagram\_records}: prefactor + $\Sym(\Gamma)$ per diagram \\
{}[7] phase J    & \code{compute\_correction\_td()} per $\ell$ & \code{C\_tau\_by\_ell}, \code{total\_C\_by\_ell} \\
\bottomrule
\end{longtable}
\end{center}

The data threads through linearly: \code{ft} from [1] feeds [2], [3], and [5];
the \code{prop} dict from [2] is filled in place by [4] and consumed by [7];
\code{num\_params} from [3] is the universal substitution dictionary every
numerical phase needs. We trace that flow explicitly in
Section~\ref{sec:co-dataflow}.

\section{Argument validation and the Taylor budget}\label{sec:co-args}

Before any phase runs, the function validates its arguments and picks a few
derived quantities. The full signature
(\file{pipeline/compute.py:108}) is long because the function is the single
public surface for a great many run options:

\begin{lstlisting}
def compute_cumulants(
    model: dict,
    k: int,
    max_ell: int = 0,
    fundamental: dict = None,
    external_fields: list[tuple[str, int]] = None,
    *,
    tau_max: float = 50.0,
    tau_step: float = 0.5,
    tau_grid=None,                   # explicit tau grid; overrides tau_max/tau_step
    spatial_grid=None,
    taylor_order: int = None,
    origin_leaf_idx: int = 0,
    output_npz: str = None,
    output_csv: str = None,
    use_cache: bool = True,
    parallel: bool = True,
    spatial_parallel: bool = True,
    spatial_n_q: int = 64,
    spatial_points=None,
    n_workers: int = None,
    use_grouped_phase_j: bool = False,
    fixed_point_index: int = 0,
    mf_dae_n_starts: int = 64,
    mf_dae_seed_box: dict | None = None,
    verbose: bool = True,
) -> dict[str, Any]:
\end{lstlisting}

The load-bearing arguments are the first five: \code{model}, \code{k},
\code{max\_ell}, \code{fundamental}, and \code{external\_fields}. The keyword-only
block (\code{*}) after them tunes the run --- grid extents, caching, parallelism,
the spatial APIs, the DAE multi-root solver --- and all have defaults.

\textbf{Validation} (\file{pipeline/compute.py:253}). The function defaults a
missing \code{fundamental} to \code{\{\}}, then enforces the one structural
invariant it can check up front: \code{external\_fields} must be present and must
have exactly \code{k} entries.

\begin{lstlisting}
if fundamental is None:
    fundamental = {}
if external_fields is None:
    raise ValueError('external_fields is required')
if len(external_fields) != k:
    raise ValueError(
        f'external_fields has {len(external_fields)} entries but k={k}'
    )
\end{lstlisting}

\subsection{Why \texorpdfstring{$k + 2\ell$}{k + 2 ell} vertices, and the floor of 2}

The nonlinear functions in an action (a sigmoid firing rate, a quartic
restoring force) are Taylor-expanded to a finite order. How many terms do we
need? A connected diagram with $k$ external legs at $\ell$ loops needs vertices
of degree at most $k + 2\ell$: every loop you add introduces two extra legs that
have to terminate somewhere, and the worst case is that they all terminate on a
single vertex. So the smallest sufficient Taylor order is

\begin{equation}
  \texttt{taylor\_order} \;=\; \max(k + 2\,\texttt{max\_ell},\ 2),
  \qquad\text{(\file{pipeline/compute.py:279})}
  \label{eq:co-taylor}
\end{equation}

implemented as:

\begin{lstlisting}
if taylor_order is None:
    taylor_order = max(k + 2 * max_ell, 2)
\end{lstlisting}

The floor of $2$ is structural, not a heuristic. Order $2$ is the minimum that
still produces the bilinear $(1,1)$ propagator kernel and the
$(0,0)/(1,0)/(0,1)$ mean-field saddle sectors --- and downstream code reads those
\emph{unconditionally}, even in the degenerate $k=2,\ \texttt{max\_ell}=0$ case
where the tree-level pair correlator is just the bare propagator with no
interaction vertices at all.

\begin{note}
The docstring (\file{pipeline/compute.py:179}) records that this floor was
\emph{historically 4}, chosen for a now-obsolete cache layout that kept all
1-loop 2-cumulant runs in one directory. Lowering it from $4$ to $2$ saved about
$90$ minutes on heavy Bernoulli theories at $k=2,\ \texttt{max\_ell}=0$, the one
case where the old floor exceeded the mathematical minimum. If you are probing
higher-order vertices for cache-invalidation testing you must now pass an
explicit \code{taylor\_order} override.
\end{note}

\subsection{Natural names to internal names}

Users pass external fields by their natural names, e.g. \code{('n', 1)} for the
first population's physical field. The action and propagator-typing code, by
contrast, work with internal \emph{fluctuation} names, e.g. \code{('dn', 1)}. The
translation happens at \file{pipeline/compute.py:288}:

\begin{lstlisting}
naming_convention    = model.get('naming_convention')
external_fields_user = list(external_fields)
external_fields      = normalize_external_fields(
    external_fields, naming_convention=naming_convention)
\end{lstlisting}

\code{normalize\_external\_fields} (in \file{pipeline/access.py}) reads the
model's declared \code{naming\_convention} --- written by the
\code{TheoryBuilder} --- to map \code{('n',1)} to \code{('dn',1)}, and falls back
to a classic n/v/m map when the model declares none. Notice that the function
keeps the \emph{user's} original form in \code{external\_fields\_user}; both forms
end up in the result's \code{config} block (the internal form under
\code{external\_fields}, the user form under \code{external\_fields\_in}) so a
caller can always recover what it asked for.

\subsection{The per-phase wall-time tracker}

The last thing set up before phase [1] is the timing apparatus
(\file{pipeline/compute.py:296}):

\begin{lstlisting}
phase_walls: dict[str, float] | None = {} if verbose else None

def _phase_time(label: str, t0: float) -> None:
    if phase_walls is None:
        return
    dt = time.perf_counter() - t0
    phase_walls[label] = dt
    print(f'      [{label}] done in {dt:.2f}s')
\end{lstlisting}

\code{phase\_walls} is a dict mapping a phase label to its elapsed seconds, or
\code{None} when \code{verbose=False} --- so the bookkeeping is opt-in and costs
nothing in a quiet run. \code{time.perf\_counter()} is Python's highest-resolution
monotonic clock (it does not jump if the system clock is adjusted), which is why
it is the right tool for measuring elapsed wall-time rather than \code{time.time()}.
The closure \code{\_phase\_time(label, t0)} is called at the end of each phase
with that phase's start time; the per-phase summary at the very end of the
function (Section~\ref{sec:co-perf}) reads this same dict.

\section{Phase [1] --- expand, with the on-disk cache}\label{sec:co-expand}

\textbf{Code:} \file{pipeline/compute.py:305}--\code{:351}.

The first phase constructs the \code{FieldTheory} object and expands the action.
Expansion --- the multivariate Taylor pass that turns the symbolic action into
the bigraded sector dictionary --- is one of the more expensive symbolic
operations in the pipeline, so it is cache-aware:

\begin{lstlisting}
ft = FieldTheory(model, taylor_order=taylor_order)

from pipeline import _expand_cache as _ec
cache_hit = False
if use_cache:
    cached_order = _ec.find_best_cached_order(model, taylor_order)
    if cached_order is not None:
        _ec.prepare_for_load(ft)
        cache_hit = _ec.load_expand(
            model, ft,
            target_order=taylor_order,
            cached_order=cached_order,
            verbose=verbose,
        )
if not cache_hit:
    ft.expand()
    if use_cache:
        try:
            _ec.save_expand(model, ft, verbose=verbose)
        except Exception as e:
            ...
\end{lstlisting}

The logic is worth spelling out because it is more clever than a plain
hit-or-miss cache. \code{find\_best\_cached\_order} looks for \emph{any} cached
expansion at an order \emph{$\ge$} the one requested --- because an order-6
expansion already contains every term of an order-4 expansion. If it finds one,
\code{prepare\_for\_load} and \code{load\_expand} downgrade-filter the cached
sectors to the requested total degree and load them. Only on a genuine miss does
it pay for \code{ft.expand()} and then \code{save\_expand} the result. The
\code{try/except} around the save is deliberate: a failed cache write is logged
and ignored, never fatal --- you would rather finish the run uncached than crash
on a read-only disk.

After expansion the code runs a self-check and extracts the vertex/source types:

\begin{lstlisting}
sanity_ok = ft.sanity_check(verbose=verbose)
if not sanity_ok:
    raise RuntimeError(
        'FieldTheory.sanity_check() failed -- see printout above.')

vtypes = extract_vertex_types(ft)
stypes = extract_source_types(ft)
n_noise = sum(1 for s in stypes if isinstance(s, NoiseSourceType))
\end{lstlisting}

\code{vtypes} are the interaction vertices the enumerator will draw from;
\code{stypes} are the source (noise) vertices, and \code{n\_noise} counts how
many are genuine \code{NoiseSourceType}s. These three are produced here and
carried all the way to phase [5]. When verbose, the phase ends by calling
\code{\_print\_action\_sectors(ft)} (the diagnostic of
Section~\ref{sec:co-msrjd}) and records its wall-time with
\code{\_phase\_time('expand', \_t\_phase)}.

\section{Phase [2] --- the propagator}\label{sec:co-prop}

\textbf{Code:} \file{pipeline/compute.py:353}--\code{:358}.

The propagator is the inverse of the free action's $(1,1)$ kernel. This phase is
a single call:

\begin{lstlisting}
prop = build_propagator(ft, model, use_cache=use_cache, verbose=verbose)
\end{lstlisting}

\code{build\_propagator} (\file{pipeline/\_propagator.py:524}) builds (or loads
from \file{saved\_theories/<tag>/propagator.sobj}) the symbolic propagator
dictionary \code{prop}. Its keys are \code{K\_ker}, \code{K\_ft}, \code{G\_ft},
\code{adj\_ft}, \code{D\_omega}, \code{D\_delta}, \code{t\_var}, \code{omega},
\code{nf}, and \code{ring\_gen\_names}, plus \code{pole\_vals} and \code{C\_mats}
placeholders that phase [4] will fill. The complete anatomy of that dict is in
the propagator chapter and Section~\ref{sec:co-datastructures}; the orchestration
view is just: \emph{this is where the bare two-point function comes from.}

\subsection{The math the propagator chapter spells out}

The free action's $(1,1)$ sector is a matrix \code{K\_ker}$(t)$ of time-domain
kernels (containing Dirac $\delta(t)$, its derivative $\delta'(t)$, and any model
kernel symbols). Fourier-transforming gives \code{K\_ft}$(\omega)$, a matrix of
rational functions of $\omega$. The bare retarded propagator is its matrix
inverse:
\begin{equation}
  \text{\code{G\_ft}}(\omega) \;=\; \text{\code{K\_ft}}(\omega)^{-1}
  \qquad\text{(rows = physical, cols = response).}
  \label{eq:co-Gft}
\end{equation}
Its poles in the upper-half $\omega$-plane are the retarded relaxation modes; the
residue matrices \code{C\_mats} at those poles give the time-domain propagator as
a sum of decaying exponentials,
\begin{equation}
  G(t)
  \;=\;
  \sum_k C_k\, \ee^{-\ii \omega_k t}\,\Theta(t)
  \;+\;
  \texttt{D\_delta}\,\delta(t),
  \label{eq:co-Gt}
\end{equation}
where \code{D\_delta}$= \lim_{\omega\to\infty}$\code{G\_ft}$(\omega)$ is the
instantaneous (white) part. Phase [2] produces the symbolic \code{G\_ft} and
\code{D\_delta}; phase [4] produces the numbers $\omega_k$ (in \code{pole\_vals})
and $C_k$ (in \code{C\_mats}). For ``rich'' kernels --- large matrices, many
symbols --- \code{build\_propagator} deliberately skips the symbolic inverse and
leaves \code{G\_ft = None}, deferring the whole job to the numerical pole-finder
in phase [4].

\section{Phase [3] --- the mean-field saddle}\label{sec:co-mf}

\textbf{Code:} \file{pipeline/compute.py:360}--\code{:380}.

The loop expansion is taken about the deterministic mean-field steady state. This
phase finds it and --- more importantly for the orchestrator --- assembles the
substitution dict that makes every later phase numeric:

\begin{lstlisting}
if model.get('equations'):
    mf = solve_mean_field_dae_compat(
        ft, model, fundamental,
        fixed_point_index=fixed_point_index,
        n_starts=mf_dae_n_starts,
        seed_box=mf_dae_seed_box,
        verbose=verbose,
    )
else:
    mf = solve_mean_field(ft, model, fundamental, verbose=verbose)
num_params = mf['num_params']
\end{lstlisting}

There are two solvers, chosen by whether the model declares \code{equations}:

\begin{itemize}
  \item \code{solve\_mean\_field} (\file{pipeline/\_mean\_field.py:14}) is the
        legacy / single-population path. It assembles the parameter substitutions
        from \code{fundamental}, reads the symbolic background $v^{*}$ from the
        model's \code{mf\_bg\_conditions}, and solves the self-consistency
        $n^{*}_i = \langle\,\text{RHS}\,\rangle(n^{*}_j)$ numerically.
  \item \code{solve\_mean\_field\_dae\_compat}
        (\file{pipeline/\_mean\_field\_dae.py}) is the DAE-based multi-root
        solver, used when the theory declares its dynamics via
        \code{.equation(lhs=..., rhs=..., population=...)}. It layers a
        multi-start global root search on top, returning \emph{every} fixed point
        with a stability annotation, and \code{fixed\_point\_index} selects among
        the stable ones.
\end{itemize}

\begin{defn}
\textbf{\code{num\_params}.} \code{num\_params} is a dictionary
$\{$\code{SR.var} $\to$ \code{float}$\}$. It
maps every symbolic parameter, every solved saddle value (\code{nstar\_i},
\code{vstar\_i}, \code{mstar\_i}), and every $\varphi$-derivative (\code{phiN\_i})
to a concrete number. It is the \emph{universal substitution dictionary} ---
every downstream phase substitutes it into its symbolic expressions to get
numbers. If any symbol is left out of \code{num\_params}, the corresponding
numerical evaluation silently collapses (see Section~\ref{sec:co-gotchas}).
\end{defn}

\subsection{What \code{fsolve} actually does}

The numerical root-finding inside \code{solve\_mean\_field} uses SciPy. SciPy is
the standard scientific-computing library for Python; its \code{scipy.optimize}
submodule provides numerical root-finders and optimizers. The call
(\file{pipeline/\_mean\_field.py:174}) is:

\begin{lstlisting}
from scipy.optimize import fsolve              # _mean_field.py:11
attempt = fsolve(mf_residual, x0, full_output=True)
\end{lstlisting}

\code{fsolve} finds a root of the multivariate residual
$\texttt{mf\_residual}(n^{*}) = n^{*} - \text{RHS}(n^{*})$ --- the saddle
self-consistency written as ``input minus output''. \code{full\_output=True}
makes it return the tuple \code{(x, infodict, ier, msg)}, where \code{ier == 1}
signals clean convergence; the solver tries a sequence of initial guesses and
keeps the first that converges. This is the only place a Newton-type solver
appears in the temporal pipeline.

\section{Phase [3.5] --- the spatial-versus-temporal fork}\label{sec:co-spatial}

\textbf{Code:} \file{pipeline/compute.py:382}--\code{:626}.

Between the saddle and the pole-finder sits the single biggest branch in the
function. A spatial (PDE) model cannot go down the temporal road at all: its
propagator \code{G\_ft} carries an \emph{inert} \code{Laplacian} symbol that the
$\omega$-pole-finder and the time-domain Phase~J simply cannot consume. So
spatial models short-circuit here, routing to the momentum-space integrator and
returning early with a \emph{different} result dict.

\subsection{The branch logic}

First, two guard cases:

\begin{lstlisting}
if spatial_grid is not None and not model.get('spatial'):
    import warnings as _warnings
    _warnings.warn(
        'spatial_grid was provided but the model is not spatial ...', stacklevel=2)
\end{lstlisting}

Passing a \code{spatial\_grid} to a \emph{non}-spatial model is a no-op: the code
warns (\file{pipeline/compute.py:390}) and falls through to the temporal path.
The substantive branch is the converse --- the model \emph{is} spatial \emph{and}
a spatial input (\code{spatial\_grid} or \code{spatial\_points}) was given:

\begin{lstlisting}
if model.get('spatial') and (spatial_grid is not None
                             or spatial_points is not None):
    ...
\end{lstlisting}

Inside, the path splits on \code{k}.

\textbf{The $k=2$ grid path.} This is the workhorse spatial case: a two-point
correlator $C(x,\tau)$ on a grid of $x$-points. It requires \code{spatial\_grid}
(passing \code{spatial\_points} at $k=2$ raises a \code{ValueError} at
\file{pipeline/compute.py:414}, because \code{spatial\_points} is the $k\ge3$
event API). It guards \code{max\_ell $\le$ 2} and the stationary initial
condition, builds the $\tau$-grid and the $x$-grid, and then dispatches by loop
order: for \code{max\_ell $\ge$ 1} it calls
\code{compute\_spatial\_correlator\_generic} (the full-diagram momentum
integrator), and for tree level (\code{max\_ell == 0}) it calls
\code{compute\_spatial\_correlator\_via\_pipeline}. Crucially, the loop path
passes \code{parallel=(parallel and spatial\_parallel)}
(\file{pipeline/compute.py:511}) --- the two flags are \emph{and}-ed, and the
spatial parallelism is thread-based, not fork-based (Section~\ref{sec:co-parallel}).

It then builds a \code{total\_C} closure that evaluates the real-space correlator
at an arbitrary point, applying the 1-loop tadpole mass-shift correction
$\delta C = \Sigma\cdot\partial C_0/\partial A$ by finite difference for $d=1$
and by a radial inverse Fourier transform for $d\ge2$, and \textbf{returns early}
(\file{pipeline/compute.py:584}) a spatial dict whose shape we give in
Section~\ref{sec:co-result}.

\textbf{The $k\ge3$ event path.} For three or more legs the spatial output is not
a grid but a set of explicit \emph{events}, so it requires \code{spatial\_points}
--- an \code{(n\_pts, k-1, 2)} array of $(x_j, \tau_j)$ offsets. It calls
\code{compute\_spatial\_kpoint}; if that raises a \code{NotImplementedError}
mentioning \code{'single-field'}, the code falls through to the
\emph{coupled multi-field} driver, which builds reaction-diffusion matrices
$M, D, V$ from \code{prop['K\_ft']}, checks that the drift $V$ is zero
(\file{pipeline/compute.py:450}; nonzero drift raises), splits the reference
diffusion, and calls \code{compute\_coupled\_kpoint}. This path \textbf{returns
early} (\file{pipeline/compute.py:463}) a dict keyed by \code{C\_kpoint}.

\subsection{The clear-failure guard}

If a spatial model reaches the end of the branch without \emph{any} spatial input,
it would otherwise fall into the temporal pole-finder and die deep in Sage with a
cryptic \code{TypeError}. The code refuses to let that happen
(\file{pipeline/compute.py:619}):

\begin{lstlisting}
if model.get('spatial'):
    raise ValueError(
        f"model {model.get('name', '<unnamed>')!r} is spatial ...: "
        f"compute_cumulants requires spatial_grid=... to route to the "
        f"spatial integrator. ...")
\end{lstlisting}

\begin{gotcha}
The two spatial APIs are not interchangeable. \code{spatial\_grid} is the $k=2$
grid API and \code{spatial\_points} is the $k\ge3$ event API; mixing them raises.
A spatial $k\ge3$ run returns a different dict shape than a $k=2$ run (events, not
a $\tau$-grid), so any caller that handles both must branch on \code{k}. And a
spatial model with \emph{no} spatial input raises a clear \code{ValueError} rather
than a cryptic Sage error --- if you see that message, you forgot to pass
\code{spatial\_grid} (try the theory's \code{METADATA['spatial\_grid']}).
\end{gotcha}

The remainder of this chapter follows the \emph{temporal} path; the spatial
integrator has its own chapters --- this one is the junction, not the destination.

\section{Phase [4] --- numerical poles and residues}\label{sec:co-poles}

\textbf{Code:} \file{pipeline/compute.py:628}--\code{:633}.

With the saddle solved and \code{num\_params} in hand, the symbolic propagator
\code{G\_ft}$(\omega)$ can be made numeric and its poles extracted:

\begin{lstlisting}
compute_poles_and_residues(prop, num_params, verbose=verbose)
\end{lstlisting}

\code{compute\_poles\_and\_residues} (\file{pipeline/\_propagator.py:851}) uses a
three-tier strategy --- an exact polynomial fraction-field over
\code{CyclotomicField(4)} = $\mathbb{Q}[\ii]$, a numpy-cofactor fallback, and a
legacy symbolic path --- to find the retarded poles ($\operatorname{Im}\omega > 0$)
and the residue matrices. The exact path computes
$C_k[i,j] = \ii\, P_{ij}(\omega_k)/Q'(\omega_k)$.

\begin{gotcha}
\code{compute\_poles\_and\_residues} \emph{mutates \code{prop} in place}
(\file{pipeline/compute.py:632}). There is no return value used: the orchestrator
relies entirely on the side effect of filling \code{prop['pole\_vals']} and
\code{prop['C\_mats']} (and \code{D\_delta}, if phase [2] left it \code{None}).
The very same \code{prop} object is later sliced into \code{propagator\_data} for
phase [7] and returned in the result dict. If you are reading the code top to
bottom, do not expect an assignment here --- the data appears inside \code{prop}.
\end{gotcha}

\section{Phase [5] --- enumerating the diagrams}\label{sec:co-diagrams}

\textbf{Code:} \file{pipeline/compute.py:635}--\code{:664}.

This phase produces the list of Feynman diagrams to integrate. First it builds
the field-index maps that tell the enumerator which ring variable is which field:

\begin{lstlisting}
ring_var_names = list(ft._ns._ring_var_names)
n_tilde = ft._n_tilde
resp_idx, phys_idx = build_field_index_map(ring_var_names, n_tilde)
for f in external_fields:
    if f not in phys_idx:
        raise ValueError(
            f'external field {f} not in phys_idx ...')
\end{lstlisting}

\code{build\_field\_index\_map} splits the ring's generator names into the
response block (the first \code{n\_tilde}) and the physical block, returning two
dicts \code{resp\_idx} and \code{phys\_idx}. The validation loop catches a mistyped
external field early --- before enumeration --- with a helpful message listing the
valid keys. Then the enumerator runs:

\begin{lstlisting}
unique_by_ell, multiplicity_by_ell, all_unique = enumerate_unique_diagrams(
    ft, model,
    k=k, max_ell=max_ell,
    external_fields=external_fields,
    G_ft=prop['G_ft'],
    resp_idx=resp_idx, phys_idx=phys_idx,
    vtypes=vtypes, stypes=stypes,
    use_cache=use_cache, parallel=parallel,
    n_workers=n_workers, verbose=verbose,
)
\end{lstlisting}

\code{enumerate\_unique\_diagrams} (\file{pipeline/\_diagrams.py:54}) returns the
dicts \code{unique\_by\_ell} ($\{\ell:$ list of \code{TypedDiagram}$\}$) and
\code{multiplicity\_by\_ell} ($\{\ell:$ list of int$\}$) plus the flat
concatenation \code{all\_unique}. For each loop order $\ell$ it checks a disk
cache and, on a miss, runs the four-stage diagram pipeline:

\begin{defn}
\textbf{Prediagram $\to$ typed $\to$ causal $\to$ unique.}
The four-stage enumeration is: \term{prediagram} enumeration (topological
skeletons), \code{enumerate\_all\_typed} (assigning a definite field type to every
leg --- this is the parallelizable stage, fork-based when \code{parallel=True}),
\code{filter\_causal} (dropping acausal diagrams), and
\code{deduplicate\_with\_multiplicities} (collapsing isomorphic graphs and
recording how many collapsed into each class).
\end{defn}

\begin{gotcha}
The cache stage name embeds a version suffix: \code{unique\_typed\_mult\_v3\_...}
(\file{pipeline/\_diagrams.py:150}). The \code{v3} is not decoration. Three
documented bumps lie behind it --- multiplicity-aware dedup, the
\code{NoiseSourceType} promotion, and the \emph{complete}-isomorphism signature.
Loading an older cache would silently resurrect a real bug (non-isomorphic
diagram classes collided and had their integrals dropped at $k\ge3$), so the
version suffix forces a rebuild rather than trusting stale data. The expand cache
and the propagator cache live separately, under \file{saved\_theories/<tag>/}.
\end{gotcha}

\section{Phase [6] --- symmetry factors and prefactors}\label{sec:co-classify}

\textbf{Code:} \file{pipeline/compute.py:666}--\code{:705}.

Each diagram needs a scalar weight before it can be integrated. This phase walks
the diagrams \emph{by loop order} (so each record carries its $\ell$ tag, which
phase [7] needs) and classifies each one:

\begin{lstlisting}
for ell in sorted(unique_by_ell.keys()):
    mults = multiplicity_by_ell.get(ell, [1] * len(unique_by_ell[ell]))
    for td, mult in zip(unique_by_ell[ell], mults):
        info = classify_coefficient_factors(
            td, time_dep_params, noise_structure)
        combined_prefactor = SR(info['scalar_prefactor'])
        kernel_groups.append({
            'diagrams':           [td],
            'combined_prefactor': combined_prefactor,
        })
        diagram_records.append({
            'typed_diagram':      td,
            'classify':           info,
            'combined_prefactor': combined_prefactor,
            'multiplicity':       mult,
            'ell':                ell,
        })
\end{lstlisting}

\code{classify\_coefficient\_factors} (\file{msrjd/diagrams/symmetry.py:618})
returns a dict with the keys \code{Scal}, \code{scalar\_prefactor},
\code{vertex\_time\_factors}, \code{source\_time\_info}, and \code{is\_stationary}.
The orchestrator pulls \code{scalar\_prefactor} out, wraps it with \code{SR(...)}
to coerce it to a symbolic expression, and appends to two parallel lists:
\code{kernel\_groups} (the legacy Phase-J input format) and \code{diagram\_records}
(the richer per-diagram record carried in the result).

\subsection{The symmetry factor \texorpdfstring{$\Sym(\Gamma)$}{S(Gamma)}}

The mathematical content of this phase is the combinatorial weight of each
diagram. A connected $k$-point cumulant is a weighted sum of diagrams; each
diagram $\Gamma$ contributes
\begin{equation}
  \mathrm{weight}(\Gamma)
  \;=\;
  \Sym(\Gamma)\,
  \prod_v c_{\mathrm{user},v}\,
  \prod_e P_e\,
  \int \prod_v \dd t_v,
  \label{eq:co-weight}
\end{equation}
where $c_{\mathrm{user},v}$ is the literal coefficient written in the action (no
implicit $1/n!$), $P_e$ is the propagator on edge $e$, the integral runs over the
internal vertex times, and $\Sym(\Gamma)$ is the combinatorial symmetry factor.
\Daedalus{} uses the canonical Feynman rule (``Path~A'' in this codebase):
\begin{equation}
  \Sym(\Gamma)
  \;=\;
  \frac{\displaystyle\prod_v \prod_\ell n_{v,\ell}!}
       {\bigl|\Aut_{\text{fixed-ext}}(\Gamma)\bigr|}.
  \label{eq:co-Sym}
\end{equation}
The numerator is the per-vertex Wick combinatorial --- factorials of
identical-leg multiplicities at each vertex, for both response and physical legs.
The denominator is the order of the automorphism group of the colour-preserving
incidence digraph, with the external leaves \emph{pinned} at their canonical
positions. This single ratio accounts for every Feynman-rule symmetry at once:
same-type vertex swaps, parallel-edge swaps, self-loop leg swaps.

The denominator is computed by Sage. \code{combinatorial\_factor}
(\file{msrjd/diagrams/symmetry.py:403}) computes the Wick numerator via
\code{\_wick\_leg\_factor} (\file{msrjd/diagrams/symmetry.py:115}) and the
automorphism order via \code{\_automorphism\_order}
(\file{msrjd/diagrams/symmetry.py:328}), which in turn calls Sage's
\code{D.automorphism\_group(partition=...)}
(\file{msrjd/diagrams/symmetry.py:339}). Sage routes that call through its bundled
graph-canonicalization backend (the \code{bliss} / \code{nauty}-family
algorithms), and \code{.order()} of the returned permutation group is
$|\Aut_{\text{fixed-ext}}(\Gamma)|$.

\begin{note}
\code{nauty} (``No AUTomorphisms, Yes?'') and \code{bliss} are C libraries that
compute graph automorphism groups and canonical forms --- the gold-standard tools
for graph isomorphism. You give them a graph with an optional vertex colouring and
they return either the automorphism group or a canonical labeling such that two
graphs are isomorphic if and only if their canonical forms are byte-identical.
\Daedalus{} never calls them directly; Sage's \code{DiGraph.automorphism\_group}
and \code{DiGraph.canonical\_label} (\file{msrjd/diagrams/symmetry.py:517}) wrap
them. The same machinery powers the dedup signature in phase [5]: the
\emph{canonical form} of the coloured incidence digraph is a complete isomorphism
invariant.
\end{note}

\begin{gotcha}
The dedup \code{multiplicity} from phase [5] is recorded in each record but
\emph{not multiplied} into the prefactor (\file{pipeline/compute.py:686}--\code{:692}).
Under Path~A the symmetry factor $\Sym(\Gamma)$ in \eqref{eq:co-Sym} already
carries each equivalence class's full weight via the orbit--stabilizer count, so
multiplying by the class size would double-count. \code{multiplicity} is kept
purely for diagnostics. (Historically a caller-side multiplication compensated for
an \emph{incomplete} signature that merged non-isomorphic diagrams; now that the
signature is a complete invariant, that compensation is obsolete and removed.)
\end{gotcha}

\subsection{What \code{classify\_coefficient\_factors} does internally}

For completeness, the per-diagram classifier (\file{msrjd/diagrams/symmetry.py:618})
does the following for one \code{TypedDiagram}:

\begin{enumerate}
  \item Computes \code{Scal = combinatorial\_factor(td)} --- the symmetry factor
        $\Sym(\Gamma)$.
  \item Walks every vertex; each vertex coefficient acquires a $(-1)$ because the
        weight is $\int \exp(-S)$ (\code{coeff = -SR(vtype.coefficient)},
        \file{msrjd/diagrams/symmetry.py:643}).
  \item For \emph{source (noise) vertices}, decides whether the amplitude can be
        pulled out of the time integral --- a time-independent white amplitude
        pulls out into the scalar prefactor; otherwise it stays in
        \code{source\_time\_info} for Phase~J.
  \item For \emph{interaction vertices}, splits the coefficient into a
        time-independent constant part (joining \code{scalar\_parts}) and a
        time-dependent part (kept in \code{vertex\_time\_factors}).
  \item Multiplies all the constant parts into \code{scalar\_prefactor}.
  \item Sets \code{is\_stationary} \code{True} iff there are no per-vertex time
        factors and no time-dependent noise amplitudes --- which tells Phase~J it
        may use the fast time-translation-invariant path.
\end{enumerate}

\section{Phase [7] --- Phase J and the per-\texorpdfstring{$\ell$}{ell} assembly}\label{sec:co-phasej}

\textbf{Code:} \file{pipeline/compute.py:707}--\code{:799}.

The final phase performs the actual loop integrals. It first builds the
$\tau$-grid and the \code{propagator\_data} dict --- precisely the subset of
\code{prop} that Phase~J consumes (\file{pipeline/compute.py:714}):

\begin{lstlisting}
tau_grid = (np.asarray(tau_grid, dtype=float) if tau_grid is not None
            else np.arange(-tau_max, tau_max + tau_step * 0.5, tau_step))

propagator_data = {
    'K_ker':   prop['K_ker'],   'K_ft':    prop['K_ft'],
    'G_ft':    prop['G_ft'],    'adj_ft':  prop['adj_ft'],
    'D_omega': prop['D_omega'], 'D_delta': prop['D_delta'],
    't_var':   prop['t_var'],   'omega':   prop['omega'],
    'nf':      prop['nf'],
    'pole_vals': prop['pole_vals'],  'C_mats': prop['C_mats'],
}
\end{lstlisting}

Note that \code{tau\_grid} runs from $-\texttt{tau\_max}$ to $+\texttt{tau\_max}$;
the \code{tau\_step * 0.5} nudge in \code{np.arange}'s upper bound is the standard
trick to make the endpoint inclusive despite floating-point round-off.

\subsection{Choosing the \texorpdfstring{$\tau$}{tau}-evaluation pattern}

The grid is evaluated differently depending on the number of external legs
(\file{pipeline/compute.py:729}):

\begin{lstlisting}
if k == 1:
    tau_points = [(float(t),) for t in tau_grid]
elif k == 2:
    tau_points = [(0.0, float(t)) for t in tau_grid]  # leaf 0 pinned, leaf 1 swept
else:
    tau_points = None   # k>=3: caller handles grid evaluation
\end{lstlisting}

For a one-point function each grid point is a single time; for a two-point
function leaf $0$ is pinned to $0$ and leaf $1$ is swept over the grid (the
familiar single-axis correlation slice $C(\tau)$); for $k\ge3$ no automatic grid
eval is done --- \code{tau\_points = None} and the caller supplies the evaluation
points.

\subsection{Once per loop order}

The heart of the phase is a loop over $\ell$. Running Phase~J \emph{separately for
each loop order} is what gives the per-order decomposition the notebooks plot as
``tree'', ``1-loop'', ``2-loop'':

\begin{lstlisting}
for ell in sorted(unique_by_ell.keys()):
    records_ell = [r for r in diagram_records if r['ell'] == ell]
    if not records_ell:
        total_C_by_ell[ell] = (lambda *t: 0.0 + 0.0j)
        phase_j_by_ell[ell] = None
        C_tau_by_ell[ell] = (np.zeros(len(tau_grid), dtype=complex)
                             if tau_points is not None else None)
        continue
    ...
    td_result_ell = compute_correction_td(
        typed_diagrams = [r['typed_diagram'] for r in records_ell],
        prefactors     = [r['combined_prefactor'] for r in records_ell],
        k=k, propagator_data=propagator_data,
        external_fields=external_fields, num_params=num_params,
        origin_leaf_idx=origin_leaf_idx,
    )
    total_C_by_ell[ell] = td_result_ell['total_C']
    phase_j_by_ell[ell] = td_result_ell
    if tau_points is not None:
        C_tau_by_ell[ell] = np.array(
            td_result_ell['total_C_batch'](
                tau_points, parallel=parallel, n_workers=n_workers),
            dtype=complex)
\end{lstlisting}

\code{compute\_correction\_td} (\file{msrjd/integration/time\_domain/pipeline.py:202})
is the per-$\ell$ workhorse: for each typed diagram it numerically integrates the
vertex times against the pole-residue propagator and accumulates a contribution
callable. It returns a dict with \code{total\_C} (a callable summing the diagrams
serially) and \code{total\_C\_batch} (a callable that fans the $\tau$-grid out
over a fork-based process pool, or degrades to serial under the notebook fork
guard --- Section~\ref{sec:co-parallel}). The orchestrator stores the per-$\ell$
callable, the full per-$\ell$ result dict, and --- for $k\in\{1,2\}$ --- the
batch-evaluated array over the grid.

\begin{gotcha}
When a loop order has \emph{no} diagrams, the code still inserts a zero callable
and a zero array (\file{pipeline/compute.py:746}). This is necessary, not
defensive padding: some $(k,\ell)$ configurations genuinely have no diagrams at a
given order, and without the placeholder the master \code{total\_C} sum below
would be missing a key.
\end{gotcha}

\begin{note}
Setting \code{use\_grouped\_phase\_j=True} routes this call to
\code{compute\_correction\_td\_grouped} (\file{pipeline/compute.py:761}) instead
of \code{compute\_correction\_td}. The grouped integrator is a prototype that
sums integrands grouped by parent prediagram before quadrature; see
\file{pipeline/\_grouped\_phase\_j.py} for its math and caveats. The default path
is the ungrouped \code{compute\_correction\_td}.
\end{note}

\subsection{Summing across loop orders}

After the per-$\ell$ loop, the master correlator is assembled by summation
(\file{pipeline/compute.py:802}):

\begin{lstlisting}
def total_C(*ext_time_values):
    return sum(complex(fn(*ext_time_values))
               for fn in total_C_by_ell.values())

if tau_points is not None:
    C_tau = sum(C_tau_by_ell[ell] for ell in C_tau_by_ell)
else:
    C_tau = None
\end{lstlisting}

The master callable \code{total\_C} is the additive sum of the per-$\ell$
callables, and the master grid array \code{C\_tau} is the elementwise sum of the
per-$\ell$ grid arrays:
\begin{equation}
  \texttt{total\_C}(\tau) \;=\; \sum_\ell C_\ell(\tau),
  \qquad
  \texttt{C\_tau} \;=\; \sum_\ell \texttt{C\_tau\_by\_ell}[\ell].
  \label{eq:co-sum}
\end{equation}

\begin{note}
\textbf{On the phrase ``set-partition.''}
You may see the framing ``moment / cumulant / central-moment set-partition
assembly'' elsewhere. The orchestrator itself does \emph{not} contain any
moment$\leftrightarrow$cumulant set-partition (M\"obius-inversion) bookkeeping.
\code{compute\_cumulants} returns the \emph{connected} cumulant directly --- the
diagrammatic sum \emph{is} the cumulant. The only assembly this file performs is
the additive per-$\ell$ sum in \eqref{eq:co-sum}. Moment and central-moment
conversions, when needed, live in downstream consumers (the engine-API and
cumulants/moments chapters), not in \file{pipeline/compute.py}.
\end{note}

\section{The adaptive mean-field dictionary}\label{sec:co-mfdict}

\textbf{Code:} \file{pipeline/compute.py:812}--\code{:879}.

Before packing the result, the function discovers which model parameters are
mean-field saddles and reads their solved values into a tidy
\code{\{internal\_name: [v\_pop1, ...]\}} dict. It tries three sources in order:
the \code{naming\_convention['mf\_parameters']} list (preferred, written by the
\code{TheoryBuilder} when \code{mean\_field=True} is set on a parameter), the
per-parameter \code{mean\_field} flag, and finally the classic
\code{nstar}/\code{vstar}/\code{mstar} fallback for hand-written models that
pre-date the declarative API.

For each discovered saddle name it maps the name to the range of population
indices it owns --- via the closure \code{\_saddle\_indices}, which handles
heterogeneous theories that declare \code{indexed\_by=['<pop>']} on each saddle ---
and reads each value out of \code{num\_params}.

\begin{gotcha}
Entries that are \emph{entirely} NaN are dropped (\file{pipeline/compute.py:878}).
A model may declare a saddle symbol it never actually substitutes --- e.g.
\code{mstar} in a non-GTaS model that nonetheless ships an \code{mstar}
declaration. Without this filter those would surface in the result as all-NaN
vectors. The check \code{not all(v != v for v in vals)} is the idiomatic
``not every entry is NaN'' test (recall \code{nan != nan} is \code{True}).
\end{gotcha}

\section{The result dict}\label{sec:co-result}

\textbf{Code:} \file{pipeline/compute.py:881}--\code{:908}.

Everything the run produced is packed into one dictionary. For the temporal path
the keys are:

\begin{center}
\small
\begin{longtable}{@{}p{0.27\textwidth}p{0.20\textwidth}p{0.45\textwidth}@{}}
\toprule
\textbf{key} & \textbf{type} & \textbf{meaning} \\
\midrule
\endhead
\code{total\_C} & callable & $\sum_\ell C_\ell(*\tau)$ \\
\code{total\_C\_by\_ell} & $\{\ell:$ callable$\}$ & per-loop-order callable \\
\code{C\_tau} & ndarray or \code{None} & total on the $\tau$-grid ($k\in\{1,2\}$; \code{None} for $k\ge3$) \\
\code{C\_tau\_by\_ell} & $\{\ell:$ ndarray$\}$ & per-loop-order grid; \code{C\_tau} is their sum \\
\code{tau\_grid} & ndarray & the $\tau$ values \\
\code{mf\_values} & $\{$name: $[v_1,\dots]\}$ & adaptive saddle values \\
\code{mf} & \code{MeanField} & natural-name saddle accessor \\
\code{params} & \code{Parameters} & natural-name parameter accessor \\
\code{num\_params} & $\{$\code{SR}: float$\}$ & numerical substitution dict \\
\code{propagator} & dict & the \code{prop} dict (poles/residues filled) \\
\code{diagrams} & list & the \code{diagram\_records} \\
\code{kernel\_groups} & list & per-diagram Phase-J input format \\
\code{phase\_j\_by\_ell} & $\{\ell:$ dict$\}$ & per-$\ell$ Phase-J result dicts \\
\code{phase\_walls} & $\{$label: sec$\}$ or \code{None} & per-phase wall-time \\
\code{config} & dict & echoed args (internal + user external fields, taylor order, \dots) \\
\bottomrule
\end{longtable}
\end{center}

The two natural-name accessors deserve a sentence. \code{mf} is a
\code{MeanField} object (\file{pipeline/access.py}): \code{mf['v', 1]} returns
$v^{*}_1$ and \code{mf['n']} returns the whole \code{nstar} vector.
\code{params} is a \code{Parameters} object: \code{params['E', 1]} and
\code{params['w', 1, 2]} index into vector and matrix parameters with 1-based
indexing. These spare downstream code from ever touching the raw symbolic dict.

\subsection{Two early-return shapes}\label{sec:co-datastructures}

The spatial paths each return a \emph{different} dict, earlier in the function.
The spatial $k=2$ early return (\file{pipeline/compute.py:584}) adds
\code{C\_tau\_x} (the full $(n_\tau, n_x)$ real-space correlator),
\code{C\_tau\_x\_by\_order} (cumulative per-loop-order grids),
\code{spatial\_grid}, and \code{spatial\_info}; its \code{total\_C} is the spatial
closure (1 arg $\to$ $C(\tau)$ at $x=0$, 2 args $\to$ $C(x,\tau)$) and its
\code{total\_C\_by\_ell} is \code{\{0: total\_C\}}. The spatial $k\ge3$ early
return (\file{pipeline/compute.py:463}) is keyed instead by \code{C\_kpoint} (an
\code{(n\_pts,)} array), \code{C\_kpoint\_by\_ell}, \code{points} (the echoed
\code{spatial\_points}), \code{spatial\_info}, \code{k}, \code{max\_ell}, and
\code{mf}.

For reference, the \code{prop} dict that phases [2] and [4] build and fill has the
following keys (the propagator chapter is the authority):

\begin{center}
\small
\begin{longtable}{@{}p{0.18\textwidth}p{0.30\textwidth}p{0.44\textwidth}@{}}
\toprule
\textbf{key} & \textbf{type} & \textbf{meaning} \\
\midrule
\endhead
\code{K\_ker}   & Sage \code{matrix(SR)} & bilinear kernel in time ($\delta$, $\delta'$) \\
\code{K\_ft}    & Sage \code{matrix(SR)} & Fourier image of \code{K\_ker} (rational in $\omega$) \\
\code{G\_ft}    & Sage matrix or \code{None} & propagator \code{K\_ft$^{-1}$}; \code{None} for rich matrices \\
\code{adj\_ft}  & Sage matrix or \code{None} & adjugate $=$ \code{G\_ft}$\cdot$det \\
\code{D\_omega} & SR or \code{None} & $\det(\text{\code{K\_ft}})$ \\
\code{D\_delta} & Sage \code{matrix(SR)} & instantaneous part $\lim_{\omega\to\infty}$\code{G\_ft} \\
\code{t\_var}, \code{omega} & SR vars & the time and frequency symbols \\
\code{nf}       & int & number of field components \\
\code{pole\_vals} & list[complex] & retarded poles, filled by phase [4] \\
\code{C\_mats}  & list[Sage \code{matrix(CDF)}] & residue matrices, filled by phase [4] \\
\bottomrule
\end{longtable}
\end{center}

\subsection{DAE extras, and the save}

\textbf{Code:} \file{pipeline/compute.py:910}--\code{:944}. When the DAE
equation-based solver ran in phase [3], the result picks up multi-root extras:
\code{mf\_all\_roots} (every distinct root with a stability annotation),
\code{mf\_stable\_roots}, \code{mf\_unstable\_roots}, \code{mf\_index\_used},
\code{mf\_state\_var\_order}, and \code{mf\_stability}.

\begin{gotcha}
\code{fixed\_point\_index} indexes the \emph{stable} subset only
(\file{pipeline/compute.py:910} comment). \code{mf\_stable\_roots} is that subset
in sort order; unstable roots are surfaced via \code{mf\_unstable\_roots} for
inspection but are \emph{not} selectable as expansion points.
\end{gotcha}

Finally, if \code{output\_npz} or \code{output\_csv} was given, the function calls
\code{save\_npz} / \code{save\_csv} from \file{pipeline/save.py} (that module owns
the on-disk schema --- per-loop-order curves \code{C\_tree}, \code{C\_1\_loop},
\dots, the adaptive mean-field arrays, the parameter keys).

\section{Parallelism: fork versus thread, and the macOS hazard}\label{sec:co-parallel}

\Daedalus{} has two parallelizable stages on the temporal path --- per-prediagram
type assignment in phase [5] and per-$\tau$ Phase~J evaluation in phase [7] ---
and both are governed by the single \code{parallel} flag. There is also a
separate \code{spatial\_parallel} flag for the spatial loop integrator. The
distinction between them is not cosmetic; it is a hard-won safety boundary.

\begin{defn}
\textbf{Fork versus thread parallelism.}
\term{Fork} parallelism copies the entire parent process (it is fast to start,
but on macOS it crashes a Jupyter kernel that has already initialised Cocoa, BLAS,
or matplotlib). \term{Thread} parallelism shares one process's memory (it is safe,
and numpy releases the GIL during its heavy linear algebra so threads genuinely
run in parallel). The temporal stages \emph{fork}; the spatial stages
\emph{thread}.
\end{defn}

\begin{gotcha}
\textbf{Spatial parallelism is thread-based, never fork-based.} The default
\code{spatial\_parallel=True} routes the loop integrator through a
\code{ThreadPoolExecutor}. Forking after Cocoa/BLAS/matplotlib initialisation
inside a macOS Jupyter kernel \emph{hard-crashes the kernel and the OS}, and
setting \code{OBJC\_DISABLE\_INITIALIZE\_FORK\_SAFETY=YES} makes it \emph{worse},
not better. The temporal \code{parallel} flag \emph{is} fork-based, but it is
guarded: \code{fork\_unsafe\_in\_notebook} (\file{msrjd/fork\_safety.py:27}) trips
only on \code{darwin} \emph{and} a live ZMQ/Jupyter kernel \emph{and} the
\code{fork} start method, degrading that single case to serial with a one-time
warning. Under pytest, on Linux, in a terminal, or from a script, fork is kept.
This is why the example notebooks in Chapter~\ref{ch:quickstart} pass
\code{parallel=False} explicitly: it is the safe default in a notebook, and the
tiny example models are fast enough serially.
\end{gotcha}

\section{Performance: the phase-wall summary}\label{sec:co-perf}

\textbf{Code:} \file{pipeline/compute.py:946}--\code{:955}. When verbose, the
function closes with a one-line \code{Done.} and, if it tracked timings, a
per-phase breakdown:

\begin{lstlisting}
if verbose:
    print(f'\nDone.  k={k}, max_ell={max_ell}, '
          f'{len(all_unique)} unique diagrams, {len(tau_grid)} tau points.')
    if phase_walls:
        total = sum(phase_walls.values())
        print(f'\n  Phase wall summary (Sigma = {total:.1f}s):')
        for label, dt in phase_walls.items():
            pct = 100.0 * dt / total if total > 0 else 0.0
            print(f'    {label:12s}  {dt:6.2f}s  ({pct:4.1f}%)')
\end{lstlisting}

This reads back the \code{phase\_walls} dict the \code{\_phase\_time} closure has
been filling all along, prints each phase's seconds and its percentage of the
total, and returns the \code{result}. In nearly every nontrivial run the phase
that dominates the clock is \code{phase\_j} (phase [7]) --- the loop integration
--- which is why so much of this manual is devoted to it. The summary is your
first diagnostic when a run is slow: it tells you immediately whether you are
paying for the symbolic expand, the pole-finder, the diagram enumeration, or the
integration.

\begin{gotcha}
The phase-wall summary is opt-in: \code{phase\_walls} is \code{None} when
\code{verbose=False} (\file{pipeline/compute.py:296}), so the timing dict in the
result is only populated in a verbose run. A quiet run returns
\code{result['phase\_walls'] = None}.
\end{gotcha}

\section{The data flow, threaded between phases}\label{sec:co-dataflow}

It is worth seeing the whole thread of data in one diagram, because the
orchestrator's real subject is \emph{what each phase hands the next}:

\begin{lstlisting}[style=console]
  model --+--------------------------------------------------------------+
          v                                                              |
  [1] FieldTheory(model, taylor_order).expand() -> ft (_ns, _n_tilde,    |
          |                                            sectors, vtypes)   |
          v                                                              |
  [2] build_propagator(ft, model) -> prop {K_ft, G_ft, D_delta, w, nf}   |
          |                                                              |
          v                                                              |
  [3] solve_mean_field(ft, model, fundamental) -> num_params {SR:float}  |
          |                                                              |
          v  (spatial? -> early return via the momentum integrator)      |
  [4] compute_poles_and_residues(prop, num_params) -> prop.pole_vals,    |
          |                                              prop.C_mats      |
          v                                                              |
  [5] enumerate_unique_diagrams(ft, model, k, ell, ext, G_ft,           |
          |   resp/phys_idx, vtypes, stypes) -> unique_by_ell            |
          v                                                              |
  [6] classify_coefficient_factors(td, time_dep, noise) ->              |
          |   diagram_records [{td, classify, prefactor, mult, ell}]     |
          v                                                              |
  [7] compute_correction_td(td_list[ell], prefactors[ell],             |
          |   propagator_data, k, external_fields, num_params)           |
          |   -> td_result {total_C, total_C_batch, ...}                 |
          v                                                              |
      total_C = sum_ell ;  C_tau = sum_ell  ->  result dict -------------+
\end{lstlisting}

\begin{example}[A Hawkes-like $k=2$ tree run]
Take \code{external\_fields = [('n',1),('n',2)]}. At validation
(\file{pipeline/compute.py:288}) it is normalized to \code{[('dn',1),('dn',2)]}.
With \code{max\_ell = 0} the Taylor budget \eqref{eq:co-taylor} is
$\max(2,2)=2$. Phases [1]--[6] run; phase [5] finds the tree-level diagrams; phase
[7] sets \code{tau\_points = [(0.0, t) for t in tau\_grid]} (leaf 0 pinned to 0,
leaf 1 swept). Each $\tau$-point is passed as \code{total\_C(0.0, t)}, producing
\code{C\_tau\_by\_ell[0]} as a complex array. With no loops, \code{C\_tau} equals
\code{C\_tau\_by\_ell[0]} exactly. The notebook then plots \code{tau\_grid}
against \code{C\_tau.real} --- the autocorrelation you saw in
Chapter~\ref{ch:quickstart}.
\end{example}

\section{The two helper functions}\label{sec:co-helpers}

The file defines two small helpers above the entry point, both pure formatting
with no effect on the computation.

\code{\_trunc(s, maxlen=200)} (\file{pipeline/compute.py:45}) returns
\code{str(s)} truncated to \code{maxlen} characters with a trailing \code{'...'}.
It keeps long symbolic coefficients readable in the verbose printout.

\code{\_print\_action\_sectors(ft)} (\file{pipeline/compute.py:50}) pretty-prints
the three action sectors after expansion: the mean-field sectors $(0,0)$,
$(1,0)$, $(0,1)$ (read from \code{ft.\_mf\_sector\_raw}), the free $(1,1)$
bilinear kernel (\code{ft.free\_action()}), and every interaction sector of total
degree $\ge2$ excluding $(1,1)$ (iterating \code{ft.sectors()}). Its inner helper
\code{\_fmt\_poly} walks a Sage polynomial's \code{.dict()} --- the
exponent-vector $\to$ coefficient map --- and reconstructs each monomial string
from the ring generator names. This is the ``here is your action'' diagnostic; it
prints, and nothing else.

\section{Gotchas and failure modes}\label{sec:co-gotchas}

We have flagged the local gotchas inline; here is the consolidated list of the
ways a \code{compute\_cumulants} run can surprise you.

\begin{itemize}
  \item \textbf{\code{num\_params} must be fully numeric.} If the mean-field phase
        leaves \emph{any} kernel symbol or saddle value symbolic, the pole-finder
        silently returns zero roots and the whole correlator collapses to $0$ ---
        a subtle bug where, for instance, a missing \code{phi0\_i} (resolved from
        \code{mf\_bg\_conditions}, \file{pipeline/\_mean\_field.py:272}) zeroes the
        result. The pole-finder warns when it finds no roots
        (\file{pipeline/\_propagator.py:976}), so a sudden flat-zero correlator is
        the symptom to watch.
  \item \textbf{\code{prop} is mutated in place by phase [4].} The poles and
        residues appear \emph{inside} the \code{prop} object, not as a return
        value (Section~\ref{sec:co-poles}).
  \item \textbf{Spatial models must be given the right spatial input.}
        \code{spatial\_grid} for $k=2$, \code{spatial\_points} for $k\ge3$; the
        wrong one (or neither) raises a clear error rather than a cryptic Sage one
        (Section~\ref{sec:co-spatial}).
  \item \textbf{Stale caches.} The diagram cache version suffix
        (\code{...\_v3\_...}) forces a rebuild after the dedup-signature fixes;
        the expand and propagator caches are separate
        (Section~\ref{sec:co-diagrams}).
  \item \textbf{The Taylor floor of $2$.} Probing higher-order vertices for
        cache-invalidation testing needs an explicit \code{taylor\_order} override
        (Section~\ref{sec:co-args}).
  \item \textbf{The grouped Phase~J is a prototype.}
        \code{use\_grouped\_phase\_j=True} routes to an experimental integrator;
        the default and trusted path is \code{compute\_correction\_td}
        (Section~\ref{sec:co-phasej}).
\end{itemize}

\section{What you now understand}\label{sec:co-summary}

\code{compute\_cumulants} is the one function that turns a model declaration plus
a parameter point into a numerical correlation function, and it does so by
sequencing seven phases: expand the action, build the propagator, solve the
saddle, find the poles, enumerate the diagrams, classify their prefactors, and
integrate them in the time domain --- with a hard fork between the temporal and
spatial routes after the saddle. It performs almost no arithmetic itself; its
real content is the \emph{contracts} between phases and the \emph{glue} that
threads \code{ft}, \code{prop}, \code{num\_params}, and \code{diagram\_records}
from one to the next, the additive per-$\ell$ assembly that produces the
loop-resolved correlator, and the disciplined parallelism that keeps a macOS
notebook from crashing. Every phase named here has its own chapter; this one is
the map that shows how they connect. The next chapters move outward from the
engine to the user surface --- the \code{daedalus} front desk that you actually
type at --- and then to the cumulant/moment conversions and the simulators that
validate everything the engine produces.

% ---------------------------------------------------------------------
%  Chapter 18 : The daedalus Engine API --- the user-facing front desk
% ---------------------------------------------------------------------
\chapter{The \texttt{daedalus} Engine API}\label{ch:engine-api}

The previous chapters opened up the engine itself --- the symbolic algebra, the
diagram enumeration, the loop integrals. This chapter is about the thin layer of
Python that sits \emph{in front} of all that and is the only thing a notebook
ever touches directly: the module \file{notebooks/daedalus.py}, imported
everywhere as \code{dd}. You met its four verbs already in
Chapter~\ref{ch:quickstart} --- \code{load\_theory}, \code{Config}, \code{run},
\code{plot\_cumulant} --- and watched them drive a complete calculation. Here we
slow down and document the module exhaustively: every function, every
\code{Config} field, the exact rules by which a friendly configuration is turned
into the engine's many-keyword call, the way the heterogeneous result is reshaped
for a plot, and the handful of sharp edges worth knowing about.

The single most important fact about this module sets the tone for the whole
chapter, so we state it up front and return to it at the end: \emph{\code{daedalus}
does almost no physics.} It imports exactly two non-standard libraries --- NumPy
and Matplotlib --- and everything genuinely hard happens behind one line,
\code{from pipeline import compute\_cumulants}. SageMath, the \code{nauty} graph
generator, \code{sympy}, \code{numba}, \code{networkx}: none of them are imported
here. They live entirely inside \code{compute\_cumulants}. So this chapter is, in
effect, a careful reading of a \emph{translation and dispatch} layer, not of the
machinery it drives.

\section{What this module is, in one breath}

A physicist in a notebook does not want to remember the exact signature of the
heavy engine function, nor how to reshape its output for a plot. They want to
write three or four lines. The module's own opening docstring
(\file{notebooks/daedalus.py:8-12}) shows the intended shape:

\begin{lstlisting}
import daedalus as dd
model, mod = dd.load_theory('kpz_1d')          # from theories/*.theory.py
cfg = dd.Config(k=2, max_ell=1, spatial_grid=(-6, 6, 49))
res = dd.run(model, cfg, mod)                  # k/ell/Dyson all here
dd.plot_cumulant(res, cfg, model)              # adaptable, auto-dispatched
\end{lstlisting}

The design intent, quoted directly from that docstring, is to
``centralise the \textbf{load $\to$ run $\to$ plot} flow so every demo notebook
is thin and uniform regardless of its group (temporal / spatial $\times$ single /
multi-field).'' The word \term{group} recurs throughout the module: it is the
informal label for which \emph{shape} of problem you are solving, because that
shape decides which plot is natural (a $C(\tau)$ curve, a $C(x,\tau)$ surface, a
$\kappa_3(\tau_1,\tau_2)$ heatmap, a bar chart of equal-time scalars). The whole
job of \code{daedalus} is to hide the group behind one uniform set of verbs.

So the module gives you three verbs and a noun:

\begin{description}
  \item[\code{load\_theory(name)}] read a \file{theories/<name>.theory.py} file
        and return the pair \code{(model, module)};
  \item[\code{Config(...)}] a \code{@dataclass} holding \emph{every} choice a
        notebook can make;
  \item[\code{run(model, cfg, module)}] resolve the config against the theory's
        defaults, call \code{compute\_cumulants}, post-process the result, and
        staple the config and model back on;
  \item[\code{plot\_cumulant(result, cfg, model, sim=None)}] auto-dispatch to the
        plotter natural to the group.
\end{description}

Around those four are a set of introspection helpers
(\code{describe\_model}, \code{field\_names}, \code{is\_spatial},
\code{spatial\_dim}, \code{is\_multifield}, \code{summary},
\code{list\_theories}) and an output-conversion layer that turns connected
cumulants into full or central moments (\code{Config.output} together with
\code{\_assemble\_moment\_temporal}). The rest of this chapter walks every one of
them.

\subsection{Where it sits in the pipeline}

It helps to keep one picture in mind. A \term{theory file} (covered in
Chapters~\ref{ch:theory-spec} and~\ref{ch:theory-files}) is the declarative
description of a model; a notebook's \code{Config} is the set of run choices.
\code{daedalus} marries the two and calls the engine:

\begin{center}
\begin{tikzpicture}[
    node distance=6mm,
    box/.style={draw, rounded corners, align=center, inner sep=4pt,
                font=\small, fill=codebg},
    eng/.style={draw, rounded corners, align=center, inner sep=4pt,
                font=\small\bfseries, fill=notebg},
    >={Stealth[]}]
  \node[box] (theory) {\file{theories/<name>.theory.py}\\[2pt]
    \footnotesize\code{build()}, \code{DEFAULT\_FUNDAMENTAL}, \code{METADATA}};
  \node[box, below=of theory] (pair) {\code{(model: dict, module)}\quad+\quad
    \code{dd.Config}};
  \node[eng, below=of pair] (engine) {\code{pipeline.compute\_cumulants(\ldots)}\\[2pt]
    \footnotesize the \msrjd{} engine --- Phases 1--7};
  \node[box, below=of engine] (post) {\code{run()} post-processes:\\[2pt]
    \footnotesize $k\!\ge\!3$ slices/grid \,$\cdot$\, moment assembly \,$\cdot$\,
    \code{\_cfg}/\code{\_model}/\code{\_resolved}};
  \node[box, below=of post] (plot) {\code{plot\_cumulant()} $\to$ one of seven
    plotters $\to$ Matplotlib \code{Figure}};
  \draw[->] (theory) -- node[right, font=\scriptsize] {\code{load\_theory()}} (pair);
  \draw[->] (pair)   -- node[right, font=\scriptsize] {\code{run()} resolves $k$, $\ell$, legs, grids, Dyson} (engine);
  \draw[->] (engine) -- node[right, font=\scriptsize] {result dict} (post);
  \draw[->] (post)   -- (plot);
\end{tikzpicture}
\end{center}

\begin{itemize}
  \item \textbf{What feeds it:} a theory file plus a \code{Config}.
  \item \textbf{What it calls:} \code{pipeline.compute\_cumulants}, the full
        \msrjd{} diagrammatic engine. Everything hard --- the field-theory Taylor
        expansion, the propagator, the mean-field saddle solve, diagram
        enumeration, the Symanzik or Phase-J integration --- happens \emph{inside}
        that one call.
  \item \textbf{What consumes its output:} the demo notebooks, which take the
        result \code{res} and a matched simulator dict \code{sim} and hand both to
        \code{plot\_cumulant}; they also read the result directly
        (\code{res['\_resolved']}, \code{res['mf']}, \code{res['spatial\_info']},
        and so on).
\end{itemize}

\section{The math the knobs control}

You are assumed to know \msrjd{} field theory already, so this section is short:
its only purpose is to tie each \emph{code} knob to the \emph{physics} object it
controls. All math here is elementary; the detailed treatment is in the chapters
on the field-theory core and on Phase~J.

\subsection{Connected cumulants are the native output}

The engine computes \term{connected $k$-point cumulants} of the physical
field(s). For a single field $\phi(x,t)$ the order-$k$ connected cumulant is
\begin{equation}
  \kappa_k(x_1,t_1;\,\ldots\,;x_k,t_k)
  \;=\; \avg{\phi(x_1,t_1)\cdots\phi(x_k,t_k)}_c ,
\end{equation}
the connected (cluster) part of the $k$-point correlation. You ask for a chosen
$k$ (the number of external legs) and a chosen \code{max\_ell} (the loop order).
The diagrammatic expansion is organised by \term{loop order} $\ell$,
\begin{equation}
  \kappa_k \;=\; \underbrace{\kappa_k^{(0)}}_{\text{tree}}
              + \underbrace{\kappa_k^{(1)}}_{\text{1-loop}}
              + \underbrace{\kappa_k^{(2)}}_{\text{2-loop}}
              + \cdots,
\end{equation}
and \code{Config.max\_ell}~$=L$ truncates the sum at $\ell\le L$. The engine
returns the \emph{per-order} pieces as \code{*\_by\_ell} dictionaries
\code{\{ell: array\_or\_callable\}}, and the running cumulative sum
$\sum_{\ell\le L}$ as the ``total''. The module's
\code{cumulative\_curves} (\file{notebooks/daedalus.py:713}) and
\code{\_order\_label} (\file{notebooks/daedalus.py:707}) exist precisely to render
this $\ell$-progression --- the labels \code{'tree'}, \code{'tree+1loop'}, and so
on.

\subsection{The \texorpdfstring{$k=2$}{k=2} case: the correlation function}

For $k=2$ the connected cumulant is just the (connected) two-point function. In
the \term{temporal} (time-only) case the engine returns it as a curve over a lag
grid $\tau$,
\begin{equation}
  C(\tau) \;=\; \avg{\phi(0)\,\phi(\tau)}_c
  \qquad\text{(key \code{'C\_tau'} over \code{'tau\_grid'}),}
\end{equation}
while in the \term{spatial} case (a field with \code{spatial\_dim}~$\ge 1$) it
returns a real-space correlator over separations $x$ and lags $\tau$,
\begin{equation}
  C(x,\tau) \;=\; \avg{\phi(0,0)\,\phi(x,\tau)}_c
  \qquad\text{(key \code{'C\_tau\_x'}, shape $(n_\tau, n_x)$).}
\end{equation}
The equal-time slice $C(x,0)$ is the headline spatial plot.

\subsection{The \texorpdfstring{$k\ge 3$}{k>=3} case: a function of \texorpdfstring{$k-1$}{k-1} time differences}

By time-translation invariance the connected $k$-point cumulant depends only on
the $k-1$ time differences relative to one anchored leg,
\begin{equation}
  \tau_j \;=\; t_j - t_0, \qquad j = 1,\ldots,k-1 \quad\text{(leg 0 pinned at $\tau=0$).}
\end{equation}
So $\kappa_k$ is a function $\kappa_k(\tau_1,\ldots,\tau_{k-1})$: for $k=2$ it is
the one variable $C(\tau)$; for $k=3$ it is a 2-D surface $\kappa_3(\tau_1,\tau_2)$;
for $k\ge 4$ it is a $(k-1)$-dimensional tensor that cannot be drawn directly.
\code{compute\_cumulants} does \emph{not} materialise this whole tensor for
$k\ge 3$. Instead it returns \code{C\_tau = None} together with a \emph{callable}
\code{total\_C(*tau)} and per-loop callables \code{total\_C\_by\_ell = \{ell:
callable\}}. It is \code{run} that synthesises two human-readable views from those
callables (Section~\ref{eng:kge3}): the axis-parallel \emph{slices} (the default)
and the optional full \emph{grid}. We return to exactly how it does this once we
reach \code{run}.

\subsection{Cumulants to moments: the set-partition expansion}\label{sec:math-moments}

The engine produces \emph{connected} cumulants $\kappa$. The full (raw) and
central \term{moments} are reconstructed by the \term{set-partition} (cluster)
expansion: the raw $k$-point moment is the sum over all set partitions $\pi$ of
$\{1,\ldots,k\}$ of products of cumulants over the blocks,
\begin{equation}
  M_k(x_1,\ldots,x_k) \;=\; \sum_{\pi}\ \prod_{B\in\pi} \kappa(B).
  \label{eq:setpart}
\end{equation}
This identity is stated verbatim in the \code{Config.output} comment
(\file{notebooks/daedalus.py:261-266}) and in the
\code{\_assemble\_moment\_temporal} docstring
(\file{notebooks/daedalus.py:390-408}). Two physics subtleties the code handles
explicitly are worth flagging now, because they shape the implementation:

\begin{enumerate}
  \item \textbf{A single shared loop budget.} A naive ``product of fully-dressed
        cumulants'' would smuggle in spurious cross-terms --- a
        $\text{1-loop}\cdot\text{1-loop}$ product is really a 2-loop object. The
        code instead enforces one \emph{shared} loop budget $L=\code{max\_ell}$
        across all multi-element blocks of a partition,
        \begin{equation}
          M_k \;=\; \sum_\pi \;
            \sum_{\substack{(\ell_B):\ \sum_B \ell_B \le L}}\;
            \prod_{B\in\pi} \kappa(B)^{(\ell_B)},
          \label{eq:sharedbudget}
        \end{equation}
        so every surviving term has \emph{total} loop order $\le L$. The two
        definitions agree only at tree level ($L=0$) and diverge at $L\ge 1$ for
        any multi-block partition.
  \item \textbf{Singletons are the tree mean.} A singleton block $\{i\}$
        contributes the field mean $\avg{\phi}$, which at tree level is the
        mean-field saddle $\phi^{*}$ (read by \code{\_external\_mean},
        \file{notebooks/daedalus.py:373}). For a \emph{central} moment all
        singletons are dropped (only no-singleton partitions survive). For a
        \emph{raw} moment singletons contribute $\mu^{\#\text{singletons}}$, and
        if the mean is zero those terms die. The 1-loop tadpole shift of the mean
        is \emph{not} yet folded in --- a documented refinement
        (\file{notebooks/daedalus.py:376-377}, \code{:404-405}).
\end{enumerate}

\code{Config.output}~$\in$~\code{\{'cumulant', 'moment', 'central\_moment'\}}
selects this. Producing a moment of order $k$ costs up to $k-1$ \emph{extra}
backend runs (one per cumulant order $2..k$), because each block size needs its
own \code{compute\_cumulants} call.

\subsection{The mean-field saddle, ``fundamental'', and Dyson dressing}

Two more knobs round out the physics surface. First, the pipeline solves a
\term{mean-field saddle} --- the deterministic fixed point $\phi^{*}$ about which
the fluctuation expansion is taken --- before any diagrams are evaluated. The
numeric values of the model's parameters ($\mu, D, \lambda, \ldots$) are passed as
the \code{fundamental} dict (the engine's name), which is the same thing as
\code{Config.parameters} (the friendly name). Saddle quantities (names ending in
\code{star}) are \emph{solved}, never supplied. Some theories --- a double well,
$\mu<0$ --- have several stable roots, and \code{Config.fixed\_point\_index},
\code{mf\_dae\_n\_starts}, and \code{mf\_dae\_seed\_box} choose which root the
expansion centres on.

Second, for coupled multi-field spatial theories whose fields have \emph{unequal}
bare diffusivities, the tree-level propagator is dressed by resumming a geometric
self-energy series --- the \term{Dyson series} --- to a finite order.
\code{Config.dyson\_order} overrides the model's built-in Dyson policy at run
time, optionally with a reference diffusivity $D_0$ around which the dressing is
organised. The detailed treatment of both is in Chapters~\ref{ch:mean-field}
and~\ref{ch:spatial-coupled}; here they are just \code{Config} fields that
\code{run} forwards.

\section{Loading a theory and finding the repository}

The discovery and loading helpers are the first block of the module
(\file{notebooks/daedalus.py:40-80}). They are small, but understanding them
explains why the demo notebooks run unchanged from any subdirectory.

\subsection{\texttt{repo\_root} --- depth-robust path discovery}

\code{repo\_root() -> str} (\file{notebooks/daedalus.py:40}) returns the absolute
path of the repository root. It starts from the directory of \emph{this file} if
\code{\_\_file\_\_} is defined, else from the current working directory (the
fallback covers code pasted into a bare notebook cell,
\file{notebooks/daedalus.py:44-45}), then walks \emph{up} the directory tree until
it finds a directory containing a \code{pipeline/} subfolder:

\begin{lstlisting}
root = here
while root != os.path.dirname(root):
    if os.path.isdir(os.path.join(root, 'pipeline')):
        return root
    root = os.path.dirname(root)
return here
\end{lstlisting}

If the walk reaches the filesystem root without finding \code{pipeline/}, it
returns the starting directory as a last resort. The same upward-walk idiom
appears in the example notebooks' own setup cell (you saw it in
Chapter~\ref{ch:quickstart}), and it is what makes them robust to being launched
from the repository top, from \code{notebooks/}, or from
\code{notebooks/examples/}. At import time the module records the consequences as
module-level constants (\file{notebooks/daedalus.py:54-58}): \code{REPO\_ROOT =
repo\_root()}, \code{THEORIES\_DIR = REPO\_ROOT/theories}, and \code{REPO\_ROOT}
is prepended to \code{sys.path} so that \code{import pipeline} resolves no matter
where you started.

\subsection{\texttt{list\_theories} --- the menu of loadable models}

\code{list\_theories() -> list[str]} (\file{notebooks/daedalus.py:61}) lists
\code{THEORIES\_DIR}, keeps the files ending in \code{.theory.py}, strips that
suffix, and sorts. So \file{kpz\_1d.theory.py} becomes the loadable name
\code{'kpz\_1d'}. It serves both as the menu of theories and as the body of the
error message in \code{load\_theory}.

\subsection{\texttt{load\_theory} --- dynamic import of a \texttt{.theory.py} file}

\code{load\_theory(name)} (\file{notebooks/daedalus.py:68}) takes the bare theory
name (no \code{.theory.py}) and returns the pair \code{(model, module)}, where
\code{model = module.build()} is the built model dictionary and \code{module} is
the live imported theory module. Step by step it builds the path
\code{THEORIES\_DIR/<name>.theory.py}; raises \code{FileNotFoundError} (listing
the available theories) if it does not exist; then dynamically imports the file.

The import mechanism deserves a paragraph because it is the one piece of standard
library here that is not obvious.

\begin{defn}
\textbf{\code{importlib.util} dynamic import.} \code{importlib} is Python's
programmatic import system. Its \code{util} submodule can import a module
\emph{from an arbitrary file path} --- not just by name from \code{sys.path}. The
canonical three-step idiom is: describe the module with
\code{spec\_from\_file\_location}, create an empty module object with
\code{module\_from\_spec}, then actually \emph{run} the file's top-level code with
\code{spec.loader.exec\_module}.
\end{defn}

That is exactly what the function does (\file{notebooks/daedalus.py:77-80}):

\begin{lstlisting}
spec = importlib.util.spec_from_file_location(f'theories.{name}', path)
mod = importlib.util.module_from_spec(spec)
spec.loader.exec_module(mod)
return mod.build(), mod
\end{lstlisting}

Running the file's top-level code is what defines \code{build},
\code{DEFAULT\_FUNDAMENTAL}, and \code{METADATA}; the function then returns the
built model dict and the live module. There are \emph{two} return values for a
real reason: the model dict drives the run, but the run \emph{recommendations}
(\code{k\_default}, \code{ell\_default}, \code{tau\_max},
\code{recommended\_external\_fields}, $\ldots$) and the numeric parameter
defaults (\code{DEFAULT\_FUNDAMENTAL}) live on the \emph{module}, not in the model
dict. You need both to run a theory the way its author intended.

\begin{note}
To see every theory you can load, call \code{dd.list\_theories()}. If you mistype
a name, \code{load\_theory} raises a \code{FileNotFoundError} whose message
includes that same list --- so the menu of valid names is always one error away.
\end{note}

\section{Introspecting a model}

Before running, you usually want to confirm the engine is about to compute what
you think it is. The introspection helpers read \emph{only} the model dict (and,
optionally, the module's docstring and \code{METADATA}), so they stay
model-independent.

\subsection{The routing predicates}

Four tiny predicates classify a model into its group:

\begin{description}
  \item[\code{is\_spatial(model) -> bool}] (\file{notebooks/daedalus.py:85}) ---
        \code{True} iff \code{model['spatial']} exists and has a truthy
        \code{dim}. This single predicate routes everything spatial versus
        temporal.
  \item[\code{spatial\_dim(model) -> int}] (\file{notebooks/daedalus.py:90}) ---
        the number of spatial dimensions $d$, or $0$.
  \item[\code{is\_multifield(model) -> bool}] (\file{notebooks/daedalus.py:95}) ---
        \code{True} iff the theory has more than one independent field channel:
        either more than one entry in \code{model['physical\_fields']}, or any
        population with \code{size > 1}.
  \item[\code{field\_names(model) -> list[str]}] (\file{notebooks/daedalus.py:105})
        --- the list \code{[f['name'] for f in model['physical\_fields']]}, used
        to build default external legs and to locate the mean field.
\end{description}

\subsection{\texttt{describe\_model} --- the ``what is this model?'' cell}

\code{describe\_model(model, module=None, show\_doc=True) -> str}
(\file{notebooks/daedalus.py:138}) is the canonical introspection cell. It builds
a boxed, multi-line report and \emph{also prints it} (\file{notebooks/daedalus.py:241}).
Reading the report top to bottom (\file{notebooks/daedalus.py:147-242}) you get,
in order: a header bar with \code{model['name']}; a \emph{Domain} line that is
either a spatial PDE description ($d=\ldots$, boundary, initial condition) or
``temporal ODE (time-only)''; a \emph{Fields} line (natural name, population tag,
$x\in\mathbb{R}^{d}$ annotation, description); response fields if any;
populations (skipping the auto scalar population named \code{'pop'}); the
\emph{Parameters} block split into \emph{numeric} (default-valued) and
\emph{saddle} (names ending \code{star}, ``solved by the pipeline''); any kernels
(non-Markovian temporal convolutions) and transfer functions; the governing
equation(s) in $\text{lhs}=\text{rhs}$ form; a \emph{Suggested run} line from
\code{module.METADATA}; and finally, when \code{show\_doc} is set, the theory's
trimmed physics docstring.

Two small helpers do supporting work. \code{\_fmt\_default(v, maxlen=52)}
(\file{notebooks/daedalus.py:109}) is a \code{repr()} truncated to 52 characters
with a trailing ellipsis --- purely cosmetic, so a long default value does not
blow up the parameter table. And \code{\_doc\_physics(doc)}
(\file{notebooks/daedalus.py:125}) trims a theory's docstring down to its
\emph{physics prose}: it walks the lines, \emph{stops} at any line containing a
\code{\_DOC\_STOP} marker (\code{'status:'}, \code{'VERBATIM'}, the
\code{'Phase N status'} lines, \ldots), and \emph{skips} any line containing a
\code{\_DOC\_DROP} substring (\code{'Generated by'}, \code{'.theory.py'},
\code{'docs/'}, \ldots). The marker lists are module constants at
\file{notebooks/daedalus.py:116} and \code{:121}. The effect is that you see the
model, not its changelog.

\begin{gotcha}
\code{describe\_model} both \emph{prints} the report and \emph{returns} it as a
string (\file{notebooks/daedalus.py:241-242}). In a notebook cell that means you
will see the text twice --- once from the print, once as the cell's echoed return
value --- unless you assign or suppress the return. Harmless, just a touch noisy.
\end{gotcha}

\section{The \texttt{Config} object, field by field}

\code{Config} (\file{notebooks/daedalus.py:247}) is the single most important
object in this module. The \code{@dataclass} decorator means Python
auto-generates the constructor from the field declarations, so you set only what
you care about and inherit theory defaults for everything else. Its own docstring
states the rule plainly: ``leave a field \code{None} to inherit the theory file's
\code{METADATA}~/~\code{DEFAULT\_FUNDAMENTAL} default.''

\begin{defn}
\textbf{\code{@dataclass}.} A standard-library Python decorator
(\code{from dataclasses import dataclass}) that reads a class's annotated field
declarations and auto-generates an \code{\_\_init\_\_} (and \code{\_\_repr\_\_},
equality, \ldots) for you. Every field below becomes a constructor keyword with
the listed default.
\end{defn}

The full field list, in declaration order, follows. The groups mirror the comment
banners in the source.

\begin{longtable}{@{}p{0.27\textwidth} p{0.12\textwidth} p{0.50\textwidth}@{}}
\toprule
\textbf{Field} & \textbf{Default} & \textbf{Meaning (source line)} \\
\midrule
\endfirsthead
\toprule
\textbf{Field} & \textbf{Default} & \textbf{Meaning (source line)} \\
\midrule
\endhead
\multicolumn{3}{@{}l}{\emph{What to compute}}\\[2pt]
\code{k} & \code{None} & Correlator order (number of external legs). If \code{None}, inferred from \code{external\_fields}, else from \code{METADATA['k\_default']} (\code{:255}). \\
\code{max\_ell} & \code{None} & Loop order $L$ (0=tree). \code{None} $\Rightarrow$ \code{METADATA['ell\_default']} (\code{:257}). \\
\code{external\_fields} & \code{None} & The $k$ legs, e.g.\ \code{[('x', 1), ('x', 1)]}; may name a response leg (\code{:258}). \\
\code{parameters} & \code{None} & Numeric parameter overrides by canonical name (\code{:259}). \\
\code{fundamental} & \code{None} & \emph{Deprecated alias} for \code{parameters} (\code{:260}); see \code{\_\_post\_init\_\_}. \\
\code{output} & \code{'cumulant'} & \code{'cumulant'}\,/\,\code{'moment'}\,/\,\code{'central\_moment'} --- see Section~\ref{sec:math-moments} (\code{:261}). \\[4pt]
\multicolumn{3}{@{}l}{\emph{Temporal grid}}\\[2pt]
\code{tau\_max} & \code{None} & Temporal grid extent; \code{None} $\Rightarrow$ METADATA or $10.0$ (\code{:269}). \\
\code{tau\_step} & \code{None} & Temporal grid spacing; \code{None} $\Rightarrow$ METADATA or $0.5$ (\code{:270}). \\[4pt]
\multicolumn{3}{@{}l}{\emph{Spatial}}\\[2pt]
\code{spatial\_grid} & \code{None} & Separations: an array, an \code{(lo, hi, n)} tuple, or \code{None} (\code{:273}). \\
\code{spatial\_points} & \code{None} & \emph{$k\!\ge\!3$ spatial only}: \code{(n\_pts, k-1, 2)} of $(x_j,\tau_j)$ per non-anchor leg (\code{:274}). \\[4pt]
\multicolumn{3}{@{}l}{\emph{Temporal $k\!\ge\!3$ slicing}}\\[2pt]
\code{kpoint\_base\_lags} & \code{None} & Length $k-1$: the fixed $\tau$ for the non-swept legs (default all 0) (\code{:280}). \\
\code{kpoint\_full\_grid} & \code{False} & Compute the full $(k\!-\!1)$-dim tensor instead of slices (\code{:283}). \\[4pt]
\multicolumn{3}{@{}l}{\emph{Dyson dressing}}\\[2pt]
\code{dyson\_order} & \code{None} & \code{None}=leave model policy; \code{int}$\ge 0$=override (\code{:288}). \\
\code{reference\_diffusion} & \code{None} & $D_0$ reference for the Dyson dressing (\code{:289}). \\[4pt]
\multicolumn{3}{@{}l}{\emph{Mean-field root selection}}\\[2pt]
\code{fixed\_point\_index} & \code{None} & Which stable saddle root (0, 1, \ldots) (\code{:296}). \\
\code{mf\_dae\_n\_starts} & \code{None} & Multi-start Newton count for the DAE saddle solve (\code{:297}). \\
\code{mf\_dae\_seed\_box} & \code{None} & \code{\{field: (lo, hi)\}} start range for the saddle solve (\code{:298}). \\[4pt]
\multicolumn{3}{@{}l}{\emph{Execution}}\\[2pt]
\code{parallel} & \code{False} & Execution parallelism flag forwarded to the engine (\code{:301}). \\
\code{verbose} & \code{False} & Print engine progress (\code{:302}). \\[4pt]
\multicolumn{3}{@{}l}{\emph{Plotting options}}\\[2pt]
\code{show\_orders} & \code{'cumulative'} & \code{'cumulative'}\,/\,\code{'incremental'}\,/\,\code{'total'} (\code{:305}). \\
\code{logy} & \code{False} & Log $y$-axis (\code{:306}). \\
\code{components} & \code{None} & Which $(i,j)$/slice to draw; \code{None}=auto (\code{:307}). \emph{Declared but never read --- see Section~\ref{sec:gotchas}.} \\
\code{figsize} & \code{None} & Matplotlib figure size tuple (\code{:308}). \\
\code{title} & \code{None} & Plot title override (\code{:309}). \\
\code{save} & \code{None} & Path to \code{savefig}, or \code{None} (\code{:310}). \\
\bottomrule
\caption{Every field of \code{Config}, in declaration order
  (\file{notebooks/daedalus.py:247-310}).}
\label{tab:config}
\end{longtable}

\subsection{The two \texttt{Config} methods}

\paragraph{\texttt{\_\_post\_init\_\_}} (\file{notebooks/daedalus.py:312-319})
keeps \code{parameters} and \code{fundamental} in sync. \code{parameters} is the
canonical name and \code{fundamental} the deprecated alias; whichever the caller
set is mirrored onto the other, with \code{parameters} winning if somehow both
are given. So an old notebook passing \code{fundamental=\{\ldots\}} and a new one
passing \code{parameters=\{\ldots\}} both work:

\begin{lstlisting}
if self.parameters is None and self.fundamental is not None:
    self.parameters = self.fundamental
elif self.fundamental is None and self.parameters is not None:
    self.fundamental = self.parameters
\end{lstlisting}

\paragraph{\texttt{resolved\_grid}} (\file{notebooks/daedalus.py:321-328})
materialises \code{spatial\_grid} into the array form the engine wants. If it is
\code{None} it returns \code{None}; if it is a 3-tuple \code{(lo, hi, n)} it
returns \code{np.linspace(lo, hi, int(n))}; otherwise it coerces whatever was
given to a float array with \code{np.asarray(g, dtype=float)}. This is the bridge
from the friendly tuple form \code{(-6, 6, 49)} to the 49-point NumPy array the
spatial path consumes.

\section{How \texttt{run} resolves a \texttt{Config}}\label{sec:run}

\code{run(model, cfg, module=None) -> dict} (\file{notebooks/daedalus.py:487}) is
the central orchestration function. It resolves the (mostly \code{None})
\code{Config} against the theory's defaults, builds the keyword dict the engine
wants, calls \code{compute\_cumulants}, post-processes the result, and stamps it.
We walk it in the order it executes.

\subsection{Resolving \texorpdfstring{$k$}{k}, the loop order, and the legs}

The first job is to pin down \emph{what} to compute. The loop order is resolved
first (\file{notebooks/daedalus.py:494-495}): \code{cfg.max\_ell} if set, else
\code{METADATA['ell\_default']}, else \code{0}.

Then $k$, which is the subtle one (\file{notebooks/daedalus.py:499-515}). The
rule is a small decision with one deliberate hard failure:

\begin{itemize}
  \item If \code{cfg.k} is set, $k=\code{cfg.k}$ --- but if \code{external\_fields}
        is \emph{also} explicit and its length disagrees with $k$, \code{run}
        raises a \code{ValueError}. A $k$/legs mismatch is treated as a
        contradiction, not silently patched.
  \item Else if \code{external\_fields} is explicit, $k = \code{len(ext)}$ --- a
        $k$-point correlator has exactly $k$ legs.
  \item Else $k = \code{METADATA['k\_default']}$ (or 2).
\end{itemize}

The contradiction guard is worth quoting, because it is the kind of thing that
saves an hour of confusion (\file{notebooks/daedalus.py:506-511}):

\begin{lstlisting}
if cfg.k is not None:
    k = cfg.k
    if ext_is_explicit and len(ext) != k:
        raise ValueError(
            f"Config mismatch: k={k} but external_fields lists "
            f"{len(ext)} leg(s) {ext}.  A k-point correlator needs exactly "
            f"k legs.  Drop k (it is inferred from external_fields), set "
            f"k={len(ext)}, or pass {k} legs.")
\end{lstlisting}

Next the external legs themselves (\file{notebooks/daedalus.py:517-533}). If you
gave none and a module exists, \code{run} first tries
\code{METADATA['recommended\_external\_fields']}. It then \emph{auto-builds} $k$
copies of the first physical field's leg \code{(f0, 1)} --- but \emph{only when
the caller gave none explicitly}, and only when the candidate is absent, has the
wrong length, or names a field that does not exist (the check uses
\code{field\_names(model)} and a tiny \code{\_nm} helper that pulls the name out
of a tuple or string). An explicit \code{Config.external\_fields} is used
verbatim, because it may legitimately name a response leg; any $k$-mismatch on it
was already caught by the guard above.

\subsection{The layered \texttt{fundamental}}\label{sec:layered}

Numeric parameters are resolved by a three-layer override
(\file{notebooks/daedalus.py:534-542}), each layer winning over the one before:

\begin{enumerate}
  \item \textbf{Model defaults.} \code{fundamental\_from\_model(model)}
        (\file{notebooks/daedalus.py:335}) builds a dict from every non-saddle
        parameter that carries a declared \code{default}. Saddle params
        (\code{*star}) and defaultless params are skipped. This is what lets a
        freshly loaded theory run out of the box even when its
        \code{DEFAULT\_FUNDAMENTAL} is empty. (The function has a canonical alias
        \code{parameters\_from\_model} at \file{notebooks/daedalus.py:352}, kept
        for existing imports.)
  \item \textbf{The module's \code{DEFAULT\_FUNDAMENTAL}}, updated on top.
  \item \textbf{Your \code{cfg.parameters}}, updated last.
\end{enumerate}

\begin{lstlisting}
fundamental = fundamental_from_model(model)
fundamental.update(getattr(module, 'DEFAULT_FUNDAMENTAL', {}) or {}
                   if module else {})
if cfg.parameters:
    fundamental.update(cfg.parameters)
\end{lstlisting}

So a loaded theory runs with sensible numbers, and your \code{Config} can
override any subset of them by name.

\subsection{The Dyson override and the MF-root forwards}

If \code{cfg.dyson\_order} is set \emph{and} the model is spatial, \code{run}
writes the override into the model's spatial policy
(\file{notebooks/daedalus.py:545-550}):

\begin{lstlisting}
model['spatial']['dyson'] = {'mode': 'fixed', 'order': int(cfg.dyson_order)}
if cfg.reference_diffusion is not None:
    model['spatial']['reference_diffusion'] = float(cfg.reference_diffusion)
\end{lstlisting}

\begin{gotcha}
This \emph{mutates the model dict in place}. If you reuse the same \code{model}
object across runs with different \code{dyson\_order}, the last write persists.
Re-\code{load\_theory} for a clean model.
\end{gotcha}

The mean-field root-selection overrides
(\file{notebooks/daedalus.py:555-564}) are forwarded into the keyword dict
\emph{only when set}, so the engine's own defaults (\code{fixed\_point\_index=0},
\code{mf\_dae\_n\_starts=64}, \code{mf\_dae\_seed\_box=None}) are preserved
otherwise. This matters for multi-root theories like the double-well $\mu<0$
regime.

\subsection{Wiring the grid, spatial versus temporal}

With the engine keyword dict \code{kw} assembled (\code{model, k, max\_ell,
fundamental, external\_fields, parallel, verbose}, at
\file{notebooks/daedalus.py:552-553}), \code{run} wires the grid according to the
group (\file{notebooks/daedalus.py:566-584}):

\begin{itemize}
  \item \textbf{Spatial.} If $k\ne 2$ and no \code{spatial\_points} were given,
        raise --- $k\ge 3$ spatial needs explicit events. If $k\ne 2$, set
        \code{kw['spatial\_points']}. Else ($k=2$) materialise the separations
        grid (defaulting to \code{np.linspace(-6, 6, 49)}) and set \code{tau\_max}
        and \code{tau\_step}, which default to $0.0$ and $1.0$ --- i.e.\
        equal-time-centred, one $\tau$ row.
  \item \textbf{Temporal.} \code{tau\_max} and \code{tau\_step} come from the
        \code{Config} or from \code{METADATA}, defaulting to $10.0$ and $0.5$.
\end{itemize}

\begin{gotcha}
The spatial and temporal grid \emph{centrings differ}. A spatial $k=2$ run
defaults to a single equal-time row ($\code{tau\_max}=0.0$), so \code{C\_tau\_x}
has shape $(1, n_x)$ unless you set \code{tau\_max}. A temporal run instead spans
$\tau\in[-10,10]$. Also note that some theories declare a \code{spatial\_grid} in
their \code{METADATA}, but \code{run} \emph{ignores} it and falls back to
\code{np.linspace(-6, 6, 49)}; pass \code{spatial\_grid} explicitly to use the
theory's recommendation.
\end{gotcha}

Finally the engine is called (\file{notebooks/daedalus.py:586}):
\code{res = compute\_cumulants(**kw)}.

\subsection{The result stamps}

After the engine returns (and after the $k\ge 3$ and moment post-processing
described next), \code{run} stamps three convenience keys onto the result
(\file{notebooks/daedalus.py:694-697}) so downstream code never has to thread the
config around separately:

\begin{lstlisting}
res['_cfg'] = cfg
res['_model'] = model
res['_resolved'] = dict(k=k, max_ell=max_ell, external_fields=ext,
                        parameters=fundamental, fundamental=fundamental)
\end{lstlisting}

\code{res['\_resolved']} records the resolved $k$, loop order, legs, and the final
layered parameter dict --- the values the engine actually used. As you saw in
Chapter~\ref{ch:quickstart}, the simulator reads
\code{res['\_resolved']['parameters']} so that it runs exactly the same physics.

\section{The \texorpdfstring{$k\ge 3$}{k>=3} slice and grid synthesis}\label{eng:kge3}

This is the most intricate part of \code{run}, and it is worth its own section.
When the run is temporal, $k\ge 3$, and the engine came back with
\code{C\_tau = None} and a callable \code{total\_C}
(\file{notebooks/daedalus.py:595-672}), \code{run} synthesises the human-readable
views the plotters expect.

First it chooses a $\tau$-grid, capping it to 41 points to bound the number of
1-D evaluations (\file{notebooks/daedalus.py:597-603}). Then it resolves the
\term{base} --- the fixed point through which the axis-parallel slices pass ---
from \code{cfg.kpoint\_base\_lags}, validating that its length is $k-1$, and
defaulting to all zeros (\file{notebooks/daedalus.py:608-616}). Two local helpers
do the geometry: \code{\_args(vals)} builds the $k$-length time vector with leg 0
pinned at 0, and \code{\_slice(fn, j)} sweeps leg $j$ over the grid with the other
legs held at the base (\file{notebooks/daedalus.py:618-630}).

With those it builds the $k-1$ axis-parallel slices and their per-$\ell$ versions
(\file{notebooks/daedalus.py:632-646}), stores them under \code{C\_tau\_slices}
and \code{C\_tau\_slices\_by\_ell}, and --- crucially --- sets the canonical
single slice \code{C\_tau = slices[1]} and \code{C\_tau\_by\_ell =
slices\_by\_ell[1]}. That last move is what lets a $k\ge 3$ run \emph{plot and
sim-compare exactly like a $k=2$ run}: a curve over \code{tau\_grid}. It also
records \code{\_kpoint\_base} and a human-readable \code{\_kpoint\_slice}
description.

\begin{note}
For a single symmetric field, all $k-1$ slices coincide --- which is a free,
visible check of the cumulant's permutation symmetry
(\file{notebooks/daedalus.py:856-867}). When the external legs name different
fields, the slices genuinely differ.
\end{note}

If \code{cfg.kpoint\_full\_grid} is set, \code{run} additionally builds the full
$(k-1)$-dimensional tensor \code{C\_tau\_grid} and its per-$\ell$ version, with
each axis downsampled so the total $n^{k-1}$ evaluation count stays bounded
(\file{notebooks/daedalus.py:650-670}). The downsample target is
\code{min(tau.size, max(2, int(4000**(1/(k-1)))))} points per axis; it records
\code{tau\_axes} and, if it downsampled, a \code{\_kpoint\_grid\_note}. For $k=3$
this tensor becomes the $\kappa_3(\tau_1,\tau_2)$ heatmap.

\begin{gotcha}
The whole synthesis block is wrapped in \code{try/except Exception: pass}
(\file{notebooks/daedalus.py:632}, \code{:671-672}). If slice synthesis fails, the
result silently retains the \code{C\_tau=None}/scalar form and the plotter falls
back to bars --- no error surfaces. This is by design (so a plot always renders)
but it can mask a real failure inside \code{total\_C}. If a $k\ge 3$ plot
unexpectedly shows bars instead of curves, suspect this swallow.
\end{gotcha}

The matching simulator estimates the same slice with
\code{lag\_bins=[0, None, 0, \ldots]} (only the swept leg free); the contract is
documented at \file{notebooks/daedalus.py:588-594} and revisited in the
simulators chapter.

\section{Cumulants to moments: \texttt{\_assemble\_moment\_temporal}}\label{sec:moments}

When \code{cfg.output} is \code{'moment'} or \code{'central\_moment'}, \code{run}
assembles the moment as post-processing on top of the cumulants
(\file{notebooks/daedalus.py:674-692}); \code{compute\_cumulants} itself stays
pure. The spatial $k=2$ case is a one-liner --- $M(x) = \kappa_2(x)$ for a central
moment, plus $\mu^2$ for a raw one (\file{notebooks/daedalus.py:681-684}). Spatial
$k\ge 3$ raises \code{NotImplementedError} (\file{notebooks/daedalus.py:685-688}).
The temporal any-$k$ case is the workhorse \code{\_assemble\_moment\_temporal}
(\file{notebooks/daedalus.py:390}), which implements the shared-loop-budget
formula~\eqref{eq:sharedbudget} exactly. It returns a complex array $M$ over
\code{res['tau\_grid']}, computed as follows.

\begin{enumerate}
  \item Read the $\tau$-grid, the first external leg \code{f0}, the loop budget
        $L$, and the singleton mean \code{mu} (0 if central, else
        \code{\_external\_mean}; \file{notebooks/daedalus.py:411-414}).
  \item Build \code{\_by\_ell\_2pt(r)} (\file{notebooks/daedalus.py:416-429}): a
        dict \code{\{ell: interpolator\}} for the 2-point cumulant's $\ell$-loop
        piece. It sorts by $|\text{lag}|$ and wraps each loop order in a
        \code{np.interp} closure (or a zero closure if that order is absent), so
        the assembly can evaluate $\kappa_2^{(\ell)}$ at any lag.
  \item Get $\kappa_2$ per order --- reusing \code{res} if $k=2$, else one extra
        \code{compute\_cumulants} run with \code{k=2}
        (\file{notebooks/daedalus.py:431-436}).
  \item Get $\kappa_j$ per order for $j=3..k$ --- reusing the order-$k$ run's
        \code{total\_C\_by\_ell} callables when $j=k$, else one fresh run per
        block size (\file{notebooks/daedalus.py:438-448}).
  \item Define \code{\_kappa(block, ell, ti)} (\file{notebooks/daedalus.py:452-458}):
        the $\ell$-loop piece of the $|\text{block}|$-point cumulant at the slice
        times --- leg 1 swept to \code{tau[ti]}, the rest pinned at 0. Two-blocks
        use the interpolators; larger blocks call the $j$-point callable.
  \item The partition sum (\file{notebooks/daedalus.py:460-483}): for each
        $\tau$-index, loop over all set partitions; split each into singletons and
        multi-blocks. Skip the partition if \code{central} and there are
        singletons. Compute \code{mu\_factor = mu**len(singles)}; if there are
        singletons but the mean is zero, the term dies. If there are no
        multi-blocks (all-singleton, raw), add \code{mu\_factor}. Otherwise
        distribute the loop budget across the multi-blocks with
        \code{itertools.product(range(L+1), repeat=len(multis))}, keep only
        assignments summing to $\le L$, multiply the per-block $\ell$-pieces by
        \code{mu\_factor}, and accumulate.
\end{enumerate}

The combinatorial engine for ``loop over all set partitions'' is
\code{\_set\_partitions(items)} (\file{notebooks/daedalus.py:357}), a small
recursive generator that yields every partition of a list as a list of blocks
(the standard recursion: partition the rest, then either insert the first element
into an existing block or start a new singleton). For $k=4$ it yields the
$B_4=15$ partitions (the Bell number).

\begin{gotcha}
A temporal moment of order $k$ triggers up to $k-1$ \emph{additional}
\code{compute\_cumulants} calls (one per block size $2..k$), each with its own
diagram enumeration. Moments are therefore meaningfully more expensive than the
plain cumulant. And because singletons use only the \emph{tree} mean $\phi^{*}$,
for a symmetric saddle ($\phi^{*}=0$) raw and central moments coincide and all
singleton terms vanish.
\end{gotcha}

\section{Plotting and dispatch}\label{sec:plot}

The output side of the module is a dispatcher and seven plotters. The dispatcher
is \code{plot\_cumulant(result, cfg=None, model=None, sim=None)}
(\file{notebooks/daedalus.py:725}); \code{cfg} and \code{model} default to the
stamped \code{\_cfg}/\code{\_model} if omitted, so \code{dd.plot\_cumulant(res)}
alone works. Every plotter returns a Matplotlib \code{Figure}.

\begin{defn}
\textbf{Matplotlib \code{Figure} / \code{Axes}.} Matplotlib is the standard
Python 2-D plotting library; \code{plt} is its stateful MATLAB-style interface. A
\code{Figure} is the whole canvas; an \code{Axes} is one subplot. The methods
\code{ax.plot}, \code{ax.errorbar}, \code{ax.bar}, \code{ax.imshow}, and
\code{ax.pcolormesh} draw curves, error-bar points, bars, image heatmaps, and
mesh heatmaps respectively. Every plotter here opens with
\code{fig, ax = plt.subplots(...)} and ends by returning \code{fig}.
\end{defn}

\subsection{The dispatch tree}

The dispatch is a single ordered chain of checks
(\file{notebooks/daedalus.py:738-756}). Read top to bottom, the first match wins:

\begin{enumerate}
  \item \code{output\_kind} is a moment and a \code{moment} array is present
        $\to$ \code{\_plot\_moment}.
  \item \code{'C\_kpoint'} in result (spatial $k\ge 3$ events) $\to$
        \code{plot\_kpoint}.
  \item \code{is\_spatial(model)} or \code{'C\_tau\_x'} in result $\to$
        \code{plot\_spatial}.
  \item \code{result['C\_tau'] is None} (temporal $k\ge 3$ scalar) $\to$
        \code{plot\_temporal\_kpoint} (bars).
  \item \code{result['C\_tau\_grid'] is not None} $\to$
        \code{plot\_temporal\_kpoint\_grid} ($k=3$ heatmap).
  \item \code{len(result['C\_tau\_slices']) >= 2} $\to$
        \code{plot\_temporal\_kpoint\_slices} (one panel per $\tau_j$).
  \item Otherwise $\to$ \code{plot\_temporal} (the $C(\tau)$ curve).
\end{enumerate}

Notice how the order makes the synthesised $k\ge 3$ slice ``just work'': because
\code{run} set \code{C\_tau = slices[1]}, a sliced $k\ge 3$ result has a
non-\code{None} \code{C\_tau} and falls through check 4, landing on check 6
(slices) or, if a full grid was built, check 5 (heatmap).

\subsection{The per-order machinery shared by the line plotters}

Two helpers feed every line plot. \code{cumulative\_curves(by\_ell)}
(\file{notebooks/daedalus.py:713}) turns a per-order dict \code{\{ell: array\}}
into running cumulative sums \code{\{ell: $\Sigma_{0..\ell}$\}}; it is the basis
for the \code{'cumulative'} plot mode. \code{\_orders\_to\_draw(by\_ell, cfg)}
(\file{notebooks/daedalus.py:801}) then resolves \code{Config.show\_orders} into a
list of \code{(label, curve)} pairs: \code{'incremental'} gives the per-order
curves, \code{'total'} gives only the top cumulative curve, and
\code{'cumulative'} (the default) gives every cumulative curve. The labels come
from \code{\_order\_label(ell)} (\file{notebooks/daedalus.py:707}): \code{'tree'}
for $\ell=0$, else \code{'tree+1loop+\ldots'}. Curves are drawn in a rotating
colour from the seven-entry palette \code{\_ORDER\_COLORS}
(\file{notebooks/daedalus.py:703}), indexed modulo its length.

\subsection{The seven plotters}

\begin{description}
  \item[\code{plot\_temporal}] (\file{notebooks/daedalus.py:815}) --- the $C(\tau)$
        curve with a per-loop-order overlay. Builds curves via
        \code{\_orders\_to\_draw} (falling back to a single \code{'total'} curve),
        then overlays \code{sim} (handling a 2-D $k\ge 3$ sim by taking slice 1).
        If \code{C\_tau} is \code{None} it redirects to
        \code{plot\_temporal\_kpoint}.
  \item[\code{plot\_temporal\_kpoint\_slices}] (\file{notebooks/daedalus.py:856})
        --- $k\ge 3$: one panel per time-difference $\tau_j$, from
        \code{C\_tau\_slices} / \code{C\_tau\_slices\_by\_ell}, each with the
        per-order overlay and (optionally) the matching sim row. Axes labelled
        $\tau_j = t_j - t_0$ and $\kappa_k$.
  \item[\code{plot\_temporal\_kpoint\_grid}] (\file{notebooks/daedalus.py:923})
        --- $k=3$: a \code{pcolormesh} heatmap of $\kappa_3(\tau_1,\tau_2)$ from
        \code{C\_tau\_grid} and \code{tau\_axes}. For $k\ge 4$ (no 2-D view) it
        falls back to the slices plotter.
  \item[\code{plot\_temporal\_kpoint}] (\file{notebooks/daedalus.py:948}) ---
        temporal $k\ge 3$ equal-time: a single scalar per loop order (the $\tau=0$
        value), drawn as tree/+loop \emph{bars} honouring \code{show\_orders}. A
        scalar sim overlays as a dashed \code{axhline} with an \code{axhspan}
        error band.
  \item[\code{plot\_spatial}] (\file{notebooks/daedalus.py:1000}) --- spatial
        $k=2$: the equal-time $C(x,0)$ slice (panel 0) plus, when a $\tau$-grid is
        present, a $C(x,\tau)$ \code{imshow} heatmap (panel 1). The $\tau\approx 0$
        row is picked via \code{argmin(|tau|)}; the per-order overlay reads
        \code{spatial\_info['C\_by\_order']}.
  \item[\code{plot\_kpoint}] (\file{notebooks/daedalus.py:1054}) --- spatial
        $k\ge 3$ at explicit events: a per-event bar chart from \code{C\_kpoint},
        with the tree/+loop decomposition as grouped bars unless
        \code{show\_orders == 'total'}.
  \item[\code{\_plot\_moment}] (\file{notebooks/daedalus.py:759}) --- the
        assembled moment as a \emph{single} curve (moments mix loop orders, so no
        per-order overlay): temporal $M_k(\tau)$, or spatial $k=2$ $M(x,0)$.
        $y$-label $M_k$.
\end{description}

\begin{gotcha}
\code{plot\_kpoint} (spatial $k\ge 3$) is the one plotter that takes \emph{no}
\code{sim} argument (\file{notebooks/daedalus.py:1054}), so a simulator overlay of
spatial $k\ge 3$ events is not supported through \code{plot\_cumulant}.
\end{gotcha}

\subsection{The \texttt{sim} overlay contract}

The optional \code{sim} dict overlays an independent simulator on the theory plot.
Its keys depend on the group (\file{notebooks/daedalus.py:734-736}, \code{:866-867},
\code{:980-986}):

\begin{itemize}
  \item temporal $C(\tau)$: \code{\{'tau', 'C', 'C\_err'\}};
  \item spatial: \code{\{'x', 'C', 'C\_err'\}};
  \item temporal $k\ge 3$ multi-slice: \code{'C'} / \code{'C\_err'} may be 2-D
        \code{(k-1, n\_tau)} (row $p$ is panel $p$) or 1-D (applied to panel 0);
  \item temporal $k\ge 3$ scalar: \code{\{'C': <scalar>, 'C\_err': <scalar or
        None>\}}.
\end{itemize}

Where \code{C\_err} is given it is drawn as error bars (or, for the scalar case,
an \code{axhspan} band); where it is absent the points are drawn plain. This is
the same contract the example in Chapter~\ref{ch:quickstart} used.

\section{Summarising a run}

\code{summary(result) -> str} (\file{notebooks/daedalus.py:1087}) is the
one-glance digest every notebook prints. It reads the \code{\_resolved} and
\code{\_model} stamps and reports the theory name, the resolved $k$ and
\code{max\_ell}, the field names, the spatial dimension, the Dyson order if the
model carries a fixed-mode Dyson policy, and the live-diagram count
\code{spatial\_info['n\_live\_diagrams']} when present:

\begin{lstlisting}[style=console]
theory : 'OU Quartic (white noise)'
k      : 2    max_ell : 2
fields : ['dx']   spatial_dim : 0
\end{lstlisting}

The live-diagram count is the only place a \code{nauty}-derived number surfaces to
the user; it is the count of distinct (deduplicated) diagrams the engine actually
evaluated. Everything else \code{summary} prints comes straight from the stamps.

\section{What lives behind \texttt{compute\_cumulants}}

It is worth being explicit about the boundary, because it explains why this
chapter never had to teach you SageMath. \code{run} and
\code{\_assemble\_moment\_temporal} both do \code{from pipeline import
compute\_cumulants} (\file{notebooks/daedalus.py:410}, \code{:492}); everything
below happens \emph{inside} that call, and \code{daedalus} is a \emph{client}, not
a user, of these libraries:

\begin{description}
  \item[SageMath] a Python computer-algebra system; its symbolic ring \code{SR},
        symbolic matrices, and exact algebra build the propagator, find
        $\omega$-poles and residues, and solve the mean-field self-consistency.
        Result keys like \code{num\_params} and the \code{propagator} matrices are
        Sage objects --- but \code{daedalus} only ever reads \emph{NumPy} arrays
        out of the result, so it never has to know Sage exists.
  \item[\code{nauty}] a graph canonicalisation and isomorphism library; the
        pipeline uses it to deduplicate enumerated diagrams and compute symmetry
        factors. It surfaces only as the diagram counts in \code{spatial\_info}.
  \item[\code{sympy}] a lighter pure-Python CAS, used for symbolic manipulation
        and \code{lambdify} (compiling a symbolic expression into a fast numeric
        function).
  \item[\code{scipy}] numerical routines on top of NumPy --- \code{quad}
        integration as a 1-D fallback, Bessel functions for the spatial radial
        transform.
  \item[\code{numba}] a just-in-time compiler for the engine's hot inner loops.
  \item[\code{networkx}] a pure-Python graph library used in enumeration and
        connectivity checks.
\end{description}

The single takeaway: \emph{\code{daedalus.py} is pure orchestration; the only
non-stdlib libraries it imports are NumPy and Matplotlib.} All of the
computer-algebra, graph, and JIT machinery is on the far side of one
\code{from pipeline import compute\_cumulants}. The chapters of Parts~III--V are
the deep dive into what that one line does; this chapter was the deep dive into
everything around it.

\section{Gotchas and limitations}\label{sec:gotchas}

A consolidated list of the sharp edges, several already flagged above. Most are by
design; a few are dead wiring worth knowing about so you do not go looking for
behaviour that is not there.

\begin{enumerate}
  \item \textbf{\code{run} mutates the model dict in place}
        (\file{notebooks/daedalus.py:545-550}): the Dyson override writes into
        \code{model['spatial']}. Reusing one \code{model} across runs with
        different \code{dyson\_order} leaves the last write in place;
        re-\code{load\_theory} for a clean model.
  \item \textbf{$k$ versus \code{external\_fields} contradiction is fatal}
        (\file{notebooks/daedalus.py:506-511}): passing both with disagreeing leg
        counts raises rather than silently choosing.
  \item \textbf{Explicit \code{external\_fields} is used verbatim}
        (\file{notebooks/daedalus.py:519-533}): only an \emph{absent} value (or a
        stale METADATA recommendation) is auto-rebuilt. An explicit list naming a
        missing field is passed straight to the engine and fails \emph{there}, not
        in \code{run}.
  \item \textbf{Spatial $k\ge 3$ needs \code{spatial\_points}}
        (\file{notebooks/daedalus.py:566-570}) --- there is no grid form --- and
        \textbf{spatial $k\ge 3$ moments are \code{NotImplemented}}
        (\code{:685-688}).
  \item \textbf{The $k\ge 3$ synthesis swallows all exceptions}
        (\file{notebooks/daedalus.py:632}, \code{:671-672}): on failure the result
        silently keeps its scalar form and the plotter falls back to bars.
  \item \textbf{$\tau$-grid caps for $k\ge 3$}: slices cap the grid to 41 points
        (\file{notebooks/daedalus.py:601-603}); the full grid downsamples each
        axis (\code{:652}). A $k\ge 3$ plot may be coarser than the same theory's
        $k=2$ plot.
  \item \textbf{Tadpole approximation in moments}: singleton blocks use the
        \emph{tree} mean $\phi^{*}$; the 1-loop tadpole shift is not folded in
        (\file{notebooks/daedalus.py:376-377}, \code{:404-405}).
  \item \textbf{\code{Config.components} is declared but never read} in this file
        --- the docstring advertises it (\file{notebooks/daedalus.py:23-25}) and
        the field exists (\code{:307}), but no plotter consults it.
  \item \textbf{\code{is\_multifield} exists but \code{plot\_cumulant} never
        branches on it} --- the dispatcher routes on spatial/$k$ only, even though
        the docstring mentions a multi-field axis. Multi-field plotting is not
        wired in this file.
  \item \textbf{\code{from dataclasses import \ldots field} is imported but
        unused} (\file{notebooks/daedalus.py:31}) --- every \code{Config} field
        uses a plain default. Harmless dead import.
\end{enumerate}

With the front desk fully mapped, the next chapters turn to the two pieces this
one only gestured at: the cumulant-to-moment conversion in its own right
(Chapter~\ref{ch:cumulants-moments}, for the combinatorics in full), and the
independent simulators that produce the \code{sim} overlays we kept drawing on
top of the theory.

% ---------------------------------------------------------------------
%  Chapter 19 : From Cumulants to Moments, and the k>=3 Synthesis
% ---------------------------------------------------------------------
\chapter{From Cumulants to Moments, and the \texorpdfstring{$k\!\ge\!3$}{k>=3} Synthesis}\label{ch:cumulants-moments}

The engine's native output is a \term{connected cumulant}. Everything the
seven-phase pipeline computes --- every diagram it enumerates, every loop
integral Phase~J performs --- adds up to one object: the connected $k$-point
function $\kappa_k$. That is the right primitive for a diagrammatic expansion,
because the connected diagrams \emph{are} the cumulant; nothing further has to be
done to ``connect'' them. But it is not always the object a physicist wants to
read off a plot. Sometimes you want a full \term{moment} $\avg{\phi\cdots\phi}$,
or a \term{central moment} of the fluctuation $\delta\phi = \phi - \avg{\phi}$.
And for a correlator of three or more legs, $\kappa_k$ is not a curve at all ---
it is a function of several time differences that the engine refuses to
materialise as a grid, handing you a \emph{callable} instead.

This chapter is about the thin but conceptually dense layer that sits between the
engine's raw cumulant output and the plot. It does three jobs, all inside
\file{notebooks/daedalus.py} and never inside the engine:

\begin{enumerate}
  \item \textbf{Cumulant $\to$ moment assembly.} The set-partition (cluster)
        expansion reconstructs full and central moments from cumulants, with two
        physics subtleties handled explicitly: a single \emph{shared loop budget}
        across the blocks of a partition, and the treatment of singleton blocks
        as the mean-field mean. This is the \code{Config.output} switch.
  \item \textbf{The $k\!\ge\!3$ slice / grid synthesis.} When the engine returns
        a callable instead of an array, \code{run} synthesises the natural
        axis-parallel slices (and, on request, the full multidimensional grid) so
        that a three-point cumulant plots and sim-compares exactly like a
        two-point one.
  \item \textbf{Dispatching the result.} A short decision tree in
        \code{plot\_cumulant} routes each of these shapes --- moment, $k\!\ge\!3$
        scalar, $k\!\ge\!3$ slices, $k\!\ge\!3$ heatmap, ordinary curve --- to the
        right plotter.
\end{enumerate}

A recurring theme is worth stating once, up front: the heavy engine,
\code{compute\_cumulants} (\file{pipeline/compute.py:108}), stays \emph{pure}. It
returns the connected cumulant and nothing else; it contains no
moment$\leftrightarrow$cumulant bookkeeping whatsoever. Everything in this chapter
is post-processing layered on top of that pure output by the front-desk module.
Keeping the boundary sharp is what lets the same engine serve cumulants, moments,
slices, and grids without a single conditional in its core.

\section{Why cumulants are the native object}

Recall the structure of the loop expansion. The connected $k$-point cumulant of
a field $\phi$ is
\begin{equation}
  \kappa_k(x_1,t_1;\ldots;x_k,t_k)
  \;=\; \avg{\,\phi(x_1,t_1)\cdots\phi(x_k,t_k)\,}_{\mathrm c},
  \label{eq:cumulant-def}
\end{equation}
the \emph{connected} (cluster) part of the $k$-point correlation. In \msrjd{}
field theory this is exactly the sum of all \emph{connected} Feynman diagrams
with $k$ external legs, organised by loop order $\ell$:
\begin{equation}
  \kappa_k \;=\; \kappa_k^{(0)} + \kappa_k^{(1)} + \kappa_k^{(2)} + \cdots,
  \qquad
  \underbrace{\kappa_k^{(0)}}_{\text{tree}}\;,\;
  \underbrace{\kappa_k^{(1)}}_{\text{1-loop}}\;,\;\ldots
  \label{eq:cumulant-loops}
\end{equation}
\code{Config.max\_ell} $=L$ truncates the series at $\ell\le L$. The engine
returns the per-order pieces as \code{*\_by\_ell} dictionaries
$\{\ell:\text{array or callable}\}$ and their running sum $\sum_{\ell\le L}$ as
the ``total''. This per-$\ell$ split is produced inside the engine by running
Phase~J once for each loop order and summing the master closure across orders:

\begin{lstlisting}[style=console]
total_C(*tau) = sum over ell of total_C_by_ell[ell](*tau)   # pipeline/compute.py:802
C_tau         = sum over ell of C_tau_by_ell[ell]           # pipeline/compute.py:808
\end{lstlisting}

The orchestrator's brief is blunt about the scope of what the engine assembles:
\code{compute\_cumulants} returns the \emph{connected} cumulant directly (the
diagrammatic sum \emph{is} the cumulant); moment and central-moment conversions,
if needed, live in downstream consumers, not in that file. That downstream
consumer is the subject of this chapter.

\begin{defn}
A \term{connected cumulant} $\kappa_k$ is the connected part of a $k$-point
correlation function --- diagrammatically, the sum of all \emph{connected}
diagrams. A \term{moment} is the full correlation $\avg{\phi\cdots\phi}$ (raw) or
the same quantity for the centred field $\delta\phi=\phi-\avg{\phi}$ (central).
Moments and cumulants are related by the set-partition (M\"obius-inversion)
identity of \S\ref{sec:setpart}; the engine emits cumulants, and \Daedalus{}
reconstructs moments on demand.
\end{defn}

For $k=2$ the connected cumulant is just the (connected) two-point function. In
the \textbf{temporal} case the engine returns it as a curve over a lag grid
$\tau$,
\begin{equation}
  C(\tau) \;=\; \avg{\phi(0)\,\phi(\tau)}_{\mathrm c}
  \qquad\text{(key \code{C\_tau} over \code{tau\_grid})},
\end{equation}
and in the \textbf{spatial} case as a real-space correlator over separations $x$
and lags $\tau$,
\begin{equation}
  C(x,\tau) \;=\; \avg{\phi(0,0)\,\phi(x,\tau)}_{\mathrm c}
  \qquad\text{(key \code{C\_tau\_x}, shape $(n_\tau,n_x)$)}.
\end{equation}
The equal-time slice $C(x,0)$ is the headline spatial plot. These two-point cases
are the easy ones: the result is already a NumPy array, and the plotter draws it
directly. The two harder cases --- $k\!\ge\!3$, where the cumulant is a callable,
and moments, where blocks of cumulants must be multiplied together --- are what
the rest of the chapter builds.

\section{The set-partition (cluster) expansion}\label{sec:setpart}

\subsection{The combinatorial identity}

Moments and cumulants are two encodings of the same statistical information,
related by the \term{set-partition} (or \term{cluster}) expansion. For the
centred field, the raw $k$-point moment is the sum, over every \emph{set
partition} $\pi$ of the index set $\{1,\ldots,k\}$, of the product of cumulants
over the blocks of that partition:
\begin{equation}
  M_k(x_1,\ldots,x_k)
  \;=\; \sum_{\pi}\;\prod_{B\in\pi}\kappa\!\left(B\right).
  \label{eq:moment-setpart}
\end{equation}
A \emph{set partition} of $\{1,\ldots,k\}$ is a way of splitting those indices
into disjoint, nonempty, unordered \emph{blocks} whose union is the whole set.
Each block $B$ contributes the connected cumulant of the legs it contains:
$\kappa(\{i\})=\avg{\phi_i}$ is the mean, $\kappa(\{i,j\})=\avg{\phi_i\phi_j}_{\mathrm c}$
is the connected two-point function, and so on. The number of partitions of a
$k$-element set is the Bell number $B_k$: $B_2=2$, $B_3=5$, $B_4=15$ --- so a
four-point moment is a sum of fifteen products of cumulants. This identity is the
exact inverse of the more familiar statement that the cumulant is the connected
part: it rebuilds the full (disconnected-included) moment by summing over every
way the legs can cluster.

\begin{example}\label{ex:k2-moment}
For $k=2$ there are exactly two partitions of $\{1,2\}$: the singletons
$\{\{1\},\{2\}\}$ and the pair $\{\{1,2\}\}$. So
\[
  M_2 \;=\; \kappa(\{1\})\,\kappa(\{2\}) \;+\; \kappa(\{1,2\})
       \;=\; \avg{\phi}^2 \;+\; \avg{\phi\,\phi}_{\mathrm c}.
\]
The raw second moment is the squared mean plus the connected variance; the
central second moment drops the $\avg{\phi}^2$ singleton term and is just the
connected variance. This is exactly the spatial $k=2$ case the code handles in
one line --- see \S\ref{sec:moment-run}.
\end{example}

This identity is stated verbatim in two places in the code. The
\code{Config.output} comment writes it as
\code{M = Sigma\_pi prod\_B kappa(B)}
(\file{notebooks/daedalus.py:261--266}), and the docstring of the assembly
workhorse \code{\_assemble\_moment\_temporal} writes the loop-resolved form we
turn to next (\file{notebooks/daedalus.py:390--408}).

\subsection{Subtlety 1: the shared loop budget}

Equation~\eqref{eq:moment-setpart} is written in terms of \emph{fully dressed}
cumulants $\kappa(B)$. If you took that literally at finite loop order you would
introduce spurious cross-terms. Consider a two-block partition at requested loop
order $L=1$: the naive product $\kappa(B_1)\,\kappa(B_2)$ would include the term
$\kappa^{(1)}(B_1)\cdot\kappa^{(1)}(B_2)$, a product of two one-loop objects. But
that term is \emph{two loops} of total order --- it has no business appearing in a
one-loop moment. A perturbatively consistent $L$-loop moment must contain only
terms whose \emph{total} loop order is $\le L$.

The code therefore enforces a \textbf{single shared loop budget} $L=\code{max\_ell}$
across all the multi-element blocks of each partition. The loop-resolved moment is
\begin{equation}
  M_k
  \;=\;
  \sum_{\pi}
  \;\sum_{\substack{(\ell_B)_{B\in\pi}\\[1pt]\sum_B \ell_B\,\le\,L}}
  \;\prod_{B\in\pi}\kappa(B)^{(\ell_B)},
  \label{eq:moment-budget}
\end{equation}
where the inner sum runs over every assignment of a loop order $\ell_B$ to each
block, subject to the constraint that the orders \emph{add up} to at most $L$.
Every surviving term then has total loop order $\le L$ by construction. The two
definitions --- naive product of dressed cumulants versus the
shared-budget~\eqref{eq:moment-budget} --- agree only at tree level ($L=0$, where
the only assignment is all-zero) and diverge at $L\ge1$ for any multi-block
partition. The code's docstring states this precisely
(\file{notebooks/daedalus.py:398--408}): the shared-budget form is the
perturbatively-consistent $L$-loop moment, in which the total loop order of every
surviving term is $\le L$, so a one-loop moment never smuggles in a partial
two-loop $\ell_{B_1}\!\cdot\ell_{B_2}$ cross-term.

\begin{note}
The shared-budget constraint $\sum_B \ell_B\le L$ is what distinguishes a
\emph{consistent} perturbative moment from a careless one. It is the moment-space
echo of the same bookkeeping that the loop expansion does everywhere: you may
spend your $L$ loops however you like across the diagram, but you may not spend
more than $L$ of them in total. Spreading the budget across the blocks of a
partition is the only honest way to multiply truncated cumulants.
\end{note}

\subsection{Subtlety 2: singletons and the tree mean}

A singleton block $\{i\}$ contributes the one-point cumulant, which is just the
field mean $\avg{\phi}$. At tree level the mean is the \emph{mean-field saddle}
$\phi^{*}$ --- the deterministic fixed point the whole fluctuation expansion is
built around (Chapter~\ref{ch:quickstart}). The code reads it from the result
dict's \code{mf\_values} via the helper \code{\_external\_mean}
(\file{notebooks/daedalus.py:373}):

\begin{lstlisting}
def _external_mean(model: dict, res: dict) -> float:
    """<phi> for the (single) external field -- the mean-field saddle phi*."""
    mfv = res.get('mf_values') or {}
    names = field_names(model)
    f0 = names[0] if names else None
    cands = [f0] + ([f0[1:]] if f0 and f0.startswith('d') else [])
    for key in cands:
        if key in mfv:
            try:
                return float(np.real(np.asarray(mfv[key]).ravel()[0]))
            except Exception:
                pass
    return 0.0
\end{lstlisting}

Two details earn the over-explanation. First, the candidate-key dance
(\code{cands = [f0] + ...}): the external leg name may be the \code{d}-prefixed
fluctuation name (e.g.\ \code{dh} for the response/fluctuation of \code{h}), so
the helper tries both the prefixed name and the stripped one when looking up the
saddle value. Second, the fallback to \code{0.0}: for a \emph{symmetric} saddle
(the common $\phi^{*}=0$ case, as in the single-well OU example) the mean is
zero, and the helper returns zero whether or not the key is found.

How singletons enter the moment depends on whether you asked for a raw or a
central moment:
\begin{itemize}
  \item For a \textbf{central} moment, every singleton is \emph{dropped}: only
        no-singleton partitions survive. This is the field-theory statement that
        the central moment is built from the fluctuation $\delta\phi=\phi-\avg{\phi}$,
        whose own mean is zero. In the code this is the early
        \code{continue} when \code{central and singles}
        (\file{notebooks/daedalus.py:466}).
  \item For a \textbf{raw} moment, a partition with $s$ singletons contributes a
        factor $\mu^{s}$, where $\mu=\avg{\phi}$. If the mean is zero those terms
        die automatically --- the code skips them explicitly with the
        \code{mu\_factor == 0.0} guard (\file{notebooks/daedalus.py:468--472}).
        So for a symmetric saddle the raw and central moments coincide and every
        singleton term vanishes.
\end{itemize}

\begin{gotcha}
The singleton mean is the \emph{tree-level} saddle $\phi^{*}$. The one-loop
\term{tadpole} shift of the mean --- the correction
$\avg{\phi}=\phi^{*}+\avg{\delta\phi}^{(1)}+\cdots$ --- is \textbf{not} folded into
the singleton blocks. This is a documented refinement, flagged twice in the source
(\file{notebooks/daedalus.py:376--377} and \file{:404--405}). The practical
consequence: for a symmetric saddle ($\phi^{*}=0$) it makes no difference (the
singleton terms vanish either way), but for an \emph{asymmetric} saddle a raw
moment uses the bare $\phi^{*}$ for its singleton factors and is therefore only
exact for those factors at tree level.
\end{gotcha}

\subsection{The cost: extra backend runs}

There is a price for assembling a moment. Equation~\eqref{eq:moment-budget} needs
the per-order cumulant $\kappa_j^{(\ell)}$ for \emph{every} block size $j$ that
appears in any partition --- that is, every $j$ from $2$ up to $k$. Each block
size is a different correlator order, and a different correlator order needs its
own diagram enumeration and its own Phase~J integration. So producing a $k$-point
moment costs up to $k-1$ \emph{additional} \code{compute\_cumulants} calls (one
per order $2,3,\ldots,k$), beyond the order-$k$ run you already did. The order-$k$
run is reused; the smaller orders are fresh. Each fresh run delivers \emph{all}
its loop orders at once (the \code{*\_by\_ell} dict), so it is one run per block
size, not one run per (size, loop) pair. This cost is recorded in the
\code{Config.output} comment (\file{notebooks/daedalus.py:264--266}) and realised
in the assembly function.

\section{Walking \texttt{\_assemble\_moment\_temporal}}\label{sec:assemble}

The temporal moment is built by \code{\_assemble\_moment\_temporal(model, res,
kw, k, central)} (\file{notebooks/daedalus.py:390}). Its arguments: \code{model}
and \code{res} are the model dict and the result dict from the order-$k$ run;
\code{kw} is the keyword dictionary already assembled for
\code{compute\_cumulants} (so the smaller-order reruns inherit identical
parameters, grid, and external-leg conventions); \code{k} is the order; and
\code{central} is the raw-versus-central flag. It returns a complex array $M$ over
\code{res['tau\_grid']} --- the moment evaluated along the same one-dimensional
slice that \code{C\_tau} uses, so it plots like any other curve.

It proceeds in five movements.

\subsection*{(a) Setup and the local imports}

\begin{lstlisting}
import itertools
from pipeline import compute_cumulants
tau = np.asarray(res['tau_grid'], dtype=float)
f0  = kw['external_fields'][0]
L   = int(kw.get('max_ell', 0) or 0)
mu  = 0.0 if central else _external_mean(model, res)
\end{lstlisting}

\code{itertools} is Python's standard iterator-combinatorics library; here it
supplies \code{itertools.product}, the Cartesian product used to enumerate
loop-order assignments. \code{compute\_cumulants} is imported locally (inside the
function, not at module top) so that \code{daedalus} can be imported even in an
environment where the heavy engine is unavailable, and so the engine is only
pulled in when a moment is actually requested. \code{L} is the shared loop budget;
\code{mu} is the singleton mean, forced to zero for a central moment.

\subsection*{(b) The two-point cumulant interpolators}

The assembly needs to evaluate $\kappa_2^{(\ell)}$ at an arbitrary lag, but the
engine only returns it on the grid. The nested helper \code{\_by\_ell\_2pt}
(\file{notebooks/daedalus.py:416--429}) turns each per-order grid array into a
one-dimensional interpolating closure:

\begin{lstlisting}
def _by_ell_2pt(r):
    """{ell: even-lag interpolator} of the 2-point cumulant's ell-loop piece."""
    be = r.get('C_tau_by_ell') or {0: r.get('C_tau')}
    t = np.asarray(r['tau_grid']); a = np.abs(t); o = np.argsort(a); ax = a[o]
    out = {}
    for e in range(L + 1):
        arr = be.get(e)
        if arr is None:
            out[e] = (lambda lag: 0.0 + 0.0j)
        else:
            ay = np.real(np.asarray(arr))[o]
            out[e] = (lambda lag, ax=ax, ay=ay:
                      complex(np.interp(abs(lag), ax, ay)))
    return out
\end{lstlisting}

\code{np.argsort(a)} returns the index order that sorts the absolute lags, so
\code{np.interp} (NumPy's piecewise-linear interpolator, which requires a
monotonically increasing $x$-axis) sees a clean ascending grid. The interpolator
is built against the absolute lag because the two-point correlator is even in
$\tau$. The defaulted keyword arguments \code{ax=ax, ay=ay} are the standard
Python idiom for \emph{capturing} the current arrays into the closure, so each
loop iteration's lambda binds its own data rather than the loop variable's final
value. A loop order that the engine did not populate (\code{arr is None}) gets a
zero closure, so the assembly can ask for any $\ell\in\{0,\ldots,L\}$
unconditionally.

\subsection*{(c) Gathering per-order cumulants for every block size}

The two-point piece comes either from the order-$k$ run already in hand (when
$k=2$) or from one fresh two-leg run; the higher block sizes $j=3,\ldots,k$ each
get their callable \code{total\_C\_by\_ell} (reusing the order-$k$ run when $j=k$):

\begin{lstlisting}
# kappa_2 per loop order (reuse the order-k run when k==2, else one 2-pt run).
if k == 2:
    k2 = _by_ell_2pt(res)
else:
    kw2 = dict(kw); kw2.update(k=2, external_fields=[f0] * 2)
    k2 = _by_ell_2pt(compute_cumulants(**kw2))

# kappa_j per loop order for j = 3..k (reuse the order-k *_by_ell).
kj = {}
for j in range(3, k + 1):
    if j == k and res.get('total_C_by_ell'):
        be = res['total_C_by_ell']
    else:
        kwj = dict(kw); kwj.update(k=j, external_fields=[f0] * j)
        be = compute_cumulants(**kwj).get('total_C_by_ell') or {}
    kj[j] = {e: (be.get(e) if callable(be.get(e))
                 else (lambda *t: 0.0 + 0.0j))
             for e in range(L + 1)}
\end{lstlisting}

This is where the ``$k-1$ extra runs'' are spent: \code{dict(kw)} clones the base
keyword dict, then \code{.update(k=j, external\_fields=[f0]*j)} retargets it to a
$j$-leg correlator of the first external field with itself. Each block size larger
than two is stored as a dict of callables (one per loop order), again with a
zero-callable fallback for missing orders.

\subsection*{(d) Evaluating one cumulant block at the slice times}

The slice the moment is computed along sweeps leg~$1$ over $\tau$ and pins the
other legs at zero. The nested \code{\_kappa(block, ell, ti)} returns the
$\ell$-loop piece of the $|\text{block}|$-point cumulant at those slice times
(\file{notebooks/daedalus.py:452--458}):

\begin{lstlisting}
def _kappa(block, ell, ti):
    """ell-loop piece of the |block|-point cumulant at the slice times."""
    b = sorted(block)
    if len(b) == 2:                       # leg 1 swept (tau); rest pinned (0)
        return k2[ell](float(tau[ti]) if 1 in b else 0.0)
    times = [float(tau[ti]) if idx == 1 else 0.0 for idx in b]
    return complex(kj[len(b)][ell](*times))
\end{lstlisting}

A two-element block uses the interpolators; the lag fed in is $\tau[ti]$ if leg
index $1$ is in the block and zero otherwise (the swept-leg convention). A larger
block calls its $j$-point callable with a time vector that is $\tau[ti]$ in the
swept-leg slot and zero elsewhere.

\subsection*{(e) The partition sum}

The heart of the function (\file{notebooks/daedalus.py:460--483}) assembles
equation~\eqref{eq:moment-budget} index by index over the $\tau$-grid:

\begin{lstlisting}
parts = list(_set_partitions(list(range(k))))
M = np.empty(tau.size, dtype=complex)
for ti in range(tau.size):
    tot = 0.0 + 0.0j
    for part in parts:
        singles = [bl for bl in part if len(bl) == 1]
        multis  = [bl for bl in part if len(bl) > 1]
        if central and singles:           # central drops every singleton
            continue
        mu_factor = mu ** len(singles)    # singletons: tree mean, ell=0 only
        if singles and mu_factor == 0.0:  # zero mean -> singleton terms die
            continue
        if not multis:                    # all-singleton (raw): pure mu^k
            tot += mu_factor
            continue
        # share the loop budget L across the multi-blocks (sum ell_B <= L):
        for assign in itertools.product(range(L + 1), repeat=len(multis)):
            if sum(assign) > L:
                continue
            prod = mu_factor
            for bl, e in zip(multis, assign):
                prod *= _kappa(bl, e, ti)
            tot += prod
    M[ti] = tot
return M
\end{lstlisting}

Read it against the formula. For each grid index \code{ti} we sum over partitions
\code{parts}. Each partition splits into its singletons and its multi-element
blocks. The central/singleton skip and the zero-mean skip implement the singleton
rules of \S\ref{sec:setpart}. An all-singleton partition (no multi-blocks)
contributes the pure mean product $\mu^{k}$. Otherwise the
\code{itertools.product(range(L+1), repeat=len(multis))} enumerates every
loop-order assignment $(\ell_B)$ to the multi-blocks, the
\code{if sum(assign) > L: continue} enforces the shared-budget constraint
$\sum_B\ell_B\le L$ exactly, and the inner product multiplies the per-block
$\ell$-pieces (each from \code{\_kappa}) times the singleton factor $\mu^{s}$.
That is equation~\eqref{eq:moment-budget}, line for line.

\subsection{The set-partition generator}\label{sec:setpart-gen}

The partitions themselves come from \code{\_set\_partitions}
(\file{notebooks/daedalus.py:357}), a small recursive generator:

\begin{lstlisting}
def _set_partitions(items):
    """Yield every set partition of items (a list) as a list of blocks."""
    items = list(items)
    if not items:
        yield []
        return
    if len(items) == 1:
        yield [items]
        return
    first, rest = items[0], items[1:]
    for sub in _set_partitions(rest):
        for i in range(len(sub)):
            yield sub[:i] + [[first] + sub[i]] + sub[i + 1:]
        yield [[first]] + sub
\end{lstlisting}

It is the textbook recursion: to partition a set, partition everything except the
first element, then for each such partition either drop the first element into one
of the existing blocks (the inner \code{for i} loop) or start a fresh singleton
block for it (the trailing \code{yield}). A \term{generator} is a Python function
that \code{yield}s values lazily one at a time rather than building a list, so the
partitions stream out without ever all living in memory at once --- though the
caller in \code{\_assemble\_moment\_temporal} does materialise them with
\code{list(...)} because it iterates them once per $\tau$-index. For $k=4$ it
yields all $B_4=15$ partitions.

\section{The \texttt{Config.output} switch and where it fires}\label{sec:moment-run}

The user-facing control is one \code{Config} field
(\file{notebooks/daedalus.py:266}):

\begin{lstlisting}
output: str = 'cumulant'                # 'cumulant'|'moment'|'central_moment'
\end{lstlisting}

It defaults to \code{'cumulant'}, the engine's native output, in which case none
of this chapter's assembly runs at all. Setting it to \code{'moment'} or
\code{'central\_moment'} triggers the conversion inside \code{run}, \emph{after}
the engine call, as a post-processing step (\file{notebooks/daedalus.py:674--692}):

\begin{lstlisting}
out_kind = getattr(cfg, 'output', 'cumulant') or 'cumulant'
if out_kind in ('moment', 'central_moment'):
    central = out_kind == 'central_moment'
    if is_spatial(model):
        if k == 2:                      # M(x) = kappa_2(x) [+ mu^2 for raw]
            mu = 0.0 if central else _external_mean(model, res)
            res['moment'] = (np.real(np.asarray(res['C_tau_x']))
                             + (0.0 if central else mu * mu))
        else:
            raise NotImplementedError(
                f"Config.output={out_kind!r} for spatial k>=3 is not "
                "implemented yet (temporal any-k and spatial k=2 are).")
    else:
        res['moment'] = _assemble_moment_temporal(
            model, res, kw, k, central)
    res['output_kind'] = out_kind
\end{lstlisting}

There are three branches, matching the three cases the framework supports:

\begin{description}
  \item[Spatial $k=2$] The whole set-partition sum collapses to
        Example~\ref{ex:k2-moment}: $M(x)=\kappa_2(x)+\mu^{2}$ (raw) or
        $M(x)=\kappa_2(x)$ (central). The code reads the connected correlator
        straight out of \code{C\_tau\_x} and adds $\mu^2$ for the raw case --- no
        extra backend runs needed, because a two-point moment needs only the
        two-point cumulant plus the mean.
  \item[Spatial $k\!\ge\!3$] Raises \code{NotImplementedError}. The multivariate
        moment assembly is wired for temporal-any-$k$ and spatial $k=2$ only;
        spatial $k\!\ge\!3$ moments are an open gap.
  \item[Temporal (any $k$)] Delegates to \code{\_assemble\_moment\_temporal} of
        \S\ref{sec:assemble}, the full shared-budget cluster expansion.
\end{description}

In all surviving cases the assembled curve is stored under \code{res['moment']}
and the chosen kind is recorded under \code{res['output\_kind']} --- the two keys
the dispatcher checks first (\S\ref{sec:dispatch}).

\begin{note}
The engine call itself is identical whether you ask for a cumulant or a moment;
the moment is computed entirely from the cumulant output afterward. This is the
clean-boundary principle in action: \code{compute\_cumulants} never learns that a
moment was requested. If you want both a cumulant plot and a moment plot for the
same parameter point, you pay the engine cost once and the extra moment-assembly
runs once.
\end{note}

\section{The \texorpdfstring{$k\!\ge\!3$}{k>=3} synthesis}\label{sec:kge3}

\subsection{Why a three-point cumulant is not a curve}

By time-translation invariance, a connected $k$-point cumulant depends only on the
$k-1$ time differences relative to one anchored leg:
\begin{equation}
  \tau_j \;=\; t_j - t_0,\qquad j=1,\ldots,k-1
  \qquad\text{(leg $0$ pinned at $\tau=0$)}.
  \label{eq:taudiffs}
\end{equation}
So $\kappa_k$ is a function $\kappa_k(\tau_1,\ldots,\tau_{k-1})$. For $k=2$ that is
a single variable (the familiar $C(\tau)$); for $k=3$ it is a two-dimensional
surface $\kappa_3(\tau_1,\tau_2)$; for $k\ge4$ it is a $(k-1)$-dimensional tensor
that cannot be drawn directly. Materialising that whole tensor on a grid would cost
$n^{k-1}$ evaluations --- prohibitive even for modest $n$ and $k$. So the engine
\textbf{does not} build it. For $k\ge3$ it returns \code{C\_tau = None} and hands
back a \emph{callable} \code{total\_C(*tau)} plus the per-loop callables
\code{total\_C\_by\_ell = \{ell: callable\}}. This is exactly what we saw at the
aggregation step: \code{C\_tau} is \code{None} when the per-$\tau$ evaluation
pattern was skipped for $k\ge3$ (\file{pipeline/compute.py:807--810}).

\Daedalus{}'s \code{run} then synthesises two derived views from those callables so
that a $k\ge3$ cumulant can still be plotted and sim-compared. The whole block
lives at \file{notebooks/daedalus.py:588--672} and fires only when the model is
temporal, $k\ge3$, the engine returned \code{C\_tau is None}, and \code{total\_C}
is callable:

\begin{lstlisting}
if (not is_spatial(model) and k >= 3 and res.get('C_tau') is None
        and callable(res.get('total_C'))):
\end{lstlisting}

\subsection{Choosing the evaluation grid and the base point}

First the synthesis bounds the cost by capping the $\tau$-grid to $41$ points
(each evaluation of \code{total\_C} is a full diagram sum, so the count matters):

\begin{lstlisting}
tau = np.asarray(res.get('tau_grid'))
if tau is None or tau.size == 0:
    tau = np.array([0.0])
    res['tau_grid'] = tau
elif tau.size > 41:                      # bound the # of 1-D evaluations
    tau = np.linspace(float(tau.min()), float(tau.max()), 41)
    res['tau_grid'] = tau
\end{lstlisting}

Then it resolves the \term{base point} --- the fixed values the non-swept legs sit
at while one leg is swept. By default all non-swept legs sit at $\tau=0$, but
\code{Config.kpoint\_base\_lags} can move them, and its length is validated to be
exactly $k-1$ (\file{notebooks/daedalus.py:608--616}):

\begin{lstlisting}
base = cfg.kpoint_base_lags
if base is not None:
    base = [float(b) for b in base]
    if len(base) != k - 1:
        raise ValueError(
            f"kpoint_base_lags must have k-1 = {k-1} entries (one per "
            f"non-anchor leg); got {len(base)}.")
else:
    base = [0.0] * (k - 1)
\end{lstlisting}

Two small closures translate between the $(k-1)$ time differences and the
$k$-length time vector the callable wants. \code{\_args(vals)} pads leg~$0$ to zero
and fills legs $1..k-1$ from \code{vals}; \code{\_slice(fn, j)} sweeps leg $j$ over
the grid while the others sit at \code{base} (\file{notebooks/daedalus.py:618--630}):

\begin{lstlisting}
def _args(vals):                # vals: tau for legs 1..k-1 (leg 0 = 0)
    a = [0.0] * k
    for leg in range(1, k):
        a[leg] = float(vals[leg - 1])
    return a

def _slice(fn, j):              # sweep leg j over tau, others at base
    out = np.empty(tau.size, dtype=complex)
    for i in range(tau.size):
        v = list(base)
        v[j - 1] = float(tau[i])
        out[i] = complex(fn(*_args(v)))
    return out
\end{lstlisting}

\subsection{The axis-parallel slices (the default)}

The default view is the set of $k-1$ \term{axis-parallel slices}: for each leg
$j$, sweep $\tau_j$ over the grid while pinning the other legs at the base point,
\begin{equation}
  \mathrm{slice}_j(\tau)
  \;=\;
  \kappa_k\!\left(\tau_1=\text{base}_1,\;\ldots,\;\tau_j=\tau,\;\ldots,\;\tau_{k-1}=\text{base}_{k-1}\right).
\end{equation}
The code wraps the whole construction in \code{try/except Exception: pass}
(\file{notebooks/daedalus.py:632}, \file{:671--672}) so that any failure inside the
callables leaves the \code{None}/scalar form for the bar-chart plotter to handle
(see the gotcha below):

\begin{lstlisting}
try:
    tcb = {e: f for e, f in (res.get('total_C_by_ell') or {}).items()
           if callable(f)}
    # The k-1 axis-parallel slices through base.
    slices = {j: _slice(res['total_C'], j) for j in range(1, k)}
    slices_by_ell = {j: {e: _slice(f, j) for e, f in tcb.items()}
                     for j in range(1, k)}
    res['C_tau_slices'] = slices
    res['C_tau_slices_by_ell'] = slices_by_ell
    res['C_tau'] = slices[1]                    # canonical single slice
    res['C_tau_by_ell'] = slices_by_ell[1]
    res['_kpoint_base'] = base
    res['_kpoint_slice'] = (
        f'{k-1} slices: leg j swept over tau_grid (tau_j = t_j - t_0), '
        f'legs != j fixed at base={base}, for j = 1..{k-1}')
\end{lstlisting}

The crucial line is \code{res['C\_tau'] = slices[1]}: the first slice is promoted
to the canonical \code{C\_tau} key (and its per-$\ell$ version to
\code{C\_tau\_by\_ell}). That single move is what makes a $k\ge3$ run plot and
sim-compare \emph{exactly} like a $k=2$ run --- downstream code that reaches for
\code{C\_tau} finds a curve over \code{tau\_grid}, regardless of $k$. The full set
of slices is kept under \code{C\_tau\_slices} for the per-panel plotter.

\begin{note}
For a single, permutation-symmetric field all $k-1$ slices \emph{coincide} ---
sweeping leg $1$ gives the same curve as sweeping leg $2$, because the cumulant is
symmetric under exchange of identical legs. Seeing the panels overlap is thus a
free visual check of the cumulant's permutation symmetry, noted in the source at
\file{notebooks/daedalus.py:856--867}. If the legs name \emph{different} fields,
the slices genuinely differ and each panel carries distinct information.
\end{note}

The matching simulator estimates the same slice by freeing only the swept leg ---
\code{lag\_bins = [0, None, 0, ...]}, with \code{None} in the swept slot ---
documented at \file{notebooks/daedalus.py:588--594}.

\subsection{The full grid (opt-in)}

Setting \code{Config.kpoint\_full\_grid = True} additionally builds the whole
$(k-1)$-dimensional tensor $\kappa_k(\tau_1,\ldots,\tau_{k-1})$. Because the cost
is $n^{k-1}$, each axis is downsampled so the total number of evaluations stays
bounded near $4000$ (\file{notebooks/daedalus.py:650--670}):

\begin{lstlisting}
if cfg.kpoint_full_grid:
    import itertools
    n = min(tau.size, max(2, int(4000 ** (1.0 / (k - 1)))))
    gtau = (tau if n >= tau.size else
            np.linspace(float(tau.min()), float(tau.max()), n))

    def _grid(fn):
        out = np.empty((gtau.size,) * (k - 1), dtype=complex)
        for idx in itertools.product(range(gtau.size), repeat=k - 1):
            out[idx] = complex(fn(*_args([gtau[m] for m in idx])))
        return out

    res['C_tau_grid'] = _grid(res['total_C'])
    res['C_tau_grid_by_ell'] = {e: _grid(f) for e, f in tcb.items()}
    res['tau_axes'] = [gtau] * (k - 1)
    if gtau.size < tau.size:
        res['_kpoint_grid_note'] = (
            f'tau downsampled to {gtau.size} pts/axis '
            f'({gtau.size ** (k - 1)} evals) to bound grid cost')
\end{lstlisting}

The per-axis point count is \code{int(4000 ** (1.0 / (k - 1)))} --- the
$(k-1)$-th root of the evaluation budget --- floored at $2$ and capped at the grid
size. So for $k=3$ that is about $\sqrt{4000}\approx63$ points per axis (then
capped to the $41$-point $\tau$-grid), giving the $\kappa_3(\tau_1,\tau_2)$
heatmap; for $k=4$ it is about $4000^{1/3}\approx16$ points per axis. The
nested-loop \code{itertools.product(range(gtau.size), repeat=k-1)} walks every cell
of the $(k-1)$-dimensional grid. If downsampling happened, a human-readable
\code{\_kpoint\_grid\_note} records it.

\begin{gotcha}
The entire $k\ge3$ synthesis block is wrapped in
\code{try/except Exception: pass}. If evaluating \code{total\_C} fails for any
reason, the result silently retains the engine's \code{C\_tau=None}/scalar form,
and the dispatcher falls back to the equal-time bar-chart plotter
(\S\ref{sec:dispatch}) --- no error surfaces. This is deliberate (a coarse
equal-time number is better than a crash), but it can mask a genuine failure inside
the callable. If a $k\ge3$ plot unexpectedly shows bars instead of slices, a
swallowed exception in \code{total\_C} is the first thing to suspect.
\end{gotcha}

\begin{gotcha}
A $k\ge3$ plot is generally \emph{coarser} than the same theory's $k=2$ plot.
Slices cap the $\tau$-grid at $41$ points; the full grid downsamples each axis to
keep $n^{k-1}$ bounded. Both caps are deliberate cost controls, but they mean you
should not read fine structure off a $k\ge3$ figure that you would trust on a
two-point curve.
\end{gotcha}

\section{Dispatching the result}\label{sec:dispatch}

All of the above produces a result dict whose \emph{shape} encodes which case you
are in: a moment curve, a $k\ge3$ scalar, a set of slices, a heatmap, or an
ordinary curve. \code{plot\_cumulant} (\file{notebooks/daedalus.py:725}) reads that
shape and routes to the right plotter. The decision tree, in order
(\file{notebooks/daedalus.py:738--756}):

\begin{enumerate}
  \item \code{output\_kind} is a moment \emph{and} a \code{moment} array is present
        $\to$ \code{\_plot\_moment}. (Checked first, so a requested moment always
        wins over the cumulant shape it was assembled from.)
  \item \code{'C\_kpoint' in result} $\to$ \code{plot\_kpoint} (spatial $k\ge3$
        events).
  \item \code{is\_spatial(model)} or \code{'C\_tau\_x' in result} $\to$
        \code{plot\_spatial} (the $C(x,0)$ slice, optionally a $C(x,\tau)$ heatmap).
  \item \code{result['C\_tau'] is None} $\to$ \code{plot\_temporal\_kpoint} (the
        temporal $k\ge3$ equal-time scalar, drawn as per-order bars --- the fallback
        when slice synthesis was skipped or failed).
  \item \code{result['C\_tau\_grid'] is not None} $\to$
        \code{plot\_temporal\_kpoint\_grid} (the $k=3$ $\kappa_3(\tau_1,\tau_2)$
        heatmap).
  \item \code{len(result['C\_tau\_slices']) >= 2} $\to$
        \code{plot\_temporal\_kpoint\_slices} (one panel per $\tau_j$).
  \item otherwise $\to$ \code{plot\_temporal} (the ordinary $C(\tau)$ curve).
\end{enumerate}

The moment plotter \code{\_plot\_moment} (\file{notebooks/daedalus.py:759}) draws
the assembled moment as a \emph{single} curve --- there is no per-order overlay,
because a moment deliberately mixes loop orders (recall the shared-budget sum runs
over assignments of different $\ell_B$ to different blocks, so ``the one-loop
moment'' is not a sum of separable per-order curves). For a temporal model it plots
$M_k(\tau)$ over \code{tau\_grid}; for spatial $k=2$ it plots $M(x,0)$, selecting
the $\tau\approx0$ row with \code{argmin(|tau|)} if the moment array is
two-dimensional. The $y$-axis is labelled $M_k$, and it honours the usual
\code{logy}, \code{title}, and \code{save} options, overlaying \code{sim} as points
or error bars when supplied.

\begin{note}
The dispatch is on \emph{shape}, not on a stored ``mode'' flag. Because the
$k\ge3$ synthesis stamps \code{C\_tau} with slice~$1$ and (optionally)
\code{C\_tau\_grid} with the full tensor, the same dict can satisfy several
branches; the \emph{order} of the checks resolves it. A full-grid $k=3$ run, for
instance, has both \code{C\_tau\_slices} and \code{C\_tau\_grid} populated, and the
grid check (branch~5) comes before the slices check (branch~6), so it draws the
heatmap. A slices-only run skips branch~5 and lands on branch~6.
\end{note}

\section{Putting it together: a temporal \texorpdfstring{$k=3$}{k=3} run}

The data flow for a three-point temporal run ties the whole chapter together.
Starting from

\begin{lstlisting}
cfg = dd.Config(k=3, max_ell=1)
res = dd.run(model, cfg, mod)
\end{lstlisting}

\code{run} resolves $k=3$, auto-builds \code{external\_fields = [(f0, 1)] * 3}
(three copies of the first physical field's leg, since none was passed), and calls
the engine. The engine returns \code{C\_tau = None} together with the callables
\code{total\_C} and \code{total\_C\_by\_ell}. The synthesis block of
\S\ref{sec:kge3} then builds the $k-1=2$ axis-parallel slices
\code{C\_tau\_slices = \{1: ..., 2: ...\}}, promotes slice~$1$ to \code{C\_tau},
and (because \code{kpoint\_full\_grid} is left \code{False}) stops there. Finally
\code{plot\_cumulant} walks its decision tree: not a moment, no \code{C\_kpoint},
not spatial, \code{C\_tau} is \emph{not} \code{None} (we set it to slice~$1$),
\code{C\_tau\_grid} is absent, and \code{len(C\_tau\_slices) == 2 >= 2} --- so it
lands on branch~6, \code{plot\_temporal\_kpoint\_slices}, and draws two panels, one
per time difference.

Add \code{kpoint\_full\_grid=True} and the same run additionally fills
\code{C\_tau\_grid}; the dispatcher now hits branch~5 first and draws the
$\kappa_3(\tau_1,\tau_2)$ heatmap instead. Add \code{output='central\_moment'} and
\code{run} also calls \code{\_assemble\_moment\_temporal} --- which spends two extra
backend runs (orders $2$ and $3$, the order-$3$ run reused) to build the
shared-budget cluster sum --- and stamps \code{res['moment']} and
\code{res['output\_kind']}, so the dispatcher takes branch~1 and draws the single
moment curve. One \code{Config}, three different figures, all from the same pure
cumulant engine.

\section{Summary and caveats}

\begin{itemize}
  \item The engine emits \emph{connected cumulants} only; every moment and every
        $k\ge3$ view in this chapter is post-processing inside
        \file{notebooks/daedalus.py}, never inside \file{pipeline/compute.py}.
  \item Moments are reconstructed by the set-partition expansion
        \eqref{eq:moment-setpart}, evaluated under a single \emph{shared loop
        budget} $\sum_B\ell_B\le L$ \eqref{eq:moment-budget} so that no term
        exceeds the requested total loop order. Singletons contribute the
        \emph{tree} mean $\phi^{*}$ (raw) or are dropped (central); the one-loop
        tadpole shift of the mean is not folded in.
  \item A temporal moment of order $k$ costs up to $k-1$ extra
        \code{compute\_cumulants} runs (one per block size $2..k$); spatial $k=2$
        moments are free (one line from \code{C\_tau\_x}); spatial $k\ge3$ moments
        raise \code{NotImplementedError}.
  \item For $k\ge3$ the cumulant is a callable, not a curve. \code{run} synthesises
        $k-1$ axis-parallel slices (capped to $41$ $\tau$-points) and, on request,
        the full $(k-1)$-dimensional grid (downsampled to bound $n^{k-1}$). Slice~$1$
        is promoted to \code{C\_tau} so $k\ge3$ plots like $k=2$.
  \item The $k\ge3$ synthesis swallows exceptions and degrades to an equal-time bar
        chart; a $k\ge3$ figure is intentionally coarser than its $k=2$ counterpart.
  \item \code{plot\_cumulant} dispatches on result \emph{shape} in a fixed order:
        moment $\to$ spatial-$k\ge3$ $\to$ spatial-$k=2$ $\to$ temporal-$k\ge3$
        scalar $\to$ $k=3$ grid $\to$ $k\ge3$ slices $\to$ ordinary curve.
\end{itemize}

The next chapter turns to the simulators --- the independent stochastic
integrations that produce the \code{sim} overlays this chapter's plotters consume,
and against which every cumulant and moment in the manual is checked.

% ---------------------------------------------------------------------
%  Chapter 20 : Validation --- Independent Simulators & Cumulant Estimators
% ---------------------------------------------------------------------
\chapter{Validation: Independent Simulators and Cumulant Estimators}\label{ch:simulators}

Everything in the preceding chapters builds \emph{one} half of the repository: the
analytic pipeline that takes an \msrjd{} action, enumerates Feynman diagrams,
assigns propagators and vertices, and computes connected correlation functions
(cumulants) order by order in the loop expansion. This chapter documents the
\emph{other} half --- the half whose entire job is to not believe the first one.

The files described here are a set of completely independent numerical
simulators that solve the \emph{same} stochastic equations of motion directly,
by brute-force time-stepping, plus a set of estimators that turn the raw
simulated time series into empirical estimates of the very same cumulant slices
the pipeline predicts. The two are then plotted on top of each other. If the
diagrammatic machinery is correct, the curves overlap, within Monte-Carlo error
bars. This is the \term{trust chain}: a closed loop in which an analytic
prediction is checked against a numerical experiment that shares \emph{no code}
with the prediction. You met one link of that chain already, at the end of
Chapter~\ref{ch:quickstart}, where the two-loop value $C_{xx}(0) = 0.9496$ was
confirmed by a simulation that measured $0.9464 \pm 0.0009$. This chapter is the
full story behind that one line of output.

\begin{note}
\textbf{These files are not part of the pipeline.} They never run during a normal
cumulant computation. They exist solely so that a skeptical reader can confirm
the diagrams give the right answer. The repository's confidence in every analytic
result rests on these brute-force checks agreeing. When you read ``the pipeline''
elsewhere in this manual, the simulators are \emph{not} included; they are the
external referee.
\end{note}

The two primary files are
\file{models/ou\_langevin\_sim\_numba.py} (the Euler--Maruyama SDE integrators)
and \file{models/cumulant\_estimator.py} (the $k$-point connected-cumulant
estimator). A wider supporting cast --- the Hawkes point-process simulators, the
spatial-field simulators, and the coupled reaction--diffusion oracle --- is
characterised at the end of the chapter. We open every tool from the ground up:
unlike the pipeline, this subsystem touches no SageMath, no \code{nauty}, no
\code{sympy}. It is deliberately lean, and that cleanness is the whole point ---
the validators share no symbolic machinery with the code they validate.

\section{Why a separate validation layer?}

A number that comes out of a diagrammatic calculation is the end of a long chain
of symbolic manipulations: a Taylor expansion of the action, a graph enumeration,
an automorphism-counting symmetry factor, a multidimensional loop integral. A bug
anywhere in that chain produces a wrong number that still \emph{looks} like a
correlation function --- it is smooth, it decays, it is symmetric in $\tau$. There
is nothing in the output itself that flags the error. The only way to catch such
a bug is to compute the same quantity a \emph{second}, completely different way,
and compare.

The cheapest fully independent second computation for a stochastic differential
equation is to simply \emph{simulate} it: discretise time, take small steps,
add the right amount of random noise each step, and record the trajectory. Run it
long enough and the time-averaged statistics of that trajectory converge (by
ergodicity) to the stationary correlation functions the theory predicts. This
shares nothing with the diagrammatics --- there is no action, no graph, no
integral. A human has hand-coded the equation of motion as ordinary numeric
arithmetic. If both routes land on the same number, the probability that they
both contain a compensating bug is negligible, and the analytic result is
licensed.

The phrase ``within error bars'' is load-bearing. A simulation is a Monte-Carlo
estimate: it has statistical noise that shrinks like $1/\sqrt{N}$ in the number of
independent runs (or the length of one run), plus a few systematic biases from the
finite step size and the finite bin width that we will quantify below. The
comparison is therefore never an exact equality; it is ``the theory value sits
inside the simulation's confidence interval.'' When the example notebook reports
$0.9464 \pm 0.0009$ against a theory value of $0.9496$, the residual $0.003$ gap
is about three standard errors --- accounted for by the two-loop truncation plus
the estimator's finite-bin and finite-$T$ biases.

\section{The two simulator families}

There are two physically distinct kinds of dynamics the repository validates,
and correspondingly two families of simulator.

\begin{description}
  \item[Langevin / SDE simulators] integrate a stochastic \emph{differential}
        equation, $\dot x = f(x) + \text{noise}$, where the state $x$ is a
        continuous real variable (or field). These cover the OU-type theories
        (\file{ou\_langevin\_sim\_numba.py}), the spatial scalar fields
        (\file{spatial\_field\_*\_sim.py}), and the coupled reaction--diffusion
        systems (\file{coupled\_rd\_1d\_sim.py}). The workhorse is the
        Euler--Maruyama scheme; stiffer linear parts use an exact
        Ornstein--Uhlenbeck step, a spectral exponential-Euler step (ETD1), or a
        semi-implicit Crank--Nicolson step.
  \item[Point-process / Hawkes simulators] integrate a self-exciting Poisson
        point process (\file{hawkes\_sim\_*numba.py}). The state is a discrete
        spike count: each timestep draws $\mathrm{Poisson}(\lambda_i\,\mathrm{d}t)$
        spikes and feeds them back into the rate. This is a discrete-event
        process, not an SDE, and it is what makes the \emph{factorial-cumulant}
        machinery of \S\ref{sec:factorial} necessary.
\end{description}

The same estimator file consumes both: a simulator returns a binned trajectory
of shape \code{(npop, n\_bins)}, and the estimator cross-correlates that to
produce a cumulant slice. Which leg is treated as a smooth continuous field and
which as a discrete spike train is selected by a per-leg \emph{field-type} label
(\S\ref{sec:fieldtypes}).

\subsection{Where it sits in the end-to-end pipeline}

The wiring is a diamond: one theory file feeds \emph{both} the pipeline and the
simulator, and the two output streams meet in a comparison notebook.

\begin{center}
\begin{tikzpicture}[
    box/.style={draw,rounded corners,align=center,font=\small,inner sep=5pt},
    >=Stealth, node distance=7mm]
  \node[box] (theory) {theory file\\ (action $S$, params)};
  \node[box, right=28mm of theory] (pipe) {\textsc{pipeline}\\ enumerate\\$\to$ assign $\to$ integrate};
  \node[box, below=16mm of pipe] (sim) {\textsc{simulator}\\ Euler--Maruyama\\ / Poisson};
  \node[box, right=24mm of pipe] (res) {\code{res['C\_tau']}\\ analytic cumulants};
  \node[box, right=24mm of sim] (est) {\textsc{estimator}\\ sim cumulant\\ slices};
  \node[box, right=22mm of res, yshift=-8mm] (cmp) {sim-compare\\ notebook\\ (overlay plot)};
  \draw[->] (theory) -- (pipe);
  \draw[->] (theory.south) |- (sim.west);
  \draw[->] (pipe) -- (res);
  \draw[->] (sim) -- (est);
  \draw[->] (res.east) -- (cmp.north west);
  \draw[->] (est.east) -- (cmp.south west);
\end{tikzpicture}
\end{center}

Three facts about that diagram matter for the rest of the chapter:

\begin{itemize}
  \item \textbf{The simulator is handed the pipeline's resolved parameters.} In
        the sim-compare notebooks the call is literally
        \code{fp = res['\_resolved']['parameters']} followed by
        \code{mu = float(fp['mu'])}, and so on. The simulator gets the
        \emph{same} numbers the theory used, cast to plain Python floats. It does
        \emph{not} read the action --- a human has hand-coded each equation of
        motion to match a specific theory file, and the docstring of each
        \code{sim\_*} function names the matching \file{theories/*.theory.py}.
  \item \textbf{The simulator output feeds the estimator.} A simulator returns a
        binned trajectory array of shape \code{(npop, n\_bins)}; the estimator
        cross-correlates that to produce a cumulant slice \code{(n\_tau,)} or a
        stack \code{(k-1, n\_tau)}.
  \item \textbf{The estimator output feeds the comparison plot.} The notebooks
        overlay the simulation's \code{sim['C']} against the theory's
        \code{res['C\_tau\_by\_ell']} (per loop order) or \code{res['C\_tau']}
        (loop-summed). \emph{That overlap is the trust-chain verdict.}
\end{itemize}

\section{Euler--Maruyama from scratch}

The scalar simulators integrate an It\^o stochastic differential equation
\begin{equation}
  \dot x \;=\; f(x) \;+\; \sqrt{2D}\,\eta(t),
  \qquad
  \avg{\eta(t)\,\eta(t')} \;=\; \delta(t-t'),
  \label{eq:sde}
\end{equation}
where $f(x)$ is the deterministic \term{drift} and $\eta$ is unit-strength
Gaussian white noise. For the flagship example --- the quartic double well in
\code{sim\_ou\_quartic\_numba} --- the drift is
\[
  f(x) \;=\; -\mu x \;-\; \varepsilon x^{3},
\]
so the equation of motion is $\dot x = -\mu x - \varepsilon x^{3} + \sqrt{2D}\,
\eta$. This is exactly the \msrjd{} action quoted in the docstring
(\file{ou\_langevin\_sim\_numba.py:29}),
\[
  S \;=\; \int \dd t\; \tilde x\,\bigl((\Dt + \mu)\,x + \varepsilon x^{3}\bigr)
        \;-\; D\,\tilde x^{2},
\]
where $\tilde x$ (written \code{xt} in the theory text) is the \msrjd{} response
field. The $-D\,\tilde x^{2}$ term is the noise vertex; its coefficient $D$ fixes
the noise strength $\avg{\xi\,\xi} = 2D\,\delta(\tau)$ with $\xi = \sqrt{2D}\,\eta$.

\subsection{The defining feature: $\sqrt{\dd t}$ noise scaling}

\term{Euler--Maruyama} is the stochastic analogue of the forward-Euler scheme for
ordinary differential equations. Over a step $\dd t$ the deterministic part
advances by $\dd t\,f(x)$, exactly as in a deterministic Euler step. The new
ingredient is the integrated noise increment $\int_t^{t+\dd t}\sqrt{2D}\,\eta\,
\dd t'$, which is a Gaussian with mean $0$ and variance $2D\,\dd t$. (The variance
follows from $\mathrm{Var}\!\left[\int \eta\right] = \int\!\!\int \delta = \dd t$:
the double integral of the delta-correlation over the step collapses to the step
length.) Therefore one Euler--Maruyama step is
\begin{equation}
  x(t+\dd t) \;=\; x(t) \;+\; \dd t\, f(x) \;+\; \sqrt{2D\,\dd t}\;N(0,1).
  \label{eq:em}
\end{equation}
The crucial detail is the $\sqrt{\dd t}$ scaling of the noise term --- not $\dd t$.
A deterministic Euler step would scale \emph{everything} by $\dd t$; the noise
must scale by $\sqrt{\dd t}$ because variance is what adds linearly in time, so
the standard deviation grows like $\sqrt{\dd t}$. This single fact is the whole
difference between Euler and Euler--Maruyama.

In code, the noise prefactor is precomputed once outside the loop
(\file{ou\_langevin\_sim\_numba.py:74}) and the step itself is three lines
(\file{ou\_langevin\_sim\_numba.py:80--82}):

\begin{lstlisting}
sqrt_2D_dt = np.sqrt(2.0 * D * dt_sim)
...
x = (x
     + dt_sim * (-mu * x - eps * x * x * x)
     + sqrt_2D_dt * np.random.randn())
\end{lstlisting}

\code{np.random.randn()} draws one $N(0,1)$ sample; multiplying by
\code{sqrt\_2D\_dt} turns it into the correctly-scaled noise increment. Note the
cubic is written \code{eps * x * x * x} rather than \code{eps * x**3}: repeated
multiplication is cheaper than a power call inside a hot loop, and (as we will see)
this code runs millions of times per simulation.

\subsection{Weak order one and the $\dd t \ll 1/\mu$ rule}

Euler--Maruyama is a \term{weak order-1.0} scheme. ``Weak'' means we care about the
accuracy of \emph{averages} and \emph{moments} of $x$ (which is what a correlation
function is), not about the pathwise accuracy of any single trajectory. ``Order
$1.0$'' means the bias in those stationary statistics scales as $O(\dd t)$ ---
halve the step, halve the bias. The practical consequence, stated in the docstring
(\file{ou\_langevin\_sim\_numba.py:36}), is that the step must be small relative to
the relaxation time $1/\mu$: $\dd t_{\mathrm{sim}} \ll 1/\mu$. If you step too
coarsely, the simulation systematically misrepresents the fast relaxation of $x$
and biases every measured statistic.

\begin{defn}
\textbf{Stationary variance sanity check.} For the \emph{pure} OU process
($\varepsilon = 0$, $f = -\mu x$) the stationary distribution is Gaussian with
variance $D/\mu$, so the equal-time correlation is $C_{xx}(0) = \avg{x^{2}} =
D/\mu$. This is the analytic target at tree level; a finite $\varepsilon$ shifts it
by $O(\varepsilon D^{2}/\mu^{3})$ (\file{ou\_langevin\_sim\_numba.py:64--66}). It is
exactly the number the pipeline's loop expansion should reproduce, and a quick way
to confirm a simulation is behaving: set $\varepsilon = 0$ and check the measured
variance lands on $D/\mu$.
\end{defn}

The sign conventions follow the theory file
(\file{ou\_langevin\_sim\_numba.py:60--63}): positive $\mu$ gives a stable origin
(sub-critical regime), $\mu = 0$ is a pitchfork bifurcation, and $\mu < 0$ is a
genuine double well with minima at $x = \pm\sqrt{|\mu|/\varepsilon}$. The
equilibrium Boltzmann potential is $U(x) = \tfrac{\mu}{2}x^{2} +
\tfrac{\varepsilon}{4}x^{4}$ (and $+\tfrac{\gamma}{6}x^{6}$ for the sextic
variant), bounded below --- hence stable and normalisable --- when all coefficients
are positive.

\subsection{The bin accumulator}

All five functions in \file{ou\_langevin\_sim\_numba.py} share one skeleton: seed
the random generator, initialise the state, loop \code{n\_steps} doing an Euler
step plus a bin accumulation, and emit a bin-averaged trajectory. The simulation
runs at a fine step \code{dt\_sim} but records output only every
\code{bin\_size\_steps} steps, averaging $x$ over that window into one output bin
(\file{ou\_langevin\_sim\_numba.py:84--89}):

\begin{lstlisting}
steps_in_bin += 1
if steps_in_bin >= bin_size_steps:
    if cur_bin < n_bins:
        x_bins[0, cur_bin] = accum / bin_size_steps
        accum = 0.0
    cur_bin += 1
    steps_in_bin = 0
\end{lstlisting}

The guard \code{if cur\_bin < n\_bins} stops writing once the output buffer is
full; any extra Euler steps beyond \code{bin\_size\_steps}\,$\times$\,\code{n\_bins}
are simply ignored. The output shape is always \code{(npop, n\_bins)}, matching
the Hawkes \code{voltage\_bins} convention, so the estimator can consume an OU
trajectory with \code{field\_types=['dv', ...]} exactly as it consumes a Hawkes
voltage trace.

\begin{gotcha}
\textbf{Bin finely, or bias $C_{xx}(0)$ low.} A coarse bin (one whose width
\code{dt\_bin} is not $\ll 1/\mu$) averages the field over the bin and smooths away
the fastest fluctuations, so the bin-averaged variance \emph{underestimates} the
instantaneous $C_{xx}(0)$ the theory reports. The example notebook sets
\code{dt\_bin}\,$= 0.02$ for $\mu \sim 1$ precisely to keep this bias below the
statistical error. This is a separate requirement from the $\dd t_{\mathrm{sim}}
\ll 1/\mu$ rule: one governs the integration step, the other the recording window.
\end{gotcha}

\section{Beyond plain Euler: removing bias from the linear part}

Plain Euler--Maruyama carries an $O(\dd t)$ bias in \emph{everything}, including
the linear relaxation. Three of the simulator families improve on this by treating
the \emph{linear} part of the dynamics exactly, leaving only the genuinely
nonlinear pieces (and, for colored noise, the $x$ step) at $O(\dd t)$. Each trick
removes a specific bias.

\subsection{Exact Ornstein--Uhlenbeck discretisation (colored noise)}
\label{sec:colored}

Some theories use temporally correlated (\term{colored}, Lorentzian) noise rather
than white noise. \code{sim\_ou\_quartic\_colored\_numba} realises this by running
a \emph{second} SDE --- an auxiliary Ornstein--Uhlenbeck process $\xi$ --- and using
its output as the forcing of $x$:
\begin{align}
  \dot\xi &= -\xi/\tau_c + (\sqrt{2D}/\tau_c)\,\eta(t)
            &&\Longrightarrow&
            \avg{\xi(t)\,\xi(t')} &= (2D/\tau_c)\,\ee^{-|t-t'|/\tau_c}, \\
  \dot x  &= -\mu x - \varepsilon x^{3} + \xi(t). &&&&
\end{align}
The key trick is that a \emph{linear} OU SDE has a closed-form one-step update ---
it does not need Euler--Maruyama at all. The exact update is
\begin{equation}
  \xi(t+\dd t) \;=\; \text{decay}\cdot\xi(t) \;+\; \sigma\,N(0,1),
  \qquad
  \text{decay} = \ee^{-\dd t/\tau_c},
  \qquad
  \sigma^{2} = (2D/\tau_c)(1 - \text{decay}^{2}).
  \label{eq:exactou}
\end{equation}
This preserves the stationary variance for \emph{any} ratio $\dd t/\tau_c$: there
is no $O(\dd t)$ bias in the noise statistics at all, only in the $x$ Euler step.
In code (\file{ou\_langevin\_sim\_numba.py:146--148}):

\begin{lstlisting}
decay = np.exp(-dt_sim / tauc)
var_factor = 1.0 - decay * decay
sigma_xi = np.sqrt((2.0 * D / tauc) * var_factor)
\end{lstlisting}

A subtle implementation detail: at each step the code draws the new noise
\code{xi\_new} but steps $x$ with the \emph{previous}-step \code{xi}, then commits
both (\file{ou\_langevin\_sim\_numba.py:160--166}):

\begin{lstlisting}
xi_new = decay * xi + sigma_xi * eta
drift = -mu * x - eps * x * x * x
x_new = x + dt_sim * (drift + xi)   # NB: previous-step xi, not xi_new
x = x_new
xi = xi_new
\end{lstlisting}

\begin{gotcha}
\textbf{Colored-noise factor of 2.} In the white limit $\tau_c \to 0$ this colored
driver tends to $\avg{\xi\xi} \to 4D\,\delta(\tau)$ --- note \textbf{4D}, not 2D ---
because of the $2D/\tau_c$ coefficient convention. To compare the colored
simulator against the white simulator you must use $D_{\mathrm{white}} = 2\,
D_{\mathrm{colored}}$ (equivalently, halve $D$ when crossing over). The docstring
spells this out (\file{ou\_langevin\_sim\_numba.py:127--132}); forgetting it
produces a clean-looking curve that is wrong by a constant factor.
\end{gotcha}

\subsection{Spectral exponential-Euler (ETD1) for stiff spatial fields}

The spatial simulators (\file{spatial\_field\_1d\_sim.py} and its 2D/3D siblings)
integrate a \emph{field} SDE on a periodic ring or torus, e.g.\ the
non-conserved $\phi^{4}$ dynamics
\[
  \Dt \phi(x,t) \;=\; -\mu\phi + D\,\partial_x^{2}\phi - \lambda\phi^{3} + \eta(x,t),
  \qquad
  \avg{\eta(x,t)\,\eta(x',t')} = 2T\,\delta(x-x')\,\delta(t-t').
\]
A naive Euler--Maruyama in real space is unstable here because the diffusion
operator $D\,\partial_x^{2}$ makes the system \term{stiff}: its largest eigenvalue
grows like $D\,k_{\max}^{2}$, and an explicit step must resolve that fastest mode
or blow up. The fix is \term{ETD1} (exponential time-differencing, order 1),
working in Fourier space:

\begin{itemize}
  \item Each Fourier mode $\hat\phi_k$ obeys a \emph{linear} OU SDE with rate
        $\omega_k = \mu + D k^{2}$ (the lattice version is $\omega_k = \mu +
        (2D/\dd x^{2})(1 - \cos(k\,\dd x))$), which is integrated \emph{exactly}
        with the same OU closed form as \eqref{eq:exactou}. Consequently the
        $\lambda = 0$ stationary spectrum is unbiased in $\dd t$.
  \item The nonlinear forcing $-\lambda\phi^{3}$ is added through the ETD1
        integrating factor $(1 - \ee^{-\omega\,\dd t})/\omega$
        (\file{spatial\_field\_1d\_sim.py:77}).
\end{itemize}

The per-mode stationary variance is $\avg{|\hat\phi_k|^{2}} = T N^{2}/(L\,\omega_k)$
(\file{spatial\_field\_1d\_sim.py:78}), and the per-step OU increment carries
variance $(1 - \ee^{-2\omega\,\dd t})$ times that
(\file{spatial\_field\_1d\_sim.py:79}). Real-mode versus interior-complex-mode
bookkeeping (\file{spatial\_field\_1d\_sim.py:83--87} and~\code{:114--118})
handles the Hermitian structure of the real FFT. Derivative vertices --- Model~B's
conserved $g_{\mathrm{lap}}\,\partial_x^{2}(\phi^{2})$, the Burgers
$-(\lambda/2)\partial_x(\phi^{2})$, the KPZ $+(\lambda/2)(\partial_x\phi)^{2}$ ---
enter through spectral multipliers \code{lap\_eig} and \code{ik\_eig}
(\file{spatial\_field\_1d\_sim.py:70--75} and~\code{:98--107}). The spatial
pipeline chapters describe the matching theory side; here it is enough to know the
simulator realises the \emph{lattice} dispersion, a point we return to in the
caveats.

\subsection{Semi-implicit Crank--Nicolson (coupled reaction--diffusion)}

The coupled-RD simulator (\file{coupled\_rd\_1d\_sim.py}) integrates an
$N$-species reaction--diffusion system with cross-noise and unequal diffusion
constants. Its default scheme is \term{semi-implicit Crank--Nicolson}: a
trapezoidal-rule implicit step that, in this linear-stochastic setting, gives the
stationary covariance \emph{exactly to} $O(\dd t^{2})$ via a Pad\'e(1,1)
propagator (\file{coupled\_rd\_1d\_sim.py:30--44}). This is the most accurate of
the three improvements for the linear part. We return to this simulator, and its
\emph{analytic} oracle, in \S\ref{sec:coupled}.

\section{Correlated and multi-field noise (Cholesky)}

The two-field correlated simulators
(\code{sim\_ou\_quartic\_two\_dim\_corr\_numba} and
\code{sim\_ou\_quartic\_two\_dim\_color\_corr\_numba}) drive two fields with
jointly Gaussian white noise of correlation $\rho$:
\[
  \avg{\eta_\alpha(t)\,\eta_\beta(t')} \;=\; C_{\alpha\beta}\,\delta(t-t'),
  \qquad
  C \;=\; \begin{pmatrix} 1 & \rho \\ \rho & 1 \end{pmatrix}.
\]
Correlated Gaussians are produced by \term{Cholesky factorization} of the
covariance matrix $C$: writing $C = L L^{\mathsf T}$ with $L$ lower-triangular,
multiplying a vector of i.i.d.\ unit Gaussians by $L$ yields a Gaussian vector
with covariance exactly $C$. For the $2\times 2$ case the Cholesky factor is
\[
  L \;=\; \begin{pmatrix} 1 & 0 \\ \rho & \sqrt{1-\rho^{2}} \end{pmatrix},
\]
so the correlated drivers are (\file{ou\_langevin\_sim\_numba.py:351--352}):

\begin{lstlisting}
u = np.random.randn()
v = np.random.randn()
eta_x = u
eta_y = rho * u + rho_perp * v   # rho_perp = sqrt(1 - rho^2)
\end{lstlisting}

The correlation must satisfy $\rho \in [-1, 1]$ strictly; outside that range $C$
loses positive-semidefiniteness and no real Cholesky factor exists. The code
clamps $1-\rho^{2}$ to $\ge 0$ defensively
(\file{ou\_langevin\_sim\_numba.py:332--335}):

\begin{lstlisting}
one_minus_rho_sq = 1.0 - rho * rho
if one_minus_rho_sq < 0.0:
    one_minus_rho_sq = 0.0
rho_perp = np.sqrt(one_minus_rho_sq)
\end{lstlisting}

\begin{gotcha}
\textbf{An out-of-range $\rho$ yields NaN, not an error.} The clamp keeps the
square root from raising, but an out-of-range correlation then silently propagates
\code{NaN} through the trajectory rather than crashing
(\file{ou\_langevin\_sim\_numba.py:332--335}). If a two-field comparison plot comes
back empty, an out-of-bounds $\rho$ is the first thing to check.
\end{gotcha}

\section{The numba/Sage boundary}
\label{sec:numbasage}

This section is the single most important tooling fact in the subsystem. The
simulators are decorated with \code{numba}, and they live in plain \code{.py}
files for one specific reason that interacts with the SageMath kernel. To explain
it we first explain numba.

\subsection{What numba is}

\term{Numba} is a just-in-time (JIT) compiler for a subset of Python and NumPy. The
decorator \code{@numba.njit} (``no-python JIT'') compiles the decorated function to
native machine code, via the LLVM compiler infrastructure, the \emph{first} time
the function is called, then caches that compiled version for all subsequent
calls. The \code{njit} mode (also written \code{nopython=True}) forbids any
fallback to the slow Python interpreter: if numba cannot infer a concrete machine
type for every variable, it raises a compilation error rather than silently
running slowly. For the tight scalar arithmetic in an Euler step --- a loop of
millions of iterations of $x = x + \dd t\cdot\text{drift} + \sqrt{\dd t}\cdot
\text{noise}$ --- this buys roughly a $100\times$ speedup over pure Python (the
Hawkes header quotes ``\textasciitilde 15M steps/sec on Apple M1'',
\file{hawkes\_sim\_numba.py:37}).

Every Euler-step simulator is one \code{@numba.njit} function with a plain
scalar/array signature (\file{ou\_langevin\_sim\_numba.py:17--21}):

\begin{lstlisting}
@numba.njit
def sim_ou_quartic_numba(n_steps, dt_sim, mu, eps, D, x_init,
                         bin_size_steps, n_bins, seed):
\end{lstlisting}

\subsection{Why these must be \texttt{.py} files: the Sage-Integer trap}

When a notebook runs under the \term{SageMath kernel}, Sage's \term{preparser}
rewrites the source of every notebook cell before executing it. An integer literal
\code{0} becomes \code{Integer(0)} (a SageMath ring element), and \code{1.5}
becomes \code{RealNumber('1.5')}. These are \emph{not} Python \code{int}/\code{float};
they are objects from Sage's algebra system, carrying exact-arithmetic semantics.
Numba's type inference cannot handle them --- it only knows native machine types.
The docstring of the Hawkes simulator states the rule directly
(\file{hawkes\_sim\_numba.py:5--11}):

\begin{quote}\itshape
This file MUST be a plain \code{.py} file (not executed through SageMath's
preparser) because SageMath's preparser converts integer literals like \code{0}
to \code{Integer(0)}\,\ldots which Numba's type inference cannot handle.
\end{quote}

Two disciplines follow, and they are enforced throughout the subsystem:

\begin{enumerate}
  \item \textbf{The simulators live in \code{.py} modules, never in notebook
        cells.} Only \emph{cell source} is preparsed; imported \code{.py} files
        are read as ordinary Python. So \code{from models.ou\_langevin\_sim\_numba
        import sim\_ou\_quartic\_numba} is safe --- the function body never passes
        through the preparser.
  \item \textbf{Every argument at the call site must be cast to a plain Python
        \code{int}/\code{float}} before the \code{njit} function sees it. This is
        why the notebooks read \code{mu = float(fp['mu'])} rather than passing the
        resolved parameter straight through. The non-numba spatial simulators
        defend the same way \emph{inside} the function body --- \code{N = int(N);
        n\_steps = int(n\_steps); ...} (\file{spatial\_field\_1d\_sim.py:141--145})
        --- because \code{np.random.default\_rng} also rejects a Sage
        \code{Integer}.
\end{enumerate}

\begin{gotcha}
\textbf{An uncast Sage scalar crashes numba.} This is not stylistic. If you pass a
Sage \code{Integer} or \code{RealNumber} into an \code{@numba.njit} function you
get a numba typing error at \emph{compile} time; if you pass one into
\code{np.random.default\_rng} (the spatial path) you get a runtime rejection.
Reading parameters off \code{res['\_resolved']['parameters']} and casting with
\code{float(...)} / \code{int(...)} also guarantees the simulation uses
\emph{exactly} the values the theory used.
\end{gotcha}

\subsection{The two random-number APIs}

A small but real difference between the numba and non-numba simulators is which
NumPy random API they use. The numba-jitted Langevin and Hawkes simulators use the
\emph{legacy global} generator: \code{np.random.seed(seed)}
(\file{ou\_langevin\_sim\_numba.py:68}) seeds it, then \code{np.random.randn()}
(\file{ou\_langevin\_sim\_numba.py:82}) draws a single $N(0,1)$ and
\code{np.random.poisson(lam)} draws spike counts. The numba runtime supports this
legacy interface but not the modern generator object. The spatial simulators, which
are \emph{not} jitted, use the modern \code{np.random.default\_rng(seed)} generator
(\file{spatial\_field\_1d\_sim.py:154}) with \code{rng.standard\_normal(M)}, because
they are free to use the better API --- and, as noted above, that very API is what
forces the explicit \code{int(...)} casts inside their function bodies.

\subsection{JIT warmup}

Because the first call to an \code{njit} function pays the one-time compilation
cost, the notebooks make a throwaway warmup call before the timed measurement
loop, so the compilation is not counted in the timing:

\begin{lstlisting}
_ = sim_ou_quartic_numba(1000, dt_sim, mu, eps, D, 0.0,
                         bin_size_steps, 100, 0)   # JIT warmup
\end{lstlisting}

The arguments are throwaway (1000 steps, 100 bins, seed 0); the only purpose is to
trigger compilation. Every subsequent call in the run loop uses the cached native
code.

\subsection{What this subsystem does \emph{not} use}

For calibration against the rest of the manual: this subsystem uses \emph{no}
SageMath (it only avoids it), \emph{no} \code{nauty}/\code{networkx} (there is no
graph enumeration --- there are no diagrams here), and \emph{no} \code{sympy}
(there is no symbolic algebra --- the equations of motion are hand-coded as numeric
arithmetic). The two primary files touch only \code{numpy}, \code{numba}, and the
standard-library \code{collections.Counter}; the wider \file{models/} directory
additionally uses \code{scipy.linalg} in exactly one place
(\S\ref{sec:coupled}). That cleanness is the point.

\section{The cumulant estimator: factorial moments}
\label{sec:factorial}

We now turn to \file{models/cumulant\_estimator.py}, which converts a binned
trajectory into an empirical cumulant slice. The mathematical heart of the file is
the \emph{factorial moment}, and to motivate it we must understand a problem that
arises only for discrete point-process (spike-train) data.

\subsection{The shot-noise diagonal}

For point-process data, the raw $k$-point correlation contains delta-function
\term{contact terms}: a spike is perfectly correlated with itself. Concretely, the
same-bin product of counts $n^{2}$ overcounts --- it includes ``a spike paired with
itself,'' a shot-noise term scaling like $1/\dd t_{\mathrm{bin}}$ that is \emph{not}
part of the smooth connected correlator the diagrams compute. If you naively
estimated $\avg{n^{2}}$ you would measure the true covariance \emph{plus} this
spurious diagonal spike.

The fix is the \term{falling factorial}. Replace the ordinary power $n^{m}$ (at $m$
coincident same-population fields in the same bin) with
\begin{equation}
  n^{(m)} \;=\; n\,(n-1)\,(n-2)\cdots(n-m+1).
  \label{eq:falling}
\end{equation}
For a Poisson variable the identity $\avg{n^{(m)}} = \lambda^{m}$ holds
\emph{exactly} --- the falling factorial strips off precisely the self-coincidence
terms. The resulting estimator targets the \term{factorial cumulant density}: the
smooth, off-diagonal part of the connected $k$-point cumulant, with the
$\delta$-function diagonals removed (\file{cumulant\_estimator.py:8--31}). For $k=2$
this reduces to the ordinary rate covariance, since the only diagonal is the shot
noise $\delta(\tau)$ (\file{cumulant\_estimator.py:13--16}). The reference is Daley
\& Vere-Jones, \emph{An Introduction to the Theory of Point Processes}, Ch.~5 on
factorial moment measures (\file{cumulant\_estimator.py:55--57}).

The falling factorial is built elementwise over a whole bin array
(\file{cumulant\_estimator.py:63--81}):

\begin{lstlisting}
def falling_factorial_array(n_arr, m):
    result = np.ones_like(n_arr, dtype=float)
    for j in range(m):
        result *= (n_arr - j)
    return result
\end{lstlisting}

For $m=1$ this is just $n$ itself (linear centering, no correction); for $m \ge 2$
it removes the self-coincidence shot noise.

\subsection{Rate conversion}

A spike-counting field's value $n$ is a \emph{count} per bin. To convert it to a
rate (units of Hz) it is divided by the bin width, $n/\dd t_{\mathrm{bin}}$. This is
the \code{denom = dt\_bin} that appears throughout the estimator. Smooth voltage
fields (field type \code{'dv'}) are \emph{not} counts and use \code{denom = 1.0}
(no rate conversion; units of volts). The choice is made by a small closure
(\file{cumulant\_estimator.py:218--221}):

\begin{lstlisting}
def _centering_denom(leg_idx):
    return dt_bin if _is_spike_ft(field_types[leg_idx]) else 1.0
\end{lstlisting}

\section{The estimator algorithm, step by step}

The core routine is \code{compute\_kpoint\_slice}
(\file{cumulant\_estimator.py:84}). Its signature exposes the full machinery:

\begin{lstlisting}
def compute_kpoint_slice(binned_counts, dt_bin, pop_indices, lag_bins,
                         max_lag_bins, n_fft=None,
                         field_types=None, voltage_bins=None,
                         ext_binned_counts=None):
\end{lstlisting}

The idea is: hold all but one of the $k$ external legs at fixed time lags, sweep
the remaining leg's lag across a window, and report the connected cumulant as a
function of that one sweep lag. The leg that sweeps is the \term{sweep axis}; the
others are pinned. We walk the ten steps in order.

\paragraph{1. Validate the sweep axis} (\file{cumulant\_estimator.py:139--145}).
Exactly one entry of \code{lag\_bins} must be \code{None}; that index is recorded
as \code{sweep\_idx}. An assertion enforces this.

\paragraph{2. Normalize field types} (\file{cumulant\_estimator.py:148--202}).
Each leg has a field type; the default is \code{'dn'}. The routine raises if a
\code{'dv'} leg is requested without \code{voltage\_bins}, or a \code{'dm'} leg
without \code{ext\_binned\_counts}, and maps natural names through the nested
\code{\_normalize\_ft} helper. We discuss the field-type system and the
bug \code{\_normalize\_ft} fixed in \S\ref{sec:fieldtypes}.

\paragraph{3. Define helpers} (\file{cumulant\_estimator.py:204--221}).
\code{\_data\_for(leg)} picks the counts/voltage/external array by field type;
\code{\_is\_spike\_ft} is \code{True} for \code{'dn'}/\code{'dm'}; and
\code{\_centering\_denom} returns \code{dt\_bin} for spike fields, \code{1.0} for
voltage (shown above).

\paragraph{4. Determine the valid window} (\file{cumulant\_estimator.py:225--240}).
From the fixed lags (plus an implicit reference at lag 0) the routine computes
\code{valid\_start = max(0, -min\_lag)}, \code{valid\_end = n\_bins - max(0,
max\_lag)}, and \code{n\_valid = valid\_end - valid\_start}. This is the set of
reference bins for which \emph{every} fixed leg still lands inside the array --- no
wraparound. It raises if \code{n\_valid <= 2*max\_lag\_bins}, i.e.\ if the
recording is too short for the requested lag window.

\paragraph{5. Means} (\file{cumulant\_estimator.py:242--251}). The mean of each
\code{(pop, field\_type)} combination is computed \emph{on the valid window}, for
consistency with the product that follows. These means are what we subtract to
center each field.

\paragraph{6. Build the fixed-leg product} (\file{cumulant\_estimator.py:261--300}).
The fixed legs are grouped by lag value; within each lag, legs are grouped by
\code{(pop, ft)} via a \code{Counter} (\S\ref{sec:counter}). For multiplicity
\code{m == 1} (or any voltage leg) the field is centered linearly,
$(n - \text{mean})/\text{denom}$. For spike fields with \code{m >= 2} the falling
factorial is applied and then centered:

\begin{lstlisting}
if m == 1 or ft == 'dv':
    factor = (n_arr - mean_n) / denom
    if m > 1:
        factor = factor ** m          # voltage power, no factorial
else:
    raw = falling_factorial_array(n_arr, m) / denom**m
    factor = raw - raw.mean()         # spike: factorial, then centered
product *= factor
\end{lstlisting}

Note the asymmetry: a \emph{voltage} field at multiplicity $m>1$ is raised to the
power $m$ (it is smooth, with no shot noise to remove), whereas a \emph{spike}
field at $m \ge 2$ gets the factorial correction.

\paragraph{7. Build the sweep series} (\file{cumulant\_estimator.py:302--307}).
The swept field is centered the same way: $(\text{sweep} - \text{mean})/
\text{denom}$.

\paragraph{8. Cross-correlate by zero-padded FFT}
(\file{cumulant\_estimator.py:309--320}). The sweep-axis cross-correlation
$C(L) = \sum_t \text{product}(t)\cdot\text{sweep}(t+L)$ is computed with a
zero-padded FFT, using the convolution theorem (a cross-correlation is an inverse
FFT of one transform times the conjugate of the other):

\begin{lstlisting}
F_product = np.fft.fft(product, n=n_fft)
F_sweep   = np.fft.fft(sweep_rate, n=n_fft)
raw_xcorr = np.fft.ifft(F_product * np.conj(F_sweep)).real
\end{lstlisting}

The FFT must be zero-padded to \code{n\_fft >= 2*n\_valid} (the default is the next
power of two $\ge 2\,n_{\mathrm{valid}}$, \file{cumulant\_estimator.py:310--311})
so that the \emph{circular} convolution the FFT computes coincides with the desired
\emph{linear} correlation over the valid window, with no periodic-wraparound
contamination.

\begin{gotcha}
\textbf{The FFT lag indexing is fragile.} The cross-correlation lag is read as
\code{raw\_xcorr[(-lag) \% n\_fft]}. The comment in the code is load-bearing
(\file{cumulant\_estimator.py:316--319}): \code{ifft(F\_product * conj(F\_sweep))[L]}
equals $\sum_t \text{product}(t)\cdot\text{sweep}(t-L)$, so the value we want at
$+\text{lag}$ lives at $L = -\text{lag}$. The comment warns explicitly:
\textbf{``Do NOT use \texttt{[::-1]} --- that shifts the index by 1 bin.''} A
reversed-array shortcut introduces a one-bin offset that is easy to miss and
poisons the whole slice.
\end{gotcha}

\paragraph{9. Extract and normalize} (\file{cumulant\_estimator.py:322--332}). For
each lag in $[-\code{max\_lag\_bins}, +\code{max\_lag\_bins}]$ the raw correlation
is divided by the number of \emph{overlapping} valid bin pairs,
\code{n\_overlap = n\_valid - abs(lag)}, not by a constant. This is the
\term{lag-dependent normalization} that makes the finite-window estimator
unbiased: fewer pairs contribute at larger lags, and dividing by the actual count
corrects for it.

\paragraph{10. The $k=4$ cumulant subtraction} (\file{cumulant\_estimator.py:346--433}).
Discussed next.

\subsection{$k=2$, $k=3$, $k=4$, and beyond}

For $k \le 3$ the connected cumulant \emph{equals} the centered moment, so the
centered, factorial-corrected product computed above \emph{is} the cumulant ---
there is nothing further to subtract. This is why the comment at
\file{cumulant\_estimator.py:343--345} notes that for $k \le 3$ the cumulant equals
the central moment.

For $k=4$ the connected cumulant requires subtracting the three disconnected pair
products (\file{cumulant\_estimator.py:334--359}):
\begin{equation}
  \kappa_4(X_1,X_2,X_3,X_4)
  \;=\; m_4
       \;-\; m_2(1,2)\,m_2(3,4)
       \;-\; m_2(1,3)\,m_2(2,4)
       \;-\; m_2(1,4)\,m_2(2,3),
  \label{eq:kappa4}
\end{equation}
with $m_n$ the centered moments. The cross-correlation above gave $m_4$ at each
sweep lag; the three pair products are computed by a nested helper
\code{\_two\_point(leg\_a, leg\_b, lag\_value)}
(\file{cumulant\_estimator.py:361--420}). That helper is itself
factorial-corrected when the pair is same-population, same spike field type, and
same effective lag (a genuine shot-noise coincidence to remove), and uses ordinary
linear covariance in every other case (\file{cumulant\_estimator.py:381--420}).
The three partitions of four legs into two pairs are
$\{(0,1)(2,3),\,(0,2)(1,3),\,(0,3)(1,2)\}$
(\file{cumulant\_estimator.py:357--359}), and the disconnected piece is subtracted
lag by lag (\file{cumulant\_estimator.py:422--433}).

\begin{gotcha}
\textbf{$k>4$ is not a cumulant.} For $k>4$ the partition-based subtraction is
\emph{not} implemented. The estimator emits a \code{RuntimeWarning} and returns the
centered $k$-point \emph{moment} instead of the connected cumulant
(\file{cumulant\_estimator.py:435--440}). Separately, the full diagonal/contact
$\delta$-terms are \emph{never} estimated by this code --- they are distributions,
not functions on a $\tau$ grid (\file{cumulant\_estimator.py:33--40}). The
estimator targets the smooth off-diagonal part only.
\end{gotcha}

\subsection{\texttt{collections.Counter}: counting coincident legs}
\label{sec:counter}

\code{collections.Counter} is a standard-library \code{dict} subclass that counts
hashable items: \code{Counter(['a','a','b'])} returns \code{\{'a': 2, 'b': 1\}}.
The estimator uses it to find the \emph{multiplicity} $m$ --- how many legs of the
same \code{(population, field\_type)} land in the same bin, which is exactly the
order of the falling factorial to apply (\file{cumulant\_estimator.py:279--282}):

\begin{lstlisting}
pop_ft_multiplicities = Counter(
    (pop_indices[i], field_types[i]) for i in leg_indices)
for (pop, ft), m in pop_ft_multiplicities.items():
    ...
\end{lstlisting}

The same idiom appears in the direct estimator (\file{cumulant\_estimator.py:547}).
The value $m$ then selects between linear centering ($m=1$) and factorial
correction ($m \ge 2$).

\section{Field types: \texttt{dn}, \texttt{dv}, \texttt{dm}}
\label{sec:fieldtypes}

The estimator's \code{field\_types} argument labels each leg as one of three kinds
(\file{cumulant\_estimator.py:109--121}):

\begin{description}
  \item[\code{'dn'}] a cortical spike-train fluctuation: discrete counts in
        \code{binned\_counts}, shot-noise corrected with the falling factorial,
        rate-converted by $\dd t_{\mathrm{bin}}$.
  \item[\code{'dv'}] a voltage fluctuation: a smooth continuous field in
        \code{voltage\_bins}, linearly centered, no factorial, no rate conversion.
        \emph{This is the label the OU simulators use} --- the OU variable $x$ is a
        smooth field, so it is passed as \code{voltage\_bins} with
        \code{field\_types=['dv', ...]}.
  \item[\code{'dm'}] an external GTaS rate fluctuation: spike-like, in
        \code{ext\_binned\_counts}, treated exactly like \code{'dn'} (discrete
        counts, $\dd t_{\mathrm{bin}}$ denominator, factorial correction).
\end{description}

\begin{gotcha}
\textbf{The field-type normalization fixed a $1/\dd t_{\mathrm{bin}}^{2}$ bug.}
Before the \code{\_normalize\_ft} helper (\file{cumulant\_estimator.py:166--200}),
natural-name legs like \code{'n'} or \code{'nE'} were treated as non-spike, so they
used \code{denom = 1.0} instead of \code{dt\_bin} and the estimator returned
$\text{count}^{2}$ instead of $\text{rate}^{2}$ --- wrong by a factor
$1/\dd t_{\mathrm{bin}}^{2}$ per leg (a factor of 16 at $\dd t_{\mathrm{bin}} =
0.25$). Every notebook that built \code{field\_types = [ef[0] for ef in
external\_fields]} hit it. New code should still prefer the canonical
\code{'dn'}/\code{'dv'}/\code{'dm'} labels; the normalizer is a safety net, not a
license to be sloppy.
\end{gotcha}

\section{The convenience wrapper: \texttt{estimate\_kpoint\_slices}}

Notebooks rarely call \code{compute\_kpoint\_slice} directly. They call
\code{estimate\_kpoint\_slices} (\file{cumulant\_estimator.py:445}), which produces
all $k-1$ axis-parallel slices of the connected $k$-point cumulant in one go,
laid out to line up panel-for-panel with the theory side's
\code{res['C\_tau\_slices']}:

\begin{lstlisting}
def estimate_kpoint_slices(dt_bin, pop_indices, field_types, base_bins,
                           max_lag_bins, binned_counts=None, voltage_bins=None,
                           ext_binned_counts=None):
\end{lstlisting}

Leg 0 is the \emph{anchor}, pinned at lag 0. \code{base\_bins} is a length-$(k-1)$
list of integer bin offsets --- the ``base point,'' one entry per non-anchor leg.
For each $j = 1\ldots k-1$ the wrapper builds a lag list with leg 0 at 0, every
other non-anchor leg at its \code{base\_bins} offset, and leg $j$ set to
\code{None} (the sweep), then calls \code{compute\_kpoint\_slice} and stacks the
result (\file{cumulant\_estimator.py:472--483}):

\begin{lstlisting}
for j in range(1, k):
    lag = [0] * k                       # leg 0 anchored at 0
    for i in range(1, k):
        lag[i] = int(base_bins[i - 1])  # non-anchor legs at the base...
    lag[j] = None                       # ...except leg j, which sweeps
    tau, C = compute_kpoint_slice(...)
    rows.append(np.asarray(C))
\end{lstlisting}

It returns \code{(tau\_grid, C)} with \code{tau\_grid} of shape
\code{(2*max\_lag\_bins+1,)} and \code{C} of shape
\code{(k-1, 2*max\_lag\_bins+1)} --- row $j-1$ is slice $j$. As a convenience, if
\code{binned\_counts} is \code{None} but \code{voltage\_bins} is given, it
allocates a zero \code{binned\_counts} (the pure-voltage case used by the OU
simulators, \file{cumulant\_estimator.py:469--470}). For $k=2$ this returns a
single autocorrelation slice; for $k\ge 3$ it returns one row per swept leg, which
is what makes the simulation curves line up with \code{dd.plot\_cumulant}'s $k\ge3$
slice panels.

\section{The non-FFT cross-check: \texttt{compute\_kpoint\_slice\_direct}}

\code{compute\_kpoint\_slice\_direct} (\file{cumulant\_estimator.py:486}) has the
same interface and output as \code{compute\_kpoint\_slice} but uses an explicit
double loop over \code{(sweep\_lag, t)} instead of FFTs --- $O(n_{\mathrm{bins}}
\times n_{\mathrm{lags}})$ versus $O(n_{\mathrm{bins}}\log n_{\mathrm{bins}})$. It is
slower but free of any FFT artifact (circular wrap, zero-pad, spectral leakage),
which makes it the natural way to cross-check the FFT estimator.

It is, deliberately, a \emph{reduced} re-implementation. It handles spike counts
only --- there are no \code{field\_types}, \code{voltage\_bins}, or
\code{ext\_binned\_counts} parameters --- and its $m \ge 2$ factorial mean is
approximated by the \emph{full-array} mean rather than the valid-window mean used
by the FFT path (\file{cumulant\_estimator.py:559--563}):

\begin{lstlisting}
ff_mean = falling_factorial_array(
    binned_counts[pop], m
).mean() / dt_bin**m
val *= (ff - ff_mean)
\end{lstlisting}

\begin{gotcha}
\textbf{Three estimators, one logic, no shared kernel.}
\code{compute\_kpoint\_slice} (FFT), \code{compute\_kpoint\_slice\_direct} (non-FFT),
and \code{cumulant\_direct\_numba.\_kpoint\_slice\_direct\_k3} (a numba-compiled
$k=3$ variant) implement the \emph{same} factorial-cumulant logic three times for
three different reasons --- FFT speed, FFT-free validation, and numba-compiled $k=3$
throughput. They must be kept in sync by discipline; there is no shared kernel.
Small discrepancies between the FFT path and the direct path are \emph{expected} at
the window edges, because the direct path uses the full-array factorial mean rather
than the valid-window mean.
\end{gotcha}

\section{Spatial estimators}

The spatial simulators (\file{spatial\_field\_1d\_sim.py} and its 2D/3D siblings)
carry their \emph{own} estimators, because a spatial field is averaged over the
periodic translation group rather than over a temporal lag sweep. The 1D file
exposes (\file{spatial\_field\_1d\_sim.py}):

\begin{itemize}
  \item \code{structure\_factor(snaps, meta)} --- the momentum-space two-point
        function $S(q) = \avg{|\hat\phi_q|^{2}}$ (\code{:175}), which tends to
        $T/(\mu + Dq^{2})$ at tree level.
  \item \code{equal\_time\_correlator(snaps)} --- the real-space equal-time
        correlator $C(x)$ (\code{:191}).
  \item \code{lattice\_sum\_variance(L, N, mu, D, T)} --- the exact $\lambda = 0$
        variance reference (\code{:205}), summed over the \emph{lattice} modes.
  \item \code{third\_cumulant\_x(snaps)} --- $\kappa_3(m) =
        \avg{\delta\phi(x_0)\,\delta\phi(x_0+m\,\dd x)^{2}}$ with a standard error
        (\code{:215}).
\end{itemize}

These are the spatial counterparts of \code{compute\_kpoint\_slice}: instead of a
temporal lag sweep they use the periodic translational average --- one circular
cross-correlation per snapshot via FFT.

\begin{gotcha}
\textbf{Matched-cutoff principle.} The spatial simulator realises the
\emph{lattice} dispersion $\omega = \mu + (2D/\dd x^{2})(1-\cos)$, not the continuum
$\mu + Dq^{2}$. Quantitative validation must therefore compare against a theory
computed at the \emph{same} UV cutoff $k_{\max} = \pi/\dd x$, or against the lattice
oracle (\code{lattice\_sum\_variance}, or
\code{coupled\_box\_correlator(dispersion='lattice')}). The two dispersions differ
at $O(\dd x)$ --- about 4\% on $C(0,0)$ at $L=20$, $n_x=128$.
\end{gotcha}

\section{Coupled multi-field validation}
\label{sec:coupled}

The coupled reaction--diffusion simulator
(\code{simulate\_coupled\_rd\_1d}, \file{coupled\_rd\_1d\_sim.py}) is the ground
truth for the coupled Dyson series. Uniquely in this subsystem, it ships with an
\emph{analytic} oracle as well as a simulator, and the oracle is the one place
\code{scipy} appears.

\term{SciPy} is the scientific-computing layer built on NumPy;
\code{scipy.linalg} is its dense linear-algebra module (a set of LAPACK wrappers).
The oracle \code{coupled\_box\_correlator} uses two of its routines, per Fourier
mode (\file{coupled\_rd\_1d\_sim.py:73}):

\begin{itemize}
  \item \code{scipy.linalg.expm(A)} --- the \term{matrix exponential} $\ee^{A}$,
        used to build the exact lag propagator $G(\tau) = \ee^{-A(q)\,|\tau|}$ at
        each mode (\file{coupled\_rd\_1d\_sim.py:353}).
  \item \code{scipy.linalg.solve\_continuous\_lyapunov(A, Q)} --- solves the
        \term{continuous Lyapunov equation} $A\Sigma + \Sigma A^{\mathsf T} = Q$ for
        $\Sigma$, the exact stationary covariance of the linear coupled system at
        each mode (\file{coupled\_rd\_1d\_sim.py:351}).
\end{itemize}

Together these give an \emph{analytic} (no-simulation) reference against which the
coupled Dyson series is validated. The simulator side
(\code{simulate\_coupled\_rd\_1d}) returns a dict
\code{\{'C', 'C\_err', 'C\_rep', 'x\_grid', 'taus', 'meta'\}}.

\begin{gotcha}
\textbf{Coupled-RD Courant guard.} \code{simulate\_coupled\_rd\_1d} \emph{asserts}
$\dd t\cdot\max(D)/\dd x^{2} < 0.4$ (\file{coupled\_rd\_1d\_sim.py:212--214}) --- an
\code{AssertionError}, not a graceful message --- to keep both schemes stable and
the Pad\'e lag bias negligible. If you crank the step or shrink the grid spacing
too far, this is what fires.
\end{gotcha}

\section{The rest of \texttt{models/}}

For completeness, the wider supporting cast. All are independent brute-force
validators paired with the estimator described above.

\begin{center}
\small
\begin{longtable}{@{}p{0.30\textwidth} p{0.15\textwidth} p{0.47\textwidth}@{}}
\toprule
\textbf{File} & \textbf{Family} & \textbf{What it simulates} \\
\midrule
\endhead
\file{cumulant\_direct\_numba.py} & estimator &
  \code{@numba.njit} direct (non-FFT) $k=3$ factorial-cumulant slice
  (\code{\_kpoint\_slice\_direct\_k3}); a plain \code{.py} file so numba can
  compile it. \\
\file{hawkes\_sim\_numba.py} & Hawkes &
  Single-population linear Hawkes; Poisson spikes feed the voltage; returns the
  tuple \code{(binned\_counts, voltage\_bins, total\_spikes)}. \\
\file{hawkes\_sim\_multipop\_numba.py} & Hawkes &
  Multipopulation linear Hawkes with per-pair exponential synaptic filters
  $F_{ij}$; flat E/I neuron index; matches
  \file{multipopulation\_test.theory.py}. \\
\file{hawkes\_sim\_expg\_numba.py}, \file{hawkes\_sim\_quad\_expg\_numba.py},
  \file{hawkes\_sim\_*\_gtas\_numba.py} & Hawkes &
  Exponential-$g$ synaptic-filter and quadratic-rate variants; \code{gtas} =
  ``GTaS'' external spike drive (the \code{'dm'} field, fourth return value
  \code{ext\_binned\_counts}). \\
\file{hawkes\_sage.py} and the \file{hawkes\_*\_sage} / \file{hawkes\_*\_expg}
  family & Hawkes (Sage) &
  Sage-side theory/orchestration helpers for the Hawkes models (these \emph{are}
  preparsed; not \code{njit}). \\
\file{dendritic\_linear\_sim\_numba.py} & Hawkes-like &
  Dendritic linear model simulator; imports \code{compute\_kpoint\_slice}. \\
\file{spatial\_field\_1d\_sim.py} & spatial Langevin &
  1D scalar $\phi^{4}$ on a ring, spectral ETD1; \code{simulate},
  \code{structure\_factor}, \code{equal\_time\_correlator},
  \code{lattice\_sum\_variance}, \code{third\_cumulant\_x}. \\
\file{spatial\_field\_phi6\_1d\_sim.py} & spatial Langevin &
  1D $\phi^{6}$ (cubic+quintic force); near-copy of the $\phi^{4}$ simulator. \\
\file{spatial\_field\_2d\_sim.py}, \file{spatial\_field\_3d\_sim.py} &
  spatial Langevin &
  2D scalar on a torus / 3D scalar, ETD1; radial structure-factor and correlator
  estimators. \\
\file{coupled\_rd\_1d\_sim.py} & coupled RD &
  $N$-species reaction--diffusion with cross-noise and unequal $D$; semi-implicit
  Crank--Nicolson or explicit scheme; \emph{plus} the analytic oracle
  \code{coupled\_box\_correlator} (\code{scipy.linalg.expm} /
  \code{solve\_continuous\_lyapunov}). \\
\bottomrule
\end{longtable}
\end{center}

\section{Putting it together: a sim-compare notebook}

We close with the canonical wiring, taken from cell~9 of
\file{notebooks/examples/temporal\_ou\_quartic\_white.ipynb} --- the notebook
behind the Chapter~\ref{ch:quickstart} payoff. Read it against the seven salient
points that follow:

\begin{lstlisting}
from models.ou_langevin_sim_numba import sim_ou_quartic_numba
from models.cumulant_estimator import estimate_kpoint_slices

fp = res['_resolved']['parameters']          # same physics as the theory run
mu = float(fp['mu']); eps = float(fp['eps']); D = float(fp['D'])   # CAST to float!

dt_sim, dt_bin = 0.01, 0.02
n_steps        = int(T_sim / dt_sim)
bin_size_steps = int(max(round(dt_bin / dt_sim), 1))
dt_bin_eff     = bin_size_steps * dt_sim     # the EFFECTIVE bin width (use this)
n_bins         = int(n_steps // bin_size_steps)
max_lag_bins   = int(tau_max / dt_bin_eff)

k           = int(res['_resolved']['k'])
base_bins   = [int(round(b / dt_bin_eff)) for b in base]   # length k-1
pop_indices = [0]*k
field_types = ['dv']*k                       # OU field x is a 'dv' (smooth)

_ = sim_ou_quartic_numba(1000, dt_sim, mu, eps, D, 0.0,
                         bin_size_steps, 100, 0)            # JIT warmup
for r in range(N_RUNS):
    x_bins = sim_ou_quartic_numba(n_steps, dt_sim, mu, eps, D, 0.0,
                                  bin_size_steps, n_bins, rng_base + r)
    tau_sim, Cj = estimate_kpoint_slices(
        dt_bin_eff, pop_indices, field_types, base_bins,
        max_lag_bins, voltage_bins=x_bins)                 # Cj: (k-1, n_tau)
    C_runs.append(np.asarray(Cj).real)
C_sim = np.array(C_runs).mean(axis=0)
C_err = np.array(C_runs).std(axis=0, ddof=1) / np.sqrt(N_RUNS)
\end{lstlisting}

\begin{enumerate}
  \item \textbf{The simulator is handed \code{res['\_resolved']['parameters']}} ---
        the pipeline-resolved physics --- so theory and simulation share parameters
        but \emph{not} code.
  \item \textbf{Every scalar is cast} (\code{float(fp['mu'])}, \code{int(...)})
        before the \code{njit} call: the Sage-Integer guard of \S\ref{sec:numbasage}
        in action.
  \item \textbf{A throwaway warmup call} triggers JIT compilation once so it is not
        timed in the measurement loop.
  \item \textbf{The OU field is passed as \code{voltage\_bins} with
        \code{field\_types=['dv']}} --- it is a smooth continuous field, so it uses
        linear centering, not the factorial spike correction. (Hawkes spike legs
        would use \code{'dn'} and pass \code{binned\_counts} instead.)
  \item \textbf{\code{dt\_bin\_eff}} (= \code{bin\_size\_steps}\,$\times$\,\code{dt\_sim}),
        not the nominal \code{dt\_bin}, is what the estimator receives --- the
        binning rounds \code{dt\_bin/dt\_sim} to an integer number of steps, so the
        true bin width may differ slightly from the requested one.
  \item \textbf{Multiple independent runs} (\code{N\_RUNS}, distinct seeds) give the
        error bar \code{C\_err = std/sqrt(N\_RUNS)}.
  \item \textbf{The comparison} is \code{C\_sim} versus \code{res['C\_tau\_by\_ell'][0]}
        (tree) and \code{res['C\_tau']} (loop-corrected) at the $\tau=0$ index ---
        \emph{that overlap is the trust-chain verdict.}
\end{enumerate}

For $k\ge 3$, \code{estimate\_kpoint\_slices} returns one row per swept leg, lining
up panel-for-panel with \code{dd.plot\_cumulant}'s theory slices, so the same loop
validates a third or fourth cumulant with no change beyond \code{k} and
\code{base\_bins}.

\section{Caveats checklist}

A consolidated list of the traps documented above, for quick reference when a
comparison plot misbehaves:

\begin{enumerate}
  \item \textbf{Sage \code{Integer}/\code{RealNumber} crash numba.} Keep the
        simulators in \code{.py} files; cast every call-site scalar with
        \code{int()}/\code{float()} (\S\ref{sec:numbasage}).
  \item \textbf{Colored-noise factor of 2.} The white limit of
        \code{sim\_ou\_quartic\_colored\_numba} is $\avg{\xi\xi}\to 4D\,\delta$;
        use $D_{\mathrm{white}} = 2\,D_{\mathrm{colored}}$ to compare
        (\S\ref{sec:colored}).
  \item \textbf{\code{sim\_ou\_quartic\_two\_dim\_corr\_numba} vs its theory file
        disagree on the noise coefficient.} The simulator uses
        $\avg{\xi\xi} = 2D\,\delta$; the theory file as written declares
        \code{coefficient='D1'} ($= D\,\delta$, half strength). A quantitative
        compare needs one side fixed --- edit the theory to \code{'2*D1'} or halve
        the simulator's $D$ (\file{ou\_langevin\_sim\_numba.py:412--419}). Speed
        tests are unaffected.
  \item \textbf{$\rho$ out of $[-1,1]$ yields NaN, not an error}
        (\file{ou\_langevin\_sim\_numba.py:332--335}).
  \item \textbf{FFT lag indexing is fragile.} Read as
        \code{raw\_xcorr[(-lag) \% n\_fft]}; do not use \code{[::-1]}; zero-pad to
        \code{n\_fft >= 2*n\_valid} (\file{cumulant\_estimator.py:316--319}).
  \item \textbf{Field-type normalization fixed a $1/\dd t_{\mathrm{bin}}^{2}$ bug.}
        Prefer the canonical \code{'dn'}/\code{'dv'}/\code{'dm'} labels
        (\S\ref{sec:fieldtypes}).
  \item \textbf{$k>4$ is not a cumulant} --- the estimator returns the centered
        moment with a \code{RuntimeWarning}
        (\file{cumulant\_estimator.py:435--440}).
  \item \textbf{Euler--Maruyama bias.} All \code{*\_numba} Langevin simulators are
        weak order 1; stationary stats carry $O(\dd t)$ bias, so
        $\dd t_{\mathrm{sim}}\ll 1/\mu$ is required. The colored and spatial
        simulators remove this for the \emph{linear} part via the exact OU / ETD1
        update; the coupled-RD semi-implicit scheme is $O(\dd t^{2})$-exact in the
        stationary covariance.
  \item \textbf{Coupled-RD Courant guard.} \code{simulate\_coupled\_rd\_1d}
        asserts $\dd t\cdot\max(D)/\dd x^{2} < 0.4$
        (\file{coupled\_rd\_1d\_sim.py:212--214}).
  \item \textbf{Bin finely or bias $C_{xx}(0)$ low.} A coarse bin smooths the field
        and underestimates the equal-time variance; the example uses
        \code{dt\_bin}\,$=0.02$ for $\mu\sim1$.
  \item \textbf{Matched-cutoff principle (spatial).} Compare against a theory at the
        same $k_{\max}=\pi/\dd x$, or against the lattice oracle; the dispersions
        differ at $O(\dd x)$.
\end{enumerate}

% ---------------------------------------------------------------------
%  Chapter 21 : The Theory-Builder UI
% ---------------------------------------------------------------------
\chapter{The Theory-Builder UI}\label{ch:ui}

Every other chapter of this manual assumes a \file{theories/<name>.theory.py}
file already exists --- a small Python module with a \code{build()} function that
chains a fluent builder (\code{TemporalTheoryBuilder(...)} or
\code{SpatialTheoryBuilder(...)}, then \code{.physical\_field(...)},
\code{.parameter(...)}, \code{.set\_action\_text(...)}, \code{.equation(...)},
\code{.build()}) into a model dict. The chapters on the theory specification
(Chapter~\ref{ch:theory-spec}) and on theory files (Chapter~\ref{ch:theory-files})
teach you to write those modules by hand. This chapter is about the alternative:
a point-and-click form, rendered \emph{inside the notebook}, that writes the file
for you.

The form is the class \code{TheoryUI}, and it lives in two files:
\file{pipeline/ui/main.py} (the form itself --- the ten-tab editor, the
validation sidebar, the buttons) and \file{pipeline/ui/widgets.py} (the small
reusable widget primitives the tabs are built from). The entire subsystem is
written on top of one external library, \term{ipywidgets}, the Jupyter project's
interactive-widget toolkit. If you have never used ipywidgets, this chapter
teaches you exactly as much of it as you need, from the ground up.

The single most important thing to understand before we start is what the UI
\emph{does not} do. \term{It does no physics.} It never enumerates a diagram,
never builds a propagator, never solves a saddle (with one small, optional
exception we will meet at the end). It is pure \emph{input specification}: it
collects the state of a form, assembles it into a plain Python dictionary --- the
\term{spec dict} --- and hands that dictionary to the serializer
\file{pipeline/theory\_serialize.py}, which renders it to a \code{.theory.py}
source file on disk. \textbf{Load} reverses the trip: it parses a saved file back
into a spec dict and repopulates the widgets. That is the whole job.

\begin{note}
Why a GUI at all, when the builder API is only a dozen lines? Because writing
those lines \emph{correctly} requires knowing the framework's conventions that
are easy to get wrong: the auto-generated response field \code{<name>t}, the
saddle parameter \code{<name>star}, the action grammar (\code{Dt} as a
time-derivative operator, \code{for i in pop} comprehensions, \code{Laplacian} in
spatial theories), and a pile of optional knobs (boundary conditions, Dyson
dressing order, operator-IR mode, stability analysis, run metadata). The form
surfaces every one of those as a labelled widget with help text, and a live
sidebar flags undeclared names and reserved-symbol collisions \emph{as you type}
--- so a typo is caught at authoring time instead of blowing up later with an
opaque \code{Sage} error deep in the engine.
\end{note}

\section{Where the UI sits in the pipeline}

The UI is the front door. A user opens the notebook
\file{notebooks/theory\_builder.ipynb}, runs one cell, fills in tabs, and clicks
\textbf{Save theory file}. From there the deliverable --- the \code{.theory.py}
file --- flows downstream exactly as if it had been hand-written. The data path
is short and one-directional on save, and reversible on load:

\begin{center}
\begin{tikzpicture}[
    node distance=6mm,
    box/.style={draw, rounded corners=2pt, align=center, font=\small,
                inner sep=5pt, fill=codebg},
    lbl/.style={font=\scriptsize\itshape, color=cmt},
    arr/.style={-{Stealth[length=2.4mm]}, thick}]
  \node[box] (nb) {\code{theory\_builder.ipynb}\\\code{TheoryUI().show()}};
  \node[box, below=of nb] (collect) {\code{TheoryUI.\_collect()}\\
        snapshot form $\rightarrow$ \textbf{spec dict}};
  \node[box, below=of collect] (ser) {\code{theory\_serialize}\\
        \code{save\_theory\_to\_file(spec, dir)}};
  \node[box, below=of ser] (file) {\file{theories/<slug>.theory.py}\\
        the deliverable on disk};
  \node[box, below=of file] (runner) {\code{theory\_runner.ipynb}\\
        \code{build()} $\rightarrow$ \code{compute\_cumulants(...)}};
  \draw[arr] (nb) -- (collect) node[lbl, midway, right] {user types in 10 tabs};
  \draw[arr] (collect) -- (ser) node[lbl, midway, right] {Save};
  \draw[arr] (ser) -- (file);
  \draw[arr] (file) -- (runner) node[lbl, midway, right] {the rest of the pipeline};
  % Load reverse arrow
  \draw[arr, dashed, color=Maroon] (file.west) to[out=180, in=180]
        node[lbl, midway, left, align=right] {Load:\\\code{load\_spec\_from\_file}\\
        $\rightarrow$ repopulate} (collect.west);
\end{tikzpicture}
\end{center}

\begin{itemize}
  \item \textbf{What feeds the UI:} the \emph{user} (keyboard input), and --- on
        \textbf{Load} --- an existing \file{theories/*.theory.py} file, parsed
        back into a spec dict by \code{load\_spec\_from\_file}.
  \item \textbf{What consumes the UI's output:} \code{pipeline.theory\_serialize}
        (the serializer that turns the spec dict into a \code{.theory.py} source
        string). The file in turn is consumed by the runner notebook (and the
        \code{nb\_support} engine), which calls \code{build()} to get the model
        dict and feeds it to \code{compute\_cumulants}. The UI also writes a
        one-line scratch file \file{.theories/.last\_built} so the
        \textbf{Open in runner notebook} button can hand the runner the
        just-saved name.
\end{itemize}

The boundary is clean: the UI's single output is the spec dict, and (via the
serializer) the \code{.theory.py} file. Everything downstream of that file is the
subject of the other chapters.

\section{Launching it}

The entire UI is two lines in a notebook cell. The package
\file{pipeline/ui/\_\_init\_\_.py} re-exports the class, so:

\begin{lstlisting}
from pipeline.ui import TheoryUI
TheoryUI().show()
\end{lstlisting}

That is the contents of the one cell in \file{notebooks/theory\_builder.ipynb}.
Constructing \code{TheoryUI()} builds every widget; \code{.show()} renders the
form into the cell's output area. The form then stays \emph{live}: typing in a box
or clicking a button runs Python on the kernel without re-executing the cell.

Two things happen the first time you call \code{show()}, and they are worth
knowing because they explain a small surprise later. Reading
\file{pipeline/ui/main.py:2311}:

\begin{lstlisting}
def show(self) -> None:
    global _CSS_INJECTED
    if not _CSS_INJECTED:
        display(HTML(_THEORY_BUILDER_CSS))
        _CSS_INJECTED = True
    try:
        self._refresh_validation()
    except Exception:
        pass
    display(self._root)
\end{lstlisting}

First, it injects an inline stylesheet --- one big \code{<style>} block in the
\code{tb-*} CSS namespace (\file{pipeline/ui/main.py:63}) --- but only \emph{once
per kernel session}, guarded by the module-level flag \code{\_CSS\_INJECTED}
(\file{pipeline/ui/main.py:244}). Second, it runs the validator a single time so
the sidebar shows real state on first render rather than placeholder text. Then it
\code{display}s the root container.

\begin{gotcha}
Because the CSS injection is guarded by a module-level flag, multiple
\code{TheoryUI()} instances in one kernel share the one stylesheet. If you edit
\code{\_THEORY\_BUILDER\_CSS} in the source, you must \emph{restart the kernel} to
see the change --- re-running the cell will not re-inject it. This is a deliberate
trade (the alternative is duplicate \code{<style>} blocks piling up on every
re-render), but it surprises people who are iterating on the styling.
\end{gotcha}

After you save, two attributes and one method form the public surface:
\code{ui.last\_saved} (the path of the last file written),
\code{ui.theories\_dir} (where saves land by default), and \code{ui.spec()}
(\file{pipeline/ui/main.py:2330}), which returns the current form state as a spec
dict without writing anything --- handy if you want to inspect or post-process the
dict in a later cell.

\section{ipywidgets in five minutes}

The rest of this chapter quotes ipywidgets code constantly, so here is the whole
mental model a physicist needs. Skip this section if you already know the library.

\term{ipywidgets} is a library of interactive HTML controls --- buttons, text
boxes, dropdowns, checkboxes, tabs, sliders --- that render inside a notebook
output cell and stay \emph{live}. The trick is that each widget is two halves kept
in sync: a Python object on the \emph{kernel} side and a DOM element on the
\emph{browser} side, connected over Jupyter's ``comm'' (communication) channel.
When you type in a box, the browser sends the new text to the kernel, the Python
object's \code{.value} attribute updates, and any callbacks you registered fire.
This is what lets a notebook cell host a full GUI with no web server of its own.

The conventional import alias is \code{W}, used in both source files
(\file{pipeline/ui/main.py:27}, \file{pipeline/ui/widgets.py:20}):

\begin{lstlisting}
import ipywidgets as W
from IPython.display import display, HTML
\end{lstlisting}

The handful of widget classes this code uses fall into four groups.

\paragraph{Layout containers.} \code{W.VBox([...])} and \code{W.HBox([...])} stack
their children vertically and horizontally. \code{W.Tab(children=[...])} is the
tab strip; \code{.set\_title(i, str)} names tab \code{i} and
\code{.selected\_index} reads or sets the active tab. \code{W.Layout(...)} is the
CSS-like layout object you attach to any widget --- \code{width='400px'},
\code{height='60px'}, or \code{display='none'} to hide it entirely.

\paragraph{Inputs.} \code{W.Text} (single line), \code{W.Textarea} (multi-line ---
used for the action and the seed box), \code{W.Dropdown} (single-select),
\code{W.Checkbox} (booleans), \code{W.ToggleButtons} (the Temporal/Spatial
switch), and the numeric trio \code{W.IntText} / \code{W.FloatText} /
\code{W.BoundedIntText} (where \emph{bounded} clamps to a \code{[min, max]}
range).

\paragraph{Output and display.} \code{W.HTML(value=...)} renders a static HTML
string --- every tab's help text, the readiness sidebar, and the cheat sheet are
\code{W.HTML} widgets. \code{W.Output(...)} is a \emph{capture region}: a panel
into which ordinary \code{print(...)} can be redirected. The UI's status panel is
one \code{W.Output}; the button handlers write to it with a
\code{with self.\_status:} context (\file{pipeline/ui/main.py:2642}), which
temporarily routes \code{stdout} into the panel, and \code{.clear\_output()} wipes
it. From \code{IPython.display}, \code{display(obj)} renders a rich object in the
current cell, \code{HTML(str)} wraps a raw HTML string, and \code{Javascript(str)}
ships a JS snippet to the browser (the \textbf{Open in runner} button uses it).

\paragraph{The event model.} Two registration calls appear everywhere:

\begin{lstlisting}
# fire a callback whenever a widget's .value changes:
self._w_theory_mode.observe(lambda _ch: self._on_mode_change(), names='value')
# fire a callback when a button is clicked:
self._btn_save.on_click(self._on_save)
\end{lstlisting}

\code{widget.observe(cb, names='value')} (\file{pipeline/ui/main.py:409}) calls
\code{cb} with a change-event object every time \code{.value} changes; the lambda
\code{lambda \_ch: ...} simply swallows that argument. \code{button.on\_click(cb)}
(\file{pipeline/ui/main.py:1376}) fires on click. Finally,
\code{widget.add\_class('tb-...')} attaches a CSS class to the widget's DOM node,
which is how the inline \code{tb-*} stylesheet finds its targets.

\begin{defn}
\term{observer} --- a callback registered with \code{widget.observe(cb,
names='value')} that fires whenever the widget's value changes. The UI registers
one observer, \code{\_mark\_changed}, on \emph{every} input, so any edit flips the
``unsaved changes'' flag and refreshes the validation sidebar. This is the
mechanism behind the live, type-as-you-go feedback.
\end{defn}

\section{The math behind each tab}

The UI does not implement the \msrjd{} field theory, but it is where the user
\emph{declares} it, so every tab maps onto a piece of the theory. A one-page
recap fixes what each tab is collecting and why the validator flags what it flags;
the full treatment is in Chapters~\ref{ch:theory-spec}--\ref{ch:enumeration}.

A stochastic equation of motion for a field $\phi$ (a scalar $x(t)$ for a temporal
theory, or a continuous field $\phi(x,t)$ for a spatial one) driven by noise
$\eta$,
\begin{equation}
  \Dt \phi \;=\; F[\phi] \;+\; \eta,
  \qquad \avg{\eta(t)\,\eta(t')} \;=\; 2D\,\delta(t-t'),
  \label{eq:ui-eom}
\end{equation}
is recast as a path integral by introducing a \term{response field} $\tilde\phi$
(Martin--Siggia--Rose / Janssen / De~Dominicis). The generating functional becomes
$Z = \int \mathcal{D}\phi\,\mathcal{D}\tilde\phi\,\ee^{-S[\phi,\tilde\phi]}$ with
the \msrjd{} \term{action}
\begin{equation}
  S[\phi,\tilde\phi] \;=\; \int \dd t\,\Big[\,\tilde\phi\,(\Dt\phi - F[\phi])
        \;-\; D\,\tilde\phi^{2}\,\Big].
  \label{eq:ui-action}
\end{equation}
This action is exactly the string the user types on the \textbf{Action} tab. The
placeholder text in the form (\file{pipeline/ui/main.py:1013}) reads:

\begin{lstlisting}[style=console]
phit * ((Dt + mu) * phi + eps * phi^3) - D * phit^2
\end{lstlisting}

which is $\tilde\phi\,(\Dt\phi + \mu\phi + \varepsilon\phi^{3}) - D\tilde\phi^{2}$
--- the action of $\dot\phi = -\mu\phi - \varepsilon\phi^{3} + \eta$ with
$\avg{\eta\eta} = 2D\delta$. Three syntactic pieces are visible:

\begin{itemize}
  \item \code{phit} --- the \term{response field} $\tilde\phi$. The UI never asks
        you to declare it: for every physical field \code{phi} you declare on the
        Fields tab, the framework auto-creates \code{phit} (response),
        \code{dphi} (fluctuation), and \code{phistar} (saddle). The action
        validator knows about all three (\file{pipeline/ui/main.py:1684}).
  \item \code{Dt} --- the \term{time-derivative operator} $\partial_t$, composed
        like an algebra element: \code{(Dt + mu)*phi} means
        $\partial_t\phi + \mu\phi$. It is one of the builtin identifiers the
        validator always accepts (\file{pipeline/ui/main.py:1653}).
  \item \code{- D phit\^{}2} --- the \term{Gaussian white-noise} term. Because the
        weight is $\ee^{-S}$, a term quadratic in $\tilde\phi$ corresponds to the
        \emph{second cumulant} of the noise: $-D\tilde\phi^{2}$ encodes
        $\kappa^{(2)} = 2D$ with a white ($\delta(\tau)$) kernel. A
        $-S_3\,\tilde\phi^{3}$ term would be a third white cumulant
        $\kappa^{(3)} = 3!\,S_3$.
\end{itemize}

With that, the role of each tab is:

\begin{description}
  \item[Model] the theory's name (which derives the filename), an optional
        description, and the Temporal/Spatial toggle.
  \item[Populations] a named group of \code{size N} identical units. Declaring a
        population \code{A} of size $N$ makes every field attached to it a
        length-$N$ vector and introduces a \code{for i in A} sum in the action.
        This is the ``$N$ coupled units'' case (oscillators, spins); a matrix
        parameter \code{w[i,j]} is the coupling between unit $i$ and unit $j$.
  \item[Fields] the physical fields $\phi$. The fluctuation split
        $\phi = \phi^{*} + \delta\phi$ expands the action around its saddle; you
        type only $\phi$ and the framework adds $\tilde\phi$, $\phi^{*}$,
        $\delta\phi$. In Spatial mode this tab also carries the
        \emph{spatial-structure} panel.
  \item[Parameters] the coefficients $\mu,\varepsilon,D,\dots$ appearing in
        $F[\phi]$ and the noise. Their \emph{shape} (scalar / vector indexed by
        one population / matrix indexed by two) follows the population structure.
  \item[Functions] non-polynomial transfer functions $f(\phi) = \tanh(a\phi)$ and
        the like. The framework Taylor-expands $f$ about the saddle of its first
        argument, $f(\phi) = f(\phi^{*}) + f'(\phi^{*})\delta\phi + \dots$, so you
        write \code{f[i](phi[i])} and the expansion is automatic.
  \item[Kernels] memory / convolution couplings $K * y = \int K(t-t')\,y(t')\,
        \dd t'$. The Fourier image $\tilde K(\omega)$ is preferred because the
        propagator builder works in frequency space.
  \item[Noise (CGF)] correlated, coloured, or cross-correlated noise declared as
        cumulant pieces $\avg{\eta(t)\eta(t')} = \text{coefficient}\times
        \text{kernel}(t-t')$. White $\to$ \code{dirac\_delta(tau)}; OU-coloured
        $\to$ \code{exp(-abs(tau)/tauc)}; cross-field $\to$ a coefficient like
        \code{rho*sqrt(D1*D2)} on two response legs.
  \item[Mean-field] the differential-algebraic-equation (DAE) residuals
        $\text{LHS} - \text{RHS} = 0$ solved at the saddle. \code{Dt} on the LHS
        marks a differential equation (set $\Dt\to 0$ at the saddle, kept symbolic
        for linear stability); no \code{Dt} means a purely algebraic equation. The
        solver finds \emph{all} roots and \code{fixed\_point\_index} picks which
        sorted root the diagrammatic expansion uses.
  \item[Defaults] the run metadata: \code{k} (number of external legs --- $k=2$ is
        the covariance / power spectrum), \code{max\_ell} (loop order --- \code{0}
        tree-level/LNA, \code{1} one-loop, \code{2} two-loop), the
        $(\code{tau\_max}, \code{tau\_step})$ lag grid, and the recommended
        external-field legs.
\end{description}

\begin{note}
\emph{Spatial structure} lives only on the Fields tab and only in Spatial mode.
Setting \code{spatial\_dim=1} makes $\phi$ a continuous field $\phi(x,t)$; the
action may then use the diffusion operator \code{Laplacian} multiplicatively
(\code{D*Laplacian*phi} $= D\Lap\phi$). \term{Derivative vertices} --- KPZ's
$(\partial_x h)^{2}$, Burgers, Model~B's $\Lap(\phi^{2})$ --- put a spatial
derivative \emph{inside} a nonlinearity and require the \term{Operator-IR} mode:
the action is then written with \code{Dt()} / \code{Lap()} / \code{Dx(phi, i)}
\emph{call} syntax instead of bare multiplication. We return to why this is
load-bearing in the Gotchas section.
\end{note}

None of this math is computed by the UI --- it is what the UI \emph{collects}. The
reader who wants the actual integrals should consult the theory-specification,
propagator, mean-field, and enumeration chapters.

\section{The widget primitives (\texttt{widgets.py})}

Before the form itself, look at the toolbox it is built from. The file
\file{pipeline/ui/widgets.py} (438 lines) defines a handful of small reusable
ipywidgets composites. Most of them are thin, but one --- \code{DynamicTable} ---
is the workhorse that every tabular tab is made of.

\subsection{Reading the system clipboard}

Several inputs carry a small clipboard-glyph button (rendered as a clipboard
emoji in the live form). Its reason for existing is a
genuine Jupyter wart, documented right in the source
(\file{pipeline/ui/widgets.py:23}): VS~Code's notebook renderer swallows
Ctrl/Cmd-V before it reaches an ipywidgets text input, so keyboard paste is broken
there. Button \emph{clicks} still work, and when the kernel runs on the same
machine as the editor (the normal case) Python can read the real system clipboard
--- so a clicked button reads it and fills the box directly.

\code{read\_system\_clipboard()} (\file{pipeline/ui/widgets.py:29}) returns the
local clipboard text or \code{None}. It tries \code{pyperclip} (an optional
cross-platform clipboard library) first, then falls back to OS-native command-line
tools picked by \code{sys.platform}: \code{pbpaste} on macOS;
\code{wl-paste}/\code{xclip}/\code{xsel} on Linux, in that preference order; and
PowerShell \code{Get-Clipboard} on Windows. For each candidate it skips the
command if \code{shutil.which} finds no binary, else runs it with
\code{subprocess.run(..., capture\_output=True, text=True, timeout=5)} and returns
\code{stdout} on a zero exit code.

\begin{gotcha}
The clipboard read happens on the machine the \emph{kernel} runs on, not your
browser. For a local Jupyter or VS~Code session that is your machine (correct);
for a \emph{remote} kernel it would read the server's clipboard (wrong, but
harmless --- you simply get the wrong text or a failure glyph). \code{paste\_button}
(\file{pipeline/ui/widgets.py:69}) handles the \code{None} case by briefly setting
its glyph to a cross mark and a tooltip suggesting \code{pip install pyperclip} or
a local kernel.
\end{gotcha}

\subsection{The scalar/vector/matrix and expression inputs}

Four small constructors round out the file. \code{vector\_input}
(\file{pipeline/ui/widgets.py:96}) builds a horizontal row of \code{FloatText}
widgets for a length-$N$ vector; \code{matrix\_input}
(\file{pipeline/ui/widgets.py:119}) an $n\times m$ grid; both expose a
monkey-patched \code{.get\_value()}. \code{expression\_input}
(\file{pipeline/ui/widgets.py:147}) is a single-line text box with a live
a live check-mark / cross-mark indicator driven by a \code{validator(text)}
callback.

\begin{gotcha}
\code{vector\_input}, \code{matrix\_input}, and \code{expression\_input} are
exported and importable, but the \emph{current form does not use them}. The form
takes vector and matrix defaults as a text cell on the Parameters tab, and does
its expression validation in the sidebar rather than per-field. They are kept as
reusable primitives; editing them will \emph{not} change the live UI. Of the bulk
import at \file{pipeline/ui/main.py:30}, only \code{textarea\_input},
\code{DynamicTable}, and \code{paste\_button} are actually called.
\end{gotcha}

The one composite the form \emph{does} use directly is \code{textarea\_input}
(\file{pipeline/ui/widgets.py:192}): a multi-line text area with a label and a
clipboard paste button, used for the action editor (\file{pipeline/ui/main.py:1008}).
It exposes the raw \code{Textarea} as \code{box.\_text\_w}, and the form reaches
into \code{\_text\_w} to attach observers and CSS classes.

\subsection{\texttt{DynamicTable} --- the spreadsheet primitive}

\code{DynamicTable} (\file{pipeline/ui/widgets.py:212}) is an add/remove-row
spreadsheet: each row is a dict of named values, with the columns declared up
front. Most of the form's tabs --- Populations, Fields, Parameters, Functions,
Kernels, Noise, Mean-field, and the external-fields table --- are a single
\code{DynamicTable}.

The constructor \code{\_\_init\_\_(self, columns, initial=None,
add\_label='+ add row')} (\file{pipeline/ui/widgets.py:245}) takes a list of
column-spec dicts. The recognized keys, from the class docstring
(\file{pipeline/ui/widgets.py:218}):

\begin{lstlisting}
{'name':             column key,
 'kind':             'text' | 'bool' | 'select' | 'int' | 'float',
 'options':          [...],               # static choices for 'select'
 'options_provider': callable() -> list,  # dynamic choices, re-queried on demand
 'default':          <value>,
 'width':            '120px',
 'placeholder':      'hint text',         # text columns
 'paste':            True}                # add a clipboard button beside a text cell
\end{lstlisting}

The factory \code{\_make\_widget} (\file{pipeline/ui/widgets.py:272}) dispatches on
\code{kind}: \code{text} $\to$ \code{W.Text}, \code{bool} $\to$ \code{W.Checkbox},
\code{select} $\to$ \code{W.Dropdown}, \code{int} $\to$ \code{W.IntText},
\code{float} $\to$ \code{W.FloatText}; an unknown kind raises \code{ValueError}.
Crucially, a \code{select} column's choices come from \code{options\_provider()} if
present, \emph{re-queried on demand}, else from the static \code{options} list ---
this is the hook that makes dropdowns track other tabs (see below). Every widget
gets an observer so edits propagate:

\begin{lstlisting}
w.observe(lambda _change: self._notify_change(), names='value')
\end{lstlisting}

\code{add\_row} (\file{pipeline/ui/widgets.py:309}) builds one widget per column;
for columns flagged \code{'paste': True} and \code{kind=='text'} it wraps the input
in an \code{HBox([w, paste\_button(w)])} --- but the \emph{value} widget is still
the bare \code{w}, so \code{get\_rows} reads it unchanged. Each row gets a
remove button (a cross-mark glyph).

\begin{gotcha}
The remove button is subtle. Its handler is built by
\code{\_make\_remove(len(self.\_row\_widgets))}, which \emph{captures} the row's
index at creation time --- but that captured index is effectively dead code. The
handler ignores it and instead scans \code{\_row\_boxes} for the box whose last
child \emph{is} this button (\file{pipeline/ui/widgets.py:332}), then pops that
index. This identity scan, not the stale captured index, is what makes removal
correct after earlier rows have already been deleted. If you ever refactor this,
keep the identity scan; the captured index is a trap.
\end{gotcha}

\code{get\_rows} (\file{pipeline/ui/widgets.py:365}) reads each row into a
\code{\{name: w.value\}} dict and is the table's single data-out path. It silently
\emph{drops rows whose \code{name} cell is blank}:

\begin{lstlisting}
def get_rows(self) -> list[dict]:
    out = []
    for w_dict in self._row_widgets:
        row = {name: w.value for name, w in w_dict.items()}
        if 'name' in row and not str(row['name']).strip():  # drop blank-name rows
            continue
        out.append(row)
    return out
\end{lstlisting}

\begin{gotcha}
Blank-row dropping is keyed on a column literally named \code{name}. The
external-fields table on the Defaults tab is keyed \code{field}, not \code{name},
so \code{get\_rows} does \emph{not} auto-drop its blank rows --- the metadata
collector re-filters on \code{field} instead (\file{pipeline/ui/main.py:2607}).
Any future table without a \code{name} column would silently keep its empty
starter rows unless it filters itself.
\end{gotcha}

Two more methods complete the picture. \code{set\_column\_visible(name, visible)}
(\file{pipeline/ui/widgets.py:375}) shows or hides one column --- header cell plus
every row's cell, current and future --- by toggling \code{layout.display} between
\code{''} and \code{'none'}; the form uses it to drop the \code{spatial\_dim}
column in Temporal mode. And \code{refresh\_dropdown\_options(col\_name)}
(\file{pipeline/ui/widgets.py:410}) re-queries a \code{select} column's
\code{options\_provider()} and pushes the new option list onto every existing
row's dropdown, preserving the prior selection when still valid:

\begin{lstlisting}
new_opts = list(provider())
for w_dict in self._row_widgets:
    w = w_dict.get(col_name)
    if not isinstance(w, W.Dropdown):
        continue
    prev = w.value
    try:
        w.options = new_opts
        if prev in new_opts:
            w.value = prev
        elif new_opts:
            w.value = new_opts[0]
    except Exception:
        # Some ipywidgets versions reject .options while .value is invalid:
        w.value = (new_opts[0] if new_opts else None)
        w.options = new_opts
\end{lstlisting}

This is how a population declared on the Populations tab instantly appears in the
index dropdowns on every other tab.

\begin{gotcha}
The \code{try/except} here is not defensive padding. Some ipywidgets versions raise
if you set \code{.options} while the current \code{.value} is momentarily not in
the new list; the fallback assigns \code{value} first, then \code{options}. If you
add a new population-aware dropdown anywhere, route its refresh through this method
--- never a raw \code{.options = ...} --- or you will hit that version-dependent
exception.
\end{gotcha}

Finally, \code{\_notify\_change} (\file{pipeline/ui/widgets.py:403}) calls every
registered \code{on\_change} callback inside a bare \code{except Exception: pass},
so one buggy observer cannot wedge the form. The trade is that a buggy validator
callback fails silently.

\section{The form (\texttt{main.py})}

\code{TheoryUI} (\file{pipeline/ui/main.py:289}) is the whole form. Its public
surface is small --- \code{\_\_init\_\_}, \code{show}, \code{spec}, \code{load},
the attributes \code{last\_saved} and \code{theories\_dir} --- and everything else
is private.

\subsection{Construction}

\code{\_\_init\_\_(self, theories\_dir=None)} (\file{pipeline/ui/main.py:296})
resolves the theories directory (default \code{\_DEFAULT\_THEORIES\_DIR}, which is
\code{<cwd>/../theories} --- since notebooks run from \code{notebooks/}, this
resolves to the repo's \file{theories/}), creates it if absent, and initializes
two pieces of bookkeeping state:

\begin{lstlisting}
self.last_saved: Optional[str] = None
self._dirty: bool = False          # unsaved-changes flag
self._loaded_extras: dict[str, dict[str, str]] = {
    'field_latex':    {},
    'function_latex': {},
    'kernel_latex':   {},
}
self._build_widgets()
\end{lstlisting}

\code{\_dirty} tracks unsaved edits. \code{\_loaded\_extras}
(\file{pipeline/ui/main.py:313}) is a per-kind \code{\{name: latex\}} carry-through
map: the \code{latex} column was removed from the UI on 2026-05-27, but older
theory files may still carry custom \code{latex} strings, so \code{load()} stashes
them here and \code{\_collect()} re-injects them. A load/edit/save round-trip thus
does not silently drop them.

\subsection{Population helpers --- the dropdown glue}

Three small methods make the cross-tab dropdowns work.
\code{\_pop\_names()} (\file{pipeline/ui/main.py:321}) returns the current non-blank
population names from the Populations table --- this is the \code{options\_provider}
that every population dropdown elsewhere is wired to. \code{\_pop\_size\_map()}
(\file{pipeline/ui/main.py:330}) returns \code{\{name: size\}} for populations with
a positive integer size. And \code{\_autofill\_default\_templates()}
(\file{pipeline/ui/main.py:345}) walks every Parameters row and, if its
\code{default} cell is empty and its index dropdowns name populations of known
size, fills a placeholder template (\code{'[, , ]'} for a length-3 vector,
\code{'[[, , ],[, , ]]'} for a $2\times3$ matrix) --- preserving any value the user
already typed. Note that this reaches directly into
\code{\_tbl\_parameters.\_row\_widgets} rather than \code{get\_rows()}, because it
needs to \emph{write} the widget \code{.value}.

\subsection{The ten tabs}

\code{\_build\_widgets()} (\file{pipeline/ui/main.py:381}) is the roughly
700-line constructor of the entire form. It builds, in order, the ten tab panels.
Each is either a few standalone widgets or a single \code{DynamicTable}:

\begin{longtable}{@{}>{\raggedright}p{0.16\textwidth}>{\raggedright\arraybackslash}p{0.78\textwidth}@{}}
\toprule
\textbf{Tab} & \textbf{What it builds} \\
\midrule
\endhead
1. Model & \code{\_w\_name} (\code{Text}), \code{\_w\_description}
  (\code{Textarea}), \code{\_w\_theory\_mode}
  (\code{ToggleButtons} of \code{['Temporal (ODE)', 'Spatial (PDE)']}); the
  toggle's observer calls \code{\_on\_mode\_change}. \\
2. Populations & \code{\_tbl\_populations} (name / size / description; no starter
  rows --- a scalar theory needs no populations). \\
3. Fields & \code{\_tbl\_physical} (name / population (select via
  \code{\_pop\_names}) / spatial\_dim / description; one starter row \code{phi}),
  plus the spatial-structure panel: \code{\_w\_spatial\_dim}
  (\code{BoundedIntText} 0--3), \code{\_w\_spatial\_apply} (a button that
  bulk-writes every row's \code{spatial\_dim}), \code{\_w\_boundary\_mode},
  \code{\_w\_boundary\_length}, \code{\_w\_initial\_mode}, \code{\_w\_operator\_ir}
  (\code{Checkbox}), \code{\_w\_dyson\_order}, \code{\_w\_reference\_diffusion}.
  These are grouped into \code{\_w\_spatial\_block}
  (\file{pipeline/ui/main.py:701}) so the toggle can hide them all at once. \\
4. Parameters & \code{\_tbl\_parameters} (name / index\_1 / index\_2 / domain /
  default (paste-enabled) / description; one starter row \code{mu}). The two index
  dropdowns are over \code{['\textemdash{}'] + populations} (the em-dash string is
  the ``no population'' sentinel); the default cell auto-templates when an index
  changes. \\
5. Functions & \code{\_tbl\_functions} (name / population / args / expression
  (paste) / description; no starter rows). \\
6. Kernels & \code{\_tbl\_kernels} (name / index\_1 / index\_2 / time\_expr
  (paste) / freq\_image (paste); no starter rows). \\
7. Noise & \code{\_tbl\_cgfs} (name / response\_legs / order / coefficient (paste)
  / kernel (paste); no starter rows). \\
8. Action & \code{\_w\_action = textarea\_input('Action S', ...)}
  (\file{pipeline/ui/main.py:1008}). \\
9. Mean-field & \code{\_tbl\_mfeqs} (lhs (paste) / rhs (paste) / population),
  \code{\_w\_fpi\_default} (\code{BoundedIntText} --- the default
  \code{fixed\_point\_index}), \code{\_w\_stability\_on} (\code{Checkbox}),
  \code{\_w\_seed\_box} (\code{Textarea} --- a multi-start Newton seed dict). \\
10. Defaults & \code{\_w\_k\_default}, \code{\_w\_ell\_default}
  (\code{BoundedIntText}), \code{\_w\_tau\_max}, \code{\_w\_tau\_step}
  (\code{FloatText}), and \code{\_tbl\_ext\_fields} (field / leaf\_index; two
  starter rows \code{phi[1], phi[2]}). \\
\bottomrule
\end{longtable}

The list of \code{(widget, title)} pairs is \code{\_all\_tabs}
(\file{pipeline/ui/main.py:1317}), and the \code{W.Tab} is composed from it with
numbered titles:

\begin{lstlisting}
self._all_tabs = [
    (tab_model, 'Model'), (tab_populations, 'Populations'),
    (tab_fields, 'Fields'), (tab_params, 'Parameters'),
    (tab_functions, 'Functions'), (tab_kernels, 'Kernels'),
    (tab_cgfs, 'Noise'), (tab_action, 'Action'),
    (tab_mfeqs, 'Mean-field'), (tab_defaults, 'Defaults'),
]
self._tabs = W.Tab(children=[w for w, _ in self._all_tabs])
for i, (_, t) in enumerate(self._all_tabs):
    self._tabs.set_title(i, f'{i + 1}. {t}')
\end{lstlisting}

After composing the tabs, \code{\_build\_widgets} wires the live plumbing
(\code{\_on\_populations\_changed} refreshes all dependent dropdowns; the parameter
table re-templates its defaults on an index change), builds the bottom bar,
attaches the readiness sidebar and CSS classes, registers \code{\_mark\_changed} on
every table and standalone widget, and finally calls \code{\_refresh\_save\_hint()}
and \code{\_on\_mode\_change(mark\_dirty=False)} for the initial render.

\begin{note}
The module docstring at \file{pipeline/ui/main.py:4} still calls this a
``9-section tab editor.'' That label is stale --- \code{\_all\_tabs} composes ten
tabs. When the prose and the code disagree, trust \code{\_all\_tabs}: there are
ten.
\end{note}

\section{Temporal vs.\ Spatial mode}

The Temporal/Spatial toggle on the Model tab does more than show and hide
controls; it actively rewrites form state so a saved theory carries only the
structure its kind allows. \code{\_is\_spatial\_mode()}
(\file{pipeline/ui/main.py:1522}) is simply \code{True} iff the toggle value starts
with \code{'Spatial'}. The work is in \code{\_on\_mode\_change(mark\_dirty=True)}
(\file{pipeline/ui/main.py:1525}):

\begin{itemize}
  \item \textbf{Switching to Temporal:} hide \code{\_w\_spatial\_block}, hide the
        \code{spatial\_dim} column, force every field's \code{spatial\_dim=0}, and
        \emph{clear} the operator-IR / Dyson / boundary / initial widgets, so a
        saved time-only theory carries no spatial structure.
  \item \textbf{Switching to Spatial:} show the block, show the column, bump
        \code{\_w\_spatial\_dim} to $\ge 1$, and \emph{clear the Kernels and Noise
        tables} --- spatial theories take only Gaussian white independent noise and
        have no convolution kernels, so neither can leak into a saved PDE theory.
\end{itemize}

Then \code{\_apply\_visible\_tabs()} (\file{pipeline/ui/main.py:1560}) rebuilds
\code{self.\_tabs.children} to the visible subset (Spatial hides the Kernels and
Noise tabs), renumbers the titles, and preserves the current selection when it
stays visible.

\begin{gotcha}
The clearing is \emph{intentional and silent, with no undo}. If you type kernels
or noise rows and then flip to Spatial, they are gone (\file{pipeline/ui/main.py:1548}).
If you set \code{spatial\_dim} and operator-IR and then flip to Temporal, those are
force-zeroed and cleared (\file{pipeline/ui/main.py:1534}). The rationale is that
each mode supports a strict subset of structure, but a user mid-edit can lose work.
Save before you flip if you are unsure.
\end{gotcha}

\section{The readiness sidebar}

The most valuable feature for a first-time author is the live validation sidebar.
It re-renders on every edit (via \code{\_mark\_changed}) and shows three things: a
per-tab readiness badge strip, a detailed list of findings, and a declared-names
cheat sheet. The renderer is \code{\_refresh\_validation}
(\file{pipeline/ui/main.py:2097}): a per-tab badge dot is red if any error, amber
if any warning, green otherwise --- note that \code{info}-severity findings do
\emph{not} trip a badge. The green ``No problems detected'' line only appears
once you have typed at least one field, parameter, or action, so a blank starter
form does not lie green.

\subsection{The action AST validator}

The flagship check is static validation of the action string. It is built on
Python's standard-library \code{ast} module, which parses Python \emph{source code}
into an \term{Abstract Syntax Tree} --- a tree of typed nodes (\code{ast.Name} = a
bare identifier, \code{ast.Call} = a function call, \code{ast.GeneratorExp} = a
\code{(... for i in ...)} comprehension) --- \emph{without executing it}. That last
point is the whole value: the action is analysed, never \code{eval}'d, so a typo is
caught at authoring time instead of detonating later inside Sage.

\code{\_validate\_action(self, action\_text, spec)}
(\file{pipeline/ui/main.py:1716}) first parses the string in expression mode and
catches a \code{SyntaxError} with its exact location:

\begin{lstlisting}
import ast
try:
    tree = ast.parse(text, mode='eval')
except SyntaxError as e:
    line = (e.lineno or 1); col = (e.offset or 0)
    return [('error', f'action syntax error at line {line}, col {col}: {e.msg}')]
\end{lstlisting}

It then walks the tree with a recursive \code{walk(node, bound)} that tracks which
names are \emph{bound} by comprehensions (\code{for i in A} binds \code{i}) and
lambda arguments, so an inner index is never falsely reported. Any \code{ast.Name}
that is neither a bound target nor a \emph{known} identifier is flagged as
undeclared (the first five are shown, then \code{+N more}); a \code{for i in X}
where \code{X} is not a declared population is flagged as a bad population; and
iterating over the legacy alias \code{pop} when no populations are declared is
called out as a foot-gun.

What counts as ``known'' is built by \code{\_action\_known\_names(spec)}
(\file{pipeline/ui/main.py:1663}). The set is the union of the builtins
(\code{\_ACTION\_BUILTINS}, \file{pipeline/ui/main.py:1651} --- operators
\code{Dt}, \code{Conv}; math \code{exp}, \code{log}, the trig family, \code{sqrt},
\code{abs}, \code{pi}, \code{e}, \code{I}, \code{oo}; \code{heaviside},
\code{dirac\_delta}, \code{sign}; indexing helpers \code{sum}, \code{range},
\code{len}), every declared parameter, every physical field \emph{with its three
auto companions}, every kernel and function name, and every population name (and
its \code{pop\_<name>} spelling, plus the legacy \code{pop}):

\begin{lstlisting}
for f in spec.get('physical_fields') or []:
    nm = (f.get('name') or '').strip()
    if not nm: continue
    names.add(nm)           # physical field itself
    names.add(f'{nm}t')     # response field (auto)
    names.add(f'd{nm}')     # fluctuation (auto)
    names.add(f'{nm}star')  # saddle
\end{lstlisting}

For spatial theories it additionally adds \code{Laplacian}; and \emph{only} when
the operator-IR checkbox is on, it adds the call forms \code{Lap}, \code{Dx},
\code{Gradient}, \code{GradX} (\file{pipeline/ui/main.py:1712}).

\begin{defn}
\term{AST (Abstract Syntax Tree)} --- the tree of Python syntax nodes produced by
\code{ast.parse}. Validating the action by walking this tree (rather than
executing the string) is what makes a typo like \code{hevyside} a tidy
``undeclared name'' finding in the sidebar instead of an opaque exception three
phases into the engine. Keeping \code{\_ACTION\_BUILTINS} explicit is what lets the
validator tell a real builtin from a misspelled one.
\end{defn}

\subsection{Cross-tab validation and the cheat sheet}

\code{\_validate(self)} (\file{pipeline/ui/main.py:1830}) walks the collected spec
and returns \code{\{tab\_index\_1based: [(severity, message), ...]\}} with severity
in \code{\{warn, error, info\}}. Its docstring is explicit about scope. It
\emph{does} check: undeclared-population references on fields, parameters, and
kernels (the ``silent \code{n\_populations=0} bug class''); the action AST
cross-reference; a Defaults-dict parse failure; reserved-name collisions; and
per-field mean-field-equation coverage for coupled theories. It deliberately does
\emph{not} nag about the theory name still being the default, stability being off,
or a blank action. The reserved names are \code{t}, \code{omega}, \code{Dt},
\code{delta\_D}, \code{delta\_Dp} everywhere, plus \code{k}, \code{Laplacian},
\code{x}, \code{y}, \code{z} in spatial theories. Critically, \code{\_validate}
wraps \code{\_collect()} in \code{try/except} so a malformed Defaults dict becomes
a \emph{finding} rather than a crash.

The cheat sheet, \code{\_build\_cheat\_sheet} (\file{pipeline/ui/main.py:2010}),
renders an HTML reference of all declared names by tab. Fields render as
\emph{scalars} (\code{n}, \code{nt}, \code{nstar}) when they have no valid
population, or \emph{indexed} (\code{n[i] for i in pop\_X}, response \code{nt[i]},
saddle \code{nstar[i]}) when they do --- showing you the exact spellings the action
will accept. Blank starter rows are filtered out.

\begin{gotcha}
Two blind spots are worth internalizing. First, the validator \emph{cannot} detect
a missing operator-IR toggle: it only allowlists the \code{Lap()}/\code{Dx()} call
names \emph{when the box is already on}, so a derivative-vertex action written
without ticking operator-IR validates cleanly yet computes the wrong physics (see
the next section). Second, the ``every field needs an equation'' check
(\file{pipeline/ui/main.py:1991}) is a word-boundary substring match on the LHS
text, not a parse, and only fires for $\ge 2$ fields (coupled theories) --- a field
whose name appears only inside a function call could pass spuriously.
\end{gotcha}

\section{From form to file: the spec dict}

Everything the form does on save funnels through one method. \code{\_collect(self)}
(\file{pipeline/ui/main.py:2335}) snapshots every tab into the \term{spec dict} ---
the flat dictionary the serializer consumes. It walks each table via
\code{get\_rows()}, drops blank rows, and maps the UI shapes onto the serializer's
shapes. The shape, assembled at \file{pipeline/ui/main.py:2498}, is:

\begin{lstlisting}
{
  'name':            str,                 # _w_name
  'populations':     [{'name','size','description'?}, ...],
  'n_populations':   int,                 # derived len(populations)
  'description':     str,
  'response_fields': [],                  # always empty -- auto-generated downstream
  'physical_fields': [{'name','population'?,'spatial_dim'?,'latex'?,'description'?}, ...],
  'parameters':      [{'name','indexed_by'?,'domain'?,'default'?,'description'?}, ...],
  'functions':       [{'name','args':[...],'expression','population'?,...}, ...],
  'kernels':         [{'name','indexed_by'?,'time_expr'?,'freq_image'?,'latex_name'?}, ...],
  'cgf_terms':       [{'name','response_field','response_legs'?,
                       'order','coefficient','kernel'}, ...],
  'action_text':     str,                 # the action textarea
  'equations':       [{'lhs','rhs','population'(or None)}, ...],
  'stability_analysis': bool,
  'default_fundamental': {},              # deliberately empty
  'metadata':        {...},               # see below
  # spatial-only keys, present ONLY when a field is spatial:
  'boundary':            {'mode': 'infinite'|'periodic', 'length'?: float|str},
  'initial':             {'mode': 'stationary'},
  'operator_ir':         True,            # only when ticked
  'dyson':               {'mode':'fixed','order':N},   # only when order>0
  'reference_diffusion': float,           # only when set
}
\end{lstlisting}

A small nested helper, \code{\_index\_list(row)}
(\file{pipeline/ui/main.py:2352}), translates the two index dropdowns (where
the em-dash string means none) into the \code{indexed\_by} list: \code{[]} for a scalar,
\code{[pop]} for a vector, \code{[pop1, pop2]} for a matrix. Parameter defaults are
parsed from their text cell into Python objects (with a raw-string fallback on
failure), the Noise rows are normalized by \code{\_normalize\_cgf\_rows}
(\file{pipeline/ui/main.py:249}), and the spatial blocks are emitted \emph{only
when at least one field is spatial} --- matching the serializer's
``emit-only-when-present'' rule.

\begin{gotcha}
Two spots use the \emph{raw} \code{eval} rather than \code{ast.literal\_eval} to
parse a default --- the \code{\_eval\_dict} helper inside \code{\_collect}
(\file{pipeline/ui/main.py:2342}) and the per-parameter default cell
(\file{pipeline/ui/main.py:2427}):
\begin{lstlisting}
obj = eval(text, {'__builtins__': {}}, {})
\end{lstlisting}
Passing \code{\{'\_\_builtins\_\_': \{\}\}} as globals strips access to
\code{open}, \code{\_\_import\_\_}, and friends, so a default like
\code{[[1,0],[0,1]]} evaluates fine but \code{\_\_import\_\_('os')} raises
\code{NameError}. This is \emph{weaker} than \code{ast.literal\_eval} --- it still
evaluates arbitrary operators --- but the threat model is the user's own typos on
the user's own machine, so it is acceptable. Do not mistake it for safe evaluation
of untrusted input. The Defaults seed-box, by contrast, \emph{does} use
\code{ast.literal\_eval} (\file{pipeline/ui/main.py:2586}).
\end{gotcha}

The run-metadata sub-dict comes from \code{\_collect\_metadata}
(\file{pipeline/ui/main.py:2570}):

\begin{lstlisting}
{
  'k_default': int, 'ell_default': int, 'tau_max': float, 'tau_step': float,
  'recommended_external_fields': [(field, leaf_index), ...],   # only if non-empty
  'fixed_point_index_default': int,                            # always written
  'seed_box_default': {var: (low, high), ...},                 # only if provided
}
\end{lstlisting}

Note that \code{fixed\_point\_index\_default} is written \emph{even when 0}
(\file{pipeline/ui/main.py:2620}), so a reload is unambiguous; most other optional
keys are emitted only when non-default. And \code{default\_fundamental} is
deliberately always emitted empty (\file{pipeline/ui/main.py:2544}) --- the
per-parameter \code{default=} declarations now carry the suggested values, and a
loaded file's \code{DEFAULT\_FUNDAMENTAL} is ignored on load.

\subsection{The serialized file}

When the form is filled like the placeholder --- one field \code{x} in population
\code{pop}; parameters \code{mu}, \code{eps}, \code{D}; the OU-quartic action ---
\code{\_collect()} yields a spec whose serialized form is exactly the on-disk file
the quickstart chapter loads:

\begin{lstlisting}
def build():
    return (
        TemporalTheoryBuilder('OU Quartic (white noise)')
        .population('pop', size=1)
        .physical_field('x', population='pop', description='variable')
        .parameter('mu', default=1.0, domain='positive')
        .parameter('eps', default=0.02, domain='positive')
        .parameter('D', default=1.0, domain='positive')
        .set_action_text('sum(xt[i]*((Dt+mu)*x[i] + eps*x[i]^3) - D*xt[i]^2 for i in pop)')
        .equation(lhs='(Dt+mu)*x[i]', rhs='-eps*x[i]^3', population='pop')
        .build()
    )
DEFAULT_FUNDAMENTAL = {'mu': 1.0, 'eps': 0.02, 'D': 1.0}
METADATA = {'k_default': 2, 'ell_default': 0,
            'recommended_external_fields': [('dx', 1), ('dx', 1)],
            'tau_max': 8.0, 'tau_step': 0.5, 'fixed_point_index_default': 0}
\end{lstlisting}

The serializer chooses \code{SpatialTheoryBuilder} versus
\code{TemporalTheoryBuilder} by whether any field carries \code{spatial\_dim} $\ge
1$. The saddle parameter \code{xstar} is \emph{not} emitted --- \code{build()}
auto-creates it from the field declaration. Compare this against the hand-written
file in Chapter~\ref{ch:quickstart}: they are identical, which is the point.

\begin{example}\label{ex:ui-roundtrip}
\textbf{The round trip.} Fill the form for the OU-quartic model and click
\textbf{Save}; you get the file above. Now in a fresh \code{TheoryUI()}, pick that
file in the Load dropdown and click \textbf{Load}: \code{load\_spec\_from\_file}
parses the file's \emph{source AST} (it never executes \code{build()}), the spec
dict is reconstructed, and \code{load(...)} clears and repopulates every tab. The
action text, equation RHS strings, function expressions, and kernel expressions
survive verbatim because they are stored as string literals; comments and
formatting are lost (it is an AST round-trip, not a textual one). Edit a parameter,
save again, and the file is regenerated --- closing the loop.
\end{example}

\section{Loading an existing theory}

\textbf{Load} reverses the save trip. \code{\_on\_load}
(\file{pipeline/ui/main.py:2769}) is gated by the dirty guard (below), refreshes
the dropdown so freshly-saved files appear, calls \code{load\_spec\_from\_file} to
parse the file into a spec dict, and hands it to \code{self.load(spec)}.

\code{load(self, spec\_or\_path)} (\file{pipeline/ui/main.py:2813}) is the big
repopulate-the-form routine, and most of its length is historical-wrinkle handling:

\begin{itemize}
  \item It strips the auto-prefix: a field stored as \code{dn} with
        \code{natural\_name='n'} displays as \code{n}.
  \item It stashes any \code{latex} strings into \code{\_loaded\_extras} for
        round-trip preservation.
  \item It \emph{auto-skips saddle parameters} (those with \code{mean\_field=True}
        or a name matching \code{<natural>star}) --- the framework re-creates them
        from the field declarations, so they must not be re-added as ordinary
        parameters.
  \item It maps \code{indexed\_by} back to the two index dropdowns, and converts a
        legacy \code{mf\_equations} shape (\code{\{saddle, rhs\}}) into the new
        \code{equations} shape (\code{\{lhs, rhs, population\}}) by stripping the
        trailing \code{star} and looking up the field's population.
  \item Finally it restores \code{stability\_analysis} and the spatial widgets,
        picks the Temporal/Spatial mode from the loaded fields and applies it, then
        drives every Defaults-tab widget from \code{metadata}.
\end{itemize}

\section{The bottom bar: Save, Preview, Pre-compute, and the rest}

\code{\_build\_bottom\_bar} (\file{pipeline/ui/main.py:1582}) composes the
save-path field, a derived-filename hint, the four primary action buttons, and the
destructive-action gate (two confirm checkboxes plus the Reset and Load rows). The
button handlers all write into the shared \code{\_status} panel.

\paragraph{Preview} \code{\_on\_preview} (\file{pipeline/ui/main.py:2640}) runs
\code{\_collect()} $\to$ \code{render\_theory\_file(spec)} and prints the resulting
\code{.theory.py} source into the status panel without writing anything --- a
look-before-you-leap.

\paragraph{Save} \code{\_on\_save} (\file{pipeline/ui/main.py:2650}) collects the
spec, resolves the target path (an explicit field value, or \code{theories\_dir},
joining a relative non-\code{.py} base), calls \code{save\_theory\_to\_file},
records \code{last\_saved}, clears \code{\_dirty} and both confirm checkboxes, and
prints the path plus the \code{THEORY\_NAME = <slug>} line you paste into the
runner. Both \code{\_collect()} and the write are wrapped so a failure prints a
humanized error rather than a traceback.

\paragraph{Reset} \code{\_on\_reset} (\file{pipeline/ui/main.py:2689}) is gated by
the dirty guard; on confirm it re-runs \code{\_build\_widgets()} to wipe state
cleanly, clears \code{\_dirty}, and re-\code{show()}s.

\paragraph{Open in runner} \code{\_on\_open\_in\_runner}
(\file{pipeline/ui/main.py:2231}) is the bridge to the runner notebook. If nothing
has been saved it prints a hint; otherwise it derives the slug from
\code{last\_saved}, writes it (one line) to \file{.theories/.last\_built}, and
\code{display}s a \code{Javascript} snippet:

\begin{lstlisting}
display(Javascript("window.open('theory_runner.ipynb', '_blank');"))
\end{lstlisting}

The runner reads \file{.theories/.last\_built} when its \code{THEORY\_NAME} is
blank, so the just-saved theory is pre-selected. It degrades to a printed
instruction when JS is unavailable (for example under an \code{nbconvert} export).

\subsection{The dirty guard}

Two pieces protect unsaved work. \code{\_mark\_changed}
(\file{pipeline/ui/main.py:2204}) sets \code{\_dirty=True} and refreshes the
sidebar; it is the observer registered on every input. \code{\_check\_dirty\_guard}
(\file{pipeline/ui/main.py:2217}) returns \code{True} (action allowed) only if the
form is clean, or if the matching confirm checkbox is ticked. Both Reset and Load
route through it, so they cannot silently destroy unsaved edits --- you must
explicitly tick ``discard unsaved changes'' or save first.

\subsection{Pre-compute: the one place the UI touches the engine}

\code{\_on\_precompute} (\file{pipeline/ui/main.py:2704}) is the single exception
to ``the UI does no physics.'' It lets you ask, before leaving the form, ``is this
theory structurally healthy?'' --- without writing a file. The sequence is:

\begin{lstlisting}
src = render_theory_file(spec)                 # render to source ...
ns: dict = {}
exec(compile(src, '<precompute-spec>', 'exec'), ns)  # ... exec in a fresh namespace ...
model = ns['build']()                          # ... call build() for an in-memory model
result = precompute(model, verbose=True)       # the structural sanity pass
\end{lstlisting}

\code{pipeline.precompute(model)} (a thin primitive in
\file{pipeline/\_precompute.py}) does one cheap structural pass: Taylor-expand the
action at order~2, verify the mean-field saddle, and build and cache the
propagator. It returns a result dict from which the handler prints a compact
verdict --- the mean-field check (a \code{PASS} mark when \code{mf\_check == 'PASS'}),
the solved saddle values, whether the propagator was \code{cached} or
\code{FAILED}, the cache directory, and the wall time. A green PASS means the
theory builds, has a real saddle, and has a constructible propagator --- the most
common authoring mistakes (an undeclared symbol, a saddle that does not solve, an
ill-posed bilinear part) are caught here, in seconds, before you ever open the
runner.

\begin{gotcha}
\textbf{Fork safety is not a concern here.} The macOS fork-crash documented in this
project's memory lives in the spatial and temporal \emph{integrator} paths, not the
UI. The UI runs entirely on the kernel's main thread with no multiprocessing, and
\code{Pre-compute} runs \code{precompute} in-process. You can click it freely in a
notebook on any platform.
\end{gotcha}

\section{Gotchas to keep in mind}

Most caveats have appeared in context above; this section collects the ones most
likely to bite, in rough order of severity.

\begin{itemize}
  \item \textbf{Operator-IR is load-bearing and silent.} Per the serializer
        (\file{pipeline/theory\_serialize.py:455}), without \code{.operator\_ir()}
        a re-saved KPZ/Burgers/Model-B theory \emph{silently loses} the
        \code{Dt()}/\code{Lap()}/\code{Dx()} lowering and computes the wrong
        physics. The UI emits it only when the checkbox is ticked, and the
        validator does \emph{not} detect a missing tick. If your action has a
        spatial derivative \emph{inside} a nonlinearity, you must tick Operator-IR.
  \item \textbf{Mode-switch data loss.} Temporal~$\to$~Spatial clears the Kernels
        and Noise tables; Spatial~$\to$~Temporal force-zeros every
        \code{spatial\_dim} and clears operator-IR / Dyson / boundary. Both are
        intentional and have no undo. Save before flipping.
  \item \textbf{The \code{eval} sandbox is weak.} Parameter defaults are parsed
        with \code{eval(text, \{'\_\_builtins\_\_': \{\}\}, \{\})}, which blocks
        name lookups but still evaluates operators. Fine for your own typos; not
        safe-evaluation of untrusted input.
  \item \textbf{\code{latex} carry-through is keyed by name.} A loaded \code{latex}
        string survives a round trip via \code{\_loaded\_extras}, but if you
        \emph{rename} the field it was attached to, the latex is dropped.
  \item \textbf{Clipboard reads the kernel's machine.} The clipboard buttons exist
        because VS~Code swallows Ctrl/Cmd-V; on a remote kernel they read the
        server's clipboard (or show a cross mark).
  \item \textbf{Open-in-runner needs JS and relative paths.}
        \code{window.open('theory\_runner.ipynb', ...)} assumes the runner sits next
        to the builder in the Jupyter file tree and that JS is allowed; otherwise it
        prints an instruction.
  \item \textbf{The exported \code{vector\_input} / \code{matrix\_input} /
        \code{expression\_input} are unused by the live form.} They are reusable
        primitives; editing them changes nothing you see in the UI.
\end{itemize}

\section{What you just learned, and where to go next}

The Theory-Builder UI is the only chapter in this manual that produces no numbers.
Its job is upstream of everything else: it is a typed, validated, point-and-click
authoring surface whose single output is a \code{.theory.py} file identical to one
you could write by hand. You saw the kernel\,$\leftrightarrow$\,browser widget
model that makes a live form possible inside a notebook cell; the
\code{DynamicTable} primitive that every tabular tab is made of; the ten tabs and
the \msrjd{} piece each one declares; the Temporal/Spatial toggle and what it
rewrites; the AST-based action validator that catches typos before they reach the
engine; the \code{\_collect()} method and the spec-dict shape it produces; the
serializer round trip via \code{load\_spec\_from\_file}; and the bottom-bar buttons,
including the in-process \code{Pre-compute} sanity pass.

From here the natural next step is to take a file the UI produced and run it.
Chapter~\ref{ch:quickstart} drives exactly such a file end to end; the theory-files
chapter (Chapter~\ref{ch:theory-files}) documents the \code{build()} /
\code{DEFAULT\_FUNDAMENTAL} / \code{METADATA} shapes the UI emits, from the
serializer's side; and the serializer's own internals --- \code{render\_theory\_file},
\code{save\_theory\_to\_file}, \code{load\_spec\_from\_file} --- live in
\file{pipeline/theory\_serialize.py}, the module on the far side of every Save and
Load button you just met.


% =====================================================================
%  APPENDICES
% =====================================================================
\appendix
\part{Appendices}
% ---------------------------------------------------------------------
%  Appendix A : Glossary
%  A single alphabetised glossary synthesised from the per-subsystem
%  briefs.  Every term that appears anywhere in the manual is defined
%  here from scratch, for a reader who knows the physics but not the
%  specific tooling (SageMath, nauty, sympy, numba, ...).
% ---------------------------------------------------------------------
\chapter{Glossary}\label{ch:glossary}

This appendix collects, in one alphabetised list, every piece of jargon the
manual leans on. It is deliberately written so that you can land here from a
cross-reference in any chapter and walk away understanding the term \emph{cold},
without having read the chapter it came from. Three kinds of vocabulary are mixed
together on purpose, because in practice you meet them mixed together:

\begin{itemize}
  \item \textbf{Physics and field theory} --- the \msrjd{} action, propagators,
        cumulants, loop order, symmetry factors, Symanzik polynomials. A reader
        who knows stochastic field theory will recognise these but may not know
        exactly how \Daedalus{} names or splits them.
  \item \textbf{Mathematical machinery} --- Schwinger parameters, divided
        differences, the Hermite--Genocchi formula, causal posets, chain
        simplices, Cholesky and Lyapunov solves. These are the tools the loop
        integrals actually reduce to.
  \item \textbf{Software tooling} --- SageMath, \code{nauty}, \code{sympy},
        \code{numba}, \code{scipy}, \code{ipywidgets}, and the dozens of small
        library calls (\code{fast\_callable}, \code{lambdify},
        \code{solve\_continuous\_lyapunov}, \code{importlib.util}, \ldots) that
        do the work. This is the layer the manual promises to over-explain, so it
        gets the most attention here.
\end{itemize}

\noindent
A handful of the most load-bearing entries (the \msrjd{} action, the symmetry
factor, the Symanzik polynomials, Phase~J) are given a boxed \emph{Definition}
rather than a one-line gloss, because you will meet them on nearly every page.
Where a term names a concrete function, class, or attribute in the codebase, the
defining file is cited so you can read the implementation. Identifiers are set in
\code{monospace}; the field-theory ``response'' field is written
$\tilde\varphi$ in the maths and, by the engine's own convention, with a
\code{t}- or \code{d}-prefix in code (\code{xt}, \code{dn}).

The entries are grouped by first letter. Symbols and Greek-letter terms are
filed under the letter of their romanised name (so $\Sym(\Gamma)$ appears under
\textbf{S}, for ``symmetry factor''; $\kappa^{(n)}$ under \textbf{C}, for
``cumulant'').

% =====================================================================
\section*{A}
% =====================================================================
\begin{description}[leftmargin=0pt,style=nextline]

  \item[\code{abs\_symbolic}] Sage's symbolic absolute-value function. It keeps an
        expression such as $|\tau|$ \emph{unevaluated} as a symbol, rather than
        forcing a numeric branch. The colored-noise kernels need this because the
        Lorentzian $\ee^{-|\tau|/\tau_c}$ must survive as an even function of the
        time lag all the way into the Fourier transform.

  \item[Action $S$] The functional in the exponent of the \msrjd{} path integral,
        $\ee^{-S[\varphi,\tilde\varphi]}$. It is the central object the user
        authors --- as a text string passed to \code{.set\_action\_text(...)}. Its
        stationary point is the classical (mean-field) solution; its fluctuations
        about that point generate every correlation function. The bilinear part is
        the free propagator; the cubic-and-higher part supplies the interaction
        vertices.

  \item[\code{add\_vararg}] Sage's internal $n$-ary addition operator. It is
        \emph{not} Python's \code{operator.add} (which is binary); when the
        field-theory core walks a symbolic expression tree it detects a sum by
        matching this operator by name.

  \item[Adjugate $\mathrm{adj}(K)$] The transpose of the cofactor matrix of a
        square matrix $K$. It gives the inverse through Cramer's rule,
        $K^{-1}=\mathrm{adj}(K)/\det(K)$, which is how the propagator
        $G=K^{-1}$ is formed symbolically before any numbers are substituted.

  \item[Allen--Cahn / reaction--diffusion] The prototypical first-tier spatial
        model: a field with a relaxational mass $\mu$ and a diffusion term
        $D\nabla^{2}$. Its inverse propagator has the form $A + B k^{2} + \ii\omega$.

  \item[\code{analytic-eligible}] A property of a $\delta$-edge subset (temporal)
        or a diagram (spatial): the loop integral has a closed form and needs no
        \code{scipy.nquad} fallback. A subset is analytic-eligible exactly when no
        smooth-kernel noise vertex (a \code{NoiseSourceType} with a non-rational
        kernel) is present; otherwise the integrand is non-rational and the
        adaptive numerical path is forced.

  \item[\code{ast} (Abstract Syntax Tree)] Python's standard-library module for
        parsing source text into a tree of syntax nodes \emph{without running it}.
        The UI and the theory loader use it to validate an action string --- they
        walk the parsed tree to check that only allowed names appear, never
        executing the code.

  \item[\code{ast.literal\_eval}] A safe evaluator that turns a string into a
        Python value \emph{only} if the string is a literal (number, string, list,
        dict, tuple, bool, or \code{None}). It never executes function calls or
        arbitrary code, which is why it is used for the random-seed box and similar
        user inputs.

  \item[Automorphism group $\Aut(\Gamma)$] The group of relabellings of a graph
        $\Gamma$ that map it onto itself edge-for-edge (preserving any colours).
        Two variants matter for symmetry factors: \code{Aut\_fixed\_ext} pins the
        external leaves in place, while \code{Aut\_free\_ext} (also written
        \code{Aut\_free}) lets same-field external leaves permute. Their ratio is
        the external-Wick compensation.

  \item[Auxiliary field (aux field)] An extra field injected by Markovian
        embedding to represent colored noise: the Ornstein--Uhlenbeck field
        $\xi$ (\code{xi}) together with its \msrjd{} response partner
        $\tilde\xi$ (\code{xit}). Adding it enlarges the state space so the dynamics
        become local in time.

\end{description}

% =====================================================================
\section*{B}
% =====================================================================
\begin{description}[leftmargin=0pt,style=nextline]

  \item[Betti number (first) / loop number] For a connected graph,
        $|E|-|V|+1$ --- the number of independent cycles, equivalently the loop
        order $\ell$ of the corresponding diagram.

  \item[Bigrade $(n_{\tilde{}},\,n_{\mathrm{phys}})$] The grading of an action
        monomial by $(\text{number of response factors},\ \text{number of
        physical factors})$; the total degree is their sum. This single pair is
        the primary classification axis of the expanded action: the saddle sector
        is $(0,0),(1,0),(0,1)$; the free propagator kernel is $(1,1)$; the white
        noise lives at $(2,0)$; any monomial of total degree $\ge 3$ is an
        interaction vertex. In the \code{FieldTheory} object the attribute
        \code{\_by\_tp} maps each bigrade to the sum of its monomials.

  \item[Bilinear vs vertex generator] In the operator-IR authoring path, a
        distinction between IR generators that fold into the propagator (total
        degree $\le 2$, the bilinear part) and those that carry per-leg momentum
        form factors forward into the loop integrator (total degree $\ge 3$, the
        vertices).

  \item[\code{bliss}] A graph-automorphism engine, an alternative to \code{nauty};
        Sage may call either internally for \code{automorphism\_group}. (See
        \code{nauty}.)

  \item[Bubble] A one-loop self-energy insertion with two internal vertices joined
        by two edges; its loop momentum couples to the external momentum $q$. In a
        multigraph a bubble is a double edge between a vertex pair.

\end{description}

% =====================================================================
\section*{C}
% =====================================================================
\begin{description}[leftmargin=0pt,style=nextline]

  \item[Canonical form / canonical label] A deterministic relabelling of a graph's
        vertices such that two graphs are isomorphic if and only if their canonical
        forms are \emph{identical}. This converts the expensive ``are these the
        same diagram?'' question into a cheap string comparison. Sage exposes it as
        \code{DiGraph.canonical\_label}; the routine and the original$\to$canonical
        id map (the \emph{certificate}, \code{cert}) come from \code{nauty}.

  \item[Catastrophic cancellation] The loss of significant digits that occurs when
        you subtract two nearly-equal large floating-point quantities. It is the
        reason the symbolic adjugate residues are unreliable close to a pole, and
        why the propagator code switches to an exact rational field there.

  \item[Causal chamber / chamber] One ordering region of the internal-vertex-time
        integral, corresponding to a single linear extension of the retarded
        partial order. Within a chamber every time-difference sign is fixed, so the
        integrand is smooth and the integral is closed-form. The full loop integral
        is the (disjoint) sum over all chambers. (Spatial pipeline; the temporal
        analogue is the \emph{polytope}/\emph{chain-simplex} decomposition.)

  \item[Causal filtering] Discarding the typed diagrams that violate the retarded
        \msrjd{} propagator structure --- those that would require an effect before
        its cause. Implemented in \code{filter\_causal}
        (\file{msrjd/diagrams/causality.py:126}). A causal diagram is structurally
        a directed acyclic graph (DAG).

  \item[Causal poset / causal polytope] For three or more surviving integration
        times, the set of retardation constraints $\Delta t_e > 0$ read as a
        partial order ($s_u < s_v$). The convex region it carves out is the
        \emph{causal polytope}, the integration domain; integrating over its
        linear extensions (nested \emph{chain simplices}) evaluates the polytope
        integral.

  \item[\code{CDF} (Complex Double Field)] Sage's type for machine-precision
        (hardware \code{double}) complex numbers. The pipeline drops into \code{CDF}
        for the fast numeric phases --- finding poles, evaluating residues --- after
        the exact symbolic work is done.

  \item[\code{C} edge / correlation line] An edge carrying a
        $\langle\varphi\varphi\rangle$ correlation. In the heat-kernel
        representation it is two retarded ($\code{R}$) segments glued at a
        two-point noise source, parameterised by a single Schwinger time
        $\sigma\in[0,\infty)$. Contrast the \code{R} edge.

  \item[Certificate (\code{cert})] \code{nauty}'s original-to-canonical vertex id
        map produced alongside the canonical label; equal canonical forms (and
        their certificates) certify two diagrams as isomorphic. (See canonical
        form.)

  \item[CGF / CGF term / CGF row] A Cumulant-Generating-Functional term: a single
        declared contribution to the noise's cumulants, given as a tuple
        \code{(coefficient, kernel, order, response\_legs)} via
        \code{.declare\_cgf\_term(...)}. Order $2$ is the noise covariance; orders
        $\ge 3$ are non-Gaussian noise. ``CGF'' is also the internal name of the
        UI's Noise tab.

  \item[Chain simplex] A nested-inequality region
        $L \le s_1 \le \dots \le s_N \le U$. The exponential integral over a chain
        simplex has a closed form with $2^{N}$ terms; the temporal $m\ge 3$
        integrator sums over the chain simplices that tile the causal polytope.

  \item[Characteristic polynomial / \code{D\_omega}] $\det(K_{\mathrm{ft}})$, the
        determinant of the Fourier-image kernel. The roots of its numerator are the
        system's poles $\omega_k$.

  \item[Cholesky factorization] The decomposition $C=LL^{\mathsf T}$ of a symmetric
        positive-semidefinite matrix. Multiplying a vector of independent unit
        Gaussians by $L$ yields correlated Gaussians with covariance $C$ --- the
        trick a simulator uses to draw correlated noise (the explicit $2\times2$
        form gives $\eta_y=\rho u + \sqrt{1-\rho^{2}}\,v$).

  \item[Cofactor expansion] Computing a determinant or adjugate by expanding along
        minors. The dense NumPy path forms each minor with \code{np.delete}.

  \item[Coloured incidence digraph] The bipartite graph built from a diagram so
        that an isomorphism engine can see edge types: vertices and \emph{edges}
        both become nodes, each coloured by its field / vertex / propagator type.
        This is the object fed to \code{nauty} for canonical labelling and
        automorphisms. (Also ``incidence digraph (coloured)''.)

  \item[Colored noise] Noise with a finite correlation time $\tau_c$; its
        covariance is a smooth function of the time lag (here the Lorentzian
        $c\,\ee^{-|\tau|/\tau_c}$) rather than a $\delta$-function. Contrast white
        noise.

  \item[Completeness bound $(k+3\ell-j-1)$] The proven upper bound on the number of
        degree-2 vertices that a valid topology can have; respecting it guarantees
        the tree search misses no diagram. It sets \code{v2\_max}. (An older
        $-\lfloor j/3\rfloor$ form of the lemma was false at $\ell\ge 3$; the
        bound above is the corrected, proven one.)

  \item[Composite vs per-leg vertex] Two modes of a derivative interaction vertex.
        A \code{'composite'} vertex applies the operator to the whole
        $\varphi^{n}$ composite, giving a response-leg momentum factor (e.g.
        $\nabla^{2}(\varphi^{2})$ in Model~B, $\partial_x(\varphi^{2})$ in
        Burgers). A \code{'perleg'} vertex applies the operator to each physical
        leg separately (e.g. $(\partial_x\varphi)^{2}$ in KPZ, giving
        $\prod \ii\,p_{\mathrm{leg}}$). The two produce different form factors.

  \item[Convolution \code{Conv(g, n)}] The formal time-convolution operator
        $\int g(t-s)\,n(s)\,\dd s$. It is asymmetric --- only its second argument
        carries time content --- and each downstream consumer reduces it
        differently. Reduced symbolically via Sage's \code{substitute\_function}.

  \item[Convolution theorem] $(g\star n)\widehat{\ }(\omega)=\hat g(\omega)\,\hat
        n(\omega)$: convolution in time is multiplication in frequency. This is why
        a kernel can be left as an abstract symbol and only resolved to its
        frequency image $\hat g(\omega)$ at Fourier-transform time.

  \item[Crank--Nicolson / semi-implicit] A trapezoidal-rule implicit time-stepping
        scheme. In the coupled reaction--diffusion simulator it delivers the
        $O(\dd t^{2})$-accurate stationary covariance through a Pad\'e$(1,1)$
        propagator.

  \item[Cumulant / connected correlation function] A connected $k$-point function
        $\kappa_k=\langle\varphi\cdots\varphi\rangle_c$ --- the sum of all
        \emph{connected} Feynman diagrams with $k$ external legs. This is the
        native output of \code{compute\_cumulants}. The second cumulant of the
        noise is its covariance $\langle\xi\xi\rangle$.

  \item[$\kappa^{(n)}$ (noise cumulant)] The $n$-th cumulant of a noise process:
        $\kappa^{(2)}$ is Gaussian (the covariance), $\kappa^{(3)},\kappa^{(4)},\dots$
        are higher (non-Gaussian) corrections. Correlated noise contributes a
        cumulant series $-W_m[\tilde m]$ to the action.

  \item[\code{\_cumulant\_kernels}] A namespace dict of noise-cumulant kernel
        callables. It is a \emph{runtime} side effect of
        \code{\_build\_cumulant\_action} and is deliberately not pickled --- it is
        rebuilt when a cache is loaded.

  \item[\code{CyclotomicField(4)} $=\mathbb{Q}(\ii)$] The exact number field
        obtained by adjoining the imaginary unit to the rationals. Working in
        $\mathbb{Q}(\ii)$ (rather than floating point) gives reliable polynomial
        GCDs, which is why it is used for the canonical, exact propagator
        inversion.

  \item[\code{cysignals} (\code{alarm} / \code{SignalError} / \code{AlarmInterrupt})]
        Cython tools that let Python time-bound and recover from long-running
        \emph{native} routines (Singular, Pynac). They are how the propagator code
        puts a wall-clock budget on a symbolic simplification that might otherwise
        hang. (See also \code{SIGALRM}.)

\end{description}

% =====================================================================
\section*{D}
% =====================================================================
\begin{description}[leftmargin=0pt,style=nextline]

  \item[DAE (differential-algebraic equation)] A system mixing differential
        equations (those carrying the time-derivative \code{Dt}) with algebraic
        constraints (those without). The mean-field saddle is the solution of the
        DAE obtained from \code{.equation(lhs, rhs, population)} after substituting
        \code{Dt} $\to 0$.

  \item[DAE saddle solve / fixed-point index] A multi-root Newton solver for the
        algebraic mean-field system. When a theory has several stable saddles (e.g.
        a double well), \code{fixed\_point\_index} selects which sorted root the
        loop expansion is built around.

  \item[DAG (Directed Acyclic Graph)] A directed graph with no directed cycle. A
        causal (retarded) diagram is structurally a DAG; this is the graph-theoretic
        face of causality.

  \item[\code{D\_delta} ($\delta$-coefficient / instantaneous part)] The
        $\omega\to\infty$ limit of the Fourier propagator $\hat G(\omega)$ --- the
        matrix weight of an instantaneous $\delta(t)$ term in the time-domain
        propagator. It is the ``immediate response'' piece.

  \item[DC gain $\hat g(0)$] A kernel's integral $\int g(t)\,\dd t$, i.e. its
        Fourier transform at zero frequency. For a unit-integral (normalised)
        kernel, $\hat g(0)=1$.

  \item[Dedup multiplicity] The size of a diagram's isomorphism-equivalence class.
        Under the framework's ``Path~A'' convention this number is \emph{diagnostic
        only}: it is reported but must not be multiplied back into the diagram's
        weight, because the symmetry factor already accounts for it.

  \item[Degree] The number of edge-ends meeting at a vertex; a multi-edge counts
        with its multiplicity.

  \item[$\delta$-edge / smooth-edge subset] The $2^{|E|}$-fold expansion of a
        product of retarded propagators, choosing on each edge whether to take its
        instantaneous ($\delta$) part or its smooth (exponential) part. Each choice
        is a ``subset''.

  \item[$\delta$-elimination] Using a $\delta$-edge --- which imposes $\Delta t=0$,
        a linear equation among the times --- to remove one integration variable.
        Solved symbolically with \code{sage\_solve}.

  \item[Derivative symbol] A renamed partial-derivative-at-saddle atom. Sage's
        formal differentiation produces clumsy ``derivative-of-function-at-0''
        objects; the field-theory core renames them to clean symbols such as
        \code{phi1\_2} ($=\varphi'_2(0)$) or \code{phi21\_1}
        ($=\partial^{2}_{x_1}\partial_{x_2}\varphi_1(0)$).

  \item[Derived generator] In the operator IR, a fresh polynomial-ring generator
        introduced to stand for $\mathrm{Op}(\delta\varphi)$ --- the trick
        $u=\delta\varphi,\ v=\nabla^{2}\delta\varphi$ --- so that multivariate
        Taylor expansion treats the derivative quantity as if it were an ordinary
        field.

  \item[\code{diagram\_signature}] The canonical form of a typed diagram's coloured
        incidence digraph; a \emph{complete} isomorphism invariant. Two diagrams are
        duplicates if and only if their signatures match
        (\file{msrjd/diagrams/symmetry.py:467}).

  \item[Diffusion matrix $\mathcal{D}$] The $k^{2}$ part of the linearised inverse
        propagator of a multi-field spatial theory; it encodes spatial spreading.
        $\mathcal{D}\propto I$ means scalar (equal) diffusion across fields.

  \item[\code{DiGraph} / \code{Graph}] Sage's directed and undirected graph classes.
        Diagram topologies are built and manipulated as these objects; their
        \code{canonical\_label} and \code{automorphism\_group} methods are the
        bridge to \code{nauty}.

  \item[\code{dirac\_delta(tau)}] Sage's symbolic Dirac delta. It is the white-noise
        kernel --- the time-domain covariance of $\delta$-correlated noise.

  \item[Dirty flag / dirty guard] In the UI, the boolean \code{\_dirty} that tracks
        unsaved edits. Reset and Load are gated behind a confirmation checkbox while
        the form is dirty, so you cannot silently discard changes.

  \item[{Divided difference $f[x_0,\dots,x_n]$}] A classical finite-difference object.
        In the coupled spatial path the $n$-fold time-ordered (simplex) integral
        equals $(-1)^{n}$ times the $n$-th divided difference of
        $f(z)=\ee^{-tz}$ --- which is what makes the Dyson dressing computable in
        closed form.

  \item[Downgrade (filter)] Keeping only the bigrades of total degree
        $\le \text{target}$ from a higher-order cached expansion. It is exact by the
        superset theorem, and is how a cache built at one Taylor order serves a
        cheaper run.

  \item[Drift $V$] The $k^{1}$ (advection) coefficient of a spatial inverse
        propagator, $\ii v$ for a velocity $v$; it is zero for even / Laplacian-only
        kernels.

  \item[\code{'dn'} / \code{'dv'} / \code{'dm'}] Estimator field-type labels in the
        simulators: spike-train fluctuation (factorial-corrected),
        voltage fluctuation (smooth, linearly centred), and external-drive spike
        fluctuation (shot-noise-like, treated as \code{'dn'}).

  \item[Duhamel convolution] The time-ordered nested integral that expresses a
        perturbed time-evolution as a series of insertions sandwiched between
        unperturbed evolutions. The Dyson--Duhamel series is built from these.

  \item[\code{Dt}] The formal time-derivative operator symbol, usable directly in
        action and equation text. It is substituted to $0$ at the stationary saddle
        and to the eigenvalue $\sigma$ in linear-stability analysis.

  \item[Dyson dressing / Dyson--Duhamel series] Resumming a self-energy (or residual
        diffusion) series into the propagator, order by order, to a finite
        \code{dyson\_order}. Each order is an $n$-fold Duhamel convolution. Used for
        coupled fields with \emph{unequal} diffusivities, where the residual
        diffusion $\hat{\mathcal{D}}$ is the perturbation. (Also ``Dyson--Duhamel
        truncation''.)

\end{description}

% =====================================================================
\section*{E}
% =====================================================================
\begin{description}[leftmargin=0pt,style=nextline]

  \item[\code{erf}-split closed form (Rescue A)] The antiderivative of
        $s^{-1/2}\ee^{-\beta/s-\alpha s}$, written as a sum of two error functions
        with $\ee^{\pm 2\sqrt{\alpha\beta}}$ weights. It is the spatial heat-kernel
        integrator's closed form for a Schwinger integral that would otherwise need
        quadrature.

  \item[ETD1 (exponential time-differencing, order 1)] A time-stepping scheme for
        stiff equations $\partial_t\varphi=L\varphi+N(\varphi)$: it treats the
        linear operator $L$ \emph{exactly} (mode by mode in Fourier space) and adds
        the nonlinear part $N$ through the factor $(1-\ee^{-\omega\,\dd t})/\omega$.

  \item[Euler--Maruyama] The first-order numerical scheme for stochastic
        differential equations: $x \mathrel{+}= \dd t\cdot\text{drift} +
        \sqrt{2D\,\dd t}\,\mathcal{N}(0,1)$. The $\sqrt{\dd t}$ scaling of the noise
        term is what distinguishes it from deterministic Euler stepping.

  \item[\code{eval}] Python's built-in expression evaluator. The DAE solver
        \code{eval}s equation text inside a deliberately sandboxed namespace (only
        the field and parameter symbols are exposed).

  \item[External fields / legs] The $k$ field operators whose correlator is being
        computed. In code, \code{external\_fields} is a length-$k$ list of
        \code{(field\_name, slot)} tuples; \code{k} is just its length (the
        correlator order).

  \item[External-Wick compensation \code{comp}] The orbit--stabilizer divisor
        $|\Aut(\Gamma,\text{free})| / |\Aut(\Gamma,\text{fixed})|$ that removes the
        over-counting introduced when the integrator sums over permutations of
        same-field external leaves. (See orbit--stabilizer theorem.)

\end{description}

% =====================================================================
\section*{F}
% =====================================================================
\begin{description}[leftmargin=0pt,style=nextline]

  \item[Factorial moment / falling factorial] The product
        $n^{(m)}=n(n-1)\cdots(n-m+1)$. For a Poisson count $n$,
        $\langle n^{(m)}\rangle=\lambda^{m}$. Replacing $n^{m}$ by $n^{(m)}$ at
        coincident time bins strips the self-spike (shot-noise) $\delta$-terms from
        a point-process correlator, isolating the smooth off-diagonal cumulant the
        diagrams predict.

  \item[Fan triangulation] Splitting a convex polygon into triangles all sharing
        vertex $0$. It is how the temporal $m=2$ polytope integral is decomposed.

  \item[\code{fast\_callable}] Sage's compiler that turns a symbolic expression into
        a fast Python closure (Sage's analogue of \code{lambdify}). It is the
        slow-path integrand compiler --- ``slow'' only relative to the JIT-compiled
        inner loops, not to symbolic evaluation.

  \item[\code{FieldTheory} / \code{expand()}] The Phase-1 engine that
        Taylor-expands the symbolic action about the saddle and classifies every
        term by bigrade. \code{ensure\_taylor\_order} re-invokes \code{expand()}
        only when a cached expansion's order is too low for the requested diagrams.

  \item[\code{\_FieldScalar}] A size-1 wrapper class that supports both bare
        arithmetic and indexed access, so a user can write \code{xt*x} and
        \code{xt[0]*x[0]} interchangeably for a single-component field.

  \item[Fixed-point index] See \emph{DAE saddle solve}. The integer that picks which
        sorted mean-field root the diagrammatics expand around.

  \item[Fluctuation field] The internal field $\delta\varphi=\varphi-\varphi^{*}$
        --- the actual fluctuating quantity about the mean-field background, and the
        variable the action is Taylor-expanded in. Natural code names carry a
        \code{d}-prefix: \code{dn}, \code{dv}, \code{dx}.

  \item[Fluctuation--regression (Onsager regression) theorem] The statement that an
        equilibrium fluctuation relaxes by the same law as a macroscopic
        perturbation: $C(\tau)=\ee^{-A\tau}\Sigma$. It connects the stationary
        covariance $\Sigma$ to the time-dependent correlator.

  \item[Fork vs thread parallelism] Two ways to run work in parallel. \emph{Fork}
        clones the parent process so workers inherit its memory with no input
        pickling --- fast, but it crashes a macOS Jupyter kernel once Cocoa/BLAS
        have initialised. \emph{Threads} share one process's memory and are safe
        because heavy NumPy operations release the GIL. The temporal stages fork
        (behind a notebook-detecting guard); the spatial stages thread. \emph{Spawn}
        is a third method that pickles its arguments --- it would fail here on the
        unpicklable integrand closures.

  \item[Form factor] The polynomial in the momenta that a \emph{derivative}
        interaction vertex deposits on the loop integrand:
        $\nabla^{2}\to-|k|^{2}$, $\partial_x\to \ii k$,
        $\partial_t\to-\ii\omega$. Per-leg vs composite vertices give different
        form factors; in the coupled path it is written
        $\mathcal{R}(q,\ell)=\prod_{\text{vertices}}\mathfrak{F}(v)$.

  \item[Form-factor signature] \code{vertex\_form\_factor\_signature}: a SHA-256
        fingerprint of the operator-IR vertex table
        (\file{pipeline/\_expand\_cache.py:208}). It guards the expand cache against
        a slug that under-determines the per-vertex momentum form factors --- two
        models with the same name but different derivative vertices get different
        signatures.

  \item[Formal / undefined function] A Sage symbolic function such as
        \code{function('phi')(v)} left \emph{unapplied} --- with no concrete body
        --- so that \code{taylor()} can differentiate it abstractly. Its
        derivative-at-0 atoms are afterwards renamed to clean derivative symbols.
        (Also ``formal function''.)

  \item[Free action / bilinear kernel] The Gaussian $(1,1)$ sector of the action:
        exactly one response leg and one physical leg. \code{ft.free\_action()}
        returns this bigrade-$(1,1)$ polynomial; its matrix is the kernel $K$, whose
        inverse is the free propagator.

  \item[\code{fpairs} / \code{propagator\_indices}] Per-edge
        $(\text{response index},\text{physical index})$ matrix-index pairs that
        identify which field's propagator a half-edge uses. They are threaded into
        the spectral-assignment sum on the coupled multi-field path.

  \item[\code{fsolve}] SciPy's multivariate root-finder
        (\code{scipy.optimize.fsolve}), used to solve the mean-field saddle. Its
        integer convergence flag is \code{ier}; \code{ier == 1} signals clean
        convergence.

  \item[\code{fundamental} / \code{parameters}] The dict of numeric parameter
        values keyed by parameter name. \code{parameters} is the canonical name;
        \code{fundamental} is a deprecated alias that \code{Config.\_\_post\_init\_\_}
        mirrors onto the other.

\end{description}

% =====================================================================
\section*{G}
% =====================================================================
\begin{description}[leftmargin=0pt,style=nextline]

  \item[Gauss--Legendre / Gauss--Hermite / Gauss--Laguerre quadrature] Fixed-node
        numerical integration rules, each exact for polynomials up to a degree.
        \emph{Legendre} integrates over a finite interval (vertex times, the
        Schwinger $\sigma$); \emph{Hermite} integrates against a Gaussian weight
        $\ee^{-x^{2}}$ (the form-factor loop average) and an $n$-node rule is exact
        to polynomial degree $2n-1$; \emph{Laguerre} integrates
        $\int_0^{\infty}\ee^{-x}g(x)\,\dd x$ exactly for polynomial $g$ (the
        $[0,\infty)$ Schwinger parameters in the oracle path).

  \item[Generalized eigenvalue problem] The problem $C v=\lambda D v$ for matrices
        $C,D$; it reduces to the ordinary eigenproblem when $D=I$. Linear-stability
        analysis needs the general form because the ``mass'' matrix $A$ is singular
        --- algebraic constraints contribute zero rows.

  \item[\code{@dataclass}] A Python decorator that auto-generates a class's
        constructor (and a few dunder methods) from its declared fields. \code{Config}
        is a \code{@dataclass}, which is why you set only the fields you care about
        and inherit defaults for the rest.

  \item[Gillespie] The exact stochastic-simulation algorithm for jump processes (it
        schedules events one at a time). The Hawkes simulators here use a
        \emph{fixed-step Poisson-draw} approximation rather than true Gillespie
        event scheduling.

  \item[GIL (Global Interpreter Lock)] Python's lock that serialises bytecode
        execution across threads. Heavy NumPy operations \emph{release} the GIL
        while they run in C, which is precisely why thread (not fork) parallelism
        gives a real speed-up on the spatial path.

  \item[\code{G\_ft = K\_ft$^{-1}$}] The propagator matrix in frequency: the inverse
        of the quadratic-form kernel $K_{\mathrm{ft}}$. Its rows index physical
        fields and its columns index response fields; its poles are relaxation
        modes and its residues give the time-domain exponentials.

  \item[\code{graphs.trees(n)}] A Sage generator that yields every non-isomorphic
        tree on $n$ vertices exactly once. Diagram topologies are built by starting
        from these loop-free skeletons and adding edges.

  \item[GTaS] A noise/model family --- ``Gaussian-and-shifts'', a Bernoulli +
        Gaussian external-rate process modelled by \code{GTaSNoise}. It is the
        correlated-input model the cumulant machinery was originally built for; at
        order $N=2$ all its higher cumulants are local.

\end{description}

% =====================================================================
\section*{H}
% =====================================================================
\begin{description}[leftmargin=0pt,style=nextline]

  \item[Hankel transform / $J_0$] The radial Fourier transform in two dimensions,
        which uses the Bessel function $J_0$. (Spatial $d=2$ correlators.)

  \item[Hawkes process] A self-exciting point process: each spike raises the rate of
        future spikes. The simulators realise it by per-step
        $\mathrm{Poisson}(\lambda\,\dd t)$ draws fed back into the voltage.

  \item[Heat kernel] The Green's function of the diffusion operator: in $d$
        dimensions $(4\pi B t)^{-d/2}\ee^{-|x|^{2}/4Bt}$. It is the
        momentum-to-position ($q\to x$) Fourier transform of a propagator mode, and
        it is what makes the analytic inverse Fourier transform possible without a
        momentum grid.

  \item[Heat-kernel / $(k,t)$ representation] Writing propagators in
        momentum--time space, $G_R(k,t)=\theta(t)\,\ee^{-(\mu+Dk^{2})t}$. Its key
        property: each edge time acts as a Schwinger parameter, so the loop momentum
        integral becomes Gaussian.

  \item[Heaviside $\Theta(t)$] The unit step function that enforces retardation
        (a response only after its cause). The framework adopts the convention
        $\Theta(0)=0$, so the boundary $\Delta t=0$ is strictly excluded from
        retarded support.

  \item[Hermite--Genocchi formula] An identity expressing a divided difference of a
        function $f$ as a simplex integral of $f^{(n)}$. It is the bridge that turns
        the Duhamel (time-ordered) integral of the Dyson dressing into a divided
        difference. (See also Opitz theorem.)

  \item[Horner's rule] The efficient nested form for evaluating a polynomial,
        $((c_n x + c_{n-1})x + \cdots)$, using the minimum number of
        multiplications.

\end{description}

% =====================================================================
\section*{I}
% =====================================================================
\begin{description}[leftmargin=0pt,style=nextline]

  \item[Idempotent] A property of an operation: running it twice has the same effect
        as running it once. Several builder methods are idempotent so that
        re-authoring a theory does not double-apply a change.

  \item[IFT (inverse Fourier transform)] Here, the momentum-to-position transform
        $q\to x$ (or $k\to x$). ``Analytic IFT'' means it is done in closed form
        --- because each Schwinger sample's $q$-dependence is Gaussian, the
        transform is a heat kernel --- which retires the discretised $q$-grid.

  \item[Image sum / method of images] Constructing a periodic-boundary Green's
        function as a sum of displaced copies of the infinite-domain kernel.

  \item[\code{importlib.util} dynamic import] The standard-library mechanism for
        importing a Python file from an \emph{arbitrary path} rather than by a name
        on \code{sys.path}. The canonical three-step idiom ---
        \code{spec\_from\_file\_location}, \code{module\_from\_spec},
        \code{exec\_module} --- is how \code{load\_theory} imports a
        \code{.theory.py} file.

  \item[Incidence digraph (coloured)] See \emph{coloured incidence digraph}.

  \item[Interaction vertex] Any action monomial of total degree $\ge 3$ (so it has
        $\ge 1$ physical leg). It becomes a diagram vertex with $n_{\tilde{}}$
        response legs and $n_{\mathrm{phys}}$ physical legs; it is the building
        block of every loop. In code, a \code{VertexType}. Contrast a noise
        \code{SourceType} ($n_{\mathrm{phys}}=0$).

  \item[Internal vertex] A non-leaf vertex of a diagram (degree $\ge 2$); an
        interaction point. Its time is an integration variable.

  \item[\code{ipywidgets}] Jupyter's interactive-widget library, imported in the UI
        as \code{W}. It keeps Python objects on the kernel side synced to live HTML
        controls in the notebook --- sliders, text boxes, buttons.

  \item[IPython.display] IPython's rich-output API (\code{display}, \code{HTML},
        \code{Javascript}); the UI uses it to render its controls and status panel.

  \item[Isserlis / Wick theorem] The statement that the moment of a product of
        jointly Gaussian variables equals the sum over perfect matchings of products
        of pairwise covariances. It is the closed form for the loop-momentum average
        in the coupled spatial integrator.

  \item[Isomorphism (of graphs)] A relabelling of one graph's vertices that maps it
        onto another edge-for-edge. Isomorphic diagrams are ``the same diagram'';
        coloured isomorphism additionally requires colours (leaf vs internal, field
        types) to be preserved.

  \item[Iteration vs compound saddle] Two kinds of mean-field saddle. An
        \emph{iteration} (or \emph{closure}) saddle has a right-hand side that is a
        single function call (\code{nstar = phi(vstar)}) and is substituted via the
        formal rename; a \emph{compound} saddle's right-hand side is a
        parameter/saddle expression and is substituted concretely. (In legacy
        heterogeneous code, ``iteration saddle vs compound saddle'' also names the
        variables \code{fsolve} actually iterates --- rates $n^{*}$ --- versus those
        evaluated from them each step --- voltages $v^{*}$.)

  \item[\code{is\_trivial\_zero()}] Sage's \emph{syntactic} zero test: \code{True}
        only when an expression is literally zero. It is robust to the ambiguity of
        unbound parameters, where a full simplification might wrongly claim (or fail
        to claim) zero.

\end{description}

% =====================================================================
\section*{J}
% =====================================================================
\begin{description}[leftmargin=0pt,style=nextline]

  \item[$j$ (leaf surplus)] The number of ``extra'' tree leaves that the $\ell$
        added loop-edges will consume. Trees are generated with $k+j$ leaves so that
        adding $\ell$ edges (each pairing two leaves into an internal line) leaves
        exactly the $k$ external legs.

\end{description}

% =====================================================================
\section*{K}
% =====================================================================
\begin{description}[leftmargin=0pt,style=nextline]

  \item[$k$ (correlator order) / $k$-point function] The number of external legs ---
        equivalently the order of the cumulant. A $k$-point function has diagrams
        with $k$ external legs; $k=2$ is the covariance / power spectrum.

  \item[$K_0$] The modified Bessel function of the second kind of order zero; it
        appears in the static two-dimensional spatial correlator.

  \item[Keldysh / correlation \code{C} edge vs retarded \code{R} edge] A naming of
        the two edge kinds. A \code{C} (correlation) edge carries a noise-driven
        correlation, with its Schwinger parameter integrated over $[|t|,\infty)$; an
        \code{R} (retarded) edge is a fixed-duration causal propagator. (See \code{C}
        edge, \code{R} edge.)

  \item[Kernel ($g$)] A convolution (memory) kernel, e.g. a synaptic filter,
        coupling $K\star y=\int K(t-t')\,y(t')\,\dd t'$. It is declared either by its
        time form or --- preferred --- by its Fourier image $\tilde K(\omega)$.

  \item[Kernel matrix $K$] The matrix of bilinear $(1,1)$ coefficients of the free
        action, with layout $[\text{response row},\,\text{physical column}]$. Its
        inverse is the propagator $G=K^{-1}$.

  \item[Kernel normalization] The convention that synaptic / temporal kernels
        integrate to $1$. As a result, at the stationary saddle every kernel symbol
        is simply replaced by $1$.

  \item[\code{kernel\_ft\_image} hook] A model-supplied callable
        $(\mathrm{ns},\omega)\mapsto\text{subs-dict}$ that maps each abstract kernel
        symbol $g$ to its frequency image $\hat g(\omega)$. It is how a model tells
        the propagator builder what its kernels look like in Fourier space.

  \item[\code{K\_ker} (kernel form)] The kernel $K$ written with abstract
        $\delta$/$\delta'$ symbols ($c_0\,\delta + c_1\,\delta'$) so that it
        Fourier-transforms cleanly into \code{K\_ft}.

\end{description}

% =====================================================================
\section*{L}
% =====================================================================
\begin{description}[leftmargin=0pt,style=nextline]

  \item[Lag-dependent normalization] In a correlation estimator, dividing each lag's
        sum by the number of \emph{overlapping} valid bin pairs ($n_{\mathrm{valid}}
        -|\text{lag}|$), so the finite-window estimate stays unbiased.

  \item[\code{lambdify}] SymPy's utility that compiles a symbolic expression into a
        fast NumPy-callable function. (Sage's analogue is \code{fast\_callable}.)

  \item[\code{Laplacian} / \code{Lap}] The formal spatial second-derivative
        operator $\nabla^{2}$. In plain authoring it is used multiplicatively
        (\code{D*Laplacian*phi} $=D\nabla^{2}\varphi$) and is sent to $0$ on a
        spatially uniform background; in operator-IR authoring it becomes the call
        node \code{Lap()}.

  \item[Leaf / external leg] A degree-1 vertex of a diagram --- an external argument
        of the correlation function. For $k\ge 3$, one leaf (the \emph{origin leaf})
        is pinned to $t=0$.

  \item[Leg] One field ``stub'' at a vertex or source, identified by
        $(\text{field\_base\_name},\,\text{population\_index})$ and repeated once per
        exponent. Response legs feed outgoing edges; physical legs feed incoming
        edges.

  \item[Linear extension] A topological sort of a partial order --- one total
        ordering consistent with all the constraints. Each linear extension of the
        causal poset is one chain simplex / one causal chamber; the polytope is the
        disjoint union over all of them.

  \item[Linear stability] The classification of a fixed point: it is linearly stable
        when every (finite) eigenvalue of the linearised dynamics has negative real
        part, so small perturbations decay. Opt in with
        \code{.stability\_analysis(True)}; the test is the generalized eigenvalue
        problem $(B-\ii\omega A)\,\delta x=0$.

  \item[Loop number / loop order $\ell$ (\code{max\_ell})] The number of independent
        loops in a diagram, $|E|-|V|+1$ --- equivalently the order in the
        perturbative expansion. $\ell=0$ is tree level; each additional loop is one
        more power of the small parameter. \code{max\_ell} is the run's cutoff.

  \item[Lorentzian (kernel)] The exponential $\ee^{-|\tau|/\tau_c}$, named for its
        Lorentzian Fourier transform $\propto 1/(\omega^{2}+1/\tau_c^{2})$. It is
        the stationary covariance of an Ornstein--Uhlenbeck colored-noise process.

  \item[Lyapunov equation] The matrix equation $A\Sigma+\Sigma A^{\mathsf T}=Q$ (or
        $A\Sigma+\Sigma A^{\mathsf H}=Q$); its solution $\Sigma$ is the stationary
        covariance of the linear SDE $\dd x=-A x\,\dd t+\text{noise}(Q)$. Solved by
        \code{scipy.linalg.solve\_continuous\_lyapunov} (Bartels--Stewart) in the
        coupled simulators and oracle.

\end{description}

% =====================================================================
\section*{M}
% =====================================================================
\begin{description}[leftmargin=0pt,style=nextline]

  \item[\code{m} (after delta)] The number of integration variables that survive
        $\delta$-elimination within a subset. Its value selects which closed-form
        integrator runs: $m=0$, $1$, $2$, or the general $m\ge 3$
        (chain-simplex) path.

  \item[\code{manifest.json}] The human-readable index of cached entries (their
        \code{stage}, \code{k}, \code{loop\_order}, \code{saved\_at}). It is
        best-effort and \emph{not} authoritative --- the \code{.sobj} files are the
        truth.

  \item[Markovian embedding / Markovianization] Enlarging the state space with
        auxiliary fields so that a non-Markovian (memory-carrying, colored) process
        becomes Markovian (white-driven, local in time). The observable output field
        is unchanged; only the representation grows. Toggled with
        \code{.markovianize(...)}.

  \item[Mass $A$ (a.k.a. $\mu$)] The $k^{0}$ term of a spatial inverse propagator
        --- the relaxation rate. Also the reaction (mass) matrix $M$ for a
        multi-field theory: the momentum-independent part of the linearised inverse
        propagator; a diagonal $M$ means decoupled fields.

  \item[Matched-cutoff principle] When comparing a simulation against theory,
        comparing them at the \emph{same} ultraviolet cutoff $k_{\max}=\pi/\dd x$,
        because a finite-difference lattice's dispersion differs from the continuum
        $q^{2}$ at order $\dd x$.

  \item[Mean field / saddle point ($\varphi^{*}$)] The deterministic steady state
        about which the loop expansion is taken --- the constant background that
        makes the action's field-linear (tadpole) term vanish. Saddle parameters end
        in \code{star} (e.g. \code{xstar}, \code{nstar}); they are solved by the
        engine, never supplied. The saddle (MF) sector is the bigrades
        $(0,0),(1,0),(0,1)$, which the framework verifies and zeroes.

  \item[Merged-residue tensor $B_\alpha$] The grouped path's per-pole-tuple
        coefficient, $\sum_j cp_j\prod\delta\prod C^{(j)}$ --- a single tensor that
        replaces the per-diagram Cartesian product $\prod_e C_{\alpha_e}$. Summing
        residues once, across diagrams sharing a topology, is what makes the grouped
        Phase~J fast and exact.

  \item[Mode $(C_\alpha,\lambda_\alpha)$ / mode sum] One pole's residue $C_\alpha$
        and rate $\lambda_\alpha=\ii p_\alpha$, so that the smooth part of a
        propagator is $\sum_\alpha C_\alpha\ee^{\lambda_\alpha\Delta t}$. This
        per-edge representation is the \emph{mode sum}, stored as an
        \code{EdgeModeSum}.

  \item[Moment] A raw correlation $\langle\varphi\cdots\varphi\rangle$, or the same
        for the centred field (a central moment). Reconstructed from cumulants by the
        set-partition (cluster-expansion) formula.

  \item[\code{mp.dps}] mpmath's decimal-digit precision setting (raised to $50$ in
        the precision-rescue paths). (See mpmath.)

  \item[mpmath] A Python library for arbitrary-precision floating point and special
        functions. It is used where machine \code{double} would suffer catastrophic
        cancellation --- e.g. the \code{erf}-cancellation in the heat-kernel closed
        form.

  \item[\msrjd{} action / field theory] The
        Martin--Siggia--Rose--Janssen--De~Dominicis path-integral representation of a
        stochastic dynamical system. (See the boxed definition below.)

\end{description}

\begin{defn}
The \msrjd{} construction rewrites a Langevin (stochastic differential)
equation as a path integral with weight $\ee^{-S[\varphi,\tilde\varphi]}$. Every
\textbf{physical field} $\varphi$ is paired with a conjugate \textbf{response
field} $\tilde\varphi$ (the ``tilde'' or ``hatted'' field; written with a
\code{t}- or \code{d}-prefix in code, e.g.\ \code{xt}, \code{dh}). The
deterministic drift sits in the $\tilde\varphi\cdot(\dots)$ terms; the noise sits
in the response-field-only terms (white Gaussian noise is the single vertex
$-D\,\tilde\varphi^{2}$, the bigrade-$(2,0)$ sector). Sourcing the response field
produces response functions; sourcing the physical field produces correlation
functions. Propagators are directed (causal): the retarded $\Gret$ connects a
physical leg to a response leg and is non-zero only for $t>0$.
\end{defn}

\begin{description}[leftmargin=0pt,style=nextline]
  \item[Multi-edge (multigraph)] Two or more edges between the same pair of
        vertices; a double edge forms a ``bubble'' loop. Enabled in Sage with
        \code{multiedges=True}.

  \item[Multi-start Newton] Running a Newton-type root finder from many random
        initial guesses to discover \emph{all} the roots of a nonlinear system (e.g.
        both wells of a double-well saddle), then sorting them so
        \code{fixed\_point\_index} can pick one deterministically.

  \item[\code{multiset\_permutations} (SymPy)] A SymPy generator yielding the
        $N!/\prod_r n_r!$ distinct orderings of a multiset (a list with repeats). It
        enumerates the distinct external-leaf assignments without double-counting
        identical legs.
\end{description}

% =====================================================================
\section*{N}
% =====================================================================
\begin{description}[leftmargin=0pt,style=nextline]

  \item[\code{nauty}] Brendan McKay's C library for graph canonical labelling and
        isomorphism testing --- the name stands for ``No AUTomorphisms, Yes?''. It is
        bundled inside SageMath and is invoked implicitly by every
        \code{is\_isomorphic}, \code{canonical\_label}, and \code{automorphism\_group}
        call. It is what makes diagram deduplication and symmetry-factor computation
        tractable.

  \item[\code{networkx}] A pure-Python graph library used inside the engine for some
        graph manipulations (distinct from Sage's \code{Graph}/\code{DiGraph}).

  \item[Newton refinement] Polishing an approximate root of $\det(K(\omega))$ to
        machine precision with Newton's method, after a coarser method has located
        it.

  \item[Noise kernel] The response-field-only sectors of the action --- bigrades
        $(n_{\tilde{}}\ge 2,\,0)$. These are the noise sources from which the diagram
        expansion's stochastic forcing comes.

  \item[NoiseSourceType / non-local cumulant kernel] A noise vertex whose kernel
        $\kappa(\tau)$ introduces \emph{extra} $\tau$ integration variables and a
        non-rational integrand, which disables the analytic (closed-form) integration
        path and forces \code{scipy.nquad}. Contrast a plain (local, white) noise
        source.

  \item[\code{nquad}] SciPy's nested numerical multidimensional quadrature
        (\code{scipy.integrate.nquad}); the slow adaptive fallback (about
        $10^{-8}$ precision) that the analytic machinery exists to avoid.

  \item[\code{ns} (NamespaceForFields)] The runtime object handed to the model's
        lambdas, carrying the symbolic field and parameter variables (Sage \code{SR}
        objects) as attributes. It is built by the consumer, so a model author can
        write \code{ns.mu}, \code{ns.x}, and so on.

  \item[\code{num\_params}] The substitution dictionary
        $\{\code{SR}\ \text{symbol}\to\text{float}\}$ that converts every symbolic
        propagator, vertex, and residue into a numeric one. It is the deliverable of
        the mean-field / numericisation stage. Its keys are Sage \code{SR} symbols,
        not strings.

  \item[\code{numba} / \code{@njit} / nopython mode] \code{numba} is a
        just-in-time compiler that turns annotated Python functions into native
        machine code. Its \code{@njit} decorator forces ``nopython'' mode --- no
        interpreter fallback --- giving roughly $100\times$ speed-ups on tight
        numeric loops (e.g. the chain-simplex inner loop, the Euler--Maruyama
        stepper). It cannot type a Sage ring element, which is why simulator inputs
        must be pre-cast to plain Python floats.

  \item[NumPy \code{ndarray}] The $n$-dimensional array type underlying all numeric
        data in the pipeline --- grids, correlators, residue matrices.

\end{description}

% =====================================================================
\section*{O}
% =====================================================================
\begin{description}[leftmargin=0pt,style=nextline]

  \item[\code{.operator()} / \code{.operands()}] Sage methods that return,
        respectively, an expression's outermost operation and the list of its
        arguments. Walking an expression with these is how the framework inspects a
        symbolic tree without string parsing.

  \item[Operator IR] The intermediate representation for derivative interaction
        vertices: $\nabla^{2}$ (\code{Lap}), $\partial_t$ (\code{Dt}), and
        $\partial_{x_i}$ (\code{Dx}) become inert symbolic \emph{call} nodes with
        explicit algebra passes (\file{pipeline/spatial\_operator\_ir.py}). It is
        what lets spatial theories with derivative vertices (Model~B, KPZ, Burgers)
        carry the right momentum form factors. Authoring is gated behind
        \code{.operator\_ir()}.

  \item[Opitz theorem] The fact that a divided difference equals the top-right entry
        of $f$ applied to the upper-bidiagonal ``Opitz matrix'' (the nodes on the
        diagonal, ones on the superdiagonal). It is \emph{confluent-safe} --- no
        $0/0$ at repeated nodes --- which is why the Dyson dressing uses it for the
        $\Phi_n$ primitive.

  \item[ORACLE-ONLY] A label on a module that exists solely as an \emph{independent}
        numerical cross-check, kept off the production path. The codebase
        banner-annotates such modules.

  \item[Orbit--stabilizer theorem] The group-theory identity
        $|\text{orbit}|=|\text{group}|/|\text{stabilizer}|$. It justifies the
        external-Wick compensation: the number of distinct external-leaf labellings
        is the automorphism group's size divided by the stabilizer that fixes the
        leaves.

  \item[\code{origin\_leaf\_idx} / origin pinning] The index of the external leaf
        pinned to $t=0$. By time-translation invariance a $k$-point cumulant depends
        only on time differences, so one leaf is fixed at the origin and the
        resulting callable computes a function of the remaining differences.

  \item[Ornstein--Uhlenbeck (OU) process] The simplest mean-reverting linear SDE,
        $\dd\xi/\dd t=-\xi/\tau_c+\text{white noise}$. It is Gaussian, with
        stationary variance $D/\mu$ and a single-Lorentzian (exponential)
        autocorrelation. It appears both as a \emph{theory} (the linear baseline)
        and as the auxiliary \emph{colored-noise generator} in Markovian embedding.

\end{description}

% =====================================================================
\section*{P}
% =====================================================================
\begin{description}[leftmargin=0pt,style=nextline]

  \item[Partition (\code{nauty}/Sage sense)] A vertex colouring passed to the graph
        algorithms; vertices may only be mapped within the same colour class. (Not to
        be confused with a set partition of indices.)

  \item[Per-leg vs composite] See \emph{composite vs per-leg vertex}.

  \item[Phase J] The time-domain loop-integration stage of the temporal pipeline
        (\code{compute\_correction\_td}). (See the boxed definition below.)

  \item[$\Phi_n(t;\nu)$] The confluent-safe nested-simplex time primitive of the
        Dyson dressing, with $\Phi_0=1$ and $\Phi_1=(1-\ee^{-t\nu})/\nu$. It is
        evaluated via the Opitz matrix to stay stable at repeated nodes.

  \item[$\mathcal{H}_n(t)$] The matrix factor of the $n$-th Dyson order: a sum over
        projector strings times $\ee^{-m_{\alpha_0}t}\,\Phi_n$.

  \item[$\varphi^{(k)}$ / \code{phik\_i}] The $k$-th derivative of the transfer
        function at the saddle of population $i$; \code{phi0\_i}
        $=\varphi(v^{*}_i)=n^{*}_i$. These derivatives become the diagram couplings.

  \item[\code{\_PhiCallableList}] A list of per-population transfer-function
        callables that also supports a bare scalar call \code{f(v)} when there is
        exactly one population.

  \item[Physical (fluctuation) field] The actual fluctuating quantity
        $\delta\varphi=\varphi-\varphi^{*}$, the ``in-leg'' of a diagram. Natural code
        names: \code{n}, \code{v}, \code{x}, \code{h}. (See also fluctuation field.)

  \item[Picklable] Serializable by Python's \code{pickle}, so that
        \code{multiprocessing} can ship an object between processes. This is why the
        CGF kernel callables (\code{\_CGFKernelCallable}, \code{\_CGFKernelSum}) are
        defined as module-level classes rather than closures.

  \item[PipelineCache] The low-level cache store keyed by \code{(stage, k, ell)},
        writing \code{.sobj} files with a JSON manifest
        (\file{msrjd/core/cache.py}).

  \item[Pole $\omega_k$ / pole-residue form] A retarded zero of
        $\det(K_{\mathrm{ft}})$ with $\operatorname{Im}(\omega_k)>0$, setting a
        relaxation/oscillation rate. The propagator's \emph{pole-residue form} is
        $\sum_k C_k\,\ee^{-\ii\omega_k t}+D_\delta\,\delta(t)$ --- the form Phase~J
        integrates against. In code, \code{pole\_vals} holds the poles and
        \code{C\_mats} the matrix residues.

  \item[Pole tuple $\alpha$] A choice of one pole per smooth edge. Expanding the
        product of per-edge mode sums yields a sum over all pole tuples; each tuple's
        contribution is a single exponential whose polytope integral is closed-form.

  \item[Polytope / chamber] The convex region carved out by the causal
        $\Delta t>0$ constraints --- the integration domain. (See causal poset, chain
        simplex, causal chamber.)

  \item[\code{poly.dict()}] A Sage polynomial-ring method returning
        $\{\text{exponent tuple}\to\text{coefficient}\}$ --- the pickle-stable
        ``dict form'' the expand cache stores in place of the (less stable)
        polynomial objects themselves.

  \item[\code{PolynomialRing(SR, names)} / fraction field] A polynomial ring over the
        Symbolic Ring whose generators are the named field symbols; it exposes
        monomials via \code{.dict()}. Its fraction field $\mathrm{Frac}(R[\omega])$
        reduces each propagator entry to a canonical irreducible $P/Q$.

  \item[Population] A named index set --- a group of \code{size N} identical units.
        Attaching a field to a population makes it a length-$N$ vector and introduces
        \code{for i in pop} sums in the action. Heterogeneous theories carry several
        populations.

  \item[\code{pop\_idx} (population index)] The 1-based subscript on a field
        generator (e.g. the \code{2} in \code{nt2}), distinguishing
        populations/components.

  \item[Precompute] The cheap one-time structural pass --- \code{expand} at order 2,
        the saddle sanity check, the mean-field solve, and the propagator build ---
        that populates the durable caches. Exposed as \code{pipeline.precompute(model)}
        and run by the UI's \emph{Pre-compute} button.

  \item[Prediagram] A bare, oriented graph \emph{topology} --- leaves, internal
        vertices, and arrows on edges --- carried as the tuple
        \code{(D, G, leaves, internal)}, \emph{before} any propagator types or vertex
        monomials are assigned. It is the output of the enumeration subsystem.

  \item[Preparse / SageMath preparser] Sage's source rewrite of notebook-cell text:
        integer literals become \code{Integer}, floats become \code{RealNumber}, and
        \code{\^{}} becomes \code{**}. Because this rewrite changes plain Python
        semantics, the simulators must live in (un-preparsed) \code{.py} files and
        accept pre-cast scalars.

  \item[Proper / strictly proper rational] A rational function with
        $\deg(\text{numerator})<\deg(\text{denominator})$; it vanishes at infinity
        (no $\delta$ term). A non-proper rational has a polynomial part, which becomes
        a $\delta(t)$ in the time domain.

  \item[Propagator $G=K^{-1}$] The two-point Gaussian correlation/response function
        --- the building block of every Feynman edge --- obtained as the matrix
        inverse of the bilinear kernel $K$. (See \code{G\_ft}, retarded propagator.)

\end{description}

\begin{defn}
\textbf{Phase J} is the temporal pipeline's loop-integration stage,
\code{compute\_correction\_td}. It performs the actual vertex-time quadrature for
each diagram, in the time domain, using the propagator's pole-residue form
$\sum_\alpha C_\alpha\ee^{\ii p_\alpha t}+D_\delta\,\delta(t)$. The product of
retarded propagators over the edges is expanded into $2^{|E|}$ choices of
``instantaneous $\delta$ vs smooth'' per edge; each $\delta$-edge eliminates one
integration variable, leaving $m$ variables that select a closed-form path
($m=0,1,2$) or the general chain-simplex path ($m\ge 3$). The \emph{grouped}
variant (Stage~4a) sums residues once across diagrams sharing a topology, via the
merged-residue tensor $B_\alpha$, dropping the precision floor to machine
$\varepsilon$.
\end{defn}

% =====================================================================
\section*{Q}
% =====================================================================
\begin{description}[leftmargin=0pt,style=nextline]

  \item[\code{q}-grid] A discretised momentum axis used by a \emph{numerical} Fourier
        transform. It is ``retired'' by the analytic heat-kernel inverse Fourier
        transform for the plain-vertex and $d=1$ cases, which removes the associated
        cutoff and ringing artefacts.

  \item[\code{QQ}] Sage's notation for the field of rational numbers $\mathbb{Q}$;
        used (alongside \code{CyclotomicField(4)}) for exact symbolic arithmetic.

  \item[\code{quad}] SciPy's adaptive one-dimensional quadrature
        (\code{scipy.integrate.quad}); the 1-D counterpart of \code{nquad}.

  \item[$Q_{\mathrm{eff}}$] The second Symanzik polynomial in its Schur-complement
        form, $Q-N^{\mathsf T}\Lambda^{-1}N$ --- the residual quadratic form carrying
        external-momentum dependence after the loop block is integrated out. (See
        Symanzik polynomials, Schur complement.)

\end{description}

% =====================================================================
\section*{R}
% =====================================================================
\begin{description}[leftmargin=0pt,style=nextline]

  \item[\code{R} edge / retarded line] A $\Gret$ propagator segment, carrying a
        causal $\theta(w\ge 0)$. (Contrast the \code{C} correlation edge.)

  \item[Reaction (mass) matrix $M$] See \emph{mass}.

  \item[Reference diffusion $D_0$] The scalar part split off the diffusion matrix,
        $\mathcal{D}=D_0 I+\hat{\mathcal{D}}$ (default $D_0=\operatorname{trace}/N$),
        chosen to make the residual $\hat{\mathcal{D}}$ small so the Dyson series
        converges quickly.

  \item[Regime 1 (finite cutoff)] Evaluating the spatial correlator \emph{at} a
        finite momentum cutoff $k_{\max}$; the continuum limit is
        $k_{\max}\to\infty$.

  \item[Residual diffusion $\hat{\mathcal{D}}$] $\mathcal{D}-D_0 I$, the Dyson
        perturbation. If it is zero, no series is needed and the tree propagator is
        exact.

  \item[Residue matrix \code{C\_mats[k]}] $\ii\cdot\operatorname{Res}_{\omega_k}
        G_{\mathrm{ft}}$ --- the amplitude matrix of mode $k$. Together with the poles
        \code{pole\_vals} it gives the time-domain propagator
        $\sum_k C_k\,\ee^{\ii\omega_k t}$.

  \item[Response field (tilde / hat field)] The conjugate field $\tilde\varphi$ of
        the \msrjd{} construction --- the ``out-leg'' in diagram language. Its code
        base names carry a \code{t}- or \code{d}-prefix (\code{xt}, \code{nt},
        \code{dv}, \code{dn}). A degree-2 monomial in it is white Gaussian noise; a
        degree-$\ge 3$ monomial is non-Gaussian noise.

  \item[Retarded propagator $\Gret$] The causal Green's function:
        $\Gret(t)=D_\delta\,\delta(t)+\Theta(t)\sum_k C_k\,\ee^{\ii\omega_k t}$. It is
        the linear response of a physical field to a response-field source ---
        non-zero only for $t>0$ (effect after cause). In momentum--time space,
        $\Gret(k,t)=\theta(t)\,\ee^{-(\mu+Dk^{2})t}$; for a multi-field theory,
        $\Theta(t)\,\ee^{-(M+\mathcal{D}k^{2})t}$.

  \item[Routing coefficients $(a_e,b_e)$] The integer ($\pm 1$ or $0$) coefficients
        in $k_e=\sum a_{ei}\ell_i+\sum b_{ej}q_j$ --- how the loop momenta $\ell$ and
        external momenta $q$ thread each edge. They are found by solving momentum
        conservation (e.g. with SymPy \code{linsolve}).

\end{description}

% =====================================================================
\section*{S}
% =====================================================================
\begin{description}[leftmargin=0pt,style=nextline]

  \item[Saddle / mean-field parameter] See \emph{mean field}. A \code{<field>star}
        parameter holding a field's steady-state value, auto-created with each
        \code{physical\_field} and solved (never passed in).

  \item[SageMath / \code{SR} (Symbolic Ring)] SageMath is an open-source mathematics
        system that bundles many libraries (including \code{nauty}, Singular, Maxima,
        Pynac/GiNaC) behind a single Python API. \code{SR} is its Symbolic Ring ---
        the universal type of symbolic expressions; \code{SR.var('x')} makes a
        symbolic variable, and \code{diff}, \code{.subs}, \code{.simplify\_full}, and
        \code{float} operate on \code{SR} elements. The physics-facing algebra of the
        pipeline lives in \code{SR}, and a Sage kernel is required to run it.

  \item[\code{sage\_eval}] Sage's namespaced string-to-symbolic-expression evaluator
        --- it parses a string into an \code{SR} expression within a supplied
        namespace of allowed symbols.

  \item[\code{sage\_solve}] Sage's symbolic equation solver (backed by Maxima); used,
        for instance, to solve the linear $\Delta t=0$ equations of
        $\delta$-elimination.

  \item[Schur complement] For a block matrix, the residual form
        $Q-N^{\mathsf T}\Lambda^{-1}N$ left after integrating out the $\Lambda$
        (loop-momentum) block --- which is exactly the second Symanzik polynomial
        $Q_{\mathrm{eff}}$.

  \item[Schwinger parameter] An auxiliary integration variable (here an edge time
        $w$, or a noise-source time $\sigma\in[0,\infty)$) introduced to turn a
        propagator denominator into an exponential. Doing this for every edge makes
        the loop \emph{momentum} integral Gaussian, which then collapses to the
        Symanzik polynomials. In the heat-kernel representation the Schwinger
        parameter \emph{is} the edge's time duration.

  \item[Self-consistency equation] The fixed-point condition the saddle must satisfy
        (e.g. $n^{*}=\varphi(E+w\cdot n^{*})$); solving it is the mean-field solve.

  \item[Self-energy $\Sigma$] The one-particle-irreducible (1PI) two-point
        correction; it dresses the propagator through the Dyson equation. In the
        spatial diagnostics, \code{Sigma} is the one-loop self-energy used for the
        tadpole mass-shift.

  \item[Self-loop] An edge from a vertex to itself; disabled in the enumeration
        (\code{loops=False}) because it is never physical here.

  \item[Set partition / cluster expansion] The combinatorial identity expressing
        moments as sums over partitions of products of cumulants over the blocks. It
        is how \code{output='moment'} reconstructs full moments from the engine's
        native cumulants.

  \item[Shot noise / contact / diagonal term] The $1/\dd t_{\mathrm{bin}}$-scaling
        $\delta$-function piece of a point-process correlator --- a spike correlated
        with itself. It is removed by the factorial-moment correction. The temporal
        analogue is a \emph{shot-noise $\delta$-spike}: a residual same-population
        equality $\tau=0$ left after $\delta$-elimination.

  \item[\code{.simplify\_full()}] Sage's most aggressive symbolic simplifier. It is
        powerful but slow, and it can stall on unbound parameters --- hence the
        \code{cysignals} time-budget around it and the cheaper
        \code{is\_trivial\_zero()} where a full simplify is overkill.

  \item[Slice / full grid] For $k\ge 3$, a \emph{slice} is an axis-parallel 1-D cut
        (sweep one leg's lag, hold the rest) while the \emph{full grid} is the entire
        $(k-1)$-dimensional tensor of $(x_j,\tau_j)$ events.

  \item[Slug / slugify] A filesystem-safe lowercase identifier derived from a theory
        name, \code{re.sub(r'[\^{}A-Za-z0-9]+','\_',name).strip('\_').lower()}
        (\file{pipeline/theory\_serialize.py:96}). It is the per-theory cache
        directory name and the \code{<slug>.theory.py} filename.

  \item[\code{.sobj}] ``Sage Object'': a gzip-compressed pickle written by
        \code{sage.all.save} and read by \code{sage.all.load}. It is the on-disk
        format for every cached artefact --- the expanded action, the propagator, the
        enumerated diagrams.

  \item[Source / source vertex] An internal vertex with in-degree $0$ (and not a
        leaf) --- a pure response-field monomial from the noise kernel that
        \emph{injects} noise (a \code{SourceType}, with $n_{\mathrm{phys}}=0$).

  \item[\code{spatial\_grid} / \code{spatial\_points}] How a spatial run specifies its
        evaluation points. For $k=2$, a grid of separations $x$; for $k\ge 3$, an
        explicit \code{(n\_pts, k-1, 2)} array of $(x_j,\tau_j)$ events.

  \item[\code{spatial\_info} / modes / \code{Sigma}] The spatial driver's diagnostic
        dict. \code{modes} are $(A,B,N)$ triples parameterising each propagator
        mode's mass $A$, dispersion $B$, and noise weight $N$; \code{Sigma} is the
        one-loop self-energy.

  \item[Spectral assignment / spectral projector $P_\alpha$] On the coupled
        (unequal-$D$) path, a labelling of every propagator segment with one of the
        $N_{\mathrm{modes}}$ spectral eigenmodes of the mass matrix $M=\sum_\alpha
        m_\alpha P_\alpha$; the diagram sums over all such labellings. $P_\alpha$ is
        the (rank-1 or higher) matrix projecting onto the $\alpha$-th eigenspace,
        built from the eigenvectors.

  \item[Spec dict] The flat Python dict the UI produces from its form (via
        \code{\_collect}) and the serializer renders to / parses from a
        \code{.theory.py}.

  \item[Specialization] A model-supplied substitution that replaces an abstract
        derivative symbol with a concrete value, e.g. \code{phi2\_i} $\to 2a$.

  \item[Structure factor $S(q)$] The equal-time momentum-space two-point function
        $\langle|\varphi_q|^{2}\rangle$ of a spatial field; at tree level
        $S(q)\to T/(\mu+Dq^{2})$.

  \item[Sunset] The canonical two-loop ($L=2$) test diagram: three lines run between
        two vertices.

  \item[Superset theorem] The fact that a Taylor expansion at order $N$ contains, on
        the degrees it shares, a byte-identical copy of the expansion at any lower
        order $M\le N$ (\code{by\_tp@N} $\supseteq$ \code{by\_tp@M}). It is the
        justification for the expand cache and for downgrade-filtering.

  \item[Sutherland--Hodgman] The polygon-clipping algorithm used to build the
        two-dimensional ($m=2$) integration region.

  \item[Sweep axis] In the cumulant estimator, the single external leg whose lag is
        left as \code{None}; the cumulant is reported as a function of \emph{its} lag,
        with the other legs held fixed.

  \item[Symanzik polynomials ($\Usym$, $\Fsym$/$Q_{\mathrm{eff}}$)] The two graph
        polynomials a Feynman loop integral reduces to once Schwinger/Feynman
        parameters are introduced. (See the boxed definition below.)

  \item[Symmetry factor $\Sym(\Gamma)$] The Feynman-rule combinatorial weight of a
        diagram. (See the boxed definition below.)

  \item[\code{sympy}] A lighter, pure-Python computer-algebra system (distinct from
        Sage's \code{SR}). The engine uses it for targeted jobs --- \code{lambdify} to
        compile expressions to fast NumPy functions, \code{linsolve} to route
        momenta, \code{multiset\_permutations} to enumerate leaf orderings.

  \item[\code{\_\_slots\_\_}] A Python class optimisation that removes the per-instance
        \code{\_\_dict\_\_} (saving memory and speeding attribute access). A class with
        \code{\_\_slots\_\_} needs explicit
        \code{\_\_getstate\_\_}/\code{\_\_setstate\_\_} to be picklable.

  \item[\code{substitute\_function}] A Sage method that replaces every occurrence of a
        formal function with the result of a Python callback. It is how the formal
        convolution \code{Conv} is reduced to a concrete expression.

  \item[Squarefree polynomial] A polynomial with no repeated roots. A non-squarefree
        characteristic denominator $Q$ signals a multiplicity-$\ge 2$ (Jordan-block)
        pole, which the single-pole residue formula cannot handle.

\end{description}

\begin{defn}
The \textbf{symmetry factor} $\Sym(\Gamma)$ is the combinatorial weight a diagram
$\Gamma$ carries under the Feynman rules,
\[
  \Sym(\Gamma)\;=\;\frac{\prod_{\text{legs}} n_{\text{leg}}!}{\bigl|\Aut_{\text{fixed-ext}}(\Gamma)\bigr|},
\]
the product of per-vertex Wick leg-factorials divided by the order of the
automorphism group of the coloured incidence digraph with external leaves pinned.
In the code it is carried by \code{combinatorial\_factor} (and \code{pv} in the
spatial integrator is $\Sym(\Gamma)$ times the vertex-coupling prefactor). Under
the framework's ``Path~A'' convention the dedup multiplicity is \emph{not}
multiplied back in --- $\Sym(\Gamma)$ already accounts for the size of the
isomorphism class. The automorphism orders are computed by \code{nauty} through
Sage's \code{automorphism\_group}.
\end{defn}

\begin{defn}
The \textbf{Symanzik polynomials} arise when a Feynman loop integral is rewritten
using one Schwinger parameter per edge. The first, $\Usym=\det\Lambda$, is a sum
over the diagram's spanning trees and is purely topological; the second, $\Fsym$
(equivalently $Q_{\mathrm{eff}}$, the Schur complement $Q-N^{\mathsf
T}\Lambda^{-1}N$), carries the external-momentum dependence. The Gaussian
integral over the loop momenta collapses \emph{exactly} to a ratio built from
$\Usym$ and $\Fsym$, turning a multidimensional momentum integral into an
integral over the bounded Schwinger parameters. The degenerate direction
$\det\Lambda=\Usym\to 0$ (e.g. a self-loop) is the $\Usym^{-d/2}$ singularity
that would make naive Monte-Carlo sampling infinite-variance; \Daedalus{} avoids
it by an analytic radial-Bessel reduction.
\end{defn}

% =====================================================================
\section*{T}
% =====================================================================
\begin{description}[leftmargin=0pt,style=nextline]

  \item[Tadpole] Two related senses. (1) The linear-in-fluctuation term of the
        expanded action --- the bigrade-$(1,0)$ piece --- whose vanishing \emph{is}
        the saddle equation. (2) In diagrams, a $C$ self-loop ($u=v$): a decoupled
        $\langle\varphi^{2}\rangle$-type loop that carries no external momentum.

  \item[$\tau_c$ (\code{tauc})] The noise correlation time, equivalently the
        Ornstein--Uhlenbeck relaxation time of the colored-noise generator.

  \item[$\tau_j=t_j-t_0$] The $k-1$ time differences a $k$-point cumulant depends on
        (by time-translation invariance).

  \item[Taylor order (\code{taylor\_order}, $N$)] The total-degree truncation of the
        multivariate Taylor expansion of the action in the fluctuation fields. It
        equals the maximum number of legs a vertex can have, and defaults to
        $\max(k+2\,\code{max\_ell},2)$. It must be at least the largest prediagram
        vertex degree --- the job of \code{degree\_scan.ensure\_taylor\_order}.

  \item[\code{ThreadPoolExecutor} / \code{as\_completed}] Python standard-library
        thread-pool primitives used for opt-in parallelism --- threads, never fork,
        on macOS-safety grounds. \code{as\_completed} yields futures as they finish.

  \item[Tier 1 vs Tier 2] In the spatial pipeline, Tier~1 is the diagonal
        closed-form heat-kernel path; Tier~2 is the numerical / coupled /
        higher-derivative fallback (the v2 integrator).

  \item[Topology] An \emph{undirected} prediagram --- vertices and edges but no
        arrows yet. It is the intermediate output of the topology stage, before
        orientation makes it a prediagram.

  \item[Transfer function ($\varphi$)] The (possibly nonlinear) firing-rate map from
        input (voltage/current) to output (rate); in code, the user function literally
        named \code{phi}. Its Taylor coefficients at the saddle,
        \code{phi0\_i, phi1\_i, ...}, become the vertex couplings. (The name clashes
        with the field $\varphi$; context disambiguates.)

  \item[Tree] A connected acyclic graph; the loop-free skeleton from which
        topologies are built by adding edges.

  \item[Trust chain] The closed validation loop the repository runs on: an analytic
        diagram prediction versus an \emph{independent} brute-force simulation.
        Agreement within error bars is what licenses every analytic result.

  \item[Typed diagram (\code{TypedDiagram})] A prediagram fully decorated: a definite
        field type on every leg, a propagator $(\text{resp idx},\text{phys idx})$ on
        every edge, and a coupling on every vertex. Many typed diagrams can share one
        prediagram.

\end{description}

% =====================================================================
\section*{U}
% =====================================================================
\begin{description}[leftmargin=0pt,style=nextline]

  \item[$\Usym$ (first Symanzik polynomial)] $\det\Lambda$, the sum over the graph's
        spanning trees. (See the Symanzik boxed definition under \textbf{S}.) Its
        vanishing, $\det\Lambda\to 0$, is the source of the $\Usym^{-d/2}$
        Monte-Carlo singularity, cured analytically by the radial-Bessel reduction.

\end{description}

% =====================================================================
\section*{V}
% =====================================================================
\begin{description}[leftmargin=0pt,style=nextline]

  \item[Vertex] An interaction term of the action (cubic and higher in the fields)
        that becomes a diagram vertex. (See interaction vertex.)

  \item[Vertex coloring / colored isomorphism] Assigning colours to vertices (here
        $0=$ leaf, $1=$ internal) and requiring an isomorphism to \emph{preserve}
        them, so a leaf can only map to a leaf. Implemented via Sage's
        \code{set\_vertex}.

  \item[\code{vertex\_form\_factor\_signature}] See \emph{form-factor signature}
        (\file{pipeline/\_expand\_cache.py:208}).

  \item[\code{vertex\_role\_signature}] A depth-1 graph invariant of a vertex
        (\file{msrjd/diagrams/symmetry.py:247}); it was the basis of the
        \emph{old}, now-superseded external-Wick compensation heuristic.

  \item[\code{v2\_max} / \code{v3\_max}] Search bounds on the number of degree-2 and
        degree-$\ge 3$ vertices a topology may have: $\code{v3\_max}=k+j-2$,
        $\code{v2\_max}=k+3\ell-j-1$. (See completeness bound.)

\end{description}

% =====================================================================
\section*{W}
% =====================================================================
\begin{description}[leftmargin=0pt,style=nextline]

  \item[White noise] Noise with zero correlation time; its covariance is
        $\propto\delta(\tau)$. In the \msrjd{} action it is the single local vertex
        $-D\,\tilde\varphi^{2}$ (the bigrade-$(2,0)$ sector). Contrast colored noise.

  \item[Wick theorem] See \emph{Isserlis / Wick theorem}.

\end{description}

% =====================================================================
\section*{Z}
% =====================================================================
\begin{description}[leftmargin=0pt,style=nextline]

  \item[$\Theta(0)=0$ convention] The convention that the retarded support boundary
        $\Delta t=0$ is strictly excluded, enforced by strict inequalities and the
        Heaviside filter. (Filed under Z for ``zero''; see Heaviside.)

  \item[\code{ZMQInteractiveShell}] The IPython shell class used only by a
        ZMQ/Jupyter kernel. The fork-safety guard fingerprints a notebook by
        detecting this class --- if it is the active shell, the fork-based parallel
        path is downgraded to serial (because forking a Jupyter kernel after
        Cocoa/BLAS init crashes it).

  \item[$2^{-n_C}$ (convention factor)] The factor converting the enumeration's
        $2T$ noise-vertex normalisation to the spatial integrator's unit-amplitude
        $C$ edges, with $n_C$ the number of correlation edges.

\end{description}

% =====================================================================
%  APPENDIX B : External tools, explained from scratch
% =====================================================================
\chapter{External Tools, Explained from Scratch}
\label{app:tools}

\Daedalus{} is glue around a handful of powerful third-party libraries.  If you
came to this code from the physics side, the unfamiliar part is rarely the
field theory --- it is the \emph{tooling}: what \emph{is} SageMath, why does
everything run under \code{sage -python}, what does \code{nauty} compute, why do
the simulators use \code{numba} but the symbolic pipeline does not?  This
appendix answers those questions one tool at a time.  For each library we say
\textbf{what it is}, \textbf{why \Daedalus{} needs it}, and \textbf{where in the
code it appears} (with the rough usage count found by grepping the source, so
you can gauge how central it is).

\begin{note}
You do not install any of these by hand for the symbolic pipeline.  They all
come bundled inside the SageMath distribution (see the project's
\code{MSRJD\_diagrams} conda environment).  The simulators add \code{numba},
which is a plain \code{pip}/\code{conda} package.
\end{note}

% ---------------------------------------------------------------------
\section{SageMath --- the symbolic backbone}
% ---------------------------------------------------------------------

\textbf{What it is.}  SageMath (``Sage'') is a large open-source
\term{computer algebra system} (CAS) written in Python.  A CAS manipulates
mathematics \emph{symbolically} --- it knows that $\partial_t e^{-\mu t}=-\mu
e^{-\mu t}$ as an identity, not as a floating-point approximation on a grid.
Sage is really a Python front-end that stitches together dozens of specialised
mathematical engines (Pari for number theory, Singular for polynomials, GAP for
group theory, Maxima for calculus, \code{nauty} for graphs, \code{mpmath} for
arbitrary precision, and many more) behind one uniform Python API.  Because the
front-end \emph{is} Python, a Sage program is a Python program with extra
built-in types.

\textbf{Why \Daedalus{} needs it.}  Every object in the front half of the
pipeline is a symbol, not a number: the action $S$, the inverse propagator
$K_{\mathrm{ft}}(\omega)$ and its matrix inverse $G_{\mathrm{ft}}$, the
interaction vertices, the loop kernels.  These must be differentiated, inverted,
Fourier-transformed, Taylor-expanded, and have their poles found \emph{exactly}.
That is exactly what a CAS is for.

\begin{defn}
The \term{Symbolic Ring}, written \code{SR} in the code, is the Sage ``parent''
that all symbolic expressions belong to --- the analogue of \code{float} for
numbers.  A Sage symbol like \code{var('mu')} is an element of \code{SR}, and any
formula built from it (\code{mu + I*omega}, \code{exp(-mu*t)}) is again an
\code{SR} element.
\end{defn}

\subsection{Running under \texttt{sage -python}}
Because the pipeline imports \code{sage.all}, it cannot run under a stock Python
interpreter --- it must run under the Python that ships \emph{inside} Sage.  That
is why throughout this project you launch scripts as \code{sage -python
myscript.py} rather than \code{python myscript.py}, and why the notebooks use a
Sage Jupyter kernel.

\subsection{Parsing model text: \texttt{preparse} and \texttt{sage\_eval}}
When you type an action string such as
\code{sum(xt[i]*((Dt+mu)*x[i] + eps*x[i]\textasciicircum 3) - D*xt[i]\textasciicircum 2 for i in pop)},
two Sage facilities turn that text into a live symbolic object
(see Chapter~\ref{ch:theory-spec}):
\begin{description}[leftmargin=1.2em]
  \item[\code{preparse}] (used $\sim$7 places) rewrites Sage's mathematician-friendly
        shorthand into legal Python --- most importantly it turns the caret
        \code{\textasciicircum} (power) into \code{**}, and wraps integer
        literals so that \code{1/3} is the exact rational $\tfrac13$ rather than
        floating-point $0.333\ldots$
  \item[\code{sage\_eval}] (used $\sim$10 places) evaluates the preparsed string in a
        controlled namespace where the field names, parameters, and operator
        symbols (\code{Dt}, \code{Lap}) are already defined, returning an
        \code{SR} expression.
\end{description}

\subsection{Calculus, inversion, and poles}
\begin{itemize}
  \item \textbf{Matrix inversion} gives $G_{\mathrm{ft}}=K_{\mathrm{ft}}^{-1}$
        symbolically (Chapter~\ref{ch:propagator}).
  \item \code{fourier\_transform} (used $\sim$8 places) moves kernels between the
        time and frequency domains.
  \item \code{.roots()} (used $\sim$6 places) solves the characteristic
        polynomial of the denominator to locate the \emph{retarded poles}
        (those with $\operatorname{Im}\omega>0$) that drive the time-domain
        propagator.
  \item \code{assume(...)} declares domain facts (e.g.\ \code{mu>0}) so the CAS
        can simplify $|{\cdot}|$, choose branch cuts, and evaluate integrals
        unambiguously.
\end{itemize}

\subsection{Graphs inside Sage (and the door to \texttt{nauty})}
Sage also has first-class \term{graph} objects.  \Daedalus{} builds diagram
topologies as \code{Graph} / \code{DiGraph} and then calls
\code{canonical\_label}, \code{is\_isomorphic}, and \code{automorphism\_group} on
them.  Those three methods are the bridge to \code{nauty} (next section): they
are how the enumerator decides that two diagrams are ``the same'' and how it
counts a diagram's symmetries.

\subsection{\texttt{.sobj} --- Sage's saved-object format}
A \code{.sobj} file is Sage's version of a Python \code{pickle}: a binary dump of
a live Sage object.  The enumeration cache writes the (expensive-to-compute)
list of unique typed diagrams to \file{saved\_theories/<name>/...\_.sobj} so the
next run loads them in milliseconds instead of re-enumerating
(Chapter~\ref{ch:caching}).

% ---------------------------------------------------------------------
\section{nauty --- canonical graphs and automorphisms}
% ---------------------------------------------------------------------
\textbf{What it is.}  \code{nauty} (a playful contraction of ``\emph{No
AUTomorphisms, Yes?}''), by Brendan McKay, is the standard high-performance tool
for two closely-related graph problems:
\begin{description}[leftmargin=1.2em]
  \item[Canonical labelling] relabel a graph's vertices into a unique normal
        form so that two graphs are \term{isomorphic} (the same up to renaming
        vertices) if and only if their canonical forms are byte-identical.  This
        turns the hard ``are these two diagrams the same?'' question into a
        cheap string comparison.
  \item[Automorphism group] compute $\Aut(\Gamma)$, the group of vertex
        relabellings that map the graph to itself --- i.e.\ its symmetries ---
        and in particular its size $|\Aut(\Gamma)|$.
\end{description}

\textbf{Why \Daedalus{} needs it.}  Two things in diagrammatics reduce to graph
symmetry:
\begin{enumerate}
  \item \textbf{De-duplication.}  The raw enumerator produces many topologies
        that are really the same diagram drawn differently; canonical labelling
        collapses them to one representative (Chapter~\ref{ch:enumeration}).
  \item \textbf{The symmetry factor $\Sym(\Gamma)$.}  Each Feynman diagram is
        divided by the order of its automorphism group (with external legs held
        fixed).  \Daedalus{} computes this exactly from
        $|\Aut(\Gamma)|$ rather than guessing a combinatorial fudge factor
        (Chapter~\ref{ch:type-assignment}).
\end{enumerate}

\begin{gotcha}
You never call \code{nauty} directly in this codebase --- it is reached
\emph{transitively} through Sage's graph methods.  If you grep for ``nauty'' you
will find almost nothing; look for \code{automorphism\_group} and
\code{canonical\_label} instead.
\end{gotcha}

% ---------------------------------------------------------------------
\section{NumPy --- arrays and fast numerics}
% ---------------------------------------------------------------------
\textbf{What it is.}  NumPy is the foundational numerical-array library for
Python: the \code{ndarray} (an $n$-dimensional, homogeneously-typed array) plus
vectorised operations on it.  \textbf{Where it appears.}  Everywhere the pipeline
turns symbols into numbers: the lag grid $\tau$, the correlator arrays
$C(\tau)$, the residue matrices, the spatial momentum grids, and the simulator
output bins are all \code{ndarray}s.  When this manual says ``a callable returns
an array of shape \code{(k-1, n\_tau)}'', that array is a NumPy array.

% ---------------------------------------------------------------------
\section{SciPy --- integration, linear algebra, special functions}
% ---------------------------------------------------------------------
\textbf{What it is.}  SciPy is the scientific-computing layer built on top of
NumPy.  \Daedalus{} uses four of its sub-packages:
\begin{description}[leftmargin=1.2em]
  \item[\code{scipy.integrate}] adaptive numerical integration.  \code{quad} does
        a single one-dimensional integral; \code{nquad} does a nested
        multi-dimensional one.  These are the \emph{fallback} evaluators in the
        temporal Phase~J engine (used when the analytic mode-sum does not apply)
        and appear in the spatial chamber integrals
        (Chapters~\ref{ch:phasej-temporal},~\ref{ch:spatial-core}).
  \item[\code{scipy.linalg}] dense linear algebra.  \code{expm} is the
        matrix exponential $e^{At}$ (time-domain propagation);
        \code{solve\_continuous\_lyapunov} solves the Lyapunov equation
        $AX+XA^{\!\top}=-Q$ for a stationary covariance; eigenvalue routines back
        the linear-stability test (Chapter~\ref{ch:mean-field}).
  \item[\code{scipy.special}] special functions.  \code{kv} is the
        \term{modified Bessel function of the second kind} $K_\nu(z)$, which
        emerges from the radial Schwinger-parameter integral in the spatial
        backend; \code{gamma} is Euler's $\Gamma$ function
        (Chapter~\ref{ch:spatial-heatkernel}).
  \item[\code{scipy.optimize}] root finding.  \code{root} / \code{fsolve} are
        used as a fallback for the mean-field self-consistency solve
        (Chapter~\ref{ch:mean-field}).
\end{description}

% ---------------------------------------------------------------------
\section{sympy and \texttt{lambdify}}
% ---------------------------------------------------------------------
\textbf{What it is.}  \code{sympy} is a second, \emph{pure-Python} computer
algebra system (independent of Sage).  Its \code{lambdify} function takes a
symbolic expression and returns a fast ordinary Python function that evaluates it
on NumPy arrays --- it ``compiles'' a formula into a numeric callable.
\textbf{Where it appears.}  Narrowly: a small number of modules convert a
symbolic spatial form-factor into a \code{lambdify}-ed callable so it can be
sampled cheaply on a momentum grid.  The bulk of the symbolic work stays in
Sage; \code{sympy} is used only at this symbolic-to-numeric hand-off.

% ---------------------------------------------------------------------
\section{numba --- JIT compilation for the simulators}
% ---------------------------------------------------------------------
\textbf{What it is.}  \code{numba} is a just-in-time (JIT) compiler: decorate a
numeric Python function with \code{@numba.njit} and the first call compiles it to
machine code through LLVM, after which it runs at C-like speed.
\textbf{Where it appears.}  \emph{Only} in the independent simulators
(\file{models/}, $\sim$11 files), never in the symbolic pipeline.  The
Euler--Maruyama and Hawkes integrators are tight inner loops over millions of
time-steps --- exactly the workload JIT compilation is for
(Chapter~\ref{ch:simulators}).

\begin{gotcha}
\code{numba} can only compile functions whose arguments are plain machine types
(\code{float}, \code{int}, NumPy arrays).  The resolved parameters coming out of
the pipeline are \emph{Sage ring elements}, which \code{numba} cannot type ---
so the simulator-driving notebook cells explicitly cast every scalar with
\code{float(...)} before calling an \code{@njit} simulator.  Forgetting the cast
gives an opaque numba typing error.
\end{gotcha}

% ---------------------------------------------------------------------
\section{matplotlib --- plotting}
% ---------------------------------------------------------------------
\textbf{What it is.}  The standard Python 2-D plotting library.
\textbf{Where it appears.}  \code{dd.plot\_cumulant} builds the theory-vs-lag
curves (and the $k\!\ge\!3$ slice/heat-map panels and the simulator overlays)
with matplotlib \code{Figure}/\code{Axes} objects
(Chapter~\ref{ch:engine-api}).

% ---------------------------------------------------------------------
\section{networkx --- a lightweight graph helper}
% ---------------------------------------------------------------------
\textbf{What it is.}  \code{networkx} is a pure-Python graph library (nodes,
edges, traversals).  \textbf{Where it appears.}  In a single module, as a
convenience for representing/handling diagram structure; the heavy graph work
(canonicalisation, automorphisms) is done by Sage/\code{nauty}, not
\code{networkx}.

% ---------------------------------------------------------------------
\section{multiprocessing --- parallelism, and the macOS fork hazard}
% ---------------------------------------------------------------------
\textbf{What it is.}  Python's standard library for running work in separate
\emph{processes} (sidestepping the Global Interpreter Lock).  A process pool can
\emph{fork} (clone the current process) or \emph{spawn} (start a fresh
interpreter).  \textbf{Where it appears.}  $\sim$7 modules optionally parallelise
diagram enumeration and the Phase~J $\tau$-batch with a fork-based pool.

\begin{gotcha}
On macOS, \code{fork}-ing a process that has already initialised Cocoa, BLAS, or
a Matplotlib backend --- which a Jupyter kernel has --- can hard-crash the kernel
\emph{and} the OS.  \Daedalus{} therefore guards every fork behind a
\term{fork-safety} check that degrades to serial execution inside a macOS Jupyter
kernel, while keeping the fork (the fast path) in terminals, scripts, and on
Linux.  This is covered in detail in Chapter~\ref{ch:type-assignment}.
\end{gotcha}

% ---------------------------------------------------------------------
\section{mpmath --- arbitrary-precision arithmetic}
% ---------------------------------------------------------------------
\textbf{What it is.}  \code{mpmath} is a pure-Python library for real and complex
floating-point arithmetic at \emph{arbitrary} precision (hundreds of digits if
asked).  Sage uses it internally, and a few precision-sensitive spots in the
pipeline ($\sim$3 modules) lean on it where ordinary double precision would lose
significant digits --- for example in mode-sums with closely-spaced poles
(see the $m\!\ge\!3$ precision note in Chapter~\ref{ch:phasej-temporal}).

% ---------------------------------------------------------------------
\section{The Python standard library}
% ---------------------------------------------------------------------
A few standard modules do quiet but essential work:
\begin{description}[leftmargin=1.2em]
  \item[\code{ast}] safe parsing of Python literals; \code{ast.literal\_eval}
        reads the \code{METADATA}/\code{DEFAULT\_FUNDAMENTAL} dictionaries out of
        a \file{.theory.py} file without executing arbitrary code
        (Chapter~\ref{ch:theory-files}).
  \item[\code{importlib}] dynamic import; loads a \file{.theory.py} file by path
        at run time so new theories need no registration.
  \item[\code{re}] regular expressions; light text munging of action/equation
        strings.
  \item[\code{dataclasses}] the \code{@dataclass} decorator that defines the
        \code{Config} object (Chapter~\ref{ch:engine-api}).
  \item[\code{itertools}] combinatorial iterators; \code{product} generates leg
        assignments and Wick pairings.
\end{description}

\begin{note}
\textbf{Rule of thumb for ``what library does this?''}  Anything
\emph{symbolic} (expressions, inversion, Fourier transforms, graph symmetry) is
\textbf{Sage} (with \code{nauty} under the graph hood).  Anything
\emph{numeric in the pipeline} (integrals, eigenvalues, Bessel functions) is
\textbf{SciPy/NumPy}.  Anything in a \emph{simulator} is \textbf{numba+NumPy}.
The \emph{UI} is \textbf{ipywidgets}.  Everything else is the Python standard
library.
\end{note}

% ---------------------------------------------------------------------
%  Appendix C : An Annotated File Index of the Source Tree
% ---------------------------------------------------------------------
\chapter{An Annotated File Index of the Source Tree}\label{app:file-index}

The preceding chapters walked the pipeline \emph{by stage} --- the
order in which data flows from a model description to a plotted cumulant.
This appendix walks the same code \emph{by directory}. It is a map, not a
tutorial: one line per module, what that module is responsible for, and
which chapter of this manual opens it up. If you arrive at the repository
cold, knowing the physics but not the layout, read this appendix first to
find the file you want, then jump to the chapter the table points you at.

The repository has a deliberate two-layer split, and understanding it is
the single most useful orientation you can have:

\begin{description}[style=nextline,leftmargin=1.4em]
  \item[\file{msrjd/} --- the engine.] The model-agnostic physics core. It
        knows about \msrjd{} actions, Feynman diagrams, propagators, loop
        integrals, and symmetry factors, but it knows \emph{nothing} about
        any particular model. It is pure \code{import}-able Python (running
        under SageMath) with no notebook dependencies.
  \item[\file{pipeline/} --- the driver.] The orchestration layer that
        sequences the engine's stages, threads data dictionaries between
        them, caches the expensive ones to disk, and exposes a single
        master function, \code{compute\_cumulants}. It also hosts the
        \emph{authoring} surface: the \code{TheoryBuilder} API, the
        text-to-Sage compiler, the colored-noise rewriter, and the
        \code{ipywidgets} theory-builder UI.
  \item[\file{models/} --- the validators.] Independent direct simulators
        (Euler--Maruyama SDE integrators, Poisson--Hawkes steppers,
        spectral PDE solvers) and the cumulant estimators that turn their
        output into the same correlator the engine predicts. These share
        \emph{no code} with the engine; they exist so every prediction can
        be checked against an experiment.
  \item[\file{theories/} --- the models.] The declarative
        \file{*.theory.py} files: one physical model each, the
        version-controllable single source of truth.
  \item[\file{notebooks/} --- the front desk.] The user-facing convenience
        layer \code{daedalus.py} and the example, template, and validation
        notebooks that drive it.
\end{description}

\begin{note}
\textbf{How to read the tables.} Each table lists files under one
directory. The first column is the file (always relative to the repository
root, the navy monospace convention of this manual). The second column is
its role in one phrase. The third column, ``Ch.'', is the chapter that
covers it in depth --- a dash means the file is peripheral (a stub, an
oracle-only cross-check, or a reference asset) and is documented here only.
The roles below are taken from each module's own docstring and verified
against its code; where a module's docstring and its behaviour disagree,
the table reflects the \emph{behaviour}, and the discrepancy is flagged in
Appendix~D (Known Issues and Open Questions).
\end{note}

\begin{gotcha}
\textbf{Two things named ``serialize'' and two things named
``propagator''.} The tree contains deliberate near-collisions that trip up
newcomers. There are \emph{two} serializers ---
\file{pipeline/theory\_serialize.py} (the \file{.theory.py} \emph{text}
format) and \file{msrjd/core/serialize.py} (the \code{.sobj} \emph{binary}
cache) --- and \emph{two} files called \code{propagator.py}, of which only
\file{pipeline/\_propagator.py} is live; \file{msrjd/core/propagator.py} is
a docstring-only stub. The tables call out each one. When you grep for a
name, check which of the pair you have landed on.
\end{gotcha}

% =====================================================================
\section{The engine: \texttt{msrjd/}}
% =====================================================================

\file{msrjd/} is the physics core, organised into five sub-packages:
\file{core/} (the symbolic action, vertices, propagator-time-domain, and
caching primitives), \file{enumeration/} (diagram topology generation),
\file{diagrams/} (typing, causality, symmetry), \file{integration/} (the
two loop-integral back ends, temporal and spatial), and a top-level
\file{fork\_safety.py} guard shared across the parallel paths.

\subsection{\texttt{msrjd/core/} --- the symbolic core}

This is where a model dict becomes a symbolic \msrjd{} action, gets
Taylor-expanded around the saddle, and is sorted into the free part (the
inverse propagator) and the interaction vertices. Everything here is a
SageMath symbolic-ring object; nothing is bound to a number yet.

\begin{center}
\begin{longtable}{@{}>{\ttfamily\footnotesize\raggedright}p{0.31\textwidth} >{\raggedright}p{0.50\textwidth} >{\centering}p{0.05\textwidth}@{}}
\toprule
\normalfont\textbf{file} & \textbf{role} & \textbf{Ch.} \tabularnewline
\midrule
\endhead
\bottomrule
\endfoot
msrjd/core/field\_theory.py
  & The \code{FieldTheory} expander: builds the symbolic namespace,
    evaluates the action lambda, resolves \code{Conv(...)} atoms, injects
    correlated-noise cumulants, runs the multivariate Taylor expansion, and
    bigrade-classifies the result into \code{ft.\_by\_tp} sectors. The
    front end of the whole pipeline.
  & \ref{ch:fieldtheory-core} \tabularnewline
\addlinespace
msrjd/core/vertices.py
  & Decomposes each bigrade sector polynomial into individual typed
    monomials: \code{VertexType} (interaction), \code{SourceType} (noise),
    and the kernel-carrying subclasses \code{ConvVertexType} (conductance
    $g(\tau)$ kernel) and \code{NoiseSourceType} (non-local $\kappa(\tau)$
    kernel).
  & \ref{ch:fieldtheory-core} \tabularnewline
\addlinespace
msrjd/core/convolution.py
  & The formal \code{Conv(kernel, field)} operator and its reduction rules
    --- how a registered synaptic kernel multiplied by a field is held
    symbolically until the propagator step can resolve it.
  & \ref{ch:fieldtheory-core} \tabularnewline
\addlinespace
msrjd/core/serialize.py
  & Save/load of an \emph{expanded} theory to a directory
    (\code{metadata.json} $+$ \code{symbolic\_data.sobj}). The
    \code{.sobj}-cache serializer --- not the \file{.theory.py} text format
    (that is \file{pipeline/theory\_serialize.py}).
  & \ref{ch:fieldtheory-core},\,\ref{ch:caching} \tabularnewline
\addlinespace
msrjd/core/cache.py
  & \code{PipelineCache}: the low-level stage-keyed \code{.sobj} store both
    higher-level caches build on. Keys each pipeline stage (enumeration,
    filtering, typing, dedup, grouping) by $(k,\ell)$ so a notebook restart
    skips the expensive recomputation.
  & \ref{ch:caching} \tabularnewline
\addlinespace
msrjd/core/propagator.py
  & \textbf{Stub only.} A ten-line docstring marked ``Status: Not yet
    extracted --- logic currently lives in the demo notebook''. Imported by
    nothing on the working path; the live propagator code is
    \file{pipeline/\_propagator.py}.
  & --- \tabularnewline
\end{longtable}
\end{center}

\subsection{\texttt{msrjd/enumeration/} --- diagram topology}

Given only the two integers $k$ (correlator order) and $\ell$ (loop order),
this package produces the complete, deduplicated list of bare graph
topologies --- \term{prediagrams} --- that could contribute, before any
physics is attached. It is theory-agnostic: it never looks at the action.

\begin{center}
\begin{longtable}{@{}>{\ttfamily\footnotesize\raggedright}p{0.31\textwidth} >{\raggedright}p{0.50\textwidth} >{\centering}p{0.05\textwidth}@{}}
\toprule
\normalfont\textbf{file} & \textbf{role} & \textbf{Ch.} \tabularnewline
\midrule
\endhead
\bottomrule
\endfoot
msrjd/enumeration/\newline loop\_diagram\_enumeration.py
  & The three-stage topology generator: nonisomorphic trees $\to$
    topologies (add $\ell$ edges) $\to$ prediagrams (orient every edge),
    each stage a filter-plus-deduplicate. Calls the \code{nauty} graph
    canonicaliser through Sage for the isomorphism rejection.
  & \ref{ch:enumeration} \tabularnewline
\addlinespace
msrjd/enumeration/degree\_scan.py
  & Build Phase C, the feedback bridge. Scans the enumerated prediagrams
    for their maximum vertex degree and, if a diagram needs a higher-degree
    vertex than the action was expanded to, triggers a re-expansion via
    \code{available\_degrees}.
  & \ref{ch:enumeration} \tabularnewline
\end{longtable}
\end{center}

\subsection{\texttt{msrjd/diagrams/} --- typing, causality, symmetry}

This package turns a bare oriented graph into a fully labelled Feynman
diagram you can integrate: it decorates edges with propagators and vertices
with monomials (typing), prunes the acausal ones, deduplicates, and
attaches the symmetry factor $\Sym(\Gamma)$.

\begin{center}
\begin{longtable}{@{}>{\ttfamily\footnotesize\raggedright}p{0.31\textwidth} >{\raggedright}p{0.50\textwidth} >{\centering}p{0.05\textwidth}@{}}
\toprule
\normalfont\textbf{file} & \textbf{role} & \textbf{Ch.} \tabularnewline
\midrule
\endhead
\bottomrule
\endfoot
msrjd/diagrams/filter.py
  & Build Phase D. A cheap set-membership check that discards prediagrams
    whose vertex-degree signature no vertex or source in the theory could
    supply --- run \emph{before} the expensive typing step.
  & \ref{ch:type-assignment} \tabularnewline
\addlinespace
msrjd/diagrams/\newline type\_assignment.py
  & Build Phase E, the heart of the subsystem. A backtracking
    constraint-satisfaction engine that emits every valid
    \code{TypedDiagram}: which monomial sits at each vertex, which external
    field on each leg, and which propagator $G_{ij}$ rides each directed
    edge.
  & \ref{ch:type-assignment} \tabularnewline
\addlinespace
msrjd/diagrams/causality.py
  & Build Phase F. Prunes typed diagrams that violate the retarded boundary
    condition --- structurally, the graph must be a DAG (no directed cycle);
    analytically, all propagator poles in the upper half $\omega$-plane.
  & \ref{ch:type-assignment} \tabularnewline
\addlinespace
msrjd/diagrams/symmetry.py
  & Build Phase G. Deduplicates typed diagrams via a complete
    graph-isomorphism invariant (Sage canonical labelling) and computes the
    symmetry factor
    $\Sym(\Gamma)=\big(\prod_v\prod_\ell n_{v,\ell}!\big)/|\Aut_{\text{fixed-ext}}(\Gamma)|$.
    Also hosts \code{external\_wick\_compensation} and
    \code{vertex\_role\_signature}, used by the integrators.
  & \ref{ch:type-assignment},\,\ref{ch:grouped-phasej} \tabularnewline
\addlinespace
msrjd/fork\_safety.py
  & The canonical fork-safety guard shared by every fork-based
    multiprocessing path. Detects ``am I inside a macOS Jupyter/ZMQ
    kernel?'' and, if so, silently degrades a parallel path to serial to
    avoid the fork-after-Cocoa/BLAS kernel crash.
  & \ref{ch:type-assignment} \tabularnewline
\end{longtable}
\end{center}

\subsection{\texttt{msrjd/integration/time\_domain/} --- the temporal loop integrator}

This is ``Phase J'': it takes each typed diagram's product of retarded
propagators and integrates over the internal vertex times. Because every
retarded propagator is a sum of decaying exponentials, the integral over a
causal-ordering polytope has a closed form, so the engine almost never
falls back to brute quadrature.

\begin{center}
\begin{longtable}{@{}>{\ttfamily\footnotesize\raggedright}p{0.31\textwidth} >{\raggedright}p{0.50\textwidth} >{\centering}p{0.05\textwidth}@{}}
\toprule
\normalfont\textbf{file} & \textbf{role} & \textbf{Ch.} \tabularnewline
\midrule
\endhead
\bottomrule
\endfoot
msrjd/integration/\newline time\_domain/final\_integral.py
  & The 5{,}269-line core. \code{integrate\_diagram} and the family of
    closed-form mode-sum evaluators (\code{\_integrate\_1d\_polytope\_modesum},
    \code{\_integrate\_2d\_polygon\_modesum},
    \code{\_integrate\_nd\_polytope\_poset\_modesum}), the \code{EdgeModeSum}
    representation, and the pole-tuple enumerator. The numerical heart of
    the temporal pipeline.
  & \ref{ch:phasej-temporal} \tabularnewline
\addlinespace
msrjd/integration/\newline time\_domain/pipeline.py
  & The orchestration/batching layer over the per-diagram core:
    \code{compute\_correction\_td} loops the integrator once per typed
    diagram and assembles the per-loop-order correction. Carries the
    fork-guarded \code{total\_C\_batch} parallel path.
  & \ref{ch:phasej-temporal} \tabularnewline
\addlinespace
msrjd/integration/\newline time\_domain/propagator\_td.py
  & Turns frequency-domain propagator data (poles, residues,
    \code{D\_delta}) into the time-domain matrix $G(t)$ = pole-residue sum
    $+\;\delta(t)$ part, with per-edge lookups \code{G\_t\_entry} and
    \code{G\_t\_delta\_coeff} the integrators call.
  & \ref{ch:propagator},\,\ref{ch:phasej-temporal} \tabularnewline
\addlinespace
msrjd/integration/\newline time\_domain/grouped\_integral.py
  & The grouped evaluator \code{integrate\_grouped\_diagram} and its
    closed-form helpers: evaluates all typed diagrams sharing one prediagram
    together, summing the per-edge propagator-index choices into one
    mode-sum instead of re-integrating each.
  & \ref{ch:grouped-phasej} \tabularnewline
\addlinespace
msrjd/integration/\newline time\_domain/subgraph.py
  & \textbf{Near-stub.} \code{identify\_loop\_subgraphs} returns \code{[]}
    (a never-finished loop-reduction route). Still re-exported by the
    package \code{\_\_init\_\_} and touched by a test, so it is a live import
    surface with a dead body.
  & --- \tabularnewline
\end{longtable}
\end{center}

\subsection{\texttt{msrjd/integration/spatial/} --- the spatial loop integrator}

For a stochastic \emph{partial} differential equation the engine does the
momentum integral first and analytically: every diagram, any $k$, any
$\ell$, any dimension $d$, reduces to the \emph{same} Symanzik integral and
is evaluated by the \emph{same} code. Several files in this package are
banner-annotated \emph{oracle-only}: independent cross-checks that
\code{compute\_cumulants} never calls but the production code is tested
against.

\begin{center}
\begin{longtable}{@{}>{\ttfamily\footnotesize\raggedright}p{0.31\textwidth} >{\raggedright}p{0.50\textwidth} >{\centering}p{0.05\textwidth}@{}}
\toprule
\normalfont\textbf{file} & \textbf{role} & \textbf{Ch.} \tabularnewline
\midrule
\endhead
\bottomrule
\endfoot
msrjd/integration/\newline spatial/full\_integrator.py
  & The production integrator. \code{diagram\_kinematic} and the batched
    Symanzik core (\code{\_momentum\_factor\_batch},
    \code{\_symanzik\_kernel\_batch}, \code{\_heat\_kernel\_x\_general}), the
    \code{grid}/\code{mc}/\code{bessel} back ends, and the assembly chain
    \code{diagram\_value\_x} $\to$ \code{diagram\_correlator\_x} $\to$
    \code{correlator\_2pt\_x}. ``One genuine integral evaluates every
    enumerated diagram.''
  & \ref{ch:spatial-core},\,\ref{ch:spatial-heatkernel} \tabularnewline
\addlinespace
msrjd/integration/\newline spatial/diagram\_descriptor.py
  & \code{diagram\_to\_cstack}: flattens a typed diagram into the canonical
    \code{CStackDiagram} edge list (frozen \code{CEdge} records), contracting
    each 2-point noise source into a single $C$ edge. Topology-blind by
    design --- no bubble/tadpole/sunset branch.
  & \ref{ch:spatial-core} \tabularnewline
\addlinespace
msrjd/integration/\newline spatial/momentum\_routing.py
  & \code{route\_momenta}: solves momentum conservation at each spatial
    vertex with SymPy's \code{linsolve}, returning the per-edge routing
    coefficients $(a_e,b_e)$ via \code{edge\_coeffs()}. The only SymPy in
    the spatial subsystem.
  & \ref{ch:spatial-core} \tabularnewline
\addlinespace
msrjd/integration/\newline spatial/spatial\_reduce.py
  & The textbook single-sample Symanzik reference:
    \code{symanzik\_matrices} ($\Lambda,N,Q$),
    \code{symanzik\_polynomials} ($\Usym=\det\Lambda$, $Q_{\text{eff}}$ via
    Schur complement), \code{momentum\_integral}. Exists so the batched
    production version can be checked against it.
  & \ref{ch:spatial-core} \tabularnewline
\addlinespace
msrjd/integration/\newline spatial/causal\_chambers.py
  & \code{causal\_chambers}: wraps the temporal pipeline's \code{\_CausalPoset}
    and \code{\_enumerate\_linear\_extensions} to tile the internal-time
    domain into smooth ordering chambers. Also a reference
    \code{integrate\_over\_chambers} (nested \code{scipy.quad}).
  & \ref{ch:spatial-core} \tabularnewline
\addlinespace
msrjd/integration/\newline spatial/heat\_kernel.py
  & The real-space propagator layer. \code{build\_spatial\_propagator} /
    \code{make\_g\_tx\_callables} and the analytic inverse Fourier transform
    of the momentum Gaussian to the heat kernel
    $(4\pi Bt)^{-d/2}\ee^{-x^2/4Bt}$ --- the trick that retired the
    $q$-grid.
  & \ref{ch:spatial-heatkernel} \tabularnewline
\addlinespace
msrjd/integration/\newline spatial/spatial\_correlator.py
  & The analytic $q\to x$ Fourier transforms used to land the diagram answer
    $C(q,\tau)$ in physical $(x,\tau)$ space, and the per-mode heat-kernel
    structure the bridge certifies against.
  & \ref{ch:spatial-heatkernel} \tabularnewline
\addlinespace
msrjd/integration/\newline spatial/spectral\_propagator.py
  & Builds the coupled multi-field free propagator $\ee^{-Mt}$ via the
    spectral (eigenprojector) decomposition of the reaction matrix $M$ --- the
    equal-diffusion coupled-tree path.
  & \ref{ch:spatial-coupled} \tabularnewline
\addlinespace
msrjd/integration/\newline spatial/dyson\_dressing.py
  & The Dyson--Duhamel series for unequal diffusion ($\mathcal{D}\ne D_0 I$):
    splits $\mathcal{D}=D_0 I+\hat{\mathcal{D}}$ and re-sums the residual
    $\hat{\mathcal{D}}|k|^2$ correction order by order.
  & \ref{ch:spatial-coupled} \tabularnewline
\addlinespace
msrjd/integration/\newline spatial/pipeline\_bridge.py
  & The 2{,}157-line spatial$\to$shared-pipeline bridge and dispatcher.
    ``Route through the shared pipeline, then certify'': substitutes
    $\Lap\to-q^2$, runs the same enumeration/typing/Phase-J path, certifies
    against the heat-kernel structure, then does the analytic $q\to x$
    transform. Hosts the \code{SPATIAL\_*} environment knobs and the memory
    guard.
  & \ref{ch:spatial-core},\,\ref{ch:spatial-coupled} \tabularnewline
\addlinespace
msrjd/integration/\newline spatial/generic\_evaluator.py
  & \emph{Oracle only.} An older Dyson-convolution evaluator that still
    branches on \code{is\_tadpole\_like()} --- the one surviving topology
    branch, kept as an independent cross-check. Not on the production path.
  & --- \tabularnewline
\addlinespace
msrjd/integration/\newline spatial/loop\_dyson.py
  & \emph{Oracle only.} Loop-level single-mode Dyson cross-check, reached
    only by its own test.
  & --- \tabularnewline
\addlinespace
msrjd/integration/\newline spatial/loop\_parametric.py
  & \emph{Oracle only.} A parametric (Feynman-parameter) loop cross-check.
  & --- \tabularnewline
\addlinespace
msrjd/integration/\newline spatial/temporal\_integrate.py
  & \emph{Oracle only.} A spatial-path temporal-integration cross-check.
  & --- \tabularnewline
\end{longtable}
\end{center}

\begin{note}
\textbf{The diagnostic helpers.} A few short members of this package are
inspection-only and never branched on in production:
\code{CStackDiagram.is\_tadpole\_like()} in
\file{diagram\_descriptor.py}, for example, classifies a diagram for
display but the production evaluator stays topology-blind. Whether a loop
momentum couples to the external $q$ is a property of the Symanzik $\Fsym$
\emph{downstream}, not of any descriptor flag --- which is the whole
``one integral, every diagram'' point.
\end{note}

\subsection{\texttt{msrjd/integration/} --- the two reference back ends}

Two files sit directly under \file{integration/}, outside both
\file{time\_domain/} and \file{spatial/}. Both predate the production
integrators and survive as references or stubs.

\begin{center}
\begin{longtable}{@{}>{\ttfamily\footnotesize\raggedright}p{0.31\textwidth} >{\raggedright}p{0.50\textwidth} >{\centering}p{0.05\textwidth}@{}}
\toprule
\normalfont\textbf{file} & \textbf{role} & \textbf{Ch.} \tabularnewline
\midrule
\endhead
\bottomrule
\endfoot
msrjd/integration/symbolic.py
  & A frequency-domain symbolic diagram-integral constructor for
    \emph{stationary} systems (the Helias \& Dahmen Ch.~9 method). A
    reference implementation, distinct from the time-domain production path.
  & \ref{ch:phasej-temporal} \tabularnewline
\addlinespace
msrjd/integration/numerical.py
  & \textbf{Stub only.} Docstring marked ``Status: Not started'' (an
    intended adaptive-quadrature / Monte-Carlo fallback). Imported by
    nothing; the \code{numerical} grep hits elsewhere are the English word.
  & --- \tabularnewline
\end{longtable}
\end{center}

% =====================================================================
\section{The driver: \texttt{pipeline/}}
% =====================================================================

\file{pipeline/} sequences the engine and exposes the authoring surface.
It splits into the master orchestrator (\file{compute.py}), the
underscore-prefixed phase modules it calls (\file{\_propagator.py},
\file{\_mean\_field*.py}, \file{\_diagrams.py}, the two caches), the
authoring API (\file{theory*.py}), the colored-noise rewriter, the result
accessors and serializers, and the \file{ui/} sub-package.

\subsection{Orchestration and the phase modules}

\begin{center}
\begin{longtable}{@{}>{\ttfamily\footnotesize\raggedright}p{0.31\textwidth} >{\raggedright}p{0.50\textwidth} >{\centering}p{0.05\textwidth}@{}}
\toprule
\normalfont\textbf{file} & \textbf{role} & \textbf{Ch.} \tabularnewline
\midrule
\endhead
\bottomrule
\endfoot
pipeline/compute.py
  & The spine. \code{compute\_cumulants(model, k, max\_ell, ...)} sequences
    the seven phases (expand $\to$ propagator $\to$ mean field $\to$ poles
    $\to$ enumerate $\to$ classify $\to$ Phase~J), branches spatial-vs-temporal,
    assembles the per-loop-order decomposition, and records wall-time. Does
    little arithmetic itself --- it is glue.
  & \ref{ch:compute-orchestration} \tabularnewline
\addlinespace
pipeline/\_propagator.py
  & The live propagator engine. Builds $K_{\text{ker}}\to K_{\text{ft}}\to
    G_{\text{ft}}$, finds the retarded poles, and computes the residue
    matrices \code{C\_mats} and the instantaneous matrix \code{D\_delta}.
  & \ref{ch:propagator} \tabularnewline
\addlinespace
pipeline/\_mean\_field.py
  & The legacy iteration saddle solver \code{solve\_mean\_field}: for
    theories whose steady state is a self-consistency closure
    $n^*_i=\phi(v^*_i)$. Iterates the rate-like variables, evaluates the
    rest on the fly.
  & \ref{ch:mean-field} \tabularnewline
\addlinespace
pipeline/\_mean\_field\_dae.py
  & The DAE multi-root solver \code{solve\_mean\_field\_dae}: sets
    $\partial_t\to0$ on the declared \code{.equation(...)} system,
    multi-start Newton for all roots, dedup, optional linear-stability
    classification, selects one root by index.
  & \ref{ch:mean-field} \tabularnewline
\addlinespace
pipeline/\_diagrams.py
  & The enumeration cache wrapper. \code{enumerate\_unique\_diagrams} wires
    the four-stage diagram pipeline (prediagrams $\to$ typed $\to$ causal $\to$
    unique) together with disk caching keyed by structure.
  & \ref{ch:type-assignment},\,\ref{ch:caching} \tabularnewline
\addlinespace
pipeline/\_expand\_cache.py
  & Disk cache for \code{FieldTheory.expand()} results, keyed by theory
    structure and Taylor order, written under
    \file{saved\_theories/<slug>/}. Avoids repeating the (up to
    $\sim$90-minute) Taylor step.
  & \ref{ch:caching} \tabularnewline
\addlinespace
pipeline/\_precompute.py
  & \code{precompute(model)}: the one-time structural pass
    (expand@order-2 $+$ propagator $+$ mean-field sanity check) that warms the
    caches and is invoked by the UI's \emph{Pre-compute} button.
  & \ref{ch:caching} \tabularnewline
\addlinespace
pipeline/\_grouped\_phase\_j.py
  & The per-prediagram grouping wrapper \code{compute\_correction\_td\_grouped}
    --- the dispatch site that hands one prediagram's typed-diagram family to
    the grouped evaluator.
  & \ref{ch:grouped-phasej} \tabularnewline
\end{longtable}
\end{center}

\subsection{Authoring: builders, compiler, templates, serialization}

\begin{center}
\begin{longtable}{@{}>{\ttfamily\footnotesize\raggedright}p{0.31\textwidth} >{\raggedright}p{0.50\textwidth} >{\centering}p{0.05\textwidth}@{}}
\toprule
\normalfont\textbf{file} & \textbf{role} & \textbf{Ch.} \tabularnewline
\midrule
\endhead
\bottomrule
\endfoot
pipeline/theory.py
  & The fluent builder classes \code{TemporalTheoryBuilder},
    \code{SpatialTheoryBuilder}, and the back-compatible \code{TheoryBuilder},
    plus \code{build()}. Accumulates \code{.physical\_field(...)},
    \code{.parameter(...)}, \code{.set\_action\_text(...)},
    \code{.equation(...)} declarations into the model dict.
  & \ref{ch:theory-spec} \tabularnewline
\addlinespace
pipeline/theory\_compiler.py
  & Text-expression $\to$ Sage-lambda compilation. Turns the human-typed
    action / mean-field / kernel / noise strings into the \code{SR}-returning
    lambdas the \code{FieldTheory} expander expects; hosts \code{\_FieldScalar}
    and the \code{ns} namespace.
  & \ref{ch:theory-spec} \tabularnewline
\addlinespace
pipeline/theory\_templates.py
  & Pre-canned action / kernel / noise templates that emit the Sage lambdas
    directly, so a user can pick a template instead of typing the action.
  & \ref{ch:theory-spec} \tabularnewline
\addlinespace
pipeline/theory\_serialize.py
  & The \file{.theory.py} \emph{text} format bridge.
    \code{render\_theory\_file} / \code{save\_theory\_to\_file} emit source
    text from a spec dict; \code{load\_spec\_from\_file} parses it back via
    the \code{ast} module (not by executing it). The round-trip behind the
    UI's Save/Load.
  & \ref{ch:theory-files} \tabularnewline
\addlinespace
pipeline/colored\_to\_markovian.py
  & The builder-level preprocessor that rewrites a colored-noise cumulant
    term into white noise by adding one auxiliary OU field per noisy field,
    so the analytic integrators apply at $\ell\ge1$.
  & \ref{ch:colored-markovian} \tabularnewline
\addlinespace
pipeline/spatial\_operator\_ir.py
  & The \code{Lap}/\code{Dt}/\code{Dx} operator intermediate representation:
    hosts spatial-derivative vertices (Model~B, Burgers, KPZ) as inert
    symbolic nodes, gives the operator algebra, and lowers them to ring
    generators so the Taylor machinery works unchanged.
  & \ref{ch:theory-spec},\,\ref{ch:spatial-coupled} \tabularnewline
\end{longtable}
\end{center}

\subsection{Result access, serialization, and reporting}

\begin{center}
\begin{longtable}{@{}>{\ttfamily\footnotesize\raggedright}p{0.31\textwidth} >{\raggedright}p{0.50\textwidth} >{\centering}p{0.05\textwidth}@{}}
\toprule
\normalfont\textbf{file} & \textbf{role} & \textbf{Ch.} \tabularnewline
\midrule
\endhead
\bottomrule
\endfoot
pipeline/access.py
  & Natural-name accessors for \code{compute\_cumulants} results: hides the
    internal \msrjd{} naming (\code{dn} for fluctuations, \code{nstar} for
    saddle values) behind the model's declared user-facing names, via
    \code{model['naming\_convention']}.
  & \ref{ch:cumulants-moments} \tabularnewline
\addlinespace
pipeline/save.py
  & NPZ $+$ CSV serialization of results with a $k$-/$\ell$-/model-adaptive
    schema: per-loop-order curves (\code{C\_tree}, \code{C\_<n>\_loop},
    \code{C\_total}) plus mean-field arrays.
  & \ref{ch:cumulants-moments} \tabularnewline
\addlinespace
pipeline/report.py
  & A multi-page PDF showing prediagrams, typed-diagram assignments, and
    per-diagram numerical contributions --- an intuition-building prototype.
  & \ref{ch:cumulants-moments} \tabularnewline
\addlinespace
pipeline/\_\_init\_\_.py
  & The package's public API surface: re-exports \code{compute\_cumulants}
    and the high-level helpers so \code{from pipeline import ...} works in
    scripts and notebooks.
  & \ref{ch:compute-orchestration} \tabularnewline
\end{longtable}
\end{center}

\subsection{\texttt{pipeline/ui/} --- the graphical theory builder}

\begin{center}
\begin{longtable}{@{}>{\ttfamily\footnotesize\raggedright}p{0.31\textwidth} >{\raggedright}p{0.50\textwidth} >{\centering}p{0.05\textwidth}@{}}
\toprule
\normalfont\textbf{file} & \textbf{role} & \textbf{Ch.} \tabularnewline
\midrule
\endhead
\bottomrule
\endfoot
pipeline/ui/main.py
  & The 3{,}071-line \code{TheoryUI}: a 10-tab \code{ipywidgets} editor with
    a live readiness sidebar and Save/Preview/Pre-compute/Load/Reset
    buttons; \code{\_collect()} snapshots the form into a spec dict for
    \code{theory\_serialize}.
  & \ref{ch:ui} \tabularnewline
\addlinespace
pipeline/ui/widgets.py
  & The reusable \code{ipywidgets} primitives the tabs compose from:
    \code{DynamicTable}, \code{vector\_input}, \code{matrix\_input},
    \code{expression\_input}, \code{textarea\_input}, \code{paste\_button},
    \code{read\_system\_clipboard}.
  & \ref{ch:ui} \tabularnewline
\addlinespace
pipeline/ui/\_\_init\_\_.py
  & Re-exports \code{TheoryUI}.
  & \ref{ch:ui} \tabularnewline
\end{longtable}
\end{center}

\subsection{\texttt{pipeline/theories/} and \texttt{pipeline/examples/}}

These are \emph{hand-written} model dicts and example scripts that predate
the \file{theories/*.theory.py} format. They are not superseded --- the
demo notebook and several tests still import them directly.

\begin{center}
\begin{longtable}{@{}>{\ttfamily\footnotesize\raggedright}p{0.31\textwidth} >{\raggedright}p{0.50\textwidth} >{\centering}p{0.05\textwidth}@{}}
\toprule
\normalfont\textbf{file} & \textbf{role} & \textbf{Ch.} \tabularnewline
\midrule
\endhead
\bottomrule
\endfoot
pipeline/theories/\newline linear\_hawkes\_2pop\_delta.py
  & Hand-written model dict: linear Hawkes, 2 populations, instantaneous
    ($\delta$) synaptic kernel. Equivalent to
    \file{models/hawkes\_linear\_sage.py}.
  & \ref{ch:theory-spec} \tabularnewline
\addlinespace
pipeline/theories/\newline linear\_hawkes\_2pop\_expg.py,
  \newline linear\_hawkes\_2pop\_expg\_gtas.py,
  \newline quad\_hawkes\_2pop\_expg.py,
  \newline quad\_hawkes\_2pop\_expg\_gtas.py
  & The exponential-kernel and GTaS (Gaussian-truncated-and-shifted)
    Hawkes variants, linear and quadratic. Imported by
    \file{pipeline\_demo.ipynb}, several \file{*\_sim\_compare} notebooks, and
    \file{pipeline/tests/test\_theory\_equivalence.py}.
  & \ref{ch:theory-spec} \tabularnewline
\addlinespace
pipeline/examples/\newline run\_linear\_gtas.py
  & A standalone end-to-end example script: compute the $k=2$
    cross-correlator $\langle\delta n_1(0)\,\delta n_2(\tau)\rangle$ for the
    linear-GTaS model and emit a PDF report.
  & \ref{ch:compute-orchestration} \tabularnewline
\addlinespace
pipeline/tests/\newline test\_theory\_equivalence.py
  & Asserts the hand-written \file{pipeline/theories/} dicts and the
    builder-produced theories agree.
  & --- \tabularnewline
\end{longtable}
\end{center}

% =====================================================================
\section{The validators: \texttt{models/}}
% =====================================================================

\file{models/} is the trust chain's other half: completely independent
direct simulators and the estimators that turn their time series into the
same cumulant slices the engine predicts. The two are overlaid in the
\file{*\_sim\_compare} notebooks; if the diagrammatics are correct, the
curves coincide within Monte-Carlo error.

\subsection{SDE and field simulators}

\begin{center}
\begin{longtable}{@{}>{\ttfamily\footnotesize\raggedright}p{0.31\textwidth} >{\raggedright}p{0.50\textwidth} >{\centering}p{0.05\textwidth}@{}}
\toprule
\normalfont\textbf{file} & \textbf{role} & \textbf{Ch.} \tabularnewline
\midrule
\endhead
\bottomrule
\endfoot
models/ou\_langevin\_sim\_numba.py
  & Euler--Maruyama simulators for scalar Langevin SDEs (OU-type theories),
    \code{numba}-JIT compiled. The simulator behind the quickstart's
    \code{sim\_ou\_quartic\_numba}.
  & \ref{ch:simulators} \tabularnewline
\addlinespace
models/spatial\_field\_1d\_sim.py
  & Spectral exponential-Euler (ETD1) simulator for a 1D scalar stochastic
    field on a periodic ring --- the spatial-v1 cross-check.
  & \ref{ch:simulators} \tabularnewline
\addlinespace
models/spatial\_field\_2d\_sim.py
  & 2D spectral ETD1 field simulator on a periodic square --- the $d=2$
    validation oracle.
  & \ref{ch:simulators} \tabularnewline
\addlinespace
models/spatial\_field\_3d\_sim.py
  & 3D spectral ETD1 field simulator on a periodic cube --- the $d=3$
    validation oracle.
  & \ref{ch:simulators} \tabularnewline
\addlinespace
models/spatial\_field\_phi6\_1d\_sim.py
  & ETD1 simulator for the 1D stochastic $\phi^6$ field (Allen--Cahn $+$
    quintic force) --- the cross-check for the quintic theory.
  & \ref{ch:simulators} \tabularnewline
\addlinespace
models/coupled\_rd\_1d\_sim.py
  & $N$-species coupled stochastic reaction--diffusion Langevin simulator on
    a ring --- the coupled-field cross-check.
  & \ref{ch:simulators} \tabularnewline
\addlinespace
models/dendritic\_linear\_sim\_numba.py
  & Soma$+$dendrite compartmental linear-rate simulator (\code{numba} JIT)
    for the dendritic theory.
  & \ref{ch:simulators} \tabularnewline
\end{longtable}
\end{center}

\subsection{Hawkes / point-process simulators}

\begin{center}
\begin{longtable}{@{}>{\ttfamily\footnotesize\raggedright}p{0.31\textwidth} >{\raggedright}p{0.50\textwidth} >{\centering}p{0.05\textwidth}@{}}
\toprule
\normalfont\textbf{file} & \textbf{role} & \textbf{Ch.} \tabularnewline
\midrule
\endhead
\bottomrule
\endfoot
models/hawkes\_sim\_numba.py
  & \code{numba}-JIT Euler-step simulator for the linear Hawkes process
    (fixed-step Poisson-draw realization).
  & \ref{ch:simulators} \tabularnewline
\addlinespace
models/hawkes\_sim\_multipop\_numba.py
  & The linear \emph{multipopulation} Hawkes stepper.
  & \ref{ch:simulators} \tabularnewline
\addlinespace
models/hawkes\_sim\_expg\_numba.py,
  \newline hawkes\_sim\_quad\_expg\_numba.py
  & Exponential-kernel Hawkes simulators, linear and quadratic.
  & \ref{ch:simulators} \tabularnewline
\addlinespace
models/\newline hawkes\_sim\_linear\_expg\_gtas\_numba.py,
  \newline hawkes\_sim\_quad\_expg\_gtas\_numba.py
  & The GTaS-kernel Hawkes simulators, linear and quadratic.
  & \ref{ch:simulators} \tabularnewline
\end{longtable}
\end{center}

\subsection{Estimators and hand-written Hawkes model dicts}

\begin{center}
\begin{longtable}{@{}>{\ttfamily\footnotesize\raggedright}p{0.31\textwidth} >{\raggedright}p{0.50\textwidth} >{\centering}p{0.05\textwidth}@{}}
\toprule
\normalfont\textbf{file} & \textbf{role} & \textbf{Ch.} \tabularnewline
\midrule
\endhead
\bottomrule
\endfoot
models/cumulant\_estimator.py
  & The general $k$-point connected-cumulant estimator
    (\code{estimate\_kpoint\_slices}): bins a trajectory and estimates the
    connected cumulant density versus lag, removing the discrete-spike
    shot-noise diagonal.
  & \ref{ch:simulators} \tabularnewline
\addlinespace
models/cumulant\_direct\_numba.py
  & \code{numba}-JIT direct (non-FFT) 3-point cumulant estimator --- the
    $k=3$ companion to the general estimator.
  & \ref{ch:simulators} \tabularnewline
\addlinespace
models/hawkes\_sage.py,
  \newline hawkes\_linear\_sage.py
  & Hand-written SageMath model specifications for the nonlinear and linear
    Hawkes processes.
  & \ref{ch:theory-spec} \tabularnewline
\addlinespace
models/hawkes\_linear\_expg.py,
  \newline hawkes\_quad\_expg.py,
  \newline hawkes\_linear\_expg\_gtas.py,
  \newline hawkes\_quad\_expg\_gtas.py
  & Hand-written model dicts for the exponential- and GTaS-kernel Hawkes
    families (the symbolic counterparts of the simulators above).
    \file{hawkes\_quad\_expg.py} is the canonical worked example in the
    field-theory-core chapter.
  & \ref{ch:fieldtheory-core},\,\ref{ch:theory-spec} \tabularnewline
\end{longtable}
\end{center}

% =====================================================================
\section{The models: \texttt{theories/}}
% =====================================================================

Each \file{theories/<name>.theory.py} is one physical model: an importable
module exposing exactly \code{build()} (returns the model dict),
\code{DEFAULT\_FUNDAMENTAL} (default numeric parameters), and
\code{METADATA} (run recommendations). \code{dd.load\_theory('<name>')}
imports the file and calls \code{build()}. Rather than list all
forty-plus, the table groups them by the builder they use and the
capability each exercises; the loadable name is the filename without the
\file{.theory.py} suffix.

\begin{center}
\begin{longtable}{@{}>{\raggedright}p{0.30\textwidth} >{\ttfamily\footnotesize\raggedright}p{0.40\textwidth} >{\centering}p{0.16\textwidth}@{}}
\toprule
\textbf{group (builder)} & \normalfont\textbf{representative files} & \textbf{Ch.} \tabularnewline
\midrule
\endhead
\bottomrule
\endfoot
Temporal, single field, the nonlinear baseline
  & ou\_quartic, ou\_quartic\_double\_well, ou\_sextic,
    toy\_quartic\_double\_well
  & \ref{ch:quickstart},\,\ref{ch:theory-files} \tabularnewline
\addlinespace
Temporal, colored / correlated noise
  & ou\_quartic\_colored, ou\_quartic\_two\_dim,
    ou\_quartic\_two\_dim\_corr, ou\_quartic\_two\_dim\_color\_corr
  & \ref{ch:colored-markovian} \tabularnewline
\addlinespace
Temporal Hawkes (point process)
  & linear\_hawkes, quadratic\_hawkes, quadratic\_hawkes\_alpha
  & \ref{ch:theory-spec} \tabularnewline
\addlinespace
Temporal neuron / compartmental tests
  & single\_pop\_dendritic\_linear,
    single\_population\_conductance\_test,
    single\_population\_spike\_reset\_test,
    single\_population\_quad\_exp\_test, multipopulation\_test
  & \ref{ch:theory-spec},\,\ref{ch:mean-field} \tabularnewline
\addlinespace
Spatial, single field, plain vertex
  & allen\_cahn\_1d\_subcritical\_infinite,
    allen\_cahn\_1d\_subcritical\_pbc,
    allen\_cahn\_quintic\_1d\_subcritical\_infinite,
    reaction\_diffusion\_quadratic\_1d
  & \ref{ch:spatial-core},\,\ref{ch:spatial-heatkernel} \tabularnewline
\addlinespace
Spatial, derivative vertex (operator IR)
  & kpz\_1d, kpz\_2d, burgers\_1d, edwards\_wilkinson\_1d,
    reaction\_diffusion\_conserved\_1d
  & \ref{ch:spatial-heatkernel},\,\ref{ch:spatial-coupled} \tabularnewline
\addlinespace
Spatial, coupled / combined fields
  & coupled\_rd\_2species\_1d, reaction\_diffusion\_2d,
    combined\_allencahn\_modelb\_kpz\_\{1,2,3\}d, linear\_field\_2d,
    linear\_diffusion\_test
  & \ref{ch:spatial-coupled} \tabularnewline
\end{longtable}
\end{center}

\begin{note}
\textbf{Builder-class hint.} Files using \code{TemporalTheoryBuilder} are
time-only; \code{SpatialTheoryBuilder} files are PDEs; a handful (e.g.
\file{combined\_allencahn\_modelb\_kpz\_1d.theory.py},
\file{linear\_field\_2d.theory.py}) use the back-compatible
\code{TheoryBuilder} that auto-detects at \code{build()} time. The
\code{*\_test} files are regression fixtures exercising one feature each
(spike-reset resets, conductance synapses, $\delta$- vs exponential
voltage), not headline physics.
\end{note}

% =====================================================================
\section{The front desk and notebooks: \texttt{notebooks/}}
% =====================================================================

\subsection{The convenience layer}

\begin{center}
\begin{longtable}{@{}>{\ttfamily\footnotesize\raggedright}p{0.31\textwidth} >{\raggedright}p{0.50\textwidth} >{\centering}p{0.05\textwidth}@{}}
\toprule
\normalfont\textbf{file} & \textbf{role} & \textbf{Ch.} \tabularnewline
\midrule
\endhead
\bottomrule
\endfoot
notebooks/daedalus.py
  & The 1{,}101-line front desk imported as \code{dd}. The four verbs
    (\code{load\_theory}, \code{Config}, \code{run}, \code{plot\_cumulant})
    plus introspection helpers and the \code{output='moment'} assembly
    layer. Thin orchestration --- it does almost no physics.
  & \ref{ch:quickstart},\,\ref{ch:engine-api},\,\ref{ch:cumulants-moments} \tabularnewline
\end{longtable}
\end{center}

\subsection{The notebook directories}

The notebooks are not modules to read line-by-line; they are the
runnable surface. They are grouped by purpose:

\begin{center}
\begin{longtable}{@{}>{\ttfamily\footnotesize\raggedright}p{0.31\textwidth} >{\raggedright}p{0.50\textwidth} >{\centering}p{0.05\textwidth}@{}}
\toprule
\normalfont\textbf{directory} & \textbf{role} & \textbf{Ch.} \tabularnewline
\midrule
\endhead
\bottomrule
\endfoot
notebooks/examples/
  & The nine bonafide example notebooks, one per pipeline capability,
    each with a model-info $\to$ theory $\to$ sim structure
    (\file{temporal\_ou\_quartic\_white.ipynb} is the quickstart's).
  & \ref{ch:quickstart},\,\ref{ch:engine-api} \tabularnewline
\addlinespace
notebooks/templates/
  & The four group templates (\{temporal,spatial\}$\times$\{single,multi\})
    plus \code{*\_sim\_compare} variants that every demo is generated from.
  & \ref{ch:engine-api} \tabularnewline
\addlinespace
notebooks/temporal/
  & The temporal \code{*\_sim\_compare} validation notebooks (OU quartic,
    Hawkes, conductance, spike-reset, $\dots$) --- theory overlaid on
    simulation.
  & \ref{ch:simulators} \tabularnewline
\addlinespace
notebooks/spatial/
  & The spatial \code{*\_sim\_compare} validation notebooks (Allen--Cahn,
    KPZ, Burgers, Model~B, coupled RD, $d=1,2,3$).
  & \ref{ch:simulators} \tabularnewline
\addlinespace
notebooks/legacy/
  & The pre-pipeline demo notebooks (enumeration, field-theory-Sage,
    Hawkes time-domain), kept for provenance.
  & --- \tabularnewline
\addlinespace
notebooks/\newline theory\_builder.ipynb,
  \newline theory\_runner.ipynb
  & The single-cell launchers for the \code{TheoryUI} form and for running a
    saved \file{.theory.py}.
  & \ref{ch:ui} \tabularnewline
\end{longtable}
\end{center}

% =====================================================================
\section{Supporting trees: tests, scripts, legacy, docs}
% =====================================================================

The remaining top-level directories are not part of the runtime data flow
but are essential to working in the repository.

\begin{center}
\begin{longtable}{@{}>{\ttfamily\footnotesize\raggedright}p{0.27\textwidth} >{\raggedright}p{0.54\textwidth} >{\centering}p{0.05\textwidth}@{}}
\toprule
\normalfont\textbf{directory} & \textbf{role} & \textbf{Ch.} \tabularnewline
\midrule
\endhead
\bottomrule
\endfoot
tests/
  & The pytest suite (61 \code{test\_*.py} files). One file per subsystem
    plus the oracle cross-checks: e.g. \file{test\_full\_integrator.py},
    \file{test\_spatial\_reduce.py}, \file{test\_momentum\_routing.py},
    \file{test\_symmetry.py}, \file{test\_grouped\_vs\_perdiag.py},
    \file{test\_all\_k\_boltzmann.py}, \file{test\_kpz\_burgers\_sim.py},
    \file{test\_generic\_evaluator.py} (the spatial oracle),
    \file{test\_fork\_safety.py}.
  & all \tabularnewline
\addlinespace
scripts/
  & Standalone runnable scripts; currently
    \file{scripts/profile\_pipeline\_k3.py}, a $k=3$ profiling harness.
  & --- \tabularnewline
\addlinespace
legacy/
  & The pre-pipeline \code{sympy}-era codebase (the original Enumeration
    Code, Field Theory Framework, Theory Builder, and \code{*\_sympy.py}
    modules), isolated and kept for provenance. \emph{Not} on any live
    path.
  & --- \tabularnewline
\addlinespace
saved\_theories/
  & The on-disk \code{.sobj} cache (per-theory \file{<slug>/} directories of
    pickled propagators, expanded actions, and unique-diagram sets, plus a
    \code{manifest.json}). Gitignored, regenerable.
  & \ref{ch:caching} \tabularnewline
\addlinespace
docs/
  & The prose documentation, design records, and this manual
    (\file{docs/manual/}). \file{docs/spatial\_pipeline.md} is the
    authoritative spatial-state reference; \file{docs/CHANGELOG.md} is the
    chronological fix log.
  & --- \tabularnewline
\addlinespace
scratch/, \newline pipeline\_outputs/
  & Throwaway probe scripts/logs and generated run artefacts respectively.
    Not part of the codebase; safe to delete. (\file{scratch/} also holds
    the tropical-Monte-Carlo research PDFs cited by the efficiency roadmap.)
  & --- \tabularnewline
\end{longtable}
\end{center}

% =====================================================================
\section{The shortest path to the file you want}
% =====================================================================

To close, a reverse index --- start from what you are trying to do, and
this tells you which file (and chapter) to open:

\begin{description}[style=nextline,leftmargin=1.4em]
  \item[``I want to add a new model.''] Write a
        \file{theories/<name>.theory.py} using \file{pipeline/theory.py}'s
        builder (Chapter~\ref{ch:theory-spec},
        Chapter~\ref{ch:theory-files}), or point-and-click it with
        \file{pipeline/ui/main.py} (Chapter~\ref{ch:ui}).
  \item[``I want to run a model and plot it.''] \file{notebooks/daedalus.py}
        --- the \code{load\_theory $\to$ Config $\to$ run $\to$ plot} spine
        (Chapter~\ref{ch:quickstart}, Chapter~\ref{ch:engine-api}).
  \item[``The answer is wrong and I suspect the diagrams.''] The enumeration
        is \file{msrjd/enumeration/loop\_diagram\_enumeration.py}
        (Chapter~\ref{ch:enumeration}); the typing/symmetry is
        \file{msrjd/diagrams/} (Chapter~\ref{ch:type-assignment}).
  \item[``The answer is wrong and I suspect the loop integral.''] Temporal:
        \file{msrjd/integration/time\_domain/final\_integral.py}
        (Chapter~\ref{ch:phasej-temporal}). Spatial:
        \file{msrjd/integration/spatial/full\_integrator.py}
        (Chapter~\ref{ch:spatial-core},
        Chapter~\ref{ch:spatial-heatkernel}).
  \item[``The saddle looks wrong.''] \file{pipeline/\_mean\_field\_dae.py}
        (Chapter~\ref{ch:mean-field}).
  \item[``It is too slow / it recomputes every run.''] The caches:
        \file{pipeline/\_expand\_cache.py} and \file{pipeline/\_diagrams.py}
        (Chapter~\ref{ch:caching}).
  \item[``I want to check it against a simulation.''] The matching
        \file{models/*\_sim*.py} simulator $+$
        \file{models/cumulant\_estimator.py} (Chapter~\ref{ch:simulators}),
        driven from a \file{notebooks/\{temporal,spatial\}/*\_sim\_compare.ipynb}.
  \item[``Where does it all start?''] \code{compute\_cumulants} in
        \file{pipeline/compute.py} (Chapter~\ref{ch:compute-orchestration})
        --- the one function every path runs through.
\end{description}

% =====================================================================
%  APPENDIX D : Known issues, caveats, and open questions
% =====================================================================
\chapter{Known Issues, Caveats, and Open Questions}
\label{app:known-issues}

No production scientific code is free of rough edges, and \Daedalus{} is honest
about its.  This appendix collects the limitations a \emph{user} should know
before trusting a number, and points to the developer-facing log of code-hygiene
items.

\begin{note}
A living, more granular version of this list --- including line-level code
inconsistencies found while writing this manual and the per-chapter
manual-vs-code audit results --- is kept in
\file{docs/manual/AUDIT\_FINDINGS.md}.  This appendix is the curated,
reader-facing summary.
\end{note}

% ---------------------------------------------------------------------
\section{Numerical and physical limitations}
% ---------------------------------------------------------------------

\begin{description}[leftmargin=0pt,style=nextline]

  \item[Colored-noise two-loop runs are slow]
        A finite-$\tau_c$ (colored) theory is handled by Markovian embedding
        (Chapter~\ref{ch:colored-markovian}), which \emph{doubles} the field
        content.  At two loops (\code{max\_ell=2}) this produces $\sim$66 diagrams
        whose per-$\tau$ evaluation falls back to nested numerical quadrature, so
        a single correlator can take many minutes per lag.  This is a cost of the
        embedding, \emph{not} a regression.  For a fast nonlinear two-loop
        baseline use the \emph{white}-noise \code{ou\_quartic} theory, which is a
        single field and runs in well under a second
        (Chapter~\ref{ch:quickstart}).

  \item[$\mu=1/\tau_c$ exact degeneracy]
        The v1 colored$\to$Markovian preprocessor needs the physical relaxation
        rate $\mu$ and the noise rate $1/\tau_c$ to be \emph{distinct} --- the
        embedded propagator then has simple poles at each.  At exact equality the
        pole is a double pole, which the current code does not special-case.
        Keep $\mu\tau_c\neq 1$.

  \item[Large-$|\mu|$ OU-quartic blow-up at $\ell\ge1$]
        For the OU-quartic family the Phase-J loop integrals grow rapidly in
        magnitude as $|\mu|$ increases at one loop and beyond; the perturbative
        series is only well-behaved in the mildly-perturbative window
        ($g_{\mathrm{eff}}\approx\varepsilon D/\mu^2$ small).  See
        \file{docs/ou\_quartic\_large\_mu\_phase\_j\_audit.md}.

  \item[$m\ge3$ chain-simplex precision]
        The temporal mode-sum can lose precision when three or more propagator
        poles are very closely spaced (the ``close-paired pole'' regime),
        affecting e.g.\ spike-reset correlators at $k\ge2,\ \ell\ge1$.  Documented
        with a regression fixture in
        \file{docs/m\_ge3\_precision\_bug\_audit.md}.

  \item[Spatial $d\ge2$ with derivative vertices is still numerical]
        The analytic heat-kernel inverse Fourier transform
        (Chapter~\ref{ch:spatial-heatkernel}) is exact for plain vertices in any
        dimension and for $d=1$ derivative vertices (Model~B, KPZ, Burgers).  The
        $d\ge2$ derivative-vertex transverse handling has not yet been done
        analytically, so those cases use the numerical Fourier path.

  \item[Spatial coupled fields with unequal diffusivities: tree only]
        For coupled multi-field spatial models with \emph{equal} diffusivity the
        full tree+loop machinery runs; for \emph{unequal} diffusivities only the
        tree correlator (via Dyson dressing, Chapter~\ref{ch:spatial-coupled}) is
        wired --- loop corrections are gated off.

  \item[Spatial $k\ge3$ loop path]
        The public \code{compute\_cumulants} loop path is $k=2$-gated.  The
        spatial integrator core supports general $k$, but $k\ge3$ loop records
        still need the reduced chamber quadrature before they go through the
        public API.

\end{description}

% ---------------------------------------------------------------------
\section{Code-hygiene items and open questions}
% ---------------------------------------------------------------------

While documenting the codebase we noted a number of small inconsistencies ---
dead imports, a stale module docstring, two docstring-only stub modules with no
importers, a couple of no-op code branches, and a few advertised-but-unwired
\code{Config} fields.  None affects the correctness of a computed cumulant; they
are tidy-up candidates.  The full itemised list (with file:line references and a
recommended action for each) lives in \file{docs/manual/AUDIT\_FINDINGS.md},
Part~A.

% ---------------------------------------------------------------------
\section{How this manual was checked}
% ---------------------------------------------------------------------

Each chapter was drafted from a deep read of its subsystem and then
\emph{independently audited} against the source code by a separate pass that
re-grepped every named identifier, verified file:line citations, and checked
behavioural claims.  Corrections from that pass (for example the Fourier-sign and
template-lambda-count fixes in Chapter~\ref{ch:theory-spec}) were folded back
into the chapters.  Any audit findings still open at the time of writing are
tracked in \file{docs/manual/AUDIT\_FINDINGS.md}; treat that file as the manual's
errata sheet.


\end{document}
